\documentclass[11pt,a5paper]{book}
\usepackage[utf8]{inputenc}
\usepackage[T1]{fontenc}
\usepackage[brazil]{babel}
\usepackage{geometry}
\geometry{margin=1.8cm}
\usepackage{graphicx}
\usepackage{hyperref}
\usepackage{microtype}
\usepackage{setspace}
\onehalfspacing
\title{O Livro do Porquê\\\large A nova ciência de causa e efeito}
\author{Judea Pearl e Dana Mackenzie}
\begin{document}
\maketitle
\tableofcontents
\mainmatter
\begin{titlepage}\centering
\includegraphics[width=0.95\textwidth]{../capa_br.png}
\end{titlepage}
\chapter{Prefácio}
Há quase duas décadas, quando escrevi o prefácio de meu livro \textit{Causality} (2000), fiz uma observação bastante ousada, que amigos me aconselharam a suavizar. “A causalidade sofreu uma grande transformação”, escrevi, “de um conceito envolto em mistério a um objeto matemático com semântica bem definida e lógica bem fundamentada. Paradoxos e controvérsias foram resolvidos, conceitos escorregadios foram esclarecidos, e problemas práticos que dependem de informação causal e que por muito tempo foram tidos como metafísicos ou intratáveis agora podem ser solucionados com matemática elementar. Em suma, a causalidade foi matematizada.”

Relendo essa passagem hoje, sinto que fui um tanto míope. O que descrevi como uma “transformação” revelou-se uma “revolução” que mudou o pensamento em muitas das ciências. Muitos a chamam agora de “Revolução Causal”[^1], e a empolgação que ela tem gerado nos círculos de pesquisa está se espalhando para a educação e as aplicações. Creio que chegou a hora de compartilhá-la com um público mais amplo.

Este livro busca cumprir uma missão de três frentes: primeiro, apresentar a você, em linguagem não matemática, o conteúdo intelectual da Revolução Causal e o modo como ela está afetando nossas vidas, bem como nosso futuro; segundo, partilhar algumas das jornadas heroicas — bem-sucedidas e fracassadas — que cientistas empreenderam ao se depararem com perguntas críticas de causa e efeito. Por fim, devolvendo a Revolução Causal ao seu ventre na inteligência artificial, pretendo descrever como se podem construir robôs que aprendam a se comunicar em nossa língua materna — a linguagem de causa e efeito. Essa nova geração de robôs deverá nos explicar por que as coisas aconteceram, por que reagiram da maneira que reagiram, e por que a natureza opera de um modo e não de outro. Mais ambiciosamente, deverão também nos ensinar sobre nós mesmos: por que nossa mente clica do jeito que clica e o que significa pensar racionalmente sobre causa e efeito, crédito e arrependimento, intenção e responsabilidade.

Quando escrevo equações, tenho uma ideia muito clara de quem são meus leitores. Não é o caso quando escrevo para o público geral — uma aventura inteiramente nova para mim. Estranho, mas essa experiência tem sido uma das viagens educativas mais recompensadoras de minha vida. A necessidade de moldar ideias em sua linguagem, de imaginar seu repertório, suas perguntas e suas reações, fez mais por afiar minha compreensão da causalidade do que todas as equações que escrevi antes deste livro. Por isso serei para sempre grato a você. Espero que esteja tão entusiasmado quanto eu para ver os resultados.

Judea Pearl Los Angeles, outubro de 2017

\chapter{Introdução: Mente sobre dados}
\begin{quote}Toda ciência que prosperou prosperou sobre seus próprios símbolos.   — AUGUSTUS DE MORGAN (1864)\end{quote}
Este livro conta a história de uma ciência que mudou a maneira como distinguimos fatos de ficção e que, no entanto, permaneceu sob o radar do grande público. As consequências da nova ciência já estão impactando facetas cruciais de nossas vidas e têm o potencial de afetar ainda mais — do desenvolvimento de novos medicamentos ao controle de políticas econômicas, da educação e da robótica ao controle de armas e ao aquecimento global. De modo notável, apesar da diversidade e da aparente incomensurabilidade dessas áreas problemáticas, a nova ciência as abrange todas sob um arcabouço unificado que praticamente não existia há duas décadas.

A nova ciência não tem um nome pomposo: chamo-a simplesmente de “inferência causal”[^1], como fazem muitos de meus colegas. Tampouco é particularmente de alta tecnologia. A tecnologia ideal que a inferência causal se esforça por emular reside em nossas próprias mentes. Há algumas dezenas de milhares de anos, os seres humanos começaram a perceber que certas coisas causam outras coisas e que mexer nas primeiras pode mudar as segundas. Nenhuma outra espécie compreende isso — certamente não na medida em que nós o fazemos. Dessa descoberta vieram as sociedades organizadas, depois vilas e cidades e, por fim, a civilização baseada em ciência e tecnologia de que desfrutamos hoje. Tudo porque fizemos uma pergunta simples: Por quê?

A inferência causal trata de levar essa pergunta a sério. Ela postula que o cérebro humano é a ferramenta mais avançada já concebida para gerir causas e efeitos. Nossos cérebros armazenam uma quantidade incrível de conhecimento causal que, suplementado por dados, poderíamos aproveitar para responder a algumas das questões mais urgentes de nosso tempo. Mais ambiciosamente, uma vez que compreendamos de fato a lógica por trás do pensamento causal, poderíamos emulá-la em computadores modernos e criar um “cientista artificial”. Esse robô inteligente descobriria fenômenos ainda desconhecidos, encontraria explicações para dilemas científicos pendentes, projetaria novos experimentos e extrairia continuamente mais conhecimento causal do ambiente.

Mas, antes de nos aventurarmos a especular sobre tais desenvolvimentos futuristas, é importante compreender as conquistas que a inferência causal acumulou até agora. Exploraremos o modo como ela transformou o pensamento de cientistas em quase todas as disciplinas informadas por dados e como está prestes a mudar nossas vidas.

A nova ciência aborda perguntas aparentemente diretas como estas:

\begin{itemize}
\item Quão eficaz é determinado tratamento na prevenção de uma doença?
\item A nova lei tributária fez nossas vendas subirem, ou foi a campanha publicitária?
\item Qual é o custo de saúde atribuível à obesidade?
\item Registros de contratação podem provar que um empregador é culpado de uma política de discriminação sexual?
\item Estou prestes a pedir demissão. Devo?
\end{itemize}
Essas perguntas têm em comum a preocupação com relações de causa e efeito, reconhecíveis por palavras como “prevenção”, “causar”, “atribuível a”, “política” e “devo”. Tais palavras são comuns na linguagem cotidiana, e nossa sociedade exige constantemente respostas a essas perguntas. No entanto, até muito recentemente, a ciência não nos deu meios sequer de articulá-las, quanto menos de respondê-las.

De longe, a contribuição mais importante da inferência causal à humanidade foi transformar essa negligência científica em coisa do passado. A nova ciência gerou uma linguagem matemática simples para articular as relações causais que conhecemos, bem como aquelas que desejamos descobrir. A capacidade de expressar essa informação em forma matemática desencadeou uma riqueza de métodos poderosos e principados para combinar nosso conhecimento com dados e responder a perguntas causais como as cinco acima.

Tive a sorte de participar desse desenvolvimento científico ao longo do último quarto de século. Vi seu progresso tomar forma em cubículos de estudantes e laboratórios de pesquisa, e ouvi seus avanços ressoar em sóbrias conferências científicas, longe dos holofotes da atenção pública. Agora, ao entrarmos na era da inteligência artificial (IA) forte, e enquanto muitos apregoam as possibilidades infinitas do \textit{big data}[^2] e do aprendizado profundo, acho oportuno e empolgante apresentar ao leitor alguns dos caminhos mais aventureiros que a nova ciência está trilhando, como ela impacta a ciência de dados e as muitas maneiras pelas quais mudará nossas vidas no século XXI.

Quando me ouve descrever essas conquistas como uma “nova ciência”, você pode ficar cético. Pode até perguntar: por que isso não foi feito há muito tempo? Digamos, quando Virgílio proclamou pela primeira vez: “Feliz aquele que pôde conhecer as causas das coisas” (29 a.C.). Ou quando os fundadores da estatística moderna, Francis Galton e Karl Pearson, descobriram que dados populacionais podem lançar luz sobre questões científicas. Há uma longa história por trás de sua infeliz falha em abraçar a causalidade nesse momento crucial, que as seções históricas deste livro relatarão. Mas o obstáculo mais sério, em minha opinião, tem sido a lacuna fundamental entre o vocabulário em que formulamos perguntas causais e o vocabulário tradicional em que comunicamos teorias científicas.

Para apreciar a profundidade dessa lacuna, imagine as dificuldades que um cientista enfrentaria ao tentar expressar algumas relações causais óbvias — digamos, que a leitura do barômetro \textit{B} acompanha a pressão atmosférica \textit{P}. Podemos facilmente escrever essa relação em uma equação como \textit{B} = \textit{kP}, onde \textit{k} é alguma constante de proporcionalidade. As regras da álgebra agora nos permitem reescrever essa mesma equação em uma variedade selvagem de formas — por exemplo, \textit{P} = \textit{B}/\textit{k}, \textit{k} = \textit{B}/\textit{P}, ou \textit{B} − \textit{kP} = 0. Todas significam a mesma coisa: que se conhecermos quaisquer duas das três quantidades, a terceira fica determinada. Nenhuma das letras \textit{k}, \textit{B} ou \textit{P} é, de qualquer modo matemático, privilegiada em relação às outras. Como, então, podemos expressar nossa firme convicção de que é a pressão que faz o barômetro mudar, e não o contrário? E se não conseguimos expressar sequer isso, como podemos esperar expressar as muitas outras convicções causais que não têm fórmulas matemáticas, como a de que o canto do galo não faz o sol nascer?

Meus professores da faculdade não conseguiam fazê-lo e nunca se queixaram. Apostaria que nenhum dos seus também o fez. Agora compreendemos o porquê: nunca lhes foi mostrada uma linguagem matemática de causas; tampouco lhes foram mostrados seus benefícios. É, de fato, uma acusação à ciência o fato de ela ter negligenciado o desenvolvimento de tal linguagem por tantas gerações. Todo mundo sabe que acionar um interruptor faz uma luz acender ou apagar e que uma tarde de verão quente e abafada faz as vendas subirem na sorveteria local. Por que, então, os cientistas não capturaram fatos tão óbvios em fórmulas, como fizeram com as leis básicas da óptica, da mecânica ou da geometria? Por que permitiram que esses fatos definhassem na mera intuição, privados das ferramentas matemáticas que permitiram a outros ramos da ciência florescer e amadurecer?

Parte da resposta é que as ferramentas científicas se desenvolvem para atender a necessidades científicas. Precisamente porque somos tão bons em lidar com perguntas sobre interruptores, sorvete e barômetros, nossa necessidade de uma maquinaria matemática especial para tratá-los não era óbvia. Mas, à medida que a curiosidade científica aumentou e começamos a formular perguntas causais em situações complexas de direito, negócios, medicina e formulação de políticas, descobrimos que nos faltavam as ferramentas e os princípios que uma ciência madura deveria fornecer.

Despertares tardios desse tipo não são incomuns na ciência. Por exemplo, até cerca de quatrocentos anos atrás, as pessoas estavam bastante satisfeitas com sua capacidade natural de gerir as incertezas da vida cotidiana, de atravessar uma rua a arriscar uma briga de socos. Somente depois que jogadores inventaram jogos de azar intricados, por vezes cuidadosamente desenhados para nos induzir a más escolhas, matemáticos como Blaise Pascal (1654), Pierre de Fermat (1654) e Christiaan Huygens (1657) acharam necessário desenvolver o que hoje chamamos de teoria das probabilidades. Do mesmo modo, somente quando organizações de seguros exigiram estimativas precisas de anuidades vitalícias, matemáticos como Edmond Halley (1693) e Abraham de Moivre (1725) começaram a examinar tábuas de mortalidade para calcular expectativas de vida. De forma semelhante, as demandas dos astrônomos por previsões precisas do movimento celeste levaram Jacob Bernoulli, Pierre-Simon Laplace e Carl Friedrich Gauss a desenvolver uma teoria dos erros para nos ajudar a extrair sinais do ruído. Esses métodos foram todos predecessores da estatística de hoje.

Ironicamente, a necessidade de uma teoria da causalidade começou a aflorar na mesma época em que a estatística nascia. De fato, a estatística moderna eclodiu das perguntas causais que Galton e Pearson fizeram sobre a hereditariedade e de suas tentativas engenhosas de respondê-las usando dados intergeracionais. Infelizmente, eles fracassaram nessa empreitada e, em vez de fazer uma pausa para perguntar por quê, declararam essas questões fora dos limites e se dedicaram a desenvolver um empreendimento próspero e livre de causalidade chamado estatística.

Esse foi um momento crítico na história da ciência. A oportunidade de munir as perguntas causais com uma linguagem própria esteve muito perto de se concretizar, mas foi desperdiçada. Nos anos seguintes, essas perguntas foram declaradas não científicas e foram para a clandestinidade. Apesar dos esforços heroicos do geneticista Sewall Wright (1889–1988), o vocabulário causal ficou virtualmente proibido por mais de meio século. E quando se proíbe a fala, proíbe-se o pensamento e sufocam-se princípios, métodos e ferramentas.

Os leitores não precisam ser cientistas para testemunhar essa proibição. Em Estatística 101, todo estudante aprende a entoar: “Correlação não é causalidade”. Com razão! O canto do galo está altamente correlacionado com o nascer do sol; contudo, não o causa. Infelizmente, a estatística fetichizou essa observação de bom senso. Ela nos diz que correlação não é causalidade, mas não nos diz o que a causalidade \textit{é}. Em vão se buscará no índice de um livro-texto de estatística uma entrada para “causa”. Não se permite aos estudantes dizer que \textit{X} é a causa de \textit{Y} — apenas que \textit{X} e \textit{Y} estão “relacionados” ou “associados”.

Por causa dessa proibição, as ferramentas matemáticas para gerir perguntas causais foram tidas como desnecessárias, e a estatística se concentrou exclusivamente em como resumir dados, não em como interpretá-los. Uma brilhante exceção foi a análise de caminhos, inventada pelo geneticista Sewall Wright nos anos 1920 e ancestral direta dos métodos que veremos neste livro. Contudo, a análise de caminhos foi gravemente subestimada na estatística e em suas comunidades satélites, e definhou por décadas em seu status embrionário. O que deveria ter sido o primeiro passo em direção à inferência causal permaneceu o único passo até os anos 1980. O restante da estatística, incluindo as muitas disciplinas que a ela recorriam em busca de orientação, permaneceu na era da Proibição, acreditando falsamente que as respostas a todas as perguntas científicas residem nos dados, a serem desveladas por meio de truques engenhosos de mineração de dados.

Muito dessa história centrada em dados ainda nos assombra hoje. Vivemos em uma era que presume que o \textit{big data} é a solução para todos os nossos problemas. Cursos de “ciência de dados” proliferam em nossas universidades, e os empregos de “cientistas de dados” são lucrativos nas empresas que participam da “economia de dados”. Mas espero, com este livro, convencê-lo de que os dados são profundamente burros. Os dados podem lhe dizer que as pessoas que tomaram um remédio se recuperaram mais rápido do que as que não o tomaram, mas não podem lhe dizer \textit{por quê}. Talvez quem tomou o remédio o tenha feito porque podia pagá-lo e teria se recuperado igualmente rápido sem ele.

Uma e outra vez, na ciência e nos negócios, vemos situações em que meros dados não bastam. A maioria dos entusiastas do \textit{big data}, embora um tanto ciente dessas limitações, continua a caçada por uma inteligência centrada em dados, como se ainda estivéssemos na era da Proibição.

Como mencionei antes, as coisas mudaram dramaticamente nas últimas três décadas. Hoje em dia, graças a modelos causais cuidadosamente elaborados, cientistas contemporâneos podem abordar problemas que outrora teriam sido considerados insolúveis ou mesmo além dos limites da investigação científica. Por exemplo, há apenas cem anos, a pergunta de se o fumo de cigarros causa um risco à saúde teria sido considerada não científica. A mera menção das palavras “causa” ou “efeito” criaria uma tempestade de objeções em qualquer periódico estatístico respeitável.

Mesmo há duas décadas, fazer a um estatístico uma pergunta como “Foi a aspirina que fez minha dor de cabeça passar?” teria sido como perguntar se ele acreditava em vodu. Para citar um estimado colega meu, seria “mais um assunto de conversa de coquetel do que uma investigação científica”. Mas hoje, epidemiologistas, cientistas sociais, cientistas da computação e ao menos alguns economistas e estatísticos esclarecidos formulam tais perguntas rotineiramente e as respondem com precisão matemática. Para mim, essa mudança é nada menos que uma revolução. Ouso chamá-la de Revolução Causal[^3] — um abalo científico que abraça, em vez de negar, nosso dom cognitivo inato de compreender causa e efeito.

A Revolução Causal não aconteceu no vácuo; tem por trás um segredo matemático que pode ser melhor descrito como um cálculo da causalidade, que responde a alguns dos problemas mais difíceis já formulados sobre relações de causa e efeito. Estou entusiasmado em revelar esse cálculo — não só porque a história turbulenta de seu desenvolvimento é intrigante, mas ainda mais porque espero que seu pleno potencial seja um dia desenvolvido além do que posso imaginar… talvez até por um leitor deste livro.

O cálculo da causalidade consiste em duas linguagens: diagramas causais, para expressar o que sabemos, e uma linguagem simbólica, que se assemelha à álgebra, para expressar o que queremos saber. Os diagramas causais são simplesmente figuras de pontos e setas que resumem nosso conhecimento científico existente. Os pontos representam quantidades de interesse, chamadas “variáveis”, e as setas representam relações causais conhecidas ou suspeitas entre essas variáveis — a saber, qual variável “escuta” quais outras. Esses diagramas são extremamente fáceis de desenhar, compreender e usar, e o leitor encontrará dezenas deles nas páginas deste livro. Se você consegue se orientar com um mapa de ruas de mão única, então pode compreender diagramas causais e resolver o tipo de perguntas formuladas no início desta introdução.

Embora os diagramas causais sejam minha ferramenta de escolha neste livro, como nos últimos trinta e cinco anos de minha pesquisa, não são o único tipo de modelo causal possível. Alguns cientistas (por exemplo, econometristas) preferem trabalhar com equações matemáticas; outros (por exemplo, estatísticos \textit{hard-core}) preferem uma lista de suposições que ostensivamente resume a estrutura do diagrama. Independentemente da linguagem, o modelo deve descrever, ainda que qualitativamente, o processo que gera os dados — em outras palavras, as forças de causa e efeito que operam no ambiente e moldam os dados gerados.

Lado a lado com essa “linguagem do conhecimento” diagramática, temos também uma “linguagem de consultas” simbólica para expressar as perguntas a que queremos respostas. Por exemplo, se estamos interessados no efeito de um fármaco (\textit{D}) sobre a expectativa de vida (\textit{L}), então nossa consulta pode ser escrita simbolicamente como: \textit{P}(\textit{L} | \textit{do}(\textit{D})). Em outras palavras, qual é a probabilidade (\textit{P}) de que um paciente típico sobreviva \textit{L} anos se for levado a tomar o fármaco? Essa pergunta descreve o que epidemiologistas chamariam de uma intervenção[^4] ou um tratamento e corresponde ao que medimos em um ensaio clínico. Em muitos casos, podemos também querer comparar \textit{P}(\textit{L} | \textit{do}(\textit{D})) com \textit{P}(\textit{L} | \textit{do}(não-\textit{D})); esta última descreve pacientes a quem se negou o tratamento, também chamados de pacientes “controle”. O operador \textit{do}[^5] significa que estamos lidando com uma intervenção, e não com uma observação passiva; a estatística clássica não tem nada remotamente semelhante a esse operador.

Devemos invocar um operador de intervenção \textit{do}(\textit{D}) para garantir que a mudança observada na expectativa de vida \textit{L} se deve ao fármaco em si e não está confundida com outros fatores que tendem a encurtar ou alongar a vida. Se, em vez de intervir, deixarmos o próprio paciente decidir se toma o fármaco, esses outros fatores podem influenciar sua decisão, e as diferenças de expectativa de vida entre tomar e não tomar o fármaco já não se deveriam unicamente ao fármaco. Por exemplo, suponha que apenas aqueles que estavam terminalmente doentes tomaram o fármaco. Tais pessoas certamente difeririam daquelas que não o tomaram, e uma comparação dos dois grupos refletiria diferenças na gravidade de sua doença, e não o efeito do fármaco. Em contraste, forçar os pacientes a tomar ou a se abster do fármaco, independentemente de precondições, eliminaria diferenças preexistentes e forneceria uma comparação válida.

Matematicamente, escrevemos a frequência observada da expectativa de vida \textit{L} entre pacientes que tomam voluntariamente o fármaco como \textit{P}(\textit{L} | \textit{D}), que é a probabilidade condicional padrão usada nos livros-texto de estatística. Essa expressão representa a probabilidade (\textit{P}) da expectativa de vida \textit{L} condicional a se ver o paciente tomar o fármaco \textit{D}. Note que \textit{P}(\textit{L} | \textit{D}) pode ser totalmente diferente de \textit{P}(\textit{L} | \textit{do}(\textit{D})). Essa diferença entre \textit{ver} e \textit{fazer} é fundamental e explica por que não consideramos a queda do barômetro uma causa da tempestade que se aproxima. Ver o barômetro cair aumenta a probabilidade da tempestade, enquanto forçá-lo a cair não afeta essa probabilidade.

Essa confusão entre ver e fazer resultou em uma fonte de paradoxos, alguns dos quais entretaremos neste livro. Um mundo desprovido de \textit{P}(\textit{L} | \textit{do}(\textit{D})) e governado unicamente por \textit{P}(\textit{L} | \textit{D}) seria de fato um mundo estranho. Por exemplo, pacientes evitariam ir ao médico para reduzir a probabilidade de estarem gravemente doentes; cidades dispensariam seus bombeiros para reduzir a incidência de incêndios; médicos recomendariam um fármaco a pacientes homens e mulheres, mas não a pacientes de gênero não divulgado; e assim por diante. É difícil acreditar que, há menos de três décadas, a ciência de fato operava em tal mundo: o operador \textit{do} não existia.

Uma das conquistas culminantes da Revolução Causal tem sido explicar como prever os efeitos de uma intervenção sem de fato a encenar. Isso nunca teria sido possível se não tivéssemos, em primeiro lugar, definido o operador \textit{do} de modo a podermos fazer a pergunta certa e, em segundo lugar, concebido uma maneira de emulá-lo por meios não invasivos.

Quando a pergunta científica de interesse envolve pensamento retrospectivo, recorremos a outro tipo de expressão exclusiva do raciocínio causal, chamada de \textit{contrafactual}[^6]. Por exemplo, suponha que Joe tomou o fármaco \textit{D} e morreu um mês depois; nossa pergunta de interesse é se o fármaco pode ter causado sua morte. Para responder a essa pergunta, precisamos imaginar um cenário em que Joe estava prestes a tomar o fármaco, mas mudou de ideia. Ele teria vivido? Mais uma vez, a estatística clássica apenas resume dados, de modo que não fornece sequer uma linguagem para fazer essa pergunta. A inferência causal fornece uma notação e, mais importante, oferece uma solução. Assim como na previsão do efeito de intervenções (mencionada acima), em muitos casos podemos emular o pensamento retrospectivo humano com um algoritmo que toma o que sabemos sobre o mundo observado e produz uma resposta sobre o mundo contrafactual. Essa “algoritmização dos contrafactuais” é outra joia desenterrada pela Revolução Causal.

O raciocínio contrafactual, que lida com “e se…”, pode parecer a alguns leitores não científico. De fato, a observação empírica nunca pode confirmar nem refutar as respostas a tais perguntas. No entanto, nossas mentes fazem o tempo todo juízos muito confiáveis e reproduzíveis sobre o que poderia ser ou poderia ter sido. Todos compreendemos, por exemplo, que se o galo tivesse ficado em silêncio esta manhã, o sol teria nascido igualmente. Esse consenso decorre do fato de que os contrafactuais não são produtos do capricho, mas refletem a própria estrutura de nosso modelo de mundo. Duas pessoas que compartilham o mesmo modelo causal também compartilharão todos os juízos contrafactuais.

Os contrafactuais são os blocos de construção do comportamento moral, bem como do pensamento científico. A capacidade de refletir sobre as próprias ações passadas e de imaginar cenários alternativos é a base do livre-arbítrio e da responsabilidade social. A algoritmização dos contrafactuais convida as máquinas pensantes a se beneficiar dessa capacidade e a participar desse modo (até agora) unicamente humano de pensar sobre o mundo.

Minha menção a máquinas pensantes no último parágrafo é intencional. Cheguei a este tema como cientista da computação que trabalha na área de inteligência artificial, o que implica dois pontos de partida diferentes da maioria de meus colegas na arena da inferência causal. Primeiro, no mundo da IA, você de fato não compreende um tópico até ser capaz de ensiná-lo a um robô mecânico. É por isso que você me encontrará enfatizando e reenfatizando notação, linguagem, vocabulário e gramática. Por exemplo, eu me obsedo com a questão de se podemos expressar determinada afirmação em uma dada linguagem e se uma afirmação se segue de outras. É impressionante o quanto se pode aprender apenas seguindo a gramática dos enunciados científicos. Minha ênfase na linguagem também vem de uma convicção profunda de que a linguagem molda nossos pensamentos. Você não pode responder a uma pergunta que não consegue fazer, e não pode fazer uma pergunta para a qual não tem palavras. Como estudante de filosofia e de ciência da computação, minha atração pela inferência causal foi em grande medida desencadeada pela empolgação de ver uma linguagem científica órfã passar do nascimento à maturidade.

Meu histórico em aprendizado de máquina me deu ainda outro incentivo para estudar a causalidade. No final dos anos 1980, percebi que a falta de compreensão das relações causais por parte das máquinas era talvez o maior obstáculo para lhes dar inteligência de nível humano. No último capítulo deste livro, voltarei às minhas raízes e, juntos, exploraremos as implicações da Revolução Causal para a inteligência artificial. Acredito que a IA forte é um objetivo alcançável e que não se deve temer — precisamente porque a causalidade é parte da solução. Um módulo de raciocínio causal dará às máquinas a capacidade de refletir sobre seus erros, de apontar fraquezas em seu \textit{software}, de funcionar como entidades morais e de conversar naturalmente com humanos sobre suas próprias escolhas e intenções.

\section{Uma planta da realidade}
Em nossa era, os leitores sem dúvida já ouviram termos como “conhecimento”, “informação”, “inteligência” e “dados”, e alguns podem se sentir confusos quanto às diferenças entre eles ou sobre como interagem. Agora proponho lançar mais um termo — “modelo causal” — nessa mistura, e o leitor pode, justificadamente, se perguntar se isso só aumentará a confusão.

Não aumentará! De fato, ele ancorará as noções elusivas de ciência, conhecimento e dados em um cenário concreto e significativo, e nos permitirá ver como os três trabalham juntos para produzir respostas a perguntas científicas difíceis. A Figura I.1 mostra uma planta de um “motor de inferência causal” que poderia lidar com o raciocínio causal de uma futura inteligência artificial. É importante perceber que esta não é apenas uma planta para o futuro, mas também um guia de como os modelos causais funcionam nas aplicações científicas de hoje e de como interagem com os dados.

O motor de inferência é uma máquina que aceita três tipos diferentes de entradas — Suposições, Consultas e Dados — e produz três tipos de saídas. A primeira das saídas é uma decisão Sim/Não quanto a se a consulta dada pode, em teoria, ser respondida sob o modelo causal existente, assumindo dados perfeitos e ilimitados. Se a resposta for Sim, o motor de inferência em seguida produz um \textit{Estimand} (estimando). Este é uma fórmula matemática que pode ser pensada como uma receita para gerar a resposta a partir de quaisquer dados hipotéticos, sempre que disponíveis. Por fim, depois que o motor de inferência recebeu a entrada de Dados, usará a receita para produzir uma Estimativa real para a resposta, juntamente com estimativas estatísticas da quantidade de incerteza nessa estimativa. Essa incerteza reflete o tamanho limitado do conjunto de dados, bem como possíveis erros de medição ou dados faltantes.

\textbf{FIGURA I.} Como um “motor de inferência” combina dados com conhecimento causal para produzir respostas a consultas de interesse. A caixa tracejada não faz parte do motor, mas é necessária para construí-lo. Setas também poderiam ser desenhadas das caixas 4 e 9 para a caixa 1, mas optei por manter o diagrama simples.

Para aprofundar o gráfico, rotulei as caixas de 1 a 9, que anotarei no contexto da consulta “Qual é o efeito do fármaco \textit{D} sobre a expectativa de vida \textit{L}?”

1. “Conhecimento” representa traços de experiência que o agente raciocinante teve no passado, incluindo observações passadas, ações passadas, educação e costumes culturais, que se consideram relevantes para a consulta de interesse. A caixa pontilhada em torno de “Conhecimento” indica que ele permanece implícito na mente do agente e não é explicitado formalmente no modelo.

2. A pesquisa científica sempre requer suposições simplificadoras — isto é, enunciados que o pesquisador considera dignos de tornar explícitos com base no Conhecimento disponível. Enquanto a maior parte do conhecimento do pesquisador permanece implícita em seu cérebro, apenas as Suposições veem a luz do dia e são encapsuladas no modelo. Elas podem, de fato, ser lidas a partir do modelo, o que levou alguns lógicos a concluir que um modelo não é nada mais do que uma lista de suposições. Os cientistas da computação discordam dessa afirmação, observando que o \textit{modo} como as suposições são representadas pode fazer uma diferença profunda na capacidade de especificá-las corretamente, delas extrair conclusões e até de estendê-las ou modificá-las à luz de evidências convincentes.

3. Existem várias opções de modelos causais: diagramas causais, equações estruturais, enunciados lógicos e assim por diante. Sou um forte defensor dos diagramas causais para quase todas as aplicações, principalmente por sua transparência, mas também pelas respostas explícitas que fornecem a muitas das perguntas que desejamos fazer. Para o propósito de construir o diagrama, a definição de “causalidade” é simples, se um tanto metafórica: uma variável \textit{X} é uma causa de \textit{Y} se \textit{Y} “escuta” \textit{X} e determina seu valor em resposta ao que ouve. Por exemplo, se suspeitamos que a expectativa de vida \textit{L} de um paciente “escuta” se o fármaco \textit{D} foi tomado, então chamamos \textit{D} de causa de \textit{L} e desenhamos uma seta de \textit{D} para \textit{L} em um diagrama causal. Naturalmente, a resposta à nossa consulta sobre \textit{D} e \textit{L} provavelmente dependerá de outras variáveis também, que devem ser representadas no diagrama juntamente com suas causas e efeitos. (Aqui, as denotaremos coletivamente por \textit{Z}.)

4. O padrão de escuta prescrito pelos caminhos do modelo causal geralmente resulta em padrões observáveis ou dependências nos dados. Esses padrões são chamados de “implicações testáveis” porque podem ser usados para testar o modelo. São enunciados como “Não há caminho conectando \textit{D} e \textit{L}”, que se traduz em um enunciado estatístico, “\textit{D} e \textit{L} são independentes” — isto é, encontrar \textit{D} não muda a verossimilhança de \textit{L}. Se os dados contradizem essa implicação, então precisamos revisar nosso modelo. Tais revisões requerem outro motor, que obtém suas entradas das caixas 4 e 7 e calcula o “grau de adequação” — isto é, o grau em que os Dados são compatíveis com as suposições do modelo. Por simplicidade, não mostrei esse segundo motor na Figura I.1.

5. As consultas submetidas ao motor de inferência são as perguntas científicas a que queremos responder. Devem ser formuladas em vocabulário causal. Por exemplo, qual é \textit{P}(\textit{L} | \textit{do}(\textit{D}))? Uma das principais conquistas da Revolução Causal tem sido tornar essa linguagem cientificamente transparente, bem como matematicamente rigorosa.

6. “\textit{Estimand}” vem do latim, significando “aquilo que deve ser estimado”. Trata-se de uma quantidade estatística a ser estimada a partir dos dados que, uma vez estimada, pode legitimamente representar a resposta à nossa consulta. Embora escrita como uma fórmula de probabilidade — por exemplo, \textit{P}(\textit{L} | \textit{D}, \textit{Z}) × \textit{P}(\textit{Z}) —, é de fato uma receita para responder à consulta causal a partir do tipo de dados de que dispomos, uma vez certificada pelo motor.

É muito importante perceber que, ao contrário da estimação tradicional em estatística, algumas consultas podem não ser respondíveis sob o modelo causal atual, mesmo após a coleta de qualquer quantidade de dados. Por exemplo, se nosso modelo mostra que tanto \textit{D} quanto \textit{L} dependem de uma terceira variável \textit{Z} (digamos, o estágio de uma doença), e se não temos nenhuma maneira de medir \textit{Z}, então a consulta \textit{P}(\textit{L} | \textit{do}(\textit{D})) não pode ser respondida. Nesse caso, é perda de tempo coletar dados. Em vez disso, precisamos voltar e refinar o modelo, seja adicionando novo conhecimento científico que possa nos permitir estimar \textit{Z}, seja fazendo suposições simplificadoras (com o risco de estarmos errados) — por exemplo, de que o efeito de \textit{Z} sobre \textit{D} é desprezível.

7. Os dados são os ingredientes que entram na receita do estimando. É crucial perceber que os dados são profundamente burros a respeito de relações causais. Eles nos falam de quantidades como \textit{P}(\textit{L} | \textit{D}) ou \textit{P}(\textit{L} | \textit{D}, \textit{Z}). É tarefa do estimando nos dizer como cozinhar essas quantidades estatísticas em uma expressão que, com base nas suposições do modelo, seja logicamente equivalente à consulta causal — digamos, \textit{P}(\textit{L} | \textit{do}(\textit{D})).

Observe que toda a noção de estimandos e, de fato, toda a parte superior da Figura I não existe nos métodos tradicionais de análise estatística. Lá, o estimando e a consulta coincidem. Por exemplo, se estamos interessados na proporção de pessoas entre aquelas com expectativa de vida \textit{L} que tomaram o fármaco \textit{D}, simplesmente escrevemos essa consulta como \textit{P}(\textit{D} | \textit{L}). A mesma quantidade seria nosso estimando. Isso já especifica quais proporções nos dados precisam ser estimadas e não requer nenhum conhecimento causal. Por essa razão, alguns estatísticos até hoje acham extremamente difícil compreender por que algum conhecimento se situa fora da província da estatística e por que os dados sozinhos não podem compensar a falta de conhecimento científico.

8. A estimativa é o que sai do forno. Contudo, é apenas aproximada por causa de outro fato do mundo real sobre os dados: eles são sempre apenas uma amostra finita de uma população teoricamente infinita. Em nosso exemplo recorrente, a amostra consiste nos pacientes que escolhemos estudar. Mesmo se os escolhermos ao acaso, há sempre alguma chance de que as proporções medidas na amostra não sejam representativas das proporções na população em geral. Felizmente, a disciplina da estatística, empoderada por técnicas avançadas de aprendizado de máquina, nos dá muitas, muitas maneiras de gerir essa incerteza — estimadores de máxima verossimilhança, escores de propensão, intervalos de confiança, testes de significância e assim por diante.

9. No fim, se nosso modelo estiver correto e nossos dados forem suficientes, obtemos uma resposta à nossa consulta causal, como “O fármaco \textit{D} aumenta a expectativa de vida \textit{L} de pacientes diabéticos \textit{Z} em 30 por cento, mais ou menos 20 por cento”. Viva! A resposta também se somará ao nosso conhecimento científico (caixa 1) e, se as coisas não saírem como esperávamos, poderá sugerir algumas melhorias em nosso modelo causal (caixa 3).

Esse fluxograma pode parecer complicado à primeira vista, e você pode se perguntar se é realmente necessário. De fato, em nossas vidas ordinárias, de algum modo somos capazes de fazer juízos causais sem passar conscientemente por um processo tão complicado e certamente sem recorrer à matemática de probabilidades e proporções. Nossa intuição causal sozinha é geralmente suficiente para lidar com o tipo de incerteza que encontramos nas rotinas domésticas ou mesmo em nossas vidas profissionais. Mas se queremos ensinar um robô burro a pensar causalmente, ou se estamos empurrando as fronteiras do conhecimento científico, onde não temos intuição para nos guiar, então um procedimento cuidadosamente estruturado como este é obrigatório.

Quero especialmente destacar o papel dos dados no processo acima. Primeiro, observe que coletamos dados somente \textit{depois} de postular o modelo causal, \textit{depois} de enunciar a consulta científica a que desejamos responder e \textit{depois} de derivar o estimando. Isso contrasta com a abordagem estatística tradicional, mencionada acima, que nem sequer tem um modelo causal.

Mas o mundo científico de nossos dias apresenta um novo desafio ao raciocínio sólido sobre causas e efeitos. Embora a consciência da necessidade de um modelo causal tenha crescido a passos largos entre as ciências, muitos pesquisadores em inteligência artificial gostariam de pular o passo difícil de construir ou adquirir um modelo causal e confiar unicamente nos dados para todas as tarefas cognitivas. A esperança — e, no presente, ela costuma ser silenciosa — é que os próprios dados nos guiem às respostas certas sempre que perguntas causais surgirem.

Sou um cético declarado dessa tendência porque sei o quão profundamente burros os dados são a respeito de causas e efeitos. Por exemplo, informação sobre os efeitos de ações ou intervenções simplesmente não está disponível em dados brutos, a menos que seja coletada por manipulação experimental controlada. Em contraste, se estamos de posse de um modelo causal, frequentemente podemos prever o resultado de uma intervenção a partir de dados sem intervenção, “de mãos livres”.

O caso a favor dos modelos causais se torna ainda mais convincente quando buscamos responder a consultas contrafactuais como “O que teria acontecido se tivéssemos agido de outro modo?”. Discutiremos os contrafactuais em grande detalhe porque são as consultas mais desafiadoras para qualquer inteligência artificial. Estão também no cerne dos avanços cognitivos que nos tornaram humanos e das capacidades imaginativas que tornaram a ciência possível. Explicaremos também por que qualquer consulta sobre o mecanismo pelo qual as causas transmitem seus efeitos — a mais prototípica pergunta “Por quê?” — é de fato uma pergunta contrafactual disfarçada. Assim, se algum dia quisermos que robôs respondam a perguntas “Por quê?” ou mesmo compreendam o que elas significam, devemos municiá-los com um modelo causal e ensiná-los a responder a consultas contrafactuais, como na Figura I.1.

Outra vantagem que os modelos causais têm e de que a mineração de dados e o aprendizado profundo carecem é a adaptabilidade. Observe que, na Figura I.1, o estimando é computado com base apenas no modelo causal, antes de um exame das especificidades dos dados. Isso torna o motor de inferência causal supremamente adaptável, porque o estimando computado vale para quaisquer dados compatíveis com o modelo qualitativo, independentemente das relações numéricas entre as variáveis.

Para ver por que essa adaptabilidade é importante, compare esse motor com um agente de aprendizado — neste caso um humano, mas em outros casos talvez um algoritmo de aprendizado profundo, ou talvez um humano usando um algoritmo de aprendizado profundo — tentando aprender unicamente a partir dos dados. Observando o resultado \textit{L} de muitos pacientes a quem se deu o fármaco \textit{D}, ela é capaz de prever a probabilidade de que um paciente com características \textit{Z} sobreviva \textit{L} anos. Agora ela é transferida para um hospital diferente, em outra parte da cidade, onde as características populacionais (dieta, higiene, hábitos de trabalho) são diferentes. Mesmo se essas novas características meramente modificarem as relações numéricas entre as variáveis registradas, ela ainda terá de se retreinar e aprender uma nova função de previsão do zero. É tudo o que um programa de aprendizado profundo pode fazer: ajustar uma função aos dados. Por outro lado, se ela possuísse um modelo de como o fármaco operava e sua estrutura causal permanecesse intacta no novo local, então o estimando que obteve no treinamento permaneceria válido. Poderia ser aplicado aos novos dados para gerar uma nova função de previsão específica da população.

Muitas perguntas científicas parecem diferentes “através de uma lente causal”, e tenho me deleitado em brincar com essa lente, que ao longo dos últimos vinte e cinco anos tem sido cada vez mais empoderada por novos \textit{insights} e novas ferramentas. Espero e acredito que os leitores deste livro compartilharão de meu deleite. Portanto, gostaria de encerrar esta introdução com uma prévia de algumas das atrações que virão neste livro.

O Capítulo 1 reúne os três degraus da observação, da intervenção e dos contrafactuais na Escada da Causalidade[^7], a metáfora central deste livro. Também o exporá aos fundamentos do raciocínio com diagramas causais, nossa principal ferramenta de modelagem, e o colocará bem a caminho de se tornar um raciocinador causal proficiente — de fato, você estará bem à frente de gerações de cientistas de dados que tentaram interpretar dados através de uma lente cega a modelos, alheios às distinções que a Escada da Causalidade ilumina.

O Capítulo 2 conta a bizarra história de como a disciplina da estatística infligiu a si mesma a cegueira causal, com efeitos de longo alcance para todas as ciências que dependem de dados. Conta também a história de um dos grandes heróis deste livro, o geneticista Sewall Wright, que nos anos 1920 desenhou os primeiros diagramas causais e por muitos anos foi um dos poucos cientistas que ousou levar a causalidade a sério. O Capítulo 3 relata a história igualmente curiosa de como me tornei um convertido à causalidade por meio de meu trabalho em IA e, particularmente, em redes bayesianas. Estas foram a primeira ferramenta que permitiu aos computadores pensar em “tons de cinza” — e por um tempo acreditei que elas guardavam a chave para destravar a IA. No final dos anos 1980, convenci-me de que estava errado, e este capítulo conta minha jornada de profeta a apóstata. Não obstante, as redes bayesianas permanecem uma ferramenta muito importante para a IA e ainda encapsulam grande parte do fundamento matemático dos diagramas causais. Além de uma introdução suave e orientada para a causalidade à regra de Bayes e aos métodos bayesianos de raciocínio, o Capítulo 3 entreterá o leitor com exemplos de aplicações reais de redes bayesianas.

O Capítulo 4 conta a principal contribuição da estatística à inferência causal: o ensaio clínico randomizado (ECR)[^8]. De uma perspectiva causal, o ECR é uma ferramenta feita pelo homem para desvelar a consulta \textit{P}(\textit{L} | \textit{do}(\textit{D})), que é uma propriedade da natureza. Seu principal propósito é dissociar variáveis de interesse (digamos, \textit{D} e \textit{L}) de outras variáveis (\textit{Z}) que de outro modo as afetariam a ambas. Desarmar as distorções, ou o “confundimento”[^9], produzidas por tais variáveis à espreita tem sido um problema de um século. Este capítulo conduz o leitor por uma solução surpreendentemente simples para o problema geral do confundimento, que você captará em dez minutos de brincar de traçar caminhos em um diagrama.

O Capítulo 5 dá conta de um momento seminal na história da causalidade e, de fato, na história da ciência, quando estatísticos lutaram com a pergunta de se fumar causa câncer de pulmão. Incapazes de usar sua ferramenta favorita, o ensaio clínico randomizado, lutaram para concordar em uma resposta ou mesmo em como dar sentido à pergunta. O debate sobre o fumo coloca a importância da causalidade em seu foco mais nítido. Milhões de vidas foram perdidas ou encurtadas porque os cientistas não tinham uma linguagem ou metodologia adequada para responder a perguntas causais.

O Capítulo 6 será, espero, uma divertida digressão para o leitor após os assuntos sérios do Capítulo 5. Este é um capítulo de paradoxos: o paradoxo de Monty Hall, o paradoxo de Simpson, o paradoxo de Berkson e outros. Paradoxos clássicos como esses podem ser apreciados como quebra-cabeças mentais, mas também têm um lado sério, especialmente quando vistos de uma perspectiva causal. De fato, quase todos eles representam embates com a intuição causal e, portanto, revelam a anatomia dessa intuição. Foram canários na mina de carvão que deveriam ter alertado os cientistas de que a intuição humana se funda em uma lógica causal, e não estatística. Acredito que o leitor apreciará esse novo toque em seus paradoxos antigos favoritos.

Os Capítulos 7 a 9, enfim, levam os leitores a uma emocionante ascensão da Escada da Causalidade. Começamos no Capítulo 7 com perguntas sobre intervenção e explicamos como meus estudantes e eu passamos por uma luta de vinte anos para automatizar as respostas a perguntas do tipo \textit{do}. Tivemos sucesso, e este capítulo explica as entranhas do “motor de inferência causal”, que produz a resposta sim/não e o estimando na Figura I.1. Estudar esse motor capacitará o leitor a detectar certos padrões no diagrama causal que entregam respostas imediatas à consulta causal. Esses padrões se chamam ajustamento pela “porta dos fundos”[^10], ajustamento pela “porta da frente” e variáveis instrumentais — os cavalos de batalha da inferência causal na prática.

O Capítulo 8 o leva ao topo da escada ao discutir os contrafactuais. Estes têm sido vistos como parte fundamental da causalidade pelo menos desde 1748, quando o filósofo escocês David Hume propôs a seguinte definição um tanto tortuosa de causalidade: “Podemos definir uma causa como um objeto seguido por outro, e onde todos os objetos semelhantes ao primeiro são seguidos por objetos semelhantes ao segundo. Ou, em outras palavras, onde, se o primeiro objeto não tivesse existido, o segundo nunca teria existido.” David Lewis, filósofo da Universidade de Princeton que morreu em 2001, apontou que Hume de fato deu duas definições, e não uma: a primeira de regularidade (isto é, a causa é regularmente seguida pelo efeito) e a segunda do contrafactual (“se o primeiro objeto não tivesse existido…”). Enquanto filósofos e cientistas tinham prestado atenção sobretudo à definição de regularidade, Lewis argumentou que a definição contrafactual se alinha mais de perto com a intuição humana: “Pensamos em uma causa como algo que faz diferença, e a diferença que ela faz deve ser uma diferença em relação ao que teria acontecido sem ela.”

Os leitores ficarão entusiasmados ao descobrir que agora podemos ir além dos debates acadêmicos e computar um valor real (ou probabilidade) para qualquer consulta contrafactual, por mais convoluta que seja. De especial interesse são as perguntas relativas a causas necessárias e suficientes de eventos observados. Por exemplo, quão provável é que a ação do réu tenha sido uma causa necessária da lesão do demandante? Quão provável é que a mudança climática causada pelo homem seja uma causa suficiente de uma onda de calor?

Por fim, o Capítulo 9 discute o tópico da mediação. Você pode ter se perguntado, quando falamos de desenhar setas em um diagrama causal, se deveríamos desenhar uma seta do fármaco \textit{D} para a expectativa de vida \textit{L} se o fármaco afeta a expectativa de vida apenas por meio de seu efeito sobre a pressão arterial \textit{Z} (um mediador). Em outras palavras, o efeito de \textit{D} sobre \textit{L} é direto ou indireto? E se for ambos, como avaliamos sua importância relativa? Tais perguntas não são apenas de grande interesse científico, mas também têm ramificações práticas; se compreendermos o mecanismo pelo qual um fármaco age, poderemos desenvolver outros fármacos com o mesmo efeito que sejam mais baratos ou tenham menos efeitos colaterais. O leitor ficará satisfeito ao descobrir como essa busca milenar por um mecanismo de mediação foi reduzida a um exercício algébrico e como os cientistas estão usando as novas ferramentas do kit causal para resolver tais problemas.

O Capítulo 10 fecha o livro voltando ao problema que inicialmente me levou à causalidade: o problema de automatizar a inteligência de nível humano (às vezes chamada de “IA forte”). Acredito que o raciocínio causal é essencial para que as máquinas se comuniquem conosco em nossa própria linguagem sobre políticas, experimentos, explicações, teorias, arrependimento, responsabilidade, livre-arbítrio e obrigações — e, eventualmente, para que tomem suas próprias decisões morais.

Se eu pudesse resumir a mensagem deste livro em uma frase sucinta, seria: você é mais inteligente do que seus dados. Os dados não compreendem causas e efeitos; os humanos, sim. Espero que a nova ciência da inferência causal nos permita compreender melhor \textit{como} o fazemos, porque não há melhor maneira de nos compreendermos do que emulando a nós mesmos. Na era dos computadores, essa nova compreensão traz também a perspectiva de amplificar nossas capacidades inatas, de modo que possamos dar melhor sentido aos dados, sejam eles grandes ou pequenos.

\chapter{1 A Escada da Causalidade}
No princípio… Eu tinha provavelmente seis ou sete anos quando li pela primeira vez a história de Adão e Eva no Jardim do Éden. Meus colegas e eu não nos surpreendemos em absoluto com as exigências caprichosas de Deus, que lhes proibia comer da Árvore do Conhecimento. As divindades têm suas razões, pensávamos. O que mais nos intrigava era a ideia de que, assim que comeram da Árvore do Conhecimento, Adão e Eva se tornaram conscientes, como nós, de sua nudez.

Na adolescência, nosso interesse foi se deslocando aos poucos para os aspectos mais filosóficos da história. (Estudantes israelenses leem o Gênesis várias vezes por ano.) De preocupação central para nós era a noção de que o surgimento do conhecimento humano não fora um processo alegre, e sim um processo doloroso, acompanhado de desobediência, culpa e castigo. Valeria a pena abrir mão da vida despreocupada do Éden?, perguntavam alguns. As revoluções agrícola e científica que se seguiram valeram as dificuldades econômicas, as guerras e as injustiças sociais que a vida moderna acarreta?

Não me entendam mal: não éramos criacionistas; mesmo nossos professores eram darwinistas no fundo. Sabíamos, porém, que o autor que coreografou a história do Gênesis se debatia para responder às perguntas filosóficas mais prementes de seu tempo. Suspeitávamos, da mesma forma, que essa história carregava as marcas culturais do processo real pelo qual o \textit{Homo sapiens} conquistou o domínio sobre nosso planeta.

Qual foi, então, a sequência de passos nesse processo veloz e super-evolutivo? Meu interesse por essas questões esmaeceu no início da minha carreira como professor de engenharia, mas reacendeu-se de súbito nos anos 1990, quando, ao escrever meu livro \textit{Causality}, deparei-me com a Escada da Causalidade.

Ao reler o Gênesis pela centésima vez, notei uma nuance que, de alguma forma, havia escapado à minha atenção durante todos aqueles anos. Quando Deus encontra Adão escondido no jardim, pergunta: “Comeste da árvore de que te proibi comer?” E Adão responde: “A mulher que me deste por companheira, ela me deu do fruto da árvore, e eu comi.” “Que é isso que fizeste?”, pergunta Deus a Eva. Ela replica: “A serpente me enganou, e eu comi.”

Como sabemos, esse jogo de acusações não funcionou muito bem com o Todo-Poderoso, que os expulsou a ambos do jardim. Mas eis o ponto que eu havia deixado passar: Deus perguntou “o quê”, e eles responderam “por quê”. Deus pediu os fatos, e eles responderam com explicações. Mais ainda: ambos estavam plenamente convencidos de que nomear causas de algum modo pintaria suas ações sob uma luz diferente. De onde lhes veio essa ideia?

Para mim, essas nuances carregavam três implicações profundas. Primeiro, muito cedo em nossa evolução, nós, humanos, compreendemos que o mundo não é feito apenas de fatos secos (o que hoje poderíamos chamar de dados); antes, esses fatos estão colados uns aos outros por uma intrincada teia de relações de causa e efeito. Segundo, as explicações causais, e não os fatos secos, constituem a maior parte do nosso conhecimento e deveriam ser a pedra angular da inteligência de máquina. Por fim, nossa transição de processadores de dados a elaboradores de explicações não foi gradual; foi um salto que exigiu um empurrão externo de um fruto incomum.

Isso combinava perfeitamente com o que eu havia observado em termos teóricos na Escada da Causalidade: nenhuma máquina consegue extrair explicações de dados brutos. Ela precisa de um empurrão.

Se buscarmos confirmação dessas mensagens na ciência evolutiva, não encontraremos, é claro, a Árvore do Conhecimento, mas ainda assim vemos uma grande transição inexplicada. Compreendemos agora que os humanos evoluíram de ancestrais semelhantes a símios ao longo de um período de 5 a 6 milhões de anos e que tais processos evolutivos graduais não são incomuns na vida na Terra. Mas, grosso modo, nos últimos 50 mil anos, algo único aconteceu, a que alguns chamam de Revolução Cognitiva e outros (com um toque de ironia) de Grande Salto à Frente. Os humanos adquiriram a capacidade de modificar seu ambiente e suas próprias habilidades a um ritmo dramaticamente mais rápido.

Por exemplo, ao longo de milhões de anos, águias e corujas desenvolveram uma visão verdadeiramente impressionante — e, no entanto, nunca inventaram óculos, microscópios, telescópios ou óculos de visão noturna. Os humanos produziram esses milagres em questão de séculos. Chamo esse fenômeno de “aceleração super-evolutiva”.

Alguns leitores podem objetar que estou comparando alhos com bugalhos, evolução com engenharia, mas esse é exatamente o meu ponto. A evolução nos dotou da capacidade de engenheirar nossas vidas, um dom que não concedeu a águias e corujas, e a pergunta, de novo, é “Por quê?”. Qual facilidade computacional os humanos adquiriram de repente e as águias não?

Muitas teorias foram propostas, mas uma é especialmente pertinente à ideia de causalidade. Em seu livro \textit{Sapiens}, o historiador Yuval Harari postula que a capacidade de nossos ancestrais de imaginar coisas inexistentes foi a chave de tudo, pois lhes permitiu comunicar-se melhor. Antes dessa mudança, eles só podiam confiar em pessoas de sua família imediata ou de sua tribo. Depois, sua confiança se estendeu a comunidades maiores, unidas por fantasias comuns (por exemplo, a crença em divindades invisíveis, porém imagináveis, na vida após a morte e na divindade do líder) e por expectativas.

Concorde-se ou não com a teoria de Harari, a conexão entre imaginar e relações causais é quase evidente por si mesma. É inútil perguntar pelas causas das coisas a menos que se possa imaginar suas consequências. Reciprocamente, não se pode afirmar que Eva o fez comer da árvore a menos que se possa imaginar um mundo em que, contrafactualmente, ela não lhe entregou a maçã.

De volta a nossos ancestrais \textit{Homo sapiens}: sua recém-adquirida imaginação causal lhes permitiu fazer muitas coisas de modo mais eficiente por meio de um processo astucioso a que chamamos “planejamento”. Imagine uma tribo preparando-se para caçar um mamute. O que seria preciso para que tivessem sucesso?

Minhas habilidades de caça a mamutes estão um tanto enferrujadas, confesso, mas, como estudante de máquinas pensantes, aprendi uma coisa: uma entidade pensante (computador, homem das cavernas ou professor) só consegue realizar uma tarefa de tal magnitude planejando as coisas com antecedência — decidindo quantos caçadores recrutar; avaliando, dadas as condições do vento, a direção de onde se aproximar do mamute; em suma, imaginando e comparando as consequências de várias estratégias de caça. Para isso, a entidade pensante precisa possuir, consultar e manipular um modelo mental de sua realidade.

\textbf{Figura 1.1.} Causas percebidas de uma caçada bem-sucedida a mamutes.

\begin{figure}[ht]
\centering
\includegraphics[width=0.92\textwidth]{../figuras/fig_1_1.png}
\caption{Figura 1.1. Causas percebidas de uma caçada bem-sucedida a mamutes.}
\end{figure}
\textit{Figura 1.1. Causas percebidas de uma caçada bem-sucedida a mamutes.}

A Figura 1.1 mostra como poderíamos desenhar tal modelo mental. Cada ponto na Figura 1.1 representa uma causa do sucesso. Note que há múltiplas causas e que nenhuma delas é determinística. Isto é, não podemos ter certeza de que ter mais caçadores permitirá o sucesso ou de que a chuva o impedirá, mas esses fatores de fato alteram a probabilidade de sucesso.

O modelo mental é a arena em que a imaginação se desenrola. Ele nos permite experimentar diferentes cenários fazendo alterações locais no modelo. Em algum lugar no modelo mental de nossos caçadores havia uma sub-rotina que avaliava o efeito do número de caçadores. Quando consideravam acrescentar mais, não precisavam reavaliar todos os outros fatores do zero. Podiam fazer uma mudança local no modelo, substituindo “Caçadores = 8” por “Caçadores = 9”, e reavaliar a probabilidade de sucesso. Essa modularidade é uma característica-chave dos modelos causais.

Não pretendo, é claro, insinuar que os primeiros humanos de fato desenhavam um modelo pictórico como este. Mas quando buscamos emular o pensamento humano em um computador, ou de fato quando tentamos resolver problemas científicos pouco familiares, desenhar uma figura explícita de pontos e setas é extremamente útil. Esses diagramas causais são o núcleo computacional do “motor de inferência causal” descrito na Introdução.

\section{Os três níveis da causalidade}
Até aqui posso ter dado a impressão de que a capacidade de organizar nosso conhecimento do mundo em causas e efeitos era monolítica e foi adquirida de uma só vez. Na verdade, minha pesquisa em aprendizado de máquina me ensinou que um aprendiz causal precisa dominar pelo menos três níveis distintos de habilidade cognitiva: ver, fazer e imaginar.

O primeiro, ver ou observar, implica a detecção de regularidades em nosso ambiente e é compartilhado por muitos animais, bem como pelos primeiros humanos antes da Revolução Cognitiva. O segundo, fazer, implica predizer o(s) efeito(s) de alterações deliberadas do ambiente e escolher entre essas alterações para produzir um resultado desejado. Apenas um punhado de espécies demonstrou elementos dessa habilidade. O uso de ferramentas, desde que seja intencional e não meramente acidental ou copiado dos ancestrais, pode ser tomado como sinal de se ter alcançado esse segundo nível. Ainda assim, mesmo usuários de ferramentas não necessariamente possuem uma “teoria” de sua ferramenta que lhes diga por que ela funciona e o que fazer quando não funciona. Para isso, é preciso ter alcançado um nível de compreensão que permita imaginar. Foi sobretudo esse terceiro nível que nos preparou para revoluções posteriores na agricultura e na ciência e conduziu a uma mudança súbita e drástica no impacto de nossa espécie sobre o planeta.

Não posso provar isso, mas posso provar matematicamente que os três níveis diferem fundamentalmente, cada um liberando capacidades que os de baixo não possuem. O arcabouço que uso para mostrar isso remonta a Alan Turing, o pioneiro da pesquisa em inteligência artificial (IA), que propôs classificar um sistema cognitivo em termos das consultas que ele é capaz de responder. Essa abordagem é excepcionalmente frutífera quando falamos de causalidade, porque contorna longas e improdutivas discussões sobre o que exatamente a causalidade é e se concentra, em vez disso, na pergunta concreta e respondível “O que um raciocinador causal pode fazer?”. Ou, mais precisamente, o que um organismo que possui um modelo causal pode computar que um organismo desprovido de tal modelo não pode?

\textbf{Figura 1.2.} A Escada da Causalidade, com organismos representativos em cada nível. A maioria dos animais, bem como as máquinas de aprendizado atuais, está no primeiro degrau, aprendendo por associação. Usuários de ferramentas, como os primeiros humanos, estão no segundo degrau se agem por planejamento e não meramente por imitação. Também podemos usar experimentos para aprender os efeitos de intervenções, e presume-se que seja assim que os bebês adquirem boa parte de seu conhecimento causal. Aprendizes contrafactuais, no degrau superior, podem imaginar mundos que não existem e inferir razões para fenômenos observados. (Fonte: desenho de Maayan Harel.)

\begin{figure}[ht]
\centering
\includegraphics[width=0.92\textwidth]{../figuras/fig_1_2.png}
\caption{Figura 1.2. A Escada da Causalidade, com organismos representativos em cada nível. A maioria dos animais, bem como as máquinas de aprendizado atuais, está no primeiro degrau, aprendendo por associação. Usuários de ferramentas, como os primeiros humanos, estão no segundo degrau se agem por planejamento e não meramente por imitação. Também podemos usar experimentos para aprender os efeitos de intervenções, e presume-se que seja assim que os bebês adquirem boa parte de seu conhecimento causal. Aprendizes contrafactuais, no degrau superior, podem imaginar mundos que não existem e inferir razões para fenômenos observados. (Fonte: desenho de Maayan Harel.)}
\end{figure}
\textit{Figura 1.2. A Escada da Causalidade, com organismos representativos em cada nível. A maioria dos animais, bem como as máquinas de aprendizado atuais, está no primeiro degrau, aprendendo por associação. Usuários de ferramentas, como os primeiros humanos, estão no segundo degrau se agem por planejamento e não meramente por imitação. Também podemos usar experimentos para aprender os efeitos de intervenções, e presume-se que seja assim que os bebês adquirem boa parte de seu conhecimento causal. Aprendizes contrafactuais, no degrau superior, podem imaginar mundos que não existem e inferir razões para fenômenos observados. (Fonte: desenho de Maayan Harel.)}

Enquanto Turing buscava uma classificação binária — humano ou não humano —, a nossa tem três camadas, correspondentes a consultas causais progressivamente mais poderosas. Usando esses critérios, podemos reunir os três níveis de consultas em uma única Escada da Causalidade (Figura 1.2), metáfora à qual voltaremos repetidas vezes. Vamos dedicar algum tempo a considerar cada degrau da escada em detalhe.

No primeiro nível, a \textbf{associação}, buscamos regularidades nas observações. É o que uma coruja faz ao observar como um rato se move e calcular onde o roedor provavelmente estará um instante depois, e é o que um programa de computador de Go faz ao estudar um banco de dados com milhões de partidas de Go para descobrir quais lances estão associados a uma porcentagem maior de vitórias.

Dizemos que um evento está associado a outro se observar um muda a probabilidade de observar o outro.

O primeiro degrau da escada pede predições baseadas em observações passivas. Caracteriza-se pela pergunta “E se eu vir…?”. Por exemplo, imagine um diretor de marketing de uma loja de departamentos que pergunta: “Qual a probabilidade de um cliente que comprou pasta de dente também comprar fio dental?”. Perguntas desse tipo são o pão de cada dia da estatística, e respondem-se, antes de tudo, coletando e analisando dados. Em nosso caso, a pergunta pode ser respondida pegando-se primeiro os dados do comportamento de compra de todos os clientes, selecionando apenas os que compraram pasta de dente e, concentrando-se nesse grupo, calculando a proporção dos que também compraram fio dental. Essa proporção, também conhecida como “probabilidade condicional”, mede (para grandes volumes de dados) o grau de associação entre “comprar pasta de dente” e “comprar fio dental”. Simbolicamente, podemos escrevê-la como \textit{P}(fio dental | pasta de dente). O “\textit{P}” significa “probabilidade”, e a barra vertical significa “dado que se vê”.

Os estatísticos desenvolveram muitos métodos elaborados para reduzir um grande corpo de dados e identificar associações entre variáveis. A “correlação” ou a “regressão”, medidas típicas de associação mencionadas com frequência neste livro, envolvem ajustar uma reta a um conjunto de pontos de dados e tomar a inclinação dessa reta. Algumas associações podem ter interpretações causais óbvias; outras, não. Mas a estatística sozinha não pode dizer qual é a causa e qual é o efeito, pasta de dente ou fio dental. Do ponto de vista do gerente de vendas, isso pode não importar de verdade. Boas predições não precisam ter boas explicações. A coruja pode ser uma boa caçadora sem compreender por que o rato sempre vai do ponto A ao ponto B.

Alguns leitores talvez se surpreendam ao ver que coloquei as máquinas de aprendizado atuais firmemente no degrau um da Escada da Causalidade, compartilhando a sabedoria de uma coruja. Parece que quase todos os dias ouvimos falar de avanços rápidos em sistemas de aprendizado de máquina — carros autônomos, sistemas de reconhecimento de fala e, especialmente nos últimos anos, algoritmos de aprendizado profundo (ou redes neurais profundas). Como podem ainda estar apenas no nível um?

Os sucessos do aprendizado profundo foram verdadeiramente notáveis e pegaram muitos de nós de surpresa. Não obstante, o aprendizado profundo teve êxito sobretudo ao mostrar que certas perguntas ou tarefas que pensávamos ser difíceis de fato não o são. Não abordou as perguntas verdadeiramente difíceis que continuam a nos impedir de alcançar uma IA de tipo humano. Como resultado, o público acredita que a “IA forte”, máquinas que pensam como humanos, está logo ali na esquina ou talvez já esteja entre nós. Na realidade, nada poderia estar mais longe da verdade.

Concordo plenamente com Gary Marcus, neurocientista da Universidade de Nova York, que escreveu recentemente no \textit{New York Times} que o campo da inteligência artificial está “repleto de microdescobertas” — o tipo de coisa que rende bons \textit{press releases} — mas as máquinas ainda estão decepcionantemente longe de uma cognição de tipo humano. Meu colega de ciência da computação na Universidade da Califórnia em Los Angeles, Adnan Darwiche, intitulou um artigo de posição “Inteligência de nível humano ou habilidades de tipo animal?”, o que, a meu ver, enquadra a pergunta exatamente do jeito certo. O objetivo da IA forte é produzir máquinas com inteligência de tipo humano, capazes de conversar com os humanos e de orientá-los. O aprendizado profundo, em vez disso, nos deu máquinas com habilidades verdadeiramente impressionantes, mas sem inteligência. A diferença é profunda e reside na ausência de um modelo da realidade.

Assim como faziam trinta anos atrás, os programas de aprendizado de máquina (incluindo os de redes neurais profundas) operam quase inteiramente em um modo associativo. São impulsionados por um fluxo de observações ao qual tentam ajustar uma função, mais ou menos do mesmo modo que um estatístico tenta ajustar uma reta a um conjunto de pontos. As redes neurais profundas acrescentaram muitas camadas à complexidade da função ajustada, mas os dados brutos ainda conduzem o processo de ajuste. Continuam a melhorar em precisão à medida que mais dados são ajustados, mas não se beneficiam da “aceleração super-evolutiva”. Se, por exemplo, os programadores de um carro autônomo quiserem que ele reaja de modo diferente a novas situações, precisam acrescentar explicitamente essas novas reações. A máquina não vai descobrir por si mesma que um pedestre com uma garrafa de uísque na mão provavelmente responderá de modo diferente a uma buzina. Essa falta de flexibilidade e adaptabilidade é inevitável em qualquer sistema que trabalhe no primeiro nível da Escada da Causalidade.

Subimos ao próximo nível de consultas causais quando começamos a mudar o mundo. Uma pergunta típica desse nível é “O que acontecerá com nossas vendas de fio dental se dobrarmos o preço da pasta de dente?”. Isso já exige um novo tipo de conhecimento, ausente dos dados, que encontramos no degrau dois da Escada da Causalidade, a \textbf{intervenção}.

A intervenção está acima da associação porque envolve não apenas ver, mas mudar o que é. Ver fumaça nos conta uma história totalmente diferente sobre a probabilidade de fogo do que fazer fumaça. Não podemos responder a perguntas sobre intervenções com dados coletados passivamente, por maior que seja o conjunto de dados ou por mais profunda que seja a rede neural. Muitos cientistas ficaram bastante traumatizados ao descobrir que nenhum dos métodos que aprenderam em estatística é suficiente sequer para articular, que dirá responder, a uma pergunta simples como “O que acontece se dobrarmos o preço?”. Sei disso porque, em muitas ocasiões, ajudei-os a subir ao próximo degrau da escada.

Por que não podemos responder à nossa pergunta sobre o fio dental apenas por observação? Por que não simplesmente ir ao nosso vasto banco de dados de compras anteriores e ver o que aconteceu antes quando a pasta de dente custava o dobro? A razão é que, nas ocasiões anteriores, o preço pode ter sido mais alto por motivos diferentes. Por exemplo, o produto pode ter estado em falta, e todas as outras lojas também tiveram de aumentar o preço. Mas agora você está considerando uma intervenção deliberada que fixará um novo preço independentemente das condições de mercado. O resultado pode ser bem diferente do que quando o cliente não conseguia encontrar um negócio melhor em outro lugar. Se você tivesse dados sobre as condições de mercado que existiam nas ocasiões anteriores, talvez pudesse fazer uma predição melhor… mas de quais dados você precisa? E, depois, como descobriria a resposta? Essas são exatamente as perguntas que a ciência da inferência causal nos permite responder.

Um modo muito direto de predizer o resultado de uma intervenção é experimentá-la sob condições cuidadosamente controladas. Empresas de \textit{big data} como o Facebook sabem disso e realizam constantemente experimentos para ver o que acontece se os itens na tela forem dispostos de modo diferente ou se o cliente receber um \textit{prompt} diferente (ou até um preço diferente).

Mais interessante e menos amplamente conhecido — mesmo no Vale do Silício — é que predições bem-sucedidas dos efeitos de intervenções às vezes podem ser feitas mesmo sem um experimento. Por exemplo, o gerente de vendas poderia desenvolver um modelo do comportamento do consumidor que incluísse as condições de mercado. Mesmo que não tenha dados sobre todos os fatores, pode ter dados sobre substitutos-chave suficientes para fazer a predição. Um modelo causal suficientemente forte e preciso pode nos permitir usar dados do degrau um (observacionais) para responder a consultas do degrau dois (intervencionais). Sem o modelo causal, não poderíamos ir do degrau um ao degrau dois. É por isso que os sistemas de aprendizado profundo (enquanto usarem apenas dados do degrau um e não tiverem um modelo causal) nunca serão capazes de responder a perguntas sobre intervenções, que por definição quebram as regras do ambiente em que a máquina foi treinada.

Como esses exemplos ilustram, a consulta definidora do segundo degrau da Escada da Causalidade é “E se fizermos…?”. O que acontecerá se mudarmos o ambiente? Podemos escrever esse tipo de consulta como \textit{P}(fio dental | \textit{do}(pasta de dente)), que pergunta sobre a probabilidade de vendermos fio dental a um certo preço, dado que fixamos o preço da pasta de dente em outro preço.

Outra pergunta popular no segundo nível da causalidade é “Como?”, prima de “E se fizermos…?”. Por exemplo, o gerente pode nos dizer que temos pasta de dente demais no depósito. “Como podemos vendê-la?”, pergunta ele. Isto é, a que preço deveríamos fixá-la? De novo, a pergunta se refere a uma intervenção, que queremos realizar mentalmente antes de decidir se e como fazê-la na vida real. Isso exige um modelo causal.

Realizamos intervenções o tempo todo em nossa vida cotidiana, embora normalmente não usemos um termo tão sofisticado para elas. Por exemplo, quando tomamos aspirina para curar uma dor de cabeça, estamos intervindo sobre uma variável (a quantidade de aspirina em nosso corpo) a fim de afetar outra (o status de nossa dor de cabeça). Se estivermos corretos em nossa crença causal sobre a aspirina, a variável de “resultado” responderá mudando de “dor de cabeça” para “sem dor de cabeça”.

Embora raciocinar sobre intervenções seja um passo importante na escada causal, ainda não responde a todas as perguntas de interesse. Podemos nos perguntar: Minha dor de cabeça sumiu agora, mas por quê? Foi a aspirina que tomei? A comida que comi? A boa notícia que ouvi? Essas consultas nos levam ao degrau superior da Escada da Causalidade, o nível dos \textbf{contrafactuais}, porque para respondê-las precisamos voltar no tempo, mudar a história e perguntar: “O que teria acontecido se eu não tivesse tomado a aspirina?”. Nenhum experimento no mundo pode negar o tratamento a uma pessoa já tratada e comparar os dois resultados, de modo que precisamos importar um tipo inteiramente novo de conhecimento.

Os contrafactuais têm uma relação particularmente problemática com os dados porque os dados são, por definição, fatos. Eles não podem nos dizer o que acontecerá em um mundo contrafactual ou imaginário em que alguns fatos observados são bluntamente negados. E, no entanto, a mente humana faz tais inferências em busca de explicação de modo confiável e repetível. Eva o fez quando identificou “A serpente me enganou” como a razão de sua ação. Essa capacidade é o que mais distingue a inteligência humana da animal, bem como das versões cegas a modelos de IA e de aprendizado de máquina.

Você pode ser cético de que a ciência possa fazer qualquer afirmação útil sobre “teria sido”, mundos que não existem e coisas que não aconteceram. Mas ela faz, e sempre fez. As leis da física, por exemplo, podem ser interpretadas como asserções contrafactuais, tais como “Se o peso nesta mola tivesse dobrado, seu comprimento também teria dobrado” (lei de Hooke). Essa afirmação é, é claro, respaldada por uma riqueza de evidências experimentais (degrau dois), derivadas de centenas de molas, em dezenas de laboratórios, em milhares de ocasiões diferentes. No entanto, uma vez ungida como “lei”, os físicos a interpretam como uma relação funcional que governa \textit{esta} mola, \textit{neste} momento, sob valores hipotéticos do peso. Todos esses mundos diferentes, em que o peso é \textit{x} libras e o comprimento da mola é \textit{L}<sub>\textit{x}</sub> polegadas, são tratados como objetivamente cognoscíveis e simultaneamente ativos, embora apenas um deles de fato exista.

Voltando ao exemplo da pasta de dente, uma pergunta do degrau superior seria “Qual a probabilidade de que um cliente que comprou pasta de dente ainda a tivesse comprado se tivéssemos dobrado o preço?”. Estamos comparando o mundo real (em que sabemos que o cliente comprou a pasta de dente ao preço atual) com um mundo fictício (em que o preço é o dobro).

As recompensas de ter um modelo causal capaz de responder a perguntas contrafactuais são imensas. Descobrir por que um erro ocorreu nos permite tomar as medidas corretivas certas no futuro. Descobrir por que um tratamento funcionou em algumas pessoas e não em outras pode conduzir a uma nova cura para uma doença. Responder à pergunta “E se as coisas tivessem sido diferentes?” nos permite aprender com a história e com a experiência alheia, algo que nenhuma outra espécie parece fazer. Não é de surpreender que o antigo filósofo grego Demócrito (460–370 a.C.) tenha dito: “Preferiria descobrir uma causa a ser o rei da Pérsia.”

A posição dos contrafactuais no topo da Escada da Causalidade explica por que lhes atribuo tanto ênfase como momento-chave na evolução da consciência humana. Concordo totalmente com Yuval Harari em que a representação de criaturas imaginárias foi a manifestação de uma nova capacidade, a que ele chama de Revolução Cognitiva. Seu exemplo prototípico é a escultura do Homem-Leão, encontrada na caverna de Stadel, no sudoeste da Alemanha, e hoje guardada no Museu de Ulm (ver Figura 1.3).

\begin{figure}[ht]
\centering
\includegraphics[width=0.92\textwidth]{../figuras/fig_1_3.png}
\caption{Figura 1.3. O Homem-Leão da caverna de Stadel. A mais antiga representação conhecida de uma criatura imaginária (metade homem e metade leão), é emblemática de uma capacidade cognitiva recém-desenvolvida, a de raciocinar sobre contrafactuais. (Fonte: foto de Yvonne Mühleis, cortesia do State Office for Cultural Heritage Baden-Württemberg / Museu de Ulm, Ulm, Alemanha.)}
\end{figure}
\textit{Figura 1.3. O Homem-Leão da caverna de Stadel. A mais antiga representação conhecida de uma criatura imaginária (metade homem e metade leão), é emblemática de uma capacidade cognitiva recém-desenvolvida, a de raciocinar sobre contrafactuais. (Fonte: foto de Yvonne Mühleis, cortesia do State Office for Cultural Heritage Baden-Württemberg / Museu de Ulm, Ulm, Alemanha.)}

O Homem-Leão, com cerca de 40 mil anos, é uma presa de mamute esculpida na forma de uma quimera, metade homem e metade leão. Não sabemos quem esculpiu o Homem-Leão nem qual era seu propósito, mas sabemos que humanos anatomicamente modernos o fizeram e que ele representa uma ruptura com qualquer arte ou ofício que o tivesse precedido. Antes, os humanos haviam fabricado ferramentas e arte representacional, de contas a flautas, de pontas de lança a elegantes entalhes de cavalos e outros animais. O Homem-Leão é diferente: uma criatura de pura imaginação.

\textbf{Figura 1.3.} O Homem-Leão da caverna de Stadel. A mais antiga representação conhecida de uma criatura imaginária (metade homem e metade leão), é emblemática de uma capacidade cognitiva recém-desenvolvida, a de raciocinar sobre contrafactuais. (Fonte: foto de Yvonne Mühleis, cortesia do State Office for Cultural Heritage Baden-Württemberg / Museu de Ulm, Ulm, Alemanha.)

Como manifestação de nossa recém-encontrada capacidade de imaginar coisas que nunca existiram, o Homem-Leão é o precursor de toda teoria filosófica, descoberta científica e inovação tecnológica, de microscópios a aviões e computadores. Cada uma delas teve de tomar forma na imaginação de alguém antes de se realizar no mundo físico.

Esse salto à frente na capacidade cognitiva foi tão profundo e importante para nossa espécie quanto qualquer uma das mudanças anatômicas que nos tornaram humanos. Em 10 mil anos após a criação do Homem-Leão, todos os outros hominídeos (exceto os hominídeos de Flores, muito isolados geograficamente) haviam se extinguido. E os humanos continuaram a mudar o mundo natural com velocidade incrível, usando nossa imaginação para sobreviver, adaptar-nos e, por fim, tomar o controle.

A vantagem que obtivemos ao imaginar contrafactuais era a mesma então e o é hoje: flexibilidade, a capacidade de refletir e melhorar ações passadas e, talvez ainda mais significativo, nossa disposição de assumir responsabilidade por ações passadas e presentes.

Como mostra a Figura 1.2, as consultas características do terceiro degrau da Escada da Causalidade são “E se eu tivesse feito…?” e “Por quê?”. Ambas envolvem comparar o mundo observado com um mundo contrafactual. Experimentos sozinhos não podem responder a tais perguntas. Enquanto o degrau um lida com o mundo visto, e o degrau dois lida com um bravo novo mundo que é visível, o degrau três lida com um mundo que não pode ser visto (porque contradiz o que é visto). Para preencher a lacuna, precisamos de um modelo do processo causal subjacente, às vezes chamado de “teoria” ou até (nos casos em que temos confiança extraordinária) de “lei da natureza”. Em suma, precisamos de compreensão. Esse é, é claro, o Santo Graal de qualquer ramo da ciência — o desenvolvimento de uma teoria que nos permita predizer o que acontecerá em situações que sequer imaginamos ainda. Mas vai ainda mais longe: ter tais leis nos permite violá-las seletivamente de modo a criar mundos que contradizem o nosso. A próxima seção apresenta tais violações em ação.

\section{O mini-teste de Turing}
Em 1950, Alan Turing perguntou o que significaria para um computador pensar como um humano. Sugeriu um teste prático, a que chamou de “jogo da imitação”, mas todo pesquisador de IA desde então o chama de “teste de Turing”. Para todos os efeitos práticos, um computador poderia ser chamado de máquina pensante se um humano comum, comunicando-se com o computador por máquina de escrever, não pudesse dizer se estava falando com um humano ou com um computador. Turing estava muito confiante de que isso estava no âmbito do factível. “Acredito que em cerca de cinquenta anos será possível programar computadores”, escreveu ele, “de modo a fazê-los jogar o jogo da imitação tão bem que um interrogador médio não terá mais do que 70 por cento de chance de fazer a identificação correta após cinco minutos de questionamento.”

A predição de Turing foi um pouco imprecisa. Todo ano, a competição do Prêmio Loebner identifica o “chatbot” mais semelhante a um humano no mundo, com uma medalha de ouro e 100 mil dólares oferecidos a qualquer programa que consiga enganar todos os quatro juízes fazendo-os pensar que é humano. Até 2015, em vinte e cinco anos de competição, nenhum programa enganou todos os juízes, nem sequer metade deles.

Turing não apenas sugeriu o “jogo da imitação”; também propôs uma estratégia para passá-lo. “Em vez de tentar produzir um programa para simular a mente adulta, por que não tentar produzir um que simule a da criança?”, perguntou ele. Se se pudesse fazer isso, bastaria ensiná-lo do mesmo modo que se ensina uma criança e, voilà, vinte anos depois (ou menos, dada a maior velocidade de um computador), ter-se-ia uma inteligência artificial. “Presume-se que o cérebro da criança seja algo como um caderno como se compra na papelaria”, escreveu ele. “Bastante pouco mecanismo, e muitas folhas em branco.”

Ele estava errado nisso: o cérebro da criança é rico em mecanismos e em \textit{templates} pré-armazenados. Não obstante, penso que Turing estava no caminho certo. Provavelmente não teremos sucesso em criar inteligência de tipo humano até que possamos criar inteligência de tipo infantil, e um componente-chave dessa inteligência é o domínio da causalidade.

Como as máquinas podem adquirir conhecimento causal? Esse ainda é um grande desafio que, sem dúvida, envolverá uma combinação intricada de \textit{inputs} de experimentação ativa, observação passiva e (não menos importante) o programador — mais ou menos os mesmos \textit{inputs} que uma criança recebe, com a evolução, os pais e os pares no lugar do programador. No entanto, podemos responder a uma pergunta um pouco menos ambiciosa: Como as máquinas (e as pessoas) podem representar o conhecimento causal de um modo que lhes permita acessar a informação necessária com rapidez, responder corretamente às perguntas e fazê-lo com facilidade, como uma criança de três anos é capaz? De fato, essa é a pergunta principal que abordamos neste livro.

Chamo isso de mini-teste de Turing. A ideia é pegar uma história simples, codificá-la em uma máquina de alguma forma e, em seguida, testar se a máquina consegue responder corretamente a perguntas causais que um humano consegue responder. É “mini” por duas razões. Primeiro, confina-se ao raciocínio causal, excluindo outros aspectos da inteligência humana, como visão e linguagem natural. Segundo, permitimos que o concorrente codifique a história em qualquer representação conveniente, desonerando a máquina da tarefa de adquirir a história a partir de sua própria experiência pessoal. Passar nesse mini-teste tem sido o trabalho da minha vida — conscientemente nos últimos vinte e cinco anos e subconscientemente mesmo antes disso.

Obviamente, ao nos prepararmos para o mini-teste de Turing, a questão da representação precisa preceder a da aquisição. Sem uma representação, não saberíamos como armazenar a informação para uso futuro. Mesmo que pudéssemos deixar nosso robô manipular seu ambiente à vontade, qualquer informação aprendida desse modo seria esquecida, a menos que nosso robô estivesse dotado de um \textit{template} para codificar os resultados dessas manipulações. Uma contribuição importante da IA ao estudo da cognição tem sido o paradigma “Representação primeiro, aquisição depois”. Muitas vezes a busca por uma boa representação levou a \textit{insights} sobre como o conhecimento deveria ser adquirido, seja a partir de dados ou de um programador.

Quando descrevo o mini-teste de Turing, as pessoas comumente afirmam que ele pode ser facilmente derrotado por trapaça. Por exemplo, pegue a lista de todas as perguntas possíveis, armazene suas respostas corretas e, em seguida, leia-as da memória quando perguntado. Não há como distinguir (assim vai o argumento) entre uma máquina que armazena uma lista burra de perguntas e respostas e uma que responde do jeito que você e eu fazemos — isto é, compreendendo a pergunta e produzindo uma resposta por meio de um modelo mental causal. Então, o que o mini-teste de Turing provaria, se a trapaça é tão fácil?

O filósofo John Searle introduziu essa possibilidade de trapaça, conhecida como o argumento do “Quarto Chinês”, em 1980, para desafiar a afirmação de Turing de que a capacidade de fingir inteligência equivale a ter inteligência. O desafio de Searle tem apenas uma falha: a trapaça não é fácil; de fato, é impossível. Mesmo com um pequeno número de variáveis, o número de perguntas possíveis cresce astronomicamente. Digamos que tenhamos dez variáveis causais, cada uma das quais toma apenas dois valores (0 ou 1). Poderíamos fazer cerca de 30 milhões de consultas possíveis, tais como “Qual a probabilidade de que o resultado seja 1, dado que vemos a variável \textit{X} igual a 1 e fazemos a variável \textit{Y} igual a 0 e a variável \textit{Z} igual a 1?”.

Se houvesse mais variáveis, ou mais de dois estados para cada uma, o número de possibilidades cresceria além de nossa capacidade até de imaginar. A lista de Searle precisaria de mais entradas do que o número de átomos no universo. Assim, claramente, uma lista burra de perguntas e respostas nunca pode simular a inteligência de uma criança, que dirá de um adulto.

Os humanos devem ter alguma representação compacta da informação necessária em seus cérebros, bem como um procedimento eficaz para interpretar cada pergunta adequadamente e extrair a resposta certa da representação armazenada. Para passar no mini-teste de Turing, portanto, precisamos equipar as máquinas com uma representação e um algoritmo de extração de respostas similarmente eficientes.

Tal representação não apenas existe, como tem simplicidade infantil: um diagrama causal. Já vimos um exemplo, o diagrama da caçada ao mamute. Considerando a extrema facilidade com que as pessoas conseguem comunicar seu conhecimento com diagramas de pontos e setas, acredito que nossos cérebros de fato usam uma representação como essa. Mas, mais importante para nossos propósitos, esses modelos passam no mini-teste de Turing; nenhum outro modelo conhecido o faz. Vejamos alguns exemplos.

\textbf{Figura 1.4.} Diagrama causal para o exemplo do pelotão de fuzilamento. \textit{A} e \textit{B} representam (as ações de) os soldados \textit{A} e \textit{B}.

\begin{figure}[ht]
\centering
\includegraphics[width=0.92\textwidth]{../figuras/fig_1_4.png}
\caption{Figura 1.4. Diagrama causal para o exemplo do pelotão de fuzilamento. \textit{A} e \textit{B} representam (as ações de) os soldados \textit{A} e \textit{B}.}
\end{figure}
\textit{Figura 1.4. Diagrama causal para o exemplo do pelotão de fuzilamento. }A\textit{ e }B\textit{ representam (as ações de) os soldados }A\textit{ e }B\textit{.}

Suponha que um prisioneiro esteja prestes a ser executado por um pelotão de fuzilamento. Uma certa cadeia de eventos precisa ocorrer para que isso aconteça. Primeiro, o tribunal ordena a execução. A ordem vai a um capitão, que sinaliza aos soldados do pelotão de fuzilamento (\textit{A} e \textit{B}) que disparem. Vamos assumir que eles são obedientes e atiradores experientes, de modo que só disparam sob comando e, se qualquer um deles atira, o prisioneiro morre.

A Figura 1.4 mostra um diagrama representando a história que acabei de contar. Cada uma das incógnitas (\textit{CO}, \textit{C}, \textit{A}, \textit{B}, \textit{D}) é uma variável verdadeiro/falso. Por exemplo, \textit{D} = verdadeiro significa que o prisioneiro está morto; \textit{D} = falso significa que o prisioneiro está vivo. \textit{CO} = falso significa que a ordem do tribunal não foi emitida; \textit{CO} = verdadeiro significa que foi, e assim por diante.

Usando esse diagrama, podemos começar a responder a perguntas causais de diferentes degraus da escada. Primeiro, podemos responder a perguntas de associação (isto é, o que um fato nos diz sobre outro). Se o prisioneiro está morto, isso significa que a ordem do tribunal foi dada? Nós (ou um computador) podemos inspecionar o grafo, rastrear as regras por trás de cada uma das setas e, usando a lógica padrão, concluir que os dois soldados não teriam disparado sem o comando do capitão. Da mesma forma, o capitão não teria dado o comando se não tivesse a ordem em sua posse. Portanto, a resposta à nossa consulta é sim. Alternativamente, suponha que descobrimos que \textit{A} disparou. O que isso nos diz sobre \textit{B}? Seguindo as setas, o computador conclui que \textit{B} também deve ter disparado. (\textit{A} não teria disparado se o capitão não tivesse sinalizado, de modo que \textit{B} também deve ter disparado.) Isso é verdade mesmo que \textit{A} não cause \textit{B} (não há seta de \textit{A} para \textit{B}).

Subindo a Escada da Causalidade, podemos fazer perguntas sobre intervenção. E se o soldado \textit{A} decidir por iniciativa própria disparar, sem esperar o comando do capitão? O prisioneiro estará morto ou vivo? Essa pergunta de fato já tem um sabor contraditório. Acabei de lhe dizer que \textit{A} só atira se for comandado a isso, e no entanto agora estamos perguntando o que acontece se ele disparar sem um comando. Se você estiver apenas usando as regras da lógica, como os computadores tipicamente fazem, a pergunta é sem sentido. Como o robô da série de TV de ficção científica dos anos 1960 \textit{Perdidos no Espaço} costumava dizer em tais situações: “Isso não computa.”

Se quisermos que nosso computador compreenda a causalidade, temos de ensiná-lo a quebrar as regras. Temos de ensiná-lo a diferença entre meramente observar um evento e fazê-lo acontecer. “Sempre que você faz um evento acontecer”, dizemos ao computador, “remova todas as setas que apontam para esse evento e continue a análise pela lógica ordinária, como se as setas nunca tivessem estado ali.” Assim, apagamos todas as setas que apontam para a variável intervinda (\textit{A}). Também fixamos manualmente essa variável em seu valor prescrito (verdadeiro). A racionalidade dessa peculiar “cirurgia” é simples: fazer um evento acontecer significa emancipá-lo de todas as outras influências e submetê-lo a uma e somente uma influência — aquela que força sua ocorrência.

A Figura 1.5 mostra o diagrama causal que resulta de nosso exemplo. Essa intervenção conduz inevitavelmente à morte do prisioneiro. Essa é a função causal por trás da seta que vai de \textit{A} para \textit{D}.

\begin{figure}[ht]
\centering
\includegraphics[width=0.92\textwidth]{../figuras/fig_1_5.png}
\caption{Figura 1.5. mostra o diagrama causal que resulta de nosso exemplo. Essa intervenção conduz inevitavelmente à morte do prisioneiro. Essa é a função causal por trás da seta que vai de \textit{A} para \textit{D}.}
\end{figure}
\textit{Figura 1.5. mostra o diagrama causal que resulta de nosso exemplo. Essa intervenção conduz inevitavelmente à morte do prisioneiro. Essa é a função causal por trás da seta que vai de }A\textit{ para }D\textit{.}

\textbf{Figura 1.5.} Raciocínio sobre intervenções. O soldado \textit{A} decide disparar; a seta de \textit{C} para \textit{A} é apagada, e \textit{A} recebe o valor verdadeiro.

Note que essa conclusão concorda com nosso julgamento intuitivo de que o disparo não autorizado de \textit{A} conduzirá à morte do prisioneiro, porque a cirurgia deixa intacta a seta de \textit{A} para \textit{D}. Também, nosso julgamento seria de que \textit{B} (com toda a probabilidade) não atirou; nada na decisão de \textit{A} deveria afetar variáveis no modelo que não sejam efeitos do tiro de \textit{A}.

Isso merece ser repetido. Se \textit{vemos} \textit{A} atirar, então concluímos que \textit{B} também atirou. Mas se \textit{A} \textit{decide} atirar, ou se \textit{fazemos} \textit{A} atirar, então o oposto é verdade. Essa é a diferença entre ver e fazer. Apenas um computador capaz de captar essa diferença pode passar no mini-teste de Turing.

Note também que meramente coletar \textit{Big Data} não nos teria ajudado a subir a escada e responder às perguntas acima. Suponha que você seja um repórter coletando registros de cenas de execução dia após dia. Seus dados consistirão de dois tipos de eventos: ou todas as cinco variáveis são verdadeiras, ou todas são falsas. Não há como esse tipo de dado, na ausência de uma compreensão de quem escuta a quem, lhe permitir (ou a qualquer algoritmo de aprendizado de máquina) predizer os resultados de persuadir o atirador \textit{A} a não atirar.

Por fim, para ilustrar o terceiro degrau da Escada da Causalidade, formulemos uma pergunta contrafactual. Suponha que o prisioneiro esteja deitado morto no chão. Disso podemos concluir (usando o nível um) que \textit{A} atirou, \textit{B} atirou, o capitão deu o sinal e o tribunal deu a ordem. Mas e se \textit{A} tivesse decidido não atirar? O prisioneiro estaria vivo? Essa pergunta exige que comparemos o mundo real com um mundo fictício e contraditório em que \textit{A} não atirou. No mundo fictício, a seta que aponta para \textit{A} é apagada para liberar \textit{A} de escutar \textit{C}. Em vez disso, \textit{A} é fixado em falso, deixando sua história passada igual à do mundo real. Assim, o mundo fictício se parece com a Figura 1.6.

\begin{figure}[ht]
\centering
\includegraphics[width=0.92\textwidth]{../figuras/fig_1_6.png}
\caption{Figura 1.6. \textbf{Figura 1.6.} Raciocínio contrafactual. Observamos que o prisioneiro está morto e perguntamos o que teria acontecido se o soldado \textit{A} tivesse decidido não disparar.}
\end{figure}
\textbf{Figura 1.6.} Raciocínio contrafactual. Observamos que o prisioneiro está morto e perguntamos o que teria acontecido se o soldado \textit{A} tivesse decidido não disparar.

Para passar no mini-teste de Turing, nosso computador deve concluir que o prisioneiro estaria morto no mundo fictício também, porque o tiro de \textit{B} o teria matado. Assim, a corajosa mudança de coração de \textit{A} não teria salvado sua vida. Sem dúvida, essa é uma das razões pelas quais os pelotões de fuzilamento existem: eles garantem que a ordem do tribunal será cumprida e também aliviam parte do fardo da responsabilidade dos atiradores individuais, que podem dizer com a consciência (um tanto) limpa que suas ações não causaram a morte do prisioneiro, pois “ele teria morrido de qualquer jeito”.

Pode parecer que estamos nos dando um trabalho enorme para responder a perguntas-brinquedo cuja resposta já era óbvia de qualquer forma. Concordo plenamente! O raciocínio causal é fácil para você porque você é humano, e um dia foi uma criança de três anos, e tinha um maravilhoso cérebro de três anos que compreendia a causalidade melhor do que qualquer animal ou computador. Todo o ponto do “problema do mini-teste de Turing” é tornar o raciocínio causal factível também para computadores. No processo, talvez aprendamos algo sobre como os humanos o fazem.

Como os três exemplos mostram, temos de ensinar o computador a quebrar seletivamente as regras da lógica. Os computadores não são bons em quebrar regras, habilidade na qual as crianças se destacam. (Os homens das cavernas também! O Homem-Leão não poderia ter sido criado sem uma violação das regras sobre que cabeça combina com que corpo.) No entanto, não fiquemos demasiado complacentes quanto à superioridade humana. Os humanos podem ter bem mais dificuldade em chegar a conclusões causais corretas em um grande número de situações. Por exemplo, pode haver bem mais variáveis, e elas podem não ser simples variáveis binárias (verdadeiro/falso). Em vez de predizer se um prisioneiro está vivo ou morto, podemos querer predizer quanto a taxa de desemprego subirá se elevarmos o salário mínimo. Esse tipo de raciocínio causal quantitativo está em geral além do poder de nossa intuição. Também, no exemplo do pelotão de fuzilamento, descartamos incertezas: talvez o capitão tenha dado sua ordem uma fração de segundo depois que o fuzileiro \textit{A} decidiu atirar, talvez a arma do fuzileiro \textit{B} tenha emperrado, e assim por diante. Para lidar com a incerteza, precisamos de informação sobre a probabilidade de que tais anormalidades ocorram.

Deixe-me dar um exemplo em que as probabilidades fazem toda a diferença. Ele ecoa o debate público que irrompeu na Europa quando a vacina contra a varíola foi introduzida pela primeira vez. Inesperadamente, os dados mostraram que mais pessoas morriam das inoculações de varíola do que da própria varíola. Naturalmente, algumas pessoas usaram essa informação para argumentar que a inoculação deveria ser banida, quando de fato ela estava salvando vidas ao erradicar a varíola. Vejamos alguns dados fictícios para ilustrar o efeito e resolver a disputa.

Suponha que, de 1 milhão de crianças, 99 por cento sejam vacinadas e 1 por cento não. Se uma criança é vacinada, ela tem uma chance em cem de desenvolver uma reação, e a reação tem uma chance em cem de ser fatal. Por outro lado, ela não tem nenhuma chance de desenvolver varíola. Enquanto isso, se uma criança não é vacinada, ela obviamente tem chance zero de desenvolver uma reação à vacina, mas tem uma chance em cinquenta de desenvolver varíola. Por fim, vamos assumir que a varíola é fatal em um de cada cinco casos.

Penso que você concordaria que a vacinação parece uma boa ideia. As chances de ter uma reação são menores do que as de pegar varíola, e a reação é bem menos perigosa do que a doença. Mas agora vejamos os dados. De 1 milhão de crianças, 990 mil são vacinadas, 9.900 têm a reação e 99 morrem dela. Enquanto isso, 10 mil não são vacinadas, 200 pegam varíola e 40 morrem da doença. Em resumo, mais crianças morrem da vacinação (99) do que da doença (40).

Posso empatizar com os pais que poderiam marchar até o departamento de saúde com cartazes dizendo “Vacinas matam!”. E os dados parecem estar do lado deles; as vacinações de fato causam mais mortes do que a própria varíola. Mas a lógica está do lado deles? Deveríamos banir a vacinação ou levar em conta as mortes evitadas?

A Figura 1.7 mostra o diagrama causal para esse exemplo. Quando começamos, a taxa de vacinação era de 99 por cento. Agora fazemos a pergunta contrafactual “E se tivéssemos fixado a taxa de vacinação em zero?”. Usando as probabilidades que lhe dei acima, podemos concluir que, de 1 milhão de crianças, 20 mil teriam pegado varíola e 4 mil teriam morrido. Comparando o mundo contrafactual com o mundo real, vemos que não vacinar teria custado a vida de 3.861 crianças (a diferença entre 4.000 e 139). Devemos agradecer à linguagem dos contrafactuais por nos ajudar a evitar tais custos.

\begin{figure}[ht]
\centering
\includegraphics[width=0.92\textwidth]{../figuras/fig_1_7.png}
\caption{Figura 1.7. mostra o diagrama causal para esse exemplo. Quando começamos, a taxa de vacinação era de 99 por cento. Agora fazemos a pergunta contrafactual “E se tivéssemos fixado a taxa de vacinação em zero?”. Usando as probabilidades que lhe dei acima, podemos concluir que, de 1 milhão de crianças, 20 mil teriam pegado varíola e 4 mil teriam morrido. Comparando o mundo contrafactual com o mundo real, vemos que não vacinar teria custado a vida de 3.861 crianças (a diferença entre 4.000 e 139). Devemos agradecer à linguagem dos contrafactuais por nos ajudar a evitar tais custos.}
\end{figure}
\textit{Figura 1.7. mostra o diagrama causal para esse exemplo. Quando começamos, a taxa de vacinação era de 99 por cento. Agora fazemos a pergunta contrafactual “E se tivéssemos fixado a taxa de vacinação em zero?”. Usando as probabilidades que lhe dei acima, podemos concluir que, de 1 milhão de crianças, 20 mil teriam pegado varíola e 4 mil teriam morrido. Comparando o mundo contrafactual com o mundo real, vemos que não vacinar teria custado a vida de 3.861 crianças (a diferença entre 4.000 e 139). Devemos agradecer à linguagem dos contrafactuais por nos ajudar a evitar tais custos.}

\textbf{Figura 1.7.} Diagrama causal para o exemplo da vacinação. A vacinação é benéfica ou prejudicial?

A principal lição para um estudante de causalidade é que um modelo causal implica mais do que meramente desenhar setas. Por trás das setas, há probabilidades. Quando desenhamos uma seta de \textit{X} para \textit{Y}, estamos implicitamente dizendo que alguma regra ou função de probabilidade especifica como \textit{Y} mudaria se \textit{X} mudasse. Podemos saber qual é a regra; mais provavelmente, teremos de estimá-la a partir dos dados. Uma das características mais intrigantes da Revolução Causal, contudo, é que em muitos casos podemos deixar esses detalhes matemáticos completamente não especificados. Com muita frequência, a própria estrutura do diagrama nos permite estimar todo tipo de relações causais e contrafactuais: simples ou complicadas, determinísticas ou probabilísticas, lineares ou não lineares.

Da perspectiva computacional, nosso esquema para passar no mini-teste de Turing também é notável no sentido de que usamos a mesma rotina nos três exemplos: traduzir a história em um diagrama, escutar a consulta, realizar uma cirurgia que corresponde à consulta dada (intervencional ou contrafactual; se a consulta for associacional, nenhuma cirurgia é necessária) e usar o modelo causal modificado para computar a resposta. Não precisamos treinar a máquina em uma multidão de novas consultas cada vez que mudávamos a história. A abordagem é flexível o bastante para funcionar sempre que pudermos desenhar um diagrama causal, seja a respeito de mamutes, pelotões de fuzilamento ou vacinações. É exatamente isso o que queremos de um motor de inferência causal: é o tipo de flexibilidade de que desfrutamos como humanos.

É claro que não há nada inerentemente mágico em um diagrama. Ele tem êxito porque carrega informação causal; isto é, quando construímos o diagrama, perguntamos: “Quem poderia causar diretamente a morte do prisioneiro?” ou “Quais são os efeitos diretos das vacinações?”. Se tivéssemos construído o diagrama perguntando sobre meras associações, ele não nos teria dado essas capacidades. Por exemplo, na Figura 1.7, se invertêssemos a seta Vacinação → Varíola, obteríamos as mesmas associações nos dados, mas concluiríamos erroneamente que a varíola afeta a vacinação.

Décadas de experiência com esse tipo de pergunta me convenceram de que, tanto em sentido cognitivo quanto filosófico, a ideia de causas e efeitos é muito mais fundamental do que a ideia de probabilidade. Começamos a aprender causas e efeitos antes de compreender a linguagem e antes de saber qualquer matemática. (Pesquisas mostraram que crianças de três anos já compreendem a Escada da Causalidade inteira.) Da mesma forma, o conhecimento transmitido em um diagrama causal é tipicamente muito mais robusto do que o codificado em uma distribuição de probabilidade. Por exemplo, suponha que os tempos mudaram e uma vacina bem mais segura e eficaz é introduzida. Suponha, ainda, que, graças a melhores condições de higiene e socioeconômicas, o perigo de contrair varíola diminuiu. Essas mudanças afetarão drasticamente todas as probabilidades envolvidas; e, no entanto, notavelmente, a estrutura do diagrama permanecerá invariante. Esse é o segredo-chave da modelagem causal. Mais ainda, uma vez que passamos pela análise e descobrimos como estimar o benefício da vacinação a partir dos dados, não precisamos repetir toda a análise do zero. Como discutido na Introdução, o mesmo estimando (isto é, a receita para responder à consulta) permanecerá válido e, enquanto o diagrama não mudar, poderá ser aplicado aos novos dados e produzir uma nova estimativa para nossa consulta. É por causa dessa robustez, conjecturo, que a intuição humana se organiza em torno de relações causais, e não estatísticas.

\section{Sobre probabilidades e causalidade}
O reconhecimento de que a causalidade não se reduz a probabilidades foi muito arduamente conquistado, tanto para mim pessoalmente quanto para filósofos e cientistas em geral. Compreender o significado de “causa” tem sido o foco de uma longa tradição de filósofos, de David Hume e John Stuart Mill nos séculos XVIII e XIX a Hans Reichenbach e Patrick Suppes em meados do século XX, a Nancy Cartwright, Wolfgang Spohn e Christopher Hitchcock hoje. Em particular, a partir de Reichenbach e Suppes, os filósofos tentaram definir a causalidade em termos de probabilidade, usando a noção de “elevação de probabilidade”: \textit{X} causa \textit{Y} se \textit{X} eleva a probabilidade de \textit{Y}.

Esse conceito está solidamente encravado na intuição. Dizemos, por exemplo, “A direção imprudente causa acidentes” ou “Você vai reprovar neste curso por causa de sua preguiça”, sabendo perfeitamente bem que os antecedentes meramente tendem a tornar as consequências mais prováveis, não absolutamente certas.

Esperar-se-ia, portanto, que a elevação de probabilidade se tornasse a ponte entre o degrau um e o degrau dois da Escada da Causalidade. Ai de nós, essa intuição conduziu a décadas de tentativas fracassadas.

O que impediu as tentativas de ter êxito não foi a ideia em si, mas o modo como foi articulada formalmente. Quase sem exceção, os filósofos expressaram a sentença “\textit{X} eleva a probabilidade de \textit{Y}” usando probabilidades condicionais e escreveram \textit{P}(\textit{Y} | \textit{X}) > \textit{P}(\textit{Y}). Essa interpretação é errada, como você certamente notou, porque “eleva” é um conceito causal, conotando uma influência causal de \textit{X} sobre \textit{Y}. A expressão \textit{P}(\textit{Y} | \textit{X}) > \textit{P}(\textit{Y}), por outro lado, fala apenas de observações e significa: “Se vemos \textit{X}, então a probabilidade de \textit{Y} aumenta.” Mas esse aumento pode ocorrer por outras razões, incluindo \textit{Y} ser uma causa de \textit{X} ou alguma outra variável (\textit{Z}) ser a causa de ambas. Essa é a pegadinha! Ela devolve os filósofos à estaca zero, tentando eliminar essas “outras razões”.

As probabilidades, como dadas por expressões como \textit{P}(\textit{Y} | \textit{X}), situam-se no primeiro degrau da Escada da Causalidade e nunca podem (por si sós) responder a consultas no segundo ou no terceiro degrau. Qualquer tentativa de “definir” a causalidade em termos de conceitos aparentemente mais simples, do primeiro degrau, deve fracassar. É por isso que não tentei definir a causalidade em nenhum lugar deste livro: definições exigem redução, e redução exige descer a um degrau inferior. Em vez disso, persegui o programa em última instância mais construtivo de explicar como responder a consultas causais e que informação é necessária para respondê-las. Se isso parecer estranho, considere que os matemáticos adotam exatamente a mesma abordagem em relação à geometria euclidiana. Em nenhum lugar de um livro de geometria você encontrará uma definição dos termos “ponto” e “reta”. E, no entanto, podemos responder a qualquer e a todas as consultas a respeito deles com base nos axiomas de Euclides (ou, ainda melhor, nas várias versões modernas dos axiomas de Euclides).

Mas vamos olhar com mais cuidado para esse critério de elevação de probabilidade e ver onde ele encalha. A questão de uma causa comum, ou confundidor, de \textit{X} e \textit{Y} foi das mais vexatórias para os filósofos. Se tomarmos o critério de elevação de probabilidade ao pé da letra, devemos concluir que as altas vendas de sorvete causam crime porque a probabilidade de crime é maior nos meses em que mais sorvete é vendido. Nesse caso particular, podemos explicar o fenômeno porque tanto as vendas de sorvete quanto o crime são maiores no verão, quando o tempo está mais quente. Não obstante, ainda nos resta perguntar que critério filosófico geral poderia nos dizer que o tempo, e não as vendas de sorvete, é a causa.

Os filósofos se esforçaram arduamente para reparar a definição condicionando o que chamavam de “fatores de fundo” (outra palavra para confundidores), resultando no critério \textit{P}(\textit{Y} | \textit{X}, \textit{K} = \textit{k}) > \textit{P}(\textit{Y} | \textit{K} = \textit{k}), em que \textit{K} representa algumas variáveis de fundo. De fato, esse critério funciona para nosso exemplo do sorvete se tratarmos a temperatura como uma variável de fundo. Por exemplo, se olharmos apenas para dias em que a temperatura é de noventa graus Fahrenheit (\textit{K} = 90), não encontraremos associação residual entre vendas de sorvete e crime. Só quando comparamos dias de noventa graus com dias de trinta graus obtemos a ilusão de uma elevação de probabilidade. Ainda assim, nenhum filósofo foi capaz de dar uma resposta convincentemente geral à pergunta “Quais variáveis precisam ser incluídas no conjunto de fundo \textit{K} e condicionadas?”. A razão é óbvia: o confundimento também é um conceito causal e, portanto, desafia a formulação probabilística.

Em 1983, Nancy Cartwright rompeu esse impasse e enriqueceu a descrição do contexto de fundo com um componente causal. Ela propôs que deveríamos condicionar qualquer fator que seja “causalmente relevante” para o efeito. Ao tomar emprestado um conceito do degrau dois da Escada da Causalidade, ela essencialmente desistiu da ideia de definir causas com base apenas na probabilidade. Isso foi progresso, mas abriu a porta à crítica de que estamos definindo uma causa em termos de si mesma.

Disputas filosóficas sobre o conteúdo apropriado de \textit{K} continuaram por mais de duas décadas e chegaram a um beco sem saída. De fato, veremos um critério correto no Capítulo 4, e não vou estragar a surpresa aqui. Basta, por ora, dizer que esse critério é praticamente impossível de enunciar sem diagramas causais.

Em resumo, a causalidade probabilística sempre naufragou no escolho do confundimento. Cada vez que os adeptos da causalidade probabilística tentam remendar o navio com um novo casco, o barco bate no mesmo rochedo e abre outro vazamento. Uma vez que se deturpa a “elevação de probabilidade” na linguagem das probabilidades condicionais, nenhuma quantidade de remendo probabilístico o levará ao próximo degrau da escada. Por mais estranho que possa soar, a noção de elevação de probabilidade não pode ser expressa em termos de probabilidades.

O modo adequado de resgatar a ideia de elevação de probabilidade é com o operador \textit{do}: podemos dizer que \textit{X} causa \textit{Y} se \textit{P}(\textit{Y} | \textit{do}(\textit{X})) > \textit{P}(\textit{Y}). Como a intervenção é um conceito do degrau dois, essa definição pode capturar a interpretação causal da elevação de probabilidade, e também pode ser tornada operacional por meio de diagramas causais. Em outras palavras, se tivermos um diagrama causal e dados em mãos e um pesquisador perguntar se \textit{P}(\textit{Y} | \textit{do}(\textit{X})) > \textit{P}(\textit{Y}), podemos responder à sua pergunta de modo coerente e algorítmico e, assim, decidir se \textit{X} é uma causa de \textit{Y} no sentido de elevação de probabilidade.

Normalmente presto grande atenção ao que os filósofos têm a dizer sobre conceitos escorregadios como causalidade, indução e a lógica da inferência científica. Os filósofos têm a vantagem de se manterem à parte da azáfama do debate científico e das realidades práticas de lidar com dados. Ficaram menos contaminados do que outros cientistas pelos vieses anticausais da estatística. Podem invocar uma tradição de pensamento sobre a causalidade que remonta pelo menos a Aristóteles, e podem falar de causalidade sem corar nem escondê-la sob o rótulo de “associação”.

No entanto, em seu esforço para matematizar o conceito de causalidade — em si uma ideia louvável —, os filósofos foram rápidos demais em se comprometer com a única linguagem de tratamento da incerteza que conheciam, a linguagem da probabilidade. Em grande medida, superaram esse equívoco na última década ou mais, mas, infelizmente, ideias semelhantes ainda estão sendo perseguidas na econometria mesmo agora, sob nomes como “causalidade de Granger” e “autocorrelação vetorial”.

Agora tenho uma confissão a fazer: cometi o mesmo erro. Nem sempre coloquei a causalidade em primeiro lugar e a probabilidade em segundo. Muito pelo contrário! Quando comecei a trabalhar em inteligência artificial, no início dos anos 1980, pensava que a incerteza era a coisa mais importante que faltava à IA. Mais ainda, insistia que a incerteza fosse representada por probabilidades. Assim, como explico no Capítulo 3, desenvolvi uma abordagem ao raciocínio sob incerteza, chamada redes bayesianas, que imita o modo como um cérebro idealizado e descentralizado poderia incorporar probabilidades em suas decisões. Dado que vemos certos fatos, as redes bayesianas podem computar rapidamente a probabilidade de que certos outros fatos sejam verdadeiros ou falsos. Não surpreendentemente, as redes bayesianas pegaram imediatamente na comunidade de IA e ainda hoje são consideradas um paradigma líder em inteligência artificial para o raciocínio sob incerteza.

Embora eu me delicie com o sucesso contínuo das redes bayesianas, elas não conseguiram preencher a lacuna entre a inteligência artificial e a humana. Tenho certeza de que você consegue descobrir o ingrediente que faltava: a causalidade. É verdade, fantasmas causais estavam por toda parte. As setas invariavelmente apontavam de causas para efeitos, e os praticantes frequentemente notavam que sistemas de diagnóstico se tornavam ingovernáveis quando a direção das setas era invertida. Mas, em grande medida, pensávamos que isso era um hábito cultural, ou um artefato de velhos padrões de pensamento, e não um aspecto central do comportamento inteligente. Na época, eu estava tão intoxicado com o poder das probabilidades que considerava a causalidade um conceito subserviente, mera conveniência ou abreviatura mental para expressar dependências probabilísticas e distinguir variáveis relevantes das irrelevantes. Em meu livro de 1988, \textit{Probabilistic Reasoning in Intelligent Systems}, escrevi: “A causação é uma linguagem com a qual se pode falar eficientemente sobre certas estruturas de relações de relevância.” As palavras me envergonham hoje, porque “relevância” é tão obviamente uma noção do degrau um.

Mesmo na época em que o livro foi publicado, eu sabia no fundo do coração que estava errado. Para meus colegas cientistas da computação, meu livro se tornou a bíblia do raciocínio sob incerteza, mas eu já me sentia um apóstata. As redes bayesianas habitam um mundo em que todas as perguntas se reduzem a probabilidades, ou (na terminologia deste capítulo) graus de associação entre variáveis; elas não podiam subir ao segundo ou ao terceiro degrau da Escada da Causalidade. Felizmente, precisaram apenas de duas pequenas reviravoltas para subir ao topo. Primeiro, em 1991, a ideia da cirurgia de grafo as capacitou a lidar tanto com observações quanto com intervenções. Outra reviravolta, em 1994, as levou ao terceiro nível e as tornou capazes de lidar com contrafactuais. Mas esses desenvolvimentos merecem uma discussão mais completa em um capítulo posterior.

O ponto principal é este: enquanto as probabilidades codificam nossas crenças sobre um mundo estático, a causalidade nos diz se e como as probabilidades mudam quando o mundo muda, seja por intervenção ou por ato de imaginação.

\textit{Sir Francis Galton demonstra sua “tábua de Galton” ou “quincunce” na Royal Institution. Ele via esse aparelho semelhante a um pinball como uma analogia para a herança de traços genéticos como a estatura. As bolinhas se acumulam em uma curva em forma de sino semelhante à distribuição das alturas humanas. O enigma de por que as alturas humanas não se espalham de uma geração para a seguinte, como as bolinhas fariam, levou-o à descoberta da “regressão à média”. (Fonte: desenho de Dakota Harr.)}

\chapter{Capítulo 2 — Dos bucaneiros às cobaias: a gênese da inferência causal}
E, no entanto, ela se move. — ATRIBUÍDO A GALILEU GALILEI (1564–1642)

POR quase dois séculos, um dos rituais mais duradouros da ciência britânica tem sido o \textit{Friday Evening Discourse} na Royal Institution of Great Britain, em Londres. Muitas descobertas do século XIX foram anunciadas pela primeira vez ao público nesse local: Michael Faraday e os princípios da fotografia em 1839; J. J. Thomson e o elétron em 1897; James Dewar e a liquefação do hidrogênio em 1904. O aparato cênico era parte importante da ocasião; era, literalmente, a ciência como teatro, e a plateia — a nata da sociedade britânica — vestia-se a rigor (tuxedos com gravata preta para os homens). Na hora marcada, um sino soava, e o orador da noite era conduzido ao auditório. Por tradição, ele começava a palestra imediatamente, sem introdução nem preâmbulo. Experimentos e demonstrações ao vivo faziam parte do espetáculo.

Em 9 de fevereiro de 1877, o orador da noite era Francis Galton, FRS, primo em primeiro grau de Charles Darwin, notável explorador da África, inventor da identificação por impressões digitais e o próprio modelo do cientista-cavalheiro vitoriano. O título de Galton era “Leis típicas da hereditariedade”. Seu aparato experimental da noite era um engenhoso contraptão a que chamava de \textit{quincunx}, hoje frequentemente conhecido como “tábua de Galton”. Um jogo semelhante aparece com frequência no \textit{game show} televisivo \textit{The Price Is Right}, onde é conhecido como Plinko.

A tábua de Galton consiste em um arranjo triangular de pinos ou cavilhas, no qual pequenas bolinhas de metal podem ser inseridas por uma abertura no topo. As bolinhas saltam para baixo de uma fileira à seguinte, no estilo de um \textit{pinball}, antes de se acomodarem em um dos encaixes alinhados na base (ver frontispício). Para qualquer bolinha individual, os zigue-zagues para a esquerda ou para a direita parecem completamente aleatórios. No entanto, se despejarmos um grande número de bolinhas na tábua de Galton, emerge uma regularidade assombrosa: as bolinhas acumuladas na base formarão sempre uma aproximação grosseira de uma curva em forma de sino. Os encaixes mais próximos do centro estarão empilhados de bolinhas, e o número de bolinhas em cada encaixe diminui gradualmente até zero nas bordas do \textit{quincunx}.

Esse padrão tem uma explicação matemática. A trajetória de qualquer bolinha individual assemelha-se a uma sequência de lançamentos independentes de moeda. Cada vez que uma bolinha atinge um pino, ela salta ou para a esquerda ou para a direita, e de longe sua escolha parece completamente aleatória. A soma dos resultados — digamos, o excesso de saltos à direita sobre os da esquerda — determina em qual encaixe a bolinha termina. Segundo o teorema do limite central, demonstrado em 1810 por Pierre-Simon Laplace, qualquer processo aleatório desse tipo — um que se reduz à soma de um grande número de lançamentos de moeda — conduzirá à mesma distribuição de probabilidade, chamada distribuição normal (ou curva em forma de sino). A tábua de Galton é simplesmente uma demonstração visual do teorema de Laplace.

O teorema do limite central é verdadeiramente um milagre da matemática do século XIX. Pense nisso: ainda que a trajetória de qualquer bolinha individual seja imprevisível, a trajetória de 1.000 bolinhas é extremamente previsível — um fato conveniente para os produtores de \textit{The Price Is Right}, que podem estimar com precisão quanto dinheiro os concorrentes ganharão no Plinko a longo prazo. É a mesma lei que torna as seguradoras tão lucrativas, a despeito das incertezas dos assuntos humanos.

A plateia bem-vestida da Royal Institution deve ter se perguntado o que tudo isso tinha a ver com as leis da hereditariedade, o tema prometido da palestra. Para ilustrar a conexão, Galton mostrou-lhes alguns dados coletados na França sobre as estaturas de recrutas militares. Estes também seguem uma distribuição normal: muitos homens têm estatura próxima da média, com um número gradualmente decrescente de indivíduos extremamente altos ou extremamente baixos. De fato, não importa se estamos falando de 1.000 recrutas militares ou de 1.000 bolinhas na tábua de Galton: os números em cada encaixe (ou categoria de estatura) são quase os mesmos. Assim, para Galton, o \textit{quincunx} era um modelo para a herança da estatura ou, de fato, de muitos outros traços genéticos. É um modelo causal. Em termos mais simples, Galton acreditava que as bolinhas “herdavam” sua posição no \textit{quincunx} do mesmo modo que os humanos herdam sua estatura.

Mas se aceitarmos esse modelo — provisoriamente —, ele coloca um enigma, que foi o principal tema de Galton naquela noite. A largura da curva em forma de sino depende do número de fileiras de pinos colocadas entre o topo e a base. Suponhamos que dobremos o número de fileiras. Isso criaria um modelo para duas gerações de herança, com a primeira metade das fileiras representando a primeira geração e a segunda metade representando a segunda. Inevitavelmente, encontraríamos mais variação na segunda geração do que na primeira e, nas gerações sucessivas, a curva em forma de sino se tornaria cada vez mais larga. Mas não é isso o que acontece com a estatura humana real. De fato, a largura da distribuição das estaturas humanas permanece relativamente constante ao longo do tempo. Não tínhamos humanos de quase três metros um século atrás, e ainda não temos. O que explica a estabilidade da dotação genética da população?

Galton vinha se debruçando sobre esse enigma por cerca de oito anos, desde a publicação de seu livro \textit{Hereditary Genius} em 1869. Como o título do livro sugere, o verdadeiro interesse de Galton não eram jogos de feira nem a estatura humana, e sim a inteligência humana. Como membro de uma família extensa com uma quantidade notável de genialidade científica, Galton naturalmente teria gostado de provar que o gênio corre nas famílias. E partira para fazer exatamente isso em seu livro. Compilou, com meticulosidade, as genealogias de 605 ingleses “eminentes” dos quatro séculos precedentes. Mas descobriu que os filhos e os pais desses homens eminentes eram um tanto menos eminentes, e os avós e netos ainda menos eminentes.

É fácil o bastante, agora, encontrar falhas no programa de Galton. Afinal, o que é a definição de eminência? E não seria possível que as pessoas de famílias eminentes sejam bem-sucedidas por causa de seu privilégio, e não de seu talento? Embora Galton estivesse ciente de tais dificuldades, ele perseguiu essa busca fútil por uma explicação genética com ritmo e determinação crescentes.

Ainda assim, Galton estava sobre algo, o que se tornou mais aparente quando começou a examinar características como a estatura, que são mais fáceis de medir e estão mais fortemente ligadas à hereditariedade do que a “eminência”. Os filhos de homens altos tendem a ser mais altos que a média — mas não tão altos quanto seus pais. Os filhos de homens baixos tendem a ser mais baixos que a média — mas não tão baixos quanto seus pais. Galton chamou primeiro esse fenômeno de “reversão” e, mais tarde, de “regressão à mediocridade”. Ele pode ser observado em muitos outros contextos. Se estudantes fazem dois testes padronizados diferentes sobre o mesmo material, aqueles que pontuaram alto no primeiro teste geralmente pontuarão acima da média no segundo, mas não tão alto quanto na primeira vez. Esse fenômeno de regressão à média é onipresente em todas as facetas da vida, da educação e dos negócios. Por exemplo, no beisebol, o Novato do Ano (um jogador que se sai inesperadamente bem em sua primeira temporada) frequentemente sofre um “declínio de segundo ano”, em que não se sai tão bem.

Galton não sabia de tudo isso, e pensava ter tropeçado em uma lei da hereditariedade, e não em uma lei da estatística. Acreditava que a regressão à média deveria ter alguma causa, e em sua palestra na Royal Institution ilustrava seu ponto. Mostrou à plateia um \textit{quincunx} de duas camadas (Figura 2.1). Depois de passar pelo primeiro arranjo de pinos, as bolinhas passavam por calhas inclinadas que as moviam para mais perto do centro da tábua. Em seguida, passavam por um segundo arranjo de pinos. Galton mostrou triunfalmente que as calhas compensavam exatamente a tendência da distribuição normal de se espalhar. Dessa vez, a distribuição de probabilidade em forma de sino mantinha uma largura constante de geração em geração.

\begin{figure}[ht]
\centering
\includegraphics[width=0.92\textwidth]{../figuras/fig_2_1.png}
\caption{Figura 2.1. O tabuleiro de Galton, usado por Francis Galton como analogia para a herança da estatura humana. (a) Quando muitas bolinhas caem pelo aparato semelhante a um \textit{pinball}, seus saltos aleatórios fazem com que se acumulem em uma curva em forma de sino. (b) Galton notou que, em duas passagens, A e B, pelo tabuleiro (análogo a duas gerações), a curva se alargava. (c) Para contrabalançar essa tendência, instalou calhas que deslocavam a “segunda geração” de volta para o centro. As calhas são a explicação causal de Galton para a regressão à média. (Fonte: Francis Galton, \textit{Natural Inheritance} [1889].)}
\end{figure}
\textit{Figura 2.1. O tabuleiro de Galton, usado por Francis Galton como analogia para a herança da estatura humana. (a) Quando muitas bolinhas caem pelo aparato semelhante a um }pinball\textit{, seus saltos aleatórios fazem com que se acumulem em uma curva em forma de sino. (b) Galton notou que, em duas passagens, A e B, pelo tabuleiro (análogo a duas gerações), a curva se alargava. (c) Para contrabalançar essa tendência, instalou calhas que deslocavam a “segunda geração” de volta para o centro. As calhas são a explicação causal de Galton para a regressão à média. (Fonte: Francis Galton, }Natural Inheritance\textit{ [1889].)}

\textbf{Figura 2.1.} A tábua de Galton, usada por Francis Galton como analogia para a herança das estaturas humanas. (a) Quando muitas bolinhas são deixadas cair pelo aparato semelhante a um \textit{pinball}, seus saltos aleatórios fazem com que se acumulem em uma curva em forma de sino. (b) Galton observou que em duas passagens, A e B, pela tábua de Galton (o análogo de duas gerações) a curva em forma de sino se tornava mais larga. (c) Para contrabalançar essa tendência, instalou calhas para mover a “segunda geração” de volta para mais perto do centro. As calhas são a explicação causal de Galton para a regressão à média. (Fonte: Francis Galton, \textit{Natural Inheritance} [1889].)

Assim, Galton conjecturou, a regressão à média era um processo físico, o modo da natureza de assegurar que a distribuição da estatura (ou da inteligência) permanecesse a mesma de geração em geração. “O processo de reversão coopera com a lei geral do desvio”, disse Galton à plateia. Comparou-o à lei de Hooke, a lei física que descreve a tendência de uma mola a retornar ao seu comprimento de equilíbrio.

Tenha em mente a data. Em 1877, Galton perseguia uma explicação causal e pensava que a regressão à média era um processo causal, como uma lei da física. Estava enganado, mas estava longe de estar sozinho. Muitas pessoas continuam a cometer o mesmo erro até hoje. Por exemplo, especialistas em beisebol sempre buscam explicações causais para o declínio de segundo ano de um jogador. “Ele ficou confiante demais”, reclamam, ou “os outros jogadores descobriram suas fraquezas”. Podem estar certos, mas o declínio de segundo ano não precisa de uma explicação causal. Acontecerá mais vezes do que não apenas pelas leis do acaso. A explicação estatística moderna é bastante simples. Como Daniel Kahneman a resume em seu livro \textit{Thinking, Fast and Slow}: “Sucesso = talento + sorte. Grande sucesso = um pouco mais de talento + muita sorte.” Um jogador que vence o prêmio de Novato do Ano é provavelmente mais talentoso que a média, mas também (provavelmente) teve muita sorte. Na temporada seguinte, não é provável que seja tão sortudo, e sua média de rebatidas será menor.

Em 1889, Galton havia compreendido isso e, no processo — em parte decepcionado, mas também fascinado —, deu o primeiro enorme passo no sentido de divorciar a estatística da causalidade. Seu raciocínio é sutil, mas vale o esforço de compreendê-lo. É a disciplina recém-nascida da estatística emitindo seu primeiro choro.

Galton começara a reunir uma variedade de estatísticas “antropométricas”: estatura, comprimento do antebraço, comprimento da cabeça, largura da cabeça, e assim por diante. Notou que, quando traçava a estatura contra o comprimento do antebraço, por exemplo, o mesmo fenômeno de regressão à média ocorria. Homens altos geralmente tinham antebraços mais longos que a média — mas não tão acima da média quanto sua estatura. Claramente, a estatura não é uma causa do comprimento do antebraço, nem vice-versa. Se algo, ambos são causados pela herança genética. Galton começou a usar uma nova palavra para esse tipo de relação: a estatura e o comprimento do antebraço estavam “co-relacionados”. Eventualmente, optou pela palavra inglesa mais usual “correlated” [correlacionados].

Mais tarde, percebeu um fato ainda mais assombroso: nas comparações geracionais, a ordem temporal podia ser invertida. Isto é, os pais dos filhos também revertem à média. O pai de um filho mais alto que a média tem probabilidade de ser mais alto que a média, mas mais baixo que seu filho (ver Figura 2.2). Uma vez que Galton compreendeu isso, teve de abandonar qualquer ideia de uma explicação causal para a regressão, porque não há como as estaturas dos filhos causarem as estaturas dos pais.

\begin{figure}[ht]
\centering
\includegraphics[width=0.92\textwidth]{../figuras/fig_2_2.png}
\caption{Figura 2.2. O diagrama de dispersão mostra um conjunto de dados de estaturas, em que cada ponto representa a altura do pai (eixo \textit{x}) e do filho (eixo \textit{y}). A linha tracejada coincide com o eixo maior da elipse, enquanto a linha contínua (chamada reta de regressão) conecta os pontos extremos da elipse. A diferença entre elas responde pela regressão à média. Por exemplo, a estrela preta mostra que pais de 72″ têm, em média, filhos de 71″. A faixa horizontal e a estrela branca mostram que a mesma perda de estatura ocorre na direção não causal (para trás no tempo). (Fonte: figura de Maayan Harel, com contribuição de Christopher Boucher.)}
\end{figure}
\textit{Figura 2.2. O diagrama de dispersão mostra um conjunto de dados de estaturas, em que cada ponto representa a altura do pai (eixo }x\textit{) e do filho (eixo }y\textit{). A linha tracejada coincide com o eixo maior da elipse, enquanto a linha contínua (chamada reta de regressão) conecta os pontos extremos da elipse. A diferença entre elas responde pela regressão à média. Por exemplo, a estrela preta mostra que pais de 72″ têm, em média, filhos de 71″. A faixa horizontal e a estrela branca mostram que a mesma perda de estatura ocorre na direção não causal (para trás no tempo). (Fonte: figura de Maayan Harel, com contribuição de Christopher Boucher.)}

Essa percepção pode soar paradoxal à primeira vista. “Espere!”, você está dizendo. “Você está me dizendo que pais altos geralmente têm filhos mais baixos, e filhos altos geralmente têm pais mais baixos. Como ambas as afirmações podem ser verdadeiras? Como um filho pode ser ao mesmo tempo mais alto e mais baixo que seu pai?”

\textbf{Figura 2.2.} O diagrama de dispersão mostra um conjunto de dados de estaturas, com cada ponto representando a estatura de um pai (no eixo \textit{x}) e de seu filho (no eixo \textit{y}). A reta tracejada coincide com o eixo maior da elipse, enquanto a reta sólida (chamada reta de regressão) conecta os pontos mais à direita e mais à esquerda da elipse. A diferença entre elas explica a regressão à média. Por exemplo, a estrela preta mostra que pais de 72″ têm, em média, filhos de 71″. (Isto é, a estatura média de todos os pontos de dados na faixa vertical é 71″.) A faixa horizontal e a estrela branca mostram que a mesma perda de estatura ocorre na direção não causal (para trás no tempo). (Fonte: figura de Maayan Harel, com contribuição de Christopher Boucher.)

A resposta é que não estamos falando de um pai individual e de um filho individual, e sim de duas populações. Começamos com a população de pais de um metro e oitenta e três. Como eles são mais altos que a média, seus filhos regredirão à média; digamos que seus filhos mediamente meçam um metro e oitenta. No entanto, a população de pares pai-filho com pais de um metro e oitenta e três não é a mesma que a população de pares pai-filho com filhos de um metro e oitenta. Todo pai no primeiro grupo tem, por definição, um metro e oitenta e três. Mas o segundo grupo terá alguns pais mais altos que um metro e oitenta e três e muitos pais mais baixos que um metro e oitenta e três. Sua estatura média será mais baixa que um metro e oitenta, exibindo novamente a regressão à média.

Outro modo de ilustrar a regressão é usar um diagrama chamado diagrama de dispersão (Figura 2.2). Cada par pai-filho é representado por um ponto, com a coordenada \textit{x} sendo a estatura do pai e a coordenada \textit{y} sendo a estatura do filho. Assim, um pai e um filho que ambos medem um metro e setenta e cinco (ou sessenta e nove polegadas) serão representados por um ponto em (69, 69), bem no centro do diagrama de dispersão. Um pai que mede um metro e oitenta e três (ou setenta e duas polegadas) com um filho de um metro e oitenta (ou setenta e uma polegadas) será representado por um ponto em (72, 71), no canto nordeste do diagrama de dispersão.

Note que o diagrama de dispersão tem uma forma aproximadamente elíptica — um fato que foi crucial para a análise de Galton e é característico de distribuições em forma de sino com duas variáveis. Como mostrado na Figura 2.2, os pares pai-filho com pais de setenta e duas polegadas situam-se em uma fatia vertical centrada em 72; os pares pai-filho com filhos de setenta e uma polegadas situam-se em uma fatia horizontal centrada em 71. Eis a prova visual de que se trata de duas populações diferentes. Se nos concentrarmos apenas na primeira população, os pares com pais de setenta e duas polegadas, podemos perguntar: “Qual a estatura média dos filhos?” É o mesmo que perguntar onde está o centro daquela fatia vertical, e a olho se vê que o centro está em cerca de 71. Se nos concentrarmos apenas na segunda população, com filhos de setenta e uma polegadas, podemos perguntar: “Qual a estatura média dos pais?” Isso é o mesmo que perguntar o centro da fatia horizontal, e a olho se vê que seu centro está em cerca de 70,3.

Podemos ir mais longe e pensar em fazer o mesmo procedimento para cada fatia vertical. Isso é equivalente a perguntar: “Para pais de estatura \textit{x}, qual a melhor predição da estatura do filho (\textit{y})?” Alternativamente, podemos tomar cada fatia horizontal e perguntar onde está seu centro: para filhos de estatura \textit{y}, qual a melhor “predição” (ou retrodicção) da estatura do pai? Ao pensar nessa pergunta, Galton tropeçou em um fato importante: as predições sempre caem sobre uma reta, a que chamou de reta de regressão, que é menos inclinada que o eixo maior (ou eixo de simetria) da elipse (Figura 2.3). De fato, há duas dessas retas, dependendo de qual variável está sendo predita e qual está sendo usada como evidência. Podemos predizer a estatura do filho com base na do pai, ou a do pai com base na do filho. A situação é completamente simétrica. Mais uma vez, isso mostra que, no que concerne à regressão à média, não há diferença entre causa e efeito.

\begin{figure}[ht]
\centering
\includegraphics[width=0.92\textwidth]{../figuras/fig_2_3.png}
\caption{Figura 2.3. As retas de regressão de Galton. A reta OM dá a melhor previsão da altura do filho se se conhece a do pai; a reta ON dá a melhor previsão da altura do pai se se conhece a do filho. Nenhuma das duas coincide com o eixo maior (eixo de simetria) do diagrama de dispersão. (Fonte: Francis Galton, \textit{Journal of the Anthropological Institute of Great Britain and Ireland} [1886].)}
\end{figure}
\textit{Figura 2.3. As retas de regressão de Galton. A reta OM dá a melhor previsão da altura do filho se se conhece a do pai; a reta ON dá a melhor previsão da altura do pai se se conhece a do filho. Nenhuma das duas coincide com o eixo maior (eixo de simetria) do diagrama de dispersão. (Fonte: Francis Galton, }Journal of the Anthropological Institute of Great Britain and Ireland\textit{ [1886].)}

A inclinação da regressão permite predizer o valor de uma variável, dado que se conhece o valor da outra. No contexto do problema de Galton, uma inclinação de 0,5 significaria que cada polegada extra de estatura do pai corresponderia, em média, a meia polegada extra para o filho, e vice-versa. Uma inclinação de 1 seria correlação perfeita, o que significa que cada polegada extra do pai é transmitida deterministicamente ao filho, que também seria uma polegada mais alto. A inclinação nunca pode ser maior que 1; se fosse, os filhos de pais altos seriam, em média, mais altos, e os filhos de pais baixos seriam mais baixos — e isso forçaria a distribuição das estaturas a se tornar mais larga ao longo do tempo. Depois de algumas gerações, começaríamos a ter pessoas de 2,70 m e de 60 cm, o que não se observa na natureza. Assim, desde que a distribuição das estaturas permaneça a mesma de uma geração para a seguinte, a inclinação da reta de regressão não pode exceder 1.

\textbf{Figura 2.3.} As retas de regressão de Galton. A reta OM dá a melhor predição da estatura de um filho se se conhece a estatura do pai; a reta ON dá a melhor predição da estatura de um pai se se conhece a estatura do filho. Nenhuma das duas é a mesma que o eixo maior (eixo de simetria) do diagrama de dispersão. (Fonte: Francis Galton, \textit{Journal of the Anthropological Institute of Great Britain and Ireland} [1886], 246–263, Placa X.)

A lei da regressão aplica-se mesmo quando correlacionamos duas quantidades diferentes, como estatura e QI. Se traçarmos uma quantidade contra a outra em um diagrama de dispersão e reescalarmos os dois eixos adequadamente, a inclinação da reta de melhor ajuste sempre desfruta das mesmas propriedades. Ela é igual a 1 apenas quando uma quantidade pode predizer a outra com precisão; é 0 sempre que a predição não é melhor do que um palpite aleatório. A inclinação (após o escalonamento) é a mesma, quer se trace \textit{X} contra \textit{Y} ou \textit{Y} contra \textit{X}. Em outras palavras, a inclinação é completamente agnóstica quanto a causa e efeito. Uma variável poderia causar a outra, ou ambas poderiam ser efeitos de uma terceira causa; para fins de predição, isso não importa.

Pela primeira vez, a ideia de correlação de Galton deu uma medida objetiva, independente de julgamento ou interpretação humana, de como duas variáveis se relacionam uma com a outra. As duas variáveis podem representar estatura, inteligência ou renda; podem estar em relação causal, neutra ou causal inversa. A correlação sempre refletirá o grau de previsibilidade cruzada entre as duas variáveis.

O discípulo de Galton, Karl Pearson, mais tarde derivou uma fórmula para a inclinação da reta de regressão (adequadamente reescalada) e chamou-a de coeficiente de correlação. Esse ainda é o primeiro número que estatísticos de todo o mundo calculam quando querem saber quão fortemente duas variáveis diferentes em um conjunto de dados estão relacionadas. Galton e Pearson devem ter ficado entusiasmados ao encontrar um modo tão universal de descrever as relações entre variáveis aleatórias. Para Pearson, em especial, os velhos e escorregadios conceitos de causa e efeito pareciam ultrapassados e não científicos, comparados ao conceito matematicamente claro e preciso de um coeficiente de correlação.

\section{Galton e a busca abandonada}
É uma ironia da história que Galton tenha partido em busca da causalidade e terminado descobrindo a correlação, uma relação que é indiferente à causalidade.

Ainda assim, indícios de pensamento causal permaneceram em seus escritos. “É fácil ver que a correlação [entre os tamanhos de dois órgãos] deve ser a consequência de as variações dos dois órgãos serem em parte devidas a causas comuns”, escreveu ele em 1889.

O primeiro sacrifício no altar da correlação foi a elaborada maquinaria de Galton para explicar a estabilidade da dotação genética da população. O \textit{quincunx} simulava a criação de variações na estatura e sua transmissão de uma geração para a seguinte. Mas Galton tivera de inventar as calhas inclinadas no \textit{quincunx} especificamente para refrear a diversidade sempre crescente da população. Tendo falhado em encontrar um mecanismo biológico satisfatório para dar conta dessa força restauradora, Galton simplesmente abandonou o esforço após oito anos e voltou sua atenção ao canto de sereia da correlação.

O historiador Stephen Stigler, que escreveu extensamente sobre Galton, notou essa súbita mudança nos objetivos e aspirações de Galton: “O que estava silenciosamente ausente era Darwin, as calhas e todo o ‘sobrevivência do mais apto.’… Em ironia suprema, o que começara como uma tentativa de matematizar o arcabouço de \textit{A origem das espécies} terminou com a essência daquela grande obra sendo descartada como desnecessária!”

Mas para nós, na era moderna da inferência causal, o problema original permanece. Como explicamos a estabilidade da população, a despeito das variações darwinianas que uma geração confere à seguinte?

Olhando em retrospecto para a máquina de Galton à luz dos diagramas causais, a primeira coisa que noto é que a máquina estava mal construída. A dispersão sempre crescente, que implorava a Galton por uma contrforça, nunca deveria ter estado ali em primeiro lugar. De fato, se rastrearmos uma bolinha caindo de um nível ao seguinte no \textit{quincunx}, vemos que o deslocamento no nível seguinte herda a soma total das variações conferidas a ela por todos os pinos ao longo do caminho. Isso está em gritante contradição com as equações de Kahneman:

Sucesso = talento + sorte

Grande sucesso = um pouco mais de talento + muita sorte.

Segundo essas equações, o sucesso na geração 2 não herda a sorte da geração 1. A sorte, por sua própria definição, é uma ocorrência transitória; portanto, não tem impacto sobre as gerações futuras. Mas tal comportamento transitório é incompatível com a máquina de Galton.

Para comparar essas duas concepções lado a lado, desenhemos seus diagramas causais associados. Na Figura 2.4(a) (a concepção de Galton), o sucesso é transmitido através das gerações, e as variações da sorte se acumulam indefinidamente. Isso é talvez natural se “sucesso” for equiparado a riqueza ou eminência. No entanto, para a herança de características físicas como a estatura, devemos substituir o modelo de Galton pelo da Figura 2.4(b). Aqui apenas o componente genético, mostrado como talento, é transmitido de uma geração à seguinte. A sorte afeta cada geração de modo independente, de tal forma que os fatores de acaso em uma geração não têm como afetar as gerações posteriores, seja direta ou indiretamente.

\begin{figure}[ht]
\centering
\includegraphics[width=0.92\textwidth]{../figuras/fig_2_4.png}
\caption{Figura 2.4. Dois modelos de herança. (a) O modelo do tabuleiro de Galton, em que a sorte se acumula de geração em geração, levando a uma distribuição cada vez mais ampla de sucesso. (b) Um modelo genético, em que a sorte não se acumula, levando a uma distribuição constante de sucesso.}
\end{figure}
\textit{Figura 2.4. Dois modelos de herança. (a) O modelo do tabuleiro de Galton, em que a sorte se acumula de geração em geração, levando a uma distribuição cada vez mais ampla de sucesso. (b) Um modelo genético, em que a sorte não se acumula, levando a uma distribuição constante de sucesso.}

\textbf{Figura 2.4.} Dois modelos de herança. (a) O modelo da tábua de Galton, em que a sorte se acumula de geração em geração, levando a uma distribuição de sucesso cada vez mais larga. (b) Um modelo genético, em que a sorte não se acumula, levando a uma distribuição constante de sucesso.

Ambos os modelos são compatíveis com a distribuição em forma de sino das estaturas. Mas o primeiro modelo não é compatível com a estabilidade da distribuição das estaturas (ou do sucesso). O segundo modelo, por outro lado, mostra que, para explicar a estabilidade do sucesso de uma geração para a seguinte, precisamos apenas explicar a estabilidade da dotação genética da população (talento). Essa estabilidade, hoje chamada equilíbrio de Hardy-Weinberg, recebeu uma explicação matemática satisfatória no trabalho de G. H. Hardy e Wilhelm Weinberg em 1908. E sim, eles usaram ainda outro modelo causal — a teoria mendeliana da herança.

Em retrospecto, Galton não poderia ter antecipado o trabalho de Mendel, Hardy e Weinberg. Em 1877, quando Galton deu sua palestra, o trabalho de Gregor Mendel de 1866 havia sido esquecido (só foi redescoberto em 1900), e a matemática das demonstrações de Hardy e Weinberg provavelmente estaria além dele. Mas é interessante notar quão perto ele chegou de encontrar o arcabouço correto e também como o diagrama causal torna fácil identificar sua premissa equivocada: a transmissão da sorte de uma geração para a seguinte. Infelizmente, ele foi levado a desviar-se por seu belo, porém falho, modelo causal e, mais tarde, tendo descoberto a beleza da correlação, passou a acreditar que a causalidade já não era necessária.

Como comentário pessoal final sobre a história de Galton, confesso cometer um pecado capital da escrita histórica, um de muitos pecados que cometerei neste livro. Nos anos 1960, tornou-se fora de moda escrever história do ponto de vista da ciência de nossos dias, como fiz acima. “História \textit{whig}” era o epíteto usado para ridicularizar o estilo retrospectivo de escrita histórica, que se concentrava em teorias e experimentos bem-sucedidos e dava pouco crédito a teorias fracassadas e becos sem saída. O estilo moderno de escrita histórica tornou-se mais democrático, tratando químicos e alquimistas com igual respeito e insistindo em compreender todas as teorias no contexto social de seu próprio tempo.

Quando se trata de explicar a expulsão da causalidade da estatística, contudo, aceito o manto de historiador \textit{whig} com orgulho. Simplesmente não há outro modo de compreender como a estatística se tornou um empreendimento de redução de dados cego a modelos, senão colocando nossas lentes causais e recontando as histórias de Galton e Pearson à luz da nova ciência de causa e efeito. De fato, ao fazê-lo, retifico as distorções introduzidas por historiadores mainstream que, carecendo de vocabulário causal, maravilham-se com a invenção da correlação e falham em notar sua baixa — a morte da causalidade.

\section{Pearson: a ira do zelote}
Restou ao discípulo de Galton, Karl Pearson, completar a tarefa de expurgar a causalidade da estatística. Ainda assim, nem ele teve êxito completo.

Ler \textit{Natural Inheritance}, de Galton, foi um dos momentos definidores da vida de Pearson: “Senti-me como um bucaneiro dos dias de Drake — um da ordem dos homens ‘não exatamente piratas, mas com tendências decididamente piratas’, como diz o dicionário!”, escreveu em 1934. “Interpretei… Galton como significando que havia uma categoria mais ampla que a causalidade, a saber, a correlação, da qual a causalidade era apenas o limite, e que essa nova concepção de correlação trazia a psicologia, a antropologia, a medicina e a sociologia em grande parte para o campo do tratamento matemático. Foi Galton quem primeiro me libertou do preconceito de que a matemática sólida só podia ser aplicada a fenômenos naturais sob a categoria da causalidade.”

Aos olhos de Pearson, Galton havia ampliado o vocabulário da ciência. A causalidade era reduzida a nada mais que um caso especial de correlação (a saber, o caso em que o coeficiente de correlação é 1 ou –1 e a relação entre \textit{x} e \textit{y} é determinística). Ele expressa sua visão da causalidade com grande clareza em \textit{The Grammar of Science} (1892): “Que uma certa sequência tenha ocorrido e reocorrido no passado é uma questão de experiência à qual damos expressão no conceito de causalidade.… A ciência em nenhum caso pode demonstrar qualquer necessidade inerente em uma sequência, nem provar com certeza absoluta que ela deve ser repetida.” Em resumo, a causalidade para Pearson é apenas uma questão de repetição e, no sentido determinístico, nunca pode ser provada.

Quanto à causalidade em um mundo não determinístico, Pearson era ainda mais desdenhoso: “a afirmação científica última de descrição da relação entre duas coisas pode sempre ser reduzida a… uma tabela de contingência.” Em outras palavras, os dados são tudo o que há na ciência. Ponto final. Nessa visão, as noções de intervenção e contrafactuais discutidas no Capítulo 1 não existem, e o degrau mais baixo da Escada da Causalidade é tudo o que se precisa para fazer ciência.

O salto mental de Galton a Pearson é de tirar o fôlego e de fato digno de um bucaneiro. Galton provara apenas que um fenômeno — a regressão à média — não requeria uma explicação causal. Agora Pearson estava removendo completamente a causalidade da ciência. O que o fez dar esse salto?

O historiador Ted Porter, em sua biografia \textit{Karl Pearson}, descreve como o ceticismo de Pearson em relação à causalidade antecedia sua leitura do livro de Galton. Pearson vinha se debatendo com o fundamento filosófico da física e escreveu (por exemplo): “A força como causa do movimento está exatamente no mesmo pé que um deus da árvore como causa do crescimento.” Mais geralmente, Pearson pertencia a uma escola filosófica chamada positivismo, que sustenta que o universo é um produto do pensamento humano e que a ciência é apenas uma descrição desses pensamentos. Assim, a causalidade, interpretada como um processo objetivo que acontece no mundo fora do cérebro humano, não podia ter nenhum significado científico. Pensamentos significativos só podem refletir padrões de observações, e estes podem ser completamente descritos por correlações.

Tendo decidido que a correlação era um descritor mais universal do pensamento humano do que a causalidade, Pearson estava preparado para descartar a causalidade completamente. Porter pinta um retrato vívido de Pearson ao longo de sua vida como um autodescrito \textit{Schwärmer}, uma palavra alemã que se traduz como “entusiasta”, mas que também pode ser interpretada de modo mais forte como “zelote”. Depois de se formar em Cambridge em 1879, Pearson passou um ano no exterior, na Alemanha, e se apaixonou tanto por sua cultura que prontamente mudou seu nome de Carl para Karl. Era socialista muito antes de isso se tornar popular, e escreveu a Karl Marx em 1881, oferecendo-se para traduzir \textit{O capital} para o inglês. Pearson, argumentavelmente um dos primeiros feministas da Inglaterra, fundou o Men’s and Women’s Club em Londres para discussões da “questão da mulher”. Preocupava-se com a posição subordinada das mulheres na sociedade e defendia que fossem remuneradas por seu trabalho. Era extremamente apaixonado por ideias e, ao mesmo tempo, muito cerebral em suas paixões. Levou-lhe quase meio ano persuadir sua futura esposa, Maria Sharpe, a casar-se com ele, e as cartas deles sugerem que ela estava francamente aterrorizada de não estar à altura de seus elevados ideais intelectuais.

Quando Pearson encontrou Galton e suas correlações, finalmente encontrou um foco para suas paixões: uma ideia que acreditava poder transformar o mundo da ciência e trazer rigor matemático a campos como a biologia e a psicologia. E avançou com o senso de propósito de um bucaneiro em direção à realização dessa missão. Seu primeiro artigo sobre estatística foi publicado em 1893, quatro anos após a descoberta da correlação por Galton. Em 1901, havia fundado um periódico, \textit{Biometrika}, que permanece um dos mais influentes periódicos de estatística (e, de modo um tanto herético, publicou meu primeiro artigo completo sobre diagramas causais em 1995). Em 1903, Pearson obtivera uma bolsa da Worshipful Company of Drapers para iniciar um Laboratório de Biometria no University College London. Em 1911, ele se tornou oficialmente um departamento quando Galton faleceu e deixou um legado para uma cátedra (com a estipulação de que Pearson fosse seu primeiro titular). Por pelo menos duas décadas, o Laboratório de Biometria de Pearson foi o centro mundial da estatística.

Uma vez que Pearson deteve uma posição de poder, seu zelo se revelou cada vez mais claramente. Como Porter escreve em sua biografia: “O movimento estatístico de Pearson tinha aspectos de uma seita cismática. Ele exigia a lealdade e o compromisso de seus associados e expulsava os dissidentes da igreja biométrica.” Um de seus primeiros assistentes, George Udny Yule, foi também uma das primeiras pessoas a sentir a ira de Pearson. O obituário de Yule sobre Pearson, escrito para a Royal Society em 1936, transmite bem a picada daqueles dias, embora envolta em linguagem polida.

\begin{quote}A contaminação de seu entusiasmo, é verdade, era inestimável; mas seu domínio, até mesmo sua própria ânsia de ajudar, podia ser uma desvantagem.… Esse desejo de dominação, de que tudo fosse exatamente como ele queria, manifesta-se de outras maneiras, notadamente na edição de \textit{Biometrika} — sem dúvida o periódico mais pessoalmente editado que já se publicou.… Aqueles que o deixavam e começavam a pensar por si mesmos tendiam, como aconteceu de modo doloroso em mais de um caso, a descobrir que após uma divergência de opinião a manutenção de relações amistosas se tornava difícil, e após crítica expressa, impossível.\end{quote}
Ainda assim, havia rachaduras no edifício de Pearson de uma ciência livre de causalidade, talvez ainda mais entre os fundadores do que entre os discípulos posteriores. Por exemplo, o próprio Pearson surpreendentemente escreveu vários artigos sobre “correlação espúria”, um conceito impossível de compreender sem alguma referência à causalidade. Pearson notou que é relativamente fácil encontrar correlações que são simplesmente absurdas. Por exemplo, para um exemplo divertido posterior à época de Pearson, há uma forte correlação entre o consumo de chocolate per capita de uma nação e seu número de vencedores do Prêmio Nobel. Essa correlação parece absurda porque não conseguimos imaginar nenhum modo pelo qual comer chocolate pudesse causar Prêmios Nobel. Uma explicação mais provável é que mais pessoas em países ocidentais ricos comem chocolate, e os vencedores do Prêmio Nobel também têm sido escolhidos preferencialmente nesses países. Mas essa é uma explicação causal, que, para Pearson, não é necessária ao pensamento científico. Para ele, a causalidade é apenas um “fetiche em meio aos arcanos inescrutáveis da ciência moderna.” A correlação deveria ser o objetivo da compreensão científica. Isso o coloca em uma posição embaraçosa quando tem de explicar por que uma correlação é significativa e outra é “espúria”. Ele explica que uma correlação genuína indica uma “relação orgânica” entre as variáveis, enquanto uma correlação espúria não. Mas o que é uma “relação orgânica”? Não é a causalidade por outro nome?

Juntos, Pearson e Yule compilaram vários exemplos de correlações espúrias. Um caso típico é o que hoje se chama confundimento, e a história do chocolate e do Nobel é um exemplo. (A riqueza e a localização são confundidores, ou causas comuns tanto do consumo de chocolate quanto da frequência de Nobels.) Outro tipo de “correlação sem sentido” frequentemente emerge em dados de séries temporais. Por exemplo, Yule encontrou uma correlação incrivelmente alta (0,95) entre a taxa de mortalidade da Inglaterra em um dado ano e a porcentagem de casamentos realizados naquele ano na Igreja da Inglaterra. Deus estava punindo anglicanos ávidos por casamento? Não! Duas tendências históricas separadas estavam simplesmente ocorrendo ao mesmo tempo: a taxa de mortalidade do país estava diminuindo e a filiação à Igreja da Inglaterra estava declinando. Como ambas estavam caindo ao mesmo tempo, havia uma correlação positiva entre elas, mas nenhuma conexão causal.

Pearson descobriu possivelmente o tipo mais interessante de “correlação espúria” já em 1899. Ela surge quando duas populações heterogêneas são agregadas em uma. Pearson, que, como Galton, era um colecionador fanático de dados sobre o corpo humano, obtivera medidas de 806 crânios masculinos e 340 crânios femininos das Catacumbas de Paris (Figura 2.5). Calculou a correlação entre o comprimento e a largura do crânio. Quando o cálculo era feito apenas para homens ou apenas para mulheres, as correlações eram negligenciáveis — não havia associação significativa entre comprimento e largura do crânio. Mas quando os dois grupos eram combinados, a correlação era 0,197, o que ordinariamente seria considerado significativo. Isso faz sentido, porque um comprimento de crânio pequeno agora é um indicador de que o crânio provavelmente pertencia a uma mulher e, portanto, de que a largura também será pequena. No entanto, Pearson considerou-o um artefato estatístico. O fato de a correlação ser positiva não tinha significado biológico ou “orgânico”; era apenas o resultado de combinar duas populações distintas de modo inapropriado.

\begin{figure}[ht]
\centering
\includegraphics[width=0.92\textwidth]{../figuras/fig_2_5.png}
\caption{Figura 2.5. Karl Pearson com um crânio das Catacumbas de Paris. (Fonte: desenho de Dakota Harr.)}
\end{figure}
\textit{Figura 2.5. Karl Pearson com um crânio das Catacumbas de Paris. (Fonte: desenho de Dakota Harr.)}

\textbf{Figura 2.5.} Karl Pearson com um crânio das Catacumbas de Paris. (Fonte: desenho de Dakota Harr.)

Esse exemplo é um caso de um fenômeno mais geral chamado paradoxo de Simpson. O Capítulo 6 discutirá quando é apropriado segregregar os dados em grupos separados e explicará por que correlações espúrias podem emergir da agregação. Mas vejamos o que Pearson escreveu: “Para aqueles que persistem em olhar todas as correlações como causa e efeito, o fato de que a correlação pode ser produzida entre dois caracteres bastante não correlacionados A e B tomando-se uma mistura artificial de duas raças estreitamente aliadas deve vir antes como um choque.” Como comenta Stephen Stigler: “Não resisto à especulação de que ele próprio foi o primeiro a se chocar.” Em essência, Pearson estava repreendendo a si mesmo pela tendência a pensar causalmente.

Olhando o mesmo exemplo pela lente da causalidade, só podemos dizer: Que oportunidade perdida! Em um mundo ideal, tais exemplos poderiam ter estimulado um cientista talentoso a pensar sobre a razão de seu choque e a desenvolver uma ciência para predizer quando correlações espúrias aparecem. No mínimo, ele deveria explicar quando agregar os dados e quando não. Mas a única orientação de Pearson a seus seguidores é que uma mistura “artificial” (seja lá o que isso signifique) é má. Ironicamente, usando nossa lente causal, sabemos agora que em alguns casos os dados agregados, e não os dados particionados, dão o resultado correto. A lógica da inferência causal pode de fato nos dizer em qual confiar. Oxalá Pearson estivesse aqui para desfrutá-lo!

Os estudantes de Pearson não o seguiram todos em fileira cerrada. Yule, que rompeu com Pearson por outras razões, rompeu com ele também por isso. Inicialmente, ele estava no campo intransigente que sustentava que as correlações dizem tudo o que poderíamos desejar compreender sobre a ciência. No entanto, mudou de ideia em certa medida quando precisou explicar as condições de pobreza em Londres. Em 1899, estudou a pergunta de se o “socorro externo” (isto é, a assistência entregue na casa de um indigente, em vez de em um asilo de pobres) aumentava a taxa de pobreza. Os dados mostravam que distritos com mais socorro externo tinham uma taxa de pobreza mais alta, mas Yule percebeu que a correlação era possivelmente espúria: esses distritos poderiam também ter mais idosos, que tendem a ser mais pobres. No entanto, mostrou então que mesmo em comparações de distritos com proporções iguais de idosos, a correlação permanecia. Isso o encorajou a dizer que a taxa de pobreza aumentada era \textit{devida} ao socorro externo. Mas depois de sair da linha para fazer essa asserção, voltou a se alinhar, escrevendo em uma nota de rodapé: “Em sentido estrito, por ‘devido a’ leia-se ‘associado a’.”

Isso estabeleceu o padrão para gerações de cientistas depois dele. Eles pensariam “devido a” e diriam “associado a”. Com Pearson e seus seguidores ativamente hostis à causalidade, e com dissidentes meio-coração como Yule temerosos de antagonizar seu líder, o palco estava armado para outro cientista do outro lado do oceano emitir o primeiro desafio direto à cultura de evitação da causalidade.

\section{Sewall Wright, cobaias e diagramas de caminhos}
Quando Sewall Wright chegou à Universidade Harvard em 1912, sua formação acadêmica mal sugeria o tipo de efeito duradouro que teria sobre a ciência. Havia frequentado uma pequena (e agora extinta) faculdade em Illinois, o Lombard College, formando-se em uma turma de apenas sete estudantes. Um de seus professores fora seu próprio pai, Philip Wright, um acadêmico pau-para-toda-obra que até administrava a gráfica da faculdade. Sewall e seu irmão Quincy ajudavam na gráfica e, entre outras coisas, publicaram a primeira poesia de um ainda não famoso estudante de Lombard, Carl Sandburg.

Os laços de Sewall Wright com o pai permaneceram muito estreitos muito depois de se formar. O pai Philip mudou-se para Massachusetts quando Sewall o fez. Mais tarde, quando Sewall trabalhou em Washington, D.C., Philip fez o mesmo, primeiro na Comissão Tarifária dos EUA e depois na Brookings Institution como economista. Embora seus interesses acadêmicos divergissem, encontraram no entanto modos de colaborar, e Philip foi o primeiro economista a fazer uso da invenção de seu filho, os diagramas de caminhos.

Wright veio a Harvard para estudar genética, na época um dos tópicos mais quentes da ciência porque a teoria de Gregor Mendel de genes dominantes e recessivos acabara de ser redescoberta. O orientador de Wright, William Castle, identificara oito diferentes fatores hereditários (ou genes, como os chamaríamos hoje) que afetavam a cor da pelagem em coelhos. Castle designou Wright para fazer a mesma coisa com cobaias.

Depois de obter o doutorado em 1915, Wright recebeu uma oferta para a qual estava unicamente qualificado: cuidar de cobaias no Departamento de Agricultura dos EUA (USDA). Pergunta-se se o USDA sabia o que estava obtendo ao contratar Wright. Talvez esperasse um zeloso cuidador de animais que pudesse pôr em ordem o caos de vinte anos de registros mal mantidos. Wright fez tudo isso e muito, muito mais.

As cobaias de Wright foram o trampolim de toda a sua carreira e de toda a sua teoria da evolução, de modo muito semelhante aos tentilhões das ilhas Galápagos que haviam inspirado Charles Darwin. Wright foi um dos primeiros defensores da visão de que a evolução não é gradual, como Darwin postulava, e sim ocorre em explosões relativamente súbitas. Em 1925, Wright mudou-se para um cargo de professor na Universidade de Chicago que provavelmente era mais adequado a alguém com seus amplos interesses teóricos. Ainda assim, permaneceu muito devotado a suas cobaias. Uma anedota frequentemente contada diz que certa vez ele estava segurando uma cobaia indócil sob o braço enquanto dava aula e, distraidamente, começou a usá-la para apagar o quadro-negro (ver Figura 2.6). Embora seus biógrafos concordem que essa história é provavelmente apócrifa, tais histórias frequentemente contêm mais verdade do que biografias áridas.

\begin{figure}[ht]
\centering
\includegraphics[width=0.92\textwidth]{../figuras/fig_2_6.png}
\caption{Figura 2.6. Sewall Wright foi a primeira pessoa a desenvolver um método matemático para responder perguntas causais a partir de dados, conhecido como diagramas de caminhos. Seu amor pela matemática só se rendia à paixão por cobaias. (Fonte: desenho de Dakota Harr.)}
\end{figure}
\textit{Figura 2.6. Sewall Wright foi a primeira pessoa a desenvolver um método matemático para responder perguntas causais a partir de dados, conhecido como diagramas de caminhos. Seu amor pela matemática só se rendia à paixão por cobaias. (Fonte: desenho de Dakota Harr.)}

O trabalho inicial de Wright no USDA é o que mais nos interessa aqui. A herança da cor da pelagem em cobaias teimosamente se recusava a seguir as regras mendelianas. Provou-se virtualmente impossível criar uma cobaia toda branca ou toda colorida, e mesmo as famílias mais endogâmicas (após múltiplas gerações de acasalamento irmão-irmã) ainda tinham variação pronunciada, de predominantemente branca a predominantemente colorida. Isso contradizia a predição da genética mendeliana de que um traço particular deveria se tornar “fixado” por múltiplas gerações de endogamia. Wright começou a duvidar de que a genética sozinha governasse a quantidade de branco e postulou que “fatores de desenvolvimento” no útero estavam causando parte das variações. Com o benefício do retrospecto, sabemos que ele estava correto. Diferentes genes de cor se expressam em diferentes lugares do corpo, e os padrões de cor dependem não apenas de quais genes o animal herdou, mas de onde e em quais combinações eles acontecem de ser expressos ou suprimidos.

\textbf{Figura 2.6.} Sewall Wright foi a primeira pessoa a desenvolver um método matemático para responder a perguntas causais a partir de dados, conhecido como diagramas de caminhos. Seu amor pela matemática só se rendia à sua paixão pelas cobaias. (Fonte: desenho de Dakota Harr.)

Como frequentemente acontece (pelo menos com os engenhosos!), um problema de pesquisa premente leva a novos métodos de análise, que transcendem vastamente suas origens na genética de cobaias. Ainda assim, para Sewall Wright, estimar os fatores de desenvolvimento provavelmente parecia um problema de nível universitário que ele poderia ter resolvido na aula de matemática de seu pai em Lombard. Ao procurar a magnitude de alguma quantidade desconhecida, primeiro se atribui um símbolo a essa quantidade, em seguida se expressa o que se sabe sobre essa e outras quantidades na forma de equações matemáticas e, por fim, se se tem paciência suficiente e equações suficientes, pode-se resolvê-las e encontrar a quantidade de interesse.

No caso de Wright, a quantidade desejada e desconhecida (mostrada na Figura 2.7) era \textit{d}, o efeito dos “fatores de desenvolvimento” sobre a pelagem branca. Outras quantidades causais que entravam em suas equações incluíam \textit{h}, para fatores “hereditários”, também desconhecidos. Por fim — e aqui vem a engenhosidade de Wright —, ele mostrou que se soubéssemos as quantidades causais na Figura 2.7, poderíamos predizer correlações nos dados (não mostradas no diagrama) por uma regra gráfica simples. Essa regra estabelece uma ponte do mundo profundo e oculto da causalidade para o mundo de superfície das correlações. Foi a primeira ponte jamais construída entre causalidade e probabilidade, a primeira travessia da barreira entre o degrau dois e o degrau um da Escada da Causalidade.

\begin{figure}[ht]
\centering
\includegraphics[width=0.92\textwidth]{../figuras/fig_2_7.png}
\caption{Figura 2.7. O primeiro diagrama de caminhos de Sewall Wright, ilustrando os fatores que levam à cor da pelagem em cobaias. \textit{D} = fatores de desenvolvimento (após a concepção, antes do nascimento), \textit{E} = fatores ambientais (após o nascimento), \textit{G} = fatores genéticos de cada genitor individual, \textit{H} = fatores hereditários combinados de ambos os genitores, \textit{O}, \textit{O′} = prole. O objetivo da análise era estimar a força dos efeitos de \textit{D}, \textit{E}, \textit{H} (escritos como \textit{d}, \textit{e}, \textit{h} no diagrama). (Fonte: Sewall Wright, \textit{Proceedings of the National Academy of Sciences} [1920].)}
\end{figure}
\textit{Figura 2.7. O primeiro diagrama de caminhos de Sewall Wright, ilustrando os fatores que levam à cor da pelagem em cobaias. }D\textit{ = fatores de desenvolvimento (após a concepção, antes do nascimento), }E\textit{ = fatores ambientais (após o nascimento), }G\textit{ = fatores genéticos de cada genitor individual, }H\textit{ = fatores hereditários combinados de ambos os genitores, }O\textit{, }O′\textit{ = prole. O objetivo da análise era estimar a força dos efeitos de }D\textit{, }E\textit{, }H\textit{ (escritos como }d\textit{, }e\textit{, }h\textit{ no diagrama). (Fonte: Sewall Wright, }Proceedings of the National Academy of Sciences\textit{ [1920].)}

Tendo construído essa ponte, Wright podia atravessá-la de volta, das correlações medidas nos dados (degrau um) para as quantidades causais ocultas, \textit{d} e \textit{h} (degrau dois). Fez isso resolvendo equações algébricas. Essa ideia deve ter parecido simples a Wright, mas revelou-se revolucionária porque foi a primeira prova de que o mantra “Correlação não implica causalidade” deveria ceder lugar a “Algumas correlações \textit{implicam} causalidade.”

\textbf{Figura 2.7.} O primeiro diagrama de caminhos de Sewall Wright, ilustrando os fatores que levam à cor da pelagem em cobaias. \textit{D} = fatores de desenvolvimento (após a concepção, antes do nascimento), \textit{E} = fatores ambientais (após o nascimento), \textit{G} = fatores genéticos de cada genitor individual, \textit{H} = fatores hereditários combinados de ambos os genitores, \textit{O}, \textit{O}′ = prole. O objetivo da análise era estimar a força dos efeitos de \textit{D}, \textit{E}, \textit{H} (escritos como \textit{d}, \textit{e}, \textit{h} no diagrama). (Fonte: Sewall Wright, \textit{Proceedings of the National Academy of Sciences} [1920], 320–332.)

No fim, Wright mostrou que os fatores de desenvolvimento hipotetizados eram mais importantes do que a hereditariedade. Em uma população de cobaias acasalada aleatoriamente, 42 por cento da variação no padrão da pelagem era devida à hereditariedade, e 58 por cento era de desenvolvimento. Em contraste, em uma família altamente endogâmica, apenas 3 por cento da variação na cobertura de pelagem branca era devida à hereditariedade, e 92 por cento era de desenvolvimento. Em outras palavras, vinte gerações de endogamia haviam praticamente eliminado a variação genética, mas os fatores de desenvolvimento permaneciam.

Por mais interessante que seja esse resultado, o cerne da questão para nossa história é o modo como Wright apresentou seu caso. O diagrama de caminhos na Figura 2.7 é o mapa de ruas que nos diz como navegar por essa ponte entre o degrau um e o degrau dois. É uma revolução científica em uma só figura — e vem completa com cobaias adoráveis!

Note que o diagrama de caminhos mostra todos os fatores concebíveis que poderiam afetar a pigmentação de uma cobaia bebê. As letras \textit{D}, \textit{E} e \textit{H} referem-se a fatores de desenvolvimento, ambientais e hereditários, respectivamente. Cada genitor (o reprodutor e a reprodutora) e cada filho (prole \textit{O} e \textit{O}′) tem seu próprio conjunto de fatores \textit{D}, \textit{E} e \textit{H}. Os dois filhos compartilham fatores ambientais, mas têm histórias de desenvolvimento diferentes. O diagrama incorpora os então novos \textit{insights} da genética mendeliana: a hereditariedade de um filho (\textit{H}) é determinada pelas células espermáticas e ovulares de seus genitores (\textit{G} e \textit{G}″), e estas por sua vez são determinadas a partir da hereditariedade dos genitores (\textit{H}″ e \textit{H}‴) via um processo de mistura que ainda não era compreendido (porque o DNA não havia sido descoberto). Compreendia-se, contudo, que o processo de mistura incluía um elemento de aleatoriedade (rotulado “Acaso” no diagrama).

Uma coisa que o diagrama não mostra explicitamente é a diferença entre uma família endogâmica e uma família normal. Em uma família endogâmica haveria uma forte correlação entre a hereditariedade do reprodutor e da reprodutora, que Wright indicava com uma seta de duas pontas entre \textit{H}″ e \textit{H}‴. Além disso, cada seta no diagrama é de mão única e leva de uma causa a um efeito. Por exemplo, a seta de \textit{G} para \textit{H} indica que a célula espermática do reprodutor pode ter um efeito causal direto sobre a hereditariedade da prole. A ausência de uma seta de \textit{G} para \textit{H}′ indica que a célula espermática que deu origem à prole \textit{O} não tem efeito causal sobre a hereditariedade da prole \textit{O}′. Quando se desmonta o diagrama seta por seta desse modo, penso que se encontrará que cada uma delas faz perfeito sentido.

Note também que cada seta é acompanhada de uma pequena letra (\textit{a}, \textit{b}, \textit{c}, etc.). Essas letras, chamadas coeficientes de caminho, representam a força dos efeitos causais que Wright queria resolver. Em termos gerais, um coeficiente de caminho representa a quantidade de variabilidade na variável-alvo que é explicada pela variável-fonte. Por exemplo, é bastante evidente que 50 por cento da constituição hereditária de cada filho deveria vir de cada genitor, de modo que \textit{a} deveria ser 1/2. (Por razões técnicas, Wright preferia tomar a raiz quadrada, de modo que \textit{a} = 1/√2 e \textit{a}² = 1/2.) Essa interpretação dos coeficientes de caminho, em termos da quantidade de variação explicada por uma variável, era razoável na época. A interpretação causal moderna é diferente: os coeficientes de caminho representam os resultados de uma intervenção hipotética sobre a variável-fonte. No entanto, a noção de uma intervenção teria de esperar até os anos 1940, e Wright não poderia tê-la antecipado quando escreveu seu artigo em 1920. Felizmente, nos modelos simples que ele analisou então, as duas interpretações produzem o mesmo resultado.

Quero enfatizar que o diagrama de caminhos não é apenas uma figura bonita; é um poderoso dispositivo computacional porque a regra para calcular correlações (a ponte do degrau dois ao degrau um) envolve rastrear os caminhos que conectam duas variáveis uma à outra e multiplicar os coeficientes encontrados ao longo do caminho. Note também que as setas omitidas de fato transmitem premissas mais significativas do que aquelas que estão presentes. Uma seta omitida restringe o efeito causal a zero, enquanto uma seta presente permanece totalmente agnóstica quanto à magnitude do efeito (a menos que \textit{a priori} imponhamos algum valor ao coeficiente de caminho).

O artigo de Wright foi um \textit{tour de force} e merece ser considerado um dos resultados marcos da biologia do século XX. Certamente é um marco para a história da causalidade. A Figura 2.7 é o primeiro diagrama causal jamais publicado, o primeiro passo da ciência do século XX sobre o segundo degrau da Escada da Causalidade. E não um passo tentativo, e sim ousado e decisivo!

No ano seguinte, Wright publicou um artigo bem mais geral chamado “Correlation and Causation” [Correlação e causalidade] que explicava como a análise de caminhos funcionava em outros contextos além da criação de cobaias. Não sei que tipo de reação o cientista de trinta anos esperava, mas a reação que obteve certamente deve tê-lo atordoado. Veio na forma de uma refutação publicada em 1921 por um certo Henry Niles, estudante do estatístico americano Raymond Pearl (sem parentesco), que por sua vez era estudante de Karl Pearson, o padrinho da estatística.

A academia está cheia de selvageria gentil, que tive a honra de enfrentar em momentos de minha própria carreira de resto plácida, mas mesmo assim raramente vi uma crítica tão selvagem quanto a de Niles. Ele começa com uma longa série de citações de seus heróis, Karl Pearson e Francis Galton, atestando a redundância ou mesmo a falta de sentido da palavra “causa”. Conclui: “Contrastar ‘causalidade’ e ‘correlação’ é injustificado porque a causalidade é simplesmente correlação perfeita.” Nessa sentença, ele está ecoando diretamente o que Pearson escreveu em \textit{Grammar of Science}.

Niles desdenha ainda mais a metodologia inteira de Wright. Escreve: “A falácia básica do método parece ser a premissa de que é possível estabelecer \textit{a priori} um sistema gráfico comparativamente simples que representará verdadeiramente as linhas de ação de várias variáveis umas sobre as outras, e sobre um resultado comum.” Por fim, Niles trabalha alguns exemplos e, trapalhando nos cálculos porque não se deu ao trabalho de compreender as regras de Wright, chega a conclusões opostas. Em resumo, declara: “Concluímos, portanto, que filosoficamente a base do método dos coeficientes de caminho é falha, enquanto na prática os resultados de aplicá-lo onde pode ser verificado provam que é totalmente não confiável.”

Do ponto de vista científico, uma discussão detalhada da crítica de Niles talvez não valha o tempo, mas seu artigo é muito importante para nós como historiadores da causalidade. Primeiro, ele reflete fielmente a atitude de sua geração em relação à causalidade e o domínio total que seu mentor, Karl Pearson, tinha sobre o pensamento científico de sua época. Segundo, continuamos a ouvir as objeções de Niles hoje.

É claro, às vezes os cientistas não conhecem toda a teia de relações entre suas variáveis. Nesse caso, argumentava Wright, podemos usar o diagrama em modo exploratório; podemos postular certas relações causais e calcular as correlações preditas entre as variáveis. Se estas contradisserem os dados, então temos evidência de que as relações que assumimos eram falsas. Esse modo de usar diagramas de caminhos, redescoberto em 1953 por Herbert Simon (laureado com o Nobel de Economia de 1978), inspirou muito trabalho nas ciências sociais.

Embora não precisemos conhecer toda relação causal entre as variáveis de interesse e possamos ser capazes de tirar algumas conclusões com apenas informação parcial, Wright deixa um ponto com clareza absoluta: não se pode tirar conclusões causais sem algumas hipóteses causais. Isso ecoa o que concluímos no Capítulo 1: não se pode responder a uma pergunta no degrau dois da Escada da Causalidade usando apenas dados coletados no degrau um.

Às vezes as pessoas me perguntam: “Isso não torna o raciocínio causal circular? Você não está apenas assumindo o que quer provar?” A resposta é não. Combinando premissas muito leves, qualitativas e óbvias (p.ex., a cor da pelagem do filho não influencia a dos pais) com seus vinte anos de dados de cobaias, ele obteve um resultado quantitativo e de modo algum óbvio: que 42 por cento da variação na cor da pelagem é devida à hereditariedade. Extrair o não óbvio do óbvio não é circular — é um triunfo científico e merece ser aclamado como tal.

A contribuição de Wright é única porque a informação que conduz à conclusão (de 42 por cento de herdabilidade) residia em duas linguagens matemáticas distintas, quase incompatíveis: a linguagem dos diagramas de um lado e a dos dados do outro. Essa ideia herética de casar “informação-de-setas” qualitativa com “informação-de-dados” quantitativa (duas línguas estrangeiras!) foi um dos milagres que primeiro me atraiu, como cientista da computação, a esse empreendimento.

Muitas pessoas ainda cometem o erro de Niles de pensar que o objetivo da análise causal é provar que \textit{X} é uma causa de \textit{Y} ou então encontrar a causa de \textit{Y} do zero. Esse é o problema da descoberta causal, que foi meu sonho ambicioso quando primeiro mergulhei na modelagem gráfica e ainda é uma área de pesquisa vigorosa. Em contraste, o foco da pesquisa de Wright, assim como deste livro, é representar conhecimento causal plausível em alguma linguagem matemática, combiná-lo com dados empíricos e responder a consultas causais que têm valor prático.

Wright compreendeu desde o início que a descoberta causal era muito mais difícil e talvez impossível. Em sua resposta a Niles, escreve: “O autor [isto é, o próprio Wright] nunca fez a pretensão absurda de que a teoria dos coeficientes de caminho fornece uma fórmula geral para a dedução de relações causais. Ele deseja submeter que a combinação do conhecimento de correlações com o conhecimento de relações causais para obter certos resultados é uma coisa diferente da dedução de relações causais a partir de correlações implicada pela afirmação de Niles.”

\section{\textit{E pur si muove} (E, no entanto, ela se move)}
Se eu fosse um historiador profissional, provavelmente pararia aqui. Mas como o “historiador \textit{whig}” que prometi ser, não consigo conter-me de expressar minha pura admiração pela precisão das palavras de Wright na citação que encerra a seção anterior, que não envelheceram nos noventa anos desde que ele as articulou pela primeira vez e que essencialmente definiram o novo paradigma da análise causal moderna.

Minha admiração pela precisão de Wright só é superada por minha admiração por sua coragem e determinação. Imagine a situação em 1921. Um matemático autodidata enfrenta sozinho a hegemonia do \textit{establishment} estatístico. Eles lhe dizem: “Seu método baseia-se em um completo equívoco quanto à natureza da causalidade no sentido científico.” E ele replica: “De modo algum! Meu método gera algo que é importante e vai além de qualquer coisa que vocês possam gerar.” Eles dizem: “Nossos gurus já examinaram esses problemas, duas décadas atrás, e concluíram que o que você fez é nonsense. Você apenas combinou correlações com correlações e obteve correlações. Quando crescer, vai compreender.” E ele continua: “Não estou dispensando seus gurus, mas pá é pá. Meus coeficientes de caminho não são correlações. São algo totalmente diferente: efeitos causais.”

Imagine que você está no jardim de infância, e seus amigos zombam de você por acreditar que 3 + 4 = 7, quando todo mundo sabe que 3 + 4 = 8. Depois imagine ir ao seu professor em busca de ajuda e ouvi-lo dizer, também, que 3 + 4 = 8. Você não iria para casa e se perguntaria se talvez houvesse algo errado com o modo como estava pensando? Mesmo o homem mais forte começaria a vacilar em suas convicções. Eu estive naquele jardim de infância, e sei. Mas Wright não piscou.

E isso não era apenas uma questão de aritmética, em que pode haver algum tipo de verificação independente. Apenas os filósofos ousavam expressar uma opinião sobre a natureza da causalidade. De onde Wright tirou essa convicção interna de que estava no caminho certo e o resto da turma do jardim de infância estava simplesmente errado? Talvez sua criação no Meio-Oeste e a minúscula faculdade que frequentou encorajassem sua autossuficiência e lhe ensinassem que o tipo mais seguro de conhecimento é o que se constrói a si mesmo.

Um dos primeiros livros de ciência que li na escola contava como a Inquisição forçou Galileu a se retratar de seu ensinamento de que a Terra gira em torno do Sol e como ele sussurrou baixinho: “E, no entanto, ela se move” (\textit{E pur si muove}). Não penso que haja uma criança no mundo que tenha lido essa lenda e não se tenha inspirado na coragem de Galileu em defender suas convicções. Ainda assim, por mais que o admiremos por sua posição, não posso deixar de pensar que ele ao menos tinha suas observações astronômicas em que se apoiar. Wright tinha apenas conclusões não testadas — digamos, que fatores de desenvolvimento respondem por 58 por cento, e não 3 por cento, da variação. Com nada em que se apoiar exceto sua convicção interna de que os coeficientes de caminho lhe dizem o que as correlações não dizem, ele ainda assim declarou: “E, no entanto, ela se move!”

Colegas me dizem que, quando as redes bayesianas lutaram contra o \textit{establishment} da inteligência artificial (ver Capítulo 3), eu agi de modo teimoso, monomaníaco e intransigente. De fato, recordo-me de estar totalmente convencido de minha abordagem, sem um átomo de hesitação. Mas eu tinha a teoria da probabilidade a meu lado. Wright não tinha nem um teorema em que se apoiar. Os cientistas haviam abandonado a causalidade, de modo que Wright não podia recorrer a nenhum arcabouço teórico. Tampouco podia se apoiar em autoridades, como Niles fazia, porque não havia ninguém a quem citar; os gurus já haviam pronunciado seus vereditos três décadas antes. Mas um consolo para Wright, e um sinal de que estava no caminho certo, deve ter sido sua compreensão de que podia responder a perguntas que não podem ser respondidas de nenhum outro modo.

Determinar a importância relativa de vários fatores era uma dessas perguntas. Outro belo exemplo disso pode ser encontrado em seu artigo “Correlation and Causation”, de 1921, que pergunta quanto o peso ao nascer de uma cobaia será afetado se ela passar um dia a mais no útero. Gostaria de examinar a resposta de Wright com algum detalhe para desfrutar a beleza de seu método e satisfazer leitores que gostariam de ver como a matemática da análise de caminhos funciona.

Note que não podemos responder à pergunta de Wright diretamente, porque não podemos pesar uma cobaia no útero. O que podemos fazer, contudo, é comparar os pesos ao nascer de cobaias que passam (digamos) sessenta e seis dias em gestação com aquelas que passam sessenta e sete dias. Wright notou que as cobaias que passavam um dia a mais no útero pesavam em média 5,66 gramas a mais ao nascer. Assim, poder-se-ia ingenuamente supor que um embrião de cobaia cresce a 5,66 gramas por dia logo antes de nascer.

“Errado!”, diz Wright. Os filhotes nascidos mais tarde geralmente nascem mais tarde por uma razão: têm menos irmãos de ninhada. Isso significa que tiveram um ambiente mais favorável para o crescimento ao longo da gravidez. Um filhote com apenas dois irmãos, por exemplo, já pesará mais no dia sessenta e seis do que um filhote com quatro irmãos. Assim, a diferença nos pesos ao nascer tem duas causas, e queremos desembaralhá-las. Quanto dos 5,66 gramas é devido a passar um dia adicional \textit{in utero} e quanto é devido a ter menos irmãos com quem competir?

Wright respondeu a essa pergunta estabelecendo um diagrama de caminhos (Figura 2.8). \textit{X} representa o peso ao nascer do filhote. \textit{Q} e \textit{P} representam as duas causas conhecidas do peso ao nascer: a duração da gestação (\textit{P}) e a taxa de crescimento \textit{in utero} (\textit{Q}). \textit{L} representa o tamanho da ninhada, que afeta tanto \textit{P} quanto \textit{Q} (uma ninhada maior faz o filhote crescer mais devagar e também ter menos dias \textit{in utero}). É muito importante perceber que \textit{X}, \textit{P} e \textit{L} podem ser medidos, para cada cobaia, mas \textit{Q} não. Por fim, \textit{A} e \textit{C} são causas exógenas sobre as quais não temos nenhum dado (p.ex., fatores hereditários e ambientais que controlam a taxa de crescimento e o tempo de gestação independentemente do tamanho da ninhada). A premissa importante de que esses fatores são independentes uns dos outros é transmitida pela ausência de qualquer seta entre eles, bem como de qualquer ancestral comum.

\begin{figure}[ht]
\centering
\includegraphics[width=0.92\textwidth]{../figuras/fig_2_8.png}
\caption{Figura 2.8. Diagrama causal (de caminhos) para o exemplo do peso ao nascer.}
\end{figure}
\textit{Figura 2.8. Diagrama causal (de caminhos) para o exemplo do peso ao nascer.}

\textbf{Figura 2.8.} Diagrama causal (de caminhos) para o exemplo do peso ao nascer.

Agora a pergunta que Wright enfrentava era: “Qual o efeito direto do período de gestação \textit{P} sobre o peso ao nascer \textit{X}?” Os dados (5,66 gramas por dia) não lhe dão o efeito direto; dão-lhe uma correlação, enviesada pelo tamanho da ninhada \textit{L}. Para obter o efeito direto, precisamos remover esse viés.

Na Figura 2.8, o efeito direto é representado pelo coeficiente de caminho \textit{p}, correspondente ao caminho \textit{P} → \textit{X}. O viés devido ao tamanho da ninhada corresponde ao caminho \textit{P} ← \textit{L} → \textit{Q} → \textit{X}. E agora a mágica algébrica: a quantidade de viés é igual ao produto dos coeficientes de caminho ao longo desse caminho (em outras palavras, \textit{l} vezes \textit{l}′ vezes \textit{q}). A correlação total, então, é apenas a soma dos coeficientes de caminho ao longo dos dois caminhos: algebricamente, \textit{p} + (\textit{l} × \textit{l}′ × \textit{q}) = 5,66 gramas por dia.

Se soubéssemos os coeficientes de caminho \textit{l}, \textit{l}′ e \textit{q}, então poderíamos simplesmente calcular o segundo termo e subtraí-lo de 5,66 para obter a quantidade desejada \textit{p}. Mas não os conhecemos, porque \textit{Q} (por exemplo) não é medido. Mas eis onde a engenhosidade dos coeficientes de caminho realmente brilha. Os métodos de Wright nos dizem como expressar cada uma das correlações medidas em termos dos coeficientes de caminho. Depois de fazer isso para cada um dos pares medidos (\textit{P}, \textit{X}), (\textit{L}, \textit{X}) e (\textit{L}, \textit{P}), obtemos três equações que podem ser resolvidas algebricamente para os coeficientes de caminho desconhecidos, \textit{p}, \textit{l}′ e \textit{l} × \textit{q}. Então estamos prontos, porque a quantidade desejada \textit{p} foi obtida.

Hoje podemos pular a matemática de todo e calcular \textit{p} por inspeção sumária do diagrama. Mas em 1920, essa foi a primeira vez que a matemática foi convocada a conectar causalidade e correlação. E funcionou! Wright calculou \textit{p} como 3,34 gramas por dia. Em outras palavras, se todas as outras variáveis (\textit{A}, \textit{L}, \textit{C}, \textit{Q}) fossem mantidas constantes e apenas o tempo de gestação fosse aumentado em um dia, o aumento médio no peso ao nascer seria de 3,34 gramas por dia.

Note que esse resultado é biologicamente significativo. Ele nos diz quão rapidamente os filhotes estão crescendo por dia antes do nascimento. Em contraste, o número 5,66 gramas por dia não tem nenhuma significância biológica, porque confunde dois processos separados, um dos quais não é causal e sim anticausal (ou diagnóstico) no elo \textit{P} ← \textit{L}.

Lição um desse exemplo: a análise causal nos permite quantificar processos no mundo real, não apenas padrões nos dados. Os filhotes estão crescendo a 3,34 gramas por dia, não a 5,66 gramas por dia. Lição dois, quer se tenha seguido a matemática ou não: na análise de caminhos, tira-se conclusões sobre relações causais individuais examinando o diagrama como um todo. A estrutura inteira do diagrama pode ser necessária para estimar cada parâmetro individual.

Em um mundo em que a ciência progride logicamente, a resposta de Wright a Niles deveria ter produzido uma excitação científica seguida de uma adoção entusiástica de seus métodos por outros cientistas e estatísticos. Mas não foi isso o que aconteceu. “Um dos mistérios da história da ciência de 1920 a 1960 é a virtual ausência de qualquer uso apreciável da análise de caminhos, exceto pelo próprio Wright e por estudantes de melhoramento animal”, escreveu um dos colegas geneticistas de Wright, James Crow. “Embora Wright tivesse ilustrado muitos problemas diversos aos quais o método era aplicável, nenhuma dessas pistas foi seguida.”

Crow não o sabia, mas o mistério se estendia também às ciências sociais. Em 1972, o economista Arthur Goldberger lamentava a “escandalosa negligência” do trabalho de Wright durante aquele período e notava, com o entusiasmo de um convertido, que “[a] abordagem [de Wright]… acendeu o recente surto de modelagem causal na sociologia.”

Se apenas pudéssemos voltar e perguntar aos contemporâneos de Wright: “Por que vocês não prestaram atenção?” Crow sugere uma razão: a análise de caminhos “não se presta a programas ‘enlatados’. O usuário tem de ter uma hipótese e deve conceber um diagrama apropriado de sequências causais múltiplas.” De fato, Crow pôs o dedo em um ponto essencial: a análise de caminhos exige pensamento científico, como todo exercício de inferência causal. A estatística, como frequentemente praticada, o desencoraja e encoraja procedimentos “enlatados” em seu lugar. Os cientistas sempre preferirão cálculos de rotina sobre dados a métodos que desafiem seu conhecimento científico.

R. A. Fisher, o incontestável sumo sacerdote da estatística na geração após Galton e Pearson, descreveu essa diferença de modo sucinto. Em 1925, escreveu: “A estatística pode ser considerada como… o estudo de métodos de redução de dados.” Preste atenção às palavras “métodos”, “redução” e “dados”. Wright abominava a ideia da estatística como meramente uma coleção de métodos; Fisher a abraçava. A análise causal enfaticamente não é apenas sobre dados; na análise causal, devemos incorporar alguma compreensão do processo que produz os dados, e então obtemos algo que não estava nos dados de início. Mas Fisher estava certo em um ponto: uma vez que se remove a causalidade da estatística, a redução de dados é a única coisa que resta.

Embora Crow não o tenha mencionado, o biógrafo de Wright, William Provine, aponta outro fator que pode ter afetado a falta de apoio à análise de caminhos. A partir de meados dos anos 1930, Fisher considerava Wright seu inimigo. Citei anteriormente Yule sobre como as relações com Pearson se tornavam tensas se se discordava dele e impossíveis se se o criticava. Exatamente a mesma coisa poderia ser dita sobre Fisher. Este último travava rixas desagradáveis com qualquer um com quem discordasse, incluindo Pearson, o filho de Pearson, Egon, Jerzy Neyman (mais se dirá sobre esses dois no Capítulo 8) e, é claro, Wright.

O verdadeiro foco da rivalidade Fisher-Wright não era a análise de caminhos, e sim a biologia evolutiva. Fisher discordava da teoria de Wright (chamada “deriva genética”) de que uma espécie pode evoluir rapidamente quando passa por um gargalo populacional. Os detalhes da disputa estão além do escopo deste livro, e o leitor interessado deve consultar Provine. O relevante aqui é isto: dos anos 1920 aos 1950, o mundo científico em grande medida se voltou a Fisher como seu oráculo de conhecimento estatístico. E pode-se ter certeza de que Fisher nunca disse uma palavra amável a ninguém sobre a análise de caminhos.

Nos anos 1960, as coisas começaram a mudar. Um grupo de cientistas sociais, incluindo Otis Duncan, Hubert Blalock e o economista Arthur Goldberger (mencionado anteriormente), redescobriu a análise de caminhos como um método de predizer o efeito de políticas sociais e educacionais. Em ainda outra ironia da história, Wright de fato fora convidado a falar a um influente grupo de econometristas chamado Cowles Commission em 1947, mas fracassou completamente em comunicar-lhes o que eram os diagramas de caminhos. Só quando os economistas chegaram a ideias semelhantes por si mesmos foi forjada uma conexão de curta duração.

Os destinos da análise de caminhos na economia e na sociologia seguiram trajetórias diferentes, cada uma levando a uma traição das ideias de Wright. Os sociólogos rebatizaram a análise de caminhos como modelagem de equações estruturais (SEM, do inglês \textit{structural equation modeling}), abraçaram os diagramas e os usaram extensivamente até 1970, quando um pacote de computador chamado LISREL automatizou o cálculo dos coeficientes de caminho (em alguns casos). Wright teria predito o que se seguiu: a análise de caminhos se transformou em um método de rotina, e os pesquisadores se tornaram usuários de \textit{software} com pouco interesse no que estava acontecendo sob o capô. No final dos anos 1980, um desafio público (pelo estatístico David Freedman) para explicar as premissas por trás da SEM ficou sem resposta, e alguns dos principais especialistas em SEM até renegaram que as SEMs tivessem qualquer coisa a ver com causalidade.

Na economia, a parte algébrica da análise de caminhos ficou conhecida como modelos de equações simultâneas (sem acrônimo). Os economistas essencialmente nunca usaram diagramas de caminhos e continuam a não usá-los até hoje, confiando em vez disso em equações numéricas e álgebra matricial. Uma consequência terrível disso é que, porque as equações algébricas são não direcionais (isto é, \textit{x} = \textit{y} é o mesmo que \textit{y} = \textit{x}), os economistas não tinham meios notacionais de distinguir equações causais de equações de regressão e, assim, eram incapazes de responder a perguntas relacionadas a políticas, mesmo depois de resolver as equações. Ainda em 1995, a maioria dos economistas se abstinha de atribuir explicitamente significado causal ou contrafactual a suas equações. Mesmo aqueles que usavam equações estruturais para decisões de política permaneciam incuravelmente suspeitosos de diagramas, que poderiam tê-los poupado páginas e páginas de cálculo. Não surpreendentemente, alguns economistas continuam a afirmar que “está tudo nos dados” até hoje.

Por todas essas razões, a promessa dos diagramas de caminhos permaneceu apenas parcialmente realizada, na melhor das hipóteses, até os anos 1990. Em 1983, o próprio Wright foi chamado de volta ao ringue mais uma vez para defendê-los, desta vez no \textit{American Journal of Human Genetics}. Quando escreveu esse artigo, Wright tinha mais de noventa anos. É ao mesmo tempo maravilhoso e trágico ler seu ensaio, escrito em 1983, sobre o mesmo tópico sobre o qual escrevera em 1923. Quantas vezes na história da ciência tivemos o privilégio de ouvir o criador de uma teoria sessenta anos depois de ele a ter posto no papel pela primeira vez? Seria como Charles Darwin voltando do túmulo para testemunhar no julgamento do Macaco de Scopes em 1925. Mas é também trágico, porque nos sessenta anos intervenientes sua teoria deveria ter se desenvolvido, crescido e florescido; em vez disso, pouco avançara desde os anos 1920.

A motivação para o artigo de Wright foi uma crítica da análise de caminhos, publicada no mesmo periódico, por Samuel Karlin (um matemático de Stanford e recipiente da Medalha Nacional de Ciência de 1989, que fez contribuições fundamentais à economia e à genética de populações) e dois coautores. De interesse para nós são dois dos argumentos de Karlin. Primeiro, Karlin objeta à análise de caminhos por uma razão que Niles não levantou: ela assume que todas as relações entre quaisquer duas variáveis no diagrama de caminhos são lineares. Essa premissa permite a Wright descrever as relações causais com um único número, o coeficiente de caminho. Se as equações não fossem lineares, então o efeito sobre \textit{Y} de uma mudança de uma unidade em \textit{X} poderia depender do valor atual de \textit{X}. Nem Karlin nem Wright perceberam que uma teoria geral não linear estava logo na esquina. (Ela seria desenvolvida três anos depois por um estudante estrela em meu laboratório, Thomas Verma.)

Mas a crítica mais interessante de Karlin foi também a que ele considerava a mais importante: “Por fim, e pensamos de modo mais frutífero, pode-se adotar uma abordagem essencialmente livre de modelo, buscando compreender os dados interativamente usando uma bateria de exibições, índices e contrastes. Essa abordagem enfatiza o conceito de robustez na interpretação de resultados.” Nessa única sentença, Karlin articula quão pouco mudara desde os dias de Pearson e quanta influência a ideologia de Pearson ainda tinha em 1983. Ele está dizendo que os dados em si já contêm toda a sabedoria científica; só precisam ser afagados e massageados (por “exibições, índices e contrastes”) para dispensar essas pérolas de sabedoria. Não há necessidade de nossa análise levar em conta o processo que gerou os dados. Faríamos tão bem, se não melhor, com uma “abordagem livre de modelo.” Se Pearson estivesse vivo hoje, vivendo na era do \textit{big data}, diria exatamente isso: as respostas estão todas nos dados.

É claro, a afirmação de Karlin viola tudo o que aprendemos no Capítulo 1. Para falar de causalidade, devemos ter um modelo mental do mundo real. Uma “abordagem livre de modelo” pode nos levar ao primeiro degrau da Escada da Causalidade, mas não além. Wright, para seu grande crédito, compreendeu o que estava em jogo e afirmou em termos inequívocos: “Ao tratar a abordagem livre de modelo (3) como alternativa preferida… Karlin et al. estão urgindo não meramente uma mudança de método, e sim um abandono do propósito da análise de caminhos e da avaliação da importância relativa de causas variáveis. Não pode haver tal análise sem um modelo. Seu conselho a qualquer um com o impulso de fazer tal avaliação é reprimi-lo e fazer outra coisa.”

Wright compreendeu que estava defendendo a própria essência do método científico e da interpretação de dados. Eu daria o mesmo conselho hoje aos entusiastas de \textit{big data} e livres de modelo. É claro, é correto extrair toda a informação que os dados podem fornecer, mas perguntemos até onde isso nos levará. Nunca nos levará além do primeiro degrau da Escada da Causalidade, e nunca responderá sequer a uma pergunta tão simples como “Qual a importância relativa de várias causas?”

\textit{E pur si muove!}

\section{Da objetividade à subjetividade — a conexão bayesiana}
Um outro tema na refutação de Wright pode sugerir outra razão para a resistência dos estatísticos à causalidade. Ele afirma repetidamente que não queria que a análise de caminhos se tornasse “estereotipada”. Segundo Wright: “A abordagem não estereotipada da análise de caminhos difere profundamente dos modos estereotipados de descrição concebidos para evitar quaisquer desvios da objetividade completa.”

O que ele quer dizer? Primeiro, quer dizer que a análise de caminhos deveria se basear na compreensão pessoal do usuário dos processos causais, refletida no diagrama causal. Não pode ser reduzida a rotinas mecânicas, como as estabelecidas em manuais de estatística. Para Wright, desenhar um diagrama de caminhos não é um exercício estatístico; é um exercício em genética, economia, psicologia ou qualquer que seja o campo de expertise do próprio cientista.

Segundo, Wright rastreia o fascínio dos métodos “livres de modelo” até sua objetividade. Essa de fato tem sido o Santo Graal dos estatísticos desde o primeiro dia — ou desde 15 de março de 1834, quando a Statistical Society of London foi fundada. Sua carta de fundação dizia que os dados deveriam receber prioridade em todos os casos sobre opiniões e interpretações. Dados são objetivos; opiniões são subjetivas. Esse paradigma antecede em muito Pearson. A luta pela objetividade — a ideia de raciocinar exclusivamente a partir de dados e experimentos — tem sido parte do modo como a ciência se define desde Galileu.

Ao contrário da correlação e da maioria das outras ferramentas da estatística mainstream, a análise causal exige que o usuário faça um compromisso subjetivo. Ela deve desenhar um diagrama causal que reflita sua crença qualitativa — ou, melhor ainda, a crença consensual dos pesquisadores em seu campo de expertise — sobre a topologia dos processos causais em ação. Deve abandonar o dogma secular da objetividade pela objetividade. Onde a causalidade está em jogo, um grão de subjetividade sábia nos diz mais sobre o mundo real do que qualquer quantidade de objetividade.

No parágrafo acima, disse que “a maioria” das ferramentas da estatística busca a objetividade completa. Há, contudo, uma importante exceção a essa regra. Um ramo da estatística chamado estatística bayesiana alcançou popularidade crescente ao longo dos últimos cinquenta anos ou mais. Outrora considerado quase anátema, agora se tornou completamente mainstream, e pode-se assistir a uma conferência inteira de estatística sem ouvir nenhum dos grandes debates entre “bayesianos” e “frequentistas” que costumavam trovejar nos anos 1960 e 1970.

O protótipo da análise bayesiana é assim: Crença prévia + Nova evidência → Crença revisada. Por exemplo, suponha que se lance uma moeda dez vezes e se descubra que em nove desses lançamentos a moeda deu cara. Sua crença de que a moeda é justa provavelmente é abalada, mas quanto? Um estatístico ortodoxo diria: “Na ausência de qualquer evidência adicional, eu acreditaria que essa moeda é viciada, de modo que apostaria nove a um que o próximo lançamento dá cara.” Um estatístico bayesiano, por outro lado, diria: “Espere um minuto. Também precisamos levar em conta nosso conhecimento prévio sobre a moeda.” Ela veio do armazém da esquina ou de um jogador suspeito? Se for apenas um quarto de dólar comum, a maioria de nós não deixaria a coincidência de nove caras abalar tão dramaticamente nossa crença. Por outro lado, se já suspeitássemos que a moeda era viciada, concluiríamos mais de bom grado que as nove caras forneciam evidência séria de viés.

A estatística bayesiana nos dá um modo objetivo de combinar a evidência observada com nosso conhecimento prévio (ou crença subjetiva) para obter uma crença revisada e, portanto, uma predição revisada do resultado do próximo lançamento da moeda. Ainda assim, o que os frequentistas não podiam suportar era que os bayesianos estavam permitindo que a opinião, na forma de probabilidades subjetivas, se intrometesse no reino imaculado da estatística. Os estatísticos mainstream só foram conquistados a contragosto, quando a análise bayesiana se mostrou uma ferramenta superior para uma variedade de aplicações, como a predição do tempo e o rastreamento de submarinos inimigos. Além disso, em muitos casos pode-se provar que a influência das crenças prévias desaparece à medida que o tamanho dos dados aumenta, deixando uma única conclusão objetiva no fim.

Infelizmente, a aceitação da subjetividade bayesiana na estatística mainstream nada fez para ajudar a aceitação da subjetividade causal, o tipo necessário para especificar um diagrama de caminhos. Por quê? A resposta repousa em uma grande barreira linguística. Para articular premissas subjetivas, os estatísticos bayesianos ainda usam a linguagem da probabilidade, a língua nativa de Galton e Pearson. As premissas que entram na inferência causal, por outro lado, exigem uma linguagem mais rica (p.ex., diagramas) que é estrangeira tanto a bayesianos quanto a frequentistas. A reconciliação entre bayesianos e frequentistas mostra que barreiras filosóficas podem ser superadas com boa vontade e uma linguagem comum. Barreiras linguísticas não são superadas com tanta facilidade.

Mais ainda, o componente subjetivo na informação causal não diminui necessariamente com o tempo, mesmo à medida que a quantidade de dados aumenta. Duas pessoas que acreditam em dois diagramas causais diferentes podem analisar os mesmos dados e podem nunca chegar à mesma conclusão, independentemente de quão “grandes” sejam os dados. Essa é uma perspectiva aterrorizante para os defensores da objetividade científica, o que explica sua recusa em aceitar a inevitabilidade de se apoiar em informação causal subjetiva.

No lado positivo, a inferência causal é objetiva em um sentido criticamente importante: uma vez que duas pessoas concordam em suas premissas, ela fornece um modo 100 por cento objetivo de interpretar qualquer nova evidência (ou dados). Ela compartilha essa propriedade com a inferência bayesiana. Assim, o leitor perspicaz provavelmente não se surpreenderá ao descobrir que cheguei à teoria da causalidade por um caminho tortuoso que começou com a probabilidade bayesiana e depois fez um enorme desvio pelas redes bayesianas. Contarei essa história no próximo capítulo.

\textit{Sherlock Holmes encontra seu correspondente moderno, um robô equipado com uma rede bayesiana. De modos diferentes, ambos estão enfrentando a pergunta de como inferir causas a partir de observações. A fórmula na tela do computador é a regra de Bayes. (Fonte: desenho de Maayan Harel.)}

\chapter{Capítulo 3 — Das evidências às causas: o reverendo Bayes encontra o sr. Holmes}
\begin{quote}Andarão dois juntos, se não houver entre eles acordo? Bramirá o leão no bosque, sem que tenha presa?\end{quote}
>

\begin{quote}— AMÓS 3:3\end{quote}
“ELEMENTAR, meu caro Watson.” Assim falava Sherlock Holmes (ao menos nos filmes) momentos antes de deslumbrar seu fiel assistente com uma de suas famosas deduções nada elementares. Mas, na verdade, Holmes não praticava apenas a dedução, que vai de uma hipótese a uma conclusão. Sua grande habilidade era a indução, que opera na direção oposta, da evidência à hipótese. Outra de suas citações célebres sugere seu \textit{modus operandi}: “Quando se elimina o impossível, o que resta, por mais improvável que seja, tem de ser a verdade.” Tendo induzido várias hipóteses, Holmes as eliminava uma a uma a fim de deduzir (por eliminação) a correta.

Embora indução e dedução andem de mãos dadas, a primeira é de longe a mais misteriosa. Esse fato manteve detetives como Sherlock Holmes no negócio. Contudo, em anos recentes, especialistas em inteligência artificial (IA) fizeram progresso considerável no sentido de automatizar o processo de raciocinar da evidência à hipótese e, da mesma forma, do efeito à causa. Tive a fortuna de participar dos estágios mais iniciais desse progresso ao desenvolver uma de suas ferramentas básicas, chamada redes bayesianas. Este capítulo explica o que elas são, examina algumas de suas aplicações atuais e discute o caminho tortuoso pelo qual me conduziram ao estudo da causalidade.

\section{Bonaparte, o detetive computadorizado}
Em 17 de julho de 2014, o voo 17 da Malaysia Airlines decolou do aeroporto de Schiphol, em Amsterdã, com destino a Kuala Lumpur. Infelizmente, o avião nunca chegou ao destino. Três horas depois da decolagem, ao sobrevoar o leste da Ucrânia, foi abatido por um míssil superfície-ar de fabricação russa. Todas as 298 pessoas a bordo — 283 passageiros e 15 tripulantes — morreram.

O dia 23 de julho, quando os primeiros corpos chegaram aos Países Baixos, foi declarado dia nacional de luto. Mas, para os investigadores do Instituto Forense dos Países Baixos (NFI, do inglês \textit{Netherlands Forensic Institute}), em Haia, 23 de julho foi o dia em que o relógio começou a contar. Sua tarefa era identificar os restos mortais dos falecidos o mais depressa possível e devolvê-los aos entes queridos para o sepultamento. O tempo era crucial, porque cada dia de incerteza traria nova angústia às famílias enlutadas.

Os investigadores enfrentavam muitos obstáculos. Os corpos estavam gravemente queimados, e muitos haviam sido conservados em formol, que degrada o DNA. Além disso, como o leste da Ucrânia era uma zona de guerra, os peritos forenses tinham acesso apenas esporádico ao local do acidente. Restos recém-recuperados continuaram a chegar por mais dez meses. Por fim, os investigadores não dispunham de registros prévios do DNA das vítimas, pela simples razão de que as vítimas não eram criminosas. Teriam de confiar, em vez disso, em correspondências parciais com membros da família.

Felizmente, os cientistas do NFI contavam com uma ferramenta poderosa a seu favor: um programa de ponta para identificação de vítimas de desastres chamado Bonaparte. Esse \textit{software}, desenvolvido em meados dos anos 2000 por uma equipe da Universidade Radboud, em Nijmegen, usa redes bayesianas para combinar informações de DNA obtidas de vários membros da família das vítimas. Graças em parte à precisão e à velocidade do Bonaparte, o NFI conseguiu identificar restos de 294 das 298 vítimas até dezembro de 2014. Em 2016, apenas duas vítimas do acidente (ambas cidadãos holandeses) haviam desaparecido sem deixar rastros.

As redes bayesianas, a ferramenta de raciocínio de máquina que sustenta o \textit{software} Bonaparte, afetam nossas vidas de muitas maneiras das quais a maioria das pessoas não tem consciência. São usadas em programas de reconhecimento de fala, em filtros de \textit{spam}, em previsão do tempo, na avaliação de poços de petróleo potenciais e no processo de aprovação de dispositivos médicos pela Food and Drug Administration. Se você joga videogames num Microsoft Xbox, uma rede bayesiana classifica sua habilidade. Se você tem um celular, os códigos que o aparelho usa para separar sua chamada de milhares de outras são decodificados por propagação de crenças (\textit{belief propagation}), um algoritmo concebido para redes bayesianas. Vint Cerf, o principal evangelista da internet em outra empresa de que você talvez já tenha ouvido falar, o Google, formula assim: “Somos enormes consumidores de métodos bayesianos.”

Neste capítulo, contarei a história das redes bayesianas desde suas raízes no século XVIII até seu desenvolvimento nos anos 1980, e darei mais alguns exemplos de como são usadas hoje. Elas se relacionam com os diagramas causais de um modo simples: um diagrama causal é uma rede bayesiana em que cada seta significa uma relação causal direta, ou ao menos a possibilidade de uma, na direção dessa seta. Nem todas as redes bayesianas são causais, e em muitas aplicações isso não importa. No entanto, se alguma vez quiser fazer uma consulta de degrau dois ou de degrau três sobre sua rede bayesiana, você terá de desenhá-la com escrupulosa atenção à causalidade.

\section{O reverendo Bayes e o problema da probabilidade inversa}
Thomas Bayes, em homenagem a quem nomeei as redes em 1985, jamais sonhou que uma fórmula que derivou nos anos 1750 um dia seria usada para identificar vítimas de desastres. Ele se preocupava apenas com as probabilidades de dois eventos, um (a hipótese) ocorrendo antes do outro (a evidência). Não obstante, a causalidade estava bem presente em sua mente. De fato, as aspirações causais foram a força motriz por trás de sua análise da “probabilidade inversa”.

Ministro presbiteriano que viveu de 1702 a 1761, o reverendo Thomas Bayes parece ter sido um \textit{nerd} da matemática. Como dissidente da Igreja da Inglaterra, não podia estudar em Oxford ou Cambridge e foi educado, em vez disso, na Universidade da Escócia, onde provavelmente absorveu uma boa dose de matemática. Continuou a se dedicar a ela e a organizar círculos de discussão matemática depois de voltar à Inglaterra. Num artigo publicado após sua morte (ver Figura 3.1), Bayes enfrentou um problema que era o casamento perfeito para ele, pondo a matemática contra a teologia.

\begin{figure}[ht]
\centering
\includegraphics[width=0.92\textwidth]{../figuras/fig_3_1.png}
\caption{Figura 3.1. ), Bayes enfrentou um problema que era o casamento perfeito para ele, pondo a matemática contra a teologia.}
\end{figure}
\textit{Figura 3.1. ), Bayes enfrentou um problema que era o casamento perfeito para ele, pondo a matemática contra a teologia.}

Para situar o contexto: em 1748, o filósofo escocês David Hume havia escrito um ensaio intitulado “Sobre os milagres” (\textit{On Miracles}), no qual argumentava que o testemunho ocular jamais poderia provar que um milagre havia ocorrido. O milagre que Hume tinha em mente era, é claro, a ressurreição de Cristo, embora fosse esperto o bastante para não dizê-lo. (Vinte anos antes, o teólogo Thomas Woolston fora para a prisão por blasfêmia por escrever coisas semelhantes.) O ponto principal de Hume era que evidências inerentemente falíveis não podem sobrepujar uma proposição com a força de uma lei natural, como “Os mortos permanecem mortos”.

\textbf{Figura 3.1.} Página de rosto do periódico em que foi publicado o artigo póstumo de Thomas Bayes sobre probabilidade inversa e a primeira página da introdução de Richard Price.

Para Bayes, essa afirmação provocou uma pergunta natural, poderíamos dizer holmesiana: quanta evidência seria necessária para nos convencer de que algo que consideramos improvável de fato aconteceu? Quando uma hipótese cruza a linha da impossibilidade para a improbabilidade e até para a probabilidade ou a certeza virtual? Embora a pergunta fosse formulada na linguagem da probabilidade, as implicações eram intencionalmente teológicas. Richard Price, um colega ministro que encontrou o ensaio entre os bens de Bayes após sua morte e o enviou para publicação com uma introdução entusiástica que ele próprio escreveu, deixou esse ponto abundantemente claro:

\begin{quote}O propósito a que me refiro é mostrar que razão temos para crer que há, na constituição das coisas, leis fixas segundo as quais as coisas acontecem e que, portanto, a estrutura do mundo deve ser o efeito da sabedoria e do poder de uma causa inteligente; e assim confirmar o argumento tomado das causas finais em favor da existência da Divindade. Será fácil ver que o problema inverso resolvido neste ensaio é mais diretamente aplicável a esse propósito; pois nos mostra, com nitidez e precisão, em cada caso de qualquer ordem particular ou recorrência de eventos, que razão há para pensar que tal recorrência ou ordem deriva de causas ou regulações estáveis na natureza, e não de quaisquer irregularidades do acaso.\end{quote}
O próprio Bayes não discutiu nada disso em seu artigo; Price destacou essas implicações teológicas, talvez para tornar o impacto do artigo de seu amigo mais abrangente. Mas acabou por se revelar que Bayes não precisava da ajuda. Seu artigo é lembrado e debatido 250 anos depois, não por sua teologia, mas porque mostra que se pode deduzir a probabilidade de uma causa a partir de um efeito. Se conhecemos a causa, é fácil estimar a probabilidade do efeito, que é uma probabilidade \textit{direta} (\textit{forward probability}). Ir na direção oposta — um problema conhecido na época de Bayes como “probabilidade inversa” — é mais difícil. Bayes não explicou por que é mais difícil; tomou isso como evidente por si, provou que é factível e nos mostrou como.

Para apreciar a natureza do problema, vejamos o exemplo que ele próprio sugeriu em seu artigo póstumo de 1763. Imagine que disparamos uma bola de bilhar sobre uma mesa, assegurando-nos de que ela quique muitas vezes de modo que não tenhamos ideia de onde parará. Qual é a probabilidade de que ela pare a menos de \textit{x} pés da extremidade esquerda da mesa? Se conhecemos o comprimento da mesa e esta é perfeitamente lisa e plana, trata-se de uma pergunta muito fácil (Figura 3.2, em cima). Por exemplo, numa mesa de snooker de doze pés, a probabilidade de a bola parar a menos de um pé da extremidade seria 1/12. Numa mesa de bilhar de oito pés, a probabilidade seria 1/8.

\begin{figure}[ht]
\centering
\includegraphics[width=0.92\textwidth]{../figuras/fig_3_2.png}
\caption{Figura 3.2. , em cima). Por exemplo, numa mesa de snooker de doze pés, a probabilidade de a bola parar a menos de um pé da extremidade seria 1/12. Numa mesa de bilhar de oito pés, a probabilidade seria 1/8.}
\end{figure}
\textit{Figura 3.2. , em cima). Por exemplo, numa mesa de snooker de doze pés, a probabilidade de a bola parar a menos de um pé da extremidade seria 1/12. Numa mesa de bilhar de oito pés, a probabilidade seria 1/8.}

\textbf{Figura 3.2.} O exemplo da mesa de bilhar de Thomas Bayes. Na primeira versão, uma pergunta de probabilidade direta, conhecemos o comprimento da mesa e queremos calcular a probabilidade de a bola parar a menos de \textit{x} pés da extremidade. Na segunda, uma pergunta de probabilidade inversa, observamos que a bola parou a \textit{x} pés da extremidade e queremos estimar a verossimilhança de que o comprimento da mesa seja \textit{L}. (Fonte: desenho de Maayan Harel.)

Nossa compreensão intuitiva da física nos diz que, em geral, se o comprimento da mesa é \textit{L} pés, a probabilidade de a bola parar a menos de \textit{x} pés da extremidade é \textit{x}/\textit{L}. Quanto maior o comprimento da mesa (\textit{L}), menor a probabilidade, porque há mais posições competindo pela honra de ser o local de repouso da bola. Por outro lado, quanto maior \textit{x}, maior a probabilidade, porque inclui um conjunto maior de posições de parada.

Agora considere o problema da probabilidade inversa. Observamos a posição final da bola como \textit{x} = 1 pé da extremidade, mas não nos é dado o comprimento \textit{L} (Figura 3.2, embaixo). O reverendo Bayes perguntou: qual é a probabilidade de que o comprimento tenha sido, digamos, cem pés? O senso comum nos diz que \textit{L} é mais provavelmente cinquenta pés do que cem pés, porque a mesa mais longa torna mais difícil explicar por que a bola parou tão perto da extremidade. Mas \textit{quanto} mais provavelmente? A “intuição” ou o “senso comum” não nos dão nenhuma orientação clara.

Por que a probabilidade direta (de \textit{x} dado \textit{L}) era tão mais fácil de avaliar mentalmente do que a probabilidade de \textit{L} dado \textit{x}? Neste exemplo, a assimetria vem do fato de que \textit{L} atua como a causa e \textit{x} é o efeito. Se observamos uma causa — por exemplo, Bobby arremessa uma bola em direção a uma janela —, a maioria de nós consegue predizer o efeito (a bola provavelmente quebrará a janela). A cognição humana opera nessa direção. Mas, dado o efeito (a janela está quebrada), precisamos de bem mais informação para deduzir a causa (qual menino arremessou a bola que a quebrou, ou mesmo o fato de que ela foi quebrada por uma bola em primeiro lugar). É preciso a mente de um Sherlock Holmes para acompanhar todas as causas possíveis.

Bayes propôs-se a romper essa assimetria cognitiva e explicar como até humanos comuns podem avaliar probabilidades inversas. Para ver como o método de Bayes funciona, comecemos com um exemplo simples sobre clientes de uma casa de chá, para os quais temos dados que documentam suas preferências. Os dados, como sabemos do Capítulo 1, são totalmente indiferentes a assimetrias de causa e efeito e, portanto, deveriam nos oferecer um modo de resolver o enigma da probabilidade inversa.

Suponha que dois terços dos clientes que entram na loja pedem chá, e metade dos bebedores de chá também pede \textit{scones}. Que fração da clientela pede tanto chá quanto \textit{scones}? Não há truque nessa pergunta, e espero que a resposta seja quase óbvia. Como a metade de dois terços é um terço, segue-se que um terço dos clientes pede tanto chá quanto \textit{scones}.

Para uma ilustração numérica, suponha que tabulamos os pedidos dos próximos doze clientes que entram pela porta. Como a Tabela 3.1 mostra, dois terços dos clientes (1, 5, 6, 7, 8, 9, 10, 12) pediram chá, e a metade dessas pessoas pediu \textit{scones} (1, 5, 8, 12). Assim, a proporção de clientes que pediram tanto chá quanto \textit{scones} é de fato (1/2) × (2/3) = 1/3, exatamente como predissemos antes de ver os dados específicos.

\begin{figure}[ht]
\centering
\includegraphics[width=0.92\textwidth]{../figuras/tab_3_1.png}
\caption{Tabela 3.1. mostra, dois terços dos clientes (1, 5, 6, 7, 8, 9, 10, 12) pediram chá, e a metade dessas pessoas pediu \textit{scones} (1, 5, 8, 12). Assim, a proporção de clientes que pediram tanto chá quanto \textit{scones} é de fato (1/2) × (2/3) = 1/3, exatamente como predissemos antes de ver os dados específicos.}
\end{figure}
\textit{Tabela 3.1. mostra, dois terços dos clientes (1, 5, 6, 7, 8, 9, 10, 12) pediram chá, e a metade dessas pessoas pediu }scones\textit{ (1, 5, 8, 12). Assim, a proporção de clientes que pediram tanto chá quanto }scones\textit{ é de fato (1/2) × (2/3) = 1/3, exatamente como predissemos antes de ver os dados específicos.}

\textbf{Tabela 3.1.} Dados fictícios para o exemplo do chá e dos \textit{scones}.

O ponto de partida da regra de Bayes é notar que poderíamos ter analisado os dados na ordem inversa. Isto é, poderíamos ter observado que cinco doze avos dos clientes (1, 2, 5, 8, 12) pediram \textit{scones}, e quatro quintos destes (1, 5, 8, 12) pediram chá. Assim, a proporção de clientes que pediram tanto chá quanto \textit{scones} é (4/5) × (5/12) = 1/3. É claro que não é coincidência que tenha dado o mesmo; estávamos meramente computando a mesma quantidade de dois modos diferentes. A ordem temporal em que os clientes anunciam seu pedido não faz diferença.

Para tornar isso uma regra geral, podemos deixar \textit{P}(\textit{T}) denotar a probabilidade de que um cliente peça chá e \textit{P}(\textit{S}) denotar a probabilidade de que peça \textit{scones}. Se já sabemos que um cliente pediu chá, então \textit{P}(\textit{S} | \textit{T}) denota a probabilidade de que ele peça \textit{scones}. (Lembre-se de que a barra vertical significa “dado que”.) Da mesma forma, \textit{P}(\textit{T} | \textit{S}) denota a probabilidade de que ele peça chá, dado que já sabemos que pediu \textit{scones}. Então o primeiro cálculo que fizemos diz:

\textit{P}(\textit{S} E \textit{T}) = \textit{P}(\textit{S} | \textit{T}) \textit{P}(\textit{T}).

O segundo cálculo diz:

\textit{P}(\textit{S} E \textit{T}) = \textit{P}(\textit{T} | \textit{S}) \textit{P}(\textit{S}).

Agora, como Euclides disse há 2.300 anos, duas coisas que são iguais a uma terceira também são iguais entre si. Isso significa que deve ser o caso que

\textit{P}(\textit{S} | \textit{T}) \textit{P}(\textit{T}) = \textit{P}(\textit{T} | \textit{S}) \textit{P}(\textit{S})  (3.1)

Essa equação de aparência inocente veio a ser conhecida como “regra de Bayes”. Se olharmos com cuidado o que ela diz, descobrimos que oferece uma solução geral para o problema da probabilidade inversa. Ela nos diz que, se conhecemos a probabilidade de \textit{S} dado \textit{T}, \textit{P}(\textit{S} | \textit{T}), devemos ser capazes de descobrir a probabilidade de \textit{T} dado \textit{S}, \textit{P}(\textit{T} | \textit{S}), supondo é claro que conheçamos \textit{P}(\textit{T}) e \textit{P}(\textit{S}). Este é talvez o papel mais importante da regra de Bayes na estatística: podemos estimar a probabilidade condicional diretamente numa direção, para a qual nosso julgamento é mais confiável, e usar a matemática para derivar a probabilidade condicional na outra direção, para a qual nosso julgamento é bastante nebuloso. A equação também desempenha esse papel nas redes bayesianas; dizemos ao computador as probabilidades diretas, e o computador nos diz as probabilidades inversas quando necessário.

Para ver como a regra de Bayes funciona no exemplo da casa de chá, suponha que você não se deu ao trabalho de calcular \textit{P}(\textit{T} | \textit{S}) e deixou em casa a planilha contendo os dados. Contudo, você se lembra de que metade dos que pedem chá também pedem \textit{scones}, e que dois terços dos clientes pedem chá e cinco doze avos pedem \textit{scones}. Inesperadamente, seu chefe pergunta: “Mas que proporção dos comedores de \textit{scones} pede chá?” Não há necessidade de entrar em pânico, porque você pode trabalhar a partir das outras probabilidades. A regra de Bayes diz que \textit{P}(\textit{T} | \textit{S}) (5/12) = (1/2)(2/3), de modo que sua resposta é \textit{P}(\textit{T} | \textit{S}) = 4/5, porque 4/5 é o único valor de \textit{P}(\textit{T} | \textit{S}) que torna essa equação verdadeira.

Também podemos olhar a regra de Bayes como um modo de atualizar nossa crença numa hipótese particular. Isso é extremamente importante de compreender, porque grande parte da crença humana sobre eventos futuros repousa na frequência com que eles ou eventos semelhantes ocorreram no passado. De fato, quando um cliente entra pela porta do restaurante, cremos, com base em nossos encontros passados com clientes semelhantes, que ela provavelmente quer chá. Mas se primeiro pede \textit{scones}, ficamos ainda mais certos. De fato, poderíamos até sugerir: “Presumo que queira chá com isso?” A regra de Bayes simplesmente nos permite atribuir números a esse processo de raciocínio. Da Tabela 3.1, vemos que a probabilidade \textit{a priori} de que a cliente queira chá (isto é, quando ela entra pela porta, antes de pedir qualquer coisa) é de dois terços. Mas se a cliente pede \textit{scones}, agora temos informação adicional sobre ela que não tínhamos antes. A probabilidade atualizada de que ela queira chá, dado que pediu \textit{scones}, é \textit{P}(\textit{T} | \textit{S}) = 4/5.

Matematicamente, é isso tudo que há na regra de Bayes. Parece quase trivial. Envolve nada mais do que o conceito de probabilidade condicional, mais uma pequena dose de lógica grega antiga. Você poderia, com razão, perguntar como um truque tão simples poderia tornar Bayes famoso e por que as pessoas discutem sobre sua regra há 250 anos. Afinal, fatos matemáticos deveriam resolver controvérsias, não criá-las.

Aqui devo confessar que, no exemplo da casa de chá, ao derivar a regra de Bayes a partir de dados, passei por cima de duas objeções profundas, uma filosófica e outra prática. A filosófica decorre da interpretação das probabilidades como um grau de crença, que usamos implicitamente no exemplo da casa de chá. Quem jamais disse que as crenças agem, ou deveriam agir, como proporções nos dados? O cerne do debate filosófico é se podemos legitimamente traduzir a expressão “dado que sei” para a linguagem das probabilidades. Mesmo se concordarmos que as probabilidades incondicionais \textit{P}(\textit{S}), \textit{P}(\textit{T}) e \textit{P}(\textit{S} E \textit{T}) refletem meu grau de crença nessas proposições, quem diz que meu grau revisado de crença em \textit{T} deveria ser igual à razão \textit{P}(\textit{S} E \textit{T})/\textit{P}(\textit{T}), como dita a regra de Bayes? “Dado que sei \textit{T}” é o mesmo que “entre os casos em que \textit{T} ocorreu”?

A linguagem da probabilidade, expressa em símbolos como \textit{P}(\textit{S}), destinava-se a capturar o conceito de frequências em jogos de azar. Mas a expressão “dado que sei” é epistemológica e deveria ser governada pela lógica do conhecimento, não pela das frequências e proporções.

Da perspectiva filosófica, o feito de Thomas Bayes reside em propor a primeira definição formal de probabilidade condicional como a razão \textit{P}(\textit{S} | \textit{T}) = \textit{P}(\textit{S} E \textit{T})/\textit{P}(\textit{T}). Seu ensaio era, admitidamente, nebuloso; ele não tem o termo “probabilidade condicional” e usa, em vez disso, a linguagem pesada “a probabilidade do 2º [evento] sob a suposição de que o 1º acontece”. O reconhecimento de que a relação “dado que” merece seu próprio símbolo evoluiu apenas nos anos 1880, e só em 1931 Harold Jeffreys (conhecido mais como geofísico do que como teórico da probabilidade) introduziu a agora padrão barra vertical em \textit{P}(\textit{S} | \textit{T}).

Como vimos, a regra de Bayes é formalmente uma consequência elementar de sua definição de probabilidade condicional. Mas epistemologicamente está longe de ser elementar. Age, de fato, como uma regra normativa para atualizar crenças em resposta a evidências. Em outras palavras, deveríamos ver a regra de Bayes não apenas como uma definição conveniente do novo conceito de “probabilidade condicional”, mas como uma afirmação empírica de representar fielmente a expressão em inglês “given that I know” [“dado que sei”]. Ela afirma, entre outras coisas, que a crença que uma pessoa atribui a \textit{S} depois de descobrir \textit{T} nunca é menor do que o grau de crença que essa pessoa atribui a \textit{S} E \textit{T} antes de descobrir \textit{T}. Também implica que, quanto mais surpreendente a evidência \textit{T} — isto é, quanto menor for \textit{P}(\textit{T}) —, mais convencido se deveria ficar de sua causa \textit{S}.

Não é de admirar que Bayes e seu amigo Price, como ministros episcopais, vissem nisso uma réplica eficaz a Hume. Se \textit{T} é um milagre (“Cristo ressuscitou dos mortos”) e \textit{S} é uma hipótese intimamente relacionada (“Cristo é o filho de Deus”), nosso grau de crença em \textit{S} aumenta de modo muito dramático se sabemos com certeza que \textit{T} é verdadeiro. Quanto mais milagroso o milagre, mais crível a hipótese que explica sua ocorrência. Isso explica por que os escritores do Novo Testamento ficaram tão impressionados com sua evidência de testemunhas oculares.

Agora deixe-me discutir a objeção prática à regra de Bayes — que pode ser ainda mais consequente quando saímos do reino da teologia e entramos no reino da ciência. Se tentarmos aplicar a regra ao enigma da bola de bilhar, a fim de encontrar \textit{P}(\textit{L} | \textit{x}), precisamos de uma quantidade que não está disponível para nós a partir da física das bolas de bilhar: precisamos da probabilidade \textit{a priori} do comprimento \textit{L}, que é tão difícil de estimar quanto a \textit{P}(\textit{L} | \textit{x}) desejada. Além disso, essa probabilidade variará significativamente de pessoa para pessoa, dependendo da experiência prévia de cada indivíduo com mesas de comprimentos diferentes. Uma pessoa que jamais em sua vida viu uma mesa de snooker duvidaria muito de que \textit{L} pudesse ser maior do que dez pés. Uma pessoa que só viu mesas de snooker e jamais viu uma mesa de bilhar, por outro lado, daria uma probabilidade \textit{a priori} muito baixa a \textit{L} ser menor do que dez pés.

Essa variabilidade, também conhecida como “subjetividade”, é às vezes vista como uma deficiência da inferência bayesiana. Outros a consideram uma vantagem poderosa; ela nos permite expressar nossa experiência pessoal matematicamente e combiná-la com dados de um modo principiado e transparente.

A regra de Bayes informa nosso raciocínio em casos em que a intuição ordinária nos falha ou em que a emoção poderia nos desviar. Demonstraremos esse poder numa situação familiar a todos nós. Suponha que você faça um exame médico para ver se tem uma doença, e o resultado é positivo. Qual a probabilidade de que você tenha a doença? Para concretude, digamos que a doença é câncer de mama, e o exame é uma mamografia. Neste exemplo, a probabilidade direta é a probabilidade de um teste positivo, dado que você tem a doença: \textit{P}(teste | doença). Isso é o que um médico chamaria de “sensibilidade” do teste, ou sua capacidade de detectar corretamente uma enfermidade. Em geral, ela é a mesma para todos os tipos de pacientes, porque depende apenas da capacidade técnica do instrumento de teste de detectar as anormalidades associadas à doença. A probabilidade inversa é a que você certamente mais se importa: qual a probabilidade de que eu tenha a doença, dado que o teste saiu positivo? Isso é \textit{P}(doença | teste), e representa um fluxo de informação na direção não causal, do resultado do teste para a probabilidade da doença. Essa probabilidade não é necessariamente a mesma para todos os tipos de pacientes; certamente veríamos o teste positivo com mais alarme num paciente com histórico familiar da doença do que num sem tal histórico.

Note que começamos a falar de direções causais e não causais. Não fizemos isso no exemplo da casa de chá porque não importava o que vinha primeiro, pedir chá ou pedir \textit{scones}. Só importava qual probabilidade condicional nos sentíamos mais capazes de avaliar. Mas o cenário causal esclarece por que nos sentimos menos confortáveis avaliando a “probabilidade inversa”, e o ensaio de Bayes deixa claro que esse é exatamente o tipo de problema que o interessava.

Suponha que uma mulher de quarenta anos faça uma mamografia para checar câncer de mama, e o resultado seja positivo. A hipótese, \textit{D} (de “doença”), é que ela tem câncer. A evidência, \textit{T} (de “teste”), é o resultado da mamografia. Com que força ela deveria crer na hipótese? Deveria se submeter a cirurgia?

Podemos responder a essas perguntas reescrevendo a regra de Bayes como segue:

(Probabilidade atualizada de \textit{D}) = \textit{P}(\textit{D} | \textit{T}) = (razão de verossimilhança) × (probabilidade \textit{a priori} de \textit{D})  (3.2)

em que o novo termo “razão de verossimilhança” (\textit{likelihood ratio}) é dado por \textit{P}(\textit{T} | \textit{D})/\textit{P}(\textit{T}). Ele mede quanto mais provável o teste positivo é em pessoas com a doença do que na população em geral. A Equação 3.2, portanto, nos diz que a nova evidência \textit{T} aumenta a probabilidade de \textit{D} por uma razão fixa, não importa qual fosse a probabilidade \textit{a priori}.

Façamos um exemplo para ver como esse conceito importante funciona. Para uma mulher típica de quarenta anos, a probabilidade de desenvolver câncer de mama no próximo ano é de cerca de uma em setecentos, de modo que usaremos isso como nossa probabilidade \textit{a priori}. Para computar a razão de verossimilhança, precisamos conhecer \textit{P}(\textit{T} | \textit{D}) e \textit{P}(\textit{T}). No contexto médico, \textit{P}(\textit{T} | \textit{D}) é a sensibilidade da mamografia — a probabilidade de que ela dê positivo se você tem câncer. Segundo o Breast Cancer Surveillance Consortium (BCSC), a sensibilidade das mamografias para mulheres de quarenta anos é de 73 por cento. O denominador, \textit{P}(\textit{T}), é um pouco mais traiçoeiro. Um teste positivo, \textit{T}, pode vir tanto de pacientes que têm a doença quanto de pacientes que não a têm. Assim, \textit{P}(\textit{T}) deveria ser uma média ponderada de \textit{P}(\textit{T} | \textit{D}) (a probabilidade de um teste positivo entre os que têm a doença) e \textit{P}(\textit{T} | ~\textit{D}) (a probabilidade de um teste positivo entre os que não a têm). A segunda é conhecida como taxa de falsos positivos. Segundo o BCSC, a taxa de falsos positivos para mulheres de quarenta anos é de cerca de 12 por cento.

Por que uma média ponderada? Porque há bem mais mulheres saudáveis (~\textit{D}) do que mulheres com câncer (\textit{D}). De fato, apenas 1 em 700 mulheres tem câncer, e as outras 699 não têm, de modo que a probabilidade de um teste positivo para uma mulher escolhida aleatoriamente deveria ser bem mais fortemente influenciada pelas 699 mulheres que não têm câncer do que pela uma mulher que tem. Matematicamente, computamos a média ponderada como segue:

\textit{P}(\textit{T}) = (1/700) × (73 por cento) + (699/700) × (12 por cento) ≈ 12,1 por cento.

Os pesos surgem porque apenas 1 em 700 mulheres tem 73 por cento de chance de um teste positivo, e as outras 699 têm 12 por cento de chance. Como se poderia esperar, \textit{P}(\textit{T}) saiu muito próxima da taxa de falsos positivos.

Agora que conhecemos \textit{P}(\textit{T}), finalmente podemos computar a probabilidade atualizada — as chances da mulher de ter câncer de mama depois que o teste dá positivo. A razão de verossimilhança é 73 por cento / 12,1 por cento ≈ 6. Como disse antes, este é o fator pelo qual aumentamos sua probabilidade \textit{a priori} para computar sua probabilidade atualizada de ter câncer. Como sua probabilidade \textit{a priori} era de uma em setecentos, sua probabilidade atualizada é 6 × 1/700 ≈ 1/116. Em outras palavras, ela ainda tem menos de 1 por cento de chance de ter câncer.

A conclusão é surpreendente. Penso que a maioria das mulheres de quarenta anos que têm uma mamografia positiva ficaria estupefata ao saber que ainda tem menos de 1 por cento de chance de ter câncer de mama. A Figura 3.3 pode tornar a razão mais fácil de compreender: o minúsculo número de verdadeiros positivos (isto é, mulheres com câncer de mama) é esmagado pelo número de falsos positivos. Nosso senso de surpresa com esse resultado vem da confusão cognitiva comum entre a probabilidade direta, que é bem estudada e minuciosamente documentada, e a probabilidade inversa, que é necessária para a tomada de decisão pessoal.

\begin{figure}[ht]
\centering
\includegraphics[width=0.92\textwidth]{../figuras/fig_3_3.png}
\caption{Figura 3.3. pode tornar a razão mais fácil de compreender: o minúsculo número de verdadeiros positivos (isto é, mulheres com câncer de mama) é esmagado pelo número de falsos positivos. Nosso senso de surpresa com esse resultado vem da confusão cognitiva comum entre a probabilidade direta, que é bem estudada e minuciosamente documentada, e a probabilidade inversa, que é necessária para a tomada de decisão pessoal.}
\end{figure}
\textit{Figura 3.3. pode tornar a razão mais fácil de compreender: o minúsculo número de verdadeiros positivos (isto é, mulheres com câncer de mama) é esmagado pelo número de falsos positivos. Nosso senso de surpresa com esse resultado vem da confusão cognitiva comum entre a probabilidade direta, que é bem estudada e minuciosamente documentada, e a probabilidade inversa, que é necessária para a tomada de decisão pessoal.}

O conflito entre nossa percepção e a realidade explica em parte o clamor quando a US Preventive Services Task Force, em 2009, recomendou que mulheres de quarenta anos não fizessem mamografias anuais. A força-tarefa compreendia o que muitas mulheres não compreendiam: um teste positivo nessa idade é bem mais provavelmente um alarme falso do que a detecção de câncer, e muitas mulheres estavam sendo desnecessariamente aterrorizadas (e recebendo tratamento desnecessário) como resultado.

\textbf{Figura 3.3.} Neste exemplo, baseado em taxas de falsos positivos e falsos negativos fornecidas pelo Breast Cancer Surveillance Consortium, apenas 3 em 363 mulheres de quarenta anos que testam positivo para câncer de mama de fato têm a doença. (As proporções não batem exatamente com o texto por causa de arredondamentos.) (Fonte: infográfico de Maayan Harel.)

Contudo, a história seria bem diferente se nossa paciente tivesse um gene que a colocasse em alto risco de câncer de mama — digamos, uma chance de uma em vinte no próximo ano. Então um teste positivo aumentaria a probabilidade para quase uma em três. Para uma mulher nessa situação, as chances de que o teste forneça informação salva-vidas são bem maiores. É por isso que a força-tarefa continuou a recomendar mamografias anuais para mulheres de alto risco.

Este exemplo mostra que \textit{P}(doença | teste) não é a mesma para todos; depende do contexto. Se você sabe que está em alto risco de uma doença desde o início, a regra de Bayes permite incorporar essa informação. Ou se você sabe que é imune, nem precisa se dar ao trabalho de fazer o teste! Em contraste, \textit{P}(teste | doença) não depende de você estar em alto risco ou não. É “robusta” a tais variações, o que explica em alguma medida por que os médicos organizam seu conhecimento e se comunicam com probabilidades diretas. Estas são propriedades da doença em si, de seu estágio de progressão, ou da sensibilidade dos instrumentos de detecção; daí permanecerem relativamente invariantes às razões da doença (epidemia, dieta, higiene, status socioeconômico, histórico familiar). A probabilidade inversa, \textit{P}(doença | teste), é sensível a essas condições.

O leitor com mentalidade histórica certamente se perguntará como Bayes lidou com a subjetividade de \textit{P}(\textit{L}), em que \textit{L} é o comprimento de uma mesa de bilhar. A resposta tem duas partes. Primeiro, Bayes não estava interessado no comprimento da mesa em si, mas em suas consequências futuras (isto é, a probabilidade de que a próxima bola parasse em alguma localização especificada na mesa). Segundo, Bayes assumiu que \textit{L} é determinado mecanicamente ao se disparar uma bola de bilhar de uma distância maior, digamos \textit{L}\textbackslash\{\}\textit{. Desse modo, ele conferiu objetividade a }P\textit{(}L\textit{) e transformou o problema num em que as probabilidades }a priori* são estimáveis a partir de dados, como vemos nos exemplos da casa de chá e do teste de câncer.

De muitas maneiras, a regra de Bayes é uma destilação do método científico. A descrição de livro-texto do método científico é mais ou menos assim: (1) formular uma hipótese, (2) deduzir uma consequência testável da hipótese, (3) realizar um experimento e coletar evidências, e (4) atualizar sua crença na hipótese. Em geral, os livros-texto lidam com testes e atualizações simples de sim-ou-não; a evidência ou confirma ou refuta a hipótese. Mas a vida e a ciência nunca são tão simples! Toda evidência vem com uma certa quantidade de incerteza. A regra de Bayes nos diz como realizar o passo (4) no mundo real.

\section{Da regra de Bayes às redes bayesianas}
No início dos anos 1980, o campo da inteligência artificial se metera num beco sem saída. Desde que Alan Turing pela primeira vez estabeleceu o desafio em seu artigo de 1950 “Computing Machinery and Intelligence” [“Maquinaria computacional e inteligência”], a abordagem dominante à IA havia sido os chamados sistemas baseados em regras ou sistemas especialistas, que organizam o conhecimento humano como uma coleção de fatos específicos e gerais, juntamente com regras de inferência para conectá-los. Por exemplo:

Sócrates é um homem (fato específico). Todos os homens são mortais (fato geral).

A partir dessa base de conhecimento, nós (ou uma máquina inteligente) podemos derivar o fato de que Sócrates é mortal, usando a regra universal de inferência: se todos os \textit{A}s são \textit{B}s, e \textit{x} é um \textit{A}, então \textit{x} é um \textit{B}.

A abordagem era boa em teoria, mas regras rígidas raramente capturam o conhecimento da vida real. Talvez sem nos darmos conta, lidamos com exceções a regras e incertezas em evidências o tempo todo. Em 1980, estava claro que os sistemas especialistas lutavam para fazer inferências corretas a partir de conhecimento incerto. O computador não conseguia replicar o processo inferencial de um especialista humano porque os próprios especialistas não eram capazes de articular seu processo de pensamento na linguagem fornecida pelo sistema.

O final dos anos 1970, então, foi um tempo de fermentação na comunidade de IA em torno da pergunta de como lidar com a incerteza. Não faltavam ideias. Lotfi Zadeh, de Berkeley, ofereceu a “lógica \textit{fuzzy}”, em que as afirmações não são nem verdadeiras nem falsas, mas em vez disso assumem um espectro de possíveis valores de verdade. Glen Shafer, da Universidade do Kansas, propôs “funções de crença”, que atribuem duas probabilidades a cada fato, uma indicando quão provável é que seja “possível”, a outra, quão provável é que seja “provável”. Edward Feigenbaum e seus colegas na Universidade de Stanford tentaram “fatores de certeza”, que inseriam medidas numéricas de incerteza em suas regras determinísticas de inferência.

Infelizmente, embora engenhosas, essas abordagens sofriam de um defeito comum: modelavam o especialista, não o mundo, e, portanto, tendiam a produzir resultados indesejados. Por exemplo, não conseguiam operar tanto em modos diagnósticos quanto preditivos, a especialidade incontestada da regra de Bayes. Na abordagem dos fatores de certeza, a regra “Se fogo, então fumaça (com certeza \textit{c}₁)” não podia se combinar de modo coerente com “Se fumaça, então fogo (com certeza \textit{c}₂)” sem desencadear um acúmulo descontrolado de crença. A probabilidade também foi considerada na época, mas imediatamente caiu em descrédito, pois as demandas de espaço de armazenamento e tempo de processamento se tornavam formidáveis.

Entrei na arena bastante tarde, em 1982, com uma proposta óbvia e no entanto radical: em vez de reinventar do zero uma nova teoria da incerteza, mantemos a probabilidade como guardiã do senso comum e meramente reparamos suas deficiências computacionais. Mais especificamente, em vez de representar a probabilidade em tabelas enormes, como se fazia antes, vamos representá-la com uma rede de variáveis frouxamente acopladas. Se só permitirmos que cada variável interaja com algumas variáveis vizinhas, então talvez superemos os obstáculos computacionais que haviam feito outros probabilistas tropeçar.

A ideia não me veio num sonho; veio de um artigo de David Rumelhart, um cientista cognitivo da Universidade da Califórnia em San Diego e pioneiro das redes neurais. Seu artigo sobre a leitura de crianças, publicado em 1976, deixava claro que a leitura é um processo complexo em que neurônios em muitos níveis diferentes estão ativos ao mesmo tempo (ver Figura 3.4). Alguns dos neurônios estão simplesmente reconhecendo traços individuais — círculos ou linhas. Acima deles, outra camada de neurônios está combinando essas formas e formando conjecturas sobre o que a letra poderia ser. Na Figura 3.4, a rede está lutando com uma grande ambiguidade sobre a segunda palavra. No nível da letra, poderia ser “FHP”, mas isso não faz muito sentido no nível da palavra. No nível da palavra, poderia ser “FAR” ou “CAR” ou “FAT”. Os neurônios passam essa informação para cima, ao nível sintático, que decide que, depois da palavra “THE”, espera um substantivo. Por fim, essa informação sobe todo o caminho até o nível semântico, que percebe que a sentença anterior mencionou um Volkswagen, de modo que a frase provavelmente é “THE CAR”, referindo-se a esse mesmo Volkswagen.

\begin{figure}[ht]
\centering
\includegraphics[width=0.92\textwidth]{../figuras/fig_3_4.png}
\caption{Figura 3.4. ). Alguns dos neurônios estão simplesmente reconhecendo traços individuais — círculos ou linhas. Acima deles, outra camada de neurônios está combinando essas formas e formando conjecturas sobre o que a letra poderia ser. Na Figura 3.4, a rede está lutando com uma grande ambiguidade sobre a segunda palavra. No nível da letra, poderia ser “FHP”, mas isso não faz muito sentido no nível da palavra. No nível da palavra, poderia ser “FAR” ou “CAR” ou “FAT”. Os neurônios passam essa informação para cima, ao nível sintático, que decide que, depois da palavra “THE”, espera um substantivo. Por fim, essa informação sobe todo o caminho até o nível semântico, que percebe que a sentença anterior mencionou um Volkswagen, de modo que a frase provavelmente é “THE CAR”, referindo-se a esse mesmo Volkswagen.}
\end{figure}
\textit{Figura 3.4. ). Alguns dos neurônios estão simplesmente reconhecendo traços individuais — círculos ou linhas. Acima deles, outra camada de neurônios está combinando essas formas e formando conjecturas sobre o que a letra poderia ser. Na Figura 3.4, a rede está lutando com uma grande ambiguidade sobre a segunda palavra. No nível da letra, poderia ser “FHP”, mas isso não faz muito sentido no nível da palavra. No nível da palavra, poderia ser “FAR” ou “CAR” ou “FAT”. Os neurônios passam essa informação para cima, ao nível sintático, que decide que, depois da palavra “THE”, espera um substantivo. Por fim, essa informação sobe todo o caminho até o nível semântico, que percebe que a sentença anterior mencionou um Volkswagen, de modo que a frase provavelmente é “THE CAR”, referindo-se a esse mesmo Volkswagen.}

O ponto-chave é que todos os neurônios estão passando informação de um lado para o outro, de cima para baixo e de baixo para cima e de lado a lado. É um sistema altamente paralelo, e um que é bem diferente de nossa autopercepção do cérebro como um sistema monolítico e controlado de modo central.

Lendo o artigo de Rumelhart, senti-me convencido de que qualquer inteligência artificial teria de se modelar no que sabemos sobre o processamento neural humano de informação e de que o raciocínio de máquina sob incerteza teria de ser construído com uma arquitetura semelhante de passagem de mensagens. Mas o que são as mensagens? Isso me tomou bastantes meses para descobrir. Finalmente percebi que as mensagens eram probabilidades condicionais numa direção e razões de verossimilhança na outra.

\textbf{Figura 3.4.} Esboço de David Rumelhart de como uma rede de passagem de mensagens aprenderia a ler a frase “THE CAR”. (Fonte: cortesia do Center for Brain and Cognition, Universidade da Califórnia em San Diego.)

Mais precisamente, assumi que a rede seria hierárquica, com setas apontando de neurônios superiores para inferiores, ou de “nós pais” para “nós filhos”. Cada nó enviaria uma mensagem a todos os seus vizinhos (tanto acima quanto abaixo na hierarquia) sobre seu grau atual de crença na variável que rastreava (p.ex., “Estou dois terços certo de que esta letra é um R”). O receptor processaria a mensagem de dois modos diferentes, dependendo de sua direção. Se a mensagem ia do pai para o filho, o filho atualizaria suas crenças usando probabilidades condicionais, como as que vimos no exemplo da casa de chá. Se a mensagem ia do filho para o pai, o pai atualizaria suas crenças multiplicando-as por uma razão de verossimilhança, como no exemplo da mamografia. Aplicar essas duas regras repetidamente a cada nó da rede chama-se propagação de crenças.

Em retrospecto, não há nada de arbitrário ou inventado nessas regras; elas estão em estrita conformidade com a regra de Bayes. O verdadeiro desafio era assegurar que, não importa em que ordem essas mensagens fossem enviadas, as coisas se assentassem por fim num equilíbrio confortável; mais ainda, que o equilíbrio final representasse o estado correto de crença nas variáveis. Por “correto” quero dizer como se tivéssemos conduzido a computação por métodos de livro-texto em vez de por passagem de mensagens. Esse desafio ocuparia meus alunos e a mim, bem como meus colegas, por vários anos. Mas, no final dos anos 1980, havíamos resolvido as dificuldades a ponto de as redes bayesianas se tornarem um esquema prático para o aprendizado de máquina. A década seguinte viu um aumento contínuo de aplicações no mundo real, como filtragem de \textit{spam} e reconhecimento de voz. Contudo, por essa época eu já estava tentando subir a Escada da Causalidade, enquanto confiei o lado probabilístico das redes bayesianas à custódia de outros.

\section{Redes bayesianas: o que as causas dizem sobre os dados}
Embora Bayes não o soubesse, sua regra para a probabilidade inversa representa a mais simples rede bayesiana. Vimos essa rede sob várias roupagens: Chá → \textit{Scones}, Doença → Teste ou, mais geralmente, Hipótese → Evidência. Ao contrário dos diagramas causais com que lidaremos ao longo do livro, uma rede bayesiana não carrega nenhuma suposição de que a seta tenha qualquer significado causal. A seta meramente significa que conhecemos a probabilidade “direta”, \textit{P}(\textit{scones} | chá) ou \textit{P}(teste | doença). A regra de Bayes nos diz como reverter o procedimento, especificamente multiplicando a probabilidade \textit{a priori} por uma razão de verossimilhança.

A propagação de crenças funciona formalmente exatamente do mesmo modo, quer as setas sejam não causais ou causais. Não obstante, você pode ter a sensação intuitiva de que fizemos algo mais significativo no segundo caso do que no primeiro. Isso porque nossos cérebros são dotados de uma maquinaria especial para compreender relações de causa e efeito (como câncer e mamografias). Não o mesmo para meras associações (como chá e \textit{scones}).

O passo seguinte após uma rede de dois nós com um elo é, é claro, uma rede de três nós com dois elos, a que chamarei de “junção”. Estas são os blocos de construção de todas as redes bayesianas (e das redes causais também). Há três tipos básicos de junções, com a ajuda dos quais podemos caracterizar qualquer padrão de setas na rede.

\textbf{1.} \textit{A} → \textit{B} → \textit{C}. Esta junção é o exemplo mais simples de uma “cadeia”, ou de mediação. Na ciência, muitas vezes se pensa em \textit{B} como o mecanismo, ou “mediador”, que transmite o efeito de \textit{A} a \textit{C}. Um exemplo familiar é Fogo → Fumaça → Alarme. Embora os chamemos de “alarmes de incêndio”, são na verdade alarmes de fumaça. O fogo por si só não dispara um alarme, de modo que não há seta direta de Fogo para Alarme. Tampouco o fogo dispara o alarme por meio de qualquer outra variável, como o calor. Funciona apenas liberando moléculas de fumaça no ar. Se desabilitarmos esse elo da cadeia, por exemplo sugando todas as moléculas de fumaça com uma capela de exaustão, então não haverá alarme.

Essa observação conduz a um ponto conceitual importante sobre cadeias: o mediador \textit{B} “bloqueia” (\textit{screens off}) a informação sobre \textit{A} de \textit{C}, e vice-versa. (Isso foi apontado pela primeira vez por Hans Reichenbach, um filósofo da ciência germano-americano.) Por exemplo, uma vez que conhecemos o valor de Fumaça, aprender sobre Fogo não nos dá nenhuma razão para elevar ou baixar nossa crença em Alarme. Essa estabilidade de crença é um conceito de degrau um; daí deveria também ser vista nos dados, quando disponíveis. Suponha que tivéssemos um banco de dados de todas as instâncias em que houve fogo, em que houve fumaça ou em que o alarme disparou. Se olhássemos apenas as linhas em que Fumaça = 1, esperaríamos Alarme = 1 toda vez, independentemente de Fogo = 0 ou Fogo = 1.

Esse padrão de bloqueio ainda vale se o efeito não for determinístico. Por exemplo, imagine um sistema de alarme defeituoso que falha em responder corretamente 5 por cento do tempo. Se olharmos apenas as linhas em que Fumaça = 1, descobriremos que a probabilidade de Alarme = 1 é a mesma (95 por cento), independentemente de Fogo = 0 ou Fogo = 1.

O processo de olhar apenas as linhas da tabela em que Fumaça = 1 chama-se condicionar a uma variável. Da mesma forma, dizemos que Fogo e Alarme são condicionalmente independentes, dado o valor de Fumaça. Isso é importante de saber se você está programando uma máquina para atualizar suas crenças; a independência condicional dá à máquina a licença para se concentrar na informação relevante e desconsiderar o resto. Todos precisamos desse tipo de licença em nosso pensamento cotidiano, ou então passaremos todo o tempo perseguindo sinais falsos. Mas como decidimos que informação desconsiderar, quando cada novo pedaço de informação muda a fronteira entre o relevante e o irrelevante? Para os humanos, essa compreensão vem naturalmente. Até crianças de três anos compreendem o efeito de bloqueio, embora não tenham um nome para ele. Seu instinto deve ter vindo de alguma representação mental, possivelmente semelhante a um diagrama causal. Mas as máquinas não têm esse instinto, o que é uma razão pela qual as equipamos com diagramas causais.

\textbf{2.} \textit{A} ← \textit{B} → \textit{C}. Esse tipo de junção chama-se “garfo”, e \textit{B} é frequentemente chamado de causa comum ou confundidor de \textit{A} e \textit{C}. Um confundidor fará \textit{A} e \textit{C} ficarem estatisticamente correlacionados mesmo que não haja elo causal direto entre eles. Um bom exemplo (devido a David Freedman) é Tamanho do Sapato ← Idade da Criança → Capacidade de Leitura. Crianças com sapatos maiores tendem a ler em um nível mais alto. Mas a relação não é de causa e efeito. Dar a uma criança sapatos maiores não a fará ler melhor! Em vez disso, ambas as variáveis são explicadas por uma terceira, que é a idade da criança. Crianças mais velhas têm sapatos maiores, e também são leitoras mais avançadas.

Podemos eliminar essa correlação espúria, como Karl Pearson e George Udny Yule a chamaram, condicionando à idade da criança. Por exemplo, se olharmos apenas crianças de sete anos, esperamos ver nenhuma relação entre tamanho do sapato e capacidade de leitura. Como no caso das junções em cadeia, \textit{A} e \textit{C} são condicionalmente independentes, dado \textit{B}.

Antes de passarmos à nossa terceira junção, precisamos acrescentar uma palavra de esclarecimento. As independências condicionais que acabei de mencionar se exibem sempre que olhamos essas junções isoladamente. Se caminhos causais adicionais as cercam, esses caminhos também precisam ser levados em conta. O milagre das redes bayesianas reside no fato de que os três tipos de junções que estamos agora descrevendo isoladamente são suficientes para ler todas as independências implicadas por uma rede bayesiana, por mais complicada que seja.

\textbf{3.} \textit{A} → \textit{B} ← \textit{C}. Esta é a junção mais fascinante, chamada de “colisor”. Felix Elwert e Chris Winship ilustraram essa junção usando três traços de atores de Hollywood: Talento → Celebridade ← Beleza. Aqui estamos afirmando que tanto o talento quanto a beleza contribuem para o sucesso de um ator, mas beleza e talento são completamente não relacionados um ao outro na população em geral.

Veremos agora que esse padrão de colisor funciona exatamente do modo oposto ao das cadeias ou dos garfos quando condicionamos à variável do meio. Se \textit{A} e \textit{C} são independentes no início, condicionar a \textit{B} os tornará dependentes. Por exemplo, se olharmos apenas atores famosos (em outras palavras, observamos a variável Celebridade = 1), veremos uma correlação negativa entre talento e beleza: descobrir que uma celebridade é pouco atraente aumenta nossa crença de que ela é talentosa.

Essa correlação negativa às vezes é chamada de viés de colisor ou efeito de “explicar-afastar” (\textit{explain-away}). Por simplicidade, suponha que você não precise de talento e beleza juntos para ser uma celebridade; um basta. Então, se a Celebridade \textit{A} é um ator particularmente bom, isso “explica-afasta” seu sucesso, e ele não precisa ser mais belo do que a pessoa média. Por outro lado, se a Celebridade \textit{B} é um ator realmente mau, então o único modo de explicar seu sucesso é sua boa aparência. Assim, dado o resultado Celebridade = 1, talento e beleza estão inversamente relacionados — embora não estejam relacionados na população como um todo. Mesmo numa situação mais realista, em que o sucesso é uma função complicada de beleza e talento, o efeito de explicar-afastar ainda estará presente.

Este exemplo é, admitidamente, um tanto apócrifo, porque beleza e talento são difíceis de medir objetivamente; não obstante, o viés de colisor é bem real, e veremos muitos exemplos neste livro.

Essas três junções — cadeias, garfos e colisores — são como buracos de fechadura na porta que separa o primeiro e o segundo níveis da Escada da Causalidade. Se espiarmos por elas, podemos ver os segredos do processo causal que gerou os dados que observamos; cada uma representa um padrão distinto de fluxo causal e deixa sua marca na forma de dependências e independências condicionais nos dados. Em minhas palestras públicas, muitas vezes as chamo de “dádivas dos deuses” porque nos permitem testar um modelo causal, descobrir novos modelos, avaliar efeitos de intervenções e muito mais. Ainda assim, isoladas, elas nos dão apenas um vislumbre. Precisamos de uma chave que abrirá completamente a porta e nos permitirá sair ao segundo degrau. Essa chave, que aprenderemos no Capítulo 7, envolve todas as três junções e se chama \textit{d}-separação. Esse conceito nos diz, para qualquer padrão dado de caminhos no modelo, que padrões de dependências deveríamos esperar nos dados. Essa conexão fundamental entre causas e probabilidades constitui a principal contribuição das redes bayesianas à ciência da inferência causal.

\section{Onde está minha mala? De Aachen a Zanzibar}
Até aqui enfatizei apenas um aspecto das redes bayesianas — a saber, o diagrama e suas setas que preferencialmente apontam da causa para o efeito. De fato, o diagrama é como o motor da rede bayesiana. Mas, como qualquer motor, uma rede bayesiana funciona com combustível. O combustível se chama tabela de probabilidade condicional. Outro modo de dizer isso é que o diagrama descreve a relação das variáveis de um modo qualitativo, mas se você quer respostas quantitativas, também precisa de entradas quantitativas.

Numa rede bayesiana, temos de especificar a probabilidade condicional de cada nó dados seus “pais”. (Lembre-se de que os pais de um nó são todos os nós que alimentam nele.) Estas são as probabilidades diretas, \textit{P}(evidência | hipóteses). No caso em que \textit{A} é um nó raiz, sem setas apontando para ele, precisamos apenas especificar a probabilidade \textit{a priori} para cada estado de \textit{A}. Em nossa segunda rede, Doença → Teste, Doença é um nó raiz. Portanto, especificamos a probabilidade \textit{a priori} de que uma pessoa tenha a doença (1/700 em nosso exemplo) e de que ela não a tenha (699/700 em nosso exemplo).

Ao descrever \textit{A} como um nó raiz, não queremos realmente dizer que \textit{A} não tem causas prévias. Quase nenhuma variável tem direito a tal status. Queremos realmente dizer que quaisquer causas prévias de \textit{A} podem ser adequadamente resumidas na probabilidade \textit{a priori} \textit{P}(\textit{A}) de que \textit{A} seja verdadeiro. Por exemplo, no exemplo Doença → Teste, o histórico familiar poderia ser uma causa de Doença. Mas, desde que estejamos certos de que esse histórico familiar não afetará a variável Teste (uma vez que conheçamos o status de Doença), não precisamos representá-lo como um nó no grafo. Contudo, se houver uma causa de Doença que também afete diretamente Teste, então essa causa deve ser representada explicitamente no diagrama.

No caso em que o nó \textit{A} tem um pai, \textit{A} tem de “ouvir” seu pai antes de decidir seu próprio estado. Em nosso exemplo da mamografia, o pai de Teste era Doença. Podemos mostrar esse processo de “escuta” com uma tabela 2 × 2 (ver Tabela 3.2). Por exemplo, se Teste “ouve” que \textit{D} = 0, então 88 por cento do tempo tomará o valor \textit{T} = 0, e 12 por cento do tempo tomará o valor \textit{T} = 1. Note que a segunda coluna dessa tabela contém a mesma informação que vimos antes do Breast Cancer Surveillance Consortium: a taxa de falsos positivos (canto superior direito) é 12 por cento, e a sensibilidade (canto inferior direito) é 73 por cento. As duas entradas restantes são preenchidas para que cada linha some 100 por cento.

\begin{figure}[ht]
\centering
\includegraphics[width=0.92\textwidth]{../figuras/tab_3_2.png}
\caption{Tabela 3.2. ). Por exemplo, se Teste “ouve” que \textit{D} = 0, então 88 por cento do tempo tomará o valor \textit{T} = 0, e 12 por cento do tempo tomará o valor \textit{T} = 1. Note que a segunda coluna dessa tabela contém a mesma informação que vimos antes do Breast Cancer Surveillance Consortium: a taxa de falsos positivos (canto superior direito) é 12 por cento, e a sensibilidade (canto inferior direito) é 73 por cento. As duas entradas restantes são preenchidas para que cada linha some 100 por cento.}
\end{figure}
\textit{Tabela 3.2. ). Por exemplo, se Teste “ouve” que }D\textit{ = 0, então 88 por cento do tempo tomará o valor }T\textit{ = 0, e 12 por cento do tempo tomará o valor }T\textit{ = 1. Note que a segunda coluna dessa tabela contém a mesma informação que vimos antes do Breast Cancer Surveillance Consortium: a taxa de falsos positivos (canto superior direito) é 12 por cento, e a sensibilidade (canto inferior direito) é 73 por cento. As duas entradas restantes são preenchidas para que cada linha some 100 por cento.}

\textbf{Tabela 3.2.} Uma tabela de probabilidade condicional simples.

À medida que passamos a redes mais complicadas, a tabela de probabilidade condicional fica igualmente mais complicada. Por exemplo, se temos um nó com dois pais, a tabela de probabilidade condicional tem de levar em conta os quatro estados possíveis de ambos os pais. Vejamos um exemplo concreto, sugerido por Stefan Conrady e Lionel Jouffe da BayesiaLab, Inc. É um cenário familiar a todos os viajantes: podemos chamá-lo de “Onde está minha mala?”.

Suponha que você acabou de pousar em Zanzibar depois de fazer uma conexão apertada em Aachen, e está esperando que sua mala apareça na esteira. Outros passageiros começaram a pegar suas malas, mas você continua esperando… e esperando… e esperando. Quais as chances de que sua mala de fato não tenha feito a conexão de Aachen a Zanzibar?

A resposta depende, é claro, de quanto tempo você tem esperado. Se as malas acabaram de começar a aparecer na esteira, talvez você devesse ser paciente e esperar um pouco mais. Se você tem esperado um longo tempo, então as coisas estão feias. Podemos quantificar essas ansiedades montando um diagrama causal (Figura 3.5).

\begin{figure}[ht]
\centering
\includegraphics[width=0.92\textwidth]{../figuras/fig_3_5.png}
\caption{Figura 3.5. Diagrama causal para o exemplo aeroporto/mala.}
\end{figure}
\textit{Figura 3.5. Diagrama causal para o exemplo aeroporto/mala.}

\textbf{Figura 3.5.} Diagrama causal para o exemplo aeroporto/mala.

Esse diagrama reflete a ideia intuitiva de que há duas causas para o aparecimento de qualquer mala na esteira. Primeiro, ela tinha de estar no avião desde o início; caso contrário, certamente nunca aparecerá na esteira. Segundo, a presença da mala na esteira se torna mais provável à medida que o tempo passa… desde que ela estivesse de fato no avião.

Para transformar o diagrama causal numa rede bayesiana, temos de especificar as tabelas de probabilidade condicional. Digamos que todas as malas no aeroporto de Zanzibar são descarregadas em dez minutos. (Eles são muito eficientes em Zanzibar!) Suponhamos também que a probabilidade de que sua mala tenha feito a conexão, \textit{P}(mala no avião = verdadeiro), seja 50 por cento. (Peço desculpas se isso ofende alguém que trabalha no aeroporto de Aachen. Estou apenas seguindo o exemplo de Conrady e Jouffe. Pessoalmente, preferiria assumir uma probabilidade \textit{a priori} mais alta, como 95 por cento.)

O verdadeiro motor dessa rede bayesiana é a tabela de probabilidade condicional para “Mala na Esteira” (ver Tabela 3.3). Essa tabela, embora grande, deveria ser fácil de compreender. As primeiras onze linhas dizem que, se sua mala não entrou no avião (mala no avião = falso), então, não importa quanto tempo tenha decorrido, ela não estará na esteira (esteira = falso). Isto é, \textit{P}(esteira = falso | mala no avião = falso) é 100 por cento. Esse é o significado dos 100s nas primeiras onze linhas. As outras onze linhas dizem que as malas são descarregadas do avião a uma taxa constante. Se sua mala está de fato no avião, há 10 por cento de probabilidade de que seja descarregada no primeiro minuto, 10 por cento de probabilidade no segundo minuto, e assim por diante. Por exemplo, após 5 minutos há 50 por cento de probabilidade de que tenha sido descarregada, de modo que vemos um 50 para \textit{P}(esteira = verdadeiro | mala no avião = verdadeiro, tempo = 5). Após dez minutos, todas as malas foram descarregadas, de modo que \textit{P}(esteira = verdadeiro | mala no avião = verdadeiro, tempo = 10) é 100 por cento. Assim, vemos um 100 na última entrada da tabela.

\begin{figure}[ht]
\centering
\includegraphics[width=0.92\textwidth]{../figuras/tab_3_3.png}
\caption{Tabela 3.3. ). Essa tabela, embora grande, deveria ser fácil de compreender. As primeiras onze linhas dizem que, se sua mala não entrou no avião (mala no avião = falso), então, não importa quanto tempo tenha decorrido, ela não estará na esteira (esteira = falso). Isto é, \textit{P}(esteira = falso | mala no avião = falso) é 100 por cento. Esse é o significado dos 100s nas primeiras onze linhas. As outras onze linhas dizem que as malas são descarregadas do avião a uma taxa constante. Se sua mala está de fato no avião, há 10 por cento de probabilidade de que seja descarregada no primeiro minuto, 10 por cento de probabilidade no segundo minuto, e assim por diante. Por exemplo, após 5 minutos há 50 por cento de probabilidade de que tenha sido descarregada, de modo que vemos um 50 para \textit{P}(esteira = verdadeiro | mala no avião = verdadeiro, tempo = 5). Após dez minutos, todas as malas foram descarregadas, de modo que \textit{P}(esteira = verdadeiro | mala no avião = verdadeiro, tempo = 10) é 100 por cento. Assim, vemos um 100 na última entrada da tabela.}
\end{figure}
\textit{Tabela 3.3. ). Essa tabela, embora grande, deveria ser fácil de compreender. As primeiras onze linhas dizem que, se sua mala não entrou no avião (mala no avião = falso), então, não importa quanto tempo tenha decorrido, ela não estará na esteira (esteira = falso). Isto é, }P\textit{(esteira = falso | mala no avião = falso) é 100 por cento. Esse é o significado dos 100s nas primeiras onze linhas. As outras onze linhas dizem que as malas são descarregadas do avião a uma taxa constante. Se sua mala está de fato no avião, há 10 por cento de probabilidade de que seja descarregada no primeiro minuto, 10 por cento de probabilidade no segundo minuto, e assim por diante. Por exemplo, após 5 minutos há 50 por cento de probabilidade de que tenha sido descarregada, de modo que vemos um 50 para }P\textit{(esteira = verdadeiro | mala no avião = verdadeiro, tempo = 5). Após dez minutos, todas as malas foram descarregadas, de modo que }P\textit{(esteira = verdadeiro | mala no avião = verdadeiro, tempo = 10) é 100 por cento. Assim, vemos um 100 na última entrada da tabela.}

A coisa mais interessante a fazer com essa rede bayesiana, como com a maioria das redes bayesianas, é resolver o problema da probabilidade inversa: se \textit{x} minutos se passaram e eu ainda não peguei minha mala, qual a probabilidade de que ela estivesse no avião? A regra de Bayes automatiza essa computação e revela um padrão interessante. Após um minuto, ainda há 47 por cento de chance de que estivesse no avião. (Lembre-se de que nossa suposição \textit{a priori} era de 50 por cento de probabilidade.) Após cinco minutos, a probabilidade cai para 33 por cento. Após dez minutos, é claro, cai a zero. A Figura 3.6 mostra um gráfico da probabilidade ao longo do tempo, a que se poderia chamar a “Curva do Abandono da Esperança”.

\begin{figure}[ht]
\centering
\includegraphics[width=0.92\textwidth]{../figuras/fig_3_6.png}
\caption{Figura 3.6. mostra um gráfico da probabilidade ao longo do tempo, a que se poderia chamar a “Curva do Abandono da Esperança”.}
\end{figure}
\textit{Figura 3.6. mostra um gráfico da probabilidade ao longo do tempo, a que se poderia chamar a “Curva do Abandono da Esperança”.}

Para mim, a coisa interessante é que é uma curva: penso que a maioria das pessoas esperaria que fosse uma reta. Na verdade, ela nos envia uma mensagem bastante otimista: não desista da esperança cedo demais! Segundo essa curva, você deveria abandonar apenas um terço de sua esperança na primeira metade do tempo alocado.

\textbf{Tabela 3.3.} Uma tabela de probabilidade condicional mais complicada.

\textbf{Figura 3.6.} A probabilidade de ver sua mala na esteira diminui lentamente no início, depois mais rapidamente. (Fonte: gráfico de Maayan Harel, dados de Stefan Conrady e Lionel Jouffe.)

Além de uma lição de vida, aprendemos que você não quer fazer isso à mão. Mesmo com essa minúscula rede de três nós, havia 2 × 11 = 22 estados dos pais, cada um contribuindo para a probabilidade do estado do filho. Para um computador, contudo, tais computações são elementares… até certo ponto. Se não forem feitas de modo organizado, o mero número de computações pode sobrecarregar até o supercomputador mais rápido. Se um nó tem dez pais, cada um dos quais tem dois estados, a tabela de probabilidade condicional terá mais de 1.000 linhas. E se cada um dos dez pais tem dez estados, a tabela terá 10 bilhões de linhas! Por essa razão, em geral se tem de peneirar as conexões na rede de modo que só as mais importantes permaneçam e a rede seja “esparsa”. Um avanço técnico no desenvolvimento das redes bayesianas consistiu em encontrar modos de aproveitar a esparsidade na estrutura da rede para alcançar tempos de computação razoáveis.

\section{Redes bayesianas no mundo real}
As redes bayesianas são agora uma tecnologia madura, e você pode comprar \textit{software} de redes bayesianas pronto de várias empresas. As redes bayesianas também estão embutidas em muitos dispositivos “inteligentes”. Para lhe dar uma ideia de como são usadas em aplicações do mundo real, voltemos ao \textit{software} de correspondência de DNA Bonaparte com que começamos este capítulo.

O Instituto Forense dos Países Baixos usa o Bonaparte todos os dias, sobretudo para casos de pessoas desaparecidas, investigações criminais e casos de imigração. (Candidatos a asilo devem provar que têm quinze membros da família nos Países Baixos.) Contudo, a rede bayesiana faz seu trabalho mais impressionante após um desastre maciço, como o acidente do voo 17 da Malaysia Airlines.

Poucas, se alguma, das vítimas do acidente aéreo puderam ser identificadas comparando DNA dos destroços com DNA num banco de dados central. A melhor coisa a fazer em seguida era pedir a membros da família que fornecessem swabs de DNA e procurar correspondências parciais com o DNA das vítimas. Métodos convencionais (não bayesianos) podem fazer isso e foram instrumentais na solução de vários casos arquivados nos Países Baixos, nos Estados Unidos e em outros lugares. Por exemplo, uma fórmula simples chamada “Índice de Paternidade” ou “Índice de Irmandade” pode estimar a verossimilhança de que o DNA não identificado venha do pai ou do irmão da pessoa cujo DNA foi testado. Contudo, esses índices são inerentemente limitados porque funcionam apenas para uma relação especificada e apenas para relações próximas.

A ideia por trás do Bonaparte é tornar possível usar informação de DNA de parentes mais distantes ou de múltiplos parentes. O Bonaparte faz isso convertendo o pedigree da família (ver Figura 3.7) numa rede bayesiana. Na Figura 3.8, vemos como o Bonaparte converte um pequeno pedaço de um pedigree numa rede bayesiana (causal). O problema central é que o genótipo de um indivíduo, detectado num teste de DNA, contém uma contribuição tanto do pai quanto da mãe, mas não podemos dizer qual parte é qual. Assim, essas duas contribuições (chamadas “alelos”) têm de ser tratadas como variáveis ocultas, não mensuráveis, na rede bayesiana. Parte do trabalho do Bonaparte é inferir a probabilidade da causa (o gene da vítima para olhos azuis veio do pai) a partir da evidência (p.ex., ele tem um gene de olhos azuis e um gene de olhos pretos; seus primos do lado do pai têm olhos azuis, mas seus primos do lado da mãe têm olhos pretos). Este é um problema de probabilidade inversa — exatamente para o que a regra de Bayes foi inventada.

\begin{figure}[ht]
\centering
\includegraphics[width=0.92\textwidth]{../figuras/fig_3_7.png}
\caption{Figura 3.7. ) numa rede bayesiana. Na Figura 3.8, vemos como o Bonaparte converte um pequeno pedaço de um pedigree numa rede bayesiana (causal). O problema central é que o genótipo de um indivíduo, detectado num teste de DNA, contém uma contribuição tanto do pai quanto da mãe, mas não podemos dizer qual parte é qual. Assim, essas duas contribuições (chamadas “alelos”) têm de ser tratadas como variáveis ocultas, não mensuráveis, na rede bayesiana. Parte do trabalho do Bonaparte é inferir a probabilidade da causa (o gene da vítima para olhos azuis veio do pai) a partir da evidência (p.ex., ele tem um gene de olhos azuis e um gene de olhos pretos; seus primos do lado do pai têm olhos azuis, mas seus primos do lado da mãe têm olhos pretos). Este é um problema de probabilidade inversa — exatamente para o que a regra de Bayes foi inventada.}
\end{figure}
\textit{Figura 3.7. ) numa rede bayesiana. Na Figura 3.8, vemos como o Bonaparte converte um pequeno pedaço de um pedigree numa rede bayesiana (causal). O problema central é que o genótipo de um indivíduo, detectado num teste de DNA, contém uma contribuição tanto do pai quanto da mãe, mas não podemos dizer qual parte é qual. Assim, essas duas contribuições (chamadas “alelos”) têm de ser tratadas como variáveis ocultas, não mensuráveis, na rede bayesiana. Parte do trabalho do Bonaparte é inferir a probabilidade da causa (o gene da vítima para olhos azuis veio do pai) a partir da evidência (p.ex., ele tem um gene de olhos azuis e um gene de olhos pretos; seus primos do lado do pai têm olhos azuis, mas seus primos do lado da mãe têm olhos pretos). Este é um problema de probabilidade inversa — exatamente para o que a regra de Bayes foi inventada.}

\begin{figure}[ht]
\centering
\includegraphics[width=0.92\textwidth]{../figuras/fig_3_8.png}
\caption{Figura 3.8. De testes de DNA a redes bayesianas. Na rede bayesiana, nós sem sombreamento representam alelos, e nós sombreados representam genótipos. Dados só estão disponíveis nos nós sombreados porque os genótipos não podem indicar qual alelo veio do pai e qual da mãe. A rede bayesiana possibilita a inferência sobre os nós não observados e também nos permite estimar a verossimilhança de que uma dada amostra de DNA tenha vindo da criança. (Fonte: infográfico de Maayan Harel.)}
\end{figure}
\textit{Figura 3.8. De testes de DNA a redes bayesianas. Na rede bayesiana, nós sem sombreamento representam alelos, e nós sombreados representam genótipos. Dados só estão disponíveis nos nós sombreados porque os genótipos não podem indicar qual alelo veio do pai e qual da mãe. A rede bayesiana possibilita a inferência sobre os nós não observados e também nos permite estimar a verossimilhança de que uma dada amostra de DNA tenha vindo da criança. (Fonte: infográfico de Maayan Harel.)}

\textbf{Figura 3.7.} Pedigree real de uma família com múltiplas vítimas no acidente da Malaysia Airlines. (Fonte: dados fornecidos por Willem Burgers.)

\textbf{Figura 3.8.} De testes de DNA a redes bayesianas. Na rede bayesiana, nós sem sombreamento representam alelos, e nós sombreados representam genótipos. Dados só estão disponíveis nos nós sombreados porque os genótipos não podem indicar qual alelo veio do pai e qual da mãe. A rede bayesiana possibilita a inferência sobre os nós não observados e também nos permite estimar a verossimilhança de que uma dada amostra de DNA tenha vindo da criança. (Fonte: infográfico de Maayan Harel.)

Uma vez que a rede bayesiana está montada, o passo final é inserir o DNA da vítima e computar a verossimilhança de que ele se encaixe num slot específico do pedigree. Isso é feito por propagação de crenças com a regra de Bayes. A rede começa com um grau particular de crença em cada afirmação possível sobre os nós da rede, como “o alelo paterno desta pessoa para a cor dos olhos é azul”. À medida que nova evidência é inserida na rede — em qualquer lugar da rede —, os graus de crença em cada nó, para cima e para baixo na rede, mudarão de modo em cascata. Assim, por exemplo, uma vez que descobrimos que uma dada amostra é uma correspondência provável para uma pessoa no pedigree, podemos propagar essa informação para cima e para baixo na rede. Desse modo, o Bonaparte não apenas aprende do DNA dos membros vivos da família, mas também das identificações que já fez.

Este exemplo ilustra vividamente várias vantagens das redes bayesianas. Uma vez que a rede está montada, o investigador não precisa intervir para lhe dizer como avaliar um novo pedaço de dados. A atualização pode ser feita muito rapidamente. (As redes bayesianas são especialmente boas para programação em um computador distribuído.) A rede é integrativa, o que significa que reage como um todo a qualquer nova informação. É por isso que até DNA de uma tia ou de um primo de segundo grau pode ajudar a identificar a vítima. As redes bayesianas são quase como um tecido orgânico vivo, o que não é acidente, porque esta é precisamente a imagem que eu tinha em mente quando lutava para fazê-las funcionar. Queria que as redes bayesianas operassem como os neurônios de um cérebro humano; você toca um neurônio, e a rede inteira responde propagando a informação a todos os outros neurônios do sistema.

A transparência das redes bayesianas as distingue da maioria das outras abordagens ao aprendizado de máquina, que tendem a produzir “caixas-pretas” inescrutáveis. Numa rede bayesiana, você pode seguir cada passo e compreender como e por que cada pedaço de evidência mudou as crenças da rede.

Por mais elegante que seja o Bonaparte, vale notar uma característica que ele ainda não incorpora: a intuição humana. Uma vez que termina a análise, ele fornece aos especialistas do NFI um \textit{ranking} das identificações mais prováveis para cada amostra de DNA e uma razão de verossimilhança para cada uma. Os investigadores estão então livres para combinar a evidência de DNA com outras evidências físicas recuperadas do local do acidente, bem como com sua intuição, para fazer suas determinações finais. No presente, nenhuma identificação é feita pelo computador agindo sozinho. Um objetivo da inferência causal é criar uma interface humano-máquina mais suave, que poderia permitir à intuição dos investigadores se juntar à dança da propagação de crenças.

Este exemplo de identificação de DNA com o Bonaparte apenas arranha a superfície das aplicações das redes bayesianas à genômica. Contudo, gostaria de passar a uma segunda aplicação que se tornou ubíqua na sociedade de hoje. De fato, há uma boa chance de que você tenha uma rede bayesiana no bolso agora mesmo. Chama-se celular, e cada um deles usa algoritmos de correção de erros baseados em propagação de crenças.

Para começar pelo começo: quando você fala num telefone, ele converte sua bela voz numa sequência de uns e zeros (chamados \textit{bits}) e os transmite usando um sinal de rádio. Infelizmente, nenhum sinal de rádio é recebido com fidelidade perfeita. À medida que o sinal faz seu caminho até a torre de celular e então até o telefone do seu amigo, alguns \textit{bits} aleatórios se inverterão de zero para um ou vice-versa. Para corrigir esses erros, podemos acrescentar informação redundante. Um esquema ultrassimples de correção de erros é simplesmente repetir cada \textit{bit} de informação três vezes: codificar um um como “111” e um zero como “000”. As sequências válidas “111” e “000” são chamadas palavras-código. Se o receptor ouve uma sequência inválida, como “101”, buscará a palavra-código válida mais provável para explicá-la. O zero é mais provavelmente errado do que ambos os uns, de modo que o decodificador interpretará essa mensagem como “111” e, portanto, concluirá que o \textit{bit} de informação era um um.

Infelizmente, esse código é altamente ineficiente, porque torna todas as nossas mensagens três vezes mais longas. Contudo, engenheiros de comunicação trabalharam por setenta anos para encontrar códigos de correção de erros cada vez melhores. O problema da decodificação é idêntico aos outros problemas de probabilidade inversa que discutimos, porque mais uma vez queremos inferir a probabilidade de uma hipótese (a mensagem enviada era “Hello world!”) a partir da evidência (a mensagem recebida era “Hxllo wovld!”). A situação parece madura para uma aplicação de propagação de crenças.

Em 1993, um engenheiro da France Telecom chamado Claude Berrou deixou o mundo da codificação estupefato com um código de correção de erros que alcançou desempenho quase ótimo. (Em outras palavras, a quantidade de informação redundante requerida está próxima do mínimo teórico.) Sua ideia, chamada “código turbo”, pode ser melhor ilustrada representando-a com uma rede bayesiana.

A Figura 3.9(a) mostra como um código tradicional funciona. Os \textit{bits} de informação, que você fala no telefone, são mostrados na primeira linha. São codificados, usando qualquer código que se queira — chame-o de código \textit{A} — em palavras-código (segunda linha), que são então recebidas com alguns erros (terceira linha). Esse diagrama é uma rede bayesiana, e podemos usar a propagação de crenças para inferir, a partir dos \textit{bits} recebidos, quais eram os \textit{bits} de informação. Contudo, isso de modo algum melhoraria o código \textit{A}.

\begin{figure}[ht]
\centering
\includegraphics[width=0.92\textwidth]{../figuras/fig_3_9.png}
\caption{Figura 3.9. (a) mostra como um código tradicional funciona. Os \textit{bits} de informação, que você fala no telefone, são mostrados na primeira linha. São codificados, usando qualquer código que se queira — chame-o de código \textit{A} — em palavras-código (segunda linha), que são então recebidas com alguns erros (terceira linha). Esse diagrama é uma rede bayesiana, e podemos usar a propagação de crenças para inferir, a partir dos \textit{bits} recebidos, quais eram os \textit{bits} de informação. Contudo, isso de modo algum melhoraria o código \textit{A}.}
\end{figure}
\textit{Figura 3.9. (a) mostra como um código tradicional funciona. Os }bits\textit{ de informação, que você fala no telefone, são mostrados na primeira linha. São codificados, usando qualquer código que se queira — chame-o de código }A\textit{ — em palavras-código (segunda linha), que são então recebidas com alguns erros (terceira linha). Esse diagrama é uma rede bayesiana, e podemos usar a propagação de crenças para inferir, a partir dos }bits\textit{ recebidos, quais eram os }bits\textit{ de informação. Contudo, isso de modo algum melhoraria o código }A\textit{.}

A ideia brilhante de Berrou foi codificar cada mensagem duas vezes, uma diretamente e uma depois de embaralhar a mensagem. Isso resulta na criação de duas palavras-código separadas e no recebimento de duas mensagens ruidosas (Figura 3.9b). Não há fórmula conhecida para decodificar diretamente uma mensagem dual desse tipo. Mas Berrou mostrou empiricamente que, se se aplicam as fórmulas de propagação de crenças em redes bayesianas repetidamente, duas coisas espantosas acontecem. Na maior parte do tempo (e com isso quero dizer algo como 99,999 por cento do tempo) você obtém os \textit{bits} de informação corretos. Não só isso: você pode usar palavras-código bem mais curtas. Em termos simples, duas cópias do código \textit{A} são bem melhores do que uma.

\textbf{Figura 3.9.} (a) Representação em rede bayesiana do processo ordinário de codificação. \textit{Bits} de informação são transformados em palavras-código; estes são transmitidos e recebidos no destino com ruído (erros). (b) Representação em rede bayesiana do código turbo. \textit{Bits} de informação são embaralhados e codificados duas vezes. A decodificação procede por propagação de crenças nessa rede. Cada processador na base usa informação do outro processador para melhorar seu palpite da palavra-código oculta, num processo iterativo.

Essa história em cápsula está correta, exceto por uma coisa: Berrou não sabia que estava trabalhando com redes bayesianas! Ele simplesmente descobrira o algoritmo de propagação de crenças por si mesmo. Só cinco anos depois David MacKay, de Cambridge, percebeu que era o mesmo algoritmo de que ele vinha gostando no final dos anos 1980 enquanto brincava com redes bayesianas. Isso colocou o algoritmo de Berrou num contexto teórico familiar e permitiu aos teóricos da informação afiar sua compreensão de seu desempenho.

De fato, outro engenheiro, Robert Gallager do Massachusetts Institute of Technology, havia descoberto um código que usava propagação de crenças (embora não chamada por esse nome) já em 1960, tão há tanto tempo que MacKay descreve seu código como “quase clarividente”. De qualquer modo, estava longe demais à frente de seu tempo. Gallager precisava de milhares de processadores num \textit{chip}, passando mensagens de um lado para o outro sobre seu grau de crença de que um \textit{bit} de informação particular era um um ou um zero. Em 1960 isso era impossível, e seu código foi virtualmente esquecido até que MacKay o redescobriu em 1998. Hoje, está em todo celular.

Por qualquer medida, os códigos turbo foram um sucesso estonteante. Antes da revolução turbo, os celulares 2G usavam “decodificação suave” (isto é, probabilidades), mas não propagação de crenças. Os celulares 3G usavam os códigos turbo de Berrou, e os telefones 4G usavam os códigos semelhantes a turbo de Gallager. Do ponto de vista do consumidor, isso significa que seu celular usa menos energia e a bateria dura mais, porque codificação e decodificação são os processos mais intensivos em energia do seu celular. Além disso, códigos melhores significam que você não precisa estar tão perto de uma torre de celular para obter transmissão de alta qualidade. Em outras palavras, as redes bayesianas permitiram aos fabricantes de telefones cumprir sua promessa: mais barras em mais lugares.

\section{Das redes bayesianas aos diagramas causais}
Depois de um capítulo dedicado às redes bayesianas, você pode se perguntar como elas se relacionam com o resto deste livro e em particular com os diagramas causais, o tipo que conhecemos no Capítulo 1. É claro, discuti-as em tanto detalhe em parte porque foram minha rota pessoal para a causalidade. Mas, mais importante tanto do ponto de vista teórico quanto prático, as redes bayesianas detêm a chave que permite aos diagramas causais fazer interface com os dados. Todas as propriedades probabilísticas das redes bayesianas (incluindo as junções que discutimos mais cedo neste capítulo) e os algoritmos de propagação de crenças que foram desenvolvidos para elas permanecem válidos nos diagramas causais. São de fato indispensáveis para compreender a inferência causal.

As principais diferenças entre redes bayesianas e diagramas causais residem em como são construídos e nos usos a que se destinam. Uma rede bayesiana é literalmente nada mais do que uma representação compacta de uma enorme tabela de probabilidade. As setas significam apenas que as probabilidades dos nós filhos estão relacionadas aos valores dos nós pais por uma certa fórmula (as tabelas de probabilidade condicional) e que essa relação é suficiente. Isto é, conhecer ancestrais adicionais do filho não mudará a fórmula. Da mesma forma, uma seta ausente entre quaisquer dois nós significa que eles são independentes, uma vez que conheçamos os valores de seus pais.

Vimos uma versão simples dessa afirmação mais cedo, quando discutimos o efeito de bloqueio em cadeias e elos. Numa cadeia \textit{A} → \textit{B} → \textit{C}, a seta ausente entre \textit{A} e \textit{C} significa que \textit{A} e \textit{C} são independentes uma vez que conheçamos os valores de seus pais. Como \textit{A} não tem pais, e o único pai de \textit{C} é \textit{B}, segue-se que \textit{A} e \textit{C} são independentes uma vez que conheçamos o valor de \textit{B}, o que concorda com o que dissemos antes.

Se, contudo, o mesmo diagrama foi construído como um diagrama causal, então tanto o pensamento que entra na construção quanto a interpretação do diagrama final mudam. Na fase de construção, precisamos examinar cada variável, digamos \textit{C}, e nos perguntar a quais outras variáveis ela “ouve” antes de escolher seu valor. A estrutura de cadeia \textit{A} → \textit{B} → \textit{C} significa que \textit{B} ouve apenas \textit{A}, \textit{C} ouve apenas \textit{B}, e \textit{A} não ouve ninguém; isto é, é determinada por forças externas que não fazem parte de nosso modelo.

Essa metáfora da escuta encapsula todo o conhecimento que uma rede causal transmite; o resto pode ser derivado, às vezes aproveitando dados.

Note que se invertermos a ordem das setas na cadeia, obtendo assim \textit{A} ← \textit{B} ← \textit{C}, a leitura causal da estrutura mudará drasticamente, mas as condições de independência permanecerão as mesmas. A seta ausente entre \textit{A} e \textit{C} ainda significará que \textit{A} e \textit{C} são independentes uma vez que conheçamos o valor de \textit{B}, como na cadeia original.

Isso tem duas implicações enormemente importantes. Primeiro, nos diz que suposições causais não podem ser inventadas a nosso bel-prazer; estão sujeitas ao escrutínio dos dados e podem ser falseadas. Por exemplo, se os dados observados não mostram \textit{A} e \textit{C} como independentes, condicionais a \textit{B}, então podemos concluir com segurança que o modelo de cadeia é incompatível com os dados e precisa ser descartado (ou reparado). Segundo, as propriedades gráficas do diagrama ditam quais modelos causais podem ser distinguidos por dados e quais permanecerão para sempre indistinguíveis, por maior que seja o volume de dados. Por exemplo, não podemos distinguir o garfo \textit{A} ← \textit{B} → \textit{C} da cadeia \textit{A} → \textit{B} → \textit{C} só com dados, porque os dois diagramas implicam as mesmas condições de independência.

Outro modo conveniente de pensar o modelo causal é em termos de experimentos hipotéticos. Cada seta pode ser pensada como uma afirmação sobre o resultado de um experimento hipotético. Uma seta de \textit{A} para \textit{C} significa que, se pudéssemos mexer apenas em \textit{A}, então esperaríamos ver uma mudança na probabilidade de \textit{C}. Uma seta ausente de \textit{A} para \textit{C} significa que, no mesmo experimento, não veríamos nenhuma mudança em \textit{C}, uma vez que mantivéssemos constantes os pais de \textit{C} (em outras palavras, \textit{B} no exemplo acima). Note que a expressão probabilística “uma vez que conheçamos o valor de \textit{B}” deu lugar à expressão causal “uma vez que mantenhamos \textit{B} constante”, o que implica que estamos fisicamente impedindo \textit{B} de variar e desabilitando a seta de \textit{A} para \textit{B}.

O pensamento causal que entra na construção da rede causal se recompensará, é claro, no tipo de perguntas que a rede pode responder. Enquanto uma rede bayesiana só pode nos dizer quão provável um evento é, dado que observamos outro (informação de degrau um), os diagramas causais podem responder a perguntas intervencionais e contrafactuais. Por exemplo, o garfo causal \textit{A} ← \textit{B} → \textit{C} nos diz em termos inequívocos que mexer em \textit{A} não teria efeito sobre \textit{C}, por mais intensa que fosse a mexida. Por outro lado, uma rede bayesiana não está equipada para lidar com uma “mexida”, nem para distinguir entre ver e fazer, nem de fato para distinguir um garfo de uma cadeia. Em outras palavras, tanto uma cadeia quanto um garfo prediriam que mudanças observadas em \textit{A} estão associadas a mudanças em \textit{C}, sem fazer nenhuma predição sobre o efeito de “mexer” em \textit{A}.

Chegamos agora ao segundo, e talvez mais importante, impacto das redes bayesianas sobre a inferência causal. As relações que foram descobertas entre a estrutura gráfica do diagrama e os dados que ele representa agora nos permitem emular a mexida sem fazê-la fisicamente. Especificamente, aplicar uma sequência inteligente de operações de condicionamento nos permite predizer o efeito de ações ou intervenções sem de fato conduzir um experimento. Para demonstrar, considere de novo o garfo causal \textit{A} ← \textit{B} → \textit{C}, em que proclamamos a correlação entre \textit{A} e \textit{C} como espúria. Podemos verificar isso com um experimento em que mexemos em \textit{A} e não encontramos correlação entre \textit{A} e \textit{C}. Mas podemos fazer melhor. Podemos pedir ao diagrama que emule o experimento e nos diga se alguma operação de condicionamento pode reproduzir a correlação que prevaleceria no experimento. A resposta sairia afirmativa: “A correlação entre \textit{A} e \textit{C} que seria medida após condicionar a \textit{B} seria igual à correlação vista no experimento.” Essa correlação pode ser estimada a partir dos dados e, em nosso caso, seria zero, confirmando fielmente nossa intuição de que mexer em \textit{A} não teria efeito sobre \textit{C}.

Essa capacidade de emular intervenções por observações inteligentes não poderia ter sido adquirida se as propriedades estatísticas das redes bayesianas não tivessem sido desvendadas entre 1980 e 1988. Agora podemos decidir que conjunto de variáveis devemos medir a fim de predizer os efeitos de intervenções a partir de estudos observacionais. Também podemos responder a perguntas de “Por quê?”. Por exemplo, alguém pode perguntar por que mexer em \textit{A} faz \textit{C} variar. É realmente o efeito direto de \textit{A}, ou é o efeito de uma variável mediadora \textit{B}? Se ambos, podemos avaliar que porção do efeito é mediada por \textit{B}? Para responder a tais perguntas de mediação, temos de imaginar duas intervenções simultâneas: mexer em \textit{A} e manter \textit{B} constante (a ser distinguido de condicionar a \textit{B}). Se pudermos realizar essa intervenção fisicamente, obtemos a resposta à nossa pergunta. Mas se estamos à mercê de estudos observacionais, precisamos emular as duas ações com um conjunto astucioso de observações. De novo, a estrutura gráfica do diagrama nos dirá se isso é possível.

Todas essas capacidades ainda estavam no futuro em 1988, quando comecei a pensar em como casar a causalidade com os diagramas. Eu só sabia que as redes bayesianas, como então concebidas, não podiam responder às perguntas que eu estava fazendo. A percepção de que você nem sequer consegue distinguir \textit{A} → \textit{B} → \textit{C} de \textit{A} ← \textit{B} → \textit{C} só com dados foi uma frustração dolorosa.

Sei que você, o leitor, está ansioso agora para aprender como os diagramas causais nos permitem fazer cálculos como os que acabei de descrever. E chegaremos lá — nos Capítulos 7 a 9. Mas ainda não estamos prontos, porque no momento em que começamos a falar de estudos observacionais \textit{versus} experimentais, deixamos as águas relativamente amigáveis da comunidade de IA pelas águas bem mais tempestuosas da estatística, que foram agitadas por seu divórcio infeliz da causalidade. Em retrospecto, lutar pela aceitação das redes bayesianas na IA foi um piquenique — não, um cruzeiro de luxo! — comparado com a luta que tive de travar pelos diagramas causais. Essa batalha ainda está em curso, com algumas ilhas restantes de resistência.

Para navegar nessas novas águas, teremos de compreender os modos pelos quais os estatísticos ortodoxos aprenderam a abordar a causalidade e as limitações desses métodos. As perguntas que levantamos acima, concernentes ao efeito de intervenções, incluindo efeitos diretos e indiretos, não fazem parte da estatística mainstream, principalmente porque os pais fundadores do campo a purgaram da linguagem de causa e efeito. Mas os estatísticos, não obstante, consideram permissível falar de causas e efeitos numa situação: um ensaio clínico randomizado (ECR) em que um tratamento \textit{A} é aleatoriamente atribuído a alguns indivíduos e não a outros e as mudanças observadas em \textit{B} são então comparadas. Aqui, tanto a estatística ortodoxa quanto a inferência causal concordam quanto ao significado da sentença “\textit{A} causa \textit{B}”.

Antes de nos voltarmos para a nova ciência de causa e efeito — iluminada por modelos causais —, deveríamos primeiro tentar compreender as forças e as limitações da velha ciência cega a modelos: por que a randomização é necessária para concluir que \textit{A} causa \textit{B} e a natureza da ameaça (chamada “confundimento”) que os ECRs se destinam a desarmar. O próximo capítulo retoma esses tópicos. Em minha experiência, a maioria dos estatísticos, bem como dos analistas de dados modernos, não se sente confortável com nenhuma dessas perguntas, pois não consegue articulá-las usando um vocabulário centrado em dados. De fato, frequentemente discordam sobre o que “confundimento” significa! Depois que examinarmos essas questões à luz dos diagramas causais, poderemos colocar os ensaios clínicos randomizados em seu contexto apropriado. Ou podemos vê-los como um caso especial de nosso motor de inferência, ou podemos ver a inferência causal como uma vasta extensão dos ECRs. Qualquer um dos pontos de vista está bem, e talvez pessoas treinadas a ver os ECRs como o árbitro da causalidade achem o segundo mais congenial.

A história bíblica de Daniel, frequentemente citada como o primeiro experimento controlado. Daniel (terceiro da esquerda?) percebeu que uma comparação adequada de duas dietas só poderia ser feita quando elas fossem dadas a dois grupos de indivíduos semelhantes, escolhidos de antemão. O rei Nabucodonosor (ao fundo) ficou impressionado com os resultados. (Fonte: desenho de Dakota Harr.)

\chapter{4 CONFUNDIMENTO E DESCONFUNDIMENTO: OU, COMO ABATER A VARIÁVEL À ESPREITA}
\begin{quote}Se a nossa concepção de efeitos causais tivesse algo a ver com experimentos randomizados, estes teriam sido inventados 500 anos antes de Fisher.\end{quote}
>

\begin{quote}— O AUTOR (2016)\end{quote}
ASHPENAZ, o superintendente da corte do rei Nabucodonosor, tinha um grande problema. Em 597 a.C., o rei da Babilônia saqueara o reino de Judá e trouxera de volta milhares de cativos, muitos deles da nobreza de Jerusalém. Como era costume em seu reino, Nabucodonosor queria que alguns deles servissem em sua corte, e ordenou a Ashpenaz que procurasse “meninos em quem não houvesse defeito algum, de boa aparência, e instruídos em toda a sabedoria, e doutos em ciência, e entendidos no conhecimento”. Esses afortunados meninos deveriam ser educados na língua e na cultura da Babilônia para que pudessem servir na administração do império, que se estendia do Golfo Pérsico ao mar Mediterrâneo. Como parte de sua educação, comeriam carne real e beberiam vinho real. E aí residia o problema. Um de seus favoritos, um rapaz chamado Daniel, recusou-se a tocar na comida. Por motivos religiosos, não podia comer carne que não tivesse sido preparada segundo as leis judaicas, e pediu que ele e seus amigos recebessem uma dieta de vegetais. Ashpenaz teria gostado de atender aos desejos do rapaz, mas temia que o rei notasse: “Quando ele vir vossos rostos abatidos, diferentes dos outros meninos da vossa idade, isso me custará a cabeça.”

Daniel tentou assegurar a Ashpenaz que a dieta vegetariana não diminuiria a capacidade deles de servir ao rei. Como convém a alguém “doutos em ciência, e entendidos no conhecimento”, propôs um experimento. Experimenta-nos por dez dias, disse ele. Toma quatro de nós e alimenta-nos só com vegetais; toma outro grupo de meninos e alimenta-os com a carne e o vinho do rei. Depois de dez dias, compara os dois grupos. Disse Daniel: “E, conforme vires, farás a teus servos.”

Mesmo que você não tenha lido a história, provavelmente consegue adivinhar o que aconteceu a seguir. Daniel e seus três companheiros prosperaram com a dieta vegetariana. O rei ficou tão impressionado com a sabedoria e o aprendizado deles — sem falar na aparência saudável — que lhes deu um lugar privilegiado em sua corte, onde “os achou dez vezes mais sábios do que todos os magos e astrólogos que havia em todo o seu reino”. Mais tarde, Daniel tornou-se intérprete dos sonhos do rei e sobreviveu a um memorável encontro numa cova de leões.

Por incrível que pareça, a história bíblica de Daniel encapsula de maneira profunda a conduta da ciência experimental de hoje. Ashpenaz formula uma pergunta sobre causalidade: Uma dieta vegetariana fará meus servos emagrecerem? Daniel propõe uma metodologia para lidar com qualquer pergunta desse tipo: monte dois grupos de pessoas, idênticos em todos os aspectos relevantes. Dê a um grupo um novo tratamento (uma dieta, um medicamento etc.), enquanto o outro grupo (chamado grupo de controle) recebe o tratamento antigo ou nenhum tratamento especial. Se, após um período adequado, você observar uma diferença mensurável entre os dois grupos supostamente idênticos, então o novo tratamento deve ser a causa da diferença. Hoje chamamos isso de experimento controlado.

O princípio é simples. Para compreender o efeito causal da dieta, gostaríamos de comparar o que acontece a Daniel em uma dieta com o que teria acontecido se ele tivesse permanecido na outra. Mas não podemos voltar no tempo e reescrever a história; por isso, fazemos a segunda melhor coisa: comparamos um grupo de pessoas que recebe o tratamento com um grupo de pessoas semelhantes que não o recebe. É óbvio, mas ainda assim crucial, que os grupos sejam comparáveis e representativos de alguma população. Se essas condições forem atendidas, os resultados deverão ser transferíveis para a população em geral. Em defesa de Daniel, ele parece compreender isso. Não está pedindo vegetais apenas em seu próprio benefício: se o ensaio mostrar que a dieta vegetariana é melhor, então todos os servos israelitas deveriam ter permissão para seguir essa dieta no futuro. É assim, ao menos, que interpreto a frase: “E, conforme vires, farás a teus servos.”

Daniel também compreendeu que era importante comparar grupos. Nesse aspecto, ele já era mais sofisticado do que muita gente hoje, que escolhe uma dieta da moda (por exemplo) só porque um amigo fez aquela dieta e emagreceu. Se você escolhe uma dieta com base apenas na experiência de um amigo, está essencialmente dizendo que acredita ser semelhante a esse amigo em todos os detalhes relevantes: idade, hereditariedade, ambiente doméstico, dieta anterior e assim por diante. Isso é muito a se assumir.

Outro ponto-chave do experimento de Daniel é que ele era prospectivo: os grupos foram escolhidos de antemão. Em contraste, suponha que você veja vinte pessoas num comercial que dizem ter emagrecido com uma dieta. Parece um tamanho de amostra bem grande, de modo que alguns espectadores poderiam considerá-lo uma evidência convincente. Mas isso equivaleria a basear a decisão na experiência de pessoas que já tiveram uma boa resposta. Pelo que você sabe, para cada pessoa que emagreceu, outras dez iguais a ela experimentaram a dieta e não tiveram sucesso. Mas, é claro, elas não foram escolhidas para aparecer no comercial.

O experimento de Daniel era notavelmente moderno em todos esses aspectos. Ensaios controlados prospectivos ainda são um marco da boa ciência. No entanto, Daniel não pensou em uma coisa: o viés de confundimento. Suponha que Daniel e seus amigos sejam mais saudáveis do que o grupo de controle desde o início. Nesse caso, sua aparência robusta após dez dias de dieta pode não ter nada a ver com a dieta em si; pode refletir a saúde geral deles. Talvez tivessem prosperado ainda mais se tivessem comido a carne do rei!

O viés de confundimento ocorre quando uma variável influencia tanto quem é selecionado para o tratamento quanto o desfecho do experimento. Às vezes os confundidores são conhecidos; outras vezes, são apenas suspeitos e atuam como uma “terceira variável à espreita”. Num diagrama causal, os confundidores são extremamente fáceis de reconhecer: na Figura 4.1, a variável \textit{Z} no centro do garfo é um confundidor de \textit{X} e \textit{Y}. (Veremos uma definição mais universal adiante, mas esse triângulo é a situação mais reconhecível e comum.)

\begin{figure}[ht]
\centering
\includegraphics[width=0.92\textwidth]{../figuras/fig_4_1.png}
\caption{Figura 4.1. , a variável \textit{Z} no centro do garfo é um confundidor de \textit{X} e \textit{Y}. (Veremos uma definição mais universal adiante, mas esse triângulo é a situação mais reconhecível e comum.)}
\end{figure}
\textit{Figura 4.1. , a variável }Z\textit{ no centro do garfo é um confundidor de }X\textit{ e }Y\textit{. (Veremos uma definição mais universal adiante, mas esse triângulo é a situação mais reconhecível e comum.)}

\textbf{FIGURA 4.1.} A versão mais básica de confundimento: \textit{Z} é um confundidor da relação causal proposta entre \textit{X} e \textit{Y}.

O termo “confundimento” (\textit{confounding}) significava originalmente “mistura” em inglês, e podemos compreender pelo diagrama por que esse nome foi escolhido. O verdadeiro efeito causal \textit{X} → \textit{Y} fica “misturado” com a correlação espúria entre \textit{X} e \textit{Y} induzida pelo garfo \textit{X} ← \textit{Z} → \textit{Y}. Por exemplo, se estamos testando um medicamento e o damos a pacientes que, em média, são mais jovens do que as pessoas no grupo de controle, então a idade se torna um confundidor — uma terceira variável à espreita. Se não tivermos dados sobre as idades, não seremos capazes de desvencilhar o efeito verdadeiro do efeito espúrio.

No entanto, o inverso também é verdadeiro. Se tivermos medições da terceira variável, então é muito fácil desconfundir os efeitos verdadeiro e espúrio. Por exemplo, se a variável de confundimento \textit{Z} for a idade, comparamos os grupos de tratamento e de controle em cada faixa etária separadamente. Podemos então tirar a média dos efeitos, ponderando cada faixa etária de acordo com sua porcentagem na população-alvo. Esse método de compensação é familiar a todos os estatísticos; chama-se “ajustamento por \textit{Z}” ou “controle por \textit{Z}”.

Curiosamente, os estatísticos tanto superestimam quanto subestimam a importância de ajustar por possíveis confundidores. Superestimam no sentido de que frequentemente controlam por muito mais variáveis do que precisam — e até por variáveis que não deveriam controlar. Recentemente deparei com uma citação de um blogueiro político chamado Ezra Klein que expressa com muita clareza esse fenômeno de “supercontrole”: “Você vê o tempo todo em estudos. ‘Controlamos por…’ E então a lista começa. Quanto mais longa, melhor. Renda. Idade. Raça. Religião. Altura. Cor do cabelo. Preferência sexual. Frequência na CrossFit. Amor pelos pais. Coca ou Pepsi. Quanto mais coisas você puder controlar, mais forte será o seu estudo — ou, ao menos, mais forte o seu estudo \textit{parece}. Os controles dão a sensação de especificidade, de precisão… Mas às vezes dá para controlar demais. Às vezes você acaba controlando pela coisa que está tentando medir.”

Klein levanta uma preocupação válida. Os estatísticos ficaram imensamente confusos sobre quais variáveis devem e não devem ser controladas, de modo que a prática padrão tem sido controlar por tudo o que se possa medir. A vasta maioria dos estudos conduzidos nos dias de hoje adere a essa prática. É um procedimento conveniente e simples de seguir, mas é ao mesmo tempo dispendioso e repleto de erros. Uma conquista fundamental da Revolução Causal foi pôr fim a essa confusão.

Ao mesmo tempo, os estatísticos subestimam enormemente o controle no sentido de que se recusam a falar de causalidade de modo algum, mesmo quando o controle foi feito corretamente. Isso também vai de encontro à mensagem deste capítulo: se você identificou um conjunto suficiente de desconfundidores no seu diagrama, reuniu dados sobre eles e os ajustou adequadamente, então tem todo o direito de dizer que calculou o efeito causal \textit{X} → \textit{Y} (desde que, é claro, possa defender o seu diagrama causal com fundamentos científicos).

A abordagem de livro-texto dos estatísticos ao confundimento é bem diferente e repousa numa ideia defendida com mais eficácia por R. A. Fisher: o ensaio clínico randomizado (ECR). Fisher estava exatamente certo, mas não exatamente pelas razões certas. O ensaio clínico randomizado é de fato uma invenção maravilhosa — mas até há pouco tempo as gerações de estatísticos que seguiram Fisher não conseguiam provar que o que obtinham do ECR era de fato o que buscavam obter. Não tinham uma linguagem para escrever o que procuravam — a saber, o efeito causal de \textit{X} sobre \textit{Y}. Um dos meus objetivos neste capítulo é explicar, do ponto de vista dos diagramas causais, precisamente por que os ECRs nos permitem estimar o efeito causal \textit{X} → \textit{Y} sem cair presa do viés de confundimento. Uma vez compreendido por que os ECRs funcionam, não há necessidade de colocá-los num pedestal e tratá-los como o padrão-ouro da análise causal, que todos os outros métodos deveriam emular. Muito pelo contrário: veremos que o chamado padrão-ouro, na verdade, deriva sua legitimidade de princípios mais básicos.

Este capítulo também mostrará que os diagramas causais possibilitam uma mudança de ênfase dos confundidores para os desconfundidores. Os primeiros causam o problema; os segundos o curam. Os dois conjuntos podem se sobrepor, mas não precisam. Se tivermos dados sobre um conjunto suficiente de desconfundidores, não importa se ignorarmos alguns ou mesmo todos os confundidores. Essa mudança de ênfase é uma das principais maneiras pelas quais a Revolução Causal nos permite ir além dos experimentos fisherianos e inferir efeitos causais a partir de estudos não experimentais. Ela nos capacita a determinar quais variáveis devem ser controladas para servir de desconfundidores. Essa questão tem atormentado estatísticos teóricos e práticos; tem sido o calcanhar de Aquiles do campo por décadas. Isso porque ela não tem nada a ver com dados ou estatística. O confundimento é um conceito causal — pertence ao degrau dois da Escada da Causalidade. Métodos gráficos, a partir dos anos 1990, desconfundiram totalmente o problema do confundimento.

Em particular, em breve conheceremos um método chamado critério da “porta dos fundos” (\textit{back-door criterion}), que identifica de forma inequívoca quais variáveis num diagrama causal são desconfundidores. Se a pesquisadora puder reunir dados sobre essas variáveis, ela poderá ajustá-las e, com isso, fazer previsões sobre o resultado de uma intervenção mesmo sem realizá-la. De fato, a Revolução Causal foi ainda mais longe. Em alguns casos, podemos controlar o confundimento mesmo quando não temos dados sobre um conjunto suficiente de desconfundidores. Nesses casos, podemos usar fórmulas de ajustamento diferentes — não a convencional, que é apropriada apenas para uso com o critério da porta dos fundos — e ainda assim erradicar todo o confundimento. Guardaremos esses desenvolvimentos empolgantes para o Capítulo 7.

Embora o confundimento tenha uma longa história em todas as áreas da ciência, o reconhecimento de que o problema exige soluções causais, e não estatísticas, é muito recente. Ainda em 2001, revisores repreenderam um artigo meu insistindo: “O confundimento está solidamente fundamentado na estatística padrão.” Felizmente, o número de tais revisores encolheu dramaticamente na última década. Há agora um consenso quase universal, ao menos entre epidemiologistas, filósofos e cientistas sociais, de que (1) o confundimento precisa e tem uma solução causal, e (2) os diagramas causais fornecem uma maneira completa e sistemática de encontrar essa solução. A era da confusão sobre o confundimento chegou ao fim!

\section{O TEMOR GÉLIDO DO CONFUNDIMENTO}
Em 1998, um estudo no \textit{New England Journal of Medicine} revelou uma associação entre caminhada regular e taxas de mortalidade reduzidas entre homens aposentados. Os pesquisadores usaram dados do Honolulu Heart Program, que acompanha a saúde de 8.000 homens de ascendência japonesa desde 1965. Os pesquisadores, liderados por Robert Abbott, bioestatístico da Universidade da Virgínia, queriam saber se os homens que se exercitavam mais viviam mais. Escolheram uma amostra de 707 homens do grupo maior de 8.000, todos fisicamente saudáveis o bastante para caminhar. A equipe de Abbott descobriu que a taxa de mortalidade ao longo de um período de doze anos era duas vezes maior entre homens que caminhavam menos de 1,6 km por dia (vou chamá-los de “caminhantes casuais”) do que entre homens que caminhavam mais de 3,2 km por dia (“caminhantes intensos”). Para ser preciso, 43\% dos caminhantes casuais haviam morrido, enquanto apenas 21,5\% dos caminhantes intensos haviam morrido.

No entanto, como os experimentadores não prescreveram quem seria um caminhante casual e quem seria um caminhante intenso, temos de levar em conta a possibilidade de viés de confundimento. Um confundidor óbvio poderia ser a idade: homens mais jovens poderiam estar mais dispostos a um treino vigoroso e também seriam menos propensos a morrer. Teríamos então um diagrama causal como o da Figura 4.2.

\begin{figure}[ht]
\centering
\includegraphics[width=0.92\textwidth]{../figuras/fig_4_2.png}
\caption{Figura 4.2. \textbf{FIGURA 4.2.} Diagrama causal para o exemplo da caminhada. O padrão clássico de garfo no nó “Idade” nos diz que a idade é um confundidor de caminhada e mortalidade.}
\end{figure}
\textit{Figura 4.2. \textbf{FIGURA 4.2.} Diagrama causal para o exemplo da caminhada. O padrão clássico de garfo no nó “Idade” nos diz que a idade é um confundidor de caminhada e mortalidade.}

\textbf{FIGURA 4.2.} Diagrama causal para o exemplo da caminhada. O padrão clássico de garfo no nó “Idade” nos diz que a idade é um confundidor de caminhada e mortalidade.

Tenho certeza de que você consegue pensar em outros possíveis confundidores. Talvez os caminhantes casuais estivessem se poupando por algum motivo; talvez não pudessem caminhar tanto. Assim, a condição física poderia ser um confundidor. Poderíamos continuar assim indefinidamente. E se os caminhantes leves fossem bebedores de álcool? E se comessem mais? A boa notícia é que os pesquisadores pensaram em todos esses fatores. O estudo levou em conta e ajustou por todo fator razoável — idade, condição física, consumo de álcool, dieta e vários outros. Por exemplo, é verdade que os caminhantes intensos tendiam a ser um pouco mais jovens. Então os pesquisadores ajustaram a taxa de mortalidade por idade e descobriram que a diferença entre caminhantes casuais e intensos ainda era muito grande. (A taxa de mortalidade ajustada por idade para os caminhantes casuais era de 41\%, comparada a 24\% para os caminhantes intensos.)

Ainda assim, os pesquisadores foram muito cautelosos em suas conclusões. No final do artigo, escreveram: “É claro que os efeitos sobre a longevidade de esforços intencionais para aumentar a distância caminhada por dia por homens idosos fisicamente capazes não podem ser abordados em nosso estudo.” Usando a linguagem do Capítulo 1, eles se recusam a dizer qualquer coisa sobre a sua probabilidade de sobreviver doze anos dado que você \textit{do}(exercício).

Em justiça a Abbott e ao restante de sua equipe, eles podem ter tido boas razões para a cautela. Este era um primeiro estudo, e a amostra era relativamente pequena e homogênea. Não obstante, essa cautela reflete uma atitude mais geral, que transcende questões de homogeneidade e tamanho amostral. Os pesquisadores foram ensinados a acreditar que um estudo observacional (aquele em que os sujeitos escolhem o próprio tratamento) nunca pode iluminar uma afirmação causal. Afirmo que essa cautela é exagerada. Por que mais alguém se daria ao trabalho de ajustar por todos esses confundidores, senão para se livrar da parte espúria da associação e, com isso, obter uma visão melhor da parte causal? Em vez de dizer “É claro que não podemos”, como fizeram, deveríamos proclamar que é claro que \textit{podemos} dizer algo sobre uma intervenção intencional. Se acreditamos que a equipe de Abbott identificou todos os confundidores importantes, devemos também acreditar que a caminhada intencional tende a prolongar a vida (ao menos em homens japoneses).

Essa conclusão provisória, baseada na premissa de que nenhum outro confundidor poderia desempenhar um papel importante nas relações encontradas, é um pedaço de informação extremamente valioso. Ela diz a um potencial caminhante precisamente que tipo de incerteza permanece ao se tomar a afirmação pelo valor de face. Diz-lhe que a incerteza restante não é maior do que a possibilidade de existirem confundidores adicionais que não foram levados em conta. Também é valiosa como guia para estudos futuros, que deveriam focar nesses outros fatores (se existirem), e não nos que foram neutralizados no estudo atual. Em suma, conhecer o conjunto de premissas que sustentam uma dada conclusão não é menos valioso do que tentar contornar essas premissas com um ECR, que, como veremos, tem complicações próprias.

\section{A HÁBIL INTERROGATÓRIO DA NATUREZA: POR QUE OS ECRS FUNCIONAM}
Como já mencionei, a única circunstância sob a qual os cientistas abandonam um pouco de sua reticência em falar de causalidade é quando conduziram um ensaio clínico randomizado. Você pode ler isso na Wikipédia ou em mil outros lugares: “O ECR é frequentemente considerado o padrão-ouro de um ensaio clínico.” Temos uma pessoa a quem agradecer por isso, R. A. Fisher, por isso é muito interessante ler o que uma pessoa muito próxima a ele escreveu sobre suas razões. A passagem é longa, mas vale a pena citar por completo:

\begin{quote}Toda a arte e a prática da experimentação científica consistem no hábil interrogatório da Natureza. A observação forneceu ao cientista um quadro da Natureza em algum aspecto, que tem todas as imperfeições de uma declaração voluntária. Ele deseja verificar sua interpretação dessa declaração formulando perguntas específicas voltadas a estabelecer relações causais. Suas perguntas, na forma de operações experimentais, são necessariamente particulares, e ele deve confiar na consistência da Natureza ao fazer deduções gerais a partir de sua resposta em um caso particular ou ao prever o desfecho a ser antecipado de operações semelhantes em outras ocasiões. Seu objetivo é extrair conclusões válidas de precisão e generalidade determinadas a partir da evidência que ele obtém. Longe de se comportar de maneira consistente, porém, a Natureza parece vacilante, esquiva e ambígua em suas respostas. Ela responde à forma da pergunta tal como esta é colocada no campo, e não necessariamente à pergunta na mente do experimentador; ela não interpreta por ele; não dá informação gratuita; e é exigente quanto à precisão. Em consequência, o experimentador que quer comparar dois tratamentos de adubo desperdiça o trabalho se, dividindo seu campo em duas partes iguais, adubar cada metade com um de seus adubos, cultivar uma colheita e comparar os rendimentos das duas metades. A forma de sua pergunta foi: qual é a diferença entre o rendimento do lote A sob o primeiro tratamento e o do lote B sob o segundo? Ele não perguntou se o lote A renderia o mesmo que o lote B sob tratamento uniforme, e não pode distinguir efeitos de lote de efeitos de tratamento, pois a Natureza registrou, conforme solicitado, não apenas a contribuição das diferenças de adubo para os rendimentos dos lotes, mas também as contribuições de diferenças na fertilidade do solo, textura, drenagem, exposição, microflora e inumeráveis outras variáveis.\end{quote}
A autora dessa passagem é Joan Fisher Box, filha de Ronald Aylmer Fisher, e ela é tirada de sua biografia do ilustre pai. Embora não seja estatística, ela absorveu com clareza o desafio central que os estatísticos enfrentam. Afirma em termos inequívocos que as perguntas que eles formulam estão “voltadas a estabelecer relações causais”. E o que atrapalha o caminho é o confundimento, embora ela não use essa palavra. Eles querem saber o efeito de um fertilizante (ou “tratamento de adubo”, como os fertilizantes eram chamados naquela época) — isto é, o rendimento esperado sob um fertilizante comparado com o rendimento sob uma alternativa. A Natureza, porém, lhes conta sobre o efeito do fertilizante misturado (lembre-se, este é o significado original de “confundido”) com uma variedade de outras causas.

Gosto da imagem que Fisher Box oferece na passagem acima: a Natureza é como um gênio que responde exatamente à pergunta que formulamos, não necessariamente à que pretendemos fazer. Mas temos de acreditar, como Fisher Box claramente acredita, que a resposta à pergunta que desejamos fazer existe na natureza. Nossos experimentos são um meio desajeitado de descobrir a resposta, mas de modo algum a definem. Se seguirmos sua analogia exatamente, então \textit{do}(\textit{X} = \textit{x}) deve vir primeiro, porque é uma propriedade da natureza que representa a resposta que buscamos: Qual é o efeito de usar o primeiro fertilizante em todo o campo? A randomização vem em segundo lugar, porque é apenas um meio feito pelo homem para extrair a resposta a essa pergunta. Pode-se compará-la ao ponteiro de um termômetro, que é um meio de obter a temperatura, mas não é a temperatura em si.

Nos primeiros anos na Estação Experimental de Rothamsted, Fisher geralmente adotava uma abordagem muito elaborada e sistemática para desvencilhar os efeitos do fertilizante de outras variáveis. Dividia seus campos em uma grade de sublotes e planejava com cuidado para que cada fertilizante fosse experimentado com cada combinação de tipo de solo e planta (ver Figura 4.3). Fazia isso para garantir a comparabilidade de cada amostra; na realidade, nunca poderia antecipar todos os confundidores que poderiam determinar a fertilidade de um dado lote. Um gênio esperto o bastante poderia derrotar qualquer arranjo estruturado do campo.

\begin{figure}[ht]
\centering
\includegraphics[width=0.92\textwidth]{../figuras/fig_4_3.png}
\caption{Figura 4.3. ). Fazia isso para garantir a comparabilidade de cada amostra; na realidade, nunca poderia antecipar todos os confundidores que poderiam determinar a fertilidade de um dado lote. Um gênio esperto o bastante poderia derrotar qualquer arranjo estruturado do campo.}
\end{figure}
\textit{Figura 4.3. ). Fazia isso para garantir a comparabilidade de cada amostra; na realidade, nunca poderia antecipar todos os confundidores que poderiam determinar a fertilidade de um dado lote. Um gênio esperto o bastante poderia derrotar qualquer arranjo estruturado do campo.}

Por volta de 1923 ou 1924, Fisher começou a perceber que o único desenho experimental que o gênio não poderia derrotar era um desenho aleatório. Imagine realizar o mesmo experimento cem vezes num campo com uma distribuição desconhecida de fertilidade. Cada vez, você atribui fertilizantes a sublotes aleatoriamente. Às vezes você pode ser muito azarado e usar o Fertilizante 1 em todos os sublotes menos férteis. Outras vezes pode ter sorte e aplicá-lo nos sublotes mais férteis. Mas, ao gerar uma nova atribuição aleatória cada vez que realiza o experimento, você pode garantir que na grande maioria das vezes não será nem sortudo nem azarado. Nesses casos, o Fertilizante 1 será aplicado a uma seleção de sublotes que é representativa do campo como um todo. É exatamente isso que se deseja para um ensaio controlado. Como a distribuição de fertilidade no campo é fixa ao longo de sua série de experimentos — o gênio não pode mudá-la —, ele é enganado a responder (na maior parte do tempo) à pergunta causal que você queria fazer.

\textbf{FIGURA 4.3.} R. A. Fisher com uma de suas muitas inovações: um desenho experimental de quadrado latino, destinado a garantir que um lote de cada tipo de planta apareça em cada linha (tipo de fertilizante) e coluna (tipo de solo). Tais desenhos ainda são usados na prática, mas Fisher argumentaria mais tarde de forma convincente que um desenho randomizado é ainda mais eficaz. (Fonte: Desenho de Dakota Harr.)

Do nosso ponto de vista, numa era em que os ensaios randomizados são o padrão-ouro, tudo isso pode parecer óbvio. Mas, na época, a ideia de um experimento desenhado aleatoriamente horrorizou os colegas estatísticos de Fisher. O fato de Fisher literalmente tirar cartas de um baralho para atribuir sublotes a cada fertilizante pode ter contribuído para o desconforto deles. Ciência sujeita aos caprichos do acaso? Mas Fisher percebeu que uma resposta incerta à pergunta certa é muito melhor do que uma resposta altamente certa à pergunta errada. Se você fizer ao gênio a pergunta errada, nunca descobrirá o que quer saber. Se fizer a pergunta certa, obter uma resposta que ocasionalmente está errada é um problema bem menor. Você ainda pode estimar a quantidade de incerteza na sua resposta, porque a incerteza vem do procedimento de randomização (que é conhecido) e não das características do solo (que são desconhecidas).

Assim, a randomização traz de fato dois benefícios. Primeiro, elimina o viés de confundimento (faz à Natureza a pergunta certa). Segundo, permite ao pesquisador quantificar sua incerteza. No entanto, segundo o historiador Stephen Stigler, o segundo benefício era realmente a principal razão de Fisher para defender a randomização. Ele era o mestre mundial em quantificar a incerteza, tendo desenvolvido muitos novos procedimentos matemáticos para isso. Em comparação, sua compreensão do desconfundimento era puramente intuitiva, pois carecia de uma notação matemática para articular o que buscava.

Agora, noventa anos depois, podemos usar o operador \textit{do} para preencher o que Fisher queria, mas não podia, perguntar. Vejamos, de um ponto de vista causal, como a randomização nos permite fazer ao gênio a pergunta certa. Comecemos, como de costume, desenhando um diagrama causal. O Modelo 1, mostrado na Figura 4.4, descreve como o rendimento de cada lote é determinado sob condições normais, em que o fazendeiro decide por capricho ou viés qual fertilizante é melhor para cada lote. A consulta que ele quer formular ao gênio Natureza é: “Qual é o rendimento sob uma aplicação uniforme do Fertilizante 1 (versus o Fertilizante 2) a todo o campo?” Ou, na notação do operador \textit{do}, qual é \textit{P}(rendimento | \textit{do}(fertilizante = 1))?

\begin{figure}[ht]
\centering
\includegraphics[width=0.92\textwidth]{../figuras/fig_4_4.png}
\caption{Figura 4.4. , descreve como o rendimento de cada lote é determinado sob condições normais, em que o fazendeiro decide por capricho ou viés qual fertilizante é melhor para cada lote. A consulta que ele quer formular ao gênio Natureza é: “Qual é o rendimento sob uma aplicação uniforme do Fertilizante 1 (versus o Fertilizante 2) a todo o campo?” Ou, na notação do operador \textit{do}, qual é \textit{P}(rendimento | \textit{do}(fertilizante = 1))?}
\end{figure}
\textit{Figura 4.4. , descreve como o rendimento de cada lote é determinado sob condições normais, em que o fazendeiro decide por capricho ou viés qual fertilizante é melhor para cada lote. A consulta que ele quer formular ao gênio Natureza é: “Qual é o rendimento sob uma aplicação uniforme do Fertilizante 1 (versus o Fertilizante 2) a todo o campo?” Ou, na notação do operador }do\textit{, qual é }P\textit{(rendimento | }do\textit{(fertilizante = 1))?}

\textbf{FIGURA 4.4.} Modelo 1: um experimento impropriamente controlado.

Se o fazendeiro realizar o experimento de forma ingênua — por exemplo, aplicando o Fertilizante 1 na parte alta de seu campo e o Fertilizante 2 na parte baixa —, ele provavelmente estará introduzindo a Drenagem como confundidor. Se usar o Fertilizante 1 num ano e o Fertilizante 2 no ano seguinte, provavelmente estará introduzindo o Clima como confundidor. Em qualquer dos casos, obterá uma comparação enviesada.

O mundo sobre o qual o fazendeiro quer saber é descrito pelo Modelo 2, em que todos os lotes recebem o mesmo fertilizante (ver Figura 4.5). Como explicado no Capítulo 1, o efeito do operador \textit{do} é apagar todas as setas que apontam para Fertilizante e forçar essa variável a um valor particular — digamos, Fertilizante = 1.

\begin{figure}[ht]
\centering
\includegraphics[width=0.92\textwidth]{../figuras/fig_4_5.png}
\caption{Figura 4.5. ). Como explicado no Capítulo 1, o efeito do operador \textit{do} é apagar todas as setas que apontam para Fertilizante e forçar essa variável a um valor particular — digamos, Fertilizante = 1.}
\end{figure}
\textit{Figura 4.5. ). Como explicado no Capítulo 1, o efeito do operador }do\textit{ é apagar todas as setas que apontam para Fertilizante e forçar essa variável a um valor particular — digamos, Fertilizante = 1.}

\textbf{FIGURA 4.5.} Modelo 2: o mundo sobre o qual gostaríamos de saber.

Por fim, vejamos como o mundo se parece quando aplicamos a randomização. Agora alguns lotes serão submetidos a \textit{do}(fertilizante = 1) e outros a \textit{do}(fertilizante = 2), mas a escolha de qual tratamento vai para qual lote é aleatória. O mundo criado por tal modelo é mostrado pelo Modelo 3 na Figura 4.6, mostrando a variável Fertilizante obtendo sua atribuição por um dispositivo aleatório — digamos, o baralho de Fisher. Observe que todas as setas apontando para Fertilizante foram apagadas, refletindo a premissa de que o fazendeiro escuta apenas a carta ao decidir qual fertilizante usar. É igualmente importante notar que não há seta de Carta para Rendimento, porque as plantas não sabem ler as cartas. (Essa é uma premissa bastante segura para plantas, mas para sujeitos humanos num ensaio randomizado é uma preocupação séria.) Portanto, o Modelo 3 descreve um mundo em que a relação entre Fertilizante e Rendimento está desconfundida (isto é, não há causa comum de Fertilizante e Rendimento). Isso significa que, no mundo descrito pela Figura 4.6, não há diferença entre \textit{ver} Fertilizante = 1 e \textit{fazer} Fertilizante = 1.

\begin{figure}[ht]
\centering
\includegraphics[width=0.92\textwidth]{../figuras/fig_4_6.png}
\caption{Figura 4.6. , mostrando a variável Fertilizante obtendo sua atribuição por um dispositivo aleatório — digamos, o baralho de Fisher. Observe que todas as setas apontando para Fertilizante foram apagadas, refletindo a premissa de que o fazendeiro escuta apenas a carta ao decidir qual fertilizante usar. É igualmente importante notar que não há seta de Carta para Rendimento, porque as plantas não sabem ler as cartas. (Essa é uma premissa bastante segura para plantas, mas para sujeitos humanos num ensaio randomizado é uma preocupação séria.) Portanto, o Modelo 3 descreve um mundo em que a relação entre Fertilizante e Rendimento está desconfundida (isto é, não há causa comum de Fertilizante e Rendimento). Isso significa que, no mundo descrito pela Figura 4.6, não há diferença entre \textit{ver} Fertilizante = 1 e \textit{fazer} Fertilizante = 1.}
\end{figure}
\textit{Figura 4.6. , mostrando a variável Fertilizante obtendo sua atribuição por um dispositivo aleatório — digamos, o baralho de Fisher. Observe que todas as setas apontando para Fertilizante foram apagadas, refletindo a premissa de que o fazendeiro escuta apenas a carta ao decidir qual fertilizante usar. É igualmente importante notar que não há seta de Carta para Rendimento, porque as plantas não sabem ler as cartas. (Essa é uma premissa bastante segura para plantas, mas para sujeitos humanos num ensaio randomizado é uma preocupação séria.) Portanto, o Modelo 3 descreve um mundo em que a relação entre Fertilizante e Rendimento está desconfundida (isto é, não há causa comum de Fertilizante e Rendimento). Isso significa que, no mundo descrito pela Figura 4.6, não há diferença entre }ver\textit{ Fertilizante = 1 e }fazer\textit{ Fertilizante = 1.}

\textbf{FIGURA 4.6.} Modelo 3: o mundo simulado por um ensaio clínico randomizado.

Isso nos leva à conclusão principal: a randomização é uma maneira de simular o Modelo 2. Ela desativa todos os confudidores antigos sem introduzir novos confudidores. Essa é a fonte de seu poder; não há nada de misterioso ou místico nisso. É nada mais, nada menos do que, como Joan Fisher Box disse, “o hábil interrogatório da Natureza”.

O experimento, porém, falharia em seu objetivo de simular o Modelo 2 se o experimentador tivesse permissão de usar o próprio julgamento para escolher um fertilizante, ou se os sujeitos experimentais — neste caso, as plantas — “soubessem” qual carta haviam tirado. É por isso que ensaios clínicos com sujeitos humanos se esforçam tanto para ocultar essa informação tanto dos pacientes quanto dos experimentadores (um procedimento conhecido como duplo-cego).

Acrescentarei a isso uma segunda conclusão principal: há outras maneiras de simular o Modelo 2. Uma delas, se você souber quais são todos os possíveis confudidores, é medi-los e ajustá-los. No entanto, a randomização tem uma grande vantagem: ela corta todos os elos de entrada para a variável randomizada, inclusive os que não conhecemos ou não podemos medir (por exemplo, fatores “Outros” nas Figuras 4.4 a 4.6). Em contraste, num estudo não randomizado, a experimentadora deve confiar em seu conhecimento do assunto. Se ela estiver confiante de que seu modelo causal dá conta de um número suficiente de desconfundidores e tiver reunido dados sobre eles, então poderá estimar o efeito do Fertilizante sobre o Rendimento de forma não enviesada. Mas o perigo é que ela possa ter deixado passar um fator de confundimento, e sua estimativa poderá, portanto, ser enviesada.

Tudo o mais igual, os ECRs ainda são preferidos a estudos observacionais, assim como redes de segurança são recomendadas para equilibristas. Mas nem tudo o mais é necessariamente igual. Em alguns casos, a intervenção pode ser fisicamente impossível (por exemplo, num estudo do efeito da obesidade sobre a doença cardíaca, não podemos atribuir aleatoriamente pacientes a serem obesos ou não). Ou a intervenção pode ser antiética (num estudo dos efeitos do fumo, não podemos pedir a pessoas selecionadas aleatoriamente que fumem por dez anos). Ou podemos encontrar dificuldades em recrutar sujeitos para procedimentos experimentais inconvenientes e acabar com voluntários que não representam a população pretendida.

Felizmente, o operador \textit{do} nos dá maneiras cientificamente sólidas de determinar efeitos causais a partir de estudos não experimentais, o que desafia a supremacia tradicional dos ECRs. Como discutido no exemplo da caminhada, tais estimativas causais produzidas por estudos observacionais podem ser rotuladas de “causalidade provisória” — isto é, causalidade contingente ao conjunto de premissas que nosso diagrama causal anuncia. É importante que não tratemos esses estudos como cidadãos de segunda classe: eles têm a vantagem de serem conduzidos no habitat natural da população-alvo, e não no ambiente artificial de um laboratório, e podem ser “puros” no sentido de não estarem contaminados por questões de ética ou viabilidade.

Agora que compreendemos que o principal objetivo de um ECR é eliminar o confundimento, vejamos os outros métodos que a Revolução Causal nos deu. A história começa com um artigo de 1986 de dois de meus colegas de longa data, que iniciou uma reavaliação do que o confundimento significa.

\section{O NOVO PARADIGMA DO CONFUNDIMENTO}
“Embora o confundimento seja amplamente reconhecido como um dos problemas centrais da pesquisa epidemiológica, uma revisão da literatura revelará pouca consistência entre as definições de confundimento ou de confundidor.” Com essa única frase, Sander Greenland, da Universidade da Califórnia em Los Angeles, e Jamie Robins, da Universidade Harvard, puseram o dedo na razão central pela qual o controle do confundimento não havia avançado um milímetro desde Fisher. Sem uma compreensão de princípios do confundimento, os cientistas não podiam dizer nada de significativo em estudos observacionais em que o controle físico sobre os tratamentos é inviável.

Como o confundimento era definido então, e como deveria ser definido? Armados com o que agora sabemos sobre a lógica da causalidade, a resposta à segunda pergunta é mais fácil. A quantidade que observamos é a probabilidade condicional do desfecho dado o tratamento, \textit{P}(\textit{Y} | \textit{X}). A pergunta que queremos fazer à Natureza tem a ver com a relação causal entre \textit{X} e \textit{Y}, que é capturada pela probabilidade intervencionista \textit{P}(\textit{Y} | \textit{do}(\textit{X})). O confundimento, então, deveria ser simplesmente definido como qualquer coisa que leve a uma discrepância entre as duas: \textit{P}(\textit{Y} | \textit{X}) ≠ \textit{P}(\textit{Y} | \textit{do}(\textit{X})).

Por que tanta confusão? Infelizmente, as coisas não eram tão fáceis assim antes dos anos 1990, porque o operador \textit{do} ainda não havia sido formalizado. Ainda hoje, se você parar um estatístico na rua e perguntar “O que ‘confundimento’ significa para você?”, provavelmente obterá uma das respostas mais convolutas e confusas que jamais ouviu de um cientista. Um livro recente, coescrito por estatísticos de destaque, gasta literalmente duas páginas tentando explicá-lo, e ainda não encontrei um leitor que tenha compreendido a explicação.

A razão da dificuldade é que o confundimento não é uma noção estatística. Ele representa a discrepância entre o que queremos avaliar (o efeito causal) e o que de fato avaliamos usando métodos estatísticos. Se você não consegue articular matematicamente o que quer avaliar, não pode esperar definir o que constitui uma discrepância.

Historicamente, o conceito de “confundimento” evoluiu em torno de duas concepções relacionadas: a incomparabilidade e uma terceira variável à espreita. Ambos os conceitos resistiram à formalização. Quando falamos de comparabilidade, no contexto do experimento de Daniel, dissemos que os grupos de tratamento e de controle deveriam ser idênticos em todos os aspectos relevantes. Mas isso nos obriga a distinguir atributos relevantes de irrelevantes. Como sabemos que a idade é relevante no estudo de caminhada de Honolulu? Como sabemos que a ordem alfabética do nome de um participante não é relevante? Você pode dizer que é óbvio ou senso comum, mas gerações de cientistas lutaram para articular formalmente esse senso comum, e um robô não pode confiar em nosso senso comum quando se pede que aja corretamente.

A mesma ambiguidade assola a definição da terceira variável. Um confundidor deve ser uma causa comum de \textit{X} e \textit{Y}, ou apenas correlacionado com cada um? Hoje podemos responder a tais perguntas recorrendo ao diagrama causal e verificando quais variáveis produzem uma discrepância entre \textit{P}(\textit{X} | \textit{Y}) e \textit{P}(\textit{X} | \textit{do}(\textit{Y})). Sem um diagrama ou um operador \textit{do}, cinco gerações de estatísticos e cientistas da saúde tiveram de lutar com substitutos, nenhum dos quais era satisfatório. Considerando que os medicamentos no seu armário de remédios podem ter sido desenvolvidos com base numa definição duvidosa de “confundidores”, você deveria estar um tanto preocupado.

Vejamos algumas das definições substitutas de confundimento. Elas se enquadram em duas categorias principais: declarativas e procedimentais. Uma definição declarativa típica (e errada) seria: “Um confundidor é qualquer variável que esteja correlacionada tanto com \textit{X} quanto com \textit{Y}.” Por outro lado, uma definição procedimental tentaria caracterizar um confundidor em termos de um teste estatístico. Isso atrai os estatísticos, que adoram qualquer teste que possa ser realizado diretamente nos dados sem apelar a um modelo. Aqui está uma definição procedimental que leva o nome assustador de “não colapsabilidade”. Vem de um artigo de 1996 do epidemiologista norueguês Sven Hernberg: “Formalmente, pode-se comparar o risco relativo bruto e o risco relativo resultante após o ajustamento pelo potencial confundidor. Uma diferença indica confundimento e, nesse caso, deve-se usar a estimativa de risco ajustada. Se não houver diferença ou ela for negligenciável, o confundimento não é uma questão e a estimativa bruta deve ser preferida.” Em outras palavras, se você suspeita de um confundidor, tente ajustar por ele e tente não ajustar por ele. Se houver diferença, é um confundidor, e você deve confiar no valor ajustado. Se não houver diferença, está livre de problemas.

Hernberg de modo algum foi a primeira pessoa a defender tal abordagem; ela tem desorientado um século de epidemiologistas, economistas e cientistas sociais, e ainda reina em certos círculos da estatística aplicada. Escolhi Hernberg apenas porque ele foi excepcionalmente explícito a respeito e porque escreveu isso em 1996, bem depois de a Revolução Causal já estar em andamento.

A mais popular das definições declarativas evoluiu ao longo de um período. Alfredo Morabia, autor de \textit{A History of Epidemiologic Methods and Concepts} [\textit{Uma história dos métodos e conceitos epidemiológicos}], chama-a de “a definição epidemiológica clássica de confundimento”, e ela consiste em três partes. Um confundidor de \textit{X} (o tratamento) e \textit{Y} (o desfecho) é uma variável \textit{Z} que é (1) associada a \textit{X} na população em geral, e (2) associada a \textit{Y} entre as pessoas que não foram expostas ao tratamento \textit{X}. Em anos recentes, isso foi suplementado por uma terceira condição: (3) \textit{Z} não deve estar no caminho causal entre \textit{X} e \textit{Y}. Observe que todos os termos na versão “clássica” (1 e 2) são estatísticos. Em particular, \textit{Z} é apenas assumida como associada a — e não como causa de — \textit{X} e \textit{Y}.

Edward Simpson propôs a condição bastante convoluta “\textit{Y} está associada a \textit{Z} entre os não expostos” em 1951. Do ponto de vista causal, parece que a ideia de Simpson era descontar a parte da correlação de \textit{Z} com \textit{Y} que se deve ao efeito causal de \textit{X} sobre \textit{Y}; em outras palavras, ele queria dizer que \textit{Z} tem um efeito sobre \textit{Y} independente de seu efeito sobre \textit{X}. A única maneira que ele conseguiu pensar para expressar esse desconto foi condicionar em \textit{X} focando no grupo de controle (\textit{X} = 0). O vocabulário estatístico, privado da palavra “efeito”, não lhe deu outra maneira de dizê-lo.

Se isso é um tanto confuso, \textit{deveria} ser! Quanto mais fácil teria sido se ele simplesmente tivesse escrito um diagrama causal, como a Figura 4.1, e dito: “\textit{Y} está associada a \textit{Z} por caminhos que não passam por \textit{X}.” Mas ele não tinha essa ferramenta, e não podia falar de caminhos, que eram um conceito proibido.

A “definição epidemiológica clássica” de um confundidor tem outros defeitos, como mostram os dois exemplos a seguir:

(i) \textit{X} → \textit{Z} → \textit{Y}

(ii) \textit{X} → \textit{M} → \textit{Y}, com \textit{Z} como proxy de \textit{M}

No exemplo (i), \textit{Z} satisfaz as condições (1) e (2), mas não é um confundidor. É conhecido como mediador: é a variável que explica o efeito causal de \textit{X} sobre \textit{Y}. É um desastre controlar por \textit{Z} se você está tentando encontrar o efeito causal de \textit{X} sobre \textit{Y}. Se você olhar apenas para aqueles indivíduos nos grupos de tratamento e de controle para os quais \textit{Z} = 0, então bloqueou completamente o efeito de \textit{X}, porque ele funciona mudando \textit{Z}. Assim, você concluirá que \textit{X} não tem efeito sobre \textit{Y}. É exatamente isso que Ezra Klein quis dizer quando afirmou: “Às vezes você acaba controlando pela coisa que está tentando medir.”

No exemplo (ii), \textit{Z} é um \textit{proxy} (substituto) do mediador \textit{M}. Os estatísticos muito frequentemente controlam por \textit{proxies} quando a variável causal real não pode ser medida; por exemplo, a filiação partidária pode ser usada como \textit{proxy} de crenças políticas. Como \textit{Z} não é uma medida perfeita de \textit{M}, parte da influência de \textit{X} sobre \textit{Y} pode “vazar” se você controlar por \textit{Z}. Não obstante, controlar por \textit{Z} ainda é um erro. Embora o viés possa ser menor do que se você controlasse por \textit{M}, ele ainda está lá. Por essa razão, estatísticos posteriores, notavelmente David Cox em seu livro-texto \textit{The Design of Experiments} [\textit{O desenho de experimentos}] (1958), advertiram que você só deveria controlar por \textit{Z} se tivesse uma “forte razão a priori” para acreditar que ela não é afetada por \textit{X}. Essa “forte razão a priori” não é nada mais, nada menos do que uma premissa causal. Ele acrescenta: “Tais hipóteses podem ser perfeitamente em ordem, mas o cientista deve estar sempre ciente de quando está apelando a elas.” Lembre-se de que estamos em 1958, em meio à grande proibição da causalidade. Cox está dizendo que você pode tomar um gole de pinga causal ao ajustar por confudidores, mas não conte ao padre. Uma sugestão ousada! Nunca deixo de elogiá-lo por sua bravura.

Por volta de 1980, as condições de Simpson e de Cox haviam sido combinadas no teste de três partes para o confundimento que mencionei acima. Ele é tão confiável quanto uma canoa com apenas três vazamentos. Embora faça um apelo tímido à causalidade na parte (3), cada uma das duas primeiras partes pode ser mostrada como desnecessária e insuficiente. Greenland e Robins chegaram a essa conclusão em seu artigo histórico de 1986.

Os dois adotaram uma abordagem completamente nova ao confundimento, que chamaram de “trocabilidade” (\textit{exchangeability}). Voltaram à ideia original de que o grupo de controle (\textit{X} = 0) deveria ser comparável ao grupo de tratamento (\textit{X} = 1). Mas acrescentaram um toque contrafactual. (Lembre-se, do Capítulo 1, de que os contrafactuais estão no degrau três da Escada da Causalidade e, portanto, são poderosos o bastante para detectar o confundimento.) A trocabilidade exige que o pesquisador considere o grupo de tratamento, imagine o que teria acontecido a seus membros se não tivessem recebido o tratamento, e então julgue se o desfecho seria o mesmo que o daqueles que (na realidade) não receberam o tratamento. Só então podemos dizer que não existe confundimento no estudo.

Em 1986, falar de contrafactuais a uma plateia de epidemiologistas exigia alguma coragem, porque eles ainda estavam muito sob a influência da estatística clássica, que sustenta que todas as respostas estão nos dados — não no que poderia ter sido, que permanecerá para sempre não observado. No entanto, a comunidade estatística estava de certo modo preparada para ouvir tal heresia, graças ao trabalho pioneiro de outro estatístico de Harvard, Donald Rubin. No arcabouço de “resultados potenciais” de Rubin, proposto em 1974, variáveis contrafactuais como “Pressão Arterial da Pessoa \textit{X} caso ela tivesse recebido o Medicamento D” e “Pressão Arterial da Pessoa \textit{X} caso ela não tivesse recebido o Medicamento D” são tão legítimas quanto uma variável tradicional como Pressão Arterial — apesar do fato de que uma dessas duas variáveis permanecerá para sempre não observada.

Robins e Greenland se propuseram a expressar sua concepção de confundimento em termos de resultados potenciais. Particionaram a população em quatro tipos de indivíduos: condenados, causativos, preventivos e imunes. A linguagem é sugestiva, então pensemos no tratamento \textit{X} como uma vacina contra a gripe e no desfecho \textit{Y} como contrair gripe. As pessoas condenadas são aquelas para as quais a vacina não funciona; terão gripe quer tomem a vacina ou não. O grupo causativo (que pode ser inexistente) inclui aqueles para os quais a vacina de fato causa a doença. O grupo preventivo consiste em pessoas para as quais a vacina previne a doença: terão gripe se não forem vacinadas e não terão gripe se forem vacinadas. Por fim, o grupo imune consiste em pessoas que não terão gripe em nenhum dos casos. A Tabela 4.1 resume essas considerações.

\begin{figure}[ht]
\centering
\includegraphics[width=0.92\textwidth]{../figuras/tab_4_1.png}
\caption{Tabela 4.1. resume essas considerações.}
\end{figure}
\textit{Tabela 4.1. resume essas considerações.}

Idealmente, cada pessoa teria um adesivo na testa identificando a que grupo pertencia. A trocabilidade simplesmente significa que a porcentagem de pessoas com cada tipo de adesivo (\textit{d}\%, \textit{c}\%, \textit{p}\% e \textit{i}\%, respectivamente) deveria ser a mesma nos grupos de tratamento e de controle. A igualdade entre essas proporções garante que o desfecho seria exatamente o mesmo se trocássemos tratamentos e controles. Caso contrário, os grupos de tratamento e de controle não são iguais, e nossa estimativa do efeito da vacina estará confundida. Observe que os dois grupos podem ser diferentes de muitas maneiras. Podem diferir em idade, sexo, condições de saúde e uma variedade de outras características. Apenas a igualdade entre \textit{d}, \textit{c}, \textit{p} e \textit{i} determina se são trocáveis ou não. Assim, a trocabilidade equivale à igualdade entre dois conjuntos de quatro proporções — uma vasta redução de complexidade em relação à alternativa de avaliar os inumeráveis fatores pelos quais os dois grupos podem diferir.

\textbf{TABELA 4.1.} Classificação de indivíduos segundo o tipo de resposta.

Usando essa definição de senso comum de confundimento, Greenland e Robins mostraram que as definições “estatísticas”, tanto declarativas quanto procedimentais, dão respostas incorretas. Uma variável pode satisfazer o teste de três partes dos epidemiologistas e ainda assim aumentar o viés, se for ajustada. A definição de Greenland e Robins foi uma grande conquista, porque lhes permitiu dar exemplos explícitos mostrando que as definições anteriores de confundimento eram inadequadas. No entanto, a definição não podia ser traduzida em prática. Simplificando, aqueles adesivos na testa não existem. Nem sequer temos uma contagem das proporções \textit{d}, \textit{c}, \textit{p} e \textit{i}. De fato, esse é precisamente o tipo de informação que o gênio da Natureza mantém trancado dentro de sua lâmpada mágica e não mostra a ninguém. Sem essa informação, o pesquisador fica a intuir se os grupos de tratamento e de controle são trocáveis ou não.

Neste ponto, espero que sua curiosidade esteja bem despertada. Como os diagramas causais podem transformar essa dor de cabeça maciça do confundimento num jogo divertido? O truque reside num teste operacional para o confundimento, chamado critério da porta dos fundos. Esse critério transforma o problema de definir o confundimento, identificar confudidores e ajustá-los num quebra-cabeça de rotina que não é mais desafiador do que resolver um labirinto. Assim, trouxe o espinhoso e antigo problema a uma conclusão feliz.

\section{O OPERADOR \textit{DO} E O CRITÉRIO DA PORTA DOS FUNDOS}
Para compreender o critério da porta dos fundos, ajuda primeiro ter um senso intuitivo de como a informação flui num diagrama causal. Gosto de pensar nos elos como canos que transportam informação de um ponto de partida \textit{X} a um fim \textit{Y}. Tenha em mente que o transporte de informação vai em ambas as direções, causal e não causal, como vimos no Capítulo 3.

De fato, os caminhos não causais são precisamente a fonte do confundimento. Lembre-se de que defino o confundimento como qualquer coisa que faça \textit{P}(\textit{Y} | \textit{do}(\textit{X})) diferir de \textit{P}(\textit{Y} | \textit{X}). O operador \textit{do} apaga todas as setas que chegam a \textit{X} e, dessa forma, impede que qualquer informação sobre \textit{X} flua na direção não causal. A randomização tem o mesmo efeito. O mesmo acontece com o ajustamento estatístico, se escolhermos as variáveis certas para ajustar.

No capítulo anterior, examinamos três regras que nos dizem como interromper o fluxo de informação através de qualquer junção individual. Vou repeti-las para ênfase:

\textbf{(a)} Numa junção em cadeia, \textit{A} → \textit{B} → \textit{C}, controlar por \textit{B} impede que a informação sobre \textit{A} chegue a \textit{C} ou vice-versa.

\textbf{(b)} Do mesmo modo, numa junção em garfo ou de confundimento, \textit{A} ← \textit{B} → \textit{C}, controlar por \textit{B} impede que a informação sobre \textit{A} chegue a \textit{C} ou vice-versa.

\textbf{(c)} Por fim, num colisor, \textit{A} → \textit{B} ← \textit{C}, valem exatamente as regras opostas. As variáveis \textit{A} e \textit{C} começam independentes, de modo que a informação sobre \textit{A} não lhe diz nada sobre \textit{C}. Mas se você controlar por \textit{B}, a informação começa a fluir pelo “cano”, devido ao efeito de “explicar-afastando” (\textit{explain-away}).

Devemos também ter em mente outra regra fundamental:

\textbf{(d)} Controlar por descendentes (ou \textit{proxies}) de uma variável é como “parcialmente” controlar pela variável em si. Controlar por um descendente de um mediador fecha parcialmente o cano; controlar por um descendente de um colisor abre parcialmente o cano.

Agora, e se tivermos canos mais longos com mais junções, como este: \textit{A} → \textit{B} → \textit{C} → \textit{D} ← \textit{E} ← \textit{F} → \textit{G} ← \textit{H} → \textit{I} → \textit{J}? A resposta é muito simples: se uma única junção estiver bloqueada, então \textit{J} não pode “ficar sabendo” de \textit{A} por esse caminho. Assim, temos muitas opções para bloquear a comunicação entre \textit{A} e \textit{J}: controlar por \textit{B}, controlar por \textit{C}, \textit{não} controlar por \textit{D} (porque é um colisor), controlar por \textit{E}, e assim por diante. Qualquer uma delas é suficiente. É por isso que o procedimento estatístico usual de controlar por tudo o que podemos medir é tão equivocado. De fato, este caminho particular está bloqueado se \textit{não} controlarmos por nada! Os colisores em \textit{D} e \textit{G} bloqueiam o caminho sem nenhuma ajuda externa. Controlar por \textit{D} e \textit{G} abriria esse caminho e permitiria que \textit{J} “ouvisse” \textit{A}.

Por fim, para desconfundir duas variáveis \textit{X} e \textit{Y}, precisamos apenas bloquear todo caminho não causal entre elas sem bloquear ou perturbar nenhum caminho causal. Mais precisamente, um caminho pela porta dos fundos é qualquer caminho de \textit{X} a \textit{Y} que começa com uma seta apontando para \textit{X}. \textit{X} e \textit{Y} estarão desconfundidas se bloquearmos todo caminho pela porta dos fundos (porque tais caminhos permitem correlação espúria entre \textit{X} e \textit{Y}). Se fizermos isso controlando por algum conjunto de variáveis \textit{Z}, também precisamos garantir que nenhum membro de \textit{Z} seja um descendente de \textit{X} num caminho causal; caso contrário, poderíamos fechar parcial ou completamente esse caminho.

É só isso! Com essas regras, o desconfundimento se torna tão simples e divertido que você pode tratá-lo como um jogo. Incentivo-o a experimentar alguns exemplos só para pegar o jeito e ver como é fácil. Se ainda achar difícil, fique tranquilo: existem algoritmos que podem resolver todos esses problemas em questão de nanossegundos. Em cada caso, o objetivo do jogo é especificar um conjunto de variáveis que desconfundirão \textit{X} e \textit{Y}. Em outras palavras, elas não devem ser descendentes de \textit{X}, e devem bloquear todos os caminhos pela porta dos fundos.

\textbf{JOGO 1.}

Este é fácil! Não há setas apontando para \textit{X}; portanto, não há caminhos pela porta dos fundos. Não precisamos controlar por nada. Não obstante, alguns pesquisadores considerariam \textit{B} um confundidor. Ela está associada a \textit{X} por causa da cadeia \textit{X} → \textit{A} → \textit{B}. Está associada a \textit{Y} entre indivíduos com \textit{X} = 0 porque há um caminho aberto \textit{B} ← \textit{A} → \textit{Y} que não passa por \textit{X}. E \textit{B} não está no caminho causal \textit{X} → \textit{A} → \textit{Y}. Portanto, passa na “definição epidemiológica clássica” de três passos para o confundimento, mas não passa no critério da porta dos fundos e levará a um desastre se for controlada.

\textbf{JOGO 2.}

Neste exemplo, você deve pensar em \textit{A}, \textit{B}, \textit{C} e \textit{D} como variáveis “pré-tratamento”. (O tratamento, como de costume, é \textit{X}.) Agora há um caminho pela porta dos fundos \textit{X} ← \textit{A} → \textit{B} ← \textit{D} → \textit{E} → \textit{Y}. Esse caminho já está bloqueado pelo colisor em \textit{B}, de modo que não precisamos controlar por nada. Muitos estatísticos controlariam por \textit{B} ou \textit{C}, pensando que não há mal nisso desde que ocorram antes do tratamento. Um estatístico de destaque escreveu recentemente: “Evitar condicionar em algumas covariáveis observadas… é improviso ad hoc antissientífico.” Ele está errado; condicionar em \textit{B} ou \textit{C} é uma má ideia porque abriria o caminho não causal e, portanto, confundiria \textit{X} e \textit{Y}. Observe que, neste caso, poderíamos refechar o caminho controlando por \textit{A} ou \textit{D}. Este exemplo mostra que pode haver diferentes estratégias de desconfundimento. Um pesquisador pode tomar o caminho fácil e não controlar por nada; um pesquisador mais tradicional pode controlar por \textit{C} e \textit{D}. Ambos estariam corretos e deveriam obter o mesmo resultado (desde que o modelo esteja correto e tenhamos uma amostra grande o bastante).

\textbf{JOGO 3.}

Nos Jogos 1 e 2 você não precisou fazer nada, mas desta vez precisa. Há um caminho pela porta dos fundos de \textit{X} a \textit{Y}, \textit{X} ← \textit{B} → \textit{Y}, que só pode ser bloqueado controlando por \textit{B}. Se \textit{B} for inobservável, então não há maneira de estimar o efeito de \textit{X} sobre \textit{Y} sem realizar um experimento controlado randomizado. Alguns (de fato, a maioria dos) estatísticos nessa situação controlariam por \textit{A}, como \textit{proxy} da variável inobservável \textit{B}, mas isso só elimina parcialmente o viés de confundimento e introduz um novo viés de colisor.

\textbf{JOGO 4.}

Este introduz um novo tipo de viés, chamado “viés em M” (\textit{M-bias}, nomeado pela forma do grafo). Mais uma vez, há apenas um caminho pela porta dos fundos, e ele já está bloqueado por um colisor em \textit{B}. Então não precisamos controlar por nada. Não obstante, todos os estatísticos antes de 1986 e muitos ainda hoje considerariam \textit{B} um confundidor. Ela está associada a \textit{X} (via \textit{X} ← \textit{A} → \textit{B}) e associada a \textit{Y} via um caminho que não passa por \textit{X} (\textit{B} ← \textit{C} → \textit{Y}). Não se situa num caminho causal e não é descendente de nada num caminho causal, porque não há caminho causal de \textit{X} a \textit{Y}. Portanto, \textit{B} passa no teste tradicional de três passos para um confundidor.

O viés em M põe o dedo no que está errado com a abordagem tradicional. É incorreto chamar uma variável, como \textit{B}, de confundidor meramente porque está associada tanto a \textit{X} quanto a \textit{Y}. Reiterando: \textit{X} e \textit{Y} estão desconfundidas se \textit{não} controlarmos por \textit{B}. \textit{B} só se torna um confundidor quando você controla por ela!

Quando comecei a mostrar esse diagrama a estatísticos nos anos 1990, alguns riram e disseram que tal diagrama era extremamente improvável de ocorrer na prática. Discordo! Por exemplo, o uso de cinto de segurança (\textit{B}) não tem efeito causal sobre o fumo (\textit{X}) ou a doença pulmonar (\textit{Y}); é meramente um indicador das atitudes de uma pessoa em relação a normas sociais (\textit{A}), bem como a medidas de segurança e saúde (\textit{C}). Algumas dessas atitudes podem afetar a suscetibilidade à doença pulmonar (\textit{Y}). Na prática, o uso de cinto de segurança foi encontrado correlacionado tanto com \textit{X} quanto com \textit{Y}; de fato, num estudo conduzido em 2006 como parte de um litígio do tabaco, o uso de cinto de segurança foi listado como uma das primeiras variáveis a serem controladas. Se você aceita o modelo acima, então controlar apenas por \textit{B} seria um erro. Observe que não há problema em controlar por \textit{B} se você também controlar por \textit{A} ou \textit{C}. Controlar pelo colisor \textit{B} abre o “cano”, mas controlar por \textit{A} ou \textit{C} o fecha novamente. Infelizmente, no exemplo do cinto de segurança, \textit{A} e \textit{C} são variáveis relacionadas às atitudes das pessoas e não são propensas a ser observáveis. Se você não pode observar, não pode ajustar.

\textbf{JOGO 5.}

O Jogo 5 é apenas o Jogo 4 com uma pequena ruga extra. Agora um segundo caminho pela porta dos fundos \textit{X} ← \textit{B} → \textit{C} → \textit{Y} precisa ser fechado. Se fecharmos esse caminho controlando por \textit{B}, então abrimos o caminho em forma de M \textit{X} ← \textit{A} → \textit{B} ← \textit{C} → \textit{Y}. Para fechar esse caminho, devemos controlar também por \textit{A} ou \textit{C}. No entanto, observe que poderíamos apenas controlar por \textit{C} sozinha; isso fecharia o caminho \textit{X} ← \textit{B} → \textit{C} → \textit{Y} e não afetaria o outro caminho.

Os Jogos 1 a 3 vêm de um artigo de 1993 de Clarice Weinberg, vice-chefe nos Institutos Nacionais de Saúde, intitulado “Toward a Clearer Definition of Confounding” [“Rumo a uma definição mais clara de confundimento”]. Saiu durante o período de transição entre 1986 e 1995, quando o artigo de Greenland e Robins já estava disponível, mas os diagramas causais ainda não eram amplamente conhecidos. Weinberg, portanto, passou pelo considerável exercício aritmético de verificar a trocabilidade em cada um dos casos mostrados. Embora tenha usado exibições gráficas para comunicar os cenários envolvidos, não usou a lógica dos diagramas para ajudar a distinguir confudidores de desconfundidores. É a única pessoa de quem tenho conhecimento que conseguiu esse feito. Mais tarde, em 2012, colaborou numa versão atualizada que analisa os mesmos exemplos com diagramas causais e verifica que todas as suas conclusões de 1993 estavam corretas.

Em ambos os artigos de Weinberg, a aplicação médica era estimar o efeito do fumo (\textit{X}) sobre abortos espontâneos (\textit{Y}). No Jogo 1, \textit{A} representa uma anormalidade subjacente induzida pelo fumo; essa não é uma variável observável porque não sabemos qual é a anormalidade. \textit{B} representa um histórico de abortos espontâneos anteriores. É muito, muito tentador para um epidemiologista levar em conta abortos espontâneos anteriores e ajustá-los ao estimar a probabilidade de futuros abortos. Mas isso é a coisa errada a se fazer aqui! Ao fazê-lo, estamos parcialmente inativando o mecanismo pelo qual o fumo age e, assim, subestimaremos o verdadeiro efeito do fumo.

O Jogo 2 é uma versão mais complicada em que há duas variáveis de fumo diferentes: \textit{X} representa se a mãe fuma agora (no início da segunda gravidez), enquanto \textit{A} representa se ela fumou durante a primeira gravidez. \textit{B} e \textit{E} são anormalidades subjacentes causadas pelo fumo, que são inobserváveis, e \textit{D} representa outras causas fisiológicas dessas anormalidades. Observe que esse diagrama permite o fato de que a mãe poderia ter mudado o comportamento de fumar entre as gestações, mas as outras causas fisiológicas não mudariam. Mais uma vez, muitos epidemiologistas ajustariam por abortos anteriores (\textit{C}), mas isso é uma má ideia a menos que você também ajuste pelo comportamento de fumar na primeira gravidez (\textit{A}).

Os Jogos 4 e 5 vêm de um artigo publicado em 2014 por Andrew Forbes, bioestatístico da Universidade Monash na Austrália, junto com vários colaboradores. Ele está interessado no efeito do fumo sobre a asma em adultos. No Jogo 4, \textit{X} representa o comportamento de fumar de um indivíduo, e \textit{Y} representa se a pessoa tem asma quando adulta. \textit{B} representa a asma na infância, que é um colisor porque é afetada tanto por \textit{A}, fumo dos pais, quanto por \textit{C}, uma predisposição subjacente (e inobservável) à asma. No Jogo 5, as variáveis têm os mesmos significados, mas Forbes acrescentou duas setas para maior realismo. (O Jogo 4 destinava-se apenas a introduzir o grafo em M.) De fato, o modelo completo no artigo de Forbes tem mais algumas variáveis e se parece com o diagrama da Figura 4.7.

\begin{figure}[ht]
\centering
\includegraphics[width=0.92\textwidth]{../figuras/fig_4_7.png}
\caption{Figura 4.7. Observe que o Jogo 5 está embutido nesse modelo no sentido de que as variáveis \textit{A}, \textit{B}, \textit{C}, \textit{X} e \textit{Y} têm exatamente as mesmas relações. Assim, podemos transferir nossas conclusões e concluir que temos de controlar por \textit{A} e \textit{B}, ou por \textit{C}; mas \textit{C} é uma variável inobservável e, portanto, incontrolável. Além disso, temos quatro novas variáveis de confundimento: \textit{D} = asma dos pais, \textit{E} = bronquite crônica, \textit{F} = sexo e \textit{G} = status socioeconômico. O leitor pode se divertir descobrindo que devemos controlar por \textit{E}, \textit{F} e \textit{G}, mas não há necessidade de controlar por \textit{D}. Assim, um conjunto suficiente de variáveis para o desconfundimento é \textit{A}, \textit{B}, \textit{E}, \textit{F} e \textit{G}.}
\end{figure}
\textit{Figura 4.7. Observe que o Jogo 5 está embutido nesse modelo no sentido de que as variáveis }A\textit{, }B\textit{, }C\textit{, }X\textit{ e }Y\textit{ têm exatamente as mesmas relações. Assim, podemos transferir nossas conclusões e concluir que temos de controlar por }A\textit{ e }B\textit{, ou por }C\textit{; mas }C\textit{ é uma variável inobservável e, portanto, incontrolável. Além disso, temos quatro novas variáveis de confundimento: }D\textit{ = asma dos pais, }E\textit{ = bronquite crônica, }F\textit{ = sexo e }G\textit{ = status socioeconômico. O leitor pode se divertir descobrindo que devemos controlar por }E\textit{, }F\textit{ e }G\textit{, mas não há necessidade de controlar por }D\textit{. Assim, um conjunto suficiente de variáveis para o desconfundimento é }A\textit{, }B\textit{, }E\textit{, }F\textit{ e }G\textit{.}

Observe que o Jogo 5 está embutido nesse modelo no sentido de que as variáveis \textit{A}, \textit{B}, \textit{C}, \textit{X} e \textit{Y} têm exatamente as mesmas relações. Assim, podemos transferir nossas conclusões e concluir que temos de controlar por \textit{A} e \textit{B}, ou por \textit{C}; mas \textit{C} é uma variável inobservável e, portanto, incontrolável. Além disso, temos quatro novas variáveis de confundimento: \textit{D} = asma dos pais, \textit{E} = bronquite crônica, \textit{F} = sexo e \textit{G} = status socioeconômico. O leitor pode se divertir descobrindo que devemos controlar por \textit{E}, \textit{F} e \textit{G}, mas não há necessidade de controlar por \textit{D}. Assim, um conjunto suficiente de variáveis para o desconfundimento é \textit{A}, \textit{B}, \textit{E}, \textit{F} e \textit{G}.

\textbf{FIGURA 4.7.} Modelo de Andrew Forbes do fumo (\textit{X}) e da asma (\textit{Y}).

No final, Forbes descobriu que o fumo tinha uma associação pequena e estatisticamente insignificante com a asma adulta nos dados brutos, e o efeito se tornou ainda menor e mais insignificante após o ajustamento pelos confudidores. O resultado nulo não deve, porém, ofuscar o fato de que seu artigo é um modelo do “hábil interrogatório da Natureza”.

Um comentário final sobre esses “jogos”: quando você começa a identificar as variáveis como fumo, aborto espontâneo e assim por diante, elas são bastante obviamente não jogos, mas negócios sérios. Refiro-me a eles como jogos porque a alegria de ser capaz de resolvê-los de forma rápida e significativa é semelhante ao prazer que uma criança sente ao descobrir que consegue desvendar quebra-cabeças que a desconcertavam antes. Poucos momentos numa carreira científica são tão satisfatórios quanto tomar um problema que intrigou e confundeu gerações de predecessores e reduzi-lo a um jogo ou algoritmo direto. Considero a solução completa do problema do confundimento um dos principais destaques da Revolução Causal porque encerrou uma era de confusão que provavelmente resultou em muitas decisões erradas no passado. Tem sido uma revolução silenciosa, fervilhando principalmente em laboratórios de pesquisa e reuniões científicas. Ainda assim, armada com essas novas ferramentas e insights, a comunidade científica está agora enfrentando problemas mais difíceis, tanto teóricos quanto práticos, como os capítulos seguintes mostrarão.

\chapter{5 O DEBATE ENEVOADO PELA FUMAÇA: LIMPANDO O AR}
\begin{quote}Finalmente os marinheiros disseram uns aos outros: Vinde, e lancemos sortes, para que saibamos por causa de quem nos sobreveio este mal.\end{quote}
>

\begin{quote}— JONAS 1:7\end{quote}
No final dos anos 1950 e início dos 1960, estatísticos e médicos se digladiaram em torno de uma das questões médicas de maior destaque do século: O fumo causa câncer de pulmão? Meio século depois desse debate, damos a resposta como certa. Mas, naquela época, a questão de modo algum estava clara. Cientistas, e até famílias, estavam divididos.

A família de Jacob Yerushalmy era uma dessas famílias divididas. Bioestatístico da Universidade da Califórnia em Berkeley, Yerushalmy (1904–1973) foi um dos últimos resistentes pró-tabaco na academia. “Yerushalmy opôs-se à noção de que o fumo de cigarro causava câncer até o dia de sua morte”, escreveu seu sobrinho David Lilienfeld muitos anos depois. Por outro lado, o pai de Lilienfeld, Abe Lilienfeld, era epidemiologista da Universidade Johns Hopkins e um dos mais veementes defensores da teoria de que o fumo de fato causava câncer. O jovem Lilienfeld recordava como o tio “Yak” (abreviação de Jacob) e seu pai se sentavam e debatiam os efeitos do fumo, envoltos o tempo todo numa “névoa de fumaça do cigarro de Yak e do cachimbo de Abe” (ver frontispício do capítulo).

Se ao menos Abe e Yak tivessem podido convocar a Revolução Causal para limpar o ar! Como este capítulo mostra, um dos argumentos científicos mais importantes contra a hipótese fumo-câncer era a possível existência de fatores não medidos que causam tanto a ânsia por nicotina quanto o câncer de pulmão. Acabamos de discutir tais padrões de confundimento e notamos que os diagramas causais de hoje expulsaram a ameaça do confundimento da existência. Mas estamos agora nos anos 1950 e 1960, duas décadas antes de Sander Greenland e Jamie Robins e três décadas antes de alguém ter ouvido falar do operador \textit{do}. É interessante examinar, portanto, como os cientistas daquela era lidaram com a questão e mostraram que o argumento do confundimento não passa de fumaça e espelhos.

Sem dúvida, o tema de muitos dos debates enevoados de Abe e Yak não era nem o tabaco nem o câncer. Era aquela palavra inocente: “causava”. Não era a primeira vez que médicos enfrentavam perguntas causais perplexas: alguns dos maiores marcos da história da medicina lidaram com a identificação de agentes causadores. Em meados dos anos 1700, James Lind descobriu que as frutas cítricas podiam prevenir o escorbuto, e em meados dos anos 1800 John Snow descobriu que a água contaminada com matéria fecal causava a cólera. (Pesquisas posteriores identificaram um agente causal mais específico em cada caso: deficiência de vitamina C para o escorbuto, o bacilo da cólera para a cólera.) Esses brilhantes trabalhos de detetive tinham em comum uma afortunada relação um-para-um entre causa e efeito. O bacilo da cólera é a \textit{única} causa da cólera; ou, como diríamos hoje, é tanto necessário quanto suficiente. Se você não for exposto a ele, não pegará a doença. Do mesmo modo, uma deficiência de vitamina C é necessária para produzir o escorbuto e, dado tempo suficiente, também é suficiente.

O debate fumo-câncer desafiou esse conceito monolítico de causalidade. Muitas pessoas fumam a vida inteira e nunca têm câncer de pulmão. Por outro lado, algumas pessoas têm câncer de pulmão sem jamais acender um cigarro. Algumas podem tê-lo por predisposição hereditária, outras por exposição a carcinógenos, e algumas por ambas as razões. É claro que os estatísticos já conheciam uma excelente maneira de estabelecer a causalidade num sentido mais geral: o ensaio clínico randomizado (ECR). Mas tal estudo não seria nem viável nem ético no caso do fumo. Como você poderia designar pessoas escolhidas ao acaso para fumar por décadas, possivelmente arruinando sua saúde, só para ver se teriam câncer de pulmão depois de trinta anos? É impossível imaginar alguém fora da Coreia do Norte “se voluntariando” para tal estudo.

Sem um ensaio clínico randomizado, não havia maneira de convencer céticos como Yerushalmy e R. A. Fisher, que estavam comprometidos com a ideia de que a associação observada entre fumo e câncer de pulmão era espúria. Para eles, algum terceiro fator à espreita poderia estar produzindo a associação observada. Por exemplo, poderia haver um gene do fumo que fizesse as pessoas desejarem cigarros e também, ao mesmo tempo, as tornasse mais propensas a desenvolver câncer de pulmão (talvez por causa de outras escolhas de estilo de vida). Os confudidores que sugeriam eram, na melhor das hipóteses, implausíveis. Ainda assim, o ônus recaía sobre o contingente antitabaco de provar que não havia confundidor — provar um negativo, o que Fisher e Yerushalmy bem sabiam ser quase impossível.

A ruptura final desse impasse é um relato ao mesmo tempo de um grande triunfo e de uma grande oportunidade perdida. Foi um triunfo para a saúde pública porque os epidemiologistas acertaram no final. O relatório do cirurgião-geral dos EUA, em 1964, declarou em termos inequívocos: “O fumo de cigarro está causalmente relacionado ao câncer de pulmão em homens.” Essa afirmação contundente encerrou para sempre o argumento de que o fumo “não estava comprovado” como causa de câncer. A taxa de fumo nos Estados Unidos entre homens começou a diminuir no ano seguinte e agora é menos da metade do que era em 1964. Sem dúvida, milhões de vidas foram salvas e esperanças de vida prolongadas.

Por outro lado, o triunfo é incompleto. O período que levou para se chegar à conclusão acima, grosso modo de 1950 a 1964, poderia ter sido mais curto se os cientistas tivessem podido recorrer a uma teoria mais principiada da causalidade. E, o mais significativo do ponto de vista deste livro, os cientistas dos anos 1960 não montaram de fato tal teoria. Para justificar a afirmação de que o fumo causava câncer, o comitê do cirurgião-geral baseou-se numa série informal de diretrizes, chamadas critérios de Hill, em homenagem ao estatístico da Universidade de Londres Austin Bradford Hill. Cada um desses critérios tem exceções demonstráveis, embora coletivamente tenham um valor de senso comum e até sabedoria. Do mundo excessivamente metodológico de Fisher, as diretrizes de Hill nos levam ao reino oposto, a um mundo sem metodologia em que a causalidade é decidida com base em padrões qualitativos de tendências estatísticas. A Revolução Causal constrói uma ponte entre esses dois extremos, empoderando nosso senso intuitivo de causalidade com rigor matemático. Mas esse trabalho seria deixado para a geração seguinte.

\section{TABACO: UMA EPIDEMIA FEITA PELO HOMEM}
Em 1902, os cigarros compreendiam apenas 2\% do mercado de tabaco dos EUA; cuspideiras, em vez de cinzeiros, eram o símbolo mais onipresente do consumo de tabaco. Mas duas forças poderosas trabalharam juntas para mudar os hábitos dos americanos: a automação e a publicidade. Cigarros feitos à máquina facilmente superaram charutos e cachimbos artesanais com base na disponibilidade e no custo. Enquanto isso, a indústria do tabaco inventou e aperfeiçoou muitos truques do ofício da publicidade (ver Figura 5.1). Pessoas que assistiam à TV nos anos 1960 podem facilmente lembrar qualquer número de jingles cativantes de cigarro, de “You get a lot to like in a Marlboro” a “You’ve come a long way, baby”. Em 1952, a fatia dos cigarros no mercado de tabaco havia saltado de 2\% para 81\%, e o próprio mercado havia crescido dramaticamente.

\begin{figure}[ht]
\centering
\includegraphics[width=0.92\textwidth]{../figuras/fig_5_1.png}
\caption{Figura 5.1. ). Pessoas que assistiam à TV nos anos 1960 podem facilmente lembrar qualquer número de jingles cativantes de cigarro, de “You get a lot to like in a Marlboro” a “You’ve come a long way, baby”. Em 1952, a fatia dos cigarros no mercado de tabaco havia saltado de 2\% para 81\%, e o próprio mercado havia crescido dramaticamente.}
\end{figure}
\textit{Figura 5.1. ). Pessoas que assistiam à TV nos anos 1960 podem facilmente lembrar qualquer número de jingles cativantes de cigarro, de “You get a lot to like in a Marlboro” a “You’ve come a long way, baby”. Em 1952, a fatia dos cigarros no mercado de tabaco havia saltado de 2\% para 81\%, e o próprio mercado havia crescido dramaticamente.}

Essa mudança radical nos hábitos de um país teve ramificações inesperadas para a saúde pública. Já nos primeiros anos do século XX, havia suspeitas de que o fumo era insalubre, de que “irritava” a garganta e causava tosse. Por volta de meados do século, a evidência começou a se tornar bem mais sinistra. Antes dos cigarros, o câncer de pulmão era tão raro que um médico podia deparar com ele apenas uma vez numa vida de prática. Mas entre 1900 e 1950, a doença antes rara quadruplicou em frequência, e em 1960 se tornaria a forma mais comum de câncer entre homens. Uma mudança tão enorme na incidência de uma doença letal pedia uma explicação.

\textbf{FIGURA 5.1.} Anúncios altamente manipuladores destinavam-se a tranquilizar o público de que os cigarros não eram prejudiciais à saúde — incluindo este anúncio de 1948 do \textit{Journal of the American Medical Association} visando os próprios médicos. (Fonte: Da coleção do Stanford Research into the Impact of Tobacco Advertising.)

Com o benefício do retrospecto, é fácil apontar o dedo da culpa para o fumo. Se plotarmos as taxas de câncer de pulmão e o consumo de tabaco num gráfico (ver Figura 5.2), a conexão é impossível de ignorar. Mas dados de séries temporais são evidência pobre de causalidade. Muitas outras coisas haviam mudado entre 1900 e 1950 e eram igualmente plausíveis como culpadas: a pavimentação de estradas, a inalação de gases de gasolina com chumbo e a poluição do ar em geral. O epidemiologista britânico Richard Doll disse em 1991: “Os automóveis… eram um fator novo e, se eu tivesse de apostar dinheiro em alguma coisa na época, teria apostado nos escapamentos de motores ou possivelmente no alcatrão das estradas.”

\begin{figure}[ht]
\centering
\includegraphics[width=0.92\textwidth]{../figuras/fig_5_2.png}
\caption{Figura 5.2. ), a conexão é impossível de ignorar. Mas dados de séries temporais são evidência pobre de causalidade. Muitas outras coisas haviam mudado entre 1900 e 1950 e eram igualmente plausíveis como culpadas: a pavimentação de estradas, a inalação de gases de gasolina com chumbo e a poluição do ar em geral. O epidemiologista britânico Richard Doll disse em 1991: “Os automóveis… eram um fator novo e, se eu tivesse de apostar dinheiro em alguma coisa na época, teria apostado nos escapamentos de motores ou possivelmente no alcatrão das estradas.”}
\end{figure}
\textit{Figura 5.2. ), a conexão é impossível de ignorar. Mas dados de séries temporais são evidência pobre de causalidade. Muitas outras coisas haviam mudado entre 1900 e 1950 e eram igualmente plausíveis como culpadas: a pavimentação de estradas, a inalação de gases de gasolina com chumbo e a poluição do ar em geral. O epidemiologista britânico Richard Doll disse em 1991: “Os automóveis… eram um fator novo e, se eu tivesse de apostar dinheiro em alguma coisa na época, teria apostado nos escapamentos de motores ou possivelmente no alcatrão das estradas.”}

\textbf{FIGURA 5.2.} Um gráfico da taxa de consumo de cigarros per capita nos Estados Unidos (preto) e da taxa de mortalidade por câncer de pulmão entre homens (cinza) mostra uma semelhança impressionante: a curva do câncer é quase uma réplica da curva do fumo, atrasada em cerca de trinta anos. Não obstante, essa evidência é circunstancial, não prova de causalidade. Algumas datas-chave são assinaladas aqui, incluindo a publicação do artigo de Richard Doll e Austin Bradford Hill em 1950, que primeiro alertou muitos profissionais médicos para a associação entre fumo e câncer de pulmão. (Fonte: Gráfico de Maayan Harel, usando dados da American Cancer Society, dos Centers for Disease Control e do Office of the Surgeon General.)

O trabalho da ciência é pôr de lado a suposição e olhar para os fatos. Em 1948, Doll e Austin Bradford Hill uniram forças para ver se podiam aprender algo sobre as causas da epidemia de câncer. Hill havia sido o estatístico-chefe num ensaio clínico randomizado altamente bem-sucedido, publicado no início daquele ano, que provara que a estreptomicina — um dos primeiros antibióticos — era eficaz contra a tuberculose. O estudo, um marco na história da medicina, não apenas apresentou aos médicos as “drogas milagrosas”, mas também cimentou a reputação dos ensaios clínicos randomizados, que logo se tornaram o padrão para a pesquisa clínica em epidemiologia.

É claro que Hill sabia que um ECR era impossível neste caso, mas havia aprendido as vantagens de comparar um grupo de tratamento a um grupo de controle. Assim, propôs comparar pacientes que já haviam sido diagnosticados com câncer a um grupo de controle de voluntários saudáveis. Os membros de cada grupo foram entrevistados sobre seus comportamentos passados e históricos médicos. Para evitar viés, os entrevistadores não foram informados de quem tinha câncer e quem era controle.

Os resultados do estudo foram chocantes: dos 649 pacientes com câncer de pulmão entrevistados, todos menos dois haviam sido fumantes. Isso era uma improbabilidade estatística tão extrema que Doll e Hill não resistiram a calcular as chances exatas contra ela: 1,5 milhão para 1. Além disso, os pacientes com câncer de pulmão haviam sido, em média, fumantes mais pesados do que os controles, mas (numa inconsistência que R. A. Fisher mais tarde atacaria) uma porcentagem menor relatou inalar a fumaça.

O tipo de estudo que Doll e Hill conduziram é agora chamado de estudo de caso-controle porque compara “casos” (pessoas com uma doença) a controles. É claramente uma melhoria em relação a dados de séries temporais, porque os pesquisadores podem controlar por confudidores como idade, sexo e exposição a poluentes ambientais. Não obstante, o desenho de caso-controle tem algumas desvantagens óbvias. É retrospectivo; isso significa que estudamos pessoas que se sabe terem câncer e olhamos para trás para descobrir o porquê. A lógica de probabilidade também é invertida. Os dados nos dizem a probabilidade de que um paciente com câncer seja fumante, em vez da probabilidade de que um fumante terá câncer. É esta última probabilidade que realmente importa para uma pessoa que quer saber se deve fumar ou não.

Além disso, estudos de caso-controle admitem várias possíveis fontes de viés. Uma delas é chamada de viés de memória (\textit{recall bias}): embora Doll e Hill tenham garantido que os entrevistadores não soubessem os diagnósticos, os pacientes certamente sabiam se tinham câncer ou não. Isso poderia ter afetado suas recordações. Outro problema é o viés de seleção. Pacientes hospitalizados com câncer de modo algum eram uma amostra representativa da população, ou mesmo da população fumante.

Em suma, os resultados de Doll e Hill eram extremamente sugestivos, mas não podiam ser tomados como prova de que o fumo causa câncer. Os dois pesquisadores tiveram o cuidado, a princípio, de chamar a correlação de uma “associação”. Depois de descartar vários confudidores, arriscaram uma afirmação mais forte de que “o fumo é um fator, e um fator importante, na produção do carcinoma do pulmão”.

Nos anos seguintes, dezenove estudos de caso-controle conduzidos em diferentes países chegaram basicamente à mesma conclusão. Mas, como R. A. Fisher não se cansava de apontar, repetir um estudo enviesado dezenove vezes não prova nada. Continua enviesado. Fisher escreveu em 1957 que esses estudos “eram meras repetições de evidência do mesmo tipo, e é necessário tentar examinar se esse tipo é suficiente para qualquer conclusão científica”.

Doll e Hill perceberam que, se houvesse vieses ocultos nos estudos de caso-controle, a mera replicação não os superaria. Assim, em 1951, iniciaram um estudo prospectivo, para o qual enviaram questionários a 60.000 médicos britânicos sobre seus hábitos de fumar e os acompanharam adiante no tempo. (A American Cancer Society lançou um estudo semelhante e maior por volta da mesma época.) Mesmo em apenas cinco anos, algumas diferenças dramáticas emergiram. Fumantes pesados tinham uma taxa de mortalidade por câncer de pulmão vinte e quatro vezes a dos não fumantes. No estudo da American Cancer Society, os resultados eram ainda mais sombrios: os fumantes morriam de câncer de pulmão vinte e nove vezes mais frequentemente do que os não fumantes, e os fumantes pesados morriam noventa vezes mais frequentemente. Por outro lado, pessoas que haviam fumado e então pararam reduziam seu risco por um fator de dois.

A consistência de todos esses resultados — mais fumo leva a um risco maior de câncer, parar leva a um risco menor — era outra forte peça de evidência de causalidade. Os médicos a chamam de “efeito dose-resposta”: se a substância \textit{A} causa um efeito biológico \textit{B}, então geralmente (embora nem sempre) uma dose maior de \textit{A} causa uma resposta \textit{B} mais forte.

Não obstante, céticos como Fisher e Yerushalmy não se convenceriam. Os estudos prospectivos ainda falhavam em comparar fumantes a não fumantes de outro modo idênticos. De fato, não está claro que tal comparação possa ser feita. Os fumantes se autoselecionam. Podem ser geneticamente ou “constitucionalmente” diferentes dos não fumantes de várias maneiras — mais propensos a correr riscos, mais propensos a beber em excesso. Alguns desses comportamentos podem causar efeitos adversos à saúde que de outra forma poderiam ser atribuídos ao fumo.

Esse era um argumento especialmente conveniente para um cético fazer porque a hipótese constitucional era quase intestável. Somente após o sequenciamento do genoma humano em 2000 se tornou possível procurar genes ligados ao câncer de pulmão. (Ironicamente, Fisher foi provado certo, ainda que de forma muito limitada: tais genes existem.) No entanto, em 1959 Jerome Cornfield, escrevendo com Abe Lilienfeld, publicou uma refutação ponto a ponto dos argumentos de Fisher que, aos olhos de muitos médicos, encerrou a questão.

Cornfield, que trabalhava nos Institutos Nacionais de Saúde, era um participante incomum no debate fumo-câncer. Nem estatístico nem biólogo por formação, formara-se em história e aprendera estatística no Departamento de Agricultura dos EUA. Embora de certo modo autodidata, acabou se tornando um consultor muito procurado e presidente da American Statistical Association. Também havia sido fumante de 2,5 maços por dia, mas abandonou o hábito quando começou a ver os dados sobre câncer de pulmão. (É interessante ver o quão pessoal o debate do fumo era para os cientistas envolvidos. Fisher nunca abandonou o cachimbo, e Yerushalmy nunca abandonou os cigarros.)

Cornfield mirava diretamente na hipótese constitucional de Fisher, e o fazia no próprio terreno de Fisher: a matemática. Suponha, argumentava ele, que haja um fator de confundimento, como um gene do fumo, que explique completamente o risco de câncer dos fumantes. Se os fumantes têm nove vezes o risco de desenvolver câncer de pulmão, o fator de confundimento precisa ser ao menos nove vezes mais comum em fumantes para explicar a diferença de risco. Pense no que isso significa. Se 11\% dos não fumantes tiverem o “gene do fumo”, então 99\% dos fumantes teriam de tê-lo. E se mesmo 12\% dos não fumantes acontecerem de ter o gene do câncer, então se torna matematicamente impossível que o gene do câncer explique plenamente a associação entre fumo e câncer.

Para os biólogos, esse argumento, chamado desigualdade de Cornfield, reduziu a hipótese constitucional de Fisher a ruínas fumegantes. É inconcebível que uma variação genética pudesse estar tão estreitamente ligada a algo tão complexo e imprevisível quanto a escolha de uma pessoa de fumar. A desigualdade de Cornfield era de fato um argumento causal em forma embrionária: ela nos dá um critério para adjudicar entre o Diagrama 5.1 (em que a hipótese constitucional não pode explicar plenamente a associação entre fumo e câncer de pulmão) e o Diagrama 5.2 (em que o gene do fumo explicaria plenamente a associação observada).

\textbf{DIAGRAMA 5.1.}

\textbf{DIAGRAMA 5.2.}

Como explicado acima, a associação entre fumo e câncer de pulmão era forte demais para ser explicada pela hipótese constitucional. De fato, o método de Cornfield plantou as sementes de uma técnica muito poderosa chamada “análise de sensibilidade”, que hoje suplementa as conclusões extraídas do motor de inferência descrito na Introdução. Em vez de extrair inferências assumindo a ausência de certas relações causais no modelo, o analista desafia tais premissas e avalia quão fortes as relações alternativas devem ser para explicar os dados observados. O resultado quantitativo é então submetido a um julgamento de plausibilidade, não muito diferente dos julgamentos brutos invocados ao se postular a ausência dessas relações causais. Desnecessário dizer que, se quisermos estender a abordagem de Cornfield a um modelo com mais de três ou quatro variáveis, precisamos de algoritmos e técnicas de estimação que são impensáveis sem o advento de ferramentas gráficas.

Os epidemiologistas nos anos 1950 enfrentavam a crítica de que sua evidência era “apenas estatística”. Alegadamente não havia “prova de laboratório”. Mas mesmo um olhar para a história mostra que esse argumento era falacioso. Se o padrão de “prova de laboratório” tivesse sido aplicado ao escorbuto, então os marinheiros teriam continuado a morrer até os anos 1930, porque até a descoberta da vitamina C não havia “prova de laboratório” de que as frutas cítricas preveniam o escorbuto. Além disso, nos anos 1950 alguns tipos de prova de laboratório dos efeitos do fumo de fato começaram a aparecer em periódicos médicos. Ratos pintados com alcatrão de cigarro desenvolveram câncer. A fumaça de cigarro foi comprovada como contendo benzopirenos, um carcinógeno previamente conhecido. Esses experimentos aumentaram a plausibilidade biológica da hipótese de que o fumo poderia causar câncer.

Ao final da década, a acumulação de tantos tipos diferentes de evidência havia convencido quase todos os especialistas na área de que o fumo de fato causava câncer. Notavelmente, até pesquisadores das companhias de tabaco estavam convencidos — um fato que permaneceu profundamente oculto até os anos 1990, quando litígios e delatores forçaram as companhias de tabaco a liberar muitos milhares de documentos previamente secretos. Em 1953, por exemplo, um químico da R.J. Reynolds, Claude Teague, havia escrito à alta gerência da companhia que o tabaco era “um importante fator etiológico na indução do câncer primário do pulmão”, quase uma repetição palavra por palavra da conclusão de Hill e Doll.

Em público, as companhias de cigarro cantavam uma melodia diferente. Em janeiro de 1954, as principais companhias de tabaco (incluindo a Reynolds) publicaram um anúncio em jornais de todo o país, “A Frank Statement to Cigarette Smokers” [“Uma declaração franca aos fumantes de cigarro”], que dizia: “Acreditamos que os produtos que fabricamos não são prejudiciais à saúde. Sempre cooperamos e sempre cooperaremos estreitamente com aqueles cuja tarefa é salvaguardar a saúde pública.” Num discurso proferido em março de 1954, George Weissman, vice-presidente da Philip Morris and Company, disse: “Se tivéssemos qualquer pensamento ou conhecimento de que de alguma forma estivéssemos vendendo um produto prejudicial aos consumidores, interromperíamos os negócios amanhã.” Sessenta anos depois, ainda estamos esperando a Philip Morris cumprir essa promessa.

Isso nos leva ao episódio mais triste de toda a controvérsia fumo-câncer: os esforços deliberados das companhias de tabaco para enganar o público sobre os riscos à saúde. Se a Natureza é como um gênio que responde a uma pergunta de forma verdadeira, mas apenas exatamente como ela é formulada, imagine o quanto é mais difícil para os cientistas enfrentar um adversário que pretende nos enganar. As guerras do cigarro foram o primeiro confronto da ciência com o negacionismo organizado, e ninguém estava preparado. As companhias de tabaco amplificavam qualquer farrapo de controvérsia científica que pudessem. Montaram seu próprio Comitê de Pesquisa da Indústria do Tabaco, uma organização de fachada que dava dinheiro a cientistas para estudar questões relacionadas ao câncer ou ao tabaco — mas de algum modo nunca chegava à questão central. Quando conseguiam encontrar céticos legítimos da conexão fumo-câncer — como R. A. Fisher e Jacob Yerushalmy —, as companhias de tabaco lhes pagavam honorários de consultoria.

O caso de Fisher é particularmente triste. É claro que o ceticismo tem seu lugar. Os estatísticos são pagos para serem céticos; são a consciência da ciência. Mas há uma diferença entre ceticismo razoável e não razoável. Fisher cruzou essa linha e foi além. Sempre incapaz de admitir seus próprios erros, e certamente influenciado por seu hábito vitalício de fumar cachimbo, ele não conseguia reconhecer que a maré de evidências se voltara contra ele. Seus argumentos se tornaram desesperados. Agarrava-se a um resultado contra-intuitivo no primeiro artigo de Doll e Hill — a descoberta (que mal alcançou o nível de significância estatística) de que pacientes com câncer de pulmão se descreviam como inaladores menos frequentemente do que os controles — e não a soltava. Nenhum dos estudos subsequentes encontrou qualquer efeito desse tipo. Embora Fisher soubesse tão bem quanto qualquer um que resultados “estatisticamente significativos” às vezes falham em ser replicados, ele recorria à zombaria. Argumentava que o estudo deles havia mostrado que inalar a fumaça de cigarro poderia ser benéfico e pedia mais pesquisas sobre esse “ponto extremamente importante”. Talvez a única coisa positiva que possamos dizer sobre o papel de Fisher no debate seja que é muito improvável que o dinheiro do tabaco o tenha corrompido de alguma forma. Sua própria obstinação bastava.

Por todas essas razões, a ligação entre fumo e câncer permaneceu controversa na mente do público muito depois de ter terminado entre os epidemiologistas. Mesmo os médicos, que deveriam estar mais atentos à ciência, permaneciam não convencidos: uma pesquisa conduzida pela American Cancer Society em 1960 mostrou que apenas um terço dos médicos americanos concordava com a afirmação de que o fumo era “uma causa maior de câncer de pulmão”, e 43\% dos médicos eram eles mesmos fumantes.

Embora possamos justamente culpar Fisher por sua obstinação e as companhias de tabaco por sua decepção deliberada, também devemos reconhecer que a comunidade científica estava laborando numa camisa de força ideológica. Fisher estivera certo em promover os ensaios clínicos randomizados como uma maneira altamente eficaz de avaliar um efeito causal. No entanto, ele e seus seguidores falharam em perceber que há muito que podemos aprender de estudos observacionais. Esse é o benefício de um modelo causal: ele alavanca o conhecimento científico do experimentador. Os métodos de Fisher assumem que o experimentador começa sem conhecimento prévio ou opiniões sobre a hipótese a ser testada. Eles impõem a ignorância ao cientista, uma situação da qual os negacionistas se aproveitaram avidamente.

Como os cientistas não tinham uma definição direta da palavra “causa” e nenhuma maneira de apurar um efeito causal sem um ensaio clínico randomizado, estavam mal preparados para um debate sobre se o fumo causava câncer. Foram forçados a tatear em direção a uma definição num processo que durou ao longo dos anos 1950 e alcançou uma conclusão dramática em 1964.

\section{A COMISSÃO DO CIRURGIÃO-GERAL E OS CRITÉRIOS DE HILL}
O artigo de Cornfield e Lilienfeld havia pavimentado o caminho para uma declaração definitiva das autoridades de saúde sobre os efeitos do fumo. O Royal College of Physicians no Reino Unido tomou a dianteira, emitindo um relatório em 1962 concluindo que o fumo de cigarro era um agente causador do câncer de pulmão. Pouco depois, o Cirurgião-Geral dos EUA Luther Terry (muito possivelmente a pedido do presidente John F. Kennedy) anunciou sua intenção de nomear um comitê consultivo especial para estudar a questão (ver Figura 5.3).

\begin{figure}[ht]
\centering
\includegraphics[width=0.92\textwidth]{../figuras/fig_5_3.png}
\caption{Figura 5.3. In 1963, a surgeon general’s advisory committee wrestled com the}
\end{figure}
\textit{Figura 5.3. In 1963, a surgeon general’s advisory committee wrestled com the}

O comitê foi cuidadosamente equilibrado para incluir cinco fumantes e cinco não fumantes, duas pessoas sugeridas pela indústria do tabaco, e ninguém que tivesse feito declarações públicas a favor ou contra o fumo. Por essa razão, pessoas como Lilienfeld e Cornfield eram inelegíveis. Os membros do comitê eram especialistas distintos em medicina, química ou biologia. Um deles, William Cochran da Universidade Harvard, era estatístico. De fato, as credenciais de Cochran em estatística eram as melhores possíveis: ele era aluno de um aluno de Karl Pearson.

\textbf{FIGURA 5.3.} Em 1963, um comitê consultivo do cirurgião-geral lutava com o problema de como avaliar os efeitos causais do fumo. Retratados aqui estão William Cochran (o estatístico do comitê), o Cirurgião-Geral Luther Terry e o químico Louis Fieser. (Fonte: Desenho de Dakota Harr.)

O comitê labutou por mais de um ano em seu relatório, e uma questão maior era o uso da palavra “causa”. Os membros do comitê tinham de pôr de lado concepções determinísticas oitocentistas de causalidade, e também tinham de pôr de lado a estatística. Como escreveram (provavelmente Cochran) no relatório: “Métodos estatísticos não podem estabelecer prova de uma relação causal numa associação. A significância causal de uma associação é uma questão de julgamento que vai além de qualquer declaração de probabilidade estatística. Para julgar ou avaliar a significância causal da associação entre o atributo ou agente e a doença, ou o efeito sobre a saúde, um número de critérios deve ser utilizado, nenhum dos quais é uma base por si só suficiente para o julgamento.”

O comitê listou cinco desses critérios: consistência (muitos estudos, em diferentes populações, mostram resultados semelhantes); força da associação (incluindo o efeito dose-resposta: mais fumo está associado a um risco maior); especificidade da associação (um agente particular deveria ter um efeito particular e não uma longa litania de efeitos); relação temporal (o efeito deveria seguir a causa); e coerência (plausibilidade biológica e consistência com outros tipos de evidência, como experimentos de laboratório e séries temporais).

Em 1965, Austin Bradford Hill, que não estava no comitê, tentou resumir os argumentos de uma maneira que pudesse ser aplicada a outros problemas de saúde pública e acrescentou quatro critérios mais à lista; como resultado, a lista inteira de nove critérios ficou conhecida como “critérios de Hill”. Na verdade, Hill os chamava de “pontos de vista”, não de requisitos, e enfatizava que qualquer um deles poderia faltar em qualquer caso particular. “Nenhum dos meus nove pontos de vista pode trazer evidência indiscutível a favor ou contra a hipótese de causa e efeito, e nenhum pode ser exigido como \textit{sine qua non}”, escreveu ele.

De fato, é bastante fácil encontrar argumentos contra cada um dos critérios, seja na lista de Hill ou na lista mais curta do comitê consultivo. A consistência por si só não prova nada; se trinta estudos ignoram cada um o mesmo confundidor, todos podem facilmente ser enviesados. A força da associação é vulnerável pela mesma razão; como apontado anteriormente, o tamanho dos sapatos das crianças está fortemente associado, mas não causalmente relacionado, à aptidão delas para a leitura. A especificidade sempre foi um critério particularmente controverso. Faz sentido no contexto de doenças infecciosas, em que um agente tipicamente produz uma doença, mas menos no contexto de exposição ambiental. O fumo leva a um risco aumentado de uma variedade de outras doenças, como enfisema e doença cardiovascular. Isso realmente enfraquece a evidência de que causa câncer? A relação temporal tem algumas exceções, como mencionado antes — por exemplo, o canto de um galo não causa o nascer do sol, embora sempre o preceda. Por fim, a coerência com a teoria ou os fatos estabelecidos é certamente desejável, mas a história da ciência está repleta de teorias derrubadas e achados de laboratório equivocados.

Os “pontos de vista” de Hill ainda são úteis como descrição de como uma disciplina chega a aceitar uma hipótese causal, usando uma variedade de evidências, mas vieram sem nenhuma metodologia para implementá-los. Por exemplo, a plausibilidade biológica e a consistência com experimentos são supostamente coisas boas. Mas como, precisamente, devemos pesar esses tipos de evidência? Como trazemos o conhecimento prévio para o quadro? Aparentemente, cada cientista só tem de decidir por si mesmo. Mas decisões de estômago podem estar erradas, especialmente se houver pressões políticas ou considerações monetárias, ou se o cientista for viciado na substância sendo estudada.

Nenhum desses comentários tem a intenção de denegrir o trabalho do comitê. Seus membros fizeram o melhor que puderam num ambiente que não lhes fornecia nenhum mecanismo para discutir a causalidade. Seu reconhecimento de que critérios não estatísticos eram necessários foi um grande passo adiante. E as difíceis decisões pessoais que os fumantes do comitê tomaram atestam a seriedade de suas conclusões. Luther Terry, que havia sido fumante de cigarro, mudou para o cachimbo. Leonard Schuman anunciou que estava parando. William Cochran reconheceu que poderia reduzir seu risco de câncer parando, mas sentia que o “conforto de meus cigarros” era compensação suficiente pelo risco. Mais dolorosamente, Louis Fieser, fumante de quatro maços por dia, foi diagnosticado com câncer de pulmão menos de um ano após o relatório. Escreveu ao comitê: “Vocês devem se lembrar de que, embora plenamente convencido pela evidência, continuei fumando pesadamente ao longo das deliberações de nosso comitê e invoquei todas as desculpas usuais… Meu caso me parece mais convincente do que quaisquer estatísticas.” Menos um pulmão, ele finalmente parou de fumar.

Visto da perspectiva da saúde pública, o relatório do comitê consultivo foi um marco. Em dois anos, o Congresso havia exigido que os fabricantes colocassem advertências de saúde em todos os maços de cigarro. Em 1971, os anúncios de cigarro foram banidos do rádio e da televisão. A porcentagem de adultos dos EUA que fumam declinou de seu máximo histórico de 45\% em 1965 para 19,3\% em 2010. A campanha antitabaco tem sido uma das maiores e mais bem-sucedidas, embora dolorosamente lenta e incompleta, intervenções de saúde pública da história. O trabalho do comitê também forneceu um valioso modelo para alcançar o consenso científico e serviu de modelo para futuros relatórios do cirurgião-geral sobre o fumo e muitos outros tópicos nos anos seguintes (incluindo a fumaça de segunda mão, que se tornou uma questão maior nos anos 1980).

Visto da perspectiva da causalidade, o relatório foi, na melhor das hipóteses, um sucesso modesto. Estabeleceu claramente a gravidade das perguntas causais e que os dados sozinhos não podiam respondê-las. Mas como um mapa para a descoberta futura, suas diretrizes eram incertas e frágeis. Os critérios de Hill são melhor lidos como um documento histórico, resumindo os tipos de evidência que haviam emergido nos anos 1950 e, em última análise, convencido a comunidade médica. Mas como guia para a pesquisa futura, são inadequados. Para todas as perguntas causais, exceto as mais amplas, precisamos de um instrumento mais preciso. Em retrospecto, a desigualdade de Cornfield, que plantou as sementes da análise de sensibilidade, foi um passo nessa direção.

\section{FUMAR PARA RECÉM-NASCIDOS}
Mesmo depois de o debate sobre fumo e câncer ter se acalmado, um grande paradoxo persistiu. Em meados dos anos 1960, Jacob Yerushalmy apontou que o fumo de uma mãe durante a gravidez parecia beneficiar a saúde de seu bebê recém-nascido, se o bebê por acaso nascesse com baixo peso. Esse enigma, chamado paradoxo do peso ao nascer, desafiava o consenso médico emergente sobre o fumo e não foi satisfatoriamente explicado até 2006 — mais de quarenta anos após o artigo original de Yerushalmy. Estou absolutamente convencido de que demorou tanto porque a linguagem da causalidade não estava disponível de 1960 a 1990.

Em 1959, Yerushalmy havia lançado um estudo de saúde pública de longo prazo que coletava dados pré e pós-natais de mais de 15.000 crianças na área da Baía de São Francisco. Os dados incluíam informações sobre os hábitos de fumar das mães, bem como os pesos ao nascer e as taxas de mortalidade de seus bebês no primeiro mês de vida. Vários estudos já haviam mostrado que os bebês de mães fumantes pesavam menos, em média, ao nascer do que os bebês de não fumantes, e era natural supor que isso se traduziria em pior sobrevivência. De fato, um estudo nacional de bebês de baixo peso ao nascer (definidos como aqueles que pesam menos de 2,5 kg ao nascer) havia mostrado que sua taxa de mortalidade era mais de vinte vezes maior do que a de bebês de peso normal ao nascer. Assim, os epidemiologistas postularam uma cadeia de causas e efeitos: Fumo → Baixo Peso ao Nascer → Mortalidade.

O que Yerushalmy encontrou nos dados foi inesperado até para ele. Era verdade que os bebês de fumantes pesavam menos, em média, do que os bebês de não fumantes (cerca de 200 gramas a menos). No entanto, os bebês de baixo peso ao nascer de mães fumantes tinham uma taxa de sobrevivência melhor do que os de não fumantes. Era como se o fumo da mãe de fato tivesse um efeito protetor.

Se Fisher tivesse descoberto algo assim, provavelmente o teria proclamado em voz alta como um dos benefícios do fumo. Yerushalmy, em seu crédito, não o fez. Escreveu, de forma muito mais cautelosa: “Esses achados paradoxais levantam dúvidas e argumentam contra a proposição de que o fumo de cigarro age como um fator exógeno que interfere no desenvolvimento intrauterino do feto.” Em suma, não há caminho causal de Fumo para Mortalidade.

Os epidemiologistas modernos acreditam que Yerushalmy estava errado. A maioria acredita que o fumo de fato aumenta a mortalidade neonatal — por exemplo, porque interfere na transferência de oxigênio através da placenta. Mas como podemos reconciliar essa hipótese com os dados?

Estatísticos e epidemiologistas insistiam em analisar o paradoxo em termos probabilísticos e vê-lo como uma anomalia peculiar ao peso ao nascer. Como se vê, ele tem pouco a ver com o peso ao nascer e tudo a ver com colisores. Quando visto sob essa luz, não é paradoxal, mas instrutivo. De fato, os dados de Yerushalmy eram completamente consistentes com o modelo Fumo → Baixo Peso ao Nascer → Mortalidade, uma vez que acrescentemos um pouco mais a ele. O fumo pode ser prejudicial no sentido de que contribui para o baixo peso ao nascer, mas certas outras causas de baixo peso ao nascer, como anormalidades genéticas graves ou que ameaçam a vida, são muito mais prejudiciais. Há duas possíveis explicações para o baixo peso ao nascer em um bebê particular: ele pode ter uma mãe fumante, ou pode ser afetado por uma daquelas outras causas. Se descobrirmos que a mãe é fumante, isso “explica-afasta” o baixo peso e, consequentemente, reduz a verossimilhança de um defeito congênito grave. Mas se a mãe não fuma, temos evidência mais forte de que a causa do baixo peso ao nascer é um defeito congênito, e o prognóstico do bebê se torna pior.

Como antes, um diagrama causal torna tudo mais claro. Quando incorporamos as novas premissas, o diagrama causal se parece com a Figura 5.4. Podemos ver que o paradoxo do peso ao nascer é um exemplo perfeito de viés de colisor. O colisor é Peso ao Nascer. Ao olhar apenas para bebês com baixo peso ao nascer, estamos condicionando nesse colisor. Isso abre um caminho pela porta dos fundos entre Fumo e Mortalidade que vai Fumo → Peso ao Nascer ← Defeito Congênito → Mortalidade. Esse caminho é não causal porque uma das setas vai na direção errada. Não obstante, induz uma correlação espúria entre Fumo e Mortalidade e enviesa nossa estimativa do efeito causal (direto) real, Fumo → Mortalidade. De fato, enviesa a estimativa a tal ponto que o fumo realmente parece benéfico.

\begin{figure}[ht]
\centering
\includegraphics[width=0.92\textwidth]{../figuras/fig_5_4.png}
\caption{Figura 5.4. Podemos ver que o paradoxo do peso ao nascer é um exemplo perfeito de viés de colisor. O colisor é Peso ao Nascer. Ao olhar apenas para bebês com baixo peso ao nascer, estamos condicionando nesse colisor. Isso abre um caminho pela porta dos fundos entre Fumo e Mortalidade que vai Fumo → Peso ao Nascer ← Defeito Congênito → Mortalidade. Esse caminho é não causal porque uma das setas vai na direção errada. Não obstante, induz uma correlação espúria entre Fumo e Mortalidade e enviesa nossa estimativa do efeito causal (direto) real, Fumo → Mortalidade. De fato, enviesa a estimativa a tal ponto que o fumo realmente parece benéfico.}
\end{figure}
\textit{Figura 5.4. Podemos ver que o paradoxo do peso ao nascer é um exemplo perfeito de viés de colisor. O colisor é Peso ao Nascer. Ao olhar apenas para bebês com baixo peso ao nascer, estamos condicionando nesse colisor. Isso abre um caminho pela porta dos fundos entre Fumo e Mortalidade que vai Fumo → Peso ao Nascer ← Defeito Congênito → Mortalidade. Esse caminho é não causal porque uma das setas vai na direção errada. Não obstante, induz uma correlação espúria entre Fumo e Mortalidade e enviesa nossa estimativa do efeito causal (direto) real, Fumo → Mortalidade. De fato, enviesa a estimativa a tal ponto que o fumo realmente parece benéfico.}

A beleza dos diagramas causais é que eles tornam a fonte do viés óbvia. Sem tais diagramas, os epidemiologistas discutiram o paradoxo por quarenta anos. De fato, ainda estão discutindo: a edição de outubro de 2014 do \textit{International Journal of Epidemiology} contém vários artigos sobre esse tópico. Um deles, de Tyler VanderWeele de Harvard, crava a explicação perfeitamente e contém um diagrama exatamente como o abaixo.

\textbf{FIGURA 5.4.} Diagrama causal para o paradoxo do peso ao nascer.

É claro, esse diagrama é provavelmente simples demais para capturar a história completa por trás do fumo, do peso ao nascer e da mortalidade infantil. No entanto, o princípio do viés de colisor é robusto. Neste caso, o viés foi detectado porque o fenômeno aparente era implausível demais, mas imagine quantos casos de viés de colisor passam despercebidos porque o viés não conflita com a teoria.

\section{DEBATES APAIXONADOS: CIÊNCIA VERSUS CULTURA}
Depois de começar a trabalhar neste capítulo, tive ocasião de contatar Allen Wilcox, o epidemiologista provavelmente mais identificado com esse paradoxo. Ele fez uma pergunta muito inconveniente sobre o diagrama da Figura 5.4: Como sabemos que o baixo peso ao nascer é de fato uma causa direta de mortalidade? De fato, ele acredita que os médicos vêm interpretando mal o baixo peso ao nascer o tempo todo. Porque está fortemente associado à mortalidade infantil, os médicos o interpretaram como uma causa. De fato, essa associação poderia ser inteiramente devida a confudidores (representados por “Defeito Congênito” na Figura 5.4, embora Wilcox não seja tão específico).

Dois pontos valem a pena sobre o argumento de Wilcox. Primeiro, mesmo se apagarmos a seta Peso ao Nascer → Mortalidade, o colisor permanece. Assim, o diagrama causal continua a explicar com sucesso o paradoxo do peso ao nascer. Segundo, a variável causal que Wilcox mais estudou não é o fumo, mas a raça. E a raça ainda incita debates apaixonados em nossa sociedade.

De fato, o mesmo paradoxo do peso ao nascer é observado em filhos de mães negras como em filhos de fumantes. Mulheres negras dão à luz bebês de baixo peso mais frequentemente do que mulheres brancas, e seus bebês têm uma taxa de mortalidade mais alta. No entanto, seus bebês de baixo peso ao nascer têm uma taxa de sobrevivência melhor do que os bebês de baixo peso ao nascer de mulheres brancas. Agora, que conclusões deveríamos tirar? Podemos dizer a uma fumante grávida que ela ajudaria seu bebê parando de fumar. Mas não podemos dizer a uma mulher negra grávida para parar de ser negra.

Em vez disso, deveríamos abordar as questões sociais que fazem os filhos de mães negras terem uma taxa de mortalidade mais alta. Isso certamente não é uma afirmação controversa. Mas que causas deveríamos abordar, e como deveríamos medir nosso progresso? Para o bem ou para o mal, muitos defensores da justiça racial assumiram o peso ao nascer como um passo intermediário na cadeia Raça → Peso ao Nascer → Mortalidade. Não só isso: tomaram o peso ao nascer como um \textit{proxy} da mortalidade infantil, assumindo que melhorias num levariam automaticamente a melhorias no outro. É fácil compreender por que fizeram isso. Medições de pesos médios ao nascer são mais fáceis de obter do que medições de mortalidade infantil.

Agora imagine o que acontece quando alguém como Wilcox aparece e afirma que o baixo peso ao nascer por si só não é uma condição médica e não tem relação causal com a mortalidade infantil. Isso revira o carrinho de maçãs. Wilcox foi acusado de racismo quando primeiro sugeriu essa ideia nos anos 1970, e não ousou publicá-la até 2001. Mesmo então, dois comentários acompanharam seu artigo, e um deles trouxe à tona a questão racial: “No contexto de uma sociedade cujos elementos dominantes justificam suas posições argumentando a inferioridade genética daqueles que dominam, é difícil ser neutro”, escreveu Richard David do Cook County Hospital em Chicago. “Na busca da ‘ciência pura’, um investigador bem-intencionado pode ser percebido como — e pode estar — auxiliando e encorajando uma ordem social que abomina.”

Essa acusação dura, concebida a partir das mais nobres motivações, certamente não é a primeira instância em que um cientista foi repreendido por elucidar verdades que poderiam ter consequências sociais adversas. As objeções do Vaticano às ideias de Galileu certamente surgiram de preocupações genuínas com a ordem social da época. O mesmo se pode dizer da evolução de Charles Darwin e da eugenia de Francis Galton. No entanto, os choques culturais que emanam de novos achados científicos são eventualmente resolvidos por realinhamentos culturais que acomodam esses achados — não pelo ocultamento. Um pré-requisito para esse realinhamento é que separemos a ciência da cultura antes que as opiniões se inflamem. Felizmente, a linguagem dos diagramas causais agora nos dá uma maneira de sermos desapaixonados sobre causas e efeitos não apenas quando é fácil, mas também quando é difícil.

\chapter{Capítulo 6 — Paradoxos à farta!}
\begin{quote}Aquele que confronta o paradoxal expõe-se à realidade.\end{quote}
>

\begin{quote}— FRIEDRICH DÜRRENMATT (1962)\end{quote}
O paradoxo do peso ao nascer, com o qual encerramos o Capítulo 5, é representativo de uma classe surpreendentemente ampla de paradoxos que refletem as tensões entre causalidade e associação. A tensão começa porque elas se situam em dois degraus diferentes da Escada da Causalidade e se agrava pelo fato de que a intuição humana opera sob a lógica da causalidade, enquanto os dados se conformam à lógica das probabilidades e das proporções. Os paradoxos surgem quando aplicamos indevidamente as regras que aprendemos em um domínio ao outro.

Vamos dedicar um capítulo a alguns dos paradoxos mais desconcertantes e bem conhecidos da probabilidade e da estatística porque, em primeiro lugar, são divertidos. Se você ainda não viu os paradoxos de Monty Hall e de Simpson, posso prometer que eles darão um bom exercício ao seu cérebro. E mesmo que você ache que já sabe tudo a respeito deles, creio que apreciará vê-los pela lente da causalidade, que faz tudo parecer um tanto diferente.

No entanto, estudamos paradoxos não apenas porque são divertidos e jogos de espírito. Como as ilusões de ótica, eles também revelam o modo de funcionar do cérebro, os atalhos que ele toma e as coisas que encontra em conflito. Os paradoxos causais projetam um holofote sobre padrões de raciocínio causal intuitivo que colidem com a lógica da probabilidade e da estatística. Na medida em que os estatísticos lutaram com eles — e veremos que se saíram bastante mal —, isso é um sinal de alerta de que algo pode estar errado em se enxergar o mundo sem uma lente causal.

\section{O DESCONCERTANTE PROBLEMA DE MONTY HALL}
No final dos anos 1980, uma escritora chamada Marilyn vos Savant iniciou uma coluna regular na revista \textit{Parade}, suplemento semanal do jornal de domingo em muitas cidades dos Estados Unidos. Sua coluna, “Ask Marilyn”, continua até hoje e apresenta suas respostas a diversos quebra-cabeças, enigmas e perguntas científicas enviadas por leitores. A revista a anunciava como “a mulher mais inteligente do mundo”, o que sem dúvida motivava os leitores a inventar uma pergunta que a deixasse sem resposta.

De todas as perguntas que ela já respondeu, nenhuma criou maior furor do que esta, publicada numa coluna de setembro de 1990: “Suponha que você esteja num programa de auditório e lhe seja dada a escolha entre três portas. Atrás de uma porta há um carro; atrás das outras, cabras. Você escolhe uma porta, digamos a nº 1, e o apresentador, que sabe o que há atrás das portas, abre outra porta, digamos a nº 3, que tem uma cabra. Ele diz a você: ‘Quer escolher a porta nº 2?’ É vantajoso para você trocar a escolha de porta?”

Para os leitores americanos, a pergunta baseava-se de forma óbvia num popular programa de televisão chamado \textit{Let’s Make a Deal}, cujo apresentador, Monty Hall, costumava jogar exatamente esse tipo de jogo mental com os concorrentes.

Em sua resposta, vos Savant argumentou que os concorrentes deveriam trocar de porta. Ao não trocar, teriam apenas uma probabilidade de um em três de ganhar; ao trocar, dobrariam suas chances para dois em três.

Nem mesmo a mulher mais inteligente do mundo poderia ter antecipado o que aconteceu a seguir. Nos meses seguintes, ela recebeu mais de 10.000 cartas de leitores, a maioria discordando dela, e muitas de pessoas que afirmavam ter doutorado em matemática ou estatística. Uma pequena amostra dos comentários de acadêmicos inclui: “Você errou, e errou feio!” (Scott Smith, PhD); “Posso sugerir que você obtenha e consulte um livro-texto padrão de probabilidade antes de tentar responder a uma pergunta desse tipo de novo?” (Charles Reid, PhD); “Você errou!” (Robert Sachs, PhD); e “Você está completamente incorreta” (Ray Bobo, PhD).

Em geral, os críticos argumentavam que não deveria importar se você troca de porta ou não — restam apenas duas portas no jogo, e você escolheu a sua completamente ao acaso, de modo que a probabilidade de que o carro esteja atrás da sua porta deve ser de metade de qualquer maneira.

Quem estava certo? Quem estava errado? E por que o problema incita tanta paixão? As três perguntas merecem um exame mais atento.

Vejamos primeiro como vos Savant resolveu o enigma. Sua solução é de fato assombrosa em sua simplicidade e mais convincente do que qualquer uma que eu tenha visto em muitos livros-texto. Ela fez uma lista (Tabela 6.1) das três possíveis disposições de portas e cabras, junto com os resultados correspondentes sob a estratégia “Trocar” e a estratégia “Ficar”. Os três casos assumem que você escolheu a Porta 1. Como as três possibilidades listadas são (inicialmente) igualmente prováveis, a probabilidade de ganhar se você trocar de porta é de dois terços, e a probabilidade de ganhar se ficar com a Porta 1 é de apenas um terço.

\begin{figure}[ht]
\centering
\includegraphics[width=0.92\textwidth]{../figuras/tab_6_1.png}
\caption{Tabela 6.1. ) das três possíveis disposições de portas e cabras, junto com os resultados correspondentes sob a estratégia “Trocar” e a estratégia “Ficar”. Os três casos assumem que você escolheu a Porta 1. Como as três possibilidades listadas são (inicialmente) igualmente prováveis, a probabilidade de ganhar se você trocar de porta é de dois terços, e a probabilidade de ganhar se ficar com a Porta 1 é de apenas um terço.}
\end{figure}
\textit{Tabela 6.1. ) das três possíveis disposições de portas e cabras, junto com os resultados correspondentes sob a estratégia “Trocar” e a estratégia “Ficar”. Os três casos assumem que você escolheu a Porta 1. Como as três possibilidades listadas são (inicialmente) igualmente prováveis, a probabilidade de ganhar se você trocar de porta é de dois terços, e a probabilidade de ganhar se ficar com a Porta 1 é de apenas um terço.}

Note que a tabela de vos Savant não declara explicitamente qual porta foi aberta pelo apresentador. Essa informação está implicitamente embutida nas colunas 4 e 5. Por exemplo, na segunda linha, mantivemos em mente que o apresentador deve abrir a Porta 3; portanto, trocar o fará cair na Porta 2, uma vitória. Do mesmo modo, na primeira linha, a porta aberta poderia ser a Porta 2 ou a Porta 3, mas a coluna 4 afirma corretamente que você perde em qualquer dos casos se trocar.

Ainda hoje, muitas pessoas que veem o enigma pela primeira vez acham o resultado difícil de acreditar. Por quê? Qual nervo intuitivo é ferido? Há provavelmente 10.000 razões diferentes, uma para cada leitor, mas penso que o argumento mais convincente é este: a solução de vos Savant parece forçar-nos a acreditar em telepatia mental. Se eu devo trocar, não importa qual porta tenha escolhido originalmente, então isso significa que os produtores de algum modo leram minha mente. Como mais poderiam posicionar o carro de modo que seja mais provável que esteja atrás da porta que eu \textit{não} escolhi?

\textbf{TABELA 6.1.} As três possíveis disposições de portas e cabras em \textit{Let’s Make a Deal}, mostrando que trocar de porta é duas vezes mais atraente do que não trocar.

O elemento-chave na resolução desse paradoxo é que precisamos levar em conta não apenas os dados (isto é, o fato de que o apresentador abriu uma porta particular), mas também o processo gerador dos dados — em outras palavras, as regras do jogo. Elas nos dizem algo sobre os dados que \textit{poderiam} ter sido, mas não foram observados. Não é de admirar que os estatísticos, em particular, tenham achado esse enigma difícil de compreender. Eles estão acostumados, como R. A. Fisher (1922) o colocou, à “redução de dados” e a ignorar o processo gerador dos dados.

Para começar, vamos tentar mudar um pouco as regras do jogo e ver como isso afeta nossa conclusão. Imagine um programa de auditório alternativo, chamado \textit{Let’s Fake a Deal}, em que Monty Hall abre uma das duas portas que você não escolheu, mas sua escolha é completamente aleatória. Em particular, ele pode abrir a porta que tem um carro atrás. Azar o seu!

Como antes, assumiremos que você escolheu a Porta 1 no início do jogo e que o apresentador, novamente, abre a Porta 3, revelando uma cabra, e lhe oferece a opção de trocar. Você deveria? Mostraremos que, sob as novas regras, embora o cenário seja idêntico, você não ganha nada ao trocar.

Para isso, fazemos uma tabela como a anterior, levando em conta que há dois eventos aleatórios e independentes — a localização do carro (três possibilidades) e a escolha de Monty Hall de uma porta a abrir (duas possibilidades). Assim, a tabela precisa ter seis linhas, cada uma das quais é igualmente provável porque os eventos são independentes.

Agora, o que acontece se Monty Hall abrir a Porta 3 e revelar uma cabra? Isso nos dá informação significativa: devemos estar na linha 2 ou 4 da tabela. Concentrando-nos apenas nas linhas 2 e 4, podemos ver que a estratégia de trocar já não nos oferece nenhuma vantagem; temos uma probabilidade de um em dois de ganhar de qualquer maneira. Assim, no jogo \textit{Let’s Fake a Deal}, todos os críticos de Marilyn vos Savant estariam certos! No entanto, os dados são os mesmos nos dois jogos. A lição é bastante simples: o modo como obtemos informação não é menos importante do que a informação em si.

\textbf{TABELA 6.2.} Possibilidades de \textit{Let’s Fake a Deal}.

Usemos nosso truque favorito e desenhemos um diagrama causal, que deve ilustrar imediatamente como os dois jogos diferem. Primeiro, a Figura 6.1 mostra um diagrama para o jogo real \textit{Let’s Make a Deal}, em que Monty Hall deve abrir uma porta que não tem um carro atrás. A ausência de uma seta entre Sua Porta e Localização do Carro significa que sua escolha de porta e a escolha dos produtores de onde colocar o carro são independentes. Isso significa que estamos explicitamente excluindo a possibilidade de que os produtores leiam sua mente (ou de que você leia a deles!).

\begin{figure}[ht]
\centering
\includegraphics[width=0.92\textwidth]{../figuras/fig_6_1.png}
\caption{Figura 6.1. mostra um diagrama para o jogo real \textit{Let’s Make a Deal}, em que Monty Hall deve abrir uma porta que não tem um carro atrás. A ausência de uma seta entre Sua Porta e Localização do Carro significa que sua escolha de porta e a escolha dos produtores de onde colocar o carro são independentes. Isso significa que estamos explicitamente excluindo a possibilidade de que os produtores leiam sua mente (ou de que você leia a deles!).}
\end{figure}
\textit{Figura 6.1. mostra um diagrama para o jogo real }Let’s Make a Deal\textit{, em que Monty Hall deve abrir uma porta que não tem um carro atrás. A ausência de uma seta entre Sua Porta e Localização do Carro significa que sua escolha de porta e a escolha dos produtores de onde colocar o carro são independentes. Isso significa que estamos explicitamente excluindo a possibilidade de que os produtores leiam sua mente (ou de que você leia a deles!).}

Ainda mais importantes são as duas setas que estão presentes no diagrama. Elas mostram que Porta Aberta é afetada tanto por sua escolha quanto pela escolha dos produtores. Isso porque Monty Hall deve escolher uma porta que seja diferente tanto de Sua Porta quanto de Localização do Carro; ele tem de levar em conta ambos os fatores.

\textbf{FIGURA 6.1.} Diagrama causal para \textit{Let’s Make a Deal}.

Como se pode ver na Figura 6.1, Porta Aberta é um colisor. Uma vez que obtemos informação sobre essa variável, todas as nossas probabilidades se tornam condicionais a essa informação. Mas quando condicionamos a um colisor, criamos uma dependência espúria entre seus pais. A dependência se confirma nas probabilidades: se você escolheu a Porta 1, a localização do carro tem duas vezes mais probabilidade de estar atrás da Porta 2 do que da Porta 1; se você escolheu a Porta 2, a localização do carro tem duas vezes mais probabilidade de estar atrás da Porta 1.

É uma dependência bizarra, com certeza, de um tipo ao qual a maioria de nós não está acostumada. É uma dependência que não tem causa. Não envolve comunicação física entre os produtores e nós. Não envolve telepatia mental. É puramente um artefato do condicionamento bayesiano: uma transferência mágica de informação sem causalidade. Nossas mentes se rebelam contra essa possibilidade porque, desde a mais tenra infância, aprendemos a associar correlação com causalidade. Se um carro atrás de nós faz todas as mesmas curvas que fazemos, primeiro pensamos que nos está seguindo (causalidade!). Em seguida, pensamos que estamos indo para o mesmo lugar (isto é, há uma causa comum por trás de cada uma de nossas curvas). Mas a correlação sem causa viola o nosso senso comum. Assim, o paradoxo de Monty Hall é exatamente como uma ilusão de ótica ou um truque de mágica: usa nossa própria maquinaria cognitiva para nos enganar.

Por que digo que a abertura da Porta 3 por Monty Hall foi uma “transferência de informação”? Afinal, ela não forneceu nenhuma evidência sobre se sua escolha inicial da Porta 1 estava correta. Você sabia de antemão que ele ia abrir uma porta que escondia uma cabra, e assim o fez. Ninguém deveria pedir que você mudasse suas crenças se testemunhar o inevitável. Então, por que sua crença na Porta 2 subiu de um terço para dois terços?

A resposta é que Monty não podia abrir a Porta 1 depois que você a escolheu — mas podia ter aberto a Porta 2. O fato de que ele não o fez torna mais provável que tenha aberto a Porta 3 porque foi forçado a isso. Assim, há mais evidência do que antes de que o carro está atrás da Porta 2.

Esse é um tema geral da análise bayesiana: qualquer hipótese que tenha sobrevivido a algum teste que ameaça sua validade se torna mais provável. Quanto maior a ameaça, mais provável ela se torna após sobreviver. A Porta 2 era vulnerável à refutação (isto é, Monty podia tê-la aberto), mas a Porta 1 não. Portanto, a Porta 2 se torna uma localização mais provável, enquanto a Porta 1 não. A probabilidade de que o carro esteja atrás da Porta 1 permanece de um em três.

Agora, para comparação, a Figura 6.2 mostra o diagrama causal para \textit{Let’s Fake a Deal}, o jogo em que Monty Hall escolhe uma porta que é diferente da sua, mas de resto escolhida ao acaso. Esse diagrama ainda tem uma seta apontando de Sua Porta para Porta Aberta porque ele tem de se assegurar de que sua porta seja diferente da sua. No entanto, a seta de Localização do Carro para Porta Aberta seria apagada porque ele já não se importa com o local do carro. Nesse diagrama, condicionar a Porta Aberta não tem absolutamente nenhum efeito: Sua Porta e Localização do Carro eram independentes desde o início, e permanecem independentes depois que vemos o conteúdo da porta de Monty. Assim, em \textit{Let’s Fake a Deal}, o carro tem a mesma probabilidade de estar atrás da sua porta ou da outra, como se observa na Tabela 6.2.

\begin{figure}[ht]
\centering
\includegraphics[width=0.92\textwidth]{../figuras/tab_6_2.png}
\caption{Tabela 6.2. Let’s Fake a Deal possibilities.}
\end{figure}
\textit{Tabela 6.2. Let’s Fake a Deal possibilities.}

\textbf{FIGURA 6.2.} Diagrama causal para \textit{Let’s Fake a Deal}.

Do ponto de vista bayesiano, a diferença entre os dois jogos é que em \textit{Let’s Fake a Deal} a Porta 1 é vulnerável à refutação. Monty Hall poderia ter aberto a Porta 3 e revelado o carro, o que teria provado que sua escolha de porta estava errada. Como sua porta e a Porta 2 eram igualmente vulneráveis à refutação, elas ainda têm probabilidade igual.

Embora puramente qualitativa, essa análise poderia ser tornada quantitativa usando a regra de Bayes ou pensando nos diagramas como uma rede bayesiana simples. Fazer isso coloca esse problema num quadro unificador que usamos para pensar sobre outros problemas. Não precisamos inventar um método para resolver o enigma; o esquema de propagação de crenças descrito no Capítulo 3 entregará a resposta correta: isto é, \textit{P}(Porta 2) = 2/3 para \textit{Let’s Make a Deal}, e \textit{P}(Porta 2) = 1/2 para \textit{Let’s Fake a Deal}.

Note que dei de fato duas explicações do paradoxo de Monty Hall. A primeira usa o raciocínio causal para explicar por que observamos uma dependência espúria entre Sua Porta e Localização do Carro; a segunda usa o raciocínio bayesiano para explicar por que a probabilidade da Porta 2 sobe em \textit{Let’s Make a Deal}. Ambas as explicações são valiosas. A bayesiana dá conta do fenômeno, mas não explica realmente por que o percebemos como tão paradoxal. Em minha opinião, uma verdadeira resolução de um paradoxo deveria explicar por que o vemos como um paradoxo em primeiro lugar. Por que as pessoas que leram a coluna dela acreditaram tão fortemente que vos Savant estava errada? Não eram apenas os sabichões. Paul Erdős, um dos matemáticos mais brilhantes dos tempos modernos, igualmente não conseguia acreditar na solução até que uma simulação de computador lhe mostrou que trocar é vantajoso. Que falha profunda em nossa visão intuitiva do mundo isso revela?

“Nossos cérebros simplesmente não estão programados para resolver problemas de probabilidade muito bem, de modo que não me surpreende que houvesse erros”, disse Persi Diaconis, estatístico da Universidade Stanford, numa entrevista de 1991 ao \textit{New York Times}. Verdade, mas há mais nisso. Nossos cérebros não estão programados para resolver problemas de probabilidade, mas estão programados para resolver problemas causais. E essa fiação causal produz erros probabilísticos sistemáticos, como ilusões de ótica. Como não há conexão causal entre Minha Porta e Localização do Carro, nem diretamente nem por meio de uma causa comum, achamos completamente incompreensível que haja uma associação probabilística. Nossos cérebros não estão preparados para aceitar correlações sem causa, e precisamos de treinamento especial — por meio de exemplos como o paradoxo de Monty Hall ou os discutidos no Capítulo 3 — para identificar situações em que elas podem surgir. Uma vez que tenhamos “reprogramado nossos cérebros” para reconhecer colisores, o paradoxo deixa de ser confuso.

\section{MAIS VIÉS DE COLISOR: O PARADOXO DE BERKSON}
Em 1946, Joseph Berkson, bioestatístico da Clínica Mayo, apontou uma peculiaridade dos estudos observacionais conduzidos em ambiente hospitalar: mesmo que duas doenças não tenham relação entre si na população geral, elas podem parecer associadas entre os pacientes de um hospital. Para compreender a observação de Berkson, comecemos com um diagrama causal (Figura 6.3). Também é útil pensar numa possibilidade muito extrema: nem a Doença 1 nem a Doença 2 são ordinariamente graves o bastante para causar hospitalização, mas a combinação o é. Nesse caso, esperaríamos que a Doença 1 estivesse altamente correlacionada com a Doença 2 na população hospitalizada.

\begin{figure}[ht]
\centering
\includegraphics[width=0.92\textwidth]{../figuras/fig_6_3.png}
\caption{Figura 6.3. ). Também é útil pensar numa possibilidade muito extrema: nem a Doença 1 nem a Doença 2 são ordinariamente graves o bastante para causar hospitalização, mas a combinação o é. Nesse caso, esperaríamos que a Doença 1 estivesse altamente correlacionada com a Doença 2 na população hospitalizada.}
\end{figure}
\textit{Figura 6.3. ). Também é útil pensar numa possibilidade muito extrema: nem a Doença 1 nem a Doença 2 são ordinariamente graves o bastante para causar hospitalização, mas a combinação o é. Nesse caso, esperaríamos que a Doença 1 estivesse altamente correlacionada com a Doença 2 na população hospitalizada.}

\textbf{FIGURA 6.3.} Diagrama causal para o paradoxo de Berkson.

Ao realizar um estudo em pacientes hospitalizados, estamos controlando por Hospitalização. Como sabemos, condicionar a um colisor cria uma associação espúria entre Doença 1 e Doença 2. Em muitos de nossos exemplos anteriores a associação era negativa por causa do efeito “explica-afasta”, mas aqui ela é positiva porque ambas as doenças precisam estar presentes para a hospitalização (não apenas uma).

No entanto, por muito tempo os epidemiologistas se recusaram a acreditar nessa possibilidade. Ainda não acreditavam em 1979, quando David Sackett da Universidade McMaster, especialista em todo tipo de viés estatístico, forneceu forte evidência de que o paradoxo de Berkson é real. Num exemplo (ver Tabela 6.3), ele estudou dois grupos de doenças: respiratórias e ósseas. Cerca de 7,5\% das pessoas na população geral têm uma doença óssea, e essa porcentagem é independente de terem ou não doença respiratória. Mas para pessoas hospitalizadas com doença respiratória, a frequência de doença óssea salta para 25\%! Sackett chamou esse fenômeno de “viés de taxa de admissão” ou “viés de Berkson”.

\begin{figure}[ht]
\centering
\includegraphics[width=0.92\textwidth]{../figuras/tab_6_3.png}
\caption{Tabela 6.3. ), ele estudou dois grupos de doenças: respiratórias e ósseas. Cerca de 7,5\% das pessoas na população geral têm uma doença óssea, e essa porcentagem é independente de terem ou não doença respiratória. Mas para pessoas hospitalizadas com doença respiratória, a frequência de doença óssea salta para 25\%! Sackett chamou esse fenômeno de “viés de taxa de admissão” ou “viés de Berkson”.}
\end{figure}
\textit{Tabela 6.3. ), ele estudou dois grupos de doenças: respiratórias e ósseas. Cerca de 7,5\% das pessoas na população geral têm uma doença óssea, e essa porcentagem é independente de terem ou não doença respiratória. Mas para pessoas hospitalizadas com doença respiratória, a frequência de doença óssea salta para 25\%! Sackett chamou esse fenômeno de “viés de taxa de admissão” ou “viés de Berkson”.}

\textbf{TABELA 6.3.} Dados de Sackett ilustrando o paradoxo de Berkson.

Sackett admite que não podemos atribuir definitivamente esse efeito ao viés de Berkson porque também poderia haver fatores de confundimento. O debate está, em certa medida, em andamento. No entanto, ao contrário de 1946 e 1979, os pesquisadores em epidemiologia agora compreendem os diagramas causais e os vieses que eles acarretam. A discussão passou agora a pontos mais finos sobre quão grande o viés pode ser e se ele é grande o bastante para se observar em diagramas causais com mais variáveis. Isso é progresso!

As correlações induzidas por colisores não são novas. Foram encontradas em trabalhos que remontam a um estudo de 1911 do economista inglês Arthur Cecil Pigou, que comparou filhos de pais alcoólatras e não alcoólatras. Também são encontradas, embora não sob esse nome, no trabalho de Barbara Burks (1926), Herbert Simon (1954) e, é claro, Berkson. Também não são tão esotéricas quanto podem parecer a partir dos meus exemplos. Experimente isto: lance duas moedas simultaneamente cem vezes e anote os resultados apenas quando ao menos uma delas der cara. Olhando para a sua tabela, que provavelmente conterá cerca de setenta e cinco entradas, você verá que os resultados dos dois lançamentos simultâneos de moeda não são independentes. Toda vez que a Moeda 1 deu coroa, a Moeda 2 deu cara. Como isso é possível? As moedas de algum modo se comunicaram uma com a outra à velocidade da luz? Claro que não. Na realidade, você condicionou a um colisor ao censurar todos os resultados coroa-coroa.

Em \textit{The Direction of Time}, publicado postumamente em 1956, o filósofo Hans Reichenbach fez uma ousada conjectura chamada “princípio da causa comum”. Refutando o ditado “Correlação não implica causalidade”, Reichenbach postulou uma ideia muito mais forte: “Nenhuma correlação sem causalidade.” Ele queria dizer que uma correlação entre duas variáveis, \textit{X} e \textit{Y}, não pode surgir por acidente. Ou uma das variáveis causa a outra, ou uma terceira variável, digamos \textit{Z}, precede e causa ambas.

Nosso simples experimento com moedas prova que o ditame de Reichenbach era forte demais, porque deixa de levar em conta o processo pelo qual as observações são selecionadas. Não havia causa comum do resultado das duas moedas, e nenhuma moeda comunicou seu resultado à outra. Não obstante, os resultados em nossa lista estavam correlacionados. O erro de Reichenbach foi sua falha em considerar estruturas de colisor — a estrutura por trás da seleção dos dados. O equívoco foi particularmente iluminador porque localiza a falha exata na fiação de nossos cérebros. Vivemos nossas vidas como se o princípio da causa comum fosse verdadeiro. Sempre que vemos padrões, buscamos uma explicação causal. De fato, temos fome de uma explicação, em termos de mecanismos estáveis que se situam fora dos dados. O tipo de explicação mais satisfatório é a causalidade direta: \textit{X} causa \textit{Y}. Quando isso falha, encontrar uma causa comum de \textit{X} e \textit{Y} geralmente nos satisfará.

Em comparação, os colisores são demasiado etéreos para satisfazer nossos apetites causais. Ainda queremos saber o mecanismo por meio do qual as duas moedas coordenam seu comportamento. A resposta é uma decepção esmagadora. Elas não se comunicam de modo algum. A correlação que observamos é, no sentido mais puro e literal, uma ilusão. Ou talvez até uma delusão: isto é, uma ilusão que trouxemos sobre nós mesmos ao escolher quais eventos incluir em nosso conjunto de dados e quais ignorar.

É importante perceber que nem sempre estamos conscientes de fazer essa escolha, e essa é uma das razões pelas quais o viés de colisor pode tão facilmente apanhar o desavisado. No experimento das duas moedas, a escolha foi consciente: eu lhe disse para não registrar as tentativas com duas coroas. Mas em muitas ocasiões não estamos cientes de fazer a escolha, ou a escolha é feita por nós. No paradoxo de Monty Hall, o apresentador abre a porta por nós. No paradoxo de Berkson, um pesquisador desavisado pode escolher estudar pacientes hospitalizados por razões de conveniência, sem perceber que está enviesando seu estudo.

O prisma distorcedor dos colisores é igualmente prevalente na vida cotidiana. Como Jordan Ellenberg pergunta em \textit{How Not to Be Wrong}, você já notou que, entre as pessoas com quem você sai, as atraentes tendem a ser babacas? Em vez de construir elaboradas teorias psicossociais, considere uma explicação mais simples. Sua escolha de pessoas com quem sair depende de dois fatores: atratividade e personalidade. Você arriscará sair com uma pessoa atraente e má ou com uma pessoa simpática e pouco atraente, e certamente com uma pessoa atraente e simpática, mas não com uma pessoa má e pouco atraente. É o mesmo que no exemplo das duas moedas, quando você censurou os resultados coroa-coroa. Isso cria uma correlação negativa espúria entre atratividade e personalidade. A triste verdade é que as pessoas pouco atraentes são tão más quanto as atraentes — mas você nunca perceberá isso, porque nunca sairá com alguém que seja ao mesmo tempo mau e pouco atraente.

\section{O PARADOXO DE SIMPSON}
Agora que mostramos que os produtores de TV não têm realmente habilidades telepáticas e que as moedas não podem se comunicar umas com as outras, que outros mitos podemos desfazer? Comecemos com o mito do medicamento ruim/ruim/bom, ou RRG. Imagine um médico — o Dr. Simpson, vamos chamá-lo — lendo em seu consultório sobre um novo medicamento promissor (Medicamento D) que parece reduzir o risco de infarto. Entusiasmado, ele consulta os dados dos pesquisadores on-line. Seu entusiasmo esfria um pouco quando olha os dados de pacientes homens e nota que o risco de infarto deles é de fato \textit{maior} se tomarem o Medicamento D. “Bom”, diz ele, “o Medicamento D deve ser muito eficaz para as mulheres.” Mas então vira a página seguinte, e sua decepção se transforma em perplexidade. “O que é isto?”, exclama o Dr. Simpson. “Diz aqui que as mulheres que tomaram o Medicamento D também estavam em maior risco de infarto. Devo estar perdendo o juízo! Esse medicamento parece ser ruim para as mulheres, ruim para os homens, mas bom para as pessoas.”

Você também está perplexo? Se sim, está em boa companhia. Esse paradoxo, descoberto pela primeira vez por um estatístico da vida real chamado Edward Simpson em 1951, tem inquietado estatísticos por mais de sessenta anos — e permanece vexatório até o dia de hoje. Mesmo em 2016, enquanto eu escrevia este livro, quatro novos artigos (incluindo uma dissertação de doutorado) saíram, tentando explicar o paradoxo de Simpson a partir de quatro pontos de vista diferentes.

Em 1983, Melvin Novick escreveu: “A resposta aparente é que quando sabemos que o sexo do paciente é masculino ou quando sabemos que é feminino não usamos o tratamento, mas se o sexo é desconhecido deveríamos usar o tratamento! Obviamente essa conclusão é ridícula.” Concordo plenamente. É ridículo que um medicamento seja ruim para os homens e ruim para as mulheres, mas bom para as pessoas. Assim, uma dessas três afirmações deve estar errada. Mas qual? E por quê? E como essa confusão é sequer possível?

Para responder a essas perguntas, é claro, precisamos olhar para os dados (ficcionais) que tanto perplexaram o nosso bom Dr. Simpson. O estudo era observacional, não randomizado, com sessenta homens e sessenta mulheres. Isso significa que os próprios pacientes decidiram se tomariam ou não o medicamento. A Tabela 6.4 mostra quantos de cada sexo receberam o Medicamento D e quantos foram subsequentemente diagnosticados com infarto.

\begin{figure}[ht]
\centering
\includegraphics[width=0.92\textwidth]{../figuras/tab_6_4.png}
\caption{Tabela 6.4. mostra quantos de cada sexo receberam o Medicamento D e quantos foram subsequentemente diagnosticados com infarto.}
\end{figure}
\textit{Tabela 6.4. mostra quantos de cada sexo receberam o Medicamento D e quantos foram subsequentemente diagnosticados com infarto.}

Deixe-me enfatizar onde está o paradoxo. Como se pode ver, 5\% (um em vinte) das mulheres no grupo de controle tiveram depois um infarto, em comparação com 7,5\% das mulheres que tomaram o medicamento. Assim, o medicamento está associado a um risco maior de infarto para as mulheres. Entre os homens, 30\% no grupo de controle tiveram um infarto, em comparação com 40\% no grupo de tratamento. Assim, o medicamento está associado a um risco maior de infarto entre os homens. O Dr. Simpson estava certo.

\textbf{TABELA 6.4.} Dados fictícios ilustrando o paradoxo de Simpson.

Mas agora olhe a terceira linha da tabela. Entre o grupo de controle, 22\% tiveram um infarto, mas no grupo de tratamento apenas 18\% o tiveram. Assim, se julgarmos com base na linha final, o Medicamento D parece \textit{diminuir} o risco de infarto na população como um todo. Bem-vindo ao mundo bizarro do paradoxo de Simpson!

Há quase vinte anos, venho tentando convencer a comunidade científica de que a confusão em torno do paradoxo de Simpson é resultado da aplicação incorreta de princípios causais a proporções estatísticas. Se usarmos notação e diagramas causais, podemos decidir de forma clara e inequívoca se o Medicamento D previne ou causa infartos. Fundamentalmente, o paradoxo de Simpson é um enigma sobre confundimento e pode, assim, ser resolvido pelos mesmos métodos que usamos para resolver aquele mistério. Curiosamente, três dos quatro artigos de 2016 que mencionei continuam a resistir a essa solução.

Qualquer pretensão de resolver um paradoxo (especialmente um que tem décadas de idade) deveria atender a alguns critérios básicos. Primeiro, como disse acima em conexão com o paradoxo de Monty Hall, deveria explicar por que as pessoas acham o paradoxo surpreendente ou inacreditável. Segundo, deveria identificar a classe de cenários em que o paradoxo pode ocorrer. Terceiro, deveria nos informar de cenários, se houver, em que o paradoxo \textit{não} pode ocorrer. Por fim, quando o paradoxo de fato ocorre, e temos de escolher entre duas afirmações plausíveis mas contraditórias, deveria nos dizer qual afirmação é correta.

Comecemos com a pergunta de por que o paradoxo de Simpson é surpreendente. Para explicar isso, deveríamos distinguir entre duas coisas: a \textit{reversão} de Simpson e o \textit{paradoxo} de Simpson. A reversão de Simpson é um fato puramente numérico: como se vê na Tabela 6.4, é uma reversão na frequência relativa de um evento particular em duas ou mais amostras diferentes ao se fundirem as amostras. Em nosso exemplo, vimos que 3/40 > 1/20 (essas eram as frequências de infarto entre mulheres com e sem o Medicamento D) e 8/20 > 12/40 (as frequências entre homens). No entanto, quando combinamos mulheres e homens, a desigualdade inverteu de direção: (3 + 8)/(40 + 20) < (1 + 12)/(20 + 40).

Se você pensou que tal reversão era matematicamente impossível, então provavelmente baseava sua reação em propriedades mal aplicadas ou mal lembradas de frações. Muitas pessoas parecem acreditar que se \textit{A}/\textit{B} > \textit{a}/\textit{b} e \textit{C}/\textit{D} > \textit{c}/\textit{d}, então se segue que (\textit{A} + \textit{C})/(\textit{B} + \textit{D}) > (\textit{a} + \textit{c})/(\textit{b} + \textit{d}). Mas essa sabedoria popular está simplesmente errada. O exemplo que acabamos de dar a refuta.

A reversão de Simpson pode ser encontrada em conjuntos de dados do mundo real. Para fãs de beisebol, eis um belo exemplo envolvendo dois astros do beisebol, David Justice e Derek Jeter. Em 1995, Justice teve uma média de rebatidas mais alta, 0,253 contra 0,250. Em 1996, Justice teve novamente uma média de rebatidas mais alta, 0,321 contra 0,314. E em 1997, ele teve uma média de rebatidas mais alta do que Jeter pela terceira temporada consecutiva, 0,329 contra 0,291. No entanto, ao longo das três temporadas combinadas, Jeter teve a média mais alta! A Tabela 6.5 mostra os cálculos para leitores que gostariam de verificá-los.

\begin{figure}[ht]
\centering
\includegraphics[width=0.92\textwidth]{../figuras/tab_6_5.png}
\caption{Tabela 6.5. mostra os cálculos para leitores que gostariam de verificá-los.}
\end{figure}
\textit{Tabela 6.5. mostra os cálculos para leitores que gostariam de verificá-los.}

Como pode um jogador ser um rebatedor pior do que o outro em 1995, 1996 e 1997, mas melhor ao longo do período de três anos? Essa reversão parece exatamente como o medicamento RRG. De fato, \textit{não} é possível; o problema é que usamos uma palavra excessivamente simples (“melhor”) para descrever um processo de média complexo sobre temporadas desiguais. Note que as vezes ao bastão (os denominadores) não estão distribuídas de forma uniforme de ano a ano. Jeter teve muito poucas vezes ao bastão em 1995, de modo que sua média de rebatidas bastante baixa naquele ano teve pouco efeito sobre sua média geral. Por outro lado, Justice teve muitas mais vezes ao bastão em seu ano menos produtivo, 1995, e isso puxou sua média geral de rebatidas para baixo. Uma vez que você percebe que “melhor rebatedor” se define não por uma competição real frente a frente, mas por uma média ponderada que leva em conta com que frequência cada jogador jogou, creio que a surpresa começa a esmaecer.

\textbf{TABELA 6.5.} Dados (não fictícios) ilustrando a reversão de Simpson.

Não há dúvida de que a reversão de Simpson é surpreendente para algumas pessoas, mesmo fãs de beisebol. Todo ano tenho alguns alunos que a princípio não conseguem acreditar nela. Mas então vão para casa, trabalham alguns exemplos como os dois que mostrei aqui, e chegam a um acordo com ela. Ela simplesmente se torna parte de sua compreensão nova e um pouco mais profunda de como os números (e especialmente os agregados de populações) funcionam. Não chamo a reversão de Simpson de paradoxo porque ela é, no máximo, simplesmente uma questão de corrigir uma crença equivocada sobre o comportamento das médias. Um paradoxo é mais do que isso: deveria envolver um conflito entre duas convicções profundamente arraigadas. Para estatísticos profissionais que trabalham com números todos os dias de suas vidas, há ainda menos razão para considerar a reversão de Simpson um paradoxo. Uma simples desigualdade aritmética não poderia possivelmente perplexá-los e fasciná-los a tal ponto que ainda estivessem escrevendo artigos sobre ela sessenta anos depois.

Agora voltemos ao nosso exemplo principal, o paradoxo do medicamento RRG. Expliquei por que as três afirmações “ruim para os homens”, “ruim para as mulheres” e “bom para as pessoas”, quando interpretadas como um aumento ou diminuição de proporções, não são matematicamente contraditórias. No entanto, ainda pode parecer a você que elas são fisicamente impossíveis. Um medicamento não pode simultaneamente fazer com que eu e você tenhamos um infarto e, ao mesmo tempo, nos impedir a ambos de ter infartos. Essa intuição é universal; a desenvolvemos como crianças de dois anos, muito antes de começarmos a aprender sobre números e frações. Assim, creio que você se sentirá aliviado ao descobrir que não precisa abandonar sua intuição. Um medicamento RRG de fato não existe e nunca será inventado, e podemos prová-lo matematicamente.

A primeira pessoa a chamar a atenção para esse princípio intuitivamente óbvio foi o estatístico Leonard Savage, que em 1954 o chamou de “princípio da coisa certa” (\textit{sure-thing principle}). Ele escreveu:

\begin{quote}Um homem de negócios contempla a compra de certa propriedade. Considera o resultado da próxima eleição presidencial relevante. Assim, para esclarecer a questão para si mesmo, pergunta se compraria se soubesse que o candidato democrata fosse vencer, e decide que sim. Do mesmo modo, considera se compraria se soubesse que o candidato republicano fosse vencer, e novamente conclui que sim. Vendo que compraria em qualquer dos eventos, decide que deveria comprar, embora não saiba qual evento se obtém, ou se obterá, como diríamos ordinariamente. É por demais raro que uma decisão possa ser alcançada com base nesse princípio, mas… não conheço nenhum outro princípio extralógico que governe decisões e que encontre aceitação tão pronta.\end{quote}
A última afirmação de Savage é particularmente perspicaz: ele percebe que o “princípio da coisa certa” é extralógico. De fato, quando devidamente interpretado, baseia-se na lógica causal, e não clássica. Também diz que “não conheço nenhum outro… princípio… que encontre aceitação tão pronta”. Obviamente ele falou sobre isso com muitas pessoas e elas acharam a linha de raciocínio muito convincente.

Para conectar o princípio da coisa certa de Savage à nossa discussão anterior, suponha que a escolha seja de fato entre duas propriedades, A e B. Se o democrata vencer, o homem de negócios tem 5\% de chance de ganhar \$1 na Propriedade A e 8\% de chance de ganhar \$1 na Propriedade B. Assim, B é preferida a A. Se o republicano vencer, ele tem 30\% de chance de ganhar \$1 na Propriedade A e 40\% de chance de ganhar \$1 na Propriedade B. Novamente, B é preferida a A. Segundo o princípio da coisa certa, ele deveria definitivamente comprar a Propriedade B. Mas leitores de olho aguçado podem notar que as quantidades numéricas são as mesmas da história de Simpson, e isso pode nos alertar de que comprar a Propriedade B pode ser uma decisão apressada demais.

De fato, o argumento acima tem uma falha gritante. Se a decisão do homem de negócios de comprar puder mudar o resultado da eleição (por exemplo, se a mídia observasse suas ações), então comprar a Propriedade A pode ser de seu melhor interesse. O dano de eleger o presidente errado pode superar qualquer ganho financeiro que ele pudesse extrair do negócio, uma vez que um presidente seja eleito. Para tornar o princípio da coisa certa válido, devemos insistir em que a decisão do homem de negócios \textit{não} afetará o resultado da eleição. Contanto que o homem de negócios tenha certeza de que sua decisão não afetará a verossimilhança de uma vitória democrata ou republicana, ele pode seguir em frente e comprar a Propriedade B. Caso contrário, todas as apostas estão canceladas.

Note que o ingrediente que faltava (e que Savage negligenciou em declarar explicitamente) é uma hipótese causal. Uma versão correta de seu princípio leria o seguinte: uma ação que aumenta a probabilidade de certo desfecho assumindo ou que o Evento C ocorreu ou que o Evento C não ocorreu também aumentará sua probabilidade se não soubermos se C ocorreu… \textit{desde que a ação não mude a probabilidade de C}.

Em particular, não existe tal coisa como um medicamento RRG. Essa versão corrigida do princípio da coisa certa de Savage não se segue da lógica clássica: para prová-la, você precisa de um cálculo causal invocando o operador \textit{do}. Nossa forte crença intuitiva de que um medicamento RRG é impossível sugere que os humanos (assim como as máquinas programadas para emular o pensamento humano) usam algo como o cálculo do \textit{do} para guiar sua intuição.

Segundo o princípio da coisa certa corrigido, uma das três afirmações seguintes deve ser falsa: o Medicamento D aumenta a probabilidade de infarto em homens e mulheres; o Medicamento D diminui a probabilidade de infarto na população como um todo; e o medicamento não muda o número de homens e mulheres. Como é muito implausível que um medicamento mude o sexo de um paciente, uma das duas primeiras afirmações deve ser falsa. Qual delas? Em vão buscará orientação na Tabela 6.4. Para responder à pergunta, devemos olhar além dos dados, para o processo gerador dos dados. Como sempre, é praticamente impossível discutir esse processo sem um diagrama causal.

O diagrama na Figura 6.4 codifica a informação crucial de que o sexo não é afetado pelo medicamento e, além disso, o sexo afeta o risco de infarto (homens estando em maior risco) e se o paciente escolhe tomar o Medicamento D. No estudo, as mulheres claramente tinham preferência por tomar o Medicamento D e os homens preferiam não tomá-lo. Assim, Sexo é um confundidor de Medicamento e Infarto. Para uma estimativa não enviesada do efeito de Medicamento sobre Infarto, devemos ajustar pelo confundidor. Podemos fazer isso olhando os dados de homens e mulheres separadamente e então tomando a média:

\begin{figure}[ht]
\centering
\includegraphics[width=0.92\textwidth]{../figuras/fig_6_4.png}
\caption{Figura 6.4. codifica a informação crucial de que o sexo não é afetado pelo medicamento e, além disso, o sexo afeta o risco de infarto (homens estando em maior risco) e se o paciente escolhe tomar o Medicamento D. No estudo, as mulheres claramente tinham preferência por tomar o Medicamento D e os homens preferiam não tomá-lo. Assim, Sexo é um confundidor de Medicamento e Infarto. Para uma estimativa não enviesada do efeito de Medicamento sobre Infarto, devemos ajustar pelo confundidor. Podemos fazer isso olhando os dados de homens e mulheres separadamente e então tomando a média:}
\end{figure}
\textit{Figura 6.4. codifica a informação crucial de que o sexo não é afetado pelo medicamento e, além disso, o sexo afeta o risco de infarto (homens estando em maior risco) e se o paciente escolhe tomar o Medicamento D. No estudo, as mulheres claramente tinham preferência por tomar o Medicamento D e os homens preferiam não tomá-lo. Assim, Sexo é um confundidor de Medicamento e Infarto. Para uma estimativa não enviesada do efeito de Medicamento sobre Infarto, devemos ajustar pelo confundidor. Podemos fazer isso olhando os dados de homens e mulheres separadamente e então tomando a média:}

\textbf{FIGURA 6.4.} Diagrama causal para o exemplo do paradoxo de Simpson.

\begin{itemize}
\item Para as mulheres, a taxa de infartos foi de 5\% sem o Medicamento D e de 7,5\% com o Medicamento D.
\item Para os homens, a taxa de infartos foi de 30\% sem o Medicamento D e de 40\% com.
\item Tomando a média (porque homens e mulheres são igualmente frequentes na população geral), a taxa de infartos sem o Medicamento D é de 17,5\% (a média de 5 e 30), e a taxa com o Medicamento D é de 23,75\% (a média de 7,5 e 40).
\end{itemize}
Essa é a resposta clara e inequívoca que buscávamos. O Medicamento D não é RRG: é RRR — ruim para as mulheres, ruim para os homens e ruim para as pessoas.

Não quero que você tire a impressão, a partir deste exemplo, de que agregar os dados é sempre errado ou de que particionar os dados é sempre certo. Depende do processo que gerou os dados. No paradoxo de Monty Hall, vimos que mudar as regras do jogo também mudava a conclusão. O mesmo princípio funciona aqui.

Usarei uma história diferente para demonstrar quando reunir os dados em \textit{pool} seria apropriado. Embora os dados sejam precisamente os mesmos, o papel da “terceira variável à espreita” será diferente e, portanto, também a conclusão.

Comecemos com a premissa de que a pressão arterial é conhecida como uma possível causa de infarto, e o Medicamento B é suposto reduzir a pressão arterial. Naturalmente, os pesquisadores do Medicamento B queriam ver se ele também poderia reduzir o risco de infarto, de modo que mediram a pressão arterial de seus pacientes após o tratamento, bem como se tiveram um infarto. A Tabela 6.6 mostra os dados do estudo do Medicamento B. Deveria parecer assombrosamente familiar: os números são os mesmos da Tabela 6.4! Não obstante, a conclusão é exatamente a oposta. Como se pode ver, tomar o Medicamento B conseguiu baixar a pressão arterial dos pacientes: entre as pessoas que o tomaram, o dobro tinha pressão arterial baixa depois (quarenta em sessenta, em comparação com vinte em sessenta no grupo de controle). Em outras palavras, fez exatamente o que um medicamento anti-infarto deveria fazer. Moveu pessoas da categoria de maior risco para a categoria de menor risco. Esse fator supera tudo o mais, e podemos concluir de forma justificável que a parte agregada da Tabela 6.6 nos dá o resultado correto.

\begin{figure}[ht]
\centering
\includegraphics[width=0.92\textwidth]{../figuras/tab_6_6.png}
\caption{Tabela 6.6. mostra os dados do estudo do Medicamento B. Deveria parecer assombrosamente familiar: os números são os mesmos da Tabela 6.4! Não obstante, a conclusão é exatamente a oposta. Como se pode ver, tomar o Medicamento B conseguiu baixar a pressão arterial dos pacientes: entre as pessoas que o tomaram, o dobro tinha pressão arterial baixa depois (quarenta em sessenta, em comparação com vinte em sessenta no grupo de controle). Em outras palavras, fez exatamente o que um medicamento anti-infarto deveria fazer. Moveu pessoas da categoria de maior risco para a categoria de menor risco. Esse fator supera tudo o mais, e podemos concluir de forma justificável que a parte agregada da Tabela 6.6 nos dá o resultado correto.}
\end{figure}
\textit{Tabela 6.6. mostra os dados do estudo do Medicamento B. Deveria parecer assombrosamente familiar: os números são os mesmos da Tabela 6.4! Não obstante, a conclusão é exatamente a oposta. Como se pode ver, tomar o Medicamento B conseguiu baixar a pressão arterial dos pacientes: entre as pessoas que o tomaram, o dobro tinha pressão arterial baixa depois (quarenta em sessenta, em comparação com vinte em sessenta no grupo de controle). Em outras palavras, fez exatamente o que um medicamento anti-infarto deveria fazer. Moveu pessoas da categoria de maior risco para a categoria de menor risco. Esse fator supera tudo o mais, e podemos concluir de forma justificável que a parte agregada da Tabela 6.6 nos dá o resultado correto.}

\textbf{TABELA 6.6.} Dados fictícios para o exemplo da pressão arterial.

Como de costume, um diagrama causal tornará tudo claro e nos permitirá derivar o resultado mecanicamente, sem sequer pensar nos dados ou se o medicamento baixa ou eleva a pressão arterial. Nesse caso, nossa “terceira variável à espreita” é Pressão Arterial, e o diagrama se parece com a Figura 6.5. Aqui, Pressão Arterial é um mediador, e não um confundidor. Um único olhar ao diagrama revela que não há confundidor da relação Medicamento–Infarto (isto é, nenhum caminho pela porta dos fundos), de modo que estratificar os dados é desnecessário. De fato, condicionar a Pressão Arterial desativaria um dos caminhos causais (talvez o principal caminho causal) pelos quais o medicamento funciona. Por ambas essas razões, nossa conclusão é exatamente a oposta do que era para o Medicamento D: o Medicamento B funciona, e os dados agregados revelam esse fato.

\begin{figure}[ht]
\centering
\includegraphics[width=0.92\textwidth]{../figuras/fig_6_5.png}
\caption{Figura 6.5. Aqui, Pressão Arterial é um mediador, e não um confundidor. Um único olhar ao diagrama revela que não há confundidor da relação Medicamento–Infarto (isto é, nenhum caminho pela porta dos fundos), de modo que estratificar os dados é desnecessário. De fato, condicionar a Pressão Arterial desativaria um dos caminhos causais (talvez o principal caminho causal) pelos quais o medicamento funciona. Por ambas essas razões, nossa conclusão é exatamente a oposta do que era para o Medicamento D: o Medicamento B funciona, e os dados agregados revelam esse fato.}
\end{figure}
\textit{Figura 6.5. Aqui, Pressão Arterial é um mediador, e não um confundidor. Um único olhar ao diagrama revela que não há confundidor da relação Medicamento–Infarto (isto é, nenhum caminho pela porta dos fundos), de modo que estratificar os dados é desnecessário. De fato, condicionar a Pressão Arterial desativaria um dos caminhos causais (talvez o principal caminho causal) pelos quais o medicamento funciona. Por ambas essas razões, nossa conclusão é exatamente a oposta do que era para o Medicamento D: o Medicamento B funciona, e os dados agregados revelam esse fato.}

Do ponto de vista histórico, é digno de nota que Simpson, em seu artigo de 1951 que iniciou toda a confusão, fez exatamente a mesma coisa que acabo de fazer. Apresentou duas histórias com exatamente os mesmos dados. Num exemplo, era intuitivamente claro que agregar os dados era, em suas palavras, “a interpretação sensata”; no outro exemplo, particionar os dados era mais sensato. Assim, Simpson compreendeu que havia um paradoxo, e não apenas uma reversão. No entanto, não sugeriu nenhuma resolução para o paradoxo além do senso comum. Mais importante, não sugeriu que, se a história contém informação extra que faz a diferença entre “sensato” e “não sensato”, talvez os estatísticos devessem abraçar essa informação extra em sua análise.

\textbf{FIGURA 6.5.} Diagrama causal para o exemplo do paradoxo de Simpson (segunda versão).

Dennis Lindley e Melvin Novick consideraram essa sugestão em 1981, mas não puderam se reconciliar com a ideia de que a decisão correta depende da história causal, e não dos dados. Confessaram: “Uma possibilidade seria usar a linguagem da causalidade.… Não escolhemos fazer isso; nem discutir a causalidade, porque o conceito, embora amplamente usado, não parece estar bem definido.” Com essas palavras, resumiram a frustração de cinco gerações de estatísticos, reconhecendo que a informação causal é urgentemente necessária, mas a linguagem para expressá-la é desesperadamente carente. Em 2009, quatro anos antes de sua morte aos noventa anos, Lindley confidenciou-me que não teria escrito aquelas palavras se meu livro tivesse estado disponível em 1981.

Alguns leitores de meus livros e artigos sugeriram que a regra que governa a agregação e a separação de dados repousa simplesmente na precedência temporal do tratamento e da “terceira variável à espreita”. Argumentam que deveríamos agregar os dados no caso da pressão arterial porque a medição da pressão arterial vem \textit{depois} de o paciente tomar o medicamento, mas deveríamos estratificar os dados no caso do sexo porque ele é determinado \textit{antes} de o paciente tomar o medicamento. Embora essa regra funcione num grande número de casos, não é infalível. Um caso simples é o do viés em M (Jogo 4 no Capítulo 4). Aqui B pode preceder A; no entanto, ainda assim não deveríamos condicionar a B, porque isso violaria o critério da porta dos fundos. Deveríamos consultar a estrutura causal da história, e não a informação temporal.

Por fim, você pode se perguntar se o paradoxo de Simpson ocorre no mundo real. A resposta é sim. Certamente não é comum o bastante para que estatísticos o encontrem numa base diária, mas tampouco é completamente desconhecido, e provavelmente acontece com mais frequência do que os artigos de periódicos reportam. Eis dois casos documentados:

\begin{itemize}
\item Num estudo observacional publicado em 1996, a cirurgia aberta para remover cálculos renais teve uma taxa de sucesso melhor do que a cirurgia endoscópica para cálculos renais pequenos. Também teve uma taxa de sucesso melhor para cálculos renais grandes. No entanto, teve uma taxa de sucesso \textit{menor} no geral. Exatamente como em nosso primeiro exemplo, esse foi um caso em que a escolha do tratamento estava relacionada à gravidade do caso dos pacientes: cálculos maiores eram mais propensos a levar à cirurgia aberta e também tinham um pior prognóstico.
\item Num estudo de doença da tireoide publicado em 1995, os fumantes tinham uma taxa de sobrevivência mais alta (76\%) ao longo de vinte anos do que os não fumantes (69\%). No entanto, os não fumantes tinham uma taxa de sobrevivência melhor em seis de sete faixas etárias, e a diferença era mínima na sétima. A idade era claramente um confundidor de Fumo e Sobrevivência: o fumante médio era mais jovem do que o não fumante médio (talvez porque os fumantes mais velhos já tivessem morrido). Estratificando os dados por idade, concluímos que o fumo tem um impacto negativo sobre a sobrevivência.
\end{itemize}
Como o paradoxo de Simpson tem sido tão mal compreendido, alguns estatísticos tomam precauções para evitá-lo. Com demasiada frequência, esses métodos evitam o sintoma, a reversão de Simpson, sem fazer nada sobre a doença, o confundimento. Em vez de suprimir os sintomas, deveríamos prestar-lhes atenção. O paradoxo de Simpson nos alerta para casos em que ao menos uma das tendências estatísticas (seja nos dados agregados, nos dados particionados, ou em ambos) não pode representar os efeitos causais. Há, é claro, outros sinais de alerta de confundimento. A estimativa agregada do efeito causal poderia, por exemplo, ser maior do que cada uma das estimativas em cada um dos estratos; isso igualmente não deveria acontecer se tivermos controlado adequadamente por confundidores. Em comparação com tais sinais, no entanto, a reversão de Simpson é mais difícil de ignorar precisamente porque é uma \textit{reversão}, uma mudança qualitativa no sinal do efeito. A ideia de um medicamento RRG evocaria incredulidade até de uma criança de três anos — e com razão.

\section{O PARADOXO DE SIMPSON EM IMAGENS}
Até agora, a maioria de nossos exemplos de reversão e paradoxo de Simpson envolveu variáveis binárias: um paciente ou recebeu o Medicamento D ou não, e ou teve um infarto ou não. No entanto, a reversão também pode ocorrer com variáveis contínuas e é talvez mais fácil de compreender nesse caso porque podemos desenhar uma figura.

Considere um estudo que mede o exercício semanal e os níveis de colesterol em várias faixas etárias. Quando plotamos as horas de exercício no eixo \textit{x} e o colesterol no eixo \textit{y}, como na Figura 6.6(a), vemos em cada faixa etária uma tendência descendente, indicando talvez que o exercício reduz o colesterol. Por outro lado, se usarmos o mesmo gráfico de dispersão mas não segregarmos os dados por idade, como na Figura 6.6(b), então vemos uma tendência ascendente pronunciada, indicando que quanto mais as pessoas se exercitam, mais alto se torna o seu colesterol. Mais uma vez, parecemos ter uma situação de medicamento RRG, em que Exercício é o medicamento: ele parece ter um efeito benéfico em cada faixa etária, mas um efeito prejudicial na população como um todo.

\begin{figure}[ht]
\centering
\includegraphics[width=0.92\textwidth]{../figuras/fig_6_6.png}
\caption{Figura 6.6. (a), vemos em cada faixa etária uma tendência descendente, indicando talvez que o exercício reduz o colesterol. Por outro lado, se usarmos o mesmo gráfico de dispersão mas não segregarmos os dados por idade, como na Figura 6.6(b), então vemos uma tendência ascendente pronunciada, indicando que quanto mais as pessoas se exercitam, mais alto se torna o seu colesterol. Mais uma vez, parecemos ter uma situação de medicamento RRG, em que Exercício é o medicamento: ele parece ter um efeito benéfico em cada faixa etária, mas um efeito prejudicial na população como um todo.}
\end{figure}
\textit{Figura 6.6. (a), vemos em cada faixa etária uma tendência descendente, indicando talvez que o exercício reduz o colesterol. Por outro lado, se usarmos o mesmo gráfico de dispersão mas não segregarmos os dados por idade, como na Figura 6.6(b), então vemos uma tendência ascendente pronunciada, indicando que quanto mais as pessoas se exercitam, mais alto se torna o seu colesterol. Mais uma vez, parecemos ter uma situação de medicamento RRG, em que Exercício é o medicamento: ele parece ter um efeito benéfico em cada faixa etária, mas um efeito prejudicial na população como um todo.}

Para decidir se o Exercício é benéfico ou prejudicial, como sempre, precisamos consultar a história por trás dos dados. Os dados mostram que as pessoas mais velhas em nossa população se exercitam mais. Como parece mais provável que a Idade cause o Exercício, e não o contrário, e como a Idade pode ter um efeito causal sobre o Colesterol, concluímos que a Idade pode ser um confundidor de Exercício e Colesterol. Assim, deveríamos controlar por Idade. Em outras palavras, deveríamos olhar os dados segregados por idade e concluir que o exercício é benéfico, independentemente da idade.

\textbf{FIGURA 6.6.} Paradoxo de Simpson: o exercício parece ser benéfico (inclinação descendente) em cada faixa etária, mas prejudicial (inclinação ascendente) na população como um todo.

Um primo do paradoxo de Simpson também tem espreitado na literatura estatística por décadas e se presta bem a uma interpretação visual. Frederic Lord originalmente enunciou esse paradoxo em 1967. É novamente fictício, mas exemplos fictícios (como os experimentos mentais de Einstein) sempre fornecem uma boa maneira de sondar os limites de nossa compreensão.

Lord postula uma escola que quer estudar os efeitos da dieta que fornece em seus refeitórios e, em particular, se ela tem efeitos diferentes sobre meninas e meninos. Para esse fim, o peso dos alunos é medido em setembro e novamente no junho seguinte. A Figura 6.7 plota os resultados, com as elipses representando mais uma vez um gráfico de dispersão de dados.

\begin{figure}[ht]
\centering
\includegraphics[width=0.92\textwidth]{../figuras/fig_6_7.png}
\caption{Figura 6.7. plota os resultados, com as elipses representando mais uma vez um gráfico de dispersão de dados.}
\end{figure}
\textit{Figura 6.7. plota os resultados, com as elipses representando mais uma vez um gráfico de dispersão de dados.}

A universidade contrata dois estatísticos, que olham os dados e chegam a conclusões opostas.

O primeiro estatístico olha a distribuição de peso das meninas como um todo e nota que o peso médio das meninas é o mesmo em junho e em setembro. (Isso se pode ver a partir da simetria do gráfico de dispersão em torno da linha \textit{W<sub>F</sub>} = \textit{W<sub>I</sub>}, isto é, peso final = peso inicial.) Meninas individuais podem, é claro, ganhar ou perder peso, mas o ganho de peso médio é zero. A mesma observação vale para os meninos. Portanto, o estatístico conclui que a dieta não tem efeito diferencial sobre os sexos.

\textbf{FIGURA 6.7.} Paradoxo de Lord. (As elipses representam gráficos de dispersão de dados.) Como um todo, nem meninos nem meninas ganham peso durante o ano, mas em cada estrato do peso inicial, os meninos tendem a ganhar mais do que as meninas.

O segundo estatístico, por outro lado, argumenta que, como o peso final de um aluno é fortemente influenciado por seu peso inicial, deveríamos estratificar os alunos por peso inicial. Se você fizer um corte vertical através de ambas as elipses, o que corresponde a olhar apenas os meninos e as meninas com um valor particular do peso inicial (digamos \textit{W<sub>0</sub>} na Figura 6.7), notará que a linha vertical intercepta a elipse dos Meninos mais acima do que a elipse das Meninas, embora haja certa quantidade de sobreposição. Isso significa que os meninos que começaram com peso \textit{W<sub>0</sub>} terão, em média, um peso final (\textit{W<sub>F</sub>}) mais alto do que as meninas que começaram com peso \textit{W<sub>0</sub>}. Conformemente, Lord escreve: “o segundo estatístico conclui, como é costumeiro em tais casos, que os meninos mostraram um ganho de peso significativamente maior do que as meninas quando se faz a devida consideração pelas diferenças no peso inicial entre os sexos.”

O que a nutricionista da escola deve fazer? Lord escreve: “As conclusões de cada estatístico estão visivelmente corretas.” Isto é, você não precisa esmagar nenhum número para ver que dois argumentos sólidos estão levando a duas conclusões diferentes. Você precisa apenas olhar a figura. Na Figura 6.7, podemos ver que os meninos ganham mais peso do que as meninas em cada estrato (cada seção transversal vertical). No entanto, é igualmente óbvio que tanto meninos quanto meninas não ganharam nada no geral. Como isso pode ser? Não é o ganho geral apenas uma média dos ganhos específicos de cada estrato?

Agora que somos profissionais experientes nos pontos finos do paradoxo de Simpson e do princípio da coisa certa, sabemos o que há de errado com esse argumento. O princípio da coisa certa funciona apenas em casos em que a proporção relativa de cada subpopulação (cada classe de peso) não muda de grupo a grupo. No entanto, no caso de Lord, o “tratamento” (sexo) afeta muito fortemente a porcentagem de alunos em cada classe de peso. Assim, não podemos nos apoiar no princípio da coisa certa, e isso nos traz de volta à estaca zero. Quem está certo? Há ou não há uma diferença nos ganhos de peso médios entre meninos e meninas quando se faz a devida consideração pelas diferenças no peso inicial entre os sexos?

A conclusão de Lord é muito pessimista: “O estudo de pesquisa usual desse tipo está tentando responder a uma pergunta que simplesmente não pode ser respondida de nenhuma maneira rigorosa com base nos dados disponíveis.” O pessimismo de Lord se espalhou para além da estatística e levou a uma literatura rica e bastante pessimista em epidemiologia e bioestatística sobre como comparar grupos que diferem em estatísticas de “linha de base”. Mostrarei agora por que o pessimismo de Lord é injustificado. A pergunta da nutricionista \textit{pode} ser respondida de maneira rigorosa, e como de costume o ponto de partida é desenhar um diagrama causal, como na Figura 6.8.

\begin{figure}[ht]
\centering
\includegraphics[width=0.92\textwidth]{../figuras/fig_6_8.png}
\caption{Figura 6.8. Nesse diagrama, vemos que Sexo (\textit{S}) é uma causa do peso inicial (\textit{W<sub>I</sub>}) e do peso final (\textit{W<sub>F</sub>}). Também, \textit{W<sub>I</sub>} afeta \textit{W<sub>F</sub>} independentemente do sexo, porque alunos de qualquer sexo que pesam mais no início do ano tendem a pesar mais no final do ano, como mostrado pelos gráficos de dispersão na Figura 6.7. Todas essas hipóteses causais são de senso comum; eu não esperaria que Lord discordasse delas. A variável de interesse para Lord é o ganho de peso, mostrado como \textit{Y} nesse diagrama. Note que \textit{Y} está relacionado a \textit{W<sub>I</sub>} e \textit{W<sub>F</sub>} de maneira puramente matemática, determinística: \textit{Y} = \textit{W<sub>F</sub>} − \textit{W<sub>I</sub>}. Isso significa que as correlações entre \textit{Y} e \textit{W<sub>I</sub>} (ou \textit{Y} e \textit{W<sub>F</sub>}) são iguais a −1 (ou 1), e mostrei essa informação no diagrama com os coeficientes −1 e +1.}
\end{figure}
\textit{Figura 6.8. Nesse diagrama, vemos que Sexo (}S\textit{) é uma causa do peso inicial (}W<sub>I</sub>\textit{) e do peso final (}W<sub>F</sub>\textit{). Também, }W<sub>I</sub>\textit{ afeta }W<sub>F</sub>\textit{ independentemente do sexo, porque alunos de qualquer sexo que pesam mais no início do ano tendem a pesar mais no final do ano, como mostrado pelos gráficos de dispersão na Figura 6.7. Todas essas hipóteses causais são de senso comum; eu não esperaria que Lord discordasse delas. A variável de interesse para Lord é o ganho de peso, mostrado como }Y\textit{ nesse diagrama. Note que }Y\textit{ está relacionado a }W<sub>I</sub>\textit{ e }W<sub>F</sub>\textit{ de maneira puramente matemática, determinística: }Y\textit{ = }W<sub>F</sub>\textit{ − }W<sub>I</sub>\textit{. Isso significa que as correlações entre }Y\textit{ e }W<sub>I</sub>\textit{ (ou }Y\textit{ e }W<sub>F</sub>\textit{) são iguais a −1 (ou 1), e mostrei essa informação no diagrama com os coeficientes −1 e +1.}

Nesse diagrama, vemos que Sexo (\textit{S}) é uma causa do peso inicial (\textit{W<sub>I</sub>}) e do peso final (\textit{W<sub>F</sub>}). Também, \textit{W<sub>I</sub>} afeta \textit{W<sub>F</sub>} independentemente do sexo, porque alunos de qualquer sexo que pesam mais no início do ano tendem a pesar mais no final do ano, como mostrado pelos gráficos de dispersão na Figura 6.7. Todas essas hipóteses causais são de senso comum; eu não esperaria que Lord discordasse delas. A variável de interesse para Lord é o ganho de peso, mostrado como \textit{Y} nesse diagrama. Note que \textit{Y} está relacionado a \textit{W<sub>I</sub>} e \textit{W<sub>F</sub>} de maneira puramente matemática, determinística: \textit{Y} = \textit{W<sub>F</sub>} − \textit{W<sub>I</sub>}. Isso significa que as correlações entre \textit{Y} e \textit{W<sub>I</sub>} (ou \textit{Y} e \textit{W<sub>F</sub>}) são iguais a −1 (ou 1), e mostrei essa informação no diagrama com os coeficientes −1 e +1.

\textbf{FIGURA 6.8.} Diagrama causal para o paradoxo de Lord.

O primeiro estatístico simplesmente compara a diferença no ganho de peso entre meninas e meninos. Nenhuma porta dos fundos entre \textit{S} e \textit{Y} precisa ser bloqueada, de modo que os dados observados e agregados fornecem a resposta: nenhum efeito, como o primeiro estatístico concluiu.

Em contraste, é difícil até formular a pergunta que o segundo estatístico está tentando responder (isto é, a “consulta corretamente formulada” descrita na Introdução). Ele quer assegurar que “se faça a devida consideração pelas diferenças no peso inicial entre os dois sexos”, que é a linguagem que você normalmente usaria ao controlar por um confundidor. Mas \textit{W<sub>I</sub>} não é um confundidor de \textit{S} e \textit{Y}. É de fato uma variável mediadora se considerarmos Sexo como o tratamento. Assim, a consulta respondida ao controlar por \textit{W<sub>I</sub>} não tem a interpretação usual de efeito causal. Tal controle pode, na melhor das hipóteses, fornecer uma estimativa do “efeito direto” do sexo sobre o peso, que discutiremos no Capítulo 9. No entanto, parece improvável que seja isso o que o segundo estatístico tinha em mente; mais provavelmente ele estava ajustando por hábito.

E, no entanto, seu argumento é uma armadilha tão fácil de cair: “Não é o ganho geral apenas uma média dos ganhos específicos de cada estrato?” Não se os próprios estratos estão se deslocando sob o tratamento! Lembre-se de que Sexo, e não Dieta, é o tratamento, e Sexo definitivamente muda a proporção de alunos em cada estrato de \textit{W<sub>I</sub>}.

Este último comentário levanta mais um ponto curioso sobre o paradoxo de Lord como originalmente formulado. Embora a intenção declarada da nutricionista da escola seja “determinar os efeitos da dieta”, em nenhum lugar de seu artigo original Lord menciona uma dieta de controle. Portanto, não podemos sequer dizer nada sobre os efeitos da dieta. Um artigo de 2006 de Howard Wainer e Lisa Brown tenta remediar esse defeito. Eles mudam a história de modo que a quantidade de interesse seja o efeito da dieta (não do sexo) sobre o ganho de peso, enquanto as diferenças de sexo não são consideradas. Em sua versão, os alunos comem em um de dois refeitórios com dietas diferentes. Conformemente, as duas elipses da Figura 6.7 representam dois refeitórios, cada um servindo uma dieta diferente, como se depreende da Figura 6.9(a). Note que os alunos que pesam mais no início tendem a comer no refeitório B, enquanto os que pesam menos comem no refeitório A.

\begin{figure}[ht]
\centering
\includegraphics[width=0.92\textwidth]{../figuras/fig_6_9.png}
\caption{Figura 6.9. [a]) não podem dizer a quem acreditar, e é de fato isso que Wainer e Brown concluem. No entanto, um diagrama causal (Figura 6.9[b]) pode resolver a questão. Há duas mudanças significativas entre a Figura 6.8 e a Figura 6.9(b). Primeiro, a variável causal se torna \textit{D} (de “dieta”), e não \textit{S}. Segundo, a seta que originalmente apontava de \textit{S} para \textit{W<sub>I</sub>} agora inverte de direção: o peso inicial agora afeta a dieta, de modo que a seta aponta de \textit{W<sub>I</sub>} para \textit{D}. Nesse diagrama, \textit{W<sub>I</sub>} é um confundidor de \textit{D} e \textit{W<sub>F</sub>}, e não um mediador. Portanto, o segundo estatístico estaria inequivocamente correto aqui. Controlar pelo peso inicial é essencial para desconfundir \textit{D} e \textit{W<sub>F</sub>} (bem como \textit{D} e \textit{Y}). O primeiro estatístico estaria errado, porque estaria medindo apenas associações estatísticas, e não efeitos causais.}
\end{figure}
\textit{Figura 6.9. [a]) não podem dizer a quem acreditar, e é de fato isso que Wainer e Brown concluem. No entanto, um diagrama causal (Figura 6.9[b]) pode resolver a questão. Há duas mudanças significativas entre a Figura 6.8 e a Figura 6.9(b). Primeiro, a variável causal se torna }D\textit{ (de “dieta”), e não }S\textit{. Segundo, a seta que originalmente apontava de }S\textit{ para }W<sub>I</sub>\textit{ agora inverte de direção: o peso inicial agora afeta a dieta, de modo que a seta aponta de }W<sub>I</sub>\textit{ para }D\textit{. Nesse diagrama, }W<sub>I</sub>\textit{ é um confundidor de }D\textit{ e }W<sub>F</sub>\textit{, e não um mediador. Portanto, o segundo estatístico estaria inequivocamente correto aqui. Controlar pelo peso inicial é essencial para desconfundir }D\textit{ e }W<sub>F</sub>\textit{ (bem como }D\textit{ e }Y\textit{). O primeiro estatístico estaria errado, porque estaria medindo apenas associações estatísticas, e não efeitos causais.}

O paradoxo de Lord agora surge com maior clareza, pois a consulta está bem definida como o efeito da dieta sobre o ganho. O primeiro estatístico afirma, com base em considerações de simetria, que trocar da Dieta A para a B não teria efeito sobre o ganho de peso (a diferença \textit{W<sub>F</sub>} − \textit{W<sub>I</sub>} tem a mesma distribuição em ambas as elipses). O segundo estatístico compara os pesos finais sob a Dieta A com os da Dieta B para um grupo de alunos que começam com peso \textit{W<sub>0</sub>} e conclui que os alunos na Dieta B ganham mais peso.

Como antes, os dados (Figura 6.9[a]) não podem dizer a quem acreditar, e é de fato isso que Wainer e Brown concluem. No entanto, um diagrama causal (Figura 6.9[b]) pode resolver a questão. Há duas mudanças significativas entre a Figura 6.8 e a Figura 6.9(b). Primeiro, a variável causal se torna \textit{D} (de “dieta”), e não \textit{S}. Segundo, a seta que originalmente apontava de \textit{S} para \textit{W<sub>I</sub>} agora inverte de direção: o peso inicial agora afeta a dieta, de modo que a seta aponta de \textit{W<sub>I</sub>} para \textit{D}. Nesse diagrama, \textit{W<sub>I</sub>} é um confundidor de \textit{D} e \textit{W<sub>F</sub>}, e não um mediador. Portanto, o segundo estatístico estaria inequivocamente correto aqui. Controlar pelo peso inicial é essencial para desconfundir \textit{D} e \textit{W<sub>F</sub>} (bem como \textit{D} e \textit{Y}). O primeiro estatístico estaria errado, porque estaria medindo apenas associações estatísticas, e não efeitos causais.

Para resumir, para nós a principal lição do paradoxo de Lord é que ele não é mais um paradoxo do que o de Simpson. Num paradoxo, a associação se inverte; no outro, ela desaparece. Em qualquer dos casos, o diagrama causal nos dirá que procedimento precisamos usar. No entanto, para estatísticos treinados na metodologia “convencional” (isto é, cega a modelos) e que evitam usar lentes causais, é profundamente paradoxal que a conclusão correta num caso seja incorreta no outro, embora os dados pareçam exatamente os mesmos.

\textbf{FIGURA 6.9.} Versão revisada de Wainer e Brown do paradoxo de Lord e o diagrama causal correspondente.

Agora que temos um sólido embasamento em colisores, confundidores e os perigos que ambos representam, estamos enfim preparados para colher os frutos de nosso trabalho. No próximo capítulo, começamos nossa ascensão da Escada da Causalidade, começando pelo degrau dois: a intervenção.

\textbf{Frontispício do capítulo.} Escalando o “Monte Intervenção”. Os métodos mais familiares para estimar o efeito de uma intervenção, na presença de confundidores, são o ajustamento pela porta dos fundos e as variáveis instrumentais. O método de ajustamento pela porta da frente era desconhecido antes da introdução dos diagramas causais. O cálculo do \textit{do}, que meus alunos automatizaram plenamente, torna possível adaptar o método de ajustamento a qualquer diagrama causal particular. (Fonte: Desenho de Dakota Harr.)

\section{Figuras e tabelas deste capítulo}
\begin{figure}[ht]
\centering
\includegraphics[width=0.92\textwidth]{../figuras/fig_6_2.png}
\caption{Figura 6.2. mostra o diagrama causal para \textit{Let’s Fake a Deal}, o jogo em que Monty Hall escolhe uma porta que é diferente da sua, mas de resto escolhida ao acaso. Esse diagrama ainda tem uma seta apontando de Sua Porta para Porta Aberta porque ele tem de se assegurar de que sua porta seja diferente da sua. No entanto, a seta de Localização do Carro para Porta Aberta seria apagada porque ele já não se importa com o local do carro. Nesse diagrama, condicionar a Porta Aberta não tem absolutamente nenhum efeito: Sua Porta e Localização do Carro eram independentes desde o início, e permanecem independentes depois que vemos o conteúdo da porta de Monty. Assim, em \textit{Let’s Fake a Deal}, o carro tem a mesma probabilidade de estar atrás da sua porta ou da outra, como se observa na Tabela 6.2.}
\end{figure}
\textbf{Figura 6.2.} mostra o diagrama causal para \textit{Let’s Fake a Deal}, o jogo em que Monty Hall escolhe uma porta que é diferente da sua, mas de resto escolhida ao acaso. Esse diagrama ainda tem uma seta apontando de Sua Porta para Porta Aberta porque ele tem de se assegurar de que sua porta seja diferente da sua. No entanto, a seta de Localização do Carro para Porta Aberta seria apagada porque ele já não se importa com o local do carro. Nesse diagrama, condicionar a Porta Aberta não tem absolutamente nenhum efeito: Sua Porta e Localização do Carro eram independentes desde o início, e permanecem independentes depois que vemos o conteúdo da porta de Monty. Assim, em \textit{Let’s Fake a Deal}, o carro tem a mesma probabilidade de estar atrás da sua porta ou da outra, como se observa na Tabela 6.2.

\chapter{Cap. 7: Além do ajustamento: a conquista do Monte Intervenção}
\textit{Aquele cujas ações excedem a sua teoria, a sua teoria perdurará.} — Rabi Hanina ben Dosa (século I d.C.)

Neste capítulo, enfim, empreendemos a ousada escalada até o segundo nível da Escada da Causalidade: o nível da intervenção — o Santo Graal do pensamento causal desde a Antiguidade até os nossos dias. Esse nível diz respeito à luta para prever os efeitos de ações e políticas ainda não ensaiadas, que vão de tratamentos médicos a programas sociais, de políticas econômicas a escolhas pessoais. O confundimento foi o principal obstáculo que nos levava a confundir \textit{ver} com \textit{fazer}. Removido esse obstáculo com as ferramentas do “bloqueio de caminhos” e do critério da porta dos fundos, podemos agora mapear as rotas de subida do Monte Intervenção com precisão sistemática. Para o escalador principiante, as rotas mais seguras são o ajustamento pela porta dos fundos e seus diversos primos, uns sob a rubrica do “ajustamento pela porta da frente” e outros sob a de “variáveis instrumentais”. Mas essas rotas nem sempre estão disponíveis; por isso, para o escalador experiente, este capítulo descreve uma “ferramenta de mapeamento universal” chamada \textit{cálculo do do} (\textit{do}-cálculo), que permite ao pesquisador explorar e traçar todas as rotas possíveis de subida do Monte Intervenção, por tortuosas que sejam. Uma vez mapeada a rota, e com cordas, mosquetões e pitões no lugar, o assalto à montanha resultará, com certeza, em conquista bem-sucedida!

\section{A rota mais simples: a fórmula de ajustamento pela porta dos fundos}
Para muitos pesquisadores, o método mais familiar (talvez o único) de prever o efeito de uma intervenção é “controlar” os confundidores por meio da fórmula de ajustamento. Esse é o método a usar se você tem confiança de que dispõe de dados sobre um conjunto suficiente de variáveis (chamadas \textit{desconfundidores}) para bloquear todos os caminhos pela porta dos fundos entre a intervenção e o desfecho. Para isso, medimos o efeito causal médio de uma intervenção estimando primeiro o seu efeito em cada “nível”, ou estrato, do desconfundidor. Em seguida, calculamos uma média ponderada desses estratos, em que cada estrato é ponderado de acordo com a sua prevalência na população. Se, por exemplo, o desconfundidor for o sexo, estimamos primeiro o efeito causal para homens e para mulheres. Depois fazemos a média dos dois, se a população for (como de costume) metade masculina e metade feminina. Se as proporções forem diferentes — digamos, dois terços de homens e um terço de mulheres —, então, para estimar o efeito causal médio, tomaríamos uma média ponderada correspondente. O papel que o critério da porta dos fundos desempenha nesse procedimento é garantir que o efeito causal em cada estrato do desconfundidor não seja outra coisa senão a tendência observada nesse estrato. Assim, o efeito causal pode ser estimado estrato a estrato a partir dos dados. Na ausência do critério da porta dos fundos, os pesquisadores não têm garantia de que qualquer ajustamento seja legítimo.

O exemplo fictício do medicamento no Capítulo 6 era a situação mais simples possível: uma variável de tratamento (Medicamento D), um desfecho (Infarto), um confundidor (Sexo), e as três variáveis binárias. O exemplo mostra como tomamos uma média ponderada das probabilidades condicionais \textit{P}(infarto | medicamento) em cada estrato de sexo. Mas o procedimento descrito acima pode ser facilmente adaptado a situações mais complicadas, incluindo múltiplos (des)confundidores e múltiplos estratos. Contudo, em muitos casos as variáveis \textit{X}, \textit{Y} ou \textit{Z} assumem valores numéricos — por exemplo, renda, altura ou peso ao nascer. Vimos isso no nosso exemplo visual do paradoxo de Simpson. Como a variável pode assumir (pelo menos, para todos os fins práticos) infinitos valores possíveis, não podemos montar uma tabela listando todas as possibilidades, como fizemos no Capítulo 6. Um remédio óbvio é separar os valores numéricos em um número finito e manejável de categorias. Em princípio, não há nada de errado com essa opção, mas a escolha das categorias é um tanto arbitrária. Pior: se tivermos mais do que um punhado de variáveis ajustadas, obtemos uma explosão exponencial no número de categorias. Isso tornará o procedimento computacionalmente proibitivo; pior ainda, muitos dos estratos acabarão desprovidos de amostras e, portanto, incapazes de fornecer quaisquer estimativas de probabilidade.

Os estatísticos inventaram métodos engenhosos para lidar com esse problema da “maldição da dimensionalidade”. A maioria envolve algum tipo de extrapolação, em que uma função suave é ajustada aos dados e usada para preencher os buracos criados pelos estratos vazios. A função de suavização mais amplamente usada é, naturalmente, a aproximação linear, que serviu de cavalo de batalha da maior parte do trabalho quantitativo nas ciências sociais e comportamentais no século XX. Vimos como Sewall Wright embutiu seus diagramas de caminhos no contexto de equações lineares, e notamos ali uma vantagem computacional dessa incorporação: todo efeito causal pode ser representado por um único número (o coeficiente de caminho). Uma segunda e não menos importante vantagem das aproximações lineares é a assombrosa simplicidade de calcular a fórmula de ajustamento. Já vimos a invenção, por Francis Galton, da reta de regressão, que toma uma nuvem de pontos de dados e interpola a reta de melhor ajuste por essa nuvem. No caso de uma variável de tratamento (\textit{X}) e uma variável de desfecho (\textit{Y}), a equação da reta de regressão terá a seguinte forma: \textit{Y} = \textit{aX} + \textit{b}. O parâmetro \textit{a} (muitas vezes denotado por \textit{r<sub>YX</sub>}, o coeficiente de regressão de \textit{Y} sobre \textit{X}) nos diz a tendência média observada: um aumento de uma unidade em \textit{X} produzirá, em média, um aumento de \textit{a} unidades em \textit{Y}. Se não houver confundidores de \textit{Y} e \textit{X}, podemos usar isso como nossa estimativa de uma intervenção para aumentar \textit{X} em uma unidade. Mas e se houver um confundidor, \textit{Z}? Nesse caso, o coeficiente de correlação \textit{r<sub>YX</sub>} não nos dará o efeito causal médio; dará apenas a tendência média observada. Esse foi o caso do problema de Wright sobre os pesos ao nascer de cobaias, discutido no Capítulo 2, em que o aparente benefício (5,66 gramas) de um dia extra de gestação estava enviesado porque era confundido com o efeito de uma ninhada menor.

Ainda assim, há uma saída: plotar as três variáveis juntas, com cada valor de (\textit{X}, \textit{Y}, \textit{Z}) descrevendo um ponto no espaço. Nesse caso, os dados formarão uma nuvem de pontos no espaço \textit{XYZ}. O análogo de uma reta de regressão é um plano de regressão, cuja equação tem a forma \textit{Y} = \textit{aX} + \textit{bZ} + \textit{c}. Podemos calcular facilmente \textit{a}, \textit{b}, \textit{c} a partir dos dados. Aqui acontece algo maravilhoso, que Galton não percebeu, mas Karl Pearson e George Udny Yule certamente perceberam. O coeficiente \textit{a} nos dá o coeficiente de regressão de \textit{Y} sobre \textit{X} já ajustado por \textit{Z}. (Chama-se coeficiente de regressão parcial e escreve-se \textit{r<sub>YX.Z</sub>}.)

Assim, podemos pular o procedimento trabalhoso de regredir \textit{Y} sobre \textit{X} para cada nível de \textit{Z} e calcular a média ponderada dos coeficientes de regressão. A natureza já faz toda a média por nós! Precisamos apenas calcular o plano que melhor se ajusta aos dados. Um pacote estatístico fará isso em um instante. O coeficiente \textit{a} na equação desse plano, \textit{Y} = \textit{aX} + \textit{bZ} + \textit{c}, ajustará automaticamente a tendência observada de \textit{Y} sobre \textit{X} para dar conta do confundidor \textit{Z}. Se \textit{Z} for o único confundidor, então \textit{a} é o efeito causal médio de \textit{X} sobre \textit{Y}. Uma simplificação verdadeiramente milagrosa! Você pode facilmente estender o procedimento para lidar também com múltiplas variáveis. Se o conjunto de variáveis \textit{Z} satisfizer a condição da porta dos fundos, então o coeficiente de \textit{X} na equação de regressão, \textit{a}, não será outra coisa senão o efeito causal médio de \textit{X} sobre \textit{Y}.

Por essa razão, gerações de pesquisadores passaram a acreditar que os coeficientes de regressão ajustados (ou parciais) estão de algum modo dotados de informação causal de que os coeficientes de regressão não ajustados carecem. Nada poderia estar mais longe da verdade. Coeficientes de regressão, ajustados ou não, são apenas tendências estatísticas, sem transmitir por si mesmos qualquer informação causal. \textit{r<sub>YX.Z</sub>} representa o efeito causal de \textit{X} sobre \textit{Y}, ao passo que \textit{r<sub>YX</sub>} não o representa, exclusivamente porque temos um diagrama mostrando \textit{Z} como confundidor de \textit{X} e \textit{Y}. Em suma, às vezes um coeficiente de regressão representa um efeito causal, e às vezes não — e você não pode confiar apenas nos dados para lhe dizer a diferença.

Dois ingredientes adicionais são necessários para dotar \textit{r<sub>YX.Z</sub>} de legitimidade causal. Primeiro, o diagrama de caminhos deve representar um quadro plausível da realidade; segundo, a(s) variável(is) ajustada(s) \textit{Z} deve(m) satisfazer o critério da porta dos fundos. Foi por isso que foi tão crucial que Sewall Wright distinguisse os coeficientes de caminho (que representam efeitos causais) dos coeficientes de regressão (que representam tendências de pontos de dados). Os coeficientes de caminho são fundamentalmente diferentes dos coeficientes de regressão, embora frequentemente possam ser calculados a partir destes. Wright não percebeu, porém — como tampouco o perceberam todos os analistas de caminhos e econometristas depois dele — que seus cálculos eram desnecessariamente complicados. Ele poderia ter obtido os coeficientes de caminho a partir de coeficientes de correlação parcial, se apenas soubesse que o conjunto adequado de variáveis de ajustamento pode ser identificado, por inspeção, a partir do próprio diagrama de caminhos.

Tenha em mente também que o ajustamento baseado em regressão funciona apenas para modelos lineares, o que envolve uma grande hipótese de modelagem. Com modelos lineares, perdemos a capacidade de modelar interações não lineares, como quando o efeito de \textit{X} sobre \textit{Y} depende do nível de \textit{Z}. O ajustamento pela porta dos fundos, por outro lado, ainda funciona bem mesmo quando não fazemos ideia de quais funções estão por trás das setas nos diagramas. Mas, nesse caso chamado não paramétrico, precisamos empregar outros métodos de extrapolação para lidar com a maldição da dimensionalidade.

Em resumo, a fórmula de ajustamento pela porta dos fundos e o critério da porta dos fundos são como as duas faces de uma moeda. O critério da porta dos fundos nos diz quais conjuntos de variáveis podemos usar para desconfundir nossos dados. A fórmula de ajustamento realiza de fato o desconfundimento. No caso mais simples de regressão linear, os coeficientes de regressão parciais realizam o ajustamento pela porta dos fundos de forma implícita. No caso não paramétrico, devemos fazer o ajustamento explicitamente, seja usando a fórmula de ajustamento pela porta dos fundos diretamente sobre os dados, seja sobre alguma versão extrapolada deles.

Você poderia pensar que nosso assalto ao Monte Intervenção terminaria aí com sucesso completo. Infelizmente, porém, o ajustamento não funciona de modo algum se houver um caminho pela porta dos fundos que não possamos bloquear porque não dispomos dos dados necessários. Ainda assim, podemos usar certos truques mesmo nessa situação. Contarei a seguir sobre um dos meus métodos favoritos, chamado ajustamento pela porta da frente. Embora tenha sido publicado há mais de vinte anos, apenas um punhado de pesquisadores aproveitou esse atalho de subida do Monte Intervenção, e estou convencido de que o seu pleno potencial permanece inexplorado.

\section{O critério da porta da frente}
O debate sobre o efeito causal do tabagismo ocorreu pelo menos duas gerações cedo demais para que os diagramas causais pudessem dar qualquer contribuição. Já vimos como a desigualdade de Cornfield ajudou a persuadir os pesquisadores de que o gene do fumo, ou a “hipótese constitucional”, era altamente implausível. Mas uma abordagem mais radical, usando diagramas causais, poderia ter lançado mais luz sobre o gene hipotético e possivelmente eliminá-lo de consideração ulterior.

Suponha que os pesquisadores tivessem medido os depósitos de alcatrão nos pulmões dos fumantes. Já na década de 1950, a formação de depósitos de alcatrão era suspeita como um dos possíveis estágios intermediários no desenvolvimento do câncer de pulmão. Suponha também que, tal como o comitê do \textit{Surgeon General}, queiramos descartar a hipótese de R. A. Fisher de que um gene do fumo confunde o comportamento de fumar e o câncer de pulmão. Poderíamos então chegar ao diagrama causal da Figura 7.1.

\begin{figure}[ht]
\centering
\includegraphics[width=0.92\textwidth]{../figuras/fig_7_1.png}
\caption{Figura 7.1. A Figura 7.1 incorpora duas hipóteses muito importantes, que suporemos válidas para os fins do nosso exemplo. A primeira hipótese é que o gene do fumo não tem efeito sobre a formação de depósitos de alcatrão, que se devem exclusivamente à ação física da fumaça do cigarro. (Essa hipótese é indicada pela ausência de uma seta entre Gene do Fumo e Alcatrão; não exclui, contudo, fatores aleatórios não relacionados ao Gene do Fumo.) A segunda hipótese significativa é que Fumar leva a Câncer apenas por meio do acúmulo de depósitos de alcatrão. Assim, assumimos que nenhuma seta direta aponta de Fumar para Câncer, e que não há outros caminhos indiretos.}
\end{figure}
\textit{Figura 7.1. A Figura 7.1 incorpora duas hipóteses muito importantes, que suporemos válidas para os fins do nosso exemplo. A primeira hipótese é que o gene do fumo não tem efeito sobre a formação de depósitos de alcatrão, que se devem exclusivamente à ação física da fumaça do cigarro. (Essa hipótese é indicada pela ausência de uma seta entre Gene do Fumo e Alcatrão; não exclui, contudo, fatores aleatórios não relacionados ao Gene do Fumo.) A segunda hipótese significativa é que Fumar leva a Câncer apenas por meio do acúmulo de depósitos de alcatrão. Assim, assumimos que nenhuma seta direta aponta de Fumar para Câncer, e que não há outros caminhos indiretos.}

A Figura 7.1 incorpora duas hipóteses muito importantes, que suporemos válidas para os fins do nosso exemplo. A primeira hipótese é que o gene do fumo não tem efeito sobre a formação de depósitos de alcatrão, que se devem exclusivamente à ação física da fumaça do cigarro. (Essa hipótese é indicada pela ausência de uma seta entre Gene do Fumo e Alcatrão; não exclui, contudo, fatores aleatórios não relacionados ao Gene do Fumo.) A segunda hipótese significativa é que Fumar leva a Câncer apenas por meio do acúmulo de depósitos de alcatrão. Assim, assumimos que nenhuma seta direta aponta de Fumar para Câncer, e que não há outros caminhos indiretos.

\textbf{FIGURA 7.1.} Diagrama causal hipotético para tabagismo e câncer, adequado ao ajustamento pela porta da frente.

Suponha que estejamos fazendo um estudo observacional e tenhamos coletado dados sobre Fumar, Alcatrão e Câncer para cada um dos participantes. Infelizmente, não podemos coletar dados sobre o Gene do Fumo porque não sabemos se tal gene existe. Na ausência de dados sobre a variável de confundimento, não podemos bloquear o caminho pela porta dos fundos Fumar ← Gene do Fumo → Câncer. Assim, não podemos usar o ajustamento pela porta dos fundos para controlar o efeito do confundidor. Precisamos, portanto, buscar outro caminho.

Em vez de entrar pela porta dos fundos, podemos entrar pela porta da frente! Nesse caso, a porta da frente é o caminho causal direto Fumar → Alcatrão → Câncer, para o qual dispomos de dados sobre as três variáveis. Intuitivamente, o raciocínio é o seguinte. Primeiro, podemos estimar o efeito causal médio de Fumar sobre Alcatrão, porque não há caminho pela porta dos fundos desbloqueado de Fumar para Alcatrão, pois o caminho Fumar ← Gene do Fumo → Câncer ← Alcatrão já está bloqueado pelo colisor em Câncer. Como já está bloqueado, nem sequer precisamos de ajustamento pela porta dos fundos. Podemos simplesmente observar \textit{P}(alcatrão | fumar) e \textit{P}(alcatrão | não fumar), e a diferença entre elas será o efeito causal médio de Fumar sobre Alcatrão.

Do mesmo modo, o diagrama nos permite estimar o efeito causal médio de Alcatrão sobre Câncer. Para isso, podemos bloquear o caminho pela porta dos fundos de Alcatrão para Câncer, Alcatrão ← Fumar ← Gene do Fumo → Câncer, ajustando por Fumar. As lições do Capítulo 4 vêm a calhar: precisamos apenas de dados sobre um conjunto suficiente de desconfundidores (isto é, Fumar). Então a fórmula de ajustamento pela porta dos fundos nos dará \textit{P}(câncer | \textit{do}(alcatrão)) e \textit{P}(câncer | \textit{do}(sem alcatrão)). A diferença entre elas é o efeito causal médio de Alcatrão sobre Câncer.

Agora sabemos o aumento médio na verossimilhança de depósitos de alcatrão devido ao fumo e o aumento médio de câncer devido aos depósitos de alcatrão. Podemos combinar esses de algum modo para obter o aumento médio de câncer devido ao fumo? Sim, podemos. O raciocínio é o seguinte. O câncer pode surgir de duas maneiras: na presença de Alcatrão ou na ausência de Alcatrão. Se forçarmos uma pessoa a fumar, então as probabilidades desses dois estados são \textit{P}(alcatrão | \textit{do}(fumar)) e \textit{P}(sem alcatrão | \textit{do}(não fumar)), respectivamente. Se evolui um estado de Alcatrão, a verossimilhança de causar Câncer é \textit{P}(câncer | \textit{do}(alcatrão)). Se, por outro lado, evolui um estado Sem Alcatrão, isso resultaria em uma verossimilhança de Câncer de \textit{P}(câncer | \textit{do}(sem alcatrão)). Podemos ponderar os dois cenários por suas respectivas probabilidades sob \textit{do}(fumar) e, desse modo, calcular a probabilidade total de câncer devida ao fumo. O mesmo argumento vale se impedirmos uma pessoa de fumar, \textit{do}(não fumar). A diferença entre os dois nos dá o efeito causal médio sobre o câncer de fumar versus não fumar.

Como acabei de explicar, podemos estimar cada uma das \textit{do}-probabilidades discutidas a partir dos dados. Isto é, podemos escrevê-las matematicamente em termos de probabilidades que não envolvem o operador \textit{do}. Dessa forma, a matemática faz por nós o que dez anos de debate e depoimentos no Congresso não puderam: quantificar o efeito causal do tabagismo sobre o câncer — desde que nossas hipóteses se sustentem, é claro.

O processo que acabei de descrever, expressando \textit{P}(câncer | \textit{do}(fumar)) em termos de probabilidades livres de \textit{do}, chama-se ajustamento pela porta da frente. Difere do ajustamento pela porta dos fundos em que ajustamos por duas variáveis (Fumar e Alcatrão) em vez de uma, e essas variáveis situam-se no caminho pela porta da frente de Fumar para Câncer, e não no caminho pela porta dos fundos.

Para os leitores que “falam matemática”, não resisto a mostrar a fórmula (Equação 7.1), que não se encontra nos livros-texto ordinários de estatística. Aqui \textit{X} representa Fumar, \textit{Y} representa Câncer, \textit{Z} representa Alcatrão, e \textit{U} (que está conspicuamente ausente da fórmula) representa a variável não observável, o Gene do Fumo.

\$\$P(Y \textbackslash\{\}mid do(X)) = \textbackslash\{\}sum\_z P(Z = z \textbackslash\{\}mid X) \textbackslash\{\}sum\_x P(Y \textbackslash\{\}mid X = x, Z = z)\textbackslash\{\}, P(X = x) \textbackslash\{\}tag\{7.1\}\$\$

Leitores com apetite para a matemática podem achar interessante comparar isso com a fórmula do ajustamento pela porta dos fundos, que se parece com a Equação 7.2.

\$\$P(Y \textbackslash\{\}mid do(X)) = \textbackslash\{\}sum\_z P(Y \textbackslash\{\}mid X, Z = z)\textbackslash\{\}, P(Z = z) \textbackslash\{\}tag\{7.2\}\$\$

Mesmo para leitores que não falam matemática, podemos fazer vários pontos interessantes sobre a Equação 7.1. Primeiro e mais importante, você não vê \textit{U} (o Gene do Fumo) em lugar nenhum. Esse era o objetivo inteiro. Desconfundimos \textit{U} com sucesso mesmo sem possuir quaisquer dados sobre ela. Qualquer estatístico da geração de Fisher teria visto isso como um puro milagre. Segundo, bem no início da Introdução falei de um \textit{estimando} como uma receita para calcular a quantidade de interesse em uma consulta. As Equações 7.1 e 7.2 são os estimandos mais complicados e interessantes que mostrarei neste livro. O lado esquerdo representa a consulta “Qual é o efeito de \textit{X} sobre \textit{Y}?”. O lado direito é o estimando, uma receita para responder à consulta. Note que o estimando não contém \textit{do}’s, apenas \textit{see}’s, representados pelas barras verticais, e isso significa que pode ser estimado a partir dos dados.

Neste ponto, tenho certeza de que alguns leitores se perguntam quão próximo esse cenário fictício está da realidade. A controvérsia tabagismo-câncer poderia ter sido resolvida por um único estudo observacional e um único diagrama causal? Se assumirmos que a Figura 7.1 reflete com precisão o mecanismo causal do câncer, a resposta é absolutamente sim. Contudo, precisamos agora discutir se nossas hipóteses são válidas no mundo real.

David Freedman, amigo de longa data e estatístico de Berkeley, me cobrou sobre essa questão. Ele argumentou que o modelo da Figura 7.1 é irrealista de três maneiras. Primeiro, se existe um gene do fumo, ele também poderia afetar a maneira como o corpo se livra de matéria estranha nos pulmões, de modo que pessoas com o gene seriam mais vulneráveis à formação de depósitos de alcatrão e pessoas sem ele mais resistentes. Portanto, ele desenharia uma seta de Gene do Fumo para Alcatrão, e nesse caso a fórmula da porta da frente seria inválida. Freedman também considerava improvável que Fumar afetasse Câncer apenas por meio de Alcatrão. Certamente outros mecanismos poderiam ser imaginados; talvez o fumo produza inflamação crônica que leve ao câncer. Por fim, disse ele, os depósitos de alcatrão nos pulmões de uma pessoa viva não podem ser medidos com precisão suficiente de qualquer modo — de modo que um estudo observacional como o que propus não pode ser conduzido no mundo real.

Não tenho contenda com a crítica de Freedman neste exemplo particular. Não sou especialista em câncer, e sempre teria de deferir à opinião de especialistas sobre se tal diagrama representa com precisão os processos do mundo real. De fato, uma das grandes conquistas dos diagramas causais é tornar as hipóteses transparentes, de modo que possam ser discutidas e debatidas por especialistas e formuladores de políticas.

Contudo, o ponto do meu exemplo não era propor um novo mecanismo para o efeito do tabagismo, mas demonstrar como a matemática, dada a situação certa, pode eliminar o efeito de confundidores mesmo sem dados sobre o confundidor. E a situação pode ser claramente reconhecida. Sempre que o efeito causal de \textit{X} sobre \textit{Y} é confundido por um conjunto de variáveis (\textit{C}) e mediado por outro (\textit{M}) (veja a Figura 7.2) e, além disso, as variáveis mediadoras estão blindadas contra os efeitos de \textit{C}, então você pode estimar o efeito de \textit{X} a partir de dados observacionais. Uma vez que os cientistas tomem consciência desse fato, devem buscar mediadores blindados sempre que enfrentarem confundidores incuráveis. Como disse Louis Pasteur: “A fortuna favorece a mente preparada.”

\begin{figure}[ht]
\centering
\includegraphics[width=0.92\textwidth]{../figuras/fig_7_2.png}
\caption{Figura 7.2. ) e, além disso, as variáveis mediadoras estão blindadas contra os efeitos de \textit{C}, então você pode estimar o efeito de \textit{X} a partir de dados observacionais. Uma vez que os cientistas tomem consciência desse fato, devem buscar mediadores blindados sempre que enfrentarem confundidores incuráveis. Como disse Louis Pasteur: “A fortuna favorece a mente preparada.”}
\end{figure}
\textit{Figura 7.2. ) e, além disso, as variáveis mediadoras estão blindadas contra os efeitos de }C\textit{, então você pode estimar o efeito de }X\textit{ a partir de dados observacionais. Uma vez que os cientistas tomem consciência desse fato, devem buscar mediadores blindados sempre que enfrentarem confundidores incuráveis. Como disse Louis Pasteur: “A fortuna favorece a mente preparada.”}

Felizmente, as virtudes do ajustamento pela porta da frente não permaneceram completamente desapercebidas. Em 2014, Adam Glynn e Konstantin Kashin, ambos cientistas políticos de Harvard (Glynn subsequentemente mudou-se para a Emory University), escreveram um artigo premiado que deveria ser leitura obrigatória para todos os cientistas sociais quantitativos. Eles aplicaram o novo método a um conjunto de dados bem escrutinado por cientistas sociais, chamado Estudo da Job Training Partnership Act (JTPA), conduzido de 1987 a 1989. Como resultado da JTPA de 1982, o Departamento do Trabalho criou um programa de treinamento profissional que, entre outros serviços, fornecia aos participantes habilidades ocupacionais, habilidades de busca de emprego e experiência de trabalho. Coletou dados sobre pessoas que se candidataram ao programa, pessoas que de fato usaram os serviços, e seus rendimentos ao longo dos dezoito meses subsequentes. Notavelmente, o estudo incluiu tanto um ensaio clínico randomizado (ECR), em que as pessoas eram aleatoriamente designadas a receber os serviços ou não, quanto um estudo observacional, em que as pessoas podiam escolher por si mesmas.

\textbf{FIGURA 7.2.} A configuração básica para o critério da porta da frente.

Glynn e Kashin não desenharam um diagrama causal, mas, a partir de sua descrição do estudo, eu o desenharia como mostrado na Figura 7.3. A variável Inscreveu-se registra se uma pessoa se registrou ou não no programa; a variável Compareceu registra se o inscrito de fato usou ou não os serviços. Obviamente o programa só pode afetar os rendimentos se o usuário de fato comparecer, de modo que a ausência de uma seta direta de Inscreveu-se para Rendimentos é fácil de justificar. Glynn e Kashin se abstêm de especificar a natureza dos confundidores, mas eu os resumi como Motivação. Claramente, uma pessoa altamente motivada a aumentar seus rendimentos tem maior probabilidade de se inscrever. Essa pessoa também tem maior probabilidade de ganhar mais após dezoito meses, independentemente de comparecer ou não. O objetivo do estudo é, naturalmente, desentranhar o efeito desse fator de confundimento e descobrir exatamente o quanto os próprios serviços estão ajudando.

\begin{figure}[ht]
\centering
\includegraphics[width=0.92\textwidth]{../figuras/fig_7_3.png}
\caption{Figura 7.3. A variável Inscreveu-se registra se uma pessoa se registrou ou não no programa; a variável Compareceu registra se o inscrito de fato usou ou não os serviços. Obviamente o programa só pode afetar os rendimentos se o usuário de fato comparecer, de modo que a ausência de uma seta direta de Inscreveu-se para Rendimentos é fácil de justificar. Glynn e Kashin se abstêm de especificar a natureza dos confundidores, mas eu os resumi como Motivação. Claramente, uma pessoa altamente motivada a aumentar seus rendimentos tem maior probabilidade de se inscrever. Essa pessoa também tem maior probabilidade de ganhar mais após dezoito meses, independentemente de comparecer ou não. O objetivo do estudo é, naturalmente, desentranhar o efeito desse fator de confundimento e descobrir exatamente o quanto os próprios serviços estão ajudando.}
\end{figure}
\textit{Figura 7.3. A variável Inscreveu-se registra se uma pessoa se registrou ou não no programa; a variável Compareceu registra se o inscrito de fato usou ou não os serviços. Obviamente o programa só pode afetar os rendimentos se o usuário de fato comparecer, de modo que a ausência de uma seta direta de Inscreveu-se para Rendimentos é fácil de justificar. Glynn e Kashin se abstêm de especificar a natureza dos confundidores, mas eu os resumi como Motivação. Claramente, uma pessoa altamente motivada a aumentar seus rendimentos tem maior probabilidade de se inscrever. Essa pessoa também tem maior probabilidade de ganhar mais após dezoito meses, independentemente de comparecer ou não. O objetivo do estudo é, naturalmente, desentranhar o efeito desse fator de confundimento e descobrir exatamente o quanto os próprios serviços estão ajudando.}

\textbf{FIGURA 7.3.} Diagrama causal para o Estudo JTPA.

Comparando a Figura 7.2 com a Figura 7.3, vemos que o critério da porta da frente se aplicaria se não houvesse seta de Motivação para Compareceu — a “blindagem” que mencionei antes. Em muitos casos, poderíamos justificar a ausência dessa seta. Por exemplo, se os serviços fossem oferecidos apenas com hora marcada e as pessoas só perdessem os compromissos por eventos fortuitos não relacionados à Motivação (uma greve de ônibus, um tornozelo torcido etc.), então poderíamos apagar essa seta e usar o critério da porta da frente. Nas circunstâncias reais do estudo, em que os serviços estavam disponíveis o tempo todo, tal argumento é difícil de sustentar. Contudo — e é aqui que as coisas ficam realmente interessantes — Glynn e Kashin testaram o critério da porta da frente mesmo assim. Podemos pensar nisso como um teste de sensibilidade. Se suspeitamos que a seta do meio é fraca, então o viés introduzido ao tratá-la como ausente pode ser muito pequeno. A julgar pelos resultados, esse foi o caso. Fazendo certas hipóteses razoáveis, Glynn e Kashin derivaram desigualdades dizendo se o ajustamento provavelmente seria alto ou baixo demais e em quanto. Por fim, compararam as previsões da porta da frente e da porta dos fundos com os resultados do experimento controlado randomizado que foi conduzido ao mesmo tempo.

Os resultados foram impressionantes. As estimativas do critério da porta dos fundos (controlando por confundidores conhecidos como Idade, Raça e Local) estavam wildly incorretas, diferindo dos benchmarks experimentais em centenas ou milhares de dólares. É exatamente o que se esperaria ver se houver um confundidor não observado, como Motivação. O critério da porta dos fundos não pode ajustar por ele. Por outro lado, as estimativas da porta da frente conseguiram remover quase todo o efeito da Motivação. Para os homens, as estimativas da porta da frente estavam bem dentro do erro experimental do ensaio clínico randomizado, mesmo com o pequeno viés positivo que Glynn e Kashin previram. Para as mulheres, os resultados foram ainda melhores: as estimativas da porta da frente corresponderam ao benchmark experimental quase perfeitamente, sem viés aparente. O trabalho de Glynn e Kashin fornece prova tanto empírica quanto metodológica de que, desde que o efeito de \textit{C} sobre \textit{M} (na Figura 7.2) seja fraco, o ajustamento pela porta da frente pode dar uma estimativa razoavelmente boa do efeito de \textit{X} sobre \textit{Y}. É muito melhor do que não controlar por \textit{C}.

Os resultados de Glynn e Kashin mostram por que o ajustamento pela porta da frente é uma ferramenta tão poderosa: permite-nos controlar por confundidores que não podemos observar (como Motivação), inclusive aqueles que sequer podemos nomear. Os ECRs são considerados o “padrão-ouro” da estimação de efeitos causais exatamente pela mesma razão. Como as estimativas pela porta da frente fazem a mesma coisa, com a virtude adicional de observar o comportamento das pessoas em seu próprio habitat natural em vez de em um laboratório, eu não me surpreenderia se esse método acabasse se tornando um sério concorrente dos ensaios clínicos randomizados.

\section{O cálculo do \textit{do}, ou a mente sobre a matéria}
Tanto na fórmula de ajustamento pela porta da frente quanto na da porta dos fundos, o objetivo último é calcular o efeito de uma intervenção, \textit{P}(\textit{Y} | \textit{do}(\textit{X})), em termos de dados como \textit{P}(\textit{Y} | \textit{X}, \textit{A}, \textit{B}, \textit{Z}, …) que não envolvem um operador \textit{do}. Se formos completamente bem-sucedidos em eliminar os \textit{do}’s, então podemos usar dados observacionais para estimar o efeito causal, permitindo-nos saltar do degrau um para o degrau dois da Escada da Causalidade.

O fato de termos tido sucesso nesses dois casos (porta da frente e porta dos fundos) levanta de imediato a questão de se há outras portas pelas quais possamos eliminar todos os \textit{do}’s. Pensando de modo mais geral, podemos perguntar se há alguma maneira de decidir de antemão se um dado modelo causal se presta a tal procedimento de eliminação. Se sim, podemos aplicar o procedimento e nos vermos na posse do efeito causal, sem precisar levantar um dedo para intervir. Caso contrário, ao menos saberíamos que as hipóteses embutidas no modelo não são suficientes para revelar o efeito causal a partir de dados observacionais e que, por mais espertos que sejamos, não há escape de conduzir algum tipo de experimento intervencionista.

A perspectiva de fazer essas determinações por meios puramente matemáticos deveria deslumbrar qualquer um que compreenda o custo e a dificuldade de conduzir ensaios clínicos randomizados, mesmo quando fisicamente viáveis e legalmente permissíveis. A ideia também me deslumbrou, no início dos anos 1990 — não como experimentador, mas como cientista da computação e filósofo de meio período. Certamente uma das experiências mais exultantes que se pode ter como cientista é sentar-se à escrivaninha e perceber que finalmente se pode descobrir o que é possível ou impossível no mundo real — especialmente se o problema for importante para a sociedade e tiver desnorteado aqueles que tentaram resolvê-lo antes de você.

Imagino que foi assim que Hiparco de Niceia se sentiu quando descobriu que podia calcular a altura de uma pirâmide a partir de sua sombra no chão, sem de fato escalar a pirâmide. Foi uma clara vitória da mente sobre a matéria. De fato, a abordagem que adotei foi muito inspirada nos gregos antigos (incluindo Hiparco) e em sua invenção de um sistema lógico formal para a geometria. No centro da lógica dos gregos, encontramos um conjunto de axiomas ou verdades autoevidentes, como “Entre quaisquer dois pontos pode-se traçar uma e somente uma reta”. Com a ajuda desses axiomas, os gregos podiam construir enunciados complexos, chamados teoremas, cuja verdade está longe de ser evidente. Tome, por exemplo, o enunciado de que a soma dos ângulos de um triângulo é 180 graus (ou dois ângulos retos), independentemente de seu tamanho ou forma. A verdade desse enunciado não é autoevidente de modo algum; ainda assim, os filósofos pitagóricos do século V a.C. foram capazes de provar a sua verdade universal usando esses axiomas autoevidentes como blocos de construção.

Se você se lembra da geometria do ensino médio, nem que seja só a essência, recordará que as provas de teoremas invariavelmente consistem em construções auxiliares: por exemplo, traçar uma reta paralela a uma aresta de um triângulo, marcar certos ângulos como iguais, desenhar um círculo com um dado segmento como raio, e assim por diante. Essas construções auxiliares podem ser vistas como sentenças matemáticas temporárias que fazem asserções (ou afirmações) sobre propriedades da figura desenhada. Cada nova construção é licenciada pelas anteriores, bem como pelos axiomas da geometria e talvez por alguns teoremas já derivados. Por exemplo, traçar uma reta paralela a uma aresta de um triângulo é licenciado pelo quinto axioma de Euclides, de que é possível traçar uma e somente uma paralela a uma reta dada a partir de um ponto exterior a essa reta. O ato de desenhar qualquer uma dessas construções auxiliares é apenas uma operação mecânica de “manipulação de símbolos”; toma a sentença previamente escrita (ou a figura previamente desenhada) e a reescreve em um novo formato, sempre que a reescrita for licenciada pelos axiomas. A grandeza de Euclides foi identificar uma curta lista de cinco axiomas elementares, a partir dos quais todos os outros enunciados geométricos verdadeiros podem ser derivados.

Voltemos agora à nossa questão central: quando um modelo pode substituir um experimento, ou quando uma quantidade “\textit{do}” pode ser reduzida a uma quantidade “\textit{see}”. Inspirados pelos geômetras gregos antigos, queremos reduzir o problema à manipulação de símbolos e, desse modo, arrancar a causalidade do Monte Olimpo e torná-la disponível ao pesquisador médio.

Primeiro, reformulemos a tarefa de encontrar o efeito de \textit{X} sobre \textit{Y} usando a linguagem de provas, axiomas e construções auxiliares — a linguagem de Euclides e Pitágoras. Começamos com nossa sentença-alvo, \textit{P}(\textit{Y} | \textit{do}(\textit{X})). Nossa tarefa estará completa se conseguirmos eliminar dela o operador \textit{do}, deixando apenas expressões clássicas de probabilidade, como \textit{P}(\textit{Y} | \textit{X}) ou \textit{P}(\textit{Y} | \textit{X}, \textit{Z}, \textit{W}). Não podemos, é claro, manipular nossa expressão-alvo à vontade; as operações devem se conformar ao que \textit{do}(\textit{X}) significa como uma intervenção física. Assim, devemos passar a expressão por uma sequência de manipulações legítimas, cada uma licenciada pelos axiomas e pelas hipóteses do nosso modelo. As manipulações devem preservar o significado da expressão manipulada, mudando apenas o formato em que é escrita. Um exemplo de transformação “preservadora de significado” é a transformação algébrica que transforma \textit{y} = \textit{ax} + \textit{b} em \textit{ax} = \textit{y} − \textit{b}. A relação entre \textit{x} e \textit{y} permanece intacta; só o formato muda.

Já estamos familiarizados com algumas transformações “legítimas” sobre \textit{do}-expressões. Por exemplo, a Regra 1 diz que, quando observamos uma variável \textit{W} irrelevante para \textit{Y} (possivelmente condicionada a outras variáveis \textit{Z}), a distribuição de probabilidade de \textit{Y} não mudará. Por exemplo, no Capítulo 3 vimos que a variável Fogo é irrelevante para Alarme uma vez que conhecemos o estado do mediador (Fumaça). Essa asserção de irrelevância se traduz em uma manipulação simbólica:

\$\$P(Y \textbackslash\{\}mid do(X), Z, W) = P(Y \textbackslash\{\}mid do(X), Z)\$\$

A equação enunciada vale desde que o conjunto de variáveis \textit{Z} bloqueie todos os caminhos de \textit{W} para \textit{Y} depois de termos deletado todas as setas que apontam para \textit{X}. No exemplo Fogo → Fumaça → Alarme, temos \textit{W} = Fogo, \textit{Z} = Fumaça, \textit{Y} = Alarme, e \textit{Z} bloqueia todos os caminhos de \textit{W} para \textit{Y}. (Nesse caso não temos uma variável \textit{X}.)

Outra transformação legítima nos é familiar da discussão sobre a porta dos fundos. Sabemos que, se um conjunto \textit{Z} de variáveis bloqueia todos os caminhos pela porta dos fundos de \textit{X} para \textit{Y}, então, condicionado a \textit{Z}, \textit{do}(\textit{X}) é equivalente a \textit{see}(\textit{X}). Podemos, portanto, escrever

\$\$P(Y \textbackslash\{\}mid do(X), Z) = P(Y \textbackslash\{\}mid X, Z)\$\$

se \textit{Z} satisfaz o critério da porta dos fundos. Adotamos isso como Regra 2 do nosso sistema axiomático. Embora talvez seja menos autoevidente do que a Regra 1, nos casos mais simples trata-se do princípio da causa comum de Hans Reichenbach, emendado para que não confundamos colisores com confundidores. Em outras palavras, estamos dizendo que, depois de controlarmos por um conjunto desconfundidor suficiente, qualquer correlação remanescente é um efeito causal genuíno.

A Regra 3 é bastante simples: essencialmente diz que podemos remover \textit{do}(\textit{X}) de \textit{P}(\textit{Y} | \textit{do}(\textit{X})) em qualquer caso em que não haja caminhos causais de \textit{X} para \textit{Y}. Isto é,

\$\$P(Y \textbackslash\{\}mid do(X)) = P(Y)\$\$

se não há caminho de \textit{X} para \textit{Y} apenas com setas dirigidas para a frente. Podemos parafrasear essa regra assim: se fazemos algo que não afeta \textit{Y}, então a distribuição de probabilidade de \textit{Y} não mudará.

Além de serem tão autoevidentes quanto os axiomas de Euclides, as Regras 1 a 3 também podem ser provadas matematicamente usando nossa definição do operador \textit{do} por deleção de setas e as leis básicas da probabilidade. Note que as Regras 1 e 2 incluem probabilidades condicionais envolvendo variáveis auxiliares \textit{Z} diferentes de \textit{X} e \textit{Y}. Essas variáveis podem ser pensadas como um contexto em que a probabilidade está sendo calculada. Às vezes a presença desse contexto em si licencia a transformação. A Regra 3 também pode ter variáveis auxiliares, mas eu as omiti por simplicidade.

Note que cada regra tem uma interpretação sintática simples. A Regra 1 permite a adição ou deleção de observações. A Regra 2 permite a substituição de uma intervenção por uma observação, ou vice-versa. A Regra 3 permite a deleção ou adição de intervenções. Todas essas permissões são emitidas sob condições apropriadas, que devem ser verificadas em cada caso particular a partir do diagrama causal.

Estamos prontos agora para demonstrar como as Regras 1 a 3 nos permitem transformar uma fórmula em outra até que, se formos espertos, obtenhamos uma expressão de nosso agrado. Embora seja um pouco elaborado, penso que nada substitui de fato mostrar-lhe como a fórmula da porta da frente é derivada por aplicação sucessiva das regras do cálculo do \textit{do} (Figura 7.4). Você não precisa seguir todos os passos, mas mostro a derivação para lhe dar o sabor do cálculo do \textit{do}. Começamos a jornada com uma expressão-alvo \textit{P}(\textit{Y} | \textit{do}(\textit{X})). Introduzimos variáveis auxiliares e transformamos a expressão-alvo em uma expressão livre de \textit{do} que coincide, naturalmente, com a fórmula de ajustamento pela porta da frente. Cada passo do argumento obtém sua licença do diagrama causal que relaciona \textit{X}, \textit{Y} e as variáveis auxiliares ou, em vários casos, de subdiagramas das quais setas foram apagadas para dar conta de intervenções. Essas licenças são exibidas no lado direito.

\begin{figure}[ht]
\centering
\includegraphics[width=0.92\textwidth]{../figuras/fig_7_4.png}
\caption{Figura 7.4. ). Você não precisa seguir todos os passos, mas mostro a derivação para lhe dar o sabor do cálculo do \textit{do}. Começamos a jornada com uma expressão-alvo \textit{P}(\textit{Y} | \textit{do}(\textit{X})). Introduzimos variáveis auxiliares e transformamos a expressão-alvo em uma expressão livre de \textit{do} que coincide, naturalmente, com a fórmula de ajustamento pela porta da frente. Cada passo do argumento obtém sua licença do diagrama causal que relaciona \textit{X}, \textit{Y} e as variáveis auxiliares ou, em vários casos, de subdiagramas das quais setas foram apagadas para dar conta de intervenções. Essas licenças são exibidas no lado direito.}
\end{figure}
\textit{Figura 7.4. ). Você não precisa seguir todos os passos, mas mostro a derivação para lhe dar o sabor do cálculo do }do\textit{. Começamos a jornada com uma expressão-alvo }P\textit{(}Y\textit{ | }do\textit{(}X\textit{)). Introduzimos variáveis auxiliares e transformamos a expressão-alvo em uma expressão livre de }do\textit{ que coincide, naturalmente, com a fórmula de ajustamento pela porta da frente. Cada passo do argumento obtém sua licença do diagrama causal que relaciona }X\textit{, }Y\textit{ e as variáveis auxiliares ou, em vários casos, de subdiagramas das quais setas foram apagadas para dar conta de intervenções. Essas licenças são exibidas no lado direito.}

\textbf{FIGURA 7.4.} Derivação da fórmula de ajustamento pela porta da frente a partir das regras do cálculo do \textit{do}.

Sinto um apego especial ao cálculo do \textit{do}. Com essas três humildes regras, pude derivar a fórmula da porta da frente. Foi o primeiro efeito causal estimado por meios outros que o controle de confundidores. Acreditava que ninguém poderia fazê-lo sem o cálculo do \textit{do}, de modo que o apresentei como um desafio em um seminário de estatística em Berkeley em 1993 e até ofereci um prêmio de 100 dólares a quem o resolvesse. Paul Holland, que assistiu ao seminário, escreveu que havia designado o problema como projeto de turma e me enviaria a solução quando madura. (Colegas me dizem que ele acabou apresentando uma longa solução em uma conferência em 1995, e posso dever-lhe 100 dólares se eu apenas encontrar a prova dele.) Os economistas James Heckman e Rodrigo Pinto fizeram a próxima tentativa de provar a fórmula da porta da frente usando “ferramentas padrão” em 2015. Tiveram sucesso, embora ao custo de oito páginas de duro labor.

No restaurante, na noite anterior à palestra, eu havia escrito a prova (muito parecida com a da Figura 7.4) em um guardanapo para David Freedman. Ele me escreveu mais tarde dizendo que havia perdido o guardanapo. Não conseguia reconstruir o argumento e perguntou se eu havia guardado uma cópia. No dia seguinte, Jamie Robins me escreveu de Harvard, dizendo que ouvira falar do “problema do guardanapo” por Freedman, e de imediato se ofereceu para voar à Califórnia a fim de checar a prova comigo. Fiquei entusiasmado de compartilhar com Robins os segredos do cálculo do \textit{do}, e creio que a sua viagem a Los Angeles naquele ano foi a chave para a sua aceitação entusiástica dos diagramas causais. Por meio da influência dele e de Sander Greenland, os diagramas tornaram-se uma segunda língua para os epidemiologistas. Isso explica por que sou tão afeiçoado ao “problema do guardanapo”.

A fórmula de ajustamento pela porta da frente foi uma surpresa deliciosa e uma indicação de que o cálculo do \textit{do} tinha algo importante a oferecer. Contudo, nesse ponto eu ainda me perguntava se as três regras do cálculo do \textit{do} eram suficientes. Seria possível que tivéssemos perdido uma quarta regra que nos ajudasse a resolver problemas insolúveis com apenas três?

Em 1994, quando propus pela primeira vez o cálculo do \textit{do}, selecionei essas três regras porque eram suficientes em qualquer caso de que eu tivesse conhecimento. Não fazia ideia se, como o fio de Ariadne, elas sempre me levariam para fora do labirinto, ou se um dia encontraria um labirinto de complexidade tão diabólica que eu não pudesse escapar. É claro que esperava o melhor. Conjecturei que, sempre que um efeito causal for estimável a partir dos dados, uma sequência de passos usando essas três regras eliminaria o operador \textit{do}. Mas não pude provar isso.

Esse tipo de problema tem muitos precedentes em matemática e lógica. A propriedade costuma ser chamada de “completude” em lógica matemática; um sistema de axiomas que é completo tem a propriedade de que os axiomas bastam para derivar todo enunciado verdadeiro naquela linguagem. Alguns sistemas de axiomas muito bons são incompletos: por exemplo, os axiomas de Philip Dawid que descrevem a independência condicional na teoria da probabilidade. Nesse conto moderno de labirinto, dois grupos de pesquisadores desempenharam o papel de Ariadne para o meu Teseu errante: Yiming Huang e Marco Valtorta na University of South Carolina e meu próprio aluno, Ilya Shpitser, na University of California, Los Angeles (UCLA). Ambos os grupos provaram de forma independente e simultânea que as Regras 1 a 3 bastam para sair de qualquer \textit{do}-labirinto que tenha uma saída. Não tenho certeza se o mundo aguardava ofegante o resultado de completude deles, porque a essa altura a maioria dos pesquisadores já se contentava em usar apenas os critérios da porta da frente e da porta dos fundos. Ambas as equipes, contudo, foram reconhecidas com prêmios de melhor artigo de estudante na conferência Uncertainty in Artificial Intelligence em 2006.

Confesso que fui eu quem aguardava ofegante esse resultado. Ele nos diz que, se não conseguirmos encontrar um modo de estimar \textit{P}(\textit{Y} | \textit{do}(\textit{X})) a partir das Regras 1 a 3, então uma solução não existe. Nesse caso, sabemos que não há alternativa a conduzir um ensaio clínico randomizado. Diz-nos ainda quais hipóteses ou experimentos adicionais poderiam tornar o efeito causal estimável.

Antes de declarar vitória total, devemos discutir uma questão com o cálculo do \textit{do}. Como qualquer outro cálculo, ele permite a construção de uma prova, mas não nos ajuda a encontrar uma. É um excelente verificador de uma solução, mas não um tão bom buscador de uma. Se você conhece a sequência correta de transformações, é fácil demonstrar a outros (que estejam familiarizados com as Regras 1 a 3) que o operador \textit{do} pode ser eliminado. Contudo, se você não conhece a sequência correta, não é fácil descobri-la, ou mesmo determinar se uma existe. Usando a analogia com as provas geométricas, precisamos decidir qual construção auxiliar tentar a seguir. Um círculo em torno do ponto A? Uma reta paralela a AB? O número de possibilidades é ilimitado, e os próprios axiomas não fornecem orientação sobre o que tentar a seguir. Meu professor de geometria do ensino médio costumava dizer que se precisam de “óculos matemáticos”.

Em lógica matemática, isso é conhecido como o “problema da decisão”. Muitos sistemas lógicos são atormentados por problemas de decisão intratáveis. Por exemplo, dado um monte de dominós de vários tamanhos, não temos modo tratável de decidir se podemos arranjá-los de modo a preencher um quadrado de tamanho dado. Mas, uma vez que um arranjo é proposto, não leva tempo algum verificar se constitui uma solução. Felizmente (de novo) para o cálculo do \textit{do}, o problema da decisão acaba por ser manejável. Ilya Shpitser, com base em trabalho anterior de um de meus outros alunos, Jin Tian, encontrou um algoritmo que decide se uma solução existe em “tempo polinomial”. Trata-se de um termo um tanto técnico, mas, continuando nossa analogia com a solução de um labirinto, significa que temos um modo muito mais eficiente de sair do labirinto do que caçar ao acaso por todos os caminhos possíveis.

O algoritmo de Shpitser para encontrar cada e todo efeito causal não elimina a necessidade do cálculo do \textit{do}. De fato, precisamos dele ainda mais, e por várias razões independentes. Primeiro, precisamos dele para ir além de estudos observacionais. Suponha que o pior aconteça, e nosso modelo causal não permita a estimação do efeito causal \textit{P}(\textit{Y} | \textit{do}(\textit{X})) a partir apenas de observações. Talvez também não possamos conduzir um experimento randomizado com designação aleatória de \textit{X}. Um pesquisador esperto poderia perguntar se poderíamos estimar \textit{P}(\textit{Y} | \textit{do}(\textit{X})) randomizando alguma outra variável, digamos \textit{Z}, que seja mais acessível ao controle do que \textit{X}. Por exemplo, se queremos avaliar o efeito dos níveis de colesterol (\textit{X}) sobre a doença cardíaca (\textit{Y}), poderíamos ser capazes de manipular a dieta dos sujeitos (\textit{Z}) em vez de exercer controle direto sobre os níveis de colesterol em seu sangue. Perguntamos então se podemos encontrar tal substituto \textit{Z} que nos permita responder à pergunta causal. No mundo do cálculo do \textit{do}, a questão é se podemos encontrar um \textit{Z} tal que possamos transformar \textit{P}(\textit{Y} | \textit{do}(\textit{X})) em uma expressão em que a variável \textit{Z}, mas não \textit{X}, está sujeita a um operador \textit{do}. Esse é um problema completamente diferente, não coberto pelo algoritmo de Shpitser. Felizmente, ele também tem uma resposta completa, com um novo algoritmo descoberto por Elias Bareinboim em meu laboratório em 2012.

Problemas ainda mais desse tipo surgem quando consideramos problemas de transportabilidade ou validade externa — avaliar se um resultado experimental ainda será válido quando transportado para um ambiente diferente que pode diferir em várias maneiras-chave daquele estudado. Esse conjunto mais ambicioso de questões toca o coração da metodologia científica, pois não há ciência sem generalização. Ainda assim, a questão da generalização vinha pairando por pelo menos dois séculos, sem um iota de progresso. As ferramentas para produzir uma solução simplesmente não estavam disponíveis. Em 2015, Bareinboim e eu apresentamos um artigo na National Academy of Sciences que resolve o problema, desde que se possa expressar as hipóteses sobre ambos os ambientes com um diagrama causal. Nesse caso, as regras do cálculo do \textit{do} fornecem um método sistemático para determinar se os efeitos causais encontrados no ambiente do estudo podem nos ajudar a estimar efeitos no ambiente-alvo pretendido.

Ainda outra razão pela qual o cálculo do \textit{do} permanece importante é a transparência. Enquanto eu escrevia este capítulo, Bareinboim (agora professor na Purdue) me enviou um novo quebra-cabeça: um diagrama com apenas quatro variáveis observadas, \textit{X}, \textit{Y}, \textit{Z} e \textit{W}, e duas variáveis não observáveis, \textit{U<sub>1</sub>}, \textit{U<sub>2</sub>} (veja a Figura 7.5). Ele me desafiou a descobrir se o efeito de \textit{X} sobre \textit{Y} era estimável. Não havia modo de bloquear os caminhos pela porta dos fundos e nenhuma condição da porta da frente. Tentei todos os meus atalhos favoritos e meus argumentos intuitivos de outra forma confiáveis, tanto a favor quanto contra, e não conseguia ver como fazê-lo. Não pude encontrar uma saída do labirinto. Mas assim que Bareinboim me sussurrou “Tente o cálculo do \textit{do}”, a resposta veio brilhando como o sorriso de um bebê. Cada passo era claro e significativo. Esse é agora o modelo mais simples conhecido por nós em que o efeito causal precisa ser estimado por um método que vai além dos ajustamentos pela porta da frente e pela porta dos fundos.

\begin{figure}[ht]
\centering
\includegraphics[width=0.92\textwidth]{../figuras/fig_7_5.png}
\caption{Figura 7.5. ). Ele me desafiou a descobrir se o efeito de \textit{X} sobre \textit{Y} era estimável. Não havia modo de bloquear os caminhos pela porta dos fundos e nenhuma condição da porta da frente. Tentei todos os meus atalhos favoritos e meus argumentos intuitivos de outra forma confiáveis, tanto a favor quanto contra, e não conseguia ver como fazê-lo. Não pude encontrar uma saída do labirinto. Mas assim que Bareinboim me sussurrou “Tente o cálculo do \textit{do}”, a resposta veio brilhando como o sorriso de um bebê. Cada passo era claro e significativo. Esse é agora o modelo mais simples conhecido por nós em que o efeito causal precisa ser estimado por um método que vai além dos ajustamentos pela porta da frente e pela porta dos fundos.}
\end{figure}
\textit{Figura 7.5. ). Ele me desafiou a descobrir se o efeito de }X\textit{ sobre }Y\textit{ era estimável. Não havia modo de bloquear os caminhos pela porta dos fundos e nenhuma condição da porta da frente. Tentei todos os meus atalhos favoritos e meus argumentos intuitivos de outra forma confiáveis, tanto a favor quanto contra, e não conseguia ver como fazê-lo. Não pude encontrar uma saída do labirinto. Mas assim que Bareinboim me sussurrou “Tente o cálculo do }do\textit{”, a resposta veio brilhando como o sorriso de um bebê. Cada passo era claro e significativo. Esse é agora o modelo mais simples conhecido por nós em que o efeito causal precisa ser estimado por um método que vai além dos ajustamentos pela porta da frente e pela porta dos fundos.}

\textbf{FIGURA 7.5.} Um novo problema do guardanapo?

A fim de não deixar o leitor com a impressão de que o cálculo do \textit{do} serve apenas para a teoria e como passatempo recreativo para o cérebro, terminarei esta seção com um problema prático recentemente levantado por dois estatísticos de ponta, Nanny Wermuth e David Cox. Ele demonstra como um sussurro amigável — “Tente o cálculo do \textit{do}” — pode ajudar estatísticos experientes a resolver problemas práticos difíceis.

Por volta de 2005, Wermuth e Cox interessaram-se por um problema chamado “decisões sequenciais” ou “tratamentos variantes no tempo”, comuns, por exemplo, no tratamento da AIDS. Tipicamente os tratamentos são administrados ao longo de um período de tempo, e em cada período os médicos variam a força e a dosagem de um tratamento de seguimento de acordo com a condição do paciente. A condição do paciente, por outro lado, é influenciada pelos tratamentos tomados no passado. Assim, acabamos com um cenário como o da Figura 7.6, mostrando dois períodos de tempo e dois tratamentos. O primeiro tratamento é randomizado (\textit{X}), e o segundo (\textit{Z}) é dado em resposta a uma observação (\textit{W}) que depende de \textit{X}. Dados os dados coletados sob tal regime de tratamento, a tarefa de Cox e Wermuth era prever o efeito de \textit{X} sobre o desfecho \textit{Y}, assumindo que manteriam \textit{Z} constante ao longo do tempo, independente da observação \textit{W}.

\begin{figure}[ht]
\centering
\includegraphics[width=0.92\textwidth]{../figuras/fig_7_6.png}
\caption{Figura 7.6. , mostrando dois períodos de tempo e dois tratamentos. O primeiro tratamento é randomizado (\textit{X}), e o segundo (\textit{Z}) é dado em resposta a uma observação (\textit{W}) que depende de \textit{X}. Dados os dados coletados sob tal regime de tratamento, a tarefa de Cox e Wermuth era prever o efeito de \textit{X} sobre o desfecho \textit{Y}, assumindo que manteriam \textit{Z} constante ao longo do tempo, independente da observação \textit{W}.}
\end{figure}
\textit{Figura 7.6. , mostrando dois períodos de tempo e dois tratamentos. O primeiro tratamento é randomizado (}X\textit{), e o segundo (}Z\textit{) é dado em resposta a uma observação (}W\textit{) que depende de }X\textit{. Dados os dados coletados sob tal regime de tratamento, a tarefa de Cox e Wermuth era prever o efeito de }X\textit{ sobre o desfecho }Y\textit{, assumindo que manteriam }Z\textit{ constante ao longo do tempo, independente da observação }W\textit{.}

\textbf{FIGURA 7.6.} Exemplo de Wermuth e Cox de um tratamento sequencial.

Jamie Robins trouxe-me pela primeira vez o problema dos tratamentos variantes no tempo em 1994 e, com a ajuda do cálculo do \textit{do}, pudemos derivar uma solução geral invocando uma versão sequencial da fórmula de ajustamento pela porta dos fundos. Wermuth e Cox, sem conhecimento desse método, chamaram seu problema de “confundimento indireto” e publicaram três artigos sobre sua análise (2008, 2014 e 2015). Incapazes de resolvê-lo em geral, recorreram a uma aproximação linear, e mesmo no caso linear acharam difícil de lidar, porque não é solúvel por métodos padrão de regressão. Felizmente, quando uma musa me sussurrou no ouvido “Tente o cálculo do \textit{do}”, notei que o problema deles pode ser resolvido em três linhas de cálculo. A lógica é a seguinte. Nossa quantidade-alvo é \textit{P}(\textit{Y} | \textit{do}(\textit{X}), \textit{do}(\textit{Z})), enquanto os dados de que dispomos são da forma \textit{P}(\textit{Y} | \textit{do}(\textit{X}), \textit{Z}, \textit{W}) e \textit{P}(\textit{W} | \textit{do}(\textit{X})). Isso reflete o fato de que, no estudo do qual temos dados, \textit{Z} não é controlado externamente, mas segue \textit{W} por algum protocolo (desconhecido). Assim, nossa tarefa é transformar a expressão-alvo em outra expressão, refletindo as condições do estudo em que o operador \textit{do} se aplica apenas a \textit{X} e não a \textit{Z}. Acontece que uma única aplicação das três regras do cálculo do \textit{do} pode realizar isso.

A moral da história não é senão um profundo apreço pelo poder da matemática de resolver problemas difíceis, que ocasionalmente acarretam consequências práticas.

\section{A tapeçaria da ciência, ou os jogadores ocultos na orquestra do \textit{do}}
Já mencionei o papel de alguns de meus alunos na tecelagem desta bela tapeçaria do cálculo do \textit{do}. Como qualquer tapeçaria, ela dá uma sensação de completude que pode ocultar quão penosa foi a sua feitura e quantas mãos contribuíram para o processo. Nesse caso, levou mais de vinte anos e contribuições de vários alunos e colegas.

O primeiro foi Thomas Verma, a quem conheci quando ele era um menino de dezesseis anos. Seu pai o trouxe ao meu escritório um dia e disse, essencialmente: “Dê-lhe algo para fazer.” Ele era talentoso demais para que qualquer de seus professores de matemática do ensino médio o mantivesse interessado. O que ele acabou por realizar foi verdadeiramente assombroso. Verma finalmente provou o que se tornou conhecido como a propriedade de \textit{d}-separação (isto é, o fato de que se pode usar as regras de bloqueio de caminhos para determinar quais independências devem valer nos dados). De modo assombroso, ele me contou que provou a propriedade de \textit{d}-separação pensando que era um problema de lição de casa, não uma conjectura não resolvida! Às vezes vale a pena ser jovem e ingênuo. Você ainda pode ver o legado dele na Regra 1 do cálculo do \textit{do} e em qualquer impressão que o bloqueio de caminhos deixe no degrau um da Escada da Causalidade.

O poder da prova de Verma teria permanecido apenas parcialmente apreciado sem um resultado complementar para mostrar que não pode ser melhorado. Isto é, nenhuma outra independência é implicada por um diagrama causal exceto aquelas reveladas por meio do bloqueio de caminhos. Esse passo foi completado por outro aluno, Dan Geiger. Ele havia se transferido para o meu laboratório de pesquisa de outro grupo na UCLA, depois que eu prometi dar-lhe um “PhD instantâneo” se ele pudesse provar dois teoremas. Ele o fez, e eu o fiz! Ele é agora decano de ciência da computação no Technion em Israel, minha \textit{alma mater}.

Mas Dan não foi o único aluno que “recrutei” de outro departamento. Um dia em 1997, enquanto me vestia no vestiário da piscina da UCLA, entabulei conversa com um sujeito chinês ao meu lado. Ele era um estudante de doutorado em física e, como era meu hábito usual na época, tentei convencê-lo a mudar para a inteligência artificial, onde estava a ação. Ele não ficou completamente convencido, mas no dia seguinte recebi um e-mail de um amigo dele, Jin Tian, dizendo que gostaria de mudar da física para a ciência da computação e se eu tinha um projeto de verão desafiador para ele. Dois dias depois, ele estava trabalhando no meu laboratório. Quatro anos depois, em abril de 2001, ele atordoou o mundo com um critério gráfico simples que generaliza a porta da frente, a porta dos fundos e todas as portas em que pudéssemos pensar na época. Recordo-me de apresentar o critério de Tian em uma conferência em Santa Fe. Um a um, líderes da comunidade de pesquisa olhavam para o meu pôster e balançavam a cabeça em descrença. Como poderia um critério tão simples funcionar para todos os diagramas?

Tian (agora professor na Iowa State University) veio ao nosso laboratório com um estilo de pensamento que nos era estranho então, nos anos 1990. Nossas conversas estavam sempre carregadas de metáforas selvagens e conjecturas a meio cozer. Mas Tian nunca proferia uma palavra a menos que fosse rigorosa, provada e cozida cinco vezes. A mistura dos dois estilos provou o seu mérito. O método de Tian, chamado \textit{c}-decomposição, permitiu a Ilya Shpitser desenvolver o seu algoritmo completo para o cálculo do \textit{do}. A moral: nunca subestime o poder de uma conversa de vestiário!

Ilya Shpitser entrou no final da batalha de dez anos para compreender as intervenções. Ele chegou durante um período muito difícil, em que eu tive de me afastar para estabelecer uma fundação em honra de meu filho, Daniel, vítima do terrorismo antiocidental. Sempre esperei que meus alunos fossem autossuficientes, mas para meus alunos naquela época essa expectativa foi levada ao extremo. Eles me deram o melhor de todos os possíveis presentes, pondo os toques finais, mas cruciais, na tapeçaria do cálculo do \textit{do}, o que eu não poderia ter feito sozinho. De fato, tentei desencorajar Ilya de tentar provar a completude do cálculo do \textit{do}. Provas de completude são notoriamente difíceis e são melhor evitadas por qualquer estudante que vise terminar o doutorado a tempo. Felizmente, Ilya o fez pelas minhas costas.

Colegas também exercem um efeito profundo sobre o seu pensamento em momentos cruciais. Peter Spirtes, professor de filosofia na Carnegie Mellon, precedeu-me na abordagem de redes à causalidade, e a sua influência foi pivotal. Em uma palestra dele em Uppsala, Suécia, aprendi pela primeira vez que realizar intervenções podia ser pensado como deletar setas de um diagrama causal. Até então eu labutava sob o mesmo fardo de gerações de estatísticos, tentando pensar a causalidade em termos de apenas um diagrama representando uma única distribuição de probabilidade estática. A ideia de deleção de setas tampouco era inteiramente de Spirtes. Em 1960, dois economistas suecos, Robert Strotz e Herman Wold, propuseram essencialmente a mesma ideia. No mundo da economia da época, diagramas nunca eram usados; em vez disso, os economistas se fiavam em modelos de equações estruturais, que são as equações de Sewall Wright sem os diagramas. A deleção de setas em um diagrama de caminhos corresponde a deletar uma equação de um modelo de equações estruturais. Assim, em um sentido aproximado, Strotz e Wold tiveram a ideia primeiro, a menos que queiramos ir ainda mais longe na história: eles foram precedidos por Trygve Haavelmo (economista norueguês e laureado com o Nobel), que em 1943 advogou a modificação de equações para representar intervenções.

Não obstante, a tradução de Spirtes da deleção de equações para o mundo dos diagramas causais desencadeou uma avalanche de novas intuições e novos resultados. O critério da porta dos fundos foi um dos primeiros beneficiários da tradução, enquanto o cálculo do \textit{do} veio em segundo. A avalanche, contudo, ainda não acabou. Avanços em áreas como contrafactuais, generalizabilidade, dados faltantes e aprendizado de máquina ainda estão surgindo. Se eu fosse menos modesto, fecharia aqui com o famoso ditado de Isaac Newton sobre “estar sobre os ombros de gigantes”. Mas, dado quem sou, sinto-me tentado a citar a Mishná em vez disso: “\textit{Harbe lamadeti mirabotai um’haverai yoter mehem, umitalmidai yoter mikulam}” — isto é, “Aprendi muito com meus mestres, e mais ainda com meus colegas, e mais do que todos com meus alunos” (Taanit 7a). O operador \textit{do} e o cálculo do \textit{do} não existiriam como existem hoje sem as contribuições de Verma, Geiger, Tian e Shpitser, entre outros.

\section{O(s) curioso(s) caso(s) do Dr. Snow}
Em 1853 e 1854, a Inglaterra estava nas garras de uma epidemia de cólera. Naquela era, a cólera era tão aterrorizante quanto o Ebola é hoje; uma pessoa saudável que bebe água contaminada por cólera pode morrer em vinte e quatro horas. Sabemos hoje que a cólera é causada por uma bactéria que ataca os intestinos. Espalha-se pela diarreia “água de arroz” de suas vítimas, que excretam essa diarreia em quantidades copiosas antes de morrer. Mas em 1853, germes causadores de doença ainda nunca haviam sido vistos sob um microscópio para qualquer enfermidade, quanto mais para a cólera. A sabedoria reinante sustentava que um “miasma” de ar insalubre causava a cólera, uma teoria aparentemente apoiada pelo fato de que a epidemia atingia com mais força as seções mais pobres de Londres, onde o saneamento era pior.

O Dr. John Snow, um médico que cuidara de vítimas de cólera por mais de vinte anos, sempre foi cético em relação à teoria do miasma. Argumentava, sensatamente, que, como os sintomas se manifestavam no trato intestinal, o corpo devia primeiro entrar em contato com o patógeno ali. Mas, como não podia ver o culpado, não tinha como provar isso — até a epidemia de 1854.

A história de John Snow tem dois capítulos, um muito mais famoso do que o outro. Naquilo que poderíamos chamar de versão “Hollywood”, ele vai meticulosamente de casa em casa, registrando onde as vítimas de cólera morreram, e nota um aglomerado de dezenas de vítimas perto de uma bomba d’água na Broad Street. Conversando com as pessoas que moram na área, descobre que quase todas as vítimas haviam tirado a água daquela bomba particular. Ele até fica sabendo de um caso fatal que ocorreu longe dali, em Hampstead, com uma mulher que gostava do gosto da água da bomba da Broad Street. Ela e a sobrinha beberam a água da Broad Street e morreram, enquanto ninguém mais na área dela sequer adoeceu. Juntando todas essas evidências, Snow pede às autoridades locais que removam a manivela da bomba, e em 8 de setembro elas concordam. Como escreveu o biógrafo de Snow: “A manivela da bomba foi removida, e a praga foi detida.”

Tudo isso faz uma história maravilhosa. Hoje em dia, uma Sociedade John Snow até reencena a remoção da famosa manivela da bomba todos os anos. Contudo, na verdade, a remoção da manivela da bomba mal fez um arranhão na epidemia de cólera em toda a cidade, que acabou ceifando quase 3.000 vidas.

No capítulo não hollywoodiano da história, vemos novamente o Dr. Snow caminhando pelas ruas de Londres, mas desta vez o seu verdadeiro objetivo é descobrir de onde os londrinos obtêm a sua água. Havia duas companhias principais de água na época: a Southwark and Vauxhall Company e a Lambeth Company. A diferença-chave entre as duas, como Snow sabia, era que a primeira tirava a água da área da London Bridge, que ficava a jusante dos esgotos de Londres. A última havia movido a sua captação de água vários anos antes, de modo a ficar a montante dos esgotos. Assim, os clientes da Southwark estavam recebendo água contaminada pelas fezes de vítimas de cólera. Os clientes da Lambeth, por outro lado, estavam recebendo água não contaminada. (Nada disso tem a ver com a água contaminada da Broad Street, que vinha de um poço.)

As estatísticas de mortes corroboraram a sombria hipótese de Snow. Distritos abastecidos pela Southwark and Vauxhall Company foram especialmente atingidos pela cólera e tinham uma taxa de mortalidade oito vezes maior. Ainda assim, a evidência era meramente circunstancial. Um defensor da teoria do miasma poderia argumentar que o miasma era mais forte naqueles distritos, e não haveria modo de desmenti-lo. Em termos de um diagrama causal, temos a situação diagramada na Figura 7.7. Não temos como observar o confundidor Miasma (ou outros confundidores como Pobreza), de modo que não podemos controlá-lo usando o ajustamento pela porta dos fundos.

\begin{figure}[ht]
\centering
\includegraphics[width=0.92\textwidth]{../figuras/fig_7_7.png}
\caption{Figura 7.7. Não temos como observar o confundidor Miasma (ou outros confundidores como Pobreza), de modo que não podemos controlá-lo usando o ajustamento pela porta dos fundos.}
\end{figure}
\textit{Figura 7.7. Não temos como observar o confundidor Miasma (ou outros confundidores como Pobreza), de modo que não podemos controlá-lo usando o ajustamento pela porta dos fundos.}

\textbf{FIGURA 7.7.} Diagrama causal para a cólera (antes da descoberta do bacilo da cólera).

Aqui Snow teve a sua ideia mais brilhante. Ele notou que, naqueles distritos servidos por ambas as companhias, a taxa de mortalidade ainda era muito maior nas residências que recebiam água da Southwark. Contudo, essas residências não diferiam em termos de miasma ou pobreza. “A mistura do abastecimento é da espécie mais íntima”, escreveu Snow. “Os canos de cada Companhia descem por todas as ruas, e por quase todos os pátios e becos. … Cada companhia abastece tanto ricos quanto pobres, tanto casas grandes quanto pequenas; não há diferença nem na condição nem na ocupação das pessoas que recebem a água das diferentes Companhias.”

Embora a noção de um ECR ainda estivesse no futuro, era muito como se as companhias de água tivessem conduzido um experimento randomizado com os londrinos. De fato, Snow até nota isso: “Nenhum experimento poderia ter sido concebido que testasse de modo mais completo o efeito do abastecimento de água sobre o progresso da cólera do que este, que as circunstâncias colocaram pronto diante do observador. O experimento, além disso, era na grander escala. Não menos de trezentas mil pessoas de ambos os sexos, de todas as idades e ocupações, e de todas as classes e posições, desde a gente fina até os muito pobres, foram divididas em dois grupos sem a sua escolha e, na maioria dos casos, sem o seu conhecimento.” Um grupo havia recebido água pura; o outro, água contaminada por esgoto.

As observações de Snow introduziram uma nova variável no diagrama causal, que agora se parece com a Figura 7.8. O meticuloso trabalho de detetive de Snow havia mostrado duas coisas importantes: (1) não há seta entre Miasma e Companhia de Água (as duas são independentes), e (2) há uma seta entre Companhia de Água e Pureza da Água. Deixada sem enunciação por Snow, mas igualmente importante, é uma terceira hipótese: (3) a ausência de uma seta direta de Companhia de Água para Cólera, que é bastante óbvia para nós hoje porque sabemos que as companhias de água não estavam entregando cólera aos seus clientes por alguma rota alternativa.

\begin{figure}[ht]
\centering
\includegraphics[width=0.92\textwidth]{../figuras/fig_7_8.png}
\caption{Figura 7.8. O meticuloso trabalho de detetive de Snow havia mostrado duas coisas importantes: (1) não há seta entre Miasma e Companhia de Água (as duas são independentes), e (2) há uma seta entre Companhia de Água e Pureza da Água. Deixada sem enunciação por Snow, mas igualmente importante, é uma terceira hipótese: (3) a ausência de uma seta direta de Companhia de Água para Cólera, que é bastante óbvia para nós hoje porque sabemos que as companhias de água não estavam entregando cólera aos seus clientes por alguma rota alternativa.}
\end{figure}
\textit{Figura 7.8. O meticuloso trabalho de detetive de Snow havia mostrado duas coisas importantes: (1) não há seta entre Miasma e Companhia de Água (as duas são independentes), e (2) há uma seta entre Companhia de Água e Pureza da Água. Deixada sem enunciação por Snow, mas igualmente importante, é uma terceira hipótese: (3) a ausência de uma seta direta de Companhia de Água para Cólera, que é bastante óbvia para nós hoje porque sabemos que as companhias de água não estavam entregando cólera aos seus clientes por alguma rota alternativa.}

\textbf{FIGURA 7.8.} Diagrama para a cólera após a introdução de uma variável instrumental.

Uma variável que satisfaz essas três propriedades é hoje chamada de variável instrumental. Claramente Snow pensava nessa variável como semelhante a um lançamento de moeda, que simula uma variável sem setas de entrada. Como não há confundidores da relação entre Companhia de Água e Cólera, qualquer associação observada deve ser causal. Do mesmo modo, como o efeito de Companhia de Água sobre Cólera deve passar por Pureza da Água, concluímos (como Snow) que a associação observada entre Pureza da Água e Cólera também deve ser causal.

Snow enunciou a sua conclusão em termos inequívocos: se a Southwark and Vauxhall Company tivesse movido o seu ponto de captação para montante, mais de 1.000 vidas teriam sido salvas. Poucas pessoas tomaram nota da conclusão de Snow na época. Ele imprimiu um panfleto dos resultados às próprias custas, e vendeu um total grandioso de cinquenta e seis exemplares. Hoje em dia, os epidemiologistas veem o seu panfleto como o documento seminal de sua disciplina. Mostrou que, por meio de “\textit{shoe-leather research}” (uma expressão que tomei emprestada de David Freedman) — pesquisa de campo a pé, de porta em porta — e raciocínio causal, você pode rastrear um assassino.

Embora a teoria do miasma tenha sido desacreditada, a pobreza era indubitavelmente um confundidor, assim como a localização. Mas mesmo sem medir esses fatores (porque o trabalho de detetive de porta em porta de Snow só foi até certo ponto), ainda podemos usar variáveis instrumentais para determinar quantas vidas teriam sido salvas pela purificação do abastecimento de água. Eis como o truque funciona. Por simplicidade, voltaremos aos nomes \textit{Z}, \textit{X}, \textit{Y} e \textit{U} para nossas variáveis e redesenharemos a Figura 7.8 como se vê na Figura 7.9. Incluí coeficientes de caminho (\textit{a}, \textit{b}, \textit{c}, \textit{d}) para representar a força dos efeitos causais. Isso significa que estamos assumindo que as variáveis são numéricas e as funções que as relacionam são lineares. Lembre-se de que o coeficiente de caminho \textit{a} significa que uma intervenção para aumentar \textit{Z} em uma unidade padrão fará \textit{X} aumentar em \textit{a} unidades padrão. (Omitirei os detalhes técnicos do que são as “unidades padrão”.)

\begin{figure}[ht]
\centering
\includegraphics[width=0.92\textwidth]{../figuras/fig_7_9.png}
\caption{Figura 7.9. (ou 7.11 ou 7.12), mas não é possível ir muito longe sem a orientação de diagramas causais. Por exemplo, em alguns casos um instrumento imperfeito (p.ex., um que não seja independente do confundidor) pode ser usado após condicionar a um conjunto engenhosamente escolhido de variáveis auxiliares, que bloqueiam os caminhos entre o instrumento e o confundidor. Meu antigo aluno Carlos Brito, agora professor na Universidade Federal do Ceará, Brasil, desenvolveu plenamente essa ideia de transformar variáveis não instrumentais em variáveis instrumentais. Além disso, Brito estudou muitos casos em que um conjunto de variáveis pode ser usado com sucesso como um instrumento. Embora a identificação de conjuntos instrumentais vá além do cálculo do \textit{do}, ainda usa as ferramentas dos diagramas causais. Para pesquisadores que compreendem essa linguagem, os possíveis desenhos de pesquisa são ricos e variados; não precisam se sentir constrangidos a usar apenas o modelo de quatro variáveis mostrado nas Figuras 7.9, 7.11 e 7.12. As possibilidades são limitadas apenas pela nossa imaginação.}
\end{figure}
\textit{Figura 7.9. (ou 7.11 ou 7.12), mas não é possível ir muito longe sem a orientação de diagramas causais. Por exemplo, em alguns casos um instrumento imperfeito (p.ex., um que não seja independente do confundidor) pode ser usado após condicionar a um conjunto engenhosamente escolhido de variáveis auxiliares, que bloqueiam os caminhos entre o instrumento e o confundidor. Meu antigo aluno Carlos Brito, agora professor na Universidade Federal do Ceará, Brasil, desenvolveu plenamente essa ideia de transformar variáveis não instrumentais em variáveis instrumentais. Além disso, Brito estudou muitos casos em que um conjunto de variáveis pode ser usado com sucesso como um instrumento. Embora a identificação de conjuntos instrumentais vá além do cálculo do }do\textit{, ainda usa as ferramentas dos diagramas causais. Para pesquisadores que compreendem essa linguagem, os possíveis desenhos de pesquisa são ricos e variados; não precisam se sentir constrangidos a usar apenas o modelo de quatro variáveis mostrado nas Figuras 7.9, 7.11 e 7.12. As possibilidades são limitadas apenas pela nossa imaginação.}

\textbf{FIGURA 7.9.} Configuração geral para variáveis instrumentais.

Como \textit{Z} e \textit{X} não estão confundidos, o efeito causal de \textit{Z} sobre \textit{X} (isto é, \textit{a}) pode ser estimado a partir da inclinação \textit{r<sub>XZ</sub>} da reta de regressão de \textit{X} sobre \textit{Z}. Do mesmo modo, as variáveis \textit{Z} e \textit{Y} não estão confundidas, porque o caminho \textit{Z} → \textit{X} ← \textit{U} → \textit{Y} é bloqueado pelo colisor em \textit{X}. Assim, a inclinação da reta de regressão de \textit{Z} sobre \textit{Y} (\textit{r<sub>ZY</sub>}) será igual ao efeito causal no caminho direto \textit{Z} → \textit{X} → \textit{Y}, que é o produto dos coeficientes de caminho: \textit{ab}. Assim, temos duas equações: \textit{ab} = \textit{r<sub>ZY</sub>} e \textit{a} = \textit{r<sub>ZX</sub>}. Se dividirmos a primeira equação pela segunda, obtemos o efeito causal de \textit{X} sobre \textit{Y}: \textit{b} = \textit{r<sub>ZY</sub>}/\textit{r<sub>ZX</sub>}.

Desse modo, as variáveis instrumentais nos permitem realizar o mesmo tipo de truque de mágica que fizemos com o ajustamento pela porta da frente: encontramos o efeito de \textit{X} sobre \textit{Y} mesmo sem podermos controlar por, ou coletar dados sobre, o confundidor \textit{U}. Podemos, portanto, fornecer aos decisores um argumento conclusivo de que devem mudar o seu abastecimento de água — mesmo se esses decisores ainda acreditarem na teoria do miasma.

Note também que obtivemos informação no segundo degrau da Escada da Causalidade (\textit{b}) a partir de informação sobre o primeiro degrau (as correlações, \textit{r<sub>ZY</sub>} e \textit{r<sub>ZX</sub>}). Pudemos fazer isso porque as hipóteses incorporadas no diagrama de caminhos são de natureza causal, especialmente a hipótese crucial de que não há seta entre \textit{U} e \textit{Z}. Se o diagrama causal fosse diferente — por exemplo, se \textit{Z} fosse um confundidor de \textit{X} e \textit{Y} — a fórmula \textit{b} = \textit{r<sub>ZY</sub>}/\textit{r<sub>ZX</sub>} não estimaria corretamente o efeito causal de \textit{X} sobre \textit{Y}. De fato, esses dois modelos não podem ser distinguidos por qualquer método estatístico, independentemente de quão grandes sejam os dados.

As variáveis instrumentais eram conhecidas antes da Revolução Causal, mas os diagramas causais trouxeram nova clareza sobre como elas funcionam. De fato, Snow estava usando uma variável instrumental de forma implícita, embora não tivesse uma fórmula quantitativa. Sewall Wright certamente compreendia esse uso de diagramas de caminhos; a fórmula \textit{b} = \textit{r<sub>ZY</sub>}/\textit{r<sub>ZX</sub>} pode ser derivada diretamente do seu método de coeficientes de caminho. E parece que a primeira pessoa, além de Sewall Wright, a usar variáveis instrumentais de forma deliberada foi… o pai de Sewall Wright, Philip!

Recorde que Philip Wright era um economista que trabalhava no que mais tarde se tornou a Brookings Institution. Ele estava interessado em prever como a produção de uma mercadoria mudaria se uma tarifa fosse imposta, o que elevaria o preço e, portanto, em teoria, encorajaria a produção. Em termos econômicos, ele queria conhecer a elasticidade da oferta. Em 1928, Wright escreveu uma longa monografia dedicada a calcular a elasticidade da oferta do óleo de linhaça. Em um apêndice notável, ele analisou o problema usando um diagrama de caminhos. Isso foi uma coisa corajosa a fazer: lembre-se de que nenhum economista jamais havia visto ou ouvido falar de tal coisa antes. (De fato, ele se precaveu e verificou os cálculos usando métodos mais tradicionais.)

A Figura 7.10 mostra uma versão um tanto simplificada do diagrama de Wright. Diferentemente da maioria dos diagramas neste livro, este tem setas “de mão dupla”, mas eu pediria ao leitor que não perdesse o sono demais com isso. Com algum truque matemático, poderíamos igualmente substituir a cadeia Demanda ⇄ Preço ⇄ Oferta por uma única seta Demanda → Oferta, e a figura se pareceria então com a Figura 7.9 (embora fosse menos aceitável para economistas). O ponto importante a notar é que Philip Wright deliberadamente introduziu a variável Rendimento por Acre (de semente de linhaça) como um instrumento que afeta diretamente a oferta mas não tem correlação com a demanda. Ele então usou uma análise como a que acabei de dar para deduzir tanto o efeito da oferta sobre o preço quanto o efeito do preço sobre a oferta.

\begin{figure}[ht]
\centering
\includegraphics[width=0.92\textwidth]{../figuras/fig_7_10.png}
\caption{Figura 7.10. mostra uma versão um tanto simplificada do diagrama de Wright. Diferentemente da maioria dos diagramas neste livro, este tem setas “de mão dupla”, mas eu pediria ao leitor que não perdesse o sono demais com isso. Com algum truque matemático, poderíamos igualmente substituir a cadeia Demanda ⇄ Preço ⇄ Oferta por uma única seta Demanda → Oferta, e a figura se pareceria então com a Figura 7.9 (embora fosse menos aceitável para economistas). O ponto importante a notar é que Philip Wright deliberadamente introduziu a variável Rendimento por Acre (de semente de linhaça) como um instrumento que afeta diretamente a oferta mas não tem correlação com a demanda. Ele então usou uma análise como a que acabei de dar para deduzir tanto o efeito da oferta sobre o preço quanto o efeito do preço sobre a oferta.}
\end{figure}
\textit{Figura 7.10. mostra uma versão um tanto simplificada do diagrama de Wright. Diferentemente da maioria dos diagramas neste livro, este tem setas “de mão dupla”, mas eu pediria ao leitor que não perdesse o sono demais com isso. Com algum truque matemático, poderíamos igualmente substituir a cadeia Demanda ⇄ Preço ⇄ Oferta por uma única seta Demanda → Oferta, e a figura se pareceria então com a Figura 7.9 (embora fosse menos aceitável para economistas). O ponto importante a notar é que Philip Wright deliberadamente introduziu a variável Rendimento por Acre (de semente de linhaça) como um instrumento que afeta diretamente a oferta mas não tem correlação com a demanda. Ele então usou uma análise como a que acabei de dar para deduzir tanto o efeito da oferta sobre o preço quanto o efeito do preço sobre a oferta.}

\textbf{FIGURA 7.10.} Versão simplificada do diagrama causal oferta-preço de Wright.

Os historiadores discutem sobre quem inventou as variáveis instrumentais, um método que se tornou extremamente popular na econometria moderna. Não há dúvida em minha mente de que Philip Wright tomou emprestada a ideia de coeficientes de caminho de seu filho. Nenhum economista jamais insistira antes na distinção entre coeficientes causais e coeficientes de regressão; estavam todos no campo de Karl Pearson–Henry Niles, de que a causalidade não é senão um caso-limite de correlação. Também, ninguém antes de Sewall Wright jamais dera uma receita para calcular coeficientes de regressão em termos de coeficientes de caminho, e então reverter o processo para obter os coeficientes causais a partir da regressão. Essa foi a invenção exclusiva de Sewall. Naturalmente, alguns historiadores da economia sugeriram que Sewall escreveu o apêndice matemático inteiro sozinho. Contudo, a análise estilométrica mostrou que Philip foi de fato o autor. Para mim, esse trabalho de detetive histórico torna a história mais bela. Mostra que Philip se deu ao trabalho de compreender a teoria de seu filho e articulá-la em sua própria linguagem.

Agora vamos avançar dos anos 1850 e 1920 para olhar um exemplo contemporâneo de variáveis instrumentais em ação, um entre literalmente dezenas que eu poderia ter escolhido.

\section{Colesterol bom e colesterol ruim}
Você se lembra de quando o seu médico de família começou a falar com você sobre colesterol “bom” e “ruim”? Pode ter acontecido nos anos 1990, quando os medicamentos que baixavam os níveis sanguíneos de colesterol “ruim”, a lipoproteína de baixa densidade (LDL), chegaram ao mercado. Esses medicamentos, chamados estatinas, tornaram-se geradores de receita multibilionária para as companhias farmacêuticas.

O primeiro medicamento modificador de colesterol submetido a um ensaio clínico randomizado foi a colestiramina. O Coronary Primary Prevention Trial, iniciado em 1973 e concluído em 1984, mostrou uma redução de 12,6 por cento no colesterol entre homens a quem se deu o medicamento colestiramina e uma redução de 19 por cento no risco de infarto. Como se tratava de um ensaio clínico randomizado, você poderia pensar que não precisaríamos de nenhum dos métodos deste capítulo, porque eles são especificamente concebidos para substituir ECRs em situações em que se tem apenas dados observacionais. Mas isso não é verdade. Esse ensaio, como muitos ECRs, enfrentou o problema da não adesão (\textit{noncompliance}), quando sujeitos randomizados para receber um medicamento de fato não o tomam. Isso reduzirá a aparente efetividade do medicamento, de modo que podemos querer ajustar os resultados para dar conta dos não aderentes. Mas, como sempre, o confundimento ergue a sua feia cabeça. Se os não aderentes forem diferentes dos aderentes de alguma maneira relevante (talvez estejam mais doentes para começar?), não podemos prever como teriam respondido se tivessem aderido às instruções.

Nessa situação, temos um diagrama causal que se parece com a Figura 7.11. A variável Designado (\textit{Z}) tomará o valor 1 se o paciente for aleatoriamente designado a receber o medicamento e 0 se for aleatoriamente designado a um placebo. A variável Recebeu será 1 se o paciente de fato tomou o medicamento e 0 caso contrário. Por conveniência, usaremos também uma definição binária para Colesterol, registrando um desfecho de 1 se os níveis de colesterol foram reduzidos por uma certa quantidade fixa.

\begin{figure}[ht]
\centering
\includegraphics[width=0.92\textwidth]{../figuras/fig_7_11.png}
\caption{Figura 7.11. A variável Designado (\textit{Z}) tomará o valor 1 se o paciente for aleatoriamente designado a receber o medicamento e 0 se for aleatoriamente designado a um placebo. A variável Recebeu será 1 se o paciente de fato tomou o medicamento e 0 caso contrário. Por conveniência, usaremos também uma definição binária para Colesterol, registrando um desfecho de 1 se os níveis de colesterol foram reduzidos por uma certa quantidade fixa.}
\end{figure}
\textit{Figura 7.11. A variável Designado (}Z\textit{) tomará o valor 1 se o paciente for aleatoriamente designado a receber o medicamento e 0 se for aleatoriamente designado a um placebo. A variável Recebeu será 1 se o paciente de fato tomou o medicamento e 0 caso contrário. Por conveniência, usaremos também uma definição binária para Colesterol, registrando um desfecho de 1 se os níveis de colesterol foram reduzidos por uma certa quantidade fixa.}

\textbf{FIGURA 7.11.} Diagrama causal para um ECR com não adesão.

Note que, nesse caso, nossas variáveis são binárias, não numéricas. Isso significa de imediato que não podemos usar um modelo linear e, portanto, não podemos aplicar a fórmula de variáveis instrumentais que derivamos antes. Contudo, em tais casos podemos frequentemente substituir a hipótese de linearidade por uma condição mais fraca chamada monotonicidade, que explicarei abaixo.

Mas, antes de fazê-lo, certifiquemo-nos de que nossas outras hipóteses necessárias para variáveis instrumentais são válidas. Primeiro, a variável instrumental \textit{Z} é independente do confundidor? A randomização de \textit{Z} garante que a resposta é sim. (Como vimos no Capítulo 4, a randomização é um ótimo modo de garantir que uma variável não seja afetada por quaisquer confundidores.) Há algum caminho direto de \textit{Z} para \textit{Y}? O senso comum diz que não há modo de que receber um número aleatório particular (\textit{Z}) afetasse o colesterol (\textit{Y}), de modo que a resposta é não. Por fim, há uma associação forte entre \textit{Z} e \textit{X}? Dessa vez os próprios dados devem ser consultados, e a resposta novamente é sim. Devemos sempre fazer as três perguntas acima antes de aplicar variáveis instrumentais. Aqui as respostas são óbvias, mas não devemos ficar cegos para o fato de que estamos usando intuição causal para respondê-las — intuição que é capturada, preservada e elucidada no diagrama.

A Tabela 7.1 mostra as frequências observadas dos desfechos \textit{X} e \textit{Y}. Por exemplo, 91,9 por cento das pessoas a quem não se designou o medicamento tiveram o desfecho \textit{X} = 0 (não tomaram o medicamento) e \textit{Y} = 0 (sem redução no colesterol). Isso faz sentido. Os outros 8,1 por cento tiveram o desfecho \textit{X} = 0 (não tomaram o medicamento) e \textit{Y} = 1 (tiveram de fato uma redução no colesterol). Evidentemente melhoraram por outras razões que não tomar o medicamento. Note também que há dois zeros na tabela: não houve ninguém a quem não se designou o medicamento (\textit{Z} = 0) mas que, ainda assim, o tenha obtido (\textit{X} = 1). Em um estudo randomizado bem conduzido, especialmente no campo médico em que os médicos têm acesso exclusivo ao medicamento experimental, isso tipicamente será verdade. A hipótese de que não há indivíduos com \textit{Z} = 0 e \textit{X} = 1 chama-se monotonicidade.

\begin{figure}[ht]
\centering
\includegraphics[width=0.92\textwidth]{../figuras/tab_7_1.png}
\caption{Tabela 7.1. mostra as frequências observadas dos desfechos \textit{X} e \textit{Y}. Por exemplo, 91,9 por cento das pessoas a quem não se designou o medicamento tiveram o desfecho \textit{X} = 0 (não tomaram o medicamento) e \textit{Y} = 0 (sem redução no colesterol). Isso faz sentido. Os outros 8,1 por cento tiveram o desfecho \textit{X} = 0 (não tomaram o medicamento) e \textit{Y} = 1 (tiveram de fato uma redução no colesterol). Evidentemente melhoraram por outras razões que não tomar o medicamento. Note também que há dois zeros na tabela: não houve ninguém a quem não se designou o medicamento (\textit{Z} = 0) mas que, ainda assim, o tenha obtido (\textit{X} = 1). Em um estudo randomizado bem conduzido, especialmente no campo médico em que os médicos têm acesso exclusivo ao medicamento experimental, isso tipicamente será verdade. A hipótese de que não há indivíduos com \textit{Z} = 0 e \textit{X} = 1 chama-se monotonicidade.}
\end{figure}
\textit{Tabela 7.1. mostra as frequências observadas dos desfechos }X\textit{ e }Y\textit{. Por exemplo, 91,9 por cento das pessoas a quem não se designou o medicamento tiveram o desfecho }X\textit{ = 0 (não tomaram o medicamento) e }Y\textit{ = 0 (sem redução no colesterol). Isso faz sentido. Os outros 8,1 por cento tiveram o desfecho }X\textit{ = 0 (não tomaram o medicamento) e }Y\textit{ = 1 (tiveram de fato uma redução no colesterol). Evidentemente melhoraram por outras razões que não tomar o medicamento. Note também que há dois zeros na tabela: não houve ninguém a quem não se designou o medicamento (}Z\textit{ = 0) mas que, ainda assim, o tenha obtido (}X\textit{ = 1). Em um estudo randomizado bem conduzido, especialmente no campo médico em que os médicos têm acesso exclusivo ao medicamento experimental, isso tipicamente será verdade. A hipótese de que não há indivíduos com }Z\textit{ = 0 e }X\textit{ = 1 chama-se monotonicidade.}

\textbf{TABELA 7.1.} Dados do ensaio da colestiramina.

Agora vejamos como podemos estimar o efeito do tratamento. Primeiro tomemos o cenário de pior caso: nenhum dos não aderentes teria melhorado se tivesse aderido ao tratamento. Nesse caso, as únicas pessoas que teriam tomado o medicamento e melhorado seriam os 47,3 por cento que de fato aderiram e melhoraram. Mas precisamos corrigir essa estimativa pelo efeito placebo, que está na terceira linha da tabela. Das pessoas a quem se designou o placebo e que tomaram o placebo, 8,1 por cento melhoraram. Assim, a melhoria líquida acima e além do efeito placebo é 47,3 por cento menos 8,1 por cento, ou 39,2 por cento.

E quanto ao cenário de melhor caso, em que todos os não aderentes teriam melhorado se tivessem aderido? Nesse caso, adicionamos os 31,5 por cento mais 7,3 por cento dos não aderentes à linha de base de 39,2 por cento que acabamos de calcular, para um total de 78,0 por cento. Assim, mesmo no cenário de pior caso, em que o confundimento vai completamente contra o medicamento, ainda podemos dizer que o medicamento melhora o colesterol para 39 por cento da população. No cenário de melhor caso, em que o confundimento trabalha completamente a favor do medicamento, 78 por cento da população veriam uma melhoria. Embora os limites estejam bastante afastados, devido ao grande número de não aderentes, o pesquisador pode afirmar categoricamente que o medicamento é efetivo para o propósito a que se destina.

Essa estratégia de tomar o pior caso e depois o melhor caso usualmente nos dará um intervalo de estimativas. Obviamente seria bom ter uma estimativa pontual, como fizemos no caso linear. Há modos de estreitar o intervalo se necessário, e em alguns casos é até possível obter estimativas pontuais. Por exemplo, se você estiver interessado apenas na subpopulação de aderentes (aquelas pessoas que tomarão \textit{X} se e somente se designadas), pode derivar uma estimativa pontual conhecida como Efeito Médio Local do Tratamento (LATE, do inglês \textit{Local Average Treatment Effect}). Em todo caso, espero que este exemplo mostre que nossas mãos não estão atadas quando deixamos o mundo dos modelos lineares.

Os métodos de variáveis instrumentais continuaram a se desenvolver desde 1984, e uma versão particular tornou-se extremamente popular: a randomização mendeliana. Eis um exemplo. Embora o efeito da LDL, ou colesterol “ruim”, esteja agora assentado, ainda há considerável incerteza sobre a lipoproteína de alta densidade (HDL), ou colesterol “bom”. Estudos observacionais iniciais, como o Framingham Heart Study no final dos anos 1970, sugeriam que a HDL tinha um efeito protetor contra infartos. Mas a HDL alta frequentemente anda de mãos dadas com a LDL baixa, de modo que como podemos dizer qual lipídio é o verdadeiro fator causal?

Para responder a essa pergunta, suponha que conhecêssemos um gene que fizesse as pessoas terem níveis mais altos de HDL, sem efeito sobre a LDL. Então poderíamos montar o diagrama causal da Figura 7.12, em que usei Estilo de Vida como um possível confundidor. Lembre-se de que é sempre vantajoso, como no exemplo de Snow, usar uma variável instrumental que seja randomizada. Se for randomizada, nenhuma seta causal aponta para ela. Por essa razão, um gene é uma variável instrumental perfeita. Nossos genes são randomizados no momento da concepção, de modo que é como se o próprio Gregor Mendel tivesse estendido a mão do céu e designado a algumas pessoas um gene de alto risco e a outras um gene de baixo risco. Essa é a razão do termo “randomização mendeliana”.

Poderia haver uma seta no outro sentido, de Gene da HDL para Estilo de Vida? Aqui precisamos novamente fazer “\textit{shoe-leather work}” e pensar causalmente. O gene da HDL só poderia afetar o estilo de vida das pessoas se elas soubessem qual versão tinham, a versão de HDL alta ou a de HDL baixa. Mas até 2008 nenhum gene desse tipo era conhecido, e mesmo hoje as pessoas não têm rotineiramente acesso a essa informação. Assim, é altamente provável que nenhuma seta desse tipo exista.

\textbf{FIGURA 7.12.} Diagrama causal para o exemplo de randomização mendeliana.

Pelo menos dois estudos adotaram essa abordagem de randomização mendeliana à questão do colesterol. Em 2012, um estudo colaborativo gigante liderado por Sekar Kathiresan do Massachusetts General Hospital mostrou que não havia benefício observável de níveis mais altos de HDL. Por outro lado, os pesquisadores descobriram que a LDL tem um efeito muito grande sobre o risco de infarto. De acordo com as suas cifras, diminuir a sua contagem de LDL em 34 mg/dl reduziria as suas chances de um infarto em cerca de 50 por cento. Assim, baixar os seus níveis de colesterol “ruim”, seja por dieta, exercício ou estatinas, parece ser uma ideia inteligente. Por outro lado, aumentar os seus níveis de colesterol “bom”, apesar do que alguns vendedores de óleo de peixe possam lhe dizer, não parece provável que mude o seu risco de infarto de modo algum.

Como sempre, há uma ressalva. O segundo estudo, publicado no mesmo ano, apontou que as pessoas com a variante de menor risco do gene da LDL tiveram níveis mais baixos de colesterol durante toda a vida. A randomização mendeliana nos diz que diminuir a sua LDL em trinta e quatro unidades ao longo de toda a vida diminuirá o seu risco de infarto em 50 por cento. Mas as estatinas não podem baixar o seu colesterol LDL ao longo de toda a vida — apenas a partir do dia em que você começa a tomar o medicamento. Se você tem sessenta anos, as suas artérias já sustentaram sessenta anos de dano. Por essa razão, é muito provável que a randomização mendeliana superestime os verdadeiros benefícios das estatinas. Por outro lado, começar a reduzir o seu colesterol quando se é jovem — seja por dieta, exercício ou mesmo estatinas — terá grandes efeitos mais tarde.

Do ponto de vista da análise causal, isso nos ensina uma boa lição: em qualquer estudo de intervenções, precisamos perguntar se a variável que estamos de fato manipulando (níveis de LDL ao longo da vida) é a mesma que a variável que pensamos estar manipulando (níveis atuais de LDL). Isso faz parte da “interrogação hábil da natureza”.

Em resumo, as variáveis instrumentais são uma ferramenta importante na medida em que nos ajudam a revelar informação causal que vai além do cálculo do \textit{do}. Este último insiste em estimativas pontuais em vez de desigualdades e desistiria de casos como a Figura 7.12, em que tudo o que podemos obter são desigualdades. Por outro lado, é também importante perceber que o cálculo do \textit{do} é vastamente mais flexível do que as variáveis instrumentais. No cálculo do \textit{do} não fazemos hipótese alguma sobre a natureza das funções no modelo causal. Mas se pudermos justificar uma hipótese como monotonicidade ou linearidade em bases científicas, então uma ferramenta de propósito mais especial como as variáveis instrumentais vale a pena ser considerada.

Os métodos de variáveis instrumentais podem ser estendidos além de modelos simples de quatro variáveis como a Figura 7.9 (ou 7.11 ou 7.12), mas não é possível ir muito longe sem a orientação de diagramas causais. Por exemplo, em alguns casos um instrumento imperfeito (p.ex., um que não seja independente do confundidor) pode ser usado após condicionar a um conjunto engenhosamente escolhido de variáveis auxiliares, que bloqueiam os caminhos entre o instrumento e o confundidor. Meu antigo aluno Carlos Brito, agora professor na Universidade Federal do Ceará, Brasil, desenvolveu plenamente essa ideia de transformar variáveis não instrumentais em variáveis instrumentais. Além disso, Brito estudou muitos casos em que um conjunto de variáveis pode ser usado com sucesso como um instrumento. Embora a identificação de conjuntos instrumentais vá além do cálculo do \textit{do}, ainda usa as ferramentas dos diagramas causais. Para pesquisadores que compreendem essa linguagem, os possíveis desenhos de pesquisa são ricos e variados; não precisam se sentir constrangidos a usar apenas o modelo de quatro variáveis mostrado nas Figuras 7.9, 7.11 e 7.12. As possibilidades são limitadas apenas pela nossa imaginação.

\chapter{Capítulo 8 — Contrafactuais: minerando mundos que poderiam ter sido}
\textit{Se o nariz de Cleópatra tivesse sido mais curto, toda a face do mundo teria mudado.} — BLAISE PASCAL (1669)

Ao nos prepararmos para subir ao degrau mais alto da Escada da Causalidade, recapitulamos o que aprendemos no segundo degrau. Vimos várias maneiras de averiguar o efeito de uma intervenção em diversos cenários e sob uma variedade de condições. No Capítulo 4, discutimos os ensaios clínicos randomizados, o amplamente citado "padrão-ouro" dos ensaios médicos. Vimos também métodos adequados a estudos observacionais, em que os grupos de tratamento e controle não são atribuídos ao acaso. Se pudermos medir variáveis que bloqueiem todos os caminhos da porta dos fundos, podemos usar a fórmula de ajustamento pela porta dos fundos para obter o efeito necessário. Se encontrarmos um caminho da porta da frente "protegido" de confundidores, podemos usar o ajustamento pela porta da frente. Se estivermos dispostos a conviver com a hipótese de linearidade ou de monotonicidade, podemos usar variáveis instrumentais (supondo que uma variável apropriada possa ser encontrada no diagrama ou criada por um experimento). E pesquisadores verdadeiramente aventureiros podem traçar outras rotas até o cume do Monte Intervenção usando o cálculo do \textit{do} ou sua versão algorítmica.

Em todos esses empreendimentos, lidamos com efeitos sobre uma população ou sobre um indivíduo típico selecionado de uma população de estudo (o efeito causal médio). Mas até agora falta-nos a capacidade de falar de causalidade personalizada no nível de eventos ou indivíduos particulares. Uma coisa é dizer "Fumar causa câncer"; outra é dizer que meu tio José, que fumou um maço por dia durante trinta anos, ainda estaria vivo se não tivesse fumado. A diferença é ao mesmo tempo óbvia e profunda: nenhuma das pessoas que, como o tio José, fumaram por trinta anos e morreram pode jamais ser observada no mundo alternativo em que não fumaram por trinta anos.

Responsabilidade e culpa, arrependimento e crédito: esses conceitos são a moeda de uma mente causal. Para fazer sentido deles, precisamos ser capazes de comparar o que de fato aconteceu com o que \textit{teria} acontecido sob alguma hipótese alternativa. Como se argumentou no Capítulo 1, nossa capacidade de conceber mundos alternativos e inexistentes nos separou de nossos ancestrais proto-humanos e, de fato, de qualquer outra criatura do planeta. Toda outra criatura pode ver o que \textit{é}. Nosso dom — que às vezes pode ser uma maldição — é que podemos ver o que \textit{poderia ter sido}.

Este capítulo mostra como usar dados observacionais e experimentais para extrair informação sobre cenários contrafactuais. Explica como representar causas no nível individual no contexto de um diagrama causal — tarefa que nos forçará a explicar alguns detalhes de funcionamento dos diagramas causais de que ainda não falamos. Discuto também um conceito estreitamente relacionado chamado "resultados potenciais", ou modelo causal de Neyman–Rubin, proposto inicialmente nos anos 1920 por Jerzy Neyman, estatístico polonês que mais tarde se tornou professor em Berkeley. Mas só depois que Donald Rubin começou a escrever sobre resultados potenciais em meados dos anos 1970 essa abordagem da análise causal de fato floresceu.

Mostrarei como os contrafactuais emergem naturalmente no arcabouço desenvolvido nos capítulos anteriores — os diagramas de caminhos de Sewall Wright e sua extensão aos modelos causais estruturais (MCE). Já tivemos um bom gosto disso no Capítulo 1, no exemplo do pelotão de fuzilamento, que mostrou como responder a perguntas contrafactuais como "O prisioneiro estaria vivo se o fuzileiro A não tivesse atirado?". Compararei como os contrafactuais são definidos no paradigma de Neyman–Rubin e nos MCE, onde desfrutam do benefício dos diagramas causais. Rubin manteve firmemente, ao longo dos anos, que os diagramas não servem a nenhum propósito útil. Examinaremos, portanto, como os estudantes do modelo causal de Rubin precisam navegar problemas causais de olhos vendados, sem um recurso para representar o conhecimento causal ou derivar suas implicações testáveis.

Por fim, veremos duas aplicações em que o raciocínio contrafactual é essencial. Por décadas ou mesmo séculos, advogados usaram um teste relativamente direto da culpabilidade de um réu chamado "causalidade \textit{but-for}" (ou "não fosse por"): a lesão não teria ocorrido \textit{não fosse} a ação do réu.[^1] Veremos como a linguagem dos contrafactuais captura essa noção elusiva e como estimar a probabilidade de que um réu seja culpável. Em seguida, discutirei a aplicação dos contrafactuais à mudança climática. Até recentemente, cientistas do clima achavam muito difícil e embaraçoso responder a perguntas como "O aquecimento global causou esta tempestade [ou esta onda de calor, ou esta seca]?". A resposta convencional tem sido que eventos meteorológicos individuais não podem ser atribuídos à mudança climática global. No entanto, essa resposta parece bastante evasiva e pode até contribuir para a indiferença pública em relação à mudança climática. A análise contrafactual permite que os cientistas do clima façam afirmações muito mais precisas e definidas do que antes. Exige, porém, um leve acréscimo ao nosso vocabulário cotidiano. Será útil distinguir três tipos diferentes de causalidade: causalidade necessária, causalidade suficiente e causalidade necessária e suficiente. (A causalidade necessária é a mesma que a causalidade \textit{but-for}.) Usando essas palavras, um cientista do clima pode dizer: "Há 90 por cento de probabilidade de que a mudança climática causada pelo homem tenha sido uma causa necessária desta onda de calor", ou "Há 80 por cento de probabilidade de que a mudança climática seja suficiente para produzir uma onda de calor tão forte pelo menos uma vez a cada 50 anos". A primeira frase tem a ver com atribuição: quem foi responsável pelo calor incomum? A segunda tem a ver com política pública. Diz que é melhor nos prepararmos para tais ondas de calor porque é provável que ocorram mais cedo ou mais tarde. Qualquer uma dessas afirmações é mais informativa do que encolher os ombros e não dizer nada sobre as causas de eventos meteorológicos individuais.

\section{De Tucídides e Abraão a Hume e Lewis}
Dado que o raciocínio contrafactual faz parte do aparato mental que nos torna humanos, não é surpreendente que possamos encontrar enunciados contrafactuais tão longe no tempo quanto quisermos na história humana. Por exemplo, na \textit{História da Guerra do Peloponeso} de Tucídides, o historiador grego antigo, frequentemente descrito como pioneiro de uma abordagem "científica" da história, descreve um tsunami ocorrido em 426 a.C.:

\begin{quote}Por volta da mesma época em que esses terremotos eram tão comuns, o mar em Oróbias, na Eubeia, retirando-se da então linha da costa, retornou numa onda enorme e invadiu grande parte da cidade, e recuou deixando parte dela ainda submersa; de modo que o que outrora era terra agora é mar; perecendo os habitantes que não puderam correr a tempo para o terreno mais alto.… A causa, em minha opinião, desse fenômeno deve ser buscada no terremoto. No ponto em que seu choque foi mais violento o mar é repelido e, recuando de súbito com força redobrada, causa a inundação. Sem um terremoto, não vejo como tal acidente poderia acontecer.\end{quote}
Esta é uma passagem verdadeiramente notável quando se considera a era em que foi escrita. Primeiro, a precisão das observações de Tucídides faria crédito a qualquer cientista moderno — e tanto mais porque ele trabalhava numa época em que não havia satélites, câmeras de vídeo nem organizações de notícias 24 horas transmitindo imagens do desastre enquanto se desenrolava. Segundo, ele escrevia num momento da história humana em que desastres naturais eram ordinariamente atribuídos à vontade dos deuses. Seu predecessor Homero ou seu contemporâneo Heródoto teriam, sem dúvida, atribuído o evento à ira de Poseidon ou de alguma outra divindade. No entanto, Tucídides propõe um modelo causal sem quaisquer processos sobrenaturais: o terremoto rechaça o mar, que recua e inunda a terra. A última frase da citação é especialmente interessante porque expressa a noção de causalidade necessária ou \textit{but-for}: não fosse o terremoto, o tsunami não poderia ter ocorrido. Esse juízo contrafactual promove o terremoto de mero antecedente do tsunami a causa efetiva.

Outra instância fascinante e reveladora de raciocínio contrafactual ocorre no livro do Gênesis na Bíblia. Abraão conversa com Deus sobre a intenção deste de destruir as cidades de Sodoma e Gomorra como retribuição por seus maus caminhos.

\begin{quote}E Abraão se aproximou e disse: Destruirás de fato o justo com o ímpio? Suponhamos que haja cinquenta justos dentro da cidade: destruirás também e não pouparás o lugar por causa dos cinquenta justos que nela houver?… E o Senhor disse: Se eu achar em Sodoma cinquenta justos dentro da cidade, pouparei todo o lugar por causa deles.\end{quote}
Mas a história não termina aí. Abraão não se satisfaz e pergunta ao Senhor: e se houver apenas quarenta e cinco justos? Ou quarenta? Ou trinta? Ou vinte? Ou mesmo dez? Cada vez recebe uma resposta afirmativa, e Deus finalmente o assegura de que poupará Sodoma mesmo por causa de dez justos, se puder encontrar tantos.

O que Abraão tenta realizar com essa regateação? Certamente não duvida da capacidade de Deus de contar. E, claro, Abraão sabe que Deus sabe quantos justos vivem em Sodoma. Ele é, afinal, onisciente. Conhecendo a obediência e a devoção de Abraão, é difícil crer que as perguntas visem convencer o Senhor a mudar de ideia. Em vez disso, destinam-se à própria compreensão de Abraão. Ele raciocina exatamente como um cientista moderno faria, tentando entender as leis que regem a punição coletiva. Que nível de maldade é suficiente para justificar a destruição? Trinta justos bastariam para salvar uma cidade? Vinte? Não temos um modelo causal completo sem tal informação. Um cientista moderno poderia chamá-la de curva dose-resposta ou efeito de limiar.

Enquanto Tucídides e Abraão sondavam contrafactuais por meio de casos individuais, o filósofo grego Aristóteles investigou aspectos mais genéricos da causalidade. Em seu estilo tipicamente sistemático, Aristóteles montou toda uma taxonomia da causalidade, incluindo "causas materiais", "causas formais", "causas eficientes" e "causas finais". Por exemplo, a causa material da forma de uma estátua é o bronze de que é fundida e suas propriedades; não poderíamos fazer a mesma estátua de massinha. No entanto, Aristóteles em nenhum lugar formula a causalidade como um contrafactual, de modo que sua engenhosa classificação carece da clareza simples do relato de Tucídides sobre a causa do tsunami.

Para encontrar um filósofo que colocou os contrafactuais no coração da causalidade, precisamos avançar até David Hume, o filósofo escocês e contemporâneo de Thomas Bayes. Hume rejeitou o esquema de classificação de Aristóteles e insistiu numa única definição de causalidade. Mas achou essa definição bastante elusiva e, de fato, ficou dividido entre duas definições diferentes. Mais tarde, elas se transformariam em duas ideologias incompatíveis, que ironicamente poderiam ambas citar Hume como sua fonte!

Em seu \textit{Tratado da natureza humana} (Figura 8.1), Hume nega que quaisquer dois objetos tenham qualidades ou "poderes" inatos que façam de um uma causa e do outro um efeito. Em sua visão, a relação causa-efeito é inteiramente um produto de nossa própria memória e experiência. "Assim, lembramo-nos de ter visto aquela espécie de objeto a que chamamos chama, e de ter sentido aquela espécie de sensação a que chamamos calor", escreve. "Recordamos igualmente sua conjunção constante em todas as instâncias passadas. Sem qualquer outra cerimônia, chamamos um de causa e o outro de efeito, e inferimos a existência de um a partir do outro." Esta é hoje conhecida como a definição de "regularidade" da causalidade.

\begin{figure}[ht]
\centering
\includegraphics[width=0.92\textwidth]{../figuras/fig_8_1.png}
\caption{Figura 8.1. A definição de “regularidade” de causa e efeito de Hume, proposta em 1739.}
\end{figure}
\textit{Figura 8.1. A definição de “regularidade” de causa e efeito de Hume, proposta em 1739.}

A passagem é de cortar o fôlego em sua ousadia. Hume está cortando o segundo e o terceiro degraus da Escada da Causalidade e dizendo que o primeiro degrau, a observação, é tudo de que precisamos. Uma vez que observamos chama e calor juntos um número suficiente de vezes (e note que a chama tem precedência temporal), concordamos em chamar a chama de causa do calor. Como a maioria dos estatísticos do século XX, Hume em 1739 parece contente em considerar a causalidade como mera espécie de correlação.

\textit{Figura 8.1. A definição de "regularidade" de causa e efeito de Hume, proposta em 1739.}

E, no entanto, Hume, para seu crédito, não permaneceu satisfeito com essa definição. Nove anos depois, em \textit{Investigação sobre o entendimento humano}, escreveu algo bem diferente: "Podemos definir uma causa como um objeto seguido por outro, e onde todos os objetos semelhantes ao primeiro são seguidos por objetos semelhantes ao segundo. Ou, em outras palavras, onde, se o primeiro objeto não tivesse existido, o segundo jamais teria existido" (ênfase no original).

A primeira frase, a versão em que A é consistentemente observado junto com B, simplesmente repete a definição de regularidade. Mas por 1748 ele parece ter alguns escrúpulos e acha-a em necessidade de reparo. Como historiadores \textit{whig} autorizados, podemos entender por quê. Segundo sua definição anterior, o canto do galo causaria o nascer do sol. Para remendar essa dificuldade, acrescenta uma segunda definição que nem sequer insinuara em seu livro anterior — uma definição contrafactual: "se o primeiro objeto não tivesse existido, o segundo jamais teria existido". Note que a segunda definição é exatamente a que Tucídides usou ao discutir o tsunami em Oróbias. A definição contrafactual também explica por que não consideramos o canto do galo uma causa do nascer do sol. Sabemos que se o galo estivesse doente um dia, ou se recusasse caprichosamente a cantar, o sol nasceria de qualquer modo.

Embora Hume tente passar essas duas definições por uma só, por meio de sua inocente interjeição "em outras palavras", a segunda versão é completamente diferente da primeira. Invoca explicitamente um contrafactual, de modo que se situa no terceiro degrau da Escada da Causalidade. Enquanto regularidades podem ser observadas, contrafactuais só podem ser imaginados.

Vale a pena pensar por um momento por que Hume escolhe definir causas em termos de contrafactuais, e não o contrário. Definições destinam-se a reduzir um conceito mais complicado a um mais simples. Hume supõe que seus leitores compreenderão o enunciado "se o primeiro objeto não tivesse existido, o segundo jamais teria existido" com menos ambiguidade do que "o primeiro objeto causou o segundo". Ele está absolutamente certo. Este último enunciado convida a toda sorte de especulação metafísica estéril sobre que qualidade ou poder inerente ao primeiro objeto produz o segundo. O primeiro enunciado apenas nos pede que realizemos um simples teste mental: imaginar um mundo sem o terremoto e perguntar se ele também contém um tsunami. Fazemos juízos como esse desde crianças, e a espécie humana os faz desde Tucídides (e provavelmente muito antes).

Não obstante, os filósofos ignoraram a segunda definição de Hume durante a maior parte dos séculos XIX e XX. Enunciados contrafactuais, os "teria sido", sempre pareceram demasiado moles e incertos para satisfazer acadêmicos. Em vez disso, os filósofos tentaram resgatar a primeira definição de Hume por meio da teoria da causalidade probabilística, como se discutiu no Capítulo 1.

Um filósofo que desafiou a convenção, David Lewis, pediu em seu livro de 1973 \textit{Counterfactuals} que se abandonasse de todo a conta da regularidade e se interpretasse "A causou B" como "B não teria ocorrido se não fosse por A". Lewis perguntou: "Por que não tomar os contrafactuais pelo valor de face: como enunciados sobre alternativas possíveis à situação atual?"

Como Hume, Lewis ficava evidentemente impressionado com o fato de que os humanos fazem juízos contrafactuais sem muito alarde — de forma rápida, cômoda e consistente. Podemos atribuir-lhes valores de verdade e probabilidades com não menos confiança do que o fazemos para enunciados factuais. Em sua visão, fazemos isso imaginando "mundos possíveis" nos quais os enunciados contrafactuais são verdadeiros. Quando dizemos "A dor de cabeça de José teria passado se ele tivesse tomado aspirina", estamos dizendo (segundo Lewis) que há outros mundos possíveis em que José tomou aspirina e sua dor de cabeça passou. Lewis argumentava que avaliamos contrafactuais comparando nosso mundo, em que ele não tomou aspirina, com o mundo \textit{mais semelhante} em que ele tomou aspirina. Ao não encontrar dor de cabeça naquele mundo, declaramos verdadeiro o enunciado contrafactual.

"Mais semelhante" é a chave. Pode haver alguns "mundos possíveis" em que sua dor de cabeça não passou — por exemplo, um mundo em que ele tomou a aspirina e depois bateu a cabeça na porta do banheiro. Mas esse mundo contém uma circunstância extra e adventícia. Entre todos os mundos possíveis em que José tomou aspirina, o mais semelhante ao nosso não seria aquele em que ele bateu a cabeça, e sim aquele em que sua dor de cabeça desapareceu.

Muitos críticos de Lewis atacaram a extravagância de suas afirmações sobre a existência literal de muitos outros mundos possíveis. "O sr. Lewis foi uma vez apelidado de 'realista modal cão-louco' por sua ideia de que qualquer mundo logicamente possível que se possa imaginar de fato existe", disse seu obituário no \textit{New York Times} em 2001. "Ele acreditava, por exemplo, que havia um mundo com jumentos falantes."

Mas penso que seus críticos (e talvez o próprio Lewis) perderam o ponto mais importante. Não precisamos argumentar se tais mundos existem como entidades físicas ou mesmo metafísicas. Se nosso objetivo é explicar o que as pessoas querem dizer ao afirmar "A causa B", basta postular que as pessoas são capazes de gerar mundos alternativos em suas cabeças, julgar qual mundo está "mais próximo" do nosso e, o mais importante, fazê-lo de forma coerente de modo a formar um consenso. Certamente não poderíamos comunicar-nos sobre contrafactuais se o "mais próximo" de uma pessoa fosse o "mais distante" de outra.

Nessa visão, o apelo de Lewis — "Por que não tomar os contrafactuais pelo valor de face?" — não pedia metafísica, e sim atenção à surpreendente uniformidade da arquitetura da mente humana. Como filósofo \textit{whig} licenciado, posso explicar essa consistência muito bem: ela decorre do fato de que experimentamos o mesmo mundo e compartilhamos o mesmo modelo mental de sua estrutura causal. Falamos disso já no Capítulo 1. Nossos modelos mentais compartilhados nos unem em comunidades. Podemos, portanto, julgar a proximidade não por alguma noção metafísica de "similaridade", e sim por quanto precisamos desmontar e perturbar nosso modelo compartilhado antes que ele satisfaça uma dada condição hipotética contrária aos fatos (José não tomar aspirina).

Nos modelos estruturais fazemos algo muito semelhante, embora enfeitado com mais detalhe matemático. Avaliamos expressões como "se \textit{X} tivesse sido \textit{x}" da mesma maneira que tratamos intervenções \textit{do}(\textit{X} = \textit{x}), apagando setas num diagrama causal ou equações num modelo estrutural. Podemos descrever isso como fazer a alteração mínima num diagrama causal necessária para garantir que \textit{X} seja igual a \textit{x}. Nesse aspecto, os contrafactuais estruturais são compatíveis com a ideia de Lewis do mundo possível mais semelhante.

Os modelos estruturais também oferecem uma resolução de um enigma sobre o qual Lewis se calou: como os humanos representam "mundos possíveis" em suas mentes e computam o mais próximo, quando o número de possibilidades está muito além da capacidade do cérebro humano? Cientistas da computação chamam isso de "problema da representação". Precisamos de algum código extremamente econômico para gerir tantos mundos. Poderiam os modelos estruturais, de alguma forma, ser o atalho real que usamos? Penso que é muito provável, por duas razões. Primeiro, os modelos causais estruturais são um atalho que funciona, e não há concorrentes por perto com essa propriedade milagrosa. Segundo, foram modelados sobre redes bayesianas, que por sua vez foram modeladas sobre a descrição de David Rumelhart da passagem de mensagens no cérebro. Não é exagero demais pensar que, há 40 mil anos, os humanos cooptaram a maquinaria em seu cérebro que já existia para o reconhecimento de padrões e começaram a usá-la para o raciocínio causal.

Os filósofos tendem a deixar para os psicólogos fazer afirmações sobre como a mente faz as coisas, o que explica por que as perguntas acima não foram abordadas até bem recentemente. No entanto, os pesquisadores de inteligência artificial (IA) não podiam esperar. Visavam construir robôs que pudessem comunicar-se com humanos sobre cenários alternativos, crédito e culpa, responsabilidade e arrependimento. Estas são todas noções contrafactuais que os pesquisadores de IA precisavam mecanizar antes de terem a menor chance de alcançar o que chamam de "IA forte" — inteligência semelhante à humana. Com essas motivações, entrei na análise contrafactual em 1994 (com meu aluno Alex Balke). Não surpreendentemente, a algoritmização dos contrafactuais fez mais barulho na inteligência artificial e na ciência cognitiva do que na filosofia.

Os filósofos tendiam a ver os modelos estruturais como apenas uma entre muitas implementações possíveis da lógica de mundos possíveis de Lewis. Ouso sugerir que são muito mais do que isso. Lógica vazia de representação é metafísica. Os diagramas causais, com suas regras simples de seguir e apagar setas, devem estar próximos do modo como nossos cérebros representam contrafactuais. Essa afirmação deve permanecer sem prova por enquanto, mas o resultado da longa história é que os contrafactuais deixaram de ser místicos. Compreendemos como os humanos os manejam, e estamos prontos para equipar robôs com capacidades semelhantes às que nossos ancestrais adquiriram há 40 mil anos.

\section{Resultados potenciais, equações estruturais e a algoritmização dos contrafactuais}
Apenas um ano após o lançamento do livro de Lewis, e independentemente dele, Donald Rubin (Figura 8.2) começou a escrever uma série de artigos que introduziram os resultados potenciais como linguagem para formular perguntas causais. Rubin, à época estatístico do Educational Testing Service, sozinho rompeu o silêncio sobre causalidade que persistira na estatística por setenta e cinco anos e legitimou o conceito de contrafactuais aos olhos de muitos cientistas da saúde. É impossível exagerar a importância desse desenvolvimento. Ele forneceu aos pesquisadores uma linguagem flexível para expressar quase toda pergunta causal que pudessem desejar formular, tanto no nível populacional quanto no individual.

\begin{figure}[ht]
\centering
\includegraphics[width=0.92\textwidth]{../figuras/fig_8_2.png}
\caption{Figura 8.2. Donald Rubin (à direita) com o autor em 2014. (Fonte: foto cortesia de Grace Hyun Kim.)}
\end{figure}
\textit{Figura 8.2. Donald Rubin (à direita) com o autor em 2014. (Fonte: foto cortesia de Grace Hyun Kim.)}

\textit{Figura 8.2. Donald Rubin (à direita) com o autor em 2014. (Fonte: foto cortesia de Grace Hyun Kim.)}

No modelo causal de Rubin, um resultado potencial de uma variável \textit{Y} é simplesmente "o valor que \textit{Y} teria tomado para o indivíduo \textit{u}, se \textit{X} tivesse sido atribuído o valor \textit{x}". São muitas palavras, de modo que com frequência é conveniente escrever essa quantidade de forma mais compacta como \textit{Y}<sub>\textit{X} = \textit{x}</sub>(\textit{u}). Muitas vezes abreviamos ainda mais como \textit{Y}<sub>\textit{x}</sub>(\textit{u}) se for aparente no contexto qual variável está sendo fixada no valor \textit{x}.

Para apreciar quão ousada é essa notação, é preciso afastar-se dos símbolos e pensar nas hipóteses que eles incorporam. Ao escrever o símbolo \textit{Y}<sub>\textit{x}</sub>, Rubin afirmava que \textit{Y} \textit{definitivamente} teria tomado algum valor se \textit{X} tivesse sido \textit{x}, e que isso tem tanta realidade objetiva quanto o valor que \textit{Y} de fato tomou. Se você não aceita essa hipótese (e estou bastante certo de que Heisenberg não a aceitaria), não pode usar resultados potenciais. Note também que o resultado potencial, ou contrafactual, é definido no nível de um indivíduo, não de uma população.

A primeira aparição científica de um resultado potencial veio na dissertação de mestrado de Jerzy Neyman, escrita em 1923. Neyman, descendente da nobreza polonesa, crescera no exílio na Rússia e só pôs os pés em sua terra natal em 1921, quando tinha vinte e sete anos. Recebera uma formação matemática muito sólida na Rússia e teria gostado de continuar a pesquisa em matemática pura, mas foi-lhe mais fácil encontrar emprego como estatístico. Muito como R. A. Fisher na Inglaterra, fez sua primeira pesquisa estatística num instituto agrícola — um emprego para o qual estava vastamente superqualificado. Não apenas era o único estatístico do instituto: era realmente a única pessoa no país pensando na estatística como disciplina. A primeira menção de Neyman a resultados potenciais veio no contexto de um experimento agrícola, em que a notação de subscrito representa o "rendimento potencial desconhecido da \textit{i}-ésima variedade [de uma dada semente] na respectiva parcela".

A dissertação permaneceu desconhecida e sem tradução para o inglês até 1990. No entanto, o próprio Neyman não permaneceu desconhecido. Arranjou passar um ano no laboratório estatístico de Karl Pearson no University College London, onde fez amizade com o filho de Pearson, Egon. Os dois mantiveram contato pelos sete anos seguintes, e sua colaboração deu grandes dividendos: a abordagem Neyman–Pearson de testes de hipóteses estatísticas foi um marco que todo estudante iniciante de estatística aprende. Em 1933, a longa liderança autocrática de Karl Pearson finalmente chegou ao fim com sua aposentadoria, e Egon era seu sucessor lógico — ou teria sido, se não fosse o singular problema de R. A. Fisher, então o estatístico mais famoso da Inglaterra. A universidade chegou a uma solução única e desastrosa, dividindo o cargo de Pearson numa cátedra de estatística (Egon Pearson) e numa cátedra de eugenia (Fisher). Egon não perdeu tempo em contratar seu amigo polonês. Neyman chegou em 1934 e quase imediatamente entrou em conflito com Fisher.

Fisher já estava sedento de briga. Sabia que era o principal estatístico do mundo e praticamente inventara grandes partes do assunto, e no entanto estava proibido de lecionar no departamento de estatística. As relações eram extraordinariamente tensas. "A sala comum era cuidadosamente compartilhada", escreve Constance Reid em sua biografia de Neyman. "O grupo de Pearson tomava chá às 4; e às 4h30, quando estavam bem longe, Fisher e seu grupo entravam em tropa."

Em 1935, Neyman deu uma palestra na Royal Statistical Society intitulada "Problemas estatísticos na experimentação agrícola", na qual questionou alguns dos próprios métodos de Fisher e, incidentalmente, discutiu a ideia de resultados potenciais. Depois que Neyman terminou, Fisher levantou-se e disse à sociedade que "esperava que o artigo do dr. Neyman fosse sobre um assunto com o qual o autor estivesse plenamente familiarizado". "[Neyman] afirmara que Fisher estava errado", escreveu Oscar Kempthorne anos depois sobre o incidente. "Isso era uma ofensa imperdoável — Fisher nunca estava errado e, de fato, a sugestão de que pudesse estar era tratada por ele como um assalto mortal. Qualquer um que não aceitasse os escritos de Fisher como a verdade dada por Deus era no mínimo estúpido e no máximo maligno."

Neyman e Pearson viram a extensão da fúria de Fisher poucos dias depois, quando foram ao departamento à noite e encontraram os modelos de madeira de Neyman, com os quais ilustrara a palestra, espalhados pelo chão. Concluíram que só Fisher poderia ter sido responsável pela devastação. Embora o ataque de raiva de Fisher possa parecer divertido agora, sua atitude teve consequências sérias. É claro que ele não podia engolir o orgulho e usar a notação de resultados potenciais de Neyman, ainda que isso o tivesse ajudado mais tarde com problemas de mediação. A falta de vocabulário de resultados potenciais levou a ele e a muitas outras pessoas a chamada Falácia da Mediação, que discutiremos no Capítulo 9.

Neste ponto, alguns leitores ainda podem achar o conceito de contrafactuais um tanto místico, de modo que gostaria de mostrar como alguns seguidores de Rubin infeririam resultados potenciais e contrastar essa abordagem livre de modelo com a abordagem do modelo causal estrutural. Suponha que estejamos examinando certa empresa para ver se a educação ou os anos de experiência são o fator mais importante na determinação do salário de um empregado. Coletamos alguns dados sobre os salários existentes nessa empresa, reproduzidos na Tabela 8.1. Deixamos \textit{EX} representar anos de experiência, \textit{ED} representar educação e \textit{S} representar salário. Assumimos também, por simplicidade, apenas três níveis de educação: 0 = diploma de ensino médio, 1 = diploma universitário, 2 = diploma de pós-graduação. Assim, \textit{S}<sub>\textit{ED} = 0</sub>(\textit{u}), ou \textit{S}<sub>0</sub>(\textit{u}), representa o salário do indivíduo \textit{u} se \textit{u} fosse apenas graduado no ensino médio, e \textit{S}<sub>1</sub>(\textit{u}) representa o salário de \textit{u} se \textit{u} fosse graduado na universidade. Uma pergunta contrafactual típica que poderíamos querer fazer é: "Qual seria o salário de Alice se ela tivesse um diploma universitário?" Em outras palavras, o que é \textit{S}<sub>1</sub>(Alice)?

\begin{figure}[ht]
\centering
\includegraphics[width=0.92\textwidth]{../figuras/tab_8_1.png}
\caption{Tabela 8.1. Dados fictícios para o exemplo de resultados potenciais.}
\end{figure}
\textit{Tabela 8.1. Dados fictícios para o exemplo de resultados potenciais.}

A primeira coisa a notar sobre a Tabela 8.1 é toda a ausência de dados, indicada por pontos de interrogação. Nunca podemos observar mais de um resultado potencial no mesmo indivíduo. Embora óbvia, essa afirmação é, não obstante, importante. O estatístico Paul Holland certa vez a chamou de "problema fundamental da inferência causal", nome que pegou. Se só pudéssemos preencher os pontos de interrogação, poderíamos responder a todas as nossas perguntas causais.

Nunca concordei com a caracterização de Holland dos valores ausentes na Tabela 8.1 como um "problema fundamental", talvez porque raramente tenha descrito problemas causais em termos de uma tabela. Mas, mais fundamentalmente, ver a inferência causal como um problema de dados ausentes pode ser terrivelmente enganoso, como veremos em breve. Observe que, além dos cabeçalhos decorativos das três últimas colunas, a Tabela 8.1 está totalmente desprovida de informação causal sobre \textit{ED}, \textit{EX} e \textit{S} — por exemplo, se a educação afeta o salário ou o contrário. Pior ainda: não nos permite representar tal informação quando disponível. Mas para estatísticos que percebem o "problema fundamental" como dados ausentes, tal tabela parece apresentar oportunidades infinitas. De fato, se olharmos para \textit{S}<sub>0</sub>, \textit{S}<sub>1</sub> e \textit{S}<sub>2</sub> não como resultados potenciais, e sim como variáveis ordinárias, temos dezenas de técnicas de interpolação para preencher os brancos ou, como diriam os estatísticos, "imputar os dados ausentes", de alguma maneira ótima.

\textit{Tabela 8.1. Dados fictícios para o exemplo de resultados potenciais.}

Uma abordagem comum é o \textit{pareamento} (\textit{matching}). Procuramos pares de indivíduos bem pareados em todas as variáveis exceto a de interesse e então preenchemos suas linhas para que combinem um com o outro. O caso mais claro aqui é o de Bert e Caroline, que combinam perfeitamente em experiência. Assim, assumimos que o salário de Bert, se ele tivesse um diploma de pós-graduação, seria o mesmo de Caroline (97.000 dólares), e o salário de Caroline, se ela tivesse apenas um diploma de graduação, seria o mesmo de Bert (92.500 dólares). Note que o pareamento invoca a mesma ideia que o condicionamento (ou estratificação): selecionamos para comparação grupos que compartilham uma característica observada e usamos a comparação para inferir características que parecem não compartilhar.

É difícil estimar o salário de Alice desse modo porque não há um bom pareamento para ela nos dados que dei. Não obstante, os estatísticos desenvolveram técnicas de considerável sutileza para imputar dados ausentes a partir de pareamentos aproximados, e Rubin tem sido pioneiro dessa abordagem. Infelizmente, mesmo o mais talentoso casamenteiro do mundo não pode transformar dados em resultados potenciais, nem sequer aproximadamente. Mostrarei abaixo que a resposta correta depende criticamente de se a educação afeta a experiência ou o contrário — informação que em lugar nenhum se encontra na tabela.

Um segundo método possível é a regressão linear (não confundir com equações estruturais). Nessa abordagem, fingimos que os dados vieram de alguma fonte aleatória desconhecida e usamos métodos estatísticos padrão para encontrar a reta (ou, neste caso, o plano) que melhor se ajusta aos dados. A saída de tal abordagem poderia ser uma equação que se parece com isto:

\textit{S} = 65.000 + 2.500 × \textit{EX} + 5.000 × \textit{ED}  (8.1)

A Equação 8.1 nos diz que (em média) o salário-base de um empregado sem experiência e apenas com diploma de ensino médio é 65.000 dólares. Para cada ano de experiência, o salário aumenta 2.500 dólares, e para cada grau educacional adicional (até dois), o salário aumenta 5.000 dólares. Consequentemente, um analista de regressão afirmaria, nossa estimativa do salário de Alice, se ela tivesse um diploma universitário, é 65.000 + 2.500 × 6 + 5.000 × 1 = 85.000 dólares.

A facilidade e a familiaridade de tais técnicas de imputação explicam por que a concepção de Rubin da inferência causal como um problema de dados ausentes desfrutou de ampla popularidade. Ai de nós: por mais inofensivos que esses métodos de interpolação pareçam, são fundamentalmente falhos. São guiados por dados, não por modelos. Todos os dados ausentes são preenchidos examinando outros valores na tabela. Como aprendemos com a Escada da Causalidade, qualquer método desse tipo está fadado ao fracasso desde o início; nenhum método baseado apenas em dados (degrau um) pode responder a perguntas contrafactuais (degrau três).

Antes de contrastar esses métodos com a abordagem do modelo causal estrutural, examinemos intuitivamente o que dá errado com a imputação cega a modelos. Em particular, expliquemos por que Bert e Caroline, que combinam perfeitamente em experiência, podem de fato ser bastante incomparáveis quando se trata de comparar seus resultados potenciais. Mais surpreendente ainda: uma história causal razoável (ajustada à Tabela 8.1) mostraria que o melhor pareamento para Caroline quanto ao Salário seria alguém que \textit{não} combina com ela em Experiência.

O primeiro ponto-chave a perceber é que a Experiência é provavelmente dependente da Educação. Afinal, aqueles empregados que obtiveram um grau educacional extra dedicaram quatro anos de suas vidas a isso. Assim, se Caroline tivesse apenas um grau de educação (como Bert), teria podido usar esse tempo extra para ganhar mais experiência do que tem agora. Isso lhe teria dado educação igual e experiência maior do que Bert. Podemos, portanto, concluir que \textit{S}<sub>1</sub>(Caroline) > \textit{S}<sub>1</sub>(Bert), contrariamente ao que o pareamento ingênuo preveria. Vemos que, uma vez que temos uma história causal em que a Educação afeta a Experiência, é inevitável que "parear" em Experiência crie um descompasso no Salário potencial. Ironicamente, a Experiência igual, que começou como um convite ao pareamento, agora se transformou num alto aviso \textit{contra} ele. A Tabela 8.1 continuará, é claro, em silêncio sobre tais perigos. Por essa razão, não posso compartilhar o entusiasmo de Holland por moldar a inferência causal como um problema de dados ausentes. Muito pelo contrário. Trabalho recente de Karthika Mohan, uma ex-aluna minha, revela que mesmo problemas padrão de dados ausentes exigem modelagem causal para sua solução.

Agora vejamos como um modelo causal estrutural trataria os mesmos dados. Primeiro, antes mesmo de olhar para os dados, desenhamos um diagrama causal (Figura 8.3). O diagrama codifica a história causal por trás dos dados, segundo a qual a Experiência "escuta" a Educação e o Salário escuta ambos. De fato, já podemos dizer algo muito importante só olhando para o diagrama. Se nosso modelo estivesse errado e \textit{EX} fosse uma causa de \textit{ED}, em vez do contrário, então a Experiência seria um confundidor, e parear empregados com experiência semelhante seria completamente apropriado. Com \textit{ED} como causa de \textit{EX}, a Experiência é um mediador. Como o leitor certamente já sabe, confundir um mediador com um confundidor é um dos pecados mais mortais na inferência causal e pode levar aos erros mais escandalosos. Este último convida ao ajustamento; o primeiro o proíbe.

\begin{figure}[ht]
\centering
\includegraphics[width=0.92\textwidth]{../figuras/fig_8_3.png}
\caption{Figura 8.3. Diagrama causal para o efeito da educação (\textit{ED}) e da experiência (\textit{EX}) sobre o salário (\textit{S}).}
\end{figure}
\textit{Figura 8.3. Diagrama causal para o efeito da educação (}ED\textit{) e da experiência (}EX\textit{) sobre o salário (}S\textit{).}

\textit{Figura 8.3. Diagrama causal para o efeito da educação (ED) e da experiência (EX) sobre o salário (S).}

Até agora neste livro usei uma palavra muito informal — "escutar" — para expressar o que quero dizer com as setas num diagrama causal. Mas agora é hora de pôr um pouco de carne matemática nesse conceito, e é de fato aqui que os modelos causais estruturais diferem das redes bayesianas ou dos modelos de regressão. Quando digo que o Salário escuta a Educação e a Experiência, quero dizer que ele é uma função matemática dessas variáveis: \textit{S} = \textit{f}<sub>\textit{S}</sub>(\textit{EX}, \textit{ED}). Mas precisamos permitir variações individuais, de modo que estendemos essa função para ler \textit{S} = \textit{f}<sub>\textit{S}</sub>(\textit{EX}, \textit{ED}, \textit{U}<sub>\textit{S}</sub>), onde \textit{U}<sub>\textit{S}</sub> representa "variáveis não observadas que afetam o salário". Sabemos que essas variáveis existem (por exemplo, Alice é amiga do presidente da empresa), mas são demasiado diversas e numerosas para incorporar explicitamente em nosso modelo.

Vejamos como isso se desenrolaria em nosso exemplo educação/experiência/salário, assumindo funções lineares em toda parte. Podemos usar os mesmos métodos estatísticos de antes para encontrar a equação linear de melhor ajuste. O resultado se pareceria com a Equação 8.1 com uma pequena diferença:

\textit{S} = 65.000 + 2.500 × \textit{EX} + 5.000 × \textit{ED} + \textit{U}<sub>\textit{S}</sub>  (8.2)

No entanto, a semelhança formal entre as Equações 8.1 e 8.2 é profundamente enganosa; suas interpretações diferem como o dia e a noite. O fato de termos escolhido regredir \textit{S} sobre \textit{ED} e \textit{EX} na Equação 8.1 de modo algum implica que \textit{S} escuta \textit{ED} e \textit{EX} no mundo real. Essa escolha foi puramente nossa, e nada nos dados nos impediria de regredir \textit{EX} sobre \textit{ED} e \textit{S} ou de seguir qualquer outra ordem. (Lembre-se da descoberta de Francis Galton no Capítulo 2 de que as regressões são cegas à causa.) Perdemos essa liberdade uma vez que proclamamos uma equação como "estrutural". Em outras palavras, o autor da Equação 8.2 deve comprometer-se a escrever equações que espelhem sua crença sobre quem escuta quem no mundo. Em nosso caso, ele acredita que \textit{S} de fato escuta \textit{EX} e \textit{ED}. Mais importante: a ausência de uma equação \textit{ED} = \textit{f}<sub>\textit{ED}</sub>(\textit{EX}, \textit{S}, \textit{U}<sub>\textit{ED}</sub>) do modelo significa que se acredita que \textit{ED} é indiferente a mudanças em \textit{EX} ou \textit{S}. Essa diferença de compromisso dá às equações estruturais o poder de sustentar contrafactuais — um poder negado às equações de regressão.

Em conformidade com a Figura 8.3, devemos também ter uma equação estrutural para \textit{EX}, mas agora forçaremos o coeficiente de \textit{S} a zero, para refletir a ausência de uma seta de \textit{S} para \textit{EX}. Uma vez que estimamos os coeficientes a partir dos dados, a equação poderia se parecer com algo assim:

\textit{EX} = 10 − 4 × \textit{ED} + \textit{U}<sub>\textit{EX}</sub>  (8.3)

Essa equação diz que a experiência média de pessoas sem graus avançados é de dez anos, e cada grau de educação (até dois) diminui \textit{EX} em quatro anos em média. Novamente, note a diferença-chave entre equações estruturais e de regressão: a variável \textit{S} não entra na Equação 8.3, apesar do fato de que \textit{S} e \textit{EX} são provavelmente altamente correlacionadas. Isso reflete a crença do analista de que a experiência \textit{EX} adquirida por qualquer indivíduo é totalmente não afetada por seu salário atual.

Agora demonstremos como derivar contrafactuais a partir de um modelo estrutural. Para estimar o salário de Alice se ela tivesse educação universitária, realizamos três passos:

1. \textbf{(Abdução)} Usar os dados sobre Alice e sobre os outros empregados para estimar os fatores idiossincráticos de Alice, \textit{U}<sub>\textit{S}</sub>(Alice) e \textit{U}<sub>\textit{EX}</sub>(Alice). 2. \textbf{(Ação)} Usar o operador \textit{do} para mudar o modelo de modo a refletir a hipótese contrafactual que se está fazendo — neste caso, que ela tem um diploma universitário: \textit{ED}(Alice) = 1. 3. \textbf{(Predição)} Calcular o novo salário de Alice usando o modelo modificado e a informação atualizada sobre as variáveis exógenas \textit{U}<sub>\textit{S}</sub>(Alice), \textit{U}<sub>\textit{EX}</sub>(Alice) e \textit{ED}(Alice).

Esse salário recém-calculado é igual a \textit{S}<sub>\textit{ED} = 1</sub>(Alice).

Para o passo 1, observamos a partir dos dados que \textit{EX}(Alice) = 6 e \textit{ED}(Alice) = 0. Substituímos esses valores nas Equações 8.2 e 8.3. As equações então nos dizem os fatores idiossincráticos de Alice: \textit{U}<sub>\textit{S}</sub>(Alice) = 1.000 dólares e \textit{U}<sub>\textit{EX}</sub>(Alice) = −4. Isso representa tudo o que é único, especial e maravilhoso em Alice. Seja o que for, acrescenta 1.000 dólares a seu salário previsto.

O passo 2 nos diz para usar o operador \textit{do} a fim de apagar as setas que apontam para a variável que está sendo fixada num valor contrafactual (Educação) e fixar a Educação de Alice num diploma universitário (Educação = 1). Neste exemplo, o passo 2 é trivial, porque não há setas apontando para Educação e, portanto, nenhuma seta a apagar. Em modelos mais complicados, porém, esse passo de apagar as setas não pode ser omitido, porque afeta o cálculo no passo 3. Variáveis que poderiam ter afetado o resultado por meio da variável intervenida não serão mais permitidas a fazê-lo.

Finalmente, o passo 3 diz para atualizar o modelo de modo a refletir a nova informação de que \textit{U}<sub>\textit{S}</sub> = 1.000 dólares, \textit{U}<sub>\textit{EX}</sub> = −4 e \textit{ED} = 1. Primeiro usamos a Equação 8.3 para recalcular o que seria a Experiência de Alice se ela tivesse ido à universidade: \textit{EX}<sub>\textit{ED} = 1</sub>(Alice) = 10 − 4 − 4 = 2 anos. Então usamos a Equação 8.2 para recalcular seu Salário potencial: \textit{S}<sub>\textit{ED} = 1</sub>(Alice) = 65.000 + 2.500 × 2 + 5.000 × 1 + 1.000 = 76.000 dólares.

Nosso resultado, \textit{S}<sub>1</sub>(Alice) = 76.000 dólares, é uma estimativa válida do salário hipotético de Alice; isto é, os dois coincidirão se as hipóteses do modelo forem válidas. Como este exemplo implica um modelo causal muito simples e funções muito simples (lineares), as diferenças entre ele e o método de regressão guiado por dados podem parecer bastante menores. Mas as diferenças menores na superfície refletem vastas diferenças por baixo. Qualquer resultado contrafactual (potencial) que obtemos do método estrutural segue logicamente das hipóteses exibidas no modelo, enquanto a resposta obtida pelo método guiado por dados é tão caprichosa quanto correlações espúrias, porque deixa hipóteses de modelagem importantes sem contabilização.

Este exemplo nos forçou a ir mais longe nos "parafusos e porcas" dos modelos causais do que tínhamos feito anteriormente neste livro. Mas deixe-me recuar um pouco para celebrar e apreciar o milagre que veio à existência por meio do exemplo de Alice. Usando uma combinação de dados e modelo, fomos capazes de prever o comportamento de um indivíduo (Alice) sob condições totalmente hipotéticas. É claro que não há almoço grátis: obtivemos esses resultados fortes porque fizemos hipóteses fortes. Além de afirmar as relações causais entre as variáveis observadas, também assumimos que as relações funcionais eram lineares. Mas a linearidade importa menos aqui do que saber quais são aquelas funções específicas. Isso nos capacitou a computar as idiossincrasias de Alice a partir de suas características observadas e atualizar o modelo como exigido no procedimento de três passos.

Ao risco de acrescentar uma nota sóbria à nossa celebração, tenho de dizer que essa informação funcional nem sempre estará disponível na prática. Em geral, chamamos um modelo de "completamente especificado" se as funções por trás das setas são conhecidas e de "parcialmente especificado" em caso contrário. Por exemplo, como nas redes bayesianas, podemos conhecer apenas relações probabilísticas entre pais e filhos no grafo. Se o modelo é parcialmente especificado, podemos não ser capazes de estimar o salário de Alice exatamente; em vez disso, podemos ter de fazer um enunciado de intervalo de probabilidade, como "Há de 10 a 20 por cento de chance de que seu salário seria 76.000 dólares". Mas mesmo tais respostas probabilísticas são boas o bastante para muitas aplicações. Além disso, é verdadeiramente notável quanta informação podemos extrair do diagrama causal mesmo quando não temos informação sobre as funções específicas por trás das setas, ou apenas informação muito geral, como a hipótese de "monotonicidade" que encontramos no último capítulo.

Os passos 1 a 3 acima podem ser resumidos no que chamo de "primeira lei da inferência causal": \textit{Y}<sub>\textit{x}</sub>(\textit{u}) = \textit{Y}<sub>\textit{M}<sub>\textit{x}</sub></sub>(\textit{u}). Esta é a mesma regra que usamos no exemplo do pelotão de fuzilamento no Capítulo 1, exceto que as funções são diferentes. A primeira lei diz que o resultado potencial \textit{Y}<sub>\textit{x}</sub>(\textit{u}) pode ser imputado indo ao modelo \textit{M}<sub>\textit{x}</sub> (com setas para \textit{X} apagadas) e computando o resultado \textit{Y}(\textit{u}) ali. Todas as quantidades estimáveis nos degraus dois e três da Escada da Causalidade seguem daí. Em suma, a redução dos contrafactuais a um algoritmo nos permite conquistar tanto território do degrau três quanto a matemática permitir — mas, é claro, nem um pouco mais.

\section{A virtude de ver as próprias hipóteses}
O método MCE que mostrei para computar contrafactuais não é o mesmo método que Rubin usaria. Um ponto maior de diferença entre nós é o uso de diagramas causais. Eles permitem que os pesquisadores representem hipóteses causais em termos que podem compreender e então tratem todos os contrafactuais como propriedades derivadas de seu modelo de mundo. O modelo causal de Rubin trata os contrafactuais como objetos matemáticos abstratos que são manejados por maquinaria algébrica, mas não derivados de um modelo.

Privado de um recurso gráfico, o usuário do modelo causal de Rubin é usualmente pedido a aceitar três hipóteses. A primeira, chamada "hipótese de valor estável de tratamento unitário", ou SUTVA (do inglês \textit{stable unit treatment value assumption}), é razoavelmente transparente. Diz que cada indivíduo (ou "unidade", o termo preferido dos modeladores causais) terá o mesmo efeito do tratamento independentemente de que tratamento os outros indivíduos (ou "unidades") recebam. Em muitos casos, excluindo epidemias e outras interações coletivas, isso faz perfeito sentido. Por exemplo, assumindo que a dor de cabeça não é contagiosa, minha resposta à aspirina não dependerá de se José recebe aspirina.

A segunda hipótese no modelo de Rubin, também benigna, é chamada "consistência". Diz que uma pessoa que tomou aspirina e se recuperou também se recuperaria se recebesse aspirina por desenho experimental. Essa hipótese razoável, que é um teorema no arcabouço MCE, diz em efeito que o experimento está livre de efeitos placebo e de outras imperfeições.

Mas a principal hipótese que os praticantes de resultados potenciais são invariavelmente obrigados a fazer é chamada "ignorabilidade". É mais técnica, mas é a parte crucial da transação, pois é em essência a mesma coisa que a condição de trocabilidade (\textit{exchangeability}) de Jamie Robins e Sander Greenland discutida no Capítulo 4. A ignorabilidade expressa esse mesmo requisito em termos da variável de resultado potencial \textit{Y}<sub>\textit{x}</sub>. Exige que \textit{Y}<sub>\textit{x}</sub> seja independente do tratamento realmente recebido, a saber \textit{X}, dados os valores de um certo conjunto de variáveis (des)confundidoras \textit{Z}.

Antes de explorar sua interpretação, devemos reconhecer que qualquer hipótese expressa como independência condicional herda um grande corpo de maquinaria matemática familiar desenvolvida por estatísticos para variáveis ordinárias (não contrafactuais). Por exemplo, os estatísticos rotineiramente usam regras para decidir quando uma independência condicional se segue de outra. Para crédito de Rubin, ele reconheceu as vantagens de traduzir a noção causal de "não-confundimento" para a sintaxe da teoria da probabilidade, ainda que sobre variáveis contrafactuais. A hipótese de ignorabilidade faz do modelo causal de Rubin de fato um \textit{modelo}; a Tabela 8.1 em si não é um modelo porque não contém hipóteses sobre o mundo.

Infelizmente, ainda não encontrei uma única pessoa que possa explicar o que a ignorabilidade significa numa linguagem falada por aqueles que precisam fazer essa hipótese ou avaliar sua plausibilidade num dado problema. Aqui vai minha melhor tentativa. A atribuição de pacientes a tratamento ou controle é ignorável se, dentro de qualquer estrato do confundidor \textit{Z}, pacientes que teriam um resultado potencial, \textit{Y}<sub>\textit{x}</sub> = \textit{y}, são tão prováveis de estar no grupo de tratamento ou controle quanto os pacientes que teriam um resultado potencial diferente, \textit{Y}<sub>\textit{x}</sub> = \textit{y′}. Essa definição é perfeitamente legítima para alguém na posse de uma função de probabilidade sobre contrafactuais. Mas como um biólogo ou economista, com apenas conhecimento científico como guia, deve avaliar se isso é verdadeiro ou não? Mais concretamente: como um cientista deve avaliar se a ignorabilidade vale em qualquer dos exemplos discutidos neste livro?

Para compreender a dificuldade, tentemos aplicar essa explicação ao nosso exemplo. Para determinar se \textit{ED} é ignorável (condicional em \textit{EX}), devemos julgar se empregados que teriam um salário potencial, digamos \textit{S}<sub>1</sub> = \textit{s}, são tão prováveis de ter um nível de educação quanto os empregados que teriam um salário potencial diferente, digamos \textit{S}<sub>1</sub> = \textit{s′}. Se você acha que isso soa circular, só posso concordar! Queremos determinar o salário potencial de Alice e, mesmo antes de começar — mesmo antes de termos uma pista sobre a resposta —, devemos especular se o resultado é dependente ou independente de \textit{ED}, em cada estrato de \textit{EX}. É um verdadeiro pesadelo cognitivo.

Como se verifica, \textit{ED} em nosso exemplo \textit{não} é ignorável em relação a \textit{S}, condicional em \textit{EX}, e é por isso que a abordagem de pareamento (igualando Bert e Caroline) renderia a resposta errada para seus salários potenciais. De fato, suas estimativas deveriam diferir por um montante \textit{S}<sub>1</sub>(Bert) − \textit{S}<sub>1</sub>(Caroline) = 5.000 dólares. (O leitor deve ser capaz de mostrar isso a partir dos números da Tabela 8.1 e do procedimento de três passos.)

Mostrarei agora que, com a ajuda de um diagrama causal, um estudante poderia ver imediatamente que \textit{ED} não é ignorável e não tentaria o pareamento aqui. Sem um diagrama, um estudante seria tentado a assumir que a ignorabilidade vale por padrão e cairia nessa armadilha. (Isso não é especulação. Peguei emprestada a ideia deste exemplo de um artigo na \textit{Harvard Law Review} em que a história era essencialmente a mesma da Figura 8.3 e o autor de fato usou pareamento.)

Eis como podemos usar um diagrama causal para testar a ignorabilidade (condicional). Para determinar se \textit{X} é ignorável em relação ao resultado \textit{Y}, condicional num conjunto \textit{Z} de variáveis de pareamento, precisamos apenas testar se \textit{Z} bloqueia todos os caminhos da porta dos fundos entre \textit{X} e \textit{Y} e se nenhum membro de \textit{Z} é descendente de \textit{X}. É tão simples assim! Em nosso exemplo, a variável de pareamento proposta (Experiência) bloqueia todos os caminhos da porta dos fundos (porque não há nenhum), mas falha no teste porque é descendente de Educação. Portanto \textit{ED} não é ignorável, e \textit{EX} não pode ser usado para pareamento. Não são necessárias ginásticas mentais elaboradas — apenas um olhar para um diagrama. Nunca se exige de um pesquisador que avalie mentalmente quão provável é um resultado potencial dado um tratamento ou outro.

Infelizmente, Rubin não considera que os diagramas causais "ajudem a inferência causal". Portanto, pesquisadores que seguem seu conselho serão privados desse teste de ignorabilidade e terão de realizar formidáveis ginásticas mentais para convencer-se de que a hipótese vale, ou então simplesmente aceitar a hipótese como uma "caixa-preta". De fato, um proeminente pesquisador de resultados potenciais, Marshall Joffe, escreveu em 2010 que as hipóteses de ignorabilidade são usualmente feitas porque justificam o uso de métodos estatísticos disponíveis, não porque sejam verdadeiramente acreditadas.

Estreitamente relacionada à transparência é a noção de testabilidade, que surgiu várias vezes neste livro. Um modelo lançado como diagrama causal pode facilmente ser testado quanto à compatibilidade com os dados, ao passo que um modelo lançado na linguagem de resultados potenciais carece dessa característica. O teste é este: sempre que todos os caminhos entre \textit{X} e \textit{Y} no diagrama são bloqueados por um conjunto de nós \textit{Z}, então nos dados \textit{X} e \textit{Y} devem ser independentes, condicionais em \textit{Z}. Esta é a propriedade de \textit{d}-separação mencionada no Capítulo 7, que nos permite rejeitar um modelo sempre que a independência deixa de aparecer nos dados. Em contraste, se o mesmo modelo é expresso na linguagem de resultados potenciais (isto é, como uma coleção de enunciados de ignorabilidade), carecemos da maquinaria matemática para desvelar as independências que o modelo implica, e os pesquisadores são incapazes de submeter o modelo a um teste. É difícil compreender como os pesquisadores de resultados potenciais conseguiram viver com essa deficiência sem se rebelar. Minha única explicação é que foram mantidos longe de ferramentas gráficas por tanto tempo que esqueceram que modelos causais podem e devem ser testáveis.

Agora devo aplicar os mesmos padrões de transparência a mim mesmo e dizer um pouco mais sobre as hipóteses incorporadas num modelo causal estrutural. Lembra-se da história de Abraão que relatei antes? A primeira resposta de Abraão à notícia da iminente destruição de Sodoma foi procurar uma relação dose-resposta, ou uma função de resposta, relacionando a maldade da cidade a seu castigo. Foi um instinto científico sólido, mas suspeito que poucos de nós teriam estado calmos o bastante para reagir assim. A função de resposta é o ingrediente-chave que dá aos MCE o poder de lidar com contrafactuais. Ela é implícita no paradigma de resultados potenciais de Rubin, mas um ponto maior de diferença entre os MCE e as redes bayesianas, incluindo redes bayesianas causais.

Numa rede bayesiana probabilística, as setas para \textit{Y} significam que a probabilidade de \textit{Y} é governada pelas tabelas de probabilidade condicional para \textit{Y}, dadas observações de suas variáveis-pai. O mesmo vale para redes bayesianas causais, exceto que as tabelas de probabilidade condicional especificam a probabilidade de \textit{Y} dadas intervenções nas variáveis-pai. Ambos os modelos especificam probabilidades para \textit{Y}, não um valor específico de \textit{Y}.

Num modelo causal estrutural, não há tabelas de probabilidade condicional. As setas simplesmente significam que \textit{Y} é uma função de seus pais, bem como da variável exógena \textit{U}<sub>\textit{Y}</sub>:

\textit{Y} = \textit{f}<sub>\textit{Y}</sub>(\textit{X}, \textit{A}, \textit{B}, \textit{C}, …, \textit{U}<sub>\textit{Y}</sub>)  (8.4)

Assim, o instinto de Abraão era sólido. Para transformar uma rede bayesiana não causal num modelo causal — ou, mais precisamente, para torná-la capaz de responder a consultas contrafactuais —, precisamos de uma relação dose-resposta em cada nó. Essa realização não me veio facilmente. Mesmo antes de aprofundar-me nos contrafactuais, tentei por muito tempo formular modelos causais usando tabelas de probabilidade condicional. Um obstáculo que enfrentei foram os modelos cíclicos, que eram totalmente resistentes a formulações de probabilidade condicional. Outro obstáculo foi o de inventar uma notação para distinguir redes bayesianas probabilísticas das causais. Em 1991, de repente me ocorreu que todas as dificuldades desapareceriam se tornássemos \textit{Y} uma função de suas variáveis-pai e deixássemos o termo \textit{U}<sub>\textit{Y}</sub> lidar com todas as incertezas concernentes a \textit{Y}.

Na época, parecia uma heresia contra meu próprio ensinamento. Depois de dedicar vários anos à causa das probabilidades na inteligência artificial, estava agora propondo dar um passo atrás e usar um modelo não probabilístico, quase determinístico. Ainda me lembro de meu aluno à época, Danny Geiger, perguntando incredulamente: "Equações determinísticas? Verdadeiramente determinísticas?" Era como se Steve Jobs acabasse de lhe dizer para comprar um PC em vez de um Mac. (Isso era 1990!)

Na superfície, não havia nada de revolucionário nessas equações. Economistas e sociólogos vinham usando tais modelos desde os anos 1950 e 1960 e chamando-os de modelos de equações estruturais (SEM, do inglês \textit{structural equation models}). Mas esse nome sinalizava controvérsia e confusão sobre a interpretação causal das equações. Com o tempo, os economistas perderam de vista o fato de que os pioneiros desses modelos, Trygve Haavelmo em economia e Otis Dudley Duncan em sociologia, os tinham destinado a representar relações causais. Começaram a confundir equações estruturais com retas de regressão, assim esvaziando a substância da forma. Por exemplo, em 1988, quando David Freedman desafiou onze pesquisadores de SEM a explicar como aplicar intervenções a um modelo de equações estruturais, nenhum deles pôde. Podiam dizer como estimar os coeficientes a partir dos dados, mas não podiam dizer por que alguém deveria se dar ao trabalho. Se a interpretação de função de resposta que apresentei entre 1990 e 1994 fez algo de novo, foi simplesmente restaurar e formalizar as intenções originais de Haavelmo e Duncan e colocar diante de seus discípulos as ousadas conclusões que se seguem dessas intenções se as levarmos a sério. Algumas dessas conclusões seriam consideradas assombrosas, mesmo por Haavelmo e Duncan. Tome por exemplo a ideia de que de todo SEM, por mais simples que seja, podemos computar todos os contrafactuais que se possa imaginar entre as variáveis do modelo. Nossa capacidade de computar o salário potencial de Alice, se ela tivesse educação universitária, seguiu-se dessa ideia. Mesmo hoje, economistas contemporâneos não internalizaram essa ideia.

Uma outra diferença importante entre SEM e MCE, além da letra do meio, é que a relação entre causas e efeitos num MCE não é necessariamente linear. As técnicas que emergem da análise de MCE são válidas para funções não lineares bem como lineares, variáveis discretas bem como contínuas. Modelos lineares de equações estruturais têm muitas vantagens e muitas desvantagens. Do ponto de vista da metodologia, são sedutoramente simples. Podem ser estimados a partir de dados observacionais por regressão linear, e você pode escolher entre dezenas de pacotes de \textit{software} estatístico para fazer isso por você. Por outro lado, modelos lineares não podem representar curvas dose-resposta que não sejam retas. Não podem representar efeitos de limiar, como um medicamento que tem efeitos crescentes até certa dosagem e então nenhum efeito adicional. Também não podem representar interações entre variáveis. Por exemplo, um modelo linear não pode descrever uma situação em que uma variável realça ou inibe o efeito de outra variável. (Por exemplo, a Educação poderia realçar o efeito da Experiência colocando o indivíduo num emprego de carreira mais rápida que obtém aumentos anuais maiores.)

Embora debates sobre as hipóteses apropriadas a fazer sejam inevitáveis, nossa mensagem principal é bastante simples: Alegrem-se! Com um modelo causal estrutural plenamente especificado, que implica um diagrama causal e todas as funções por trás dele, podemos responder a qualquer consulta contrafactual. Mesmo com um MCE parcial, em que algumas variáveis estão ocultas ou as relações dose-resposta são desconhecidas, ainda assim em muitos casos podemos responder à nossa consulta. As duas seções seguintes dão alguns exemplos.

\section{Contrafactuais e o direito}
Em princípio, os contrafactuais deveriam encontrar fácil aplicação no tribunal. Digo "em princípio" porque a profissão jurídica é muito conservadora e leva muito tempo a aceitar novos métodos matemáticos. Mas usar contrafactuais como modo de argumentação é de fato muito antigo e conhecido na profissão jurídica como "causalidade \textit{but-for}". O Código Penal Modelo (\textit{Model Penal Code}) expressa o teste \textit{but-for} da seguinte forma: "A conduta é a causa de um resultado quando: (a) é um antecedente sem o qual o resultado em questão não teria ocorrido."

Se o réu disparou uma arma e a bala atingiu e matou a vítima, o disparo da arma é uma causa \textit{but-for}, ou necessária, da morte, pois a vítima estaria viva se não fosse o disparo. Causas \textit{but-for} também podem ser indiretas. Se José bloqueia a saída de incêndio de um prédio com móveis, e Judite morre num incêndio depois de não conseguir alcançar a saída, então José é legalmente responsável por sua morte ainda que não tenha ateado o fogo.

Como podemos expressar causas necessárias ou \textit{but-for} em termos de resultados potenciais? Se deixamos o resultado \textit{Y} ser "a morte de Judite" (com \textit{Y} = 0 se Judite vive e \textit{Y} = 1 se Judite morre) e o tratamento \textit{X} ser "o bloqueio da saída de incêndio por José" (com \textit{X} = 0 se ele não a bloqueia e \textit{X} = 1 se bloqueia), então somos instruídos a fazer a seguinte pergunta: dado que sabemos que a saída de incêndio estava bloqueada (\textit{X} = 1) e Judite morreu (\textit{Y} = 1), qual é a probabilidade de que Judite estaria viva (\textit{Y} = 0) se \textit{X} tivesse sido 0? Simbolicamente, a probabilidade que queremos avaliar é

\textit{P}(\textit{Y}<sub>\textit{X} = 0</sub> = 0 | \textit{X} = 1, \textit{Y} = 1).

Como essa expressão é bastante pesada, mais tarde a abreviarei como "PN", a probabilidade de necessidade (isto é, a probabilidade de que \textit{X} = 1 seja uma causa necessária ou \textit{but-for} de \textit{Y} = 1). Note que a probabilidade de necessidade envolve um contraste entre dois mundos diferentes: o mundo real em que \textit{X} = 1 e o mundo contrafactual em que \textit{X} = 0 (expresso pelo subscrito \textit{X} = 0). De fato, a retrospectiva (saber o que aconteceu no mundo real) é uma distinção crítica entre contrafactuais (degrau três da Escada da Causalidade) e intervenções (degrau dois). Sem retrospectiva, não há diferença entre \textit{P}(\textit{Y}<sub>\textit{X} = 0</sub> = 0) e \textit{P}(\textit{Y} = 0 | \textit{do}(\textit{X} = 0)). Ambos expressam a probabilidade de que, sob condições normais, Judite estará viva se assegurarmos que a saída não está bloqueada; não mencionam o fogo, a morte de Judite ou a saída bloqueada. Mas a retrospectiva pode mudar nossa estimativa das probabilidades. Suponha que observemos que \textit{X} = 1 e \textit{Y} = 1 (retrospectiva). Então \textit{P}(\textit{Y}<sub>\textit{X} = 0</sub> = 0 | \textit{X} = 1, \textit{Y} = 1) não é o mesmo que \textit{P}(\textit{Y}<sub>\textit{X} = 0</sub> = 0 | \textit{X} = 1). Saber que Judite morreu (\textit{Y} = 1) nos dá informação sobre as circunstâncias que não obteríamos só por saber que a porta estava bloqueada (\textit{X} = 1). Por uma coisa, é evidência da força do fogo. De fato, pode-se mostrar que não há maneira de capturar \textit{P}(\textit{Y}<sub>\textit{X} = 0</sub> = 0 | \textit{X} = 1, \textit{Y} = 1) numa expressão \textit{do}. Embora isso possa parecer um ponto bastante obscuro, ele dá prova matemática de que os contrafactuais (degrau três) estão acima das intervenções (degrau dois) na Escada da Causalidade.

Nas últimas poucas parágrafos, quase sub-repticiamente introduzimos probabilidades em nossa discussão. Advogados há muito compreendem que a certeza matemática é um padrão de prova demasiado alto. Para casos criminais nos Estados Unidos, a Suprema Corte em 1880 estabeleceu que a culpa tem de ser provada "com exclusão de toda dúvida razoável". A corte não disse "além de toda dúvida" ou "além da sombra de uma dúvida", e sim além de dúvida razoável. A Suprema Corte nunca deu uma definição precisa desse termo, mas pode-se conjecturar que há algum limiar, talvez 99 por cento ou 99,9 por cento de probabilidade de culpa, acima do qual a dúvida se torna irrazoável e é do interesse da sociedade encarcerar o réu. Em processos cíveis, em vez de criminais, o padrão de prova é um tanto mais claro. A lei exige uma "preponderância de evidência" de que o réu causou a lesão, e parece razoável interpretar isso como significando que a probabilidade é maior que 50 por cento.

Embora a causalidade \textit{but-for} seja geralmente aceita, os advogados reconheceram que em alguns casos ela pode levar a um erro judicial. Um exemplo clássico é o cenário do "piano que cai", em que o réu dispara um tiro na vítima e erra, e no processo de fugir da cena, a vítima por acaso corre sob um piano que cai e é morta. Pelo teste \textit{but-for} o réu seria culpado de assassinato, porque a vítima não estaria em lugar nenhum perto do piano que caía se não estivesse fugindo. Mas nossa intuição diz que o réu não é culpado de assassinato (embora possa ser culpado de tentativa de assassinato), porque não havia maneira de ele ter antecipado o piano que caía. Um advogado diria que o piano, e não o tiro, é a \textit{causa próxima} da morte.

A doutrina da causa próxima é muito mais obscura do que a causa \textit{but-for}. O Código Penal Modelo diz que o resultado não deve ser "demasiado remoto ou acidental em sua ocorrência para ter uma [justa] incidência sobre a responsabilidade do ator ou a gravidade de sua ofensa". No presente, essa determinação é deixada à intuição do juiz. Eu sugeriria que é uma forma de causa suficiente. A ação do réu foi suficiente para trazer, com probabilidade alta o bastante, o evento que de fato causou a morte?

Enquanto o significado de causa próxima é muito vago, o significado de causa suficiente é bastante preciso. Usando notação contrafactual, podemos definir a probabilidade de suficiência, ou PS, como \textit{P}(\textit{Y}<sub>\textit{X} = 1</sub> = 1 | \textit{X} = 0, \textit{Y} = 0). Isso nos diz para imaginar uma situação em que \textit{X} = 0 e \textit{Y} = 0: o atirador não disparou na vítima, e a vítima não correu sob um piano. Então perguntamos quão provável é que, numa tal situação, disparar o tiro (\textit{X} = 1) resultaria no resultado \textit{Y} = 1 (correr sob um piano)? Isso exige juízo contrafactual, mas penso que a maioria de nós concordaria que a probabilidade de tal resultado seria extremamente pequena. Tanto a intuição quanto o Código Penal Modelo sugerem que se PS for demasiado pequeno, não devemos condenar o réu por causar \textit{Y} = 1.

Como a distinção entre causas necessárias e suficientes é tão importante, penso que pode ajudar ancorar esses dois conceitos em exemplos simples. A causa suficiente é a mais comum das duas, e já encontramos esse conceito no exemplo do pelotão de fuzilamento do Capítulo 1. Ali, o disparo do Soldado A ou do Soldado B é suficiente para causar a morte do prisioneiro, e nenhum (em si) é necessário. Assim, PS = 1 e PN = 0. As coisas ficam um pouco mais interessantes quando a incerteza entra em cena — por exemplo, se cada soldado tem alguma probabilidade de desobedecer a ordens ou errar o alvo. Por exemplo, se o Soldado A tem uma probabilidade \textit{p}<sub>A</sub> de errar o alvo, então seu PS seria 1 − \textit{p}<sub>A</sub>, pois esta é sua probabilidade de acertar o alvo e causar a morte. Seu PN, porém, dependeria de quão provável é que o Soldado B se abstenha de atirar ou erre o alvo. Só sob tais circunstâncias o disparo do Soldado A seria necessário; isto é, o prisioneiro estaria vivo se o Soldado A não tivesse atirado.

Um exemplo clássico que demonstra a causalidade necessária conta a história de um fogo que irrompeu depois que alguém riscou um fósforo, e a pergunta é "O que causou o fogo, riscar o fósforo ou a presença de oxigênio na sala?". Note que ambos os fatores são igualmente necessários, pois o fogo não teria ocorrido na ausência de um deles. Assim, de um ponto de vista puramente lógico, os dois fatores são igualmente responsáveis pelo fogo. Por que, então, consideramos acender o fósforo uma explicação mais razoável do fogo do que a presença de oxigênio?

Para responder, considere as duas frases:

1. A casa ainda estaria de pé se apenas o fósforo não tivesse sido riscado. 2. A casa ainda estaria de pé se apenas o oxigênio não tivesse estado presente.

Ambas as frases são verdadeiras. No entanto, a esmagadora maioria dos leitores, tenho certeza, viria com o primeiro cenário se pedida a explicar o que causou a casa a queimar — o fósforo ou o oxigênio. Então, o que explica a diferença? A resposta claramente tem algo a ver com a normalidade: ter oxigênio na casa é bastante normal, mas dificilmente podemos dizer o mesmo de riscar um fósforo. A diferença não aparece na lógica, mas aparece nas duas medidas que discutimos acima, PS e PN. Se levarmos em conta que a probabilidade de riscar um fósforo é muito menor do que a de ter oxigênio, encontramos quantitativamente que para Fósforo tanto PN quanto PS são altos, enquanto para Oxigênio PN é alto mas PS é baixo. É por isso, intuitivamente, que culpamos o fósforo e não o oxigênio? Bastante possivelmente, mas pode ser apenas parte da resposta.

Em 1982, os psicólogos Daniel Kahneman e Amos Tversky investigaram como as pessoas escolhem um culpado "se apenas" para "desfazer" um resultado indesejado e encontraram padrões consistentes em suas escolhas. Um foi que as pessoas são mais propensas a imaginar desfazer um evento raro do que um comum. Por exemplo, se estamos desfazendo um encontro perdido, é mais provável que digamos "Se apenas o trem tivesse saído no horário" do que "Se apenas o trem tivesse saído mais cedo". Outro padrão foi a tendência das pessoas a culpar suas próprias ações (por exemplo, riscar um fósforo) em vez de eventos fora de seu controle. Nossa capacidade de estimar PN e PS a partir de nosso modelo do mundo sugere uma maneira sistemática de contabilizar essas considerações e, eventualmente, ensinar robôs a produzir explicações significativas de eventos peculiares.

Vimos que PN captura a racionalidade por trás do critério \textit{but-for} num cenário jurídico. Mas PS deveria entrar em considerações legais no direito penal e civil? Creio que sim, porque a atenção à suficiência implica atenção às consequências da própria ação. A pessoa que acendeu o fósforo deveria ter antecipado a presença de oxigênio, ao passo que ninguém é geralmente esperado a bombear todo o oxigênio para fora da casa em antecipação a uma cerimônia de riscar fósforos. Que peso, então, a lei deveria atribuir aos componentes necessários versus suficientes da causalidade? Filósofos do direito não discutiram o estatuto legal dessa pergunta, talvez porque as noções de PS e PN não tenham sido formalizadas com tal precisão. No entanto, de uma perspectiva de IA, claramente PN e PS deveriam tomar parte na geração de explicações. Um robô instruído a explicar por que um fogo irrompeu não tem escolha senão considerar ambos. Focar apenas em PN renderia a conclusão insustentável de que riscar um fósforo e ter oxigênio são explicações igualmente adequadas para o fogo. Um robô que emite esse tipo de explicação rapidamente perderá a confiança de seu dono.

\section{Causas necessárias, causas suficientes e mudança climática}
Em agosto de 2003, a onda de calor mais intensa em cinco séculos atingiu a Europa ocidental, concentrando seus efeitos mais severos na França. O governo francês culpou a onda de calor por quase 15 mil mortes, muitas delas entre idosos que viviam sozinhos e não tinham ar-condicionado. Foram eles vítimas do aquecimento global ou de má sorte — de viver no lugar errado na hora errada?

Antes de 2003, os cientistas do clima evitavam especular sobre tais perguntas. A sabedoria convencional era algo como isto: "Embora este seja o tipo de fenômeno que o aquecimento global poderia tornar mais frequente, é impossível atribuir este evento particular a emissões passadas de gases de efeito estufa." Myles Allen, físico da Universidade de Oxford e autor da citação acima, sugeriu uma maneira de fazer melhor: usar uma métrica chamada fração de risco atribuível (FAR, do inglês \textit{fraction of attributable risk}) para quantificar o efeito da mudança climática. A FAR exige que conheçamos dois números: \textit{p}<sub>0</sub>, a probabilidade de uma onda de calor como a de 2003 \textit{antes} da mudança climática (por exemplo, antes de 1800), e \textit{p}<sub>1</sub>, a probabilidade \textit{depois} da mudança climática. Por exemplo, se a probabilidade dobra, então podemos dizer que metade do risco se deve à mudança climática. Se triplica, então dois terços do risco se devem à mudança climática.

Como a FAR é definida puramente a partir de dados, ela não tem necessariamente nenhum significado causal. Acontece, porém, que sob duas hipóteses causais brandas, ela é idêntica à probabilidade de necessidade. Primeiro, precisamos assumir que o tratamento (gases de efeito estufa) e o resultado (ondas de calor) não são confundidos: não há causa comum de cada um. Isso é muito razoável, porque até onde sabemos a única causa do aumento dos gases de efeito estufa somos nós mesmos. Segundo, precisamos assumir monotonicidade. Discutimos essa hipótese brevemente no último capítulo; neste contexto, significa que o tratamento nunca tem o efeito oposto do que esperamos: isto é, os gases de efeito estufa nunca podem nos proteger de uma onda de calor. Contanto que as hipóteses de ausência de confundimento e de ausência de proteção valham, a métrica de degrau um da FAR é promovida ao degrau três, onde se torna PN.

Mas Allen não conhecia a interpretação causal da FAR — ela provavelmente não é conhecimento comum entre meteorologistas — e isso o forçou a apresentar seus resultados usando linguagem um tanto tortuosa. Mas que dados podemos usar para estimar a FAR (ou PN)? Observamos apenas uma tal onda de calor. Não podemos fazer um experimento controlado, porque isso exigiria que controlássemos o nível de dióxido de carbono como se estivéssemos acionando um interruptor. Felizmente, os cientistas do clima têm uma arma secreta: podem conduzir um experimento \textit{in silico} — uma simulação computacional. Allen e Peter Stott do Met Office (o serviço meteorológico britânico) aceitaram o desafio e, em 2004, tornaram-se os primeiros cientistas do clima a comprometer-se com um enunciado causal sobre um evento meteorológico individual. Ou tornaram-se? Julgue você mesmo. Eis o que escreveram: "É muito provável que mais da metade do risco de anomalias de temperatura de verão europeias excederem um limiar de 1,6 °C seja atribuível à influência humana."

Embora eu elogie a bravura de Allen e Stott, é uma pena que sua importante descoberta tenha sido enterrada numa tal espessura de linguagem impenetrável. Deixe-me desempacotar esse enunciado e então tentar explicar por que tiveram de expressá-lo de forma tão convoluta. Primeiro, "anomalia de temperatura excedendo um limiar de 1,6 °C" era sua maneira de definir o resultado. Escolheram esse limiar porque a temperatura média na Europa naquele verão foi mais de 1,6 °C acima do normal, o que nunca havia acontecido antes na história registrada. Sua escolha equilibrava os objetivos concorrentes de escolher um resultado suficientemente extremo para capturar o efeito do aquecimento global, mas não demasiado ajustado aos detalhes específicos do evento de 2003. Em vez de usar, por exemplo, a temperatura média na França durante agosto, escolheram o critério mais amplo da temperatura média na Europa ao longo de todo o verão.

Em seguida, o que queriam dizer com "muito provável" e "metade do risco"? Em termos matemáticos, Allen e Stott queriam dizer que havia 90 por cento de chance de que a FAR fosse superior a 50 por cento. Ou, equivalentemente, há 90 por cento de chance de que verões como o de 2003 sejam mais do que duas vezes tão prováveis com os níveis atuais de dióxido de carbono quanto seriam com níveis pré-industriais. Note que há duas camadas de probabilidade aqui: estamos falando de uma probabilidade de uma probabilidade! Não é de admirar que nossa mente se confunda e nossos olhos nadem quando lemos tal enunciado. A razão do golpe duplo é que a onda de calor está sujeita a dois tipos de incerteza. Primeiro, há incerteza sobre a quantidade de mudança climática de longo prazo. Esta é a incerteza que entra na primeira figura de 90 por cento. Mesmo se soubéssemos a quantidade de mudança climática de longo prazo exatamente, há incerteza sobre o tempo em qualquer ano dado. Esse é o tipo de variabilidade embutido na fração de 50 por cento de risco atribuível. Assim, temos de conceder que Allen e Stott estavam tentando comunicar uma ideia complicada. Não obstante, uma coisa está ausente de sua conclusão: a causalidade. Seu enunciado não contém nem sequer uma pista de causalidade — ou talvez apenas uma pista, na frase ambígua e inescrutável "atribuível à influência humana".

Agora compare isso com uma versão causal da mesma conclusão: "As emissões de CO<sub>2</sub> são muito provavelmente uma causa necessária da onda de calor de 2003." Qual frase, a deles ou a nossa, você ainda lembrará amanhã? Qual você poderia explicar ao seu vizinho da porta ao lado?

Não sou pessoalmente um especialista em mudança climática, de modo que peguei este exemplo de um de meus colaboradores, Alexis Hannart do Instituto Franco-Argentino de Estudo do Clima e Seus Impactos em Buenos Aires, que tem sido um grande proponente da análise causal na ciência do clima. Hannart desenha o grafo causal na Figura 8.4. Como Gases de Efeito Estufa é um nó de nível superior no modelo climático, sem setas entrando nele, ele argumenta que não há confundimento entre ele e a Resposta Climática. Da mesma forma, ele atesta a hipótese de ausência de proteção (isto é, gases de efeito estufa não podem nos proteger de ondas de calor).

\begin{figure}[ht]
\centering
\includegraphics[width=0.92\textwidth]{../figuras/fig_8_4.png}
\caption{Figura 8.4. Diagrama causal para o exemplo da mudança climática.}
\end{figure}
\textit{Figura 8.4. Diagrama causal para o exemplo da mudança climática.}

\textit{Figura 8.4. Diagrama causal para o exemplo da mudança climática.}

Hannart vai além de Allen e Stott e usa nossas fórmulas para computar a probabilidade de suficiência (PS) e de necessidade (PN). No caso da onda de calor europeia de 2003, ele encontra que PS era extremamente baixo, cerca de 0,0072, significando que não havia maneira de prever que esse evento aconteceria neste ano particular. Por outro lado, a probabilidade de necessidade PN era 0,9, em acordo com os resultados de Allen e Stott. Isso significa que é altamente provável que, sem gases de efeito estufa, a onda de calor não teria acontecido.

O valor aparentemente baixo de PS tem de ser posto num contexto maior. Não queremos apenas saber a probabilidade de uma onda de calor \textit{este ano}; gostaríamos de saber a probabilidade de uma recorrência de uma onda de calor tão severa ao longo de um horizonte de tempo mais longo — digamos nos próximos dez ou cinquenta anos. À medida que o horizonte de tempo se alonga, PN diminui porque outros mecanismos possíveis para uma onda de calor poderiam entrar em jogo. Por outro lado, PS aumenta porque estamos em efeito dando aos dados mais chances de saírem olhos de serpente. Assim, por exemplo, Hannart computa que há 80 por cento de probabilidade de que a mudança climática seja uma causa suficiente de outra onda de calor europeia como a de 2003 (ou pior) num período de duzentos anos. Isso pode não soar demasiado aterrorizante, mas isso assume os níveis de gases de efeito estufa de \textit{hoje}. Na realidade, os níveis de CO<sub>2</sub> certamente continuarão a subir, o que só pode aumentar PS e encurtar a janela de tempo até a próxima onda de calor.

Pessoas comuns podem aprender a compreender a diferença entre causas necessárias e suficientes? Esta é uma pergunta não trivial. Mesmo cientistas às vezes lutam. De fato, dois estudos conflitantes saíram que analisaram a onda de calor de 2010 na Rússia, quando a Rússia teve seu verão mais quente de todos os tempos e fogos de turfa escureceram os céus de Moscou. Um grupo concluiu que a variabilidade natural causou a onda de calor; outro concluiu que a mudança climática a causou. Com toda probabilidade, o desacordo ocorreu porque os dois grupos definiram seu resultado de forma diferente. Um grupo aparentemente baseou seu argumento em PN e obteve alta probabilidade de que a mudança climática era a causa, enquanto o outro usou PS e obteve baixa probabilidade. O segundo grupo atribuiu a onda de calor a um padrão persistente de alta pressão ou "bloqueio" sobre a Rússia — o que me soa como uma causa suficiente — e encontrou que os gases de efeito estufa tinham pouco a ver com esse fenômeno. Mas qualquer estudo que use PS como métrica, ao longo de um período curto, está estabelecendo uma barra alta para provar causalidade.

Antes de deixar este exemplo, gostaria de comentar novamente sobre os modelos computacionais. A maioria dos outros cientistas tem de trabalhar muito duro para obter informação contrafactual — por exemplo, combinando penosamente dados de estudos observacionais e experimentais. Cientistas do clima podem obter contrafactuais muito facilmente de seus modelos computacionais: basta digitar um novo número para a concentração de dióxido de carbono e deixar o programa rodar. "Facilmente" é, é claro, relativo. Por trás do simples diagrama causal da Figura 8.4 jaz uma função de resposta fabulosamente complexa, dada pelos milhões de linhas de código de computador que entram numa simulação climática. Isso levanta uma pergunta natural: quanto podemos confiar nas simulações computacionais? A pergunta tem ramificações políticas, especialmente aqui nos Estados Unidos. No entanto, tentarei dar uma resposta apolítica. Consideraria a função de resposta neste exemplo muito mais credível do que os modelos lineares que se veem tão frequentemente nas ciências naturais e sociais. Modelos lineares são frequentemente escolhidos por nenhuma boa razão além da conveniência. Em comparação, os modelos climáticos refletem mais de um século de estudo por físicos, meteorologistas e cientistas do clima. Representam os melhores esforços de uma comunidade de cientistas para compreender os processos que governam nosso tempo e clima. Por quaisquer padrões científicos normais, os modelos climáticos são evidência forte e convincente, mas com uma ressalva. Embora sejam excelentes em prever o tempo alguns dias à frente, nunca foram verificados num ensaio prospectivo ao longo de escalas de tempo seculares, de modo que ainda poderiam conter erros sistemáticos de que não sabemos.

\section{Um mundo de contrafactuais}
Espero que até agora seja óbvio que os contrafactuais são uma parte essencial de como os humanos aprendem sobre o mundo e de como nossas ações o afetam. Embora nunca possamos caminhar pelos dois caminhos que divergem num bosque, em um grande número de casos podemos saber, com algum grau de confiança, o que jaz em cada um.

Além de qualquer dúvida, a variedade e a riqueza das consultas causais que podemos colocar a nosso "motor de inferência" são grandemente realçadas quando podemos incluir contrafactuais na mistura. Outro tipo muito popular de consulta, que não discuti aqui, chamado efeito do tratamento nos tratados (ETT, do inglês \textit{effect of treatment on the treated}), é usado para avaliar se as pessoas que ganham acesso a um tratamento são aquelas que mais se beneficiariam dele. Essa medida é em muitos casos superior à medida convencional da eficácia de um tratamento, o efeito causal médio (ACE, do inglês \textit{average causal effect}). O ACE, que se pode obter de um ensaio clínico randomizado, media a eficácia do tratamento sobre toda a população. Mas e se, na implementação real, aqueles recrutados para um programa de tratamento forem os menos propensos a se beneficiar dele? Para avaliar a eficácia global do programa, o ETT mede quão adversamente os pacientes tratados seriam afetados se não tivessem sido tratados — uma medida contrafactual de significância crítica na tomada de decisão prática. Meu ex-aluno Ilya Shpitser (agora na Johns Hopkins) fez agora para o ETT o que o cálculo do \textit{do} fez para o ACE — forneceu uma compreensão completa de quando ele é estimável a partir de dados, dado um diagrama causal.

Sem dúvida a aplicação mais popular dos contrafactuais na ciência hoje é chamada análise de mediação. Por essa razão, dedico a ela um capítulo separado (Capítulo 9). Curiosamente, muitas pessoas, especialmente se usam técnicas clássicas de análise de mediação, podem não perceber que estão falando de um efeito contrafactual. Num contexto científico, um mediador, ou variável mediadora, é uma que transmite o efeito do tratamento ao resultado. Vimos muitos exemplos de mediação neste livro, como Fumar → Alcatrão → Câncer (onde Alcatrão é o mediador). A principal pergunta de interesse em tais casos é se a variável mediadora responde por todo o efeito da variável de tratamento ou se alguma parte do efeito não exige um mediador. Representaríamos tal efeito por uma seta separada levando diretamente do tratamento ao resultado, como Fumar → Câncer. A análise de mediação visa desembaraçar o efeito direto (que não passa pelo mediador) do efeito indireto (a parte que passa pelo mediador). A importância é fácil de ver. Se fumar causa câncer de pulmão apenas pela formação de depósitos de alcatrão, então poderíamos eliminar o excesso de risco de câncer dando aos fumantes cigarros sem alcatrão, como cigarros eletrônicos. Por outro lado, se fumar causa câncer diretamente ou por meio de um mediador diferente, então os cigarros eletrônicos podem não resolver o problema. No presente essa pergunta médica está sem resolução.

Neste ponto provavelmente não é óbvio para você que efeitos diretos e indiretos envolvem enunciados contrafactuais. Definitivamente não era óbvio para mim! De fato, foi uma das maiores surpresas de minha carreira. O próximo capítulo conta essa história e dá muitas aplicações da análise de mediação na vida real.

\chapter{Capítulo 9 — Mediação: a busca de um mecanismo}
\textit{Por falta de um prego, a ferradura se perdeu. Por falta de uma ferradura, o cavalo se perdeu.… Por falta de uma batalha, o reino se perdeu. E tudo por falta de um prego.} — ANÔNIMO

Na linguagem ordinária, a pergunta “Por quê?” tem pelo menos duas versões. A primeira é direta: você vê um efeito e quer saber a causa. Seu avô está deitado no hospital, e você pergunta: “Por quê? Como ele pode ter tido um infarto se parecia tão saudável?” Mas há uma segunda versão da pergunta “Por quê?”, que formulamos quando queremos compreender melhor a conexão entre uma causa conhecida e um efeito conhecido. Por exemplo, observamos que o Medicamento B previne infartos. Ou, como James Lind, observamos que as frutas cítricas previnem o escorbuto. A mente humana é inquieta e sempre quer saber mais. Em pouco tempo começamos a fazer a segunda versão da pergunta: “Por quê? Qual é o mecanismo pelo qual as frutas cítricas previnem o escorbuto?” Este capítulo concentra-se nessa segunda versão do “porquê”.

A busca de mecanismos é crítica para a ciência, bem como para a vida cotidiana, porque mecanismos diferentes exigem ações diferentes quando as circunstâncias mudam. Suponha que acabem as laranjas. Conhecendo o mecanismo pelo qual as laranjas funcionam, ainda podemos prevenir o escorbuto. Precisamos simplesmente de outra fonte de vitamina C. Se não conhecêssemos o mecanismo, poderíamos ser tentados a experimentar bananas.

A palavra que os cientistas usam para o segundo tipo de pergunta “Por quê?” é “mediação”. Você pode ler em uma revista científica um enunciado como este: “O efeito do Medicamento B sobre os infartos é mediado pelo seu efeito sobre a pressão arterial.” Esse enunciado codifica um modelo causal simples: Medicamento → Pressão Arterial → Infarto. Nesse caso, o medicamento reduz a hipertensão, que por sua vez reduz o risco de infarto. (Os biólogos tipicamente usam um símbolo diferente, A —| B, quando a causa A inibe o efeito B, mas na literatura da causalidade é costume usar A → B tanto para causas positivas quanto negativas.) Do mesmo modo, podemos resumir o efeito das frutas cítricas sobre o escorbuto pelo modelo causal Frutas Cítricas → Vitamina C → Escorbuto.

Queremos fazer certas perguntas típicas sobre um mediador: Ele explica o efeito inteiro? O Medicamento B funciona exclusivamente por meio da pressão arterial, ou talvez também por outros mecanismos? O efeito placebo é um tipo comum de mediador na medicina: se um medicamento age apenas por meio da crença do paciente no seu benefício, a maioria dos médicos o considerará ineficaz. A mediação também é um conceito importante no direito. Se perguntamos se uma empresa discriminou mulheres ao pagar-lhes salários mais baixos, estamos formulando uma pergunta de mediação. A resposta depende de se a disparidade salarial observada é produzida diretamente em resposta ao sexo da candidata ou indiretamente, por meio de um mediador como a qualificação, sobre o qual o empregador não tem controle.

Todas as perguntas acima exigem uma sensibilidade delicada para distinguir efeitos totais, efeitos diretos (que não passam por um mediador) e efeitos indiretos (que passam). Mesmo definir esses termos foi um grande desafio para os cientistas ao longo do século passado. Inibidos pelos tabus contra pronunciar a palavra “causalidade”, alguns tentaram definir a mediação usando um vocabulário livre de causalidade. Outros descartaram de todo a análise de mediação e declararam os conceitos de efeitos diretos e indiretos “mais enganosos do que úteis para o pensamento estatístico claro”.

Para mim, também, a mediação foi uma luta — no fim, uma das mais recompensadoras da minha carreira, porque eu estava errado no começo e, ao aprender com o meu erro, cheguei a uma solução inesperada. Por um tempo, fui da opinião de que os efeitos indiretos não têm implicações operacionais porque, ao contrário dos efeitos diretos, não podem ser definidos na linguagem das intervenções. Foi um avanço pessoal quando percebi que eles \textit{podem} ser definidos em termos de contrafactuais e que também podem ter importantes implicações de política. Só podem ser quantificados depois de alcançarmos o terceiro degrau da Escada da Causalidade, e é por isso que os coloquei no final deste livro.

A mediação floresceu em seu novo habitat e nos permitiu quantificar, muitas vezes a partir dos meros dados, a porção do efeito mediada por qualquer caminho desejado. Compreensivelmente, devido ao seu traje contrafactual, os efeitos indiretos permanecem um tanto enigmáticos mesmo entre os campeões da Revolução Causal. Creio, contudo, que a sua utilidade avassaladora acabará por vencer quaisquer dúvidas remanescentes sobre a metafísica dos contrafactuais. Talvez possam ser comparados aos números irracionais e imaginários: a princípio deixaram as pessoas desconfortáveis (daí o nome “irracional”), mas por fim a sua utilidade transformou o desconforto em deleite.

Para ilustrar esse ponto, darei vários exemplos de como pesquisadores em diversas disciplinas extraíram insights úteis da análise de mediação. Uma pesquisadora estudou uma reforma educacional chamada “Álgebra para Todos”, que a princípio parecia um fracasso, mas depois se transformou em sucesso. Um estudo do uso de torniquetes nas guerras do Iraque e do Afeganistão não conseguiu mostrar que eles tinham algum benefício; a análise cuidadosa de mediação explica por que o benefício pode ter sido mascarado no estudo. Em resumo, nos últimos quinze anos a Revolução Causal descobriu regras claras e simples para quantificar quanto de um dado efeito é direto e quanto é indireto. Transformou a mediação de um conceito mal compreendido, de legitimidade duvidosa, em uma ferramenta popular e amplamente aplicável de análise científica.

\section{Escorbuto: o mediador errado}
Gostaria de começar com um exemplo histórico verdadeiramente estarrecedor que destaca a importância de compreender o mediador. Um dos primeiros exemplos de experimento controlado foi o estudo do escorbuto pelo capitão de navio James Lind, publicado em 1747. No tempo de Lind, o escorbuto era uma doença aterradora, estimada em ter matado 2 milhões de marinheiros entre 1500 e 1800. Lind estabeleceu, de forma tão conclusiva quanto alguém podia na época, que uma dieta de frutas cítricas prevenia os marinheiros de desenvolver essa temível doença. No início dos anos 1800, o escorbuto se tornara um problema do passado para a marinha britânica, pois todos os seus navios partiam para o mar com um suprimento adequado de frutas cítricas.

É geralmente nesse ponto que os livros de história encerram a narrativa, celebrando um grande triunfo do método científico. Parece então muito surpreendente que essa doença completamente prevenível tenha feito um retorno inesperado um século depois, quando expedições britânicas começaram a explorar as regiões polares. A Expedição Ártica Britânica de 1875, a Expedição Jackson-Harmsworth ao Ártico em 1894 e, de modo mais notável, as duas expedições de Robert Falcon Scott à Antártida em 1903 e 1911 sofreram todas gravemente de escorbuto.

Como isso pôde acontecer? Em duas palavras: ignorância e arrogância — sempre uma combinação potente. Por volta de 1900, os principais médicos da Grã-Bretanha haviam esquecido as lições de um século antes. O médico de Scott na expedição de 1903, o Dr. Reginald Koettlitz, atribuía o escorbuto a carne contaminada. Mais ainda, acrescentava: “o benefício dos chamados ‘antiescorbúticos’ [isto é, preventivos do escorbuto, como o suco de limão] é uma ilusão.” Em sua expedição de 1911, Scott estocou carne seca que havia sido escrupulosamente inspecionada em busca de sinais de decomposição, mas nenhuma fruta cítrica ou suco (veja a Figura 9.1). A confiança que depositou na opinião do médico pode ter contribuído para a tragédia que se seguiu. Todos os cinco homens que chegaram ao Polo Sul morreram, dois de uma doença não especificada que muito provavelmente era escorbuto. Um membro da equipe voltou antes do polo e chegou vivo, mas com um caso grave de escorbuto.

\begin{figure}[ht]
\centering
\includegraphics[width=0.92\textwidth]{../figuras/fig_9_1.png}
\caption{Figura 9.1. ). A confiança que depositou na opinião do médico pode ter contribuído para a tragédia que se seguiu. Todos os cinco homens que chegaram ao Polo Sul morreram, dois de uma doença não especificada que muito provavelmente era escorbuto. Um membro da equipe voltou antes do polo e chegou vivo, mas com um caso grave de escorbuto.}
\end{figure}
\textit{Figura 9.1. ). A confiança que depositou na opinião do médico pode ter contribuído para a tragédia que se seguiu. Todos os cinco homens que chegaram ao Polo Sul morreram, dois de uma doença não especificada que muito provavelmente era escorbuto. Um membro da equipe voltou antes do polo e chegou vivo, mas com um caso grave de escorbuto.}

Com o benefício da retrospectiva, o conselho de Koettlitz beira a imperícia criminosa. Como a lição de James Lind pôde ser tão completamente esquecida — ou pior, descartada — um século depois? A explicação, em parte, é que os médicos não compreendiam de fato como as frutas cítricas agiam contra o escorbuto. Em outras palavras, não conheciam o mediador.

\textbf{FIGURA 9.1.} Rações diárias dos homens na caminhada de Scott até o polo: chocolate, \textit{pemmican} (um prato de carne preservada), açúcar, biscoitos, manteiga, chá. Conspicuamente ausente: qualquer fruta contendo vitamina C. (Fonte: Fotografia de Herbert Ponting, cortesia do Canterbury Museum, Nova Zelândia.)

Desde o tempo de Lind, sempre se acreditara (mas nunca se provara) que as frutas cítricas preveniam o escorbuto em virtude de sua acidez. Em outras palavras, os médicos compreendiam o processo como governado pelo seguinte diagrama causal:

Frutas Cítricas → Acidez → Escorbuto

Desse ponto de vista, qualquer ácido serviria. Até Coca-Cola funcionaria (embora ainda não tivesse sido inventada). A princípio os marinheiros usavam limões espanhóis; depois, por razões econômicas, substituíram-nos por limas das Índias Ocidentais, que eram tão ácidas quanto os limões espanhóis, mas continham apenas um quarto da vitamina C. Para piorar as coisas, começaram a “purificar” o suco de lima cozinhando-o, o que pode ter decomposto a vitamina C que ainda restava. Em outras palavras, estavam desabilitando o mediador.

Quando os marinheiros da expedição ártica de 1875 adoeceram de escorbuto apesar de tomarem suco de lima, a comunidade médica foi lançada em completa confusão. Aqueles marinheiros que tinham comido carne recém-abatida não pegaram escorbuto, ao passo que os que tinham comido carne enlatada pegaram. Koettlitz e outros culparam a carne inadequadamente preservada como a culpada. Sir Almroth Wright elaborou uma teoria de que bactérias na carne (supostamente) contaminada causavam “envenenamento por ptomaina”, que então levava ao escorbuto. Enquanto isso, a teoria de que as frutas cítricas podiam prevenir o escorbuto foi relegada à lata de lixo.

A situação só se esclareceu quando o verdadeiro mediador foi descoberto. Em 1912, um bioquímico polonês chamado Casimir Funk propôs a existência de micronutrientes que chamou de “vitamines” (o \textit{e} era intencional). Por volta de 1930, Albert Szent-Györgyi isolou o nutriente particular que prevenia o escorbuto. Não era um ácido qualquer, mas um ácido em particular, agora conhecido como vitamina C ou ácido ascórbico (uma reverência ao seu passado “antiescorbútico”). Szent-Györgyi recebeu o Prêmio Nobel por sua descoberta em 1937. Graças a Szent-Györgyi, conhecemos agora o verdadeiro caminho causal:

Frutas Cítricas → Vitamina C → Escorbuto.

Penso que é justo prever que os cientistas nunca mais “esquecerão” esse caminho causal. E penso que o leitor concordará que a análise de mediação é mais do que um exercício matemático abstrato.

\section{Natureza versus criação: a tragédia de Barbara Burks}
Até onde sei, a primeira pessoa a representar explicitamente um mediador com um diagrama foi uma estudante de pós-graduação de Stanford chamada Barbara Burks, em 1926. Essa pioneira muito pouco conhecida da ciência das mulheres é uma das verdadeiras heroínas deste livro. Há razão para crer que ela de fato inventou os diagramas de caminhos independentemente de Sewall Wright. E, no que diz respeito à mediação, estava à frente de Wright e décadas à frente de seu tempo.

O principal interesse de pesquisa de Burks, ao longo de sua carreira infelizmente breve, foi o papel da natureza versus a criação na determinação da inteligência humana. Seu orientador em Stanford era Lewis Terman, um psicólogo famoso por desenvolver o teste de QI Stanford-Binet e um firme crente de que a inteligência era herdada, não adquirida. Tenha em mente que essa era a época áurea do movimento eugênico, hoje desacreditado, mas naquele tempo legitimado pela pesquisa ativa de pessoas como Francis Galton, Karl Pearson e Terman.

O debate natureza-versus-criação é, naturalmente, um debate muito antigo que continuou muito depois de Burks. Sua contribuição única foi reduzi-lo a um diagrama causal (veja a Figura 9.2), que ela usou para perguntar (e responder) à consulta: “Quanto do efeito causal se deve ao caminho direto Inteligência Parental → Inteligência da Criança (natureza), e quanto se deve ao caminho indireto Inteligência Parental → Status Social → Inteligência da Criança (criação)?”

\begin{figure}[ht]
\centering
\includegraphics[width=0.92\textwidth]{../figuras/fig_9_2.png}
\caption{Figura 9.2. ), que ela usou para perguntar (e responder) à consulta: “Quanto do efeito causal se deve ao caminho direto Inteligência Parental → Inteligência da Criança (natureza), e quanto se deve ao caminho indireto Inteligência Parental → Status Social → Inteligência da Criança (criação)?”}
\end{figure}
\textit{Figura 9.2. ), que ela usou para perguntar (e responder) à consulta: “Quanto do efeito causal se deve ao caminho direto Inteligência Parental → Inteligência da Criança (natureza), e quanto se deve ao caminho indireto Inteligência Parental → Status Social → Inteligência da Criança (criação)?”}

Nesse diagrama, Burks usou algumas setas de duas pontas, seja para representar causalidade mútua, seja simplesmente por incerteza quanto à direção da causalidade. Por simplicidade, vamos assumir que o efeito principal de ambas as setas vai da esquerda para a direita, o que torna Status Social um mediador, de modo que a inteligência dos pais eleva a sua posição social, e isso por sua vez dá à criança uma melhor oportunidade de desenvolver a sua inteligência. A variável \textit{X} representa “outras causas remotas não medidas”.

\textbf{FIGURA 9.2.} O debate natureza-versus-criação, conforme enquadrado por Barbara Burks.

Em sua dissertação, Burks coletou dados de extensas visitas domiciliares a 204 famílias com crianças adotivas, que presumivelmente receberiam apenas os benefícios da criação e nenhum dos benefícios da natureza de seus pais adotivos (veja a Figura 9.3). Aplicou testes de QI a todos eles e a um grupo de controle de 105 famílias sem crianças adotivas. Além disso, administrou questionários que usou para avaliar vários aspectos do ambiente social da criança. Usando seus dados e a análise de caminhos, calculou o efeito direto do QI parental sobre o QI das crianças e descobriu que apenas 35 por cento, ou cerca de um terço, da variação do QI é herdada. Em outras palavras, pais com um QI quinze pontos acima da média tipicamente teriam filhos cinco pontos acima da média.

\begin{figure}[ht]
\centering
\includegraphics[width=0.92\textwidth]{../figuras/fig_9_3.png}
\caption{Figura 9.3. ). Aplicou testes de QI a todos eles e a um grupo de controle de 105 famílias sem crianças adotivas. Além disso, administrou questionários que usou para avaliar vários aspectos do ambiente social da criança. Usando seus dados e a análise de caminhos, calculou o efeito direto do QI parental sobre o QI das crianças e descobriu que apenas 35 por cento, ou cerca de um terço, da variação do QI é herdada. Em outras palavras, pais com um QI quinze pontos acima da média tipicamente teriam filhos cinco pontos acima da média.}
\end{figure}
\textit{Figura 9.3. ). Aplicou testes de QI a todos eles e a um grupo de controle de 105 famílias sem crianças adotivas. Além disso, administrou questionários que usou para avaliar vários aspectos do ambiente social da criança. Usando seus dados e a análise de caminhos, calculou o efeito direto do QI parental sobre o QI das crianças e descobriu que apenas 35 por cento, ou cerca de um terço, da variação do QI é herdada. Em outras palavras, pais com um QI quinze pontos acima da média tipicamente teriam filhos cinco pontos acima da média.}

\textbf{FIGURA 9.3.} Barbara Burks (à direita) interessava-se em separar os componentes “natureza” e “criação” da inteligência. Como estudante de pós-graduação, visitou as casas de mais de duzentas crianças adotivas, aplicou-lhes testes de QI e coletou dados sobre o seu ambiente social. Foi a primeira pesquisadora além de Sewall Wright a usar diagramas de caminhos e, de certos modos, antecipou Wright. (Fonte: Desenho de Dakota Harr.)

Como discípula de Terman, Burks deve ter ficado decepcionada ao ver um efeito tão pequeno. (De fato, as suas estimativas se sustentaram bastante bem ao longo do tempo.) Por isso questionou o método de análise então aceito, que era controlar por Status Social. “A verdadeira medida da contribuição de uma causa a um efeito é mutilada”, escreveu ela, “se tornamos constantes variáveis que podem em parte ou no todo ser causadas por qualquer um dos dois fatores cuja verdadeira relação se pretende medir, ou por ainda outras causas remotas não medidas que também afetam qualquer um dos dois fatores isolados” (ênfase no original). Em outras palavras, se você está interessado no efeito total da Inteligência Parental sobre a Inteligência da Criança, não deve ajustar por (tornar constante) qualquer variável no caminho entre elas.

Mas Burks não parou aí. O seu critério em itálico, traduzido para a linguagem moderna, diz que um viés será introduzido se condicionarmos a variáveis que são (a) efeitos seja da Inteligência Parental, seja da Inteligência da Criança, ou (b) efeitos de causas não medidas seja da Inteligência Parental, seja da Inteligência da Criança (como \textit{X} na Figura 9.2).

Esses critérios estavam muito à frente de seu tempo e diferiam de tudo o que Sewall Wright havia escrito. De fato, o critério (b) é um dos primeiros exemplos de todos os tempos de viés de colisor. Se olharmos para a Figura 9.2, vemos que Status Social é um colisor (Inteligência Parental → Status Social ← \textit{X}). Portanto, controlar por Status Social abre o caminho pela porta dos fundos Inteligência Parental → Status Social ← \textit{X} → Inteligência da Criança. Qualquer estimativa resultante dos efeitos indireto e direto estará enviesada.

Como os estatísticos antes (e depois) de Burks não pensavam em termos de setas e diagramas, estavam totalmente imersos no mito de que, embora a correlação simples não tenha implicações causais, a correlação controlada (ou os coeficientes de regressão parciais; veja p. 222) é um passo na direção da explicação causal. Burks não foi a primeira pessoa a descobrir o efeito de colisor, mas pode-se argumentar que foi a primeira a caracterizá-lo de modo geral em termos gráficos. O seu critério (b) aplica-se perfeitamente aos exemplos de viés-M no Capítulo 4. O dela é o primeiro alerta de todos os tempos contra condicionar a um fator pré-tratamento, um hábito considerado seguro por todos os estatísticos do século XX e, estranhamente, ainda considerado seguro por alguns.

Agora coloque-se no lugar de Barbara Burks. Você acaba de descobrir que todos os seus colegas têm controlado pelas variáveis erradas. Você tem duas desvantagens: é apenas uma estudante e é mulher. O que faz? Baixa a cabeça, finge aceitar a sabedoria convencional e se comunica com os colegas no vocabulário inadequado deles? Não Barbara Burks! Intitulou o seu primeiro artigo publicado “On the Inadequacy of the Partial and Multiple Correlation Technique” [“Sobre a inadequação da técnica de correlação parcial e múltipla”] e começou dizendo: “Considerações lógicas levam à conclusão de que as técnicas de correlação parcial e múltipla estão repletas de perigos que restringem seriamente a sua aplicabilidade.” Palavras de combate de alguém que ainda não tem um doutorado! Como Terman escreveu: “A sua capacidade era de algum modo temperada pela sua tendência a irritar as pessoas. Penso que o problema residia em parte no fato de que ela era mais agressiva em defender as próprias ideias do que muitos professores e estudantes de pós-graduação homens gostavam.” Evidentemente, Burks estava à frente de seu tempo de mais de uma maneira.

Burks pode de fato ter inventado os diagramas de caminhos independentemente de Sewall Wright, que a precedeu por apenas seis anos. Podemos dizer com certeza que ela não os aprendeu em nenhuma aula. A Figura 9.2 é a primeira aparição de um diagrama de caminhos fora do trabalho de Sewall Wright e a primeira de todas nas ciências sociais ou comportamentais. É verdade que ela credita Wright no final do seu artigo de 1926, mas o faz de um modo que parece um acréscimo de última hora. Tenho o palpite de que ela só ficou sabendo dos diagramas de Wright depois de ter desenhado os seus, possivelmente depois de ter sido alertada por Terman ou por um revisor astuto.

É fascinante imaginar o que Burks poderia ter se tornado, se não tivesse sido vítima de sua época. Depois de obter o doutorado, nunca conseguiu um emprego como professora em uma universidade, para o qual certamente estava qualificada. Teve de contentar-se com posições de pesquisa menos seguras, por exemplo no Carnegie Institution. Em 1942 ficou noiva, o que se poderia esperar que marcasse uma melhora em sua fortuna; em vez disso, entrou em uma depressão profunda. “Estou convencida de que, certa ou não, ela tinha certeza de que alguma mudança sinistra estava ocorrendo em seu cérebro, da qual nunca poderia se recuperar”, escreveu a mãe dela, Frances Burks, a Terman. “Assim, no mais terno amor por todos nós, escolheu poupar-nos a dor de compartilhar com ela o espetáculo de um declínio tão trágico.” Em 25 de maio de 1943, aos quarenta anos, ela se atirou da Ponte George Washington em Nova York.

Mas as ideias têm um jeito de sobreviver às tragédias. Quando os sociólogos Hubert Blalock e Otis Duncan ressuscitaram a análise de caminhos nos anos 1960, o artigo de Burks serviu como fonte de inspiração deles. Duncan explicou que um de seus mentores, William Fielding Ogburn, havia mencionado brevemente os coeficientes de caminho em sua aula de 1946 sobre correlações parciais. “Ogburn tinha um relatório de um breve artigo de Wright, aquele que tratava do material de Burks, e eu adquiri essa separata”, disse Duncan. Eis aí! O artigo de 1926 de Burks despertou o interesse de Wright no uso inadequado das correlações parciais. A resposta de Wright encontrou o caminho até a aula de Ogburn vinte anos depois e se implantou na mente de Duncan. Vinte anos depois disso, quando Duncan leu o trabalho de Blalock sobre diagramas de caminhos, isso reavivou essa memória semiesquecida de seus anos de estudante. É verdadeiramente assombroso ver como essa frágil borboleta de uma ideia esvoaçou quase despercebida por duas gerações antes de reemergir triunfantemente à luz.

\section{Em busca de uma linguagem (o paradoxo das admissões de Berkeley)}
Apesar do trabalho precoce de Burks, meio século depois os estatísticos ainda lutavam mesmo para expressar a ideia de, quanto mais para estimar, efeitos diretos e indiretos. Um caso em questão é um paradoxo bem conhecido, relacionado ao paradoxo de Simpson, mas com um toque.

Em 1973, Eugene Hammel, um decano associado da University of California, notou uma tendência preocupante nas taxas de admissão da universidade para homens e mulheres. Seus dados mostravam que 44 por cento dos homens que se candidataram à pós-graduação em Berkeley tinham sido aceitos, em comparação com apenas 35 por cento das mulheres. A discriminação de gênero estava ganhando ampla atenção pública, e Hammel não queria esperar que outra pessoa começasse a fazer perguntas. Decidiu investigar as razões da disparidade.

As decisões de admissão à pós-graduação, em Berkeley como em outras universidades, são tomadas por departamentos individuais, e não pela universidade como um todo. Assim, fazia sentido olhar os dados de admissão departamento por departamento para isolar o culpado. Mas quando o fez, Hammel descobriu um fato assombroso. Departamento após departamento, as decisões de admissão eram consistentemente mais favoráveis às mulheres do que aos homens. Como isso poderia ser?

Nesse ponto Hammel fez algo inteligente: chamou um estatístico. Peter Bickel, quando solicitado a olhar os dados, reconheceu de imediato uma forma do paradoxo de Simpson. Como vimos no Capítulo 6, o paradoxo de Simpson refere-se a uma tendência que parece ir em uma direção em cada camada de uma população (as mulheres são aceitas a uma taxa mais alta em cada departamento), mas na direção oposta para a população como um todo (os homens são aceitos a uma taxa mais alta na universidade como um todo). Vimos também no Capítulo 6 que a resolução correta do paradoxo depende muito da pergunta que se quer responder. Nesse caso a pergunta é clara: A universidade (ou alguém dentro da universidade) está discriminando as mulheres?

Quando contei pela primeira vez esse exemplo à minha mulher, a reação dela foi: “É impossível. Se cada departamento discrimina de um jeito, a escola não pode discriminar do outro jeito.” E ela tem razão! O paradoxo ofende a nossa compreensão da discriminação, que é um conceito causal, envolvendo resposta preferencial ao sexo declarado de um candidato. Se todos os atores preferem um sexo ao outro, o grupo como um todo deve mostrar essa mesma preferência. Se os dados parecem dizer o contrário, isso deve significar que não estamos processando os dados adequadamente, de acordo com a lógica da causalidade. Só com tal lógica, e com uma história causal clara, podemos determinar a inocência ou a culpa da universidade.

De fato, Bickel e Hammel encontraram uma história causal que os satisfez completamente. Escreveram um artigo, publicado na revista \textit{Science} em 1975, propondo uma explicação simples: as mulheres eram rejeitadas em maior número porque se candidataram a departamentos mais difíceis de entrar. Especificamente, uma proporção mais alta de mulheres do que de homens se candidatou a departamentos de humanidades e ciências sociais. Ali enfrentavam um golpe duplo: o número de estudantes se candidatando era maior, e o número de vagas para esses estudantes era menor. Por outro lado, as mulheres não se candidatavam com tanta frequência a departamentos como engenharia mecânica, que eram mais fáceis de entrar. Esses departamentos tinham mais dinheiro e mais vagas para estudantes de pós-graduação — em suma, uma taxa de aceitação mais alta.

Por que as mulheres se candidatavam a departamentos mais difíceis de entrar? Talvez fossem desencorajadas de se candidatar a campos técnicos porque estes tinham mais requisitos de matemática ou eram percebidos como mais “masculinos”. Talvez tivessem sido discriminadas em estágios anteriores de sua educação: a sociedade tendia a empurrar as mulheres para longe dos campos técnicos, como a história de Barbara Burks mostra com demasiada clareza. Mas essas circunstâncias não estavam sob o controle de Berkeley e, portanto, não constituiriam discriminação por parte da universidade. Bickel e Hammel concluíram: “O \textit{campus} como um todo não se engajou em discriminação contra candidatas mulheres.”

Ao menos de passagem, gostaria de notar a precisão da linguagem de Bickel nesse artigo. Ele distingue cuidadosamente entre dois termos que, no inglês comum, são frequentemente tomados como sinônimos: “viés” (\textit{bias}) e “discriminação” (\textit{discrimination}). Define viés como “um padrão de associação entre uma decisão particular e um sexo particular de candidato”. Note as palavras “padrão” e “associação”. Elas nos dizem que o viés é um fenômeno no degrau um da Escada da Causalidade. Por outro lado, define discriminação como “o exercício de decisão influenciado pelo sexo do candidato quando isso é imaterial às qualificações para a entrada”. Palavras como “exercício de decisão”, “influência” e “imaterial” cheiram a causalidade, mesmo que Bickel não conseguisse pronunciar essa palavra em 1975. A discriminação, ao contrário do viés, pertence ao degrau dois ou três da Escada da Causalidade.

Em sua análise, Bickel sentiu que os dados deveriam ser estratificados por departamento porque os departamentos eram as unidades decisórias. Foi essa a decisão certa? Para responder a essa pergunta, começamos desenhando um diagrama causal (Figura 9.4).

\begin{figure}[ht]
\centering
\includegraphics[width=0.92\textwidth]{../figuras/fig_9_4.png}
\caption{Figura 9.4. , estavam certos em particionar os dados por departamentos, e o seu resultado fornece uma estimativa válida do efeito direto de Gênero sobre Desfecho. Tiveram sucesso mesmo que a linguagem de efeitos diretos e indiretos não estivesse disponível a Bickel em 1973.}
\end{figure}
\textit{Figura 9.4. , estavam certos em particionar os dados por departamentos, e o seu resultado fornece uma estimativa válida do efeito direto de Gênero sobre Desfecho. Tiveram sucesso mesmo que a linguagem de efeitos diretos e indiretos não estivesse disponível a Bickel em 1973.}

É também muito iluminador olhar a definição de discriminação na jurisprudência dos Estados Unidos. Ela usa terminologia contrafactual, um sinal claro de que subimos ao nível três da Escada da Causalidade. Em \textit{Carson v. Bethlehem Steel Corp.} (1996), o Tribunal do Sétimo Circuito escreveu: “A questão central em qualquer caso de discriminação no emprego é se o empregador teria tomado a mesma ação se o empregado tivesse sido de uma raça diferente (idade, sexo, religião, origem nacional etc.) e tudo o mais tivesse sido o mesmo.” Essa definição expressa claramente a ideia de que devemos desabilitar ou “congelar” todos os caminhos causais que levam do gênero à admissão por meio de qualquer outra variável (por exemplo, qualificação, escolha de departamento etc.). Em outras palavras, a discriminação é igual ao efeito direto do gênero sobre o desfecho de admissão.

\textbf{FIGURA 9.4.} Diagrama causal para o paradoxo das admissões de Berkeley — versão simples.

Já vimos antes que condicionar a um mediador é incorreto se queremos estimar o efeito total de uma variável sobre outra. Mas em um caso de discriminação, segundo o tribunal, não é o efeito total, e sim o efeito direto que importa. Assim, Bickel e Hammel são vindos: sob as hipóteses mostradas na Figura 9.4, estavam certos em particionar os dados por departamentos, e o seu resultado fornece uma estimativa válida do efeito direto de Gênero sobre Desfecho. Tiveram sucesso mesmo que a linguagem de efeitos diretos e indiretos não estivesse disponível a Bickel em 1973.

Contudo, a parte mais interessante dessa história não é o artigo original que Bickel e Hammel escreveram, mas a discussão que se seguiu. Depois que o artigo foi publicado, William Kruskal da University of Chicago escreveu uma carta a Bickel argumentando que a explicação deles não inocentava de fato Berkeley. De fato, Kruskal perguntava se qualquer estudo puramente observacional (em oposição a um experimento randomizado — digamos, usando formulários de candidatura falsos) poderia jamais fazê-lo.

Para mim, a troca de cartas deles é fascinante. Não é muito frequente que possamos testemunhar duas grandes mentes lutando com um conceito (causalidade) para o qual careciam de um vocabulário adequado. Bickel mais tarde viria a receber uma bolsa “gênio” da MacArthur Foundation em 1984. Mas em 1975 estava no início de sua carreira, e deve ter sido tanto uma honra quanto um desafio para ele medir forças com Kruskal, um gigante da comunidade estatística americana.

Em sua carta a Bickel, Kruskal apontou que a relação entre Departamento e Desfecho poderia ter um confundidor não medido, como Estado de Residência. Elaborou um exemplo numérico para uma universidade hipotética com dois departamentos que discriminavam por sexo e produziam exatamente os mesmos dados do exemplo de Bickel. Fez isso assumindo que ambos os departamentos aceitam todos os homens do estado e todas as mulheres de fora do estado e rejeitam todos os homens de fora do estado e todas as mulheres do estado, e que esse é o seu único critério de decisão. Claramente essa política de admissão é um exemplo flagrante e de manual de discriminação. Mas, como o número total de candidatos de cada gênero aceitos e rejeitados era exatamente o mesmo do exemplo de Bickel, Bickel teria de concluir que não havia discriminação. Segundo Kruskal, os departamentos parecem inocentes porque Bickel controlou por apenas uma variável em vez de duas.

Kruskal pôs o dedo exatamente no ponto fraco do artigo de Bickel: a falta de um critério claramente justificado para determinar quais variáveis controlar. Kruskal não ofereceu uma solução e, de fato, a sua carta desespera de jamais encontrar uma. Ao contrário de Kruskal, podemos desenhar um diagrama e ver exatamente qual é o problema. A Figura 9.5 mostra o diagrama causal representando o contraexemplo de Kruskal. Parece um pouco familiar? Deveria! É exatamente o mesmo diagrama que Barbara Burks desenhou em 1926, mas com variáveis diferentes. Fica-se tentado a dizer “Grandes mentes pensam igual”, mas talvez fosse mais apropriado dizer que grandes problemas atraem grandes mentes.

\begin{figure}[ht]
\centering
\includegraphics[width=0.92\textwidth]{../figuras/fig_9_5.png}
\caption{Figura 9.5. mostra o diagrama causal representando o contraexemplo de Kruskal. Parece um pouco familiar? Deveria! É exatamente o mesmo diagrama que Barbara Burks desenhou em 1926, mas com variáveis diferentes. Fica-se tentado a dizer “Grandes mentes pensam igual”, mas talvez fosse mais apropriado dizer que grandes problemas atraem grandes mentes.}
\end{figure}
\textit{Figura 9.5. mostra o diagrama causal representando o contraexemplo de Kruskal. Parece um pouco familiar? Deveria! É exatamente o mesmo diagrama que Barbara Burks desenhou em 1926, mas com variáveis diferentes. Fica-se tentado a dizer “Grandes mentes pensam igual”, mas talvez fosse mais apropriado dizer que grandes problemas atraem grandes mentes.}

\textbf{FIGURA 9.5.} Diagrama causal para o paradoxo das admissões de Berkeley — versão de Kruskal.

Kruskal argumentou que a análise nessa situação deveria controlar tanto pelo departamento quanto pelo estado de residência, e um olhar à Figura 9.5 explica por que isso é assim. Para desabilitar todos os caminhos exceto o direto, precisamos estratificar por departamento. Isso fecha o caminho indireto Gênero → Departamento → Desfecho. Mas, ao fazê-lo, abrimos o caminho espúrio Gênero → Departamento ← Estado de Residência → Desfecho, por causa do colisor em Departamento. Se controlarmos também por Estado de Residência, fechamos esse caminho e, portanto, qualquer correlação remanescente deve ser devida ao caminho direto (discriminatório) Gênero → Desfecho.

Na falta de diagramas, Kruskal teve de convencer Bickel com números, e de fato os seus números mostravam a mesma coisa. Se não ajustamos por nenhuma variável, então as mulheres têm uma taxa de admissão mais baixa. Se ajustamos por Departamento, então as mulheres parecem ter uma taxa de admissão mais alta. Se ajustamos tanto por Departamento quanto por Estado de Residência, então de novo os números mostram uma taxa de admissão mais baixa para as mulheres.

A partir de argumentos como esse, você pode ver por que o conceito de mediação despertou (e ainda desperta) tais suspeitas. Parece instável e difícil de fixar. Primeiro as taxas de admissão são enviesadas contra as mulheres, depois contra os homens, depois contra as mulheres. Em sua resposta a Kruskal, Bickel continuou a manter que condicionar a uma unidade decisória (Departamento) é de algum modo diferente de condicionar a um critério de decisão (Estado de Residência). Mas não soava nada confiante a esse respeito. Pergunta plangente: “Vejo aqui uma questão não estatística: O que queremos dizer com viés?” Por que o sinal do viés muda dependendo do modo como o medimos?

De fato ele tinha a ideia certa quando distinguiu entre viés e discriminação. O viés é uma noção estatística escorregadia, que pode desaparecer se você fatiar os dados de um jeito diferente. A discriminação, como conceito causal, reflete a realidade e deve permanecer estável. A frase que faltava no vocabulário de ambos era “manter constante”. Para desabilitar o caminho indireto de Gênero para Desfecho, devemos manter constante a variável Departamento e então “mexer” na variável Gênero. Quando mantemos o departamento constante, impedimos (figurativamente falando) que os candidatos escolham a qual departamento se candidatar.

Como os estatísticos não têm uma palavra para esse conceito, fazem algo superficialmente similar: condicionam a Departamento. Foi exatamente o que Bickel fez: olhou os dados departamento por departamento e concluiu que não havia evidência de discriminação contra mulheres. Esse procedimento é válido quando Departamento e Desfecho não estão confundidos; nesse caso, \textit{ver} é o mesmo que \textit{fazer}. Mas Kruskal corretamente perguntou: “E se houver um confundidor, Estado de Residência?” Provavelmente não se deu conta de que estava seguindo os passos de Burks, que havia desenhado essencialmente o mesmo diagrama.

Não posso enfatizar o bastante quão frequentemente esse equívoco foi repetido ao longo dos anos — condicionar ao mediador em vez de manter o mediador constante. Por essa razão, chamo-o de Falácia da Mediação. Admitidamente, o equívoco é inofensivo se não houver confundimento do mediador e do desfecho. Contudo, se houver confundimento, pode reverter completamente a análise, como o exemplo numérico de Kruskal mostrou. Pode levar o investigador a concluir que não há discriminação quando de fato há.

Burks e Kruskal foram incomuns em reconhecer a Falácia da Mediação como um equívoco, embora não tenham exatamente oferecido uma solução. R. A. Fisher caiu no mesmo equívoco em 1936, e oitenta anos depois os estatísticos ainda lutam com o problema. Felizmente, houve enorme progresso desde o tempo de Fisher. Os epidemiologistas, por exemplo, sabem agora que se deve ter cuidado com confundidores entre mediador e desfecho. Ainda assim, aqueles que evitam a linguagem dos diagramas (alguns economistas ainda o fazem) se queixam e confessam que é uma tortura explicar o que esse alerta significa.

Felizmente, o problema que Kruskal certa vez chamou de “talvez insolúvel” foi resolvido duas décadas atrás. Tenho este estranho sentimento de que Kruskal teria apreciado a solução e, em minha fantasia, imagino mostrando-lhe o poder do cálculo do \textit{do} e da algoritmização dos contrafactuais. Infelizmente, ele se aposentou em 1990, exatamente quando as regras do cálculo do \textit{do} estavam sendo moldadas, e morreu em 2005.

Tenho certeza de que alguns leitores estão se perguntando: O que finalmente aconteceu no caso de Berkeley? A resposta é: nada. Hammel e Bickel estavam convencidos de que Berkeley não tinha com o que se preocupar, e de fato nenhum processo judicial ou investigação federal se materializou. Os dados insinuavam discriminação reversa contra os homens e, de fato, havia evidência explícita disso: “Na maioria dos casos envolvendo status favorecido para mulheres, parece que as comissões de admissão estavam buscando superar escassez longamente estabelecida de mulheres em seus campos”, escreveu Bickel. Apenas três anos depois, um processo sobre ação afirmativa em outro \textit{campus} da University of California chegou até a Suprema Corte. Se a Suprema Corte tivesse derrubado a ação afirmativa, tal “status favorecido para mulheres” poderia ter se tornado ilegal. Contudo, a Suprema Corte manteve a ação afirmativa, e o caso de Berkeley se tornou uma nota de rodapé histórica.

Um homem sábio deixa a palavra final não com a Suprema Corte, mas com a sua mulher. Por que a minha tinha uma convicção intuitiva tão forte de que é absolutamente impossível uma escola discriminar enquanto cada um de seus departamentos age de forma justa? É um teorema do cálculo causal semelhante ao princípio da certeza (\textit{sure-thing principle}). O princípio da certeza, como Jimmie Savage o enunciou originalmente, diz respeito a efeitos totais, ao passo que esse teorema vale para efeitos diretos. A própria definição de um efeito direto em nível global se apoia na agregação de efeitos diretos nas subpopulações. Para dizê-lo de forma sucinta, a justiça local em toda parte implica justiça global. Minha mulher tinha razão.

\section{Daisy, os gatinhos e os efeitos indiretos}
Até agora discutimos os conceitos de efeitos diretos e indiretos de modo vago e intuitivo, mas não lhes dei um significado científico preciso. Já passa da hora de retificar essa omissão.

Comecemos pelo efeito direto, porque é indubitavelmente mais fácil, e podemos definir uma versão dele usando o cálculo do \textit{do} (isto é, no degrau dois da Escada da Causalidade). Consideraremos primeiro o caso mais simples, que inclui três variáveis: um tratamento \textit{X}, um desfecho \textit{Y} e um mediador \textit{M}. Obtemos o efeito direto de \textit{X} sobre \textit{Y} quando “mexemos” em \textit{X} sem permitir que \textit{M} mude. No contexto do exemplo do paradoxo das admissões de Berkeley, forçamos todo mundo a se candidatar ao departamento de história — isto é, \textit{do}(\textit{M} = 0). Designamos aleatoriamente algumas pessoas a declarar o seu sexo (na candidatura) como masculino (\textit{do}(\textit{X} = 1)) e algumas a declará-lo como feminino (\textit{do}(\textit{X} = 0)), independentemente de seus gêneros reais. Então observamos a diferença nas taxas de admissão entre os dois grupos de declaração. O resultado chama-se efeito direto controlado, ou CDE(0). Em símbolos,

\$\$\textbackslash\{\}mathrm\{CDE\}(0) = P(Y = 1 \textbackslash\{\}mid do(X = 1), do(M = 0)) - P(Y = 1 \textbackslash\{\}mid do(X = 0), do(M = 0)) \textbackslash\{\}tag\{9.1\}\$\$

O “0” em CDE(0) indica que forçamos o mediador a assumir o valor zero. Também poderíamos fazer o mesmo experimento, forçando todo mundo a se candidatar à engenharia: \textit{do}(\textit{M} = 1). Denotaríamos o efeito direto controlado resultante como CDE(1).

Já vemos uma diferença entre efeitos diretos e efeitos totais: temos duas versões diferentes do efeito direto controlado, CDE(0) e CDE(1). Qual é a certa? Uma opção é simplesmente reportar ambas as versões. De fato, não é impensável que um departamento discrimine as mulheres e o outro os homens, e seria interessante descobrir quem faz o quê. Essa era, afinal, a intenção original de Hammel.

Contudo, eu não recomendaria realizar esse experimento, e eis por quê. Imagine um candidato chamado Joe cujo sonho de vida é estudar engenharia e que por acaso foi (aleatoriamente) designado a se candidatar ao departamento de história. Tendo participado de algumas comissões de admissão, posso jurar categoricamente que a candidatura de Joe pareceria terrivelmente estranha à comissão. Seu A+ em ondas eletromagnéticas e B– em nacionalismo europeu distorceriam totalmente a decisão da comissão, independentemente de ele ter marcado “masculino” ou “feminino” em sua candidatura. A proporção de homens e mulheres admitidos sob essas distorções dificilmente refletiria a política de admissão em comparação com candidatos que normalmente se candidatam ao departamento de história.

Felizmente, uma alternativa evita as armadilhas desse experimento supercontrolado. Instruímos os candidatos a declarar um gênero randomizado, mas a se candidatar ao departamento que teriam preferido. Chamamos isso de efeito direto natural (NDE, do inglês \textit{natural direct effect}), porque cada candidato acaba em um departamento de sua escolha. A formulação “teria” é uma pista de que a definição formal do NDE exige contrafactuais. Para leitores que gostam de matemática, eis a definição expressa como fórmula:

\$\$\textbackslash\{\}mathrm\{NDE\} = P(Y\_\{M=M\_0\}=1 \textbackslash\{\}mid do(X = 1)) - P(Y\_\{M=M\_0\}=1 \textbackslash\{\}mid do(X = 0)) \textbackslash\{\}tag\{9.2\}\$\$

O termo interessante é o primeiro, que representa a probabilidade de que uma estudante selecionando um departamento de sua escolha (\textit{M} = \textit{M}<sub>0</sub>) seria admitida se falsificasse o sexo para constar “masculino” (\textit{do}(\textit{X} = 1)). Aqui a escolha do departamento é governada pelo sexo real, enquanto a admissão é decidida pelo sexo declarado (falso). Como o primeiro não pode ser mandatado, não podemos traduzir esse termo em um que envolva operadores \textit{do}; precisamos invocar o subscrito contrafactual.

Agora você sabe como definimos o efeito direto controlado e o efeito direto natural, mas como os calculamos? A tarefa é simples para o efeito direto controlado; como pode ser expresso como uma \textit{do}-expressão, precisamos apenas usar as leis do cálculo do \textit{do} para reduzir as \textit{do}-expressões a \textit{see}-expressões (isto é, probabilidades condicionais, que podem ser estimadas a partir de dados observacionais). O efeito direto natural, porém, coloca um desafio maior, porque não pode ser definido em uma \textit{do}-expressão. Exige a linguagem dos contrafactuais e, portanto, não pode ser estimado usando o cálculo do \textit{do}. Foi uma das maiores emoções da minha vida quando consegui despir a fórmula do NDE de todos os seus subscritos contrafactuais. O resultado, chamado Fórmula de Mediação, torna o NDE uma ferramenta verdadeiramente prática porque podemos estimá-lo a partir de dados observacionais.

Os efeitos indiretos, ao contrário dos efeitos diretos, não têm versão “controlada” porque não há modo de desabilitar o caminho direto mantendo alguma variável constante. Mas têm uma versão “natural”, o efeito indireto natural (NIE, do inglês \textit{natural indirect effect}), que é definido (como o NDE) usando contrafactuais. Para motivar a definição, considerarei um exemplo um tanto brincalhão sugerido pelo meu coautor.

Meu coautor e a mulher dele adotaram uma cadela chamada Daisy, uma misturada de poodle e chihuahua buliçosa, com vontade própria. Daisy não era tão fácil de adestrar para as necessidades quanto o cão anterior deles e, depois de várias semanas, ainda tinha “acidentes” ocasionais dentro de casa. Mas então algo muito estranho aconteceu. Dana e a mulher trouxeram para casa três gatinhos de abrigo do canil, e os “acidentes” pararam. Os gatinhos de abrigo ficaram com eles por três semanas, e Daisy não quebrou o adestramento uma única vez durante esse período. Foi só coincidência, ou os gatinhos de algum modo inspiraram Daisy a um comportamento civilizado?

A mulher de Dana sugeriu que os gatinhos poderiam ter dado a Daisy um senso de pertencimento a uma “matilha”, e ela não quereria sujar a área onde a matilha vivia. Essa teoria foi reforçada quando, poucos dias depois de os gatinhos voltarem ao abrigo, Daisy recomeçou a urinar na casa, como se nunca tivesse ouvido falar de boas maneiras.

Mas então ocorreu a Dana que outra coisa havia mudado quando os gatinhos chegaram e partiram. Enquanto os gatinhos estavam lá, Daisy tinha de ser ou separada deles ou cuidadosamente supervisionada. Assim, passava longos períodos na casinha de transporte ou sendo observada de perto por um humano, até amarrada a um humano. Ambas as intervenções — caixote e guia — também são métodos reconhecidos de adestramento doméstico. Quando os gatinhos partiram, os Mackenzie interromperam a supervisão intensiva, e o comportamento rude retornou. Dana hipotetizou que o efeito dos gatinhos não era direto (como na teoria da matilha), mas indireto, mediado pelo caixote e pela supervisão. A Figura 9.6 mostra um grafo causal.

\begin{figure}[ht]
\centering
\includegraphics[width=0.92\textwidth]{../figuras/fig_9_6.png}
\caption{Figura 9.6. mostra um grafo causal.}
\end{figure}
\textit{Figura 9.6. mostra um grafo causal.}

Nesse ponto, Dana e a mulher tentaram um experimento. Trataram Daisy como teriam feito com gatinhos por perto, mantendo-a em um caixote e supervisionando-a cuidadosamente fora do caixote. Se os acidentes parassem, poderiam razoavelmente concluir que o mediador era responsável. Se não parassem, então o efeito direto (a psicologia da matilha) se tornaria mais plausível.

\textbf{FIGURA 9.6.} Diagrama causal para o adestramento doméstico de Daisy.

Na hierarquia da evidência científica, o experimento deles seria considerado muito frágil — certamente não um que pudesse jamais ser publicado em uma revista científica. Um experimento real teria de ser conduzido com mais de um cão e tanto na presença quanto na ausência dos gatinhos. Não obstante, é a lógica causal por trás do experimento que nos interessa aqui. Pretendemos recriar o que teria acontecido se os gatinhos \textit{não} estivessem presentes e se o mediador tivesse sido fixado no valor que assumiria \textit{com} os gatinhos presentes. Em outras palavras, removemos os gatinhos (intervenção número um) e supervisionamos o cão como faríamos se os gatinhos estivessem presentes (intervenção número dois).

Quando se olha cuidadosamente o parágrafo acima, pode-se notar dois “teria”, que são condições contrafactuais. Os gatinhos estavam presentes quando a cadela mudou o comportamento — mas perguntamos o que teria acontecido se \textit{não} estivessem presentes. Do mesmo modo, se os gatinhos não estivessem presentes, Dana \textit{não} teria supervisionado Daisy — mas perguntamos o que teria acontecido se ele \textit{tivesse}. Você pode ver por que os estatísticos lutaram tanto tempo para definir efeitos indiretos. Se mesmo um único contrafactual era extravagante, então contrafactuais duplamente aninhados estavam completamente além do aceitável.

Não obstante, essa definição se conforma de perto com a nossa intuição natural sobre causalidade. Nossa intuição é tão convincente que a mulher de Dana, sem nenhum treinamento especial, compreendeu prontamente a lógica do experimento proposto. Para leitores à vontade com fórmulas, eis como definir o NIE que acabamos de descrever em palavras:

\$\$\textbackslash\{\}mathrm\{NIE\} = P(Y\_\{M=M\_1\}=1 \textbackslash\{\}mid do(X = 0)) - P(Y\_\{M=M\_0\}=1 \textbackslash\{\}mid do(X = 0)) \textbackslash\{\}tag\{9.3\}\$\$

O primeiro termo \textit{P} é o desfecho do experimento de Daisy: a probabilidade de adestramento doméstico bem-sucedido (\textit{Y} = 1), dado que não introduzimos outros animais de estimação (\textit{X} = 0), mas fixamos o mediador no valor que teria se os tivéssemos introduzido (\textit{M} = \textit{M}<sub>1</sub>). Contrastamos isso com a probabilidade de adestramento doméstico bem-sucedido sob condições “normais”, sem outros animais. Note que o contrafactual \textit{M}<sub>1</sub> tem de ser calculado para cada animal caso a caso: cães diferentes podem ter necessidades diferentes de Caixote/Supervisão. Isso coloca o efeito indireto fora do alcance do cálculo do \textit{do}. Também pode tornar o experimento inviável, porque o experimentador pode não conhecer \textit{M}<sub>1</sub>(\textit{u}) para um cão particular \textit{u}.

Não obstante, assumindo que não há confundimento entre \textit{M} e \textit{Y}, o efeito indireto natural ainda pode ser calculado. É possível remover todos os contrafactuais do NIE e chegar a uma Fórmula de Mediação para ele, como a do NDE. Essa quantidade, que exige informação do terceiro degrau da Escada da Causalidade, pode ainda assim ser reduzida a uma expressão que pode ser calculada com dados do degrau um. Tal redução só é possível porque fizemos uma hipótese de ausência de confundimento que, devido à natureza determinística das equações em um modelo causal estrutural, está no degrau três.

Para concluir a história de Daisy, o experimento foi inconclusivo. É questionável se Dana e a mulher monitoraram Daisy com o mesmo cuidado com que teriam feito se estivessem mantendo-a longe dos gatinhos. (Assim, não está claro se \textit{M} foi verdadeiramente fixado em \textit{M}<sub>1</sub>.) Com paciência e tempo — levou vários meses — Daisy por fim aprendeu a “fazer as suas necessidades” do lado de fora. Ainda assim, a história de Daisy encerra algumas lições úteis. Simplesmente por estar atento à possibilidade de um mediador, Dana pôde conjecturar outro mecanismo causal. Esse mecanismo teve uma consequência prática importante: ele e a mulher não precisaram manter a casa cheia de uma “matilha” de gatinhos de abrigo pelo resto da vida de Daisy.

\section{Mediação no País das Maravilhas Linear}
Quando se ouve falar de contrafactuais pela primeira vez, pode-se perguntar se uma maquinaria tão elaborada é realmente necessária para expressar um efeito indireto. Certamente, pode-se argumentar, um efeito indireto é simplesmente o que sobra depois que se tira o efeito direto. Alternativamente, poderíamos escrever:

\$\$\textbackslash\{\}text\{Efeito Total\} = \textbackslash\{\}text\{Efeito Direto\} + \textbackslash\{\}text\{Efeito Indireto\} \textbackslash\{\}tag\{9.4\}\$\$

A resposta curta é que isso não funciona em modelos que envolvem interações (às vezes chamadas de \textit{moderação}). Por exemplo, imagine um medicamento que faz o corpo secretar uma enzima que age como catalisador: combina-se com o medicamento para curar uma doença. O efeito total do medicamento é, naturalmente, positivo. Mas o efeito direto é zero, porque se desabilitarmos o mediador (por exemplo, impedindo o corpo de estimular a enzima), o medicamento não funcionará. O efeito indireto também é zero, porque se não recebermos o medicamento e obtivermos artificialmente a enzima, a doença não será curada. A enzima em si não tem poder curativo. Assim, a Equação 9.4 não se sustenta: o efeito total é positivo, mas os efeitos direto e indireto são zero.

Contudo, a Equação 9.4 se sustenta automaticamente em uma situação, sem aparente necessidade de invocar contrafactuais. Esse é o caso de um modelo causal linear, do tipo que vimos no Capítulo 8. Como discutido lá, os modelos lineares não permitem interações, o que pode ser tanto uma virtude quanto uma desvantagem. É uma virtude no sentido de que torna a análise de mediação muito mais fácil, mas é uma desvantagem se queremos descrever um processo causal do mundo real que envolve interações.

Como a análise de mediação é tão mais fácil para modelos lineares, vejamos como é feita e quais são as armadilhas. Suponha que tenhamos um diagrama causal que se parece com a Figura 9.7. Como estamos trabalhando com um modelo linear, podemos representar a força de cada efeito com um único número. Os rótulos (coeficientes de caminho) indicam que aumentar a variável Tratamento em uma unidade aumentará a variável Mediador em duas unidades. Do mesmo modo, um aumento de uma unidade no Mediador aumentará o Desfecho em três unidades, e um aumento de uma unidade no Tratamento aumentará o Desfecho em sete unidades. Esses são todos efeitos diretos. Chegamos aqui à primeira razão pela qual os modelos lineares são tão simples: os efeitos diretos não dependem do nível do mediador. Isto é, o efeito direto controlado CDE(\textit{m}) é o mesmo para todos os valores \textit{m}, e podemos simplesmente falar do efeito direto.

\begin{figure}[ht]
\centering
\includegraphics[width=0.92\textwidth]{../figuras/fig_9_7.png}
\caption{Figura 9.7. Como estamos trabalhando com um modelo linear, podemos representar a força de cada efeito com um único número. Os rótulos (coeficientes de caminho) indicam que aumentar a variável Tratamento em uma unidade aumentará a variável Mediador em duas unidades. Do mesmo modo, um aumento de uma unidade no Mediador aumentará o Desfecho em três unidades, e um aumento de uma unidade no Tratamento aumentará o Desfecho em sete unidades. Esses são todos efeitos diretos. Chegamos aqui à primeira razão pela qual os modelos lineares são tão simples: os efeitos diretos não dependem do nível do mediador. Isto é, o efeito direto controlado CDE(\textit{m}) é o mesmo para todos os valores \textit{m}, e podemos simplesmente falar do efeito direto.}
\end{figure}
\textit{Figura 9.7. Como estamos trabalhando com um modelo linear, podemos representar a força de cada efeito com um único número. Os rótulos (coeficientes de caminho) indicam que aumentar a variável Tratamento em uma unidade aumentará a variável Mediador em duas unidades. Do mesmo modo, um aumento de uma unidade no Mediador aumentará o Desfecho em três unidades, e um aumento de uma unidade no Tratamento aumentará o Desfecho em sete unidades. Esses são todos efeitos diretos. Chegamos aqui à primeira razão pela qual os modelos lineares são tão simples: os efeitos diretos não dependem do nível do mediador. Isto é, o efeito direto controlado CDE(}m\textit{) é o mesmo para todos os valores }m\textit{, e podemos simplesmente falar do efeito direto.}

Qual seria o efeito total de uma intervenção que faz o Tratamento aumentar em uma unidade? Primeiro, essa intervenção causa diretamente um aumento do Desfecho em sete unidades (se mantivermos o Mediador constante). Também faz o Mediador aumentar em duas unidades. Por fim, como cada aumento de uma unidade no Mediador causa diretamente um aumento de três unidades no Desfecho, um aumento de duas unidades no Mediador levará a um aumento adicional de seis unidades no Desfecho. Assim, o aumento líquido no Desfecho, por ambos os caminhos causais, será de treze unidades. As primeiras sete unidades correspondem ao efeito direto, e as seis unidades restantes correspondem ao efeito indireto. Fácil como torta!

\textbf{FIGURA 9.7.} Exemplo de um modelo linear (diagrama de caminhos) com mediação.

Em geral, se há mais de um caminho indireto de \textit{X} para \textit{Y}, avaliamos o efeito indireto ao longo de cada caminho tomando o produto de todos os coeficientes de caminho ao longo desse caminho. Depois obtemos o efeito indireto total somando todos os caminhos causais indiretos. Por fim, o efeito total de \textit{X} sobre \textit{Y} é a soma dos efeitos direto e indireto. Essa regra da “soma de produtos” tem sido usada desde que Sewall Wright inventou a análise de caminhos e, formalmente falando, de fato se segue da definição do efeito total pelo operador \textit{do}.

Em 1986, Reuben Baron e David Kenny articularam um conjunto de princípios para detectar e avaliar a mediação em um sistema de equações. Os princípios essenciais são, primeiro, que as variáveis estão todas relacionadas por equações lineares, que são estimadas ajustando-as aos dados. Segundo, os efeitos direto e indireto são calculados ajustando duas equações aos dados: uma com o mediador incluído e outra com o mediador excluído. Uma mudança significativa nos coeficientes quando o mediador é introduzido é tomada como evidência de mediação.

A simplicidade e a plausibilidade do método de Baron-Kenny tomaram as ciências sociais de assalto. Em 2014, o artigo deles ocupava o trigésimo terceiro lugar na lista dos artigos científicos mais citados de todos os tempos. Em 2017, o Google Scholar reportava que 73 mil artigos acadêmicos haviam citado Baron e Kenny. Pense nisso! Foram citados mais vezes do que Albert Einstein, mais do que Sigmund Freud, mais do que quase qualquer outro cientista famoso de que se possa lembrar. O artigo deles ocupa o segundo lugar entre todos os artigos de psicologia e psiquiatria, e no entanto não é sobre psicologia de modo algum. É sobre mediação não causal.

A popularidade sem precedentes da abordagem de Baron-Kenny sem dúvida decorre de dois fatores. Primeiro, a mediação está em alta demanda. Nosso desejo de compreender “como a natureza funciona” (isto é, encontrar o \textit{M} em \textit{X} → \textit{M} → \textit{Y}) é talvez ainda mais forte do que o nosso desejo de quantificá-lo. Segundo, o método se reduz facilmente a um procedimento de livro de receitas baseado em conceitos familiares da estatística, uma disciplina que há muito reivindica a propriedade exclusiva da objetividade e da validade empírica. Assim, quase ninguém notou o grande salto envolvido, o fato de que uma quantidade causal (mediação) estava sendo definida e avaliada por meios puramente estatísticos.

Contudo, rachaduras nessa edificação baseada em regressão começaram a aparecer no início dos anos 2000, quando os praticantes tentaram generalizar a regra da soma de produtos a sistemas não lineares. Essa regra envolve duas hipóteses — os efeitos ao longo de caminhos distintos são aditivos, e os coeficientes de caminho ao longo de um caminho se multiplicam — e ambas levam a respostas erradas em modelos não lineares, como veremos abaixo.

Levou um longo tempo, mas os praticantes da análise de mediação finalmente acordaram. Em 2001, meu falecido amigo e colega Rod McDonald escreveu: “Penso que o melhor modo de discutir a questão de detectar ou mostrar moderação ou mediação em uma regressão é pôr de lado a literatura inteira sobre esses tópicos e começar do zero.” A literatura mais recente sobre mediação parece atender ao conselho de McDonald; métodos contrafactuais e gráficos são perseguidos com muito mais atividade do que a abordagem de regressão. E em 2014, o pai da abordagem Baron-Kenny, David Kenny, postou uma nova seção em seu \textit{site} chamada “análise de mediação causal”. Embora eu ainda não o chamasse de convertido, Kenny reconhece claramente que os tempos estão mudando e que a análise de mediação está entrando em uma nova era.

Por ora, olhemos um exemplo muito simples de como as nossas expectativas dão errado quando deixamos o País das Maravilhas Linear. Considere a Figura 9.8, uma ligeira modificação da Figura 9.7, em que um candidato a emprego decidirá aceitar o emprego se e somente se o salário oferecido exceder um certo valor limiar, no nosso caso dez. A oferta salarial é determinada, como mostrado no diagrama, por 7 × Educação + 3 × Habilidade. Note que as funções que determinam Habilidade e Salário ainda se assumem lineares, mas a relação de Salário com Desfecho é não linear, porque tem um efeito limiar.

\begin{figure}[ht]
\centering
\includegraphics[width=0.92\textwidth]{../figuras/fig_9_8.png}
\caption{Figura 9.8. , uma ligeira modificação da Figura 9.7, em que um candidato a emprego decidirá aceitar o emprego se e somente se o salário oferecido exceder um certo valor limiar, no nosso caso dez. A oferta salarial é determinada, como mostrado no diagrama, por 7 × Educação + 3 × Habilidade. Note que as funções que determinam Habilidade e Salário ainda se assumem lineares, mas a relação de Salário com Desfecho é não linear, porque tem um efeito limiar.}
\end{figure}
\textit{Figura 9.8. , uma ligeira modificação da Figura 9.7, em que um candidato a emprego decidirá aceitar o emprego se e somente se o salário oferecido exceder um certo valor limiar, no nosso caso dez. A oferta salarial é determinada, como mostrado no diagrama, por 7 × Educação + 3 × Habilidade. Note que as funções que determinam Habilidade e Salário ainda se assumem lineares, mas a relação de Salário com Desfecho é não linear, porque tem um efeito limiar.}

Calculemos, para esse modelo, os efeitos total, direto e indireto associados a aumentar a Educação em uma unidade. O efeito total é claramente igual a um, porque quando a Educação passa de zero a um, o Salário vai de zero a (7 × 1) + (3 × 2) = 13, que está acima do limiar de dez, fazendo o Desfecho mudar de zero para um.

Lembre-se de que o efeito indireto natural é a mudança esperada no desfecho, dado que não fazemos nenhuma mudança na Educação, mas fixamos a Habilidade no nível que assumiria se tivéssemos aumentado a Educação em um. É fácil ver que, nesse caso, o Salário vai de zero a 2 × 3 = 6. Isso está abaixo do limiar de dez, de modo que o candidato recusará a oferta. Assim, NIE = 0.

\textbf{FIGURA 9.8.} Mediação combinada com um efeito limiar.

E quanto ao efeito direto? Como mencionado antes, temos o problema de descobrir em que valor manter o mediador. Se mantivermos a Habilidade no nível que tinha antes de mudarmos a Educação, então o Salário aumentará de zero para sete, fazendo Desfecho = 0. Assim, CDE(0) = 0. Por outro lado, se mantivermos a Habilidade no nível que atinge depois da mudança na Educação (a saber, dois), o Salário aumentará de seis para treze. Isso muda o Desfecho de zero para um, porque treze está acima do limiar do candidato para aceitar a oferta de emprego. Assim, CDE(2) = 1.

Portanto, o efeito direto é zero ou um, dependendo do valor constante que escolhemos para o mediador. Ao contrário do País das Maravilhas Linear, a escolha de um valor para o mediador faz diferença, e temos um dilema. Se quisermos preservar o princípio aditivo, Efeito Total = Efeito Direto + Efeito Indireto, precisamos usar CDE(2) como nossa definição do efeito causal. Mas isso parece arbitrário e até um tanto antinatural. Se estamos contemplando uma mudança na Educação e queremos saber o seu efeito direto, muito provavelmente gostaríamos de manter a Habilidade no nível que já tem. Em outras palavras, faz mais sentido intuitivo usar CDE(0) como nosso efeito direto. Não só isso: isso concorda com o efeito direto natural nesse exemplo. Mas então perdemos a aditividade: Efeito Total ≠ Efeito Direto + Efeito Indireto.

Contudo — de modo bastante surpreendente — uma versão um tanto modificada da aditividade de fato se sustenta, não só nesse exemplo, mas em geral. Leitores que não se importam de fazer um pouco de cálculo podem se interessar em calcular o NIE de voltar de \textit{X} = 1 a \textit{X} = 0. Nesse caso a oferta salarial cai de treze para sete, e o Desfecho cai de um para zero (isto é, o candidato não aceita a oferta). Assim, calculado na direção inversa, NIE = –1. O fato legal e assombroso é que

\$\$\textbackslash\{\}text\{Efeito Total \}(X = 0 \textbackslash\{\}to X = 1) = \textbackslash\{\}mathrm\{NDE\}\textbackslash\{\} (X = 0 \textbackslash\{\}to X = 1) - \textbackslash\{\}mathrm\{NIE\}\textbackslash\{\} (X = 1 \textbackslash\{\}to X = 0)\$\$

ou, neste caso, 1 = 0 – (–1). Essa é a versão dos “efeitos naturais” do princípio de aditividade, só que é um princípio de \textit{subtratividade}! Fiquei extremamente feliz de ver essa versão da aditividade emergir da análise, apesar da não linearidade das equações.

Uma quantidade assombrosa de tinta foi gasta sobre o “modo certo” de generalizar efeitos diretos e indiretos de modelos lineares para não lineares. Infelizmente, a maioria dos artigos ataca o problema de trás para a frente. Em vez de repensar do zero o que queremos dizer por efeitos diretos e indiretos, partem da suposição de que só precisamos ajustar um pouco as definições lineares. Por exemplo, no País das Maravilhas Linear vimos que o efeito indireto é dado por um produto de dois coeficientes de caminho. Assim, alguns pesquisadores tentaram definir o efeito indireto na forma de um produto de duas quantidades, uma medindo o efeito de \textit{X} sobre \textit{M}, a outra o efeito de \textit{M} sobre \textit{Y}. Isso veio a ser conhecido como o método do “produto de coeficientes”. Mas também vimos que no País das Maravilhas Linear o efeito indireto é dado pela diferença entre o efeito total e o efeito direto. Assim, outro grupo de pesquisadores, igualmente dedicado, definiu o efeito indireto como uma diferença de duas quantidades, uma medindo o efeito total, a outra o efeito direto. Isso veio a ser conhecido como o método da “diferença de coeficientes”.

Qual desses está certo? Nenhum! Ambos os grupos de pesquisadores confundiram o procedimento com o significado. O procedimento é matemático; o significado é causal. De fato, o problema vai ainda mais fundo: o efeito indireto nunca teve um significado para analistas de regressão fora da bolha dos modelos lineares. O único significado do efeito indireto era como o resultado de um procedimento algébrico (“multiplique os coeficientes de caminho”). Uma vez que esse procedimento lhes foi retirado, ficaram à deriva, como um barco sem âncora.

Uma leitora do meu livro \textit{Causality} descreveu esse sentimento de perda belamente em uma carta a mim. Melanie Wall, agora na Columbia University, costumava dar um curso de modelagem a estudantes de bioestatística e saúde pública. Certa vez, explicou aos estudantes como de costume como calcular o efeito indireto tomando o produto dos coeficientes de caminho diretos. Um estudante perguntou-lhe o que o efeito indireto significava. “Dei a resposta que sempre dou, de que o efeito indireto é o efeito que uma mudança em \textit{X} tem sobre \textit{Y} por meio de sua relação com o mediador, \textit{Z}”, Wall me contou. Mas o estudante foi persistente. Lembrou-se de como a professora havia explicado o efeito direto como o efeito remanescente depois de manter o mediador fixo, e perguntou: “Então o que está sendo mantido constante quando interpretamos um efeito indireto?” Wall não sabia o que dizer. “Não tenho certeza se tenho uma boa resposta para você”, disse ela. “Que tal eu voltar a falar com você?”

Isso foi em outubro de 2001, apenas quatro meses depois de eu ter apresentado um artigo sobre mediação causal na conferência Uncertainty in Artificial Intelligence em Seattle. Desnecessário dizer, eu estava ansioso para impressionar Melanie com a minha solução recém-adquirida para o quebra-cabeça dela, e escrevi a ela a mesma resposta que lhe dei aqui: “O efeito indireto de \textit{X} sobre \textit{Y} é o aumento que veríamos em \textit{Y} mantendo \textit{X} constante e aumentando \textit{M} para qualquer valor que \textit{M} atingiria sob um aumento unitário em \textit{X}.” Não tenho certeza se Melanie ficou impressionada com a minha resposta, mas o seu estudante inquisitivo me fez pensar, de modo bastante sério, sobre como a ciência progride em nossos tempos. Aqui estamos, pensei, quarenta anos depois de Blalock e Duncan terem introduzido a análise de caminhos nas ciências sociais. Dezenas de livros-texto e centenas de artigos de pesquisa são publicados todo ano sobre efeitos diretos e indiretos, alguns com títulos oximorônicos como “Regression-Based Approach to Mediation” [“Abordagem baseada em regressão para a mediação”]. Cada geração passa à seguinte a sabedoria recebida de que o efeito indireto é apenas o produto de dois outros efeitos, ou a diferença entre os efeitos total e direto. Ninguém ousa fazer a pergunta simples: “Mas o que o efeito indireto significa em primeiro lugar?” Assim como o menino na fábula de Hans Christian Andersen “As roupas novas do imperador”, precisávamos de um estudante inocente com \textit{chutzpah} desavergonhado para abalar a nossa fé no papel oracular do consenso científico.

\section{Abrace os “teria”}
Neste ponto devo contar a minha própria história de conversão, porque por bastante tempo fiquei bloqueado pela mesma pergunta que intrigava o estudante de Melanie Wall. Escrevi no Capítulo 4 sobre Jamie Robins (Figura 9.9), um estatístico e epidemiologista pioneiro da Harvard University que, juntamente com Sander Greenland da University of California, Los Angeles, é em grande parte responsável pela ampla adoção de modelos gráficos na epidemiologia de hoje. Colaboramos por um par de anos, de 1993 a 1995, e ele me fez pensar no problema de planos de intervenção sequenciais, que era um dos seus principais interesses de pesquisa.

\begin{figure}[ht]
\centering
\includegraphics[width=0.92\textwidth]{../figuras/fig_9_9.png}
\caption{Figura 9.9. ), um estatístico e epidemiologista pioneiro da Harvard University que, juntamente com Sander Greenland da University of California, Los Angeles, é em grande parte responsável pela ampla adoção de modelos gráficos na epidemiologia de hoje. Colaboramos por um par de anos, de 1993 a 1995, e ele me fez pensar no problema de planos de intervenção sequenciais, que era um dos seus principais interesses de pesquisa.}
\end{figure}
\textit{Figura 9.9. ), um estatístico e epidemiologista pioneiro da Harvard University que, juntamente com Sander Greenland da University of California, Los Angeles, é em grande parte responsável pela ampla adoção de modelos gráficos na epidemiologia de hoje. Colaboramos por um par de anos, de 1993 a 1995, e ele me fez pensar no problema de planos de intervenção sequenciais, que era um dos seus principais interesses de pesquisa.}

\textbf{FIGURA 9.9.} Jamie Robins, um pioneiro da inferência causal em epidemiologia. (Fonte: Fotografia de Kris Snibbe, cortesia de Harvard University Photo Services.)

Anos antes, como especialista em saúde e segurança ocupacional, Robins fora chamado a depor em juízo sobre a verossimilhança de que a exposição a produtos químicos no local de trabalho tivesse causado a morte de um trabalhador. Ficou consternado ao descobrir que estatísticos e epidemiologistas não tinham ferramentas para responder a tais perguntas. Essa ainda era a era em que a linguagem causal era tabu na estatística. Só era permitida no caso de um ensaio clínico randomizado e, por razões éticas, nunca se poderia conduzir tal ensaio sobre os efeitos da exposição ao formaldeído.

Normalmente um trabalhador de fábrica é exposto a um produto químico nocivo não apenas uma vez, mas ao longo de um longo período. Por essa razão, Robins se interessou vivamente por exposições ou tratamentos que variam ao longo do tempo. Tais exposições também podem ser benéficas: por exemplo, o tratamento da AIDS é administrado ao longo de muitos anos, com diferentes planos de ação dependendo de como a contagem de CD4 do paciente responde. Como se pode separar o efeito causal do tratamento quando ele pode ocorrer em muitos estágios e as variáveis intermediárias (que se poderia querer usar como controles) dependem de estágios anteriores do tratamento? Essa tem sido uma das perguntas definidoras da carreira de Robins.

Depois que Jamie voou à Califórnia para me encontrar ao ouvir falar do “problema do guardanapo” (Capítulo 7), interessou-se vivamente em aplicar métodos gráficos aos planos de tratamento sequenciais que eram o seu \textit{métier}. Juntos chegamos a um critério sequencial da porta dos fundos para estimar o efeito causal de tal fluxo de tratamento. Aprendi algumas lições importantes com essa colaboração. Em particular, ele me mostrou que duas ações às vezes são mais fáceis de analisar do que uma, porque uma ação corresponde a apagar setas em um grafo, o que o torna mais esparso. O nosso critério da porta dos fundos lidava com um tratamento de longo prazo consistindo em um número arbitrariamente grande de \textit{do}-operações. Mas mesmo duas operações produzem alguma matemática interessante — incluindo o efeito direto controlado, que consiste em uma ação que “mexe” no valor do tratamento, enquanto outra ação fixa o valor do mediador. Mais importante, a ideia de definir efeitos diretos em termos de \textit{do}-operações libertou-os dos confins dos modelos lineares e os ancorou no cálculo causal.

Mas eu não me interessei de verdade pela mediação até mais tarde, quando vi que as pessoas ainda cometiam erros elementares, como a Falácia da Mediação mencionada anteriormente. Também me frustrava o fato de que a definição baseada em ação do efeito direto não se estendia ao efeito indireto. Como o estudante de Melanie Wall disse, não temos variável ou conjunto de variáveis sobre as quais intervir para desabilitar o caminho direto e deixar o caminho indireto ativo. Por essa razão o efeito indireto me parecia um figmento da imaginação, desprovido de significado independente, exceto para nos lembrar de que o efeito total pode diferir do efeito direto. Eu até disse isso na primeira edição (2000) do meu livro \textit{Causality}. Esse foi um dos três maiores equívocos da minha carreira.

Em retrospecto, eu estava cego pelo sucesso do cálculo do \textit{do}, que me levara a crer que o único modo de desabilitar um caminho causal era tomar uma variável e fixá-la em um valor particular. Isso não é assim; se tenho um modelo causal, posso manipulá-lo de muitas maneiras criativas, ditando quem ouve quem, quando e como. Em particular, posso fixar a variável primária com o propósito de suprimir o seu efeito direto e, hipoteticamente e simultaneamente, energizar a variável primária com o propósito de transmitir o seu efeito por meio do mediador. Isso me permite fixar a variável de tratamento (por exemplo, gatinhos) em zero e fixar o mediador no valor que teria se eu tivesse fixado gatinhos em um. Meu modelo do processo gerador de dados me diz então como calcular o efeito da intervenção dividida.

Sou grato a um leitor da primeira edição, Jacques Hagenaars (autor de \textit{Categorical Longitudinal Data}), por me instar a não desistir do efeito indireto. “Muitos especialistas em ciências sociais concordam sobre a entrada e a saída, mas diferem exatamente com respeito ao mecanismo”, escreveu-me ele. Mas fiquei bloqueado por quase dois anos no dilema de que escrevi na seção anterior: Como posso desabilitar o efeito direto?

Todas essas lutas chegaram a uma resolução súbita, quase como uma revelação divina, quando li a definição legal de discriminação que citei anteriormente neste capítulo: “se o empregado tivesse sido de uma raça diferente… e tudo o mais tivesse sido o mesmo.” Eis aí — o cerne da questão! É um jogo de faz de conta. Lidamos com cada indivíduo por seus próprios méritos, e mantemos todas as características do indivíduo constantes no nível que tinham antes da mudança na variável de tratamento.

Como isso resolve o nosso dilema? Significa, em primeiro lugar, que temos de redefinir tanto o efeito direto quanto o efeito indireto. Para o efeito direto, deixamos o mediador escolher o valor que teria — para cada indivíduo — na ausência de tratamento, e o fixamos ali. Agora mexemos no tratamento e registramos a diferença. Isso é diferente do efeito direto controlado que discuti anteriormente, em que o mediador é fixado em um valor para todos. Como deixamos o mediador escolher o seu valor “natural”, chamei-o de efeito direto natural. Do mesmo modo, para o efeito indireto natural, primeiro nego o tratamento a todos, e então deixo o mediador escolher o valor que teria, para cada indivíduo, na presença de tratamento. Por fim registro a diferença.

Não sei se as palavras legais na definição de discriminação teriam me comovido, ou a qualquer outra pessoa, da mesma maneira. Mas por volta de 2000 eu já falava contrafactuais como um nativo. Tendo aprendido a lê-los em modelos causais, percebi que não eram nada senão quantidades calculadas por operações inocentes sobre equações ou diagramas. Como tais, estavam prontas a ser encapsuladas em uma fórmula matemática. Tudo o que eu tinha de fazer era abraçar os “teria”.

Em um segundo, percebi que todo efeito direto e indireto podia ser traduzido em uma expressão contrafactual. Uma vez que vi como fazer isso, foi moleza derivar uma fórmula que diz como estimar os efeitos direto e indireto naturais a partir dos dados e quando isso é permissível. Importante: a fórmula não faz hipóteses sobre a forma funcional específica da relação entre \textit{X}, \textit{M} e \textit{Y}. Escapamos do País das Maravilhas Linear.

Chamei a nova regra de Fórmula de Mediação, embora na verdade haja duas fórmulas, uma para o efeito direto natural e outra para o efeito indireto natural. Sujeita a hipóteses transparentes, explicitamente exibidas no grafo, ela diz como podem ser estimados a partir dos dados. Por exemplo, em uma situação como a da Figura 9.4, em que não há confundimento entre nenhuma das variáveis, e \textit{M} é o mediador entre o tratamento \textit{X} e o desfecho \textit{Y}:

\$\$\textbackslash\{\}mathrm\{NIE\} = \textbackslash\{\}sum\_m \textbackslash\{\}bigl[P(M = m \textbackslash\{\}mid X = 1) - P(M = m \textbackslash\{\}mid X = 0)\textbackslash\{\}bigr] \textbackslash\{\}times P(Y = 1 \textbackslash\{\}mid X = 0, M = m) \textbackslash\{\}tag\{9.5\}\$\$

A interpretação dessa fórmula é iluminadora. A expressão entre colchetes representa o efeito de \textit{X} sobre \textit{M}, e a expressão seguinte representa o efeito de \textit{M} sobre \textit{Y} (quando \textit{X} = 0). Assim, revela a origem da ideia do produto de coeficientes, fundida como um produto de dois efeitos não lineares. Note também que, ao contrário da Equação 9.3, a Equação 9.5 não tem subscritos nem operadores \textit{do}, de modo que pode ser estimada a partir de dados do degrau um.

Seja você um cientista em um laboratório ou uma criança andando de bicicleta, é sempre uma emoção descobrir que pode fazer hoje algo que não podia fazer ontem. E foi assim que me senti quando a Fórmula de Mediação apareceu pela primeira vez no papel. Pude ver de relance tudo sobre efeitos diretos e indiretos: o que é necessário para torná-los grandes ou pequenos, quando podemos estimá-los a partir de dados observacionais ou intervencionistas, e quando podemos considerar um mediador “responsável” por transmitir mudanças observadas à variável de desfecho. A relação entre causa e efeito pode ser linear ou não linear, numérica ou lógica. Anteriormente cada um desses casos tinha de ser tratado de um modo diferente, se é que eram discutidos. Agora uma única fórmula se aplicaria a todos eles. Dados os dados certos e o modelo certo, poderíamos determinar se um empregador era culpado de discriminação ou que tipos de confundidores nos impediriam de fazer essa determinação. A partir dos dados de Barbara Burks, poderíamos estimar quanto do QI da criança vem da natureza e quanto da criação. Poderíamos até calcular a porcentagem do efeito total explicada pela mediação e a porcentagem devida à mediação — dois conceitos complementares que colapsam em um só em modelos lineares.

Depois de escrever a definição contrafactual dos efeitos direto e indireto, soube que não fora o primeiro a chegar à ideia. Robins e Greenland chegaram antes de mim, já em 1992. Mas o artigo deles descreve o conceito do efeito natural em palavras, sem comprometê-lo a uma fórmula matemática. Mais seriamente, adotaram uma visão pessimista da ideia inteira de efeitos naturais e afirmaram que tais efeitos não podem ser estimados a partir de estudos experimentais e certamente não a partir de estudos observacionais. Essa afirmação impediu outros pesquisadores de ver o potencial dos efeitos naturais. É difícil dizer se Robins e Greenland teriam mudado para uma visão mais otimista se tivessem dado o passo extra de expressar o efeito natural como uma fórmula em linguagem contrafactual. Para mim, esse passo extra foi crucial.

Há possivelmente outra razão para a visão pessimista deles, com a qual não concordo, mas tentarei explicar. Examinaram a definição contrafactual do efeito natural e viram que ela combina informação de dois mundos diferentes, um em que se mantém o tratamento constante em zero e outro em que se muda o mediador para o que teria sido se se tivesse fixado o tratamento em um. Como não se pode replicar essa condição “entre mundos” em nenhum experimento, acreditavam que estava fora dos limites.

Essa é uma diferença filosófica entre a escola deles e a minha. Eles creem que a legitimidade da inferência causal reside em replicar um experimento randomizado o mais de perto possível, na hipótese de que essa é a única rota para a verdade científica. Eu creio que pode haver outras rotas, que derivam a sua legitimidade de uma combinação de dados e conhecimento científico estabelecido (ou assumido). Para esse fim, pode haver métodos mais poderosos do que um experimento randomizado, baseados em hipóteses do degrau três, e não hesito em usá-los. Onde eles deram aos pesquisadores um sinal vermelho, eu dei um sinal verde, que era a Fórmula de Mediação: se você se sente confortável com essas hipóteses, eis o que pode fazer!

Infelizmente, o sinal vermelho de Robins e Greenland manteve o campo da mediação parado por nove anos inteiros. Muitas pessoas acham as fórmulas intimidantes, vendo-as como um modo de ocultar em vez de revelar informação. Mas para um matemático, ou para uma pessoa adequadamente treinada no modo matemático de pensar, exatamente o inverso é verdade. Uma fórmula revela tudo: não deixa nada à dúvida ou à ambiguidade. Ao ler um artigo científico, muitas vezes me pego saltando de fórmula em fórmula, pulando as palavras por completo. Para mim, uma fórmula é uma ideia assada. As palavras são ideias no forno.

Uma fórmula serve a dois propósitos, um prático e um social. Do ponto de vista prático, estudantes ou colegas podem lê-la como leria uma receita. A receita pode ser simples ou complexa, mas no fim do dia promete que, se você seguir os passos, conhecerá os efeitos direto e indireto naturais — desde que, naturalmente, o seu modelo causal reflita com precisão o mundo real.

O segundo propósito é mais sutil. Tive um amigo em Israel que era um artista famoso. Visitei o estúdio dele para adquirir uma de suas pinturas, e as suas telas estavam por toda parte — uma centena debaixo da cama, dezenas na cozinha. Estavam precificadas entre 300 e 500 dólares cada, e eu tinha dificuldade em decidir. Por fim, apontei para uma na parede e disse: “Gosto desta.” “Esta custa 5 mil dólares”, disse ele. “Como assim?”, perguntei, em parte surpreso e em parte protestando. Ele respondeu: “Esta está emoldurada.” Levou-me alguns minutos para entender, mas então compreendi o que ele queria dizer. Não era valiosa porque estava emoldurada; estava emoldurada porque era valiosa. De todas as centenas de pinturas no apartamento dele, aquela era a sua escolha pessoal. Expressava melhor o que ele se esforçara por expressar nas outras, e assim fora ungida com um selo de completude — uma moldura.

Esse é o segundo propósito de uma fórmula. É um contrato social. Põe uma moldura em torno de uma ideia e diz: “Isto é algo que creio ser importante. Isto é algo que merece ser compartilhado.” É por isso que escolhi pôr uma moldura em torno da Fórmula de Mediação. Merece ser compartilhada porque, para mim e muitos como eu, representa o fim de um dilema milenar. E é importante, porque oferece uma ferramenta prática para identificar mecanismos e avaliar a sua importância. Essa é a promessa social que a Fórmula de Mediação expressa.

Desde então, uma vez que se firmou a percepção de que a análise de mediação não linear é possível, a pesquisa no campo decolou. Se você for a uma base de dados de artigos acadêmicos e buscar títulos com as palavras “análise de mediação” (\textit{mediation analysis}), quase não encontrará nada antes de 2004. Depois eram sete artigos por ano, depois dez, depois vinte; agora são mais de cem artigos por ano.

Gostaria de encerrar este capítulo com três exemplos, que espero ilustrem a variedade de possibilidades da análise de mediação.

\section{Estudos de caso de mediação}
\subsection{“Álgebra para Todos”: um programa e os seus efeitos colaterais}
Como muitos sistemas de escolas públicas de grandes cidades, as Escolas Públicas de Chicago enfrentam problemas que às vezes parecem intratáveis: altas taxas de pobreza, orçamentos baixos e grandes lacunas de desempenho entre estudantes negros, latinos, brancos e asiáticos. Em 1988, o então secretário de educação dos EUA William Bennett chamou as escolas públicas de Chicago de as piores da nação. Mas nos anos 1990, sob nova liderança, as Escolas Públicas de Chicago empreenderam uma série de reformas e passaram de “piores da nação” a “inovadoras da nação”. Alguns dos superintendentes responsáveis por essas mudanças ganharam proeminência nacional, como Arne Duncan, que se tornou secretário de educação sob o presidente Barack Obama.

Uma inovação que de fato antecedeu Duncan foi uma política, adotada em 1997, eliminando cursos de recuperação no ensino médio e exigindo que todos os alunos do nono ano fizessem cursos preparatórios para a faculdade como Inglês I e Álgebra I. A parte matemática dessa política foi chamada “Álgebra para Todos” (\textit{Algebra for All}).

“Álgebra para Todos” foi um sucesso? Essa pergunta, revelou-se, era surpreendentemente difícil de responder. Havia tanto boas quanto más notícias. A boa notícia era que as notas nos testes de fato melhoraram. As notas de matemática subiram 7,8 pontos em três anos, uma mudança estatisticamente significativa equivalente a cerca de 75 por cento dos estudantes pontuando acima da média que existia antes da mudança de política. Mas antes de podermos falar de causalidade, temos de descartar confundidores, e nesse caso há um importante. Por volta de 1997, as qualificações dos alunos que entravam no nono ano já estavam melhorando graças a mudanças anteriores no currículo do K–8. Assim, não estamos comparando maçãs com maçãs. Como essas crianças começaram o nono ano com melhores habilidades matemáticas do que os estudantes de 1994, as notas mais altas poderiam ser devidas às mudanças do K–8 já instituídas, e não a “Álgebra para Todos”.

Guanglei Hong, professora de desenvolvimento humano na University of Chicago, estudou os dados e não encontrou nenhuma melhora significativa nas notas dos testes uma vez que esse confundidor fosse levado em conta. Nesse ponto teria sido fácil para Hong saltar à conclusão de que “Álgebra para Todos” não foi um sucesso. Mas ela não o fez, porque havia outro fator — desta vez um mediador, não um confundidor — a levar em conta.

Como qualquer bom professor sabe, o sucesso dos estudantes depende não só do que se ensina, mas de \textit{como} se ensina. Quando a política “Álgebra para Todos” foi introduzida, mudou mais do que o currículo. Os estudantes de menor desempenho se viram em salas de aula com estudantes de maior desempenho e não conseguiam acompanhar. Isso levou a todo tipo de consequências negativas: desânimo, falta às aulas e, naturalmente, notas mais baixas nos testes. Além disso, em uma sala de aula de habilidades mistas, os estudantes de baixo desempenho podem ter recebido menos atenção dos professores do que teriam em uma classe de recuperação. Por fim, os próprios professores podem ter lutado com as novas demandas colocadas sobre eles. Os professores experientes em ensinar Álgebra I provavelmente não eram experientes em ensinar estudantes de baixa habilidade, e os professores experientes com estudantes de baixa habilidade podem não ter sido tão qualificados para ensinar álgebra. Todos esses foram efeitos colaterais não antecipados de “Álgebra para Todos”.

A análise de mediação é idealmente adequada para avaliar a influência de efeitos colaterais. Hong hipotetizou, portanto, que o ambiente da sala de aula havia mudado e afetara fortemente o desfecho da intervenção. Em outras palavras, postulou o diagrama causal mostrado na Figura 9.10. O Ambiente (que Hong mediu pelo nível mediano de habilidade de todos os estudantes na sala de aula) funciona como um mediador entre a intervenção “Álgebra para Todos” e os desfechos de aprendizagem dos estudantes. A pergunta, como de costume na análise de mediação, é quanto do efeito da política foi direto e quanto foi indireto.

\begin{figure}[ht]
\centering
\includegraphics[width=0.92\textwidth]{../figuras/fig_9_10.png}
\caption{Figura 9.10. O Ambiente (que Hong mediu pelo nível mediano de habilidade de todos os estudantes na sala de aula) funciona como um mediador entre a intervenção “Álgebra para Todos” e os desfechos de aprendizagem dos estudantes. A pergunta, como de costume na análise de mediação, é quanto do efeito da política foi direto e quanto foi indireto.}
\end{figure}
\textit{Figura 9.10. O Ambiente (que Hong mediu pelo nível mediano de habilidade de todos os estudantes na sala de aula) funciona como um mediador entre a intervenção “Álgebra para Todos” e os desfechos de aprendizagem dos estudantes. A pergunta, como de costume na análise de mediação, é quanto do efeito da política foi direto e quanto foi indireto.}

Interessantemente, os dois efeitos trabalharam em direções opostas. Hong descobriu que o efeito direto era positivo: a nova política levou diretamente a um aumento de cerca de 2,7 pontos nas notas dos testes. Isso era pelo menos uma mudança na direção certa, e era estatisticamente significativo (significando que tal melhora seria improvável de ocorrer por acaso). Contudo, por causa das mudanças no ambiente da sala de aula, o efeito indireto quase cancelara completamente essa melhora, reduzindo as notas dos testes em 2,3 pontos.

\textbf{FIGURA 9.10.} Diagrama causal para o experimento “Álgebra para Todos”.

Hong concluiu que a implementação de “Álgebra para Todos” havia minado seriamente a política. Manter a mudança curricular, mas retornar ao ambiente de sala de aula pré-política, deveria resultar em um aumento modesto nas notas dos testes dos estudantes (e, espera-se, na aprendizagem dos estudantes).

Por serendipidade, foi exatamente o que aconteceu. Em 2003 as Escolas Públicas de Chicago (agora lideradas por Duncan) instituíram uma nova reforma chamada “Álgebra em Dose Dupla” (\textit{Double-Dose Algebra}). Essa reforma ainda exigiria que todos os estudantes fizessem álgebra, mas os estudantes que pontuassem abaixo da mediana nacional no oitavo ano fariam duas aulas de álgebra por dia em vez de uma. Isso reparou o efeito colateral adverso da reforma anterior. Agora, pelo menos uma vez por dia, os estudantes de menor desempenho obtinham um ambiente de sala de aula mais próximo daquele de que desfrutavam antes da reforma “Álgebra para Todos”. A reforma “Álgebra em Dose Dupla” foi geralmente considerada um sucesso e continua até hoje.

Considero a história de “Álgebra para Todos” um sucesso também para a análise de mediação, porque a análise explica tanto os resultados pouco impressionantes da política original quanto os resultados melhorados sob a política modificada. Mesmo que a inferência causal tenha chegado tarde demais para afetar a política em tempo real, responde às nossas perguntas “Por quê?” depois do fato: Por que a reforma original teve pouco efeito? Por que a segunda reforma funcionou melhor? Desse modo pode orientar a política para o futuro.

Quero apontar outra coisa interessante sobre o trabalho de Hong. Ela estava bem ciente da abordagem de Baron-Kenny a efeitos diretos e indiretos que chamei de País das Maravilhas Linear. Em seu artigo de fato realizou a mesma análise duas vezes: uma usando uma variação da Fórmula de Mediação, a outra usando os “procedimentos convencionais” (termo dela) de Baron e Kenny. O método de Baron-Kenny falhou em detectar o efeito indireto. A razão é muito provavelmente exatamente o que discuti antes: métodos lineares não conseguem captar interações entre o tratamento e o mediador. Talvez a combinação de material mais difícil e um ambiente de sala de aula menos solidário tenha feito os estudantes de baixo desempenho se desanimarem. Isso é plausível? Penso que sim. Álgebra é um assunto difícil. Talvez a sua dificuldade tenha tornado a atenção extra dos professores sob a política de dose dupla tanto mais valiosa.

\subsection{O gene do fumo: mediação e interação}
No Capítulo 5 escrevi sobre a guerra científica e política em torno do tabagismo nos anos 1950 e 1960. Os céticos daquela era, que incluíam R. A. Fisher e Jacob Yerushalmy, argumentavam que o aparente vínculo entre fumar e câncer poderia ser um artefato estatístico devido a uma variável de confundimento. Yerushalmy pensava em termos de um tipo de personalidade fumante, enquanto Fisher sugeria a possibilidade de um gene que predisporia as pessoas tanto a fumar quanto a desenvolver câncer de pulmão.

Ironicamente, pesquisadores de genômica descobriram em 2008 que Fisher tinha razão: há um “gene do fumo” que opera exatamente do modo que ele sugeriu. Essa descoberta veio por meio de uma nova técnica de análise genômica chamada estudo de associação genômica ampla (GWAS, do inglês \textit{genome-wide association study}, pronunciado “gí-uáss”). Esse é um método prototípico de \textit{big data} que permite aos pesquisadores pentearem estatisticamente o genoma inteiro, buscando genes que por acaso aparecem com mais frequência em pessoas com uma certa doença, como diabetes, esquizofrenia ou câncer de pulmão.

É importante notar a palavra “associação” no termo GWAS. Esse método não prova causalidade; apenas identifica genes associados a uma certa doença na amostra dada. É um método orientado por dados, e não por hipóteses, e isso apresenta problemas para a inferência causal. Embora estudos de genes orientados por hipóteses anteriores não tivessem encontrado nenhuma evidência clara de genes relacionados ao fumo ou ao câncer de pulmão, as coisas mudaram da noite para o dia em 2008. Naquele ano, pesquisadores identificaram um gene localizado em uma região do décimo quinto cromossomo que codifica receptores de nicotina em células pulmonares. Tem um nome oficial, rs16969968, mas isso é um trava-línguas mesmo para especialistas em genômica. Assim, começaram a chamá-lo de “the Big One” ou “Mr. Big” por causa de sua associação extremamente forte com o câncer de pulmão. “No campo do fumo, se você disser Mr. Big, as pessoas saberão do que você está falando”, diz Laura Bierut, especialista em tabagismo da Washington University em St. Louis. Eu simplesmente o chamarei de gene do fumo.

Neste ponto, penso que ouço o fantasma rabugento de R. A. Fisher chacoalhando as correntes no porão e exigindo uma retratação de todas as coisas que escrevi no Capítulo 5. Sim, o gene do fumo está associado ao câncer de pulmão. Tem duas variantes, uma comum e outra menos comum. Pessoas que herdam duas cópias da variante menos comum (cerca de um nono da população) têm cerca de 77 por cento de risco maior de contrair câncer de pulmão. O gene do fumo também parece relacionado ao comportamento de fumar. Pessoas que têm a variante de risco parecem precisar de mais nicotina para se sentirem satisfeitas e têm mais dificuldade em parar. Contudo, há também algumas boas notícias: essas pessoas respondem melhor à terapia de reposição de nicotina do que pessoas sem o gene do fumo.

A descoberta do gene do fumo não deveria mudar a opinião de ninguém sobre o fator causal esmagadoramente mais importante no câncer de pulmão, que é o tabagismo. Sabemos que fumar está associado a um aumento de mais de dez vezes no risco de contrair câncer de pulmão. Em comparação, mesmo uma dose dupla do gene do fumo menos do que dobra o seu risco. Isso é coisa séria, sem dúvida, mas não se compara ao perigo que se enfrenta (sem nenhuma boa razão) se se é um fumante regular.

\textbf{FIGURA 9.11.} Diagrama causal para o exemplo do gene do fumo.

Como sempre, ajuda visualizar a discussão com um diagrama causal. Fisher pensava no gene do fumo (naquela época puramente hipotético) como um confundidor do fumo e do câncer (Figura 9.11). Mas como confundidor, não é nem de longe forte o bastante para explicar o efeito esmagadoramente forte do fumo sobre o risco de câncer de pulmão. Esse é, em essência, o argumento que Jerome Cornfield fez em seu artigo de 1959 que encerrou a discussão sobre a hipótese genética.

\begin{figure}[ht]
\centering
\includegraphics[width=0.92\textwidth]{../figuras/fig_9_11.png}
\caption{Figura 9.11. ). Mas como confundidor, não é nem de longe forte o bastante para explicar o efeito esmagadoramente forte do fumo sobre o risco de câncer de pulmão. Esse é, em essência, o argumento que Jerome Cornfield fez em seu artigo de 1959 que encerrou a discussão sobre a hipótese genética.}
\end{figure}
\textit{Figura 9.11. ). Mas como confundidor, não é nem de longe forte o bastante para explicar o efeito esmagadoramente forte do fumo sobre o risco de câncer de pulmão. Esse é, em essência, o argumento que Jerome Cornfield fez em seu artigo de 1959 que encerrou a discussão sobre a hipótese genética.}

Podemos facilmente reescrever o mesmo diagrama causal como mostrado na Figura 9.12. Quando olhamos o diagrama desse modo, vemos que o comportamento de fumar é um mediador entre o gene do fumo e o câncer de pulmão. Essa minúscula mudança de perspectiva reformula completamente o debate científico. Em vez de perguntar se fumar causa câncer (uma pergunta cuja resposta conhecemos), perguntamos em vez disso como o gene funciona. Ele faz as pessoas fumarem mais e inalarem mais fundo? Ou de algum modo torna as células pulmonares mais vulneráveis ao câncer? Qual é mais forte, o efeito indireto ou o efeito direto?

\begin{figure}[ht]
\centering
\includegraphics[width=0.92\textwidth]{../figuras/fig_9_12.png}
\caption{Figura 9.12. Quando olhamos o diagrama desse modo, vemos que o comportamento de fumar é um mediador entre o gene do fumo e o câncer de pulmão. Essa minúscula mudança de perspectiva reformula completamente o debate científico. Em vez de perguntar se fumar causa câncer (uma pergunta cuja resposta conhecemos), perguntamos em vez disso como o gene funciona. Ele faz as pessoas fumarem mais e inalarem mais fundo? Ou de algum modo torna as células pulmonares mais vulneráveis ao câncer? Qual é mais forte, o efeito indireto ou o efeito direto?}
\end{figure}
\textit{Figura 9.12. Quando olhamos o diagrama desse modo, vemos que o comportamento de fumar é um mediador entre o gene do fumo e o câncer de pulmão. Essa minúscula mudança de perspectiva reformula completamente o debate científico. Em vez de perguntar se fumar causa câncer (uma pergunta cuja resposta conhecemos), perguntamos em vez disso como o gene funciona. Ele faz as pessoas fumarem mais e inalarem mais fundo? Ou de algum modo torna as células pulmonares mais vulneráveis ao câncer? Qual é mais forte, o efeito indireto ou o efeito direto?}

A resposta faz diferença para o tratamento. Se o efeito for direto, então pessoas que têm o gene de alto risco deveriam talvez receber triagem extra para câncer de pulmão. Por outro lado, se o efeito for indireto, o comportamento de fumar se torna crucial. Devemos aconselhar tais pacientes sobre o seu risco aumentado e a importância de não fumar em primeiro lugar. Se já fumam, podemos precisar intervir de modo mais agressivo, talvez com terapia de reposição de nicotina.

\textbf{FIGURA 9.12.} Figura 9.11, ligeiramente rearranjada.

Tyler VanderWeele, um epidemiologista da Harvard University, leu o primeiro relatório sobre o gene do fumo na \textit{Nature} e contatou um grupo de pesquisa em Harvard liderado por David Christiani. Desde 1992, Christiani pedia a seus pacientes com câncer de pulmão, bem como a amigos e familiares deles, que preenchessem questionários e fornecessem amostras de DNA para ajudar o esforço de pesquisa. Em meados dos anos 2000, ele havia coletado dados sobre 1.800 pacientes com câncer bem como 1.400 pessoas sem câncer de pulmão, que serviam como controles. As amostras de DNA ainda estavam esfriando em um freezer quando VanderWeele ligou.

Os resultados da análise de VanderWeele foram surpreendentes a princípio. Ele descobriu que o risco aumentado de câncer de pulmão devido ao efeito indireto era de apenas 1 a 3 por cento. As pessoas com a variante de alto risco do gene fumavam em média apenas um cigarro adicional por dia, o que não era o bastante para ser clinicamente relevante. Contudo, os corpos delas respondiam de modo diferente ao fumo. O efeito do gene do fumo sobre o câncer de pulmão era grande e significativo, mas apenas para aquelas pessoas que fumavam.

Isso coloca um interessante enigma no relato dos resultados. Nesse caso, CDE(0) seria essencialmente zero: se você não fuma, o gene não o prejudicará. Por outro lado, se fixarmos o mediador em um maço por dia ou dois maços por dia, o que eu denotaria CDE(1) ou CDE(2), então o efeito do gene é forte. O efeito direto natural faz a média desses efeitos controlados. Assim, o efeito direto natural, NDE, é positivo, e é assim que VanderWeele o reportou.

Esse exemplo é um caso de manual de interação. No fim, a análise de VanderWeele prova três coisas importantes sobre o gene do fumo. Primeiro, ele não aumenta significativamente o consumo de cigarros. Segundo, não causa câncer de pulmão por um caminho independente do fumo. Terceiro, para aquelas pessoas que fumam, aumenta significativamente o risco de câncer de pulmão. A interação entre o gene e o comportamento do sujeito é tudo.

Como é o caso com qualquer novo resultado, naturalmente mais pesquisa é necessária. Bierut aponta um problema com a análise de VanderWeele e Christiani: eles tinham apenas uma medida do comportamento de fumar — o número de cigarros por dia. O gene poderia possivelmente fazer as pessoas inalarem mais profundamente para obter uma dose maior de nicotina por tragada. O estudo de Harvard simplesmente não tinha dados para testar essa teoria.

Mesmo que alguma incerteza permaneça, a pesquisa sobre o gene do fumo oferece um vislumbre do futuro da medicina personalizada. Parece bastante claro que, nesse caso, o importante é como o gene e o comportamento interagem. Ainda não sabemos com certeza se o gene muda o comportamento (como Bierut sugere) ou meramente interage com um comportamento que teria acontecido de qualquer modo (como a análise de VanderWeele sugere). Não obstante, podemos ser capazes de usar o status genético para dar às pessoas melhor informação sobre os riscos que enfrentam. No futuro, modelos causais capazes de detectar interações entre genes e comportamento ou genes e ambiente certamente serão uma ferramenta importante no kit do epidemiologista.

\subsection{Torniquetes: uma falácia oculta}
Quando John Kragh, um cirurgião do exército, chegou para o seu primeiro dia de serviço em um hospital de Bagdá em 2006, recebeu um despertar imediato para as novas realidades da medicina de guerra. Vendo uma prancheta com os casos do dia, comentou com a enfermeira de plantão: “Ei, isso é interessante — vocês tiveram um torniquete de emergência usado durante o seu turno.” A enfermeira respondeu: “Isso não é interessante. Temos um a cada turno.”

Em seus primeiros cinco minutos no trabalho, Kragh havia tropeçado em uma mudança de maré no cuidado de traumas que ocorreu durante as guerras do Iraque e do Afeganistão. Embora usados por séculos, tanto no campo de batalha quanto na sala de operação, os torniquetes sempre foram um tanto controversos. Um torniquete deixado por tempo demais levará à perda de um membro. Além disso, os torniquetes frequentemente foram improvisados sob pressão, a partir de correias ou outros materiais à mão, de modo que a sua efetividade é, sem surpresa, uma questão de acerto ou erro. Depois da Segunda Guerra Mundial foram considerados um tratamento de último recurso, e o seu uso era oficialmente desencorajado.

As guerras do Iraque e do Afeganistão mudaram radicalmente essa política. Duas coisas aconteceram: mais das lesões graves exigiam o uso de torniquete e melhores desenhos de torniquete se tornaram disponíveis. Em 2005, o cirurgião-geral do Exército dos EUA recomendou que torniquetes pré-fabricados fossem fornecidos a todos os soldados. Por volta de 2006, como Kragh notou, a chegada de soldados feridos ao hospital com um torniquete em torno de um braço ou perna era uma ocorrência cotidiana — uma situação sem precedentes na história médica. De 2002 a 2012, estima Kragh, os torniquetes salvaram 2 mil vidas militares. Os soldados na linha de frente notaram. Segundo o cirurgião do Exército dos EUA David Welling: “Tropas de combate estão, segundo se reporta, saindo em missões de patrulha perigosas com torniquetes já colocados nas extremidades, pois desejam estar plenamente prontos para responder a sangramento de extremidade, se e quando a mina ou o artefato explosivo improvisado (IED) explodir.”

A julgar pela evidência anedótica e pela popularidade dos torniquetes entre os soldados da linha de frente, o seu valor deveria estar além de questão a esta altura. Contudo, poucos, se algum, estudos de grande escala sobre o uso de torniquetes haviam sido realizados. Na vida civil, os tipos de lesões que os necessitam são demasiado raros, e na vida militar o caos da guerra torna difícil conduzir um estudo científico adequado. Mas Kragh viu a oportunidade de documentar os efeitos do seu uso. Ele e as enfermeiras coletaram dados sobre cada caso que entrava pelas portas do hospital, e o antigo novato em torniquetes se tornou conhecido como “o cara do torniquete”.

Os resultados do estudo, publicados em 2015, não foram o que Kragh esperava. Segundo os dados, os pacientes que tiveram torniquetes aplicados antes de chegar ao hospital não sobreviveram a uma taxa mais alta do que aqueles com lesões similares que não receberam torniquetes. Naturalmente, raciocinou Kragh, os que tinham torniquetes possivelmente tinham lesões mais graves para começar. Mas mesmo quando controlou por esse fator comparando casos de igual gravidade, os torniquetes não pareciam melhorar as taxas de sobrevivência (veja a Tabela 9.1).

\begin{figure}[ht]
\centering
\includegraphics[width=0.92\textwidth]{../figuras/tab_9_1.png}
\caption{Tabela 9.1. ) pode confirmar ou refutar essa hipótese.}
\end{figure}
\textit{Tabela 9.1. ) pode confirmar ou refutar essa hipótese.}

\textbf{TABELA 9.1.} Dados sobre sobrevivência com e sem torniquetes.

Isso não é uma situação de paradoxo de Simpson. Não importa se agregamos os dados ou os estratificamos; em cada categoria de gravidade, bem como no agregado, a sobrevivência foi ligeiramente maior para soldados que \textit{não} receberam torniquetes. (A diferença nas taxas de sobrevivência era, contudo, demasiado pequena para ser estatisticamente significativa.)

O que deu errado? Uma possibilidade, naturalmente, é que os torniquetes não sejam melhores. Nossa crença neles poderia ser um caso de viés de confirmação. Quando um soldado recebe um torniquete e sobrevive, os seus médicos e os seus camaradas dirão: “Aquele torniquete salvou a vida dele.” Mas se o soldado não recebe um torniquete e sobrevive, ninguém dirá: “Não colocar um torniquete salvou a vida dele.” Assim, os torniquetes podem receber mais crédito do que merecem, e a não intervenção não recebe crédito algum.

Mas havia outro possível viés nesse estudo, que o próprio Kragh apontou: os médicos só coletaram dados sobre aqueles soldados que sobreviveram tempo suficiente para chegar ao hospital em primeiro lugar. Para ver por que isso importa, desenhemos um diagrama causal (Figura 9.13).

\begin{figure}[ht]
\centering
\includegraphics[width=0.92\textwidth]{../figuras/fig_9_13.png}
\caption{Figura 9.13. Esse efeito era essencialmente zero.}
\end{figure}
\textit{Figura 9.13. Esse efeito era essencialmente zero.}

\textbf{FIGURA 9.13.} Diagrama causal para o exemplo dos torniquetes. A linha tracejada é um efeito causal hipotético (não sustentado pelos dados).

Nessa figura, podemos ver que a Gravidade da Lesão é um confundidor de todas as três variáveis: o tratamento (Uso de Torniquete), o mediador (Sobrevivência Pré-Admissão) e o desfecho (Sobrevivência Pós-Admissão). É portanto apropriado e necessário condicionar à Gravidade da Lesão, como Kragh fez em seu artigo. Contudo, como Kragh estudou apenas os pacientes que de fato sobreviveram tempo suficiente para chegar ao hospital, ele também estava condicionando ao mediador, Sobrevivência Pré-Admissão. Em efeito, estava bloqueando o caminho indireto do uso de torniquete à sobrevivência pós-admissão e, portanto, estava calculando o efeito direto, indicado pela seta tracejada na Figura 9.13. Esse efeito era essencialmente zero.

Não obstante, ainda poderia haver um efeito indireto. Se os torniquetes permitiram que mais soldados sobrevivessem até chegar ao hospital, então o torniquete seria uma intervenção muito favorável. Isso significaria que o trabalho de um torniquete é levar o paciente vivo ao hospital; uma vez feito isso, não tem valor ulterior. Infelizmente, nada nos dados (Tabela 9.1) pode confirmar ou refutar essa hipótese.

William Kruskal certa vez lamentou que não haja um Homero para cantar o louvor dos estatísticos. Gostaria de cantar o louvor de Kragh, que sob as condições mais adversas imagináveis teve a presença de espírito de coletar dados e submeter o tratamento padrão a um teste científico. O seu exemplo é uma luz brilhante para qualquer um que queira praticar medicina baseada em evidências. É uma ironia particularmente amarga que o seu estudo não pudesse ter sucesso porque ele não tinha modo de coletar dados sobre soldados que não sobreviveram até o hospital. Podemos desejar que ele tivesse podido provar de uma vez por todas que os torniquetes salvam vidas. O próprio Kragh escreveu em um e-mail: “Não tenho dúvida de que os torniquetes são uma intervenção desejável.” Mas no fim teve de reportar um “resultado nulo”, o tipo que não faz manchetes. Ainda assim, merece crédito por sólidos instintos científicos.

\chapter{Capítulo 10 — Big data, inteligência artificial e as grandes perguntas}
\begin{quote}Tudo está predeterminado, e, no entanto, a permissão é sempre concedida.   — MAIMÔNIDES (MOSHE BEN MAIMON) (1138–1204)\end{quote}
QUANDO iniciei minha jornada pela causalidade, eu seguia as pegadas de uma anomalia. Com as redes bayesianas, tínhamos ensinado as máquinas a pensar em tons de cinza, e isso era um passo importante em direção a um pensamento semelhante ao humano. Mas ainda não conseguíamos ensinar as máquinas a compreender causas e efeitos. Não podíamos explicar a um computador por que girar o seletor de um barômetro não causa chuva. Tampouco podíamos ensinar-lhe o que esperar quando um dos fuzileiros de um pelotão de fuzilamento muda de ideia e decide não atirar. Sem a capacidade de imaginar realidades alternativas e contrastá-las com a realidade atualmente existente, uma máquina não consegue passar no mini-teste de Turing; não consegue responder à pergunta mais básica que nos torna humanos: “Por quê?”

Tomei isso como uma anomalia porque não antecipava que perguntas tão naturais e intuitivas residissem além do alcance dos sistemas de raciocínio mais avançados da época. Só mais tarde percebi que a mesma anomalia afligia mais do que apenas o campo da inteligência artificial (IA). As próprias pessoas que deveriam se importar mais com as perguntas “Por quê?” — a saber, os cientistas — laboravam sob uma cultura estatística que lhes negava o direito de fazer essas perguntas. É claro que as faziam de qualquer modo, informalmente, mas tinham de reformulá-las como perguntas associacionais sempre que queriam submetê-las à análise matemática.

A busca por essa anomalia me pôs em contato com pessoas de uma variedade de campos — como Clark Glymour e sua equipe (Richard Scheines e Peter Spirtes) da filosofia, Joseph Halpern da ciência da computação, Jamie Robins e Sander Greenland da epidemiologia, Chris Winship da sociologia, e Don Rubin e Philip Dawid da estatística — que pensavam sobre o mesmo problema. Juntos acendemos a faísca de uma Revolução Causal, que se espalhou como uma corrente de fogos de artifício de uma disciplina para a seguinte: epidemiologia, psicologia, genética, ecologia, geologia, ciência do clima e assim por diante. A cada ano que passa, vejo uma disposição cada vez maior entre os cientistas de falar e escrever sobre causas e efeitos — não com desculpas e olhar baixado, mas com confiança e assertividade. Evoluiu um novo paradigma segundo o qual é legítimo basear suas alegações em suposições, desde que se tornem as suposições transparentes, de modo que você e outros possam julgar quão plausíveis elas são e quão sensíveis são suas alegações à sua violação. A Revolução Causal talvez não tenha produzido nenhum gadget particular que tenha mudado nossas vidas, mas conduziu a uma transformação de atitudes que inevitavelmente levará a uma ciência mais saudável.

Penso com frequência nessa transformação como “o segundo presente da IA à humanidade”, e ela tem sido o foco principal deste livro. Mas, ao conduzirmos a história a uma conclusão, é hora de voltarmos e inquirirmos sobre o primeiro presente, que demorou um tempo inesperadamente longo para se materializar. Estamos de fato nos aproximando do dia em que computadores ou robôs possam compreender conversas causais? Podemos criar inteligências artificiais com tanta imaginação quanto uma criança humana de três anos? Compartilho algumas reflexões, mas sem conclusões definitivas, neste capítulo final.

\section{Modelos causais e o \textit{big data}}
Em toda a ciência, os negócios, o governo e até o esporte, a quantidade de dados brutos de que dispomos sobre nosso mundo cresceu a um ritmo assombroso nos últimos anos. A mudança é talvez mais visível para aqueles de nós que usam a internet e as mídias sociais. Em 2014, o último ano para o qual vi dados, o Facebook, segundo se noticiou, armazenava 300 petabytes de dados sobre seus 2 bilhões de usuários ativos — ou 150 megabytes de dados por usuário. Os jogos que as pessoas jogam, os produtos de que gostam de comprar, os nomes de todos os seus amigos no Facebook e, é claro, todos os seus vídeos de gatos — tudo isso está lá fora, em um oceano glorioso de uns e zeros.

Menos óbvio para o público, mas igualmente importante, é o surgimento de enormes bancos de dados na ciência. Por exemplo, o Projeto 1000 Genomas coletou duzentos terabytes de informação no que chama de “o maior catálogo público de variação humana e dados de genótipo”. O Arquivo Mikulski de Telescópios Espaciais da NASA coletou 2,5 petabytes de dados de vários levantamentos do espaço profundo. Mas o \textit{big data}[^1] não afetou apenas as ciências de alto perfil; fez incursões em todas as ciências. Uma geração atrás, um biólogo marinho poderia passar meses fazendo um censo de sua espécie favorita. Agora o mesmo biólogo tem acesso imediato \textit{online} a milhões de pontos de dados sobre peixes, ovos, conteúdo estomacal ou qualquer outra coisa que deseje. Em vez de apenas fazer um censo, o biólogo pode contar uma história.

O mais relevante para nós é a pergunta sobre o que vem a seguir. Como extrair significado de todos esses números, bits e pixels? Os dados podem ser imensos, mas as perguntas que fazemos são simples. Existe um gene que causa câncer de pulmão? Que tipos de sistemas solares têm probabilidade de abrigar planetas semelhantes à Terra? Que fatores estão fazendo a população de nosso peixe favorito diminuir, e o que podemos fazer a respeito?

Em certos círculos há uma fé quase religiosa de que podemos encontrar as respostas a essas perguntas nos próprios dados, se apenas formos suficientemente engenhosos na mineração de dados. Contudo, os leitores deste livro saberão que esse alarde tem grande probabilidade de estar equivocado. As perguntas que acabei de fazer são todas causais, e perguntas causais nunca podem ser respondidas a partir dos dados sozinhos. Elas exigem que formulemos um modelo do processo que gera os dados, ou ao menos de alguns aspectos desse processo. Sempre que você vir um artigo ou um estudo que analisa os dados de maneira livre de modelo, pode ter certeza de que o produto do estudo meramente resumirá e talvez transformará, mas não interpretará os dados.

Isso não quer dizer que a mineração de dados seja inútil. Ela pode ser um primeiro passo essencial para buscar padrões interessantes de associação e formular perguntas interpretativas mais precisas. Em vez de perguntar “Existem genes causadores de câncer de pulmão?”, podemos agora começar a varrer o genoma em busca de genes com alta correlação com o câncer de pulmão (como o gene “Sr. Big” mencionado no Capítulo 9). Então podemos perguntar: “Este gene causa câncer de pulmão? (E como?)” Nunca poderíamos ter perguntado sobre o gene “Sr. Big” se não dispuséssemos de mineração de dados. Para ir mais longe, porém, precisamos desenvolver um modelo causal especificando (por exemplo) que variáveis pensamos que o gene afeta, que confundidores podem existir e que outros caminhos causais podem produzir o resultado. Interpretar dados significa formular hipóteses sobre como as coisas operam no mundo real.

Outro papel do \textit{big data} em problemas de inferência causal reside no último estágio do motor de inferência descrito na Introdução (passo 8), que nos leva do estimando à estimativa. Esse passo de estimação estatística não é trivial quando o número de variáveis é grande, e apenas o \textit{big data} e as técnicas modernas de aprendizado de máquina podem nos ajudar a superar a maldição da dimensionalidade. Do mesmo modo, o \textit{big data} e a inferência causal juntos desempenham um papel crucial na área emergente da medicina personalizada. Aqui, buscamos fazer inferências a partir do comportamento passado de um conjunto de indivíduos que são semelhantes, em tantas características quanto possível, ao indivíduo em questão. A inferência causal nos permite descartar as características irrelevantes e recrutar esses indivíduos a partir de estudos diversos, enquanto o \textit{big data} nos permite reunir informação suficiente sobre eles.

É fácil compreender por que algumas pessoas veriam a mineração de dados como o ponto de chegada, e não o primeiro passo. Ela promete uma solução usando a tecnologia disponível. Poupa-nos, bem como às máquinas futuras, o trabalho de ter de considerar e articular suposições substantivas sobre como o mundo opera. Em alguns campos, nosso conhecimento pode estar em um estado tão embrionário que não temos a menor ideia de como começar a desenhar um modelo do mundo. Mas o \textit{big data} não resolverá esse problema. A parte mais importante da resposta deve vir de tal modelo, seja esboçado por nós, seja hipotetizado e afinado por máquinas.

Para que eu não pareça excessivamente crítico do empreendimento do \textit{big data}, gostaria de mencionar uma nova oportunidade de simbiose entre o \textit{big data} e a inferência causal. Chama-se transportabilidade. Graças ao \textit{big data}, não apenas podemos acessar um número enorme de indivíduos em qualquer estudo dado, mas também podemos acessar um número enorme de estudos, conduzidos em locais diferentes e sob condições diferentes. Com frequência queremos combinar os resultados desses estudos e traduzi-los para novas populações que podem ser diferentes até de maneiras que não antecipamos.

O processo de traduzir os resultados de um estudo de um cenário para outro é fundamental para a ciência. De fato, o progresso científico se estancaria se não fosse a capacidade de generalizar resultados de experimentos de laboratório para o mundo real — por exemplo, de tubos de ensaio para animais e para humanos. Mas até recentemente cada ciência tinha de desenvolver seus próprios critérios para separar generalizações válidas de inválidas, e não havia métodos sistemáticos para abordar a “transportabilidade” em geral.

Nos últimos cinco anos, meu ex-aluno (agora colega) Elias Bareinboim e eu conseguimos fornecer um critério completo para decidir quando os resultados são transportáveis e quando não o são. Como de costume, a condição para usar esse critério é que se representem as características salientes do processo gerador de dados com um diagrama causal, marcado com os locais de potenciais disparidades. “Transportar” um resultado não significa necessariamente tomá-lo pelo valor de face e aplicá-lo ao novo ambiente. O pesquisador pode ter de recalibrá-lo para acomodar disparidades entre os dois ambientes.

Suponha que queremos saber o efeito de um anúncio \textit{online} (\textit{X}) sobre a probabilidade de que um consumidor compre o produto (\textit{Y}) — digamos, uma prancha de surfe. Temos dados de estudos em cinco lugares diferentes: Los Angeles, Boston, São Francisco, Toronto e Honolulu. Agora queremos estimar quão eficaz será o anúncio em Arkansas. Infelizmente, cada população e cada estudo difere ligeiramente. Por exemplo, a população de Los Angeles é mais jovem do que nossa população-alvo, e a população de São Francisco difere na taxa de cliques. A Figura 10.1 mostra as características únicas de cada população e de cada estudo. Podemos combinar os dados desses estudos remotos e díspares para estimar a eficácia do anúncio em Arkansas? Podemos fazê-lo sem coletar nenhum dado em Arkansas? Ou talvez medindo apenas um pequeno conjunto de variáveis ou conduzindo um estudo observacional-piloto?

\begin{figure}[ht]
\centering
\includegraphics[width=0.92\textwidth]{../figuras/fig_10_1.png}
\caption{Figura 10.1. mostra as características únicas de cada população e de cada estudo. Podemos combinar os dados desses estudos remotos e díspares para estimar a eficácia do anúncio em Arkansas? Podemos fazê-lo sem coletar nenhum dado em Arkansas? Ou talvez medindo apenas um pequeno conjunto de variáveis ou conduzindo um estudo observacional-piloto?}
\end{figure}
\textit{Figura 10.1. mostra as características únicas de cada população e de cada estudo. Podemos combinar os dados desses estudos remotos e díspares para estimar a eficácia do anúncio em Arkansas? Podemos fazê-lo sem coletar nenhum dado em Arkansas? Ou talvez medindo apenas um pequeno conjunto de variáveis ou conduzindo um estudo observacional-piloto?}

\textbf{FIGURA 10.1.} O problema da transportabilidade.

A Figura 10.2 traduz essas diferenças em forma gráfica. A variável \textit{Z} representa a idade, que é um confundidor; jovens podem ter maior probabilidade de ver o anúncio e maior probabilidade de comprar o produto mesmo se não o virem. A variável \textit{W} representa clicar em um \textit{link} para obter mais informação. Este é um mediador, um passo que deve ocorrer para converter “ver o anúncio” em “comprar o produto”. A letra \textit{S}, em cada caso, representa uma variável “produtora de diferença”, uma variável hipotética que aponta para a característica pela qual as duas populações diferem. Por exemplo, em Los Angeles (b), o indicador \textit{S} aponta para \textit{Z}, a idade. Em cada uma das outras cidades, o indicador aponta para a característica distintiva da população mencionada na Figura 10.1.

\begin{figure}[ht]
\centering
\includegraphics[width=0.92\textwidth]{../figuras/fig_10_2.png}
\caption{Figura 10.2. traduz essas diferenças em forma gráfica. A variável \textit{Z} representa a idade, que é um confundidor; jovens podem ter maior probabilidade de ver o anúncio e maior probabilidade de comprar o produto mesmo se não o virem. A variável \textit{W} representa clicar em um \textit{link} para obter mais informação. Este é um mediador, um passo que deve ocorrer para converter “ver o anúncio” em “comprar o produto”. A letra \textit{S}, em cada caso, representa uma variável “produtora de diferença”, uma variável hipotética que aponta para a característica pela qual as duas populações diferem. Por exemplo, em Los Angeles (b), o indicador \textit{S} aponta para \textit{Z}, a idade. Em cada uma das outras cidades, o indicador aponta para a característica distintiva da população mencionada na Figura 10.1.}
\end{figure}
\textit{Figura 10.2. traduz essas diferenças em forma gráfica. A variável }Z\textit{ representa a idade, que é um confundidor; jovens podem ter maior probabilidade de ver o anúncio e maior probabilidade de comprar o produto mesmo se não o virem. A variável }W\textit{ representa clicar em um }link\textit{ para obter mais informação. Este é um mediador, um passo que deve ocorrer para converter “ver o anúncio” em “comprar o produto”. A letra }S\textit{, em cada caso, representa uma variável “produtora de diferença”, uma variável hipotética que aponta para a característica pela qual as duas populações diferem. Por exemplo, em Los Angeles (b), o indicador }S\textit{ aponta para }Z\textit{, a idade. Em cada uma das outras cidades, o indicador aponta para a característica distintiva da população mencionada na Figura 10.1.}

\textbf{FIGURA 10.2.} Diferenças entre as populações estudadas, expressas em forma gráfica.

Para a agência de publicidade, a boa notícia é que um computador pode agora gerir esse complicado problema de “fusão de dados” e, guiado pelo cálculo do \textit{do}, dizer-nos quais estudos podemos usar para responder à nossa consulta e por quais meios, bem como que informação precisamos coletar em Arkansas para sustentar a conclusão. Em alguns casos, o efeito pode se transportar diretamente, sem trabalho adicional e sem sequer pormos os pés em Arkansas. Por exemplo, o efeito do anúncio em Arkansas deve ser o mesmo que em Boston, porque, segundo o diagrama, Boston (c) difere de Arkansas apenas na variável \textit{V}, que não afeta nem o tratamento \textit{X} nem o desfecho \textit{Y}. Precisamos reponderar os dados em alguns dos outros estudos — por exemplo, para dar conta da estrutura etária diferente da população no estudo de Los Angeles (b). De modo interessante, o estudo experimental em Toronto (e) é suficiente para estimar nossa consulta em Arkansas apesar da disparidade em \textit{W}, se apenas pudermos medir \textit{X}, \textit{W} e \textit{Y} em Arkansas.

De modo notável, encontramos exemplos em que nenhum transporte é viável a partir de qualquer um dos estudos disponíveis; e, no entanto, a quantidade-alvo é, ainda assim, estimável a partir de sua combinação. Além disso, mesmo estudos que não são transportáveis não são inteiramente inúteis. Tome, por exemplo, o estudo de Honolulu (f) na Figura 10.2, que não é transportável devido à seta \textit{S} → \textit{Y}. A seta \textit{X} → \textit{W}, por outro lado, não está contaminada por \textit{S}, e assim os dados disponíveis de Honolulu podem ser usados para estimar \textit{P}(\textit{W} | \textit{X}). Combinando isso com estimativas de \textit{P}(\textit{W} | \textit{X}) de outros estudos, podemos aumentar a precisão dessa subexpressão. Combinando cuidadosamente tais subexpressões, podemos ser capazes de sintetizar uma estimativa global precisa da quantidade-alvo.

Embora em casos simples esses resultados sejam intuitivamente razoáveis, quando os diagramas se tornam mais complicados, precisamos da ajuda de um método formal. O cálculo do \textit{do} fornece um critério geral para determinar a transportabilidade em tais casos. A regra é bastante simples: se você puder realizar uma sequência válida de operações \textit{do} (usando as regras do Capítulo 7) que transforme a quantidade-alvo em outra expressão na qual qualquer fator envolvendo \textit{S} esteja livre de operadores \textit{do}, então a estimativa é transportável. A lógica é simples; qualquer tal fator pode ser estimado a partir dos dados disponíveis, sem contaminação pelo fator de disparidade \textit{S}.

Elias Bareinboim conseguiu fazer para o problema da transportabilidade o mesmo que Ilya Shpitser fez para o problema das intervenções. Ele desenvolveu um algoritmo que pode determinar automaticamente se o efeito que você busca é transportável, usando apenas critérios gráficos. Em outras palavras, pode dizer-lhe se a separação requerida de \textit{S} em relação aos operadores \textit{do} pode ser realizada ou não.

Os resultados de Bareinboim são empolgantes porque mudam o que antes era visto como uma ameaça à validade em uma oportunidade de alavancar os muitos estudos em que a participação não pode ser mandatória e onde, portanto, não podemos garantir que a população do estudo seja a mesma que a população de interesse. Em vez de ver a diferença entre populações como uma ameaça à “validade externa” de um estudo, temos agora uma metodologia para estabelecer validade em situações que antes teriam parecido sem esperança. É precisamente porque vivemos na era do \textit{big data} que temos acesso a informação sobre muitos estudos e sobre muitas das variáveis auxiliares (como \textit{Z} e \textit{W}) que nos permitirão transportar resultados de uma população para outra.

Mencionarei de passagem que Bareinboim também provou resultados análogos para outro problema que há muito atormenta os estatísticos: o viés de seleção. Esse tipo de viés ocorre quando o grupo amostral em estudo difere da população-alvo de alguma maneira relevante. Isso se parece muito com o problema da transportabilidade — e de fato o é, exceto por uma modificação muito importante: em vez de desenhar uma seta da variável indicadora \textit{S} para a variável afetada, desenhamos a seta em direção a \textit{S}. Podemos pensar em \textit{S} como representando “seleção” (para o estudo). Por exemplo, se nosso estudo observa apenas pacientes hospitalizados, como no exemplo do viés de Berkson, desenharíamos uma seta de Hospitalização para \textit{S}, indicando que a hospitalização é uma causa de seleção para o nosso estudo. No Capítulo 6, vimos essa situação apenas como uma ameaça à validade do nosso estudo. Mas agora podemos olhá-la como uma oportunidade. Se compreendermos o mecanismo pelo qual recrutamos sujeitos para o estudo, podemos nos recuperar do viés coletando dados sobre o conjunto certo de desconfundidores e usando uma fórmula apropriada de reponderação ou ajustamento.

O trabalho de Bareinboim nos permite explorar a lógica causal e o \textit{big data} para realizar milagres que antes eram inconcebíveis. Palavras como “milagres” e “inconcebível” são raras no discurso científico, e o leitor pode se perguntar se estou sendo um pouco entusiasta demais. Mas as uso por uma boa razão. O conceito de validade externa como ameaça à ciência experimental existe há pelo menos meio século, desde que Donald Campbell e Julian Stanley reconheceram e definiram o termo em 1963. Conversei com dezenas de especialistas e autores proeminentes que escreveram sobre esse tópico. Para meu assombro, nenhum deles foi capaz de enfrentar qualquer um dos problemas-brinquedo apresentados na Figura 10.2. Chamo-os de “problemas-brinquedo” porque são fáceis de descrever, fáceis de resolver e fáceis de verificar se uma solução dada está correta. No presente, a cultura da “validade externa” está totalmente preocupada em listar e categorizar as ameaças à validade, em vez de combatê-las. De fato, está tão paralisada pelas ameaças que olha com suspeita e descrença a própria ideia de que as ameaças possam ser desarmadas. Os especialistas, que são noviços em modelos gráficos, acham mais fácil configurar ameaças adicionais do que tentar remediar qualquer uma delas. Uma linguagem como “milagres”, assim espero, deve sacudir meus colegas para que olhem para tais problemas como desafios intelectuais, e não como razões para o desespero.

Gostaria de poder apresentar ao leitor estudos de caso bem-sucedidos de uma tarefa complexa de transportabilidade e de recuperação do viés de seleção, mas as técnicas ainda são novas demais para terem penetrado no uso geral. Estou muito confiante, porém, de que os pesquisadores descobrirão o poder dos algoritmos de Bareinboim antes de muito tempo, e então a validade externa, como o confundimento antes dela, deixará de ter seu poder místico e aterrorizante.

\section{IA forte e livre-arbítrio}
A tinta mal havia secado no grande artigo de Alan Turing, “Computing Machinery and Intelligence” [Maquinário computacional e inteligência], quando escritores de ficção científica e futurólogos começaram a brincar com a perspectiva de máquinas que pensam. Às vezes imaginavam essas máquinas como figuras benignas ou até nobres, como o zumbidor e tagarela R2-D2 e o androide estranhamente britânico C-3PO de \textit{Guerra nas Estrelas}. Outras vezes as máquinas são bem mais sinistras, tramando a destruição da espécie humana, como nos filmes de \textit{O Exterminador do Futuro}, ou escravizando humanos em uma realidade virtual, como em \textit{Matrix}. Em todos esses casos, as IAs dizem mais sobre as ansiedades dos escritores ou sobre as capacidades do departamento de efeitos especiais do filme do que sobre a pesquisa real em inteligência artificial.

A inteligência artificial revelou-se um objetivo mais elusivo do que Turing jamais suspeitou, mesmo que o puro poder computacional de nossos computadores sem dúvida tenha excedido suas expectativas. No Capítulo 3, escrevi sobre algumas das razões para esse progresso lento. Nos anos 1970 e início dos 1980, a pesquisa em inteligência artificial era tolhida por seu foco em sistemas baseados em regras. Mas os sistemas baseados em regras provaram estar no caminho errado. Eram muito frágeis. Qualquer pequena mudança em suas suposições de trabalho exigia que fossem reescritos. Não conseguiam lidar bem com incerteza ou com dados contraditórios. Por fim, não eram cientificamente transparentes; você não podia provar matematicamente que se comportariam de certa maneira, e não podia apontar exatamente o que precisava de reparo quando não se comportavam.

Nem todos os pesquisadores de IA se opunham à falta de transparência. O campo na época estava dividido entre os “\textit{neats}” (que queriam sistemas transparentes com garantias de comportamento) e os “\textit{scruffies}” (que apenas queriam algo que funcionasse). Eu sempre fui um \textit{neat}. Tive a sorte de surgir em um momento em que o campo estava pronto para uma nova abordagem. As redes bayesianas eram probabilísticas; podiam lidar com um mundo cheio de dados conflitantes e incertos. Diferentemente dos sistemas baseados em regras, eram modulares e facilmente implementáveis em uma plataforma de computação distribuída, o que as tornava rápidas. Por fim, o que era importante para mim (e para outros \textit{neats}), as redes bayesianas lidavam com probabilidades de maneira matematicamente sólida. Isso garantia que, se algo desse errado, o \textit{bug} estava no programa, e não em nosso pensamento.

Mesmo com todas essas vantagens, as redes bayesianas ainda não conseguiam compreender causas e efeitos. Por projeto, em uma rede bayesiana, a informação flui em ambas as direções, causal e diagnóstica: a fumaça aumenta a verossimilhança de fogo, e o fogo aumenta a verossimilhança de fumaça. De fato, uma rede bayesiana nem sequer consegue dizer qual é a “direção causal”. A busca por essa anomalia — essa maravilhosa anomalia, como se revelou — me afastou do campo do aprendizado de máquina e me conduziu ao estudo da causalidade. Eu não conseguia me reconciliar com a ideia de que robôs futuros não seriam capazes de se comunicar conosco em nossa linguagem nativa de causa e efeito. Uma vez na terra da causalidade, fui naturalmente atraído para o vasto espectro de outras ciências em que a assimetria causal é da maior importância. Assim, nos últimos vinte e cinco anos, tenho sido algo como um expatriado da terra do raciocínio automatizado e do aprendizado de máquina. Não obstante, do meu ponto de observação distante, ainda posso ver as tendências e modas atuais.

Nos últimos anos, o progresso mais notável em IA ocorreu em uma área chamada “aprendizado profundo”, que usa métodos como redes neurais convolucionais. Essas redes não seguem as regras da probabilidade; não lidam com a incerteza de maneira rigorosa ou transparente. Menos ainda incorporam qualquer representação explícita do ambiente em que operam. Em vez disso, a arquitetura da rede é deixada livre para evoluir por si mesma. Ao terminar de treinar uma nova rede, o programador não tem ideia de que computações ela está realizando ou por que funcionam. Se a rede falha, ela não tem ideia de como consertá-la.

Talvez o exemplo prototípico seja o AlphaGo, um programa baseado em redes neurais convolucionais que joga o antigo jogo asiático de go, desenvolvido pela DeepMind, uma subsidiária do Google. Entre os jogos humanos de informação perfeita, o go sempre havia sido considerado o nó mais difícil para a IA. Embora os computadores tenham conquistado os humanos no xadrez em 1997, não eram considerados páreo sequer para os jogadores profissionais de go de nível mais baixo ainda em 2015. A comunidade do go pensava que os computadores ainda estavam a uma década ou mais de oferecer aos humanos uma batalha real. Isso mudou quase da noite para o dia com o advento do AlphaGo. A maioria dos jogadores de go ouviu falar do programa pela primeira vez no final de 2015, quando ele derrotou um profissional humano por 5 a 0. Em março de 2016, o AlphaGo derrotou Lee Sedol, por anos considerado o jogador humano mais forte, por 4 a 1. Alguns meses depois, jogou sessenta partidas \textit{online} contra os melhores jogadores humanos sem perder uma única, e em 2017 foi oficialmente aposentado após vencer o então campeão mundial, Ke Jie. A única partida que perdeu para Sedol é a única que jamais perderá para um humano.

Tudo isso é empolgante, e os resultados não deixam dúvida: o aprendizado profundo funciona para certas tarefas. Mas é a antítese da transparência. Mesmo os programadores do AlphaGo não conseguem dizer por que o programa joga tão bem. Sabiam, pela experiência, que redes profundas tinham tido sucesso em tarefas de visão computacional e reconhecimento de fala. Não obstante, nossa compreensão do aprendizado profundo é completamente empírica e não vem com garantias. A equipe do AlphaGo não poderia ter previsto no início que o programa derrotaria o melhor humano em um ano, ou em dois, ou em cinco. Simplesmente experimentaram, e ele o fez.

Algumas pessoas argumentarão que a transparência não é realmente necessária. Não compreendemos em detalhe como o cérebro humano funciona e, no entanto, ele funciona bem, e perdoamos nossa magra compreensão. Assim, argumentam, por que não liberar sistemas de aprendizado profundo e criar um novo tipo de inteligência sem compreender como funciona? Não posso dizer que estejam erradas. Os \textit{scruffies}, neste momento, tomaram a dianteira. Não obstante, posso dizer que eu pessoalmente não gosto de sistemas opacos, e é por isso que não escolho pesquisá-los.

Deixando de lado meu gosto pessoal, há outro fator a acrescentar a essa analogia com o cérebro humano. Sim, perdoamos nossa magra compreensão de como os cérebros humanos funcionam, mas ainda podemos nos comunicar com outros humanos, aprender com eles, instruí-los e motivá-los em nossa própria linguagem nativa de causa e efeito. Podemos fazer isso porque nossos cérebros funcionam da mesma maneira. Se todos os nossos robôs forem tão opacos quanto o AlphaGo, não seremos capazes de manter uma conversa significativa com eles, e isso seria bastante lamentável. Quando o robô da minha casa liga o aspirador de pó enquanto ainda estou dormindo (Figura 10.3) e eu lhe digo “Você não deveria ter me acordado”, quero que ele compreenda que o problema foi a aspiração, mas não quero que interprete a reclamação como uma instrução para nunca mais aspirar o andar de cima. Deve compreender o que você e eu compreendemos perfeitamente: aspiradores de pó fazem barulho, barulho acorda as pessoas, e isso deixa algumas pessoas infelizes. Em outras palavras, nosso robô terá de compreender relações de causa e efeito — de fato, relações contrafactuais, como as codificadas na frase “Você não deveria ter”. Observe, de fato, o rico conteúdo dessa curta sentença de instruções. Não deveríamos precisar dizer ao robô que o mesmo se aplica à aspiração no andar de baixo ou em qualquer outro lugar da casa, mas não quando estou acordado ou não estou em casa, quando o aspirador está equipado com um silenciador, e assim por diante. Um programa de aprendizado profundo consegue compreender a riqueza dessa instrução? É por isso que não me satisfaço com o desempenho aparentemente soberbo de sistemas opacos. A transparência possibilita a comunicação eficaz.

\begin{figure}[ht]
\centering
\includegraphics[width=0.92\textwidth]{../figuras/fig_10_3.png}
\caption{Figura 10.3. ) e eu lhe digo “Você não deveria ter me acordado”, quero que ele compreenda que o problema foi a aspiração, mas não quero que interprete a reclamação como uma instrução para nunca mais aspirar o andar de cima. Deve compreender o que você e eu compreendemos perfeitamente: aspiradores de pó fazem barulho, barulho acorda as pessoas, e isso deixa algumas pessoas infelizes. Em outras palavras, nosso robô terá de compreender relações de causa e efeito — de fato, relações contrafactuais, como as codificadas na frase “Você não deveria ter”. Observe, de fato, o rico conteúdo dessa curta sentença de instruções. Não deveríamos precisar dizer ao robô que o mesmo se aplica à aspiração no andar de baixo ou em qualquer outro lugar da casa, mas não quando estou acordado ou não estou em casa, quando o aspirador está equipado com um silenciador, e assim por diante. Um programa de aprendizado profundo consegue compreender a riqueza dessa instrução? É por isso que não me satisfaço com o desempenho aparentemente soberbo de sistemas opacos. A transparência possibilita a comunicação eficaz.}
\end{figure}
\textit{Figura 10.3. ) e eu lhe digo “Você não deveria ter me acordado”, quero que ele compreenda que o problema foi a aspiração, mas não quero que interprete a reclamação como uma instrução para nunca mais aspirar o andar de cima. Deve compreender o que você e eu compreendemos perfeitamente: aspiradores de pó fazem barulho, barulho acorda as pessoas, e isso deixa algumas pessoas infelizes. Em outras palavras, nosso robô terá de compreender relações de causa e efeito — de fato, relações contrafactuais, como as codificadas na frase “Você não deveria ter”. Observe, de fato, o rico conteúdo dessa curta sentença de instruções. Não deveríamos precisar dizer ao robô que o mesmo se aplica à aspiração no andar de baixo ou em qualquer outro lugar da casa, mas não quando estou acordado ou não estou em casa, quando o aspirador está equipado com um silenciador, e assim por diante. Um programa de aprendizado profundo consegue compreender a riqueza dessa instrução? É por isso que não me satisfaço com o desempenho aparentemente soberbo de sistemas opacos. A transparência possibilita a comunicação eficaz.}

\textbf{FIGURA 10.3.} Um robô inteligente contemplando as ramificações causais de suas ações. (Fonte: desenho de Maayan Harel.)

Um aspecto do aprendizado profundo de fato me interessa: as limitações teóricas desses sistemas, primariamente limitações que decorrem de sua incapacidade de ir além do primeiro degrau da Escada da Causalidade. Essa limitação não prejudica o desempenho do AlphaGo no mundo estreito dos jogos de go, uma vez que a descrição do tabuleiro, juntamente com as regras do jogo, constitui um modelo causal adequado do mundo-do-go. No entanto, prejudica sistemas de aprendizado que operam em ambientes governados por ricas teias de forças causais, tendo acesso meramente a manifestações de superfície dessas forças. Medicina, economia, educação, climatologia e assuntos sociais são exemplos típicos de tais ambientes. Como os prisioneiros na famosa caverna de Platão, os sistemas de aprendizado profundo exploram as sombras na parede da caverna e aprendem a predizer com precisão seus movimentos. Carecem da compreensão de que as sombras observadas são meras projeções de objetos tridimensionais movendo-se em um espaço tridimensional. A IA forte exige essa compreensão.

Os pesquisadores de aprendizado profundo não desconhecem essas limitações básicas. Por exemplo, economistas que usam aprendizado de máquina notaram que seus métodos não respondem a perguntas-chave de interesse, como estimar o impacto de políticas e ações ainda não ensaiadas. Exemplos típicos são introduzir novas estruturas de preços ou subsídios, ou mudar o salário mínimo. Em termos técnicos, os métodos de aprendizado de máquina de hoje nos fornecem uma maneira eficiente de ir de estimativas de amostras finitas a distribuições de probabilidade, e ainda precisamos ir das distribuições às relações de causa e efeito.

Quando começamos a falar de IA forte, os modelos causais passam de luxo a necessidade. Para mim, uma IA forte deveria ser uma máquina capaz de refletir sobre suas ações e aprender com erros passados. Deveria ser capaz de compreender o enunciado “Eu deveria ter agido de outro modo”, seja ela informada disso por um humano, seja chegando a essa conclusão por si mesma. A interpretação contrafactual desse enunciado se lê: “Fiz \textit{X} = \textit{x}, e o desfecho foi \textit{Y} = \textit{y}. Mas se eu tivesse agido de outro modo, digamos \textit{X} = \textit{x}′, então o desfecho teria sido melhor, talvez \textit{Y} = \textit{y}′.” Como vimos, a estimação de tais probabilidades foi completamente automatizada, dados dados suficientes e um modelo causal adequadamente especificado.

De fato, penso que um alvo muito importante para o aprendizado de máquina é a probabilidade mais simples \textit{P}(\textit{Y}<sub>\textit{X} = \textit{x}′</sub> = \textit{y}′ | \textit{X} = \textit{x}), em que a máquina observa um evento \textit{X} = \textit{x} mas não o desfecho \textit{Y}, e então pergunta pelo desfecho sob um evento alternativo \textit{X} = \textit{x}′. Se puder computar essa quantidade, a máquina pode tratar sua ação pretendida como um evento observado (\textit{X} = \textit{x}) e perguntar: “E se eu mudar de ideia e fizer \textit{X} = \textit{x}′ em vez disso?” Essa expressão é matematicamente a mesma que o efeito do tratamento nos tratados (mencionado no Capítulo 8), e dispomos de muitos resultados indicando como estimá-la.

A intenção é uma parte muito importante da tomada de decisão pessoal. Se um ex-fumante se sente tentado a acender um cigarro, deveria pensar muito a fundo sobre as razões por trás dessa intenção e perguntar se uma ação contrária poderia de fato conduzir a um desfecho melhor. A capacidade de conceber a própria intenção e então usá-la como um pedaço de evidência no raciocínio causal é um nível de autoconsciência (se não de consciência) que nenhuma máquina de que eu tenha conhecimento alcançou. Eu gostaria de ser capaz de conduzir uma máquina à tentação e fazê-la dizer “Não”.

Qualquer discussão de intenção conduz a outra questão maior para a IA forte: o livre-arbítrio. Se estamos pedindo a uma máquina que tenha a intenção de fazer \textit{X} = \textit{x}, torne-se ciente disso e escolha fazer \textit{X} = \textit{x}′ em vez disso, parecemos estar pedindo que ela tenha livre-arbítrio. Mas como um robô pode ter livre-arbítrio se apenas segue instruções armazenadas em seu programa?

O filósofo de Berkeley John Searle rotulou o problema do livre-arbítrio como “um escândalo na filosofia”, em parte devido ao progresso zero feito sobre o problema desde a Antiguidade e em parte porque não podemos descartá-lo como uma ilusão de ótica. Toda a nossa concepção de “eu” pressupõe que temos algo como escolhas. Por exemplo, parece não haver maneira de reconciliar minha vivida e inconfundível sensação de ter uma opção (digamos, de tocar ou não tocar meu nariz) com minha compreensão da realidade que pressupõe determinismo causal: todas as nossas ações são desencadeadas por sinais neurais elétricos emanados do cérebro. Enquanto muitos problemas filosóficos desapareceram ao longo do tempo à luz do progresso científico, o livre-arbítrio permanece teimosamente enigmático, tão fresco quanto pareceu a Aristóteles e a Maimônides. Além disso, enquanto o livre-arbítrio humano às vezes foi justificado em bases espirituais ou teológicas, essas explicações não se aplicariam a uma máquina programada. Assim, qualquer aparência de livre-arbítrio robótico deve ser um truque — ao menos é esse o dogma convencional.

Nem todos os filósofos estão convencidos de que realmente haja um choque entre livre-arbítrio e determinismo. Um grupo chamado “compatibilistas”, entre os quais me conto, considera-o apenas um choque aparente entre dois níveis de descrição: o nível neural em que os processos parecem determinísticos (ressalvado o indeterminismo quântico) e o nível cognitivo em que temos uma vivida sensação de opções. Tais choques aparentes não são infrequentes na ciência. Por exemplo, as equações da física são reversíveis no tempo em um nível microscópico e, no entanto, parecem irreversíveis no nível macroscópico de descrição; a fumaça nunca flui de volta para a chaminé. Mas isso abre novas perguntas: concedido que o livre-arbítrio é (ou pode ser) uma ilusão, por que é tão importante para nós, como humanos, ter essa ilusão? Por que a evolução laborou para nos dotar dessa concepção? Truque ou não, deveríamos programar a próxima geração de computadores para ter essa ilusão? Para quê? Que benefícios computacionais ela acarreta?

Penso que compreender os benefícios da ilusão de livre-arbítrio é a chave para o teimosamente enigmático problema de reconciliá-lo com o determinismo. O problema se dissolverá diante de nossos olhos uma vez que doteemos uma máquina determinística com os mesmos benefícios. Juntamente com essa questão funcional, devemos também lidar com perguntas de simulação. Se sinais neurais do cérebro desencadeiam todas as nossas ações, então nossos cérebros devem estar bastante ocupados decorando algumas ações com o título “volitiva” ou “intencional” e outras com “não intencional”. Qual é precisamente esse processo de rotulagem? Que caminho neural faria um dado sinal merecer o rótulo “volitivo”?

Em muitos casos, ações voluntárias são reconhecidas por um traço que deixam na memória de curto prazo, com o traço refletindo um propósito ou motivação. Por exemplo: “Por que você fez isso?” “Porque eu queria impressioná-lo.” Ou, como Eva inocentemente respondeu: “A serpente me enganou, e eu comi.” Mas em muitos outros casos uma ação intencional é tomada e, no entanto, nenhuma razão ou motivo vem à mente. A racionalização de ações pode ser um processo reconstrutivo, pós-ação. Por exemplo, um jogador de futebol pode explicar por que decidiu passar a bola para Joe em vez de Charlie, mas raramente é o caso de que essas razões tenham conscientemente desencadeado a ação. No calor do jogo, milhares de sinais de entrada competem pela atenção do jogador. A decisão crucial é quais sinais priorizar, e as razões dificilmente podem ser recordadas e articuladas.

Pesquisadores de IA estão, portanto, tentando responder a duas perguntas — sobre função e simulação — com a primeira impulsionando a segunda. Uma vez que compreendamos que função computacional o livre-arbítrio serve em nossas vidas, então podemos nos ocupar de equipar máquinas com tais funções. Torna-se um problema de engenharia, ainda que um difícil.

Para mim, certos aspectos da questão funcional se destacam claramente. A ilusão de livre-arbítrio nos dá a capacidade de falar sobre nossas intenções e de submetê-las ao pensamento racional, possivelmente usando lógica contrafactual. Quando o técnico nos tira de um jogo de futebol e diz “Você deveria ter passado a bola para Charlie”, considere todos os significados complexos embutidos nessas oito palavras. Primeiro, o propósito de tal instrução de “deveria ter” é transmitir rapidamente informação valiosa do técnico ao jogador: no futuro, quando enfrentar uma situação semelhante, escolha a ação \textit{B} em vez da ação \textit{A}. Mas as “situações semelhantes” são numerosas demais para listar e são pouco conhecidas até mesmo do próprio técnico. Em vez de listar as características dessas “situações semelhantes”, o técnico aponta para a ação do jogador, que é representativa de sua intenção no momento da decisão. Ao proclamar a ação inadequada, o técnico está pedindo ao jogador que identifique os pacotes de \textit{software} que conduziram à sua decisão e então reajuste as prioridades entre esses pacotes de modo que “passar para Charlie” se torne a ação preferida. Há uma sabedoria profunda nessa instrução porque quem, se não o próprio jogador, saberia as identidades desses pacotes? São caminhos neurais sem nome que não podem ser referenciados pelo técnico ou por qualquer observador externo. Pedir ao jogador que tome uma ação diferente daquela tomada equivale a encorajar uma análise específica de intenção, como a que mencionamos acima. Pensar em termos de intenções, portanto, oferece-nos uma abreviatura para converter instruções causais complicadas em simples.

Eu conjecturaria, então, que um time de robôs jogaria futebol melhor se fosse programado para se comunicar como se tivesse livre-arbítrio. Não importa quão tecnicamente proficientes sejam os robôs individuais no futebol, o desempenho de seu time melhorará quando puderem falar uns com os outros como se não fossem robôs pré-programados, mas agentes autônomos acreditando que têm opções.

Embora ainda reste ver se a ilusão de livre-arbítrio aprimora a comunicação robô-a-robô, há bem menos incerteza sobre a comunicação robô-a-humano. A fim de se comunicar naturalmente com humanos, IAs fortes certamente precisarão compreender o vocabulário de opções e intenções e, assim, precisarão emular a ilusão de livre-arbítrio. Como expliquei acima, podem também achar vantajoso “acreditar” em seu próprio livre-arbítrio, a ponto de serem capazes de observar sua intenção e agir de outro modo.

A capacidade de raciocinar sobre as próprias crenças, intenções e desejos tem sido um grande desafio para pesquisadores de IA e define a noção de “agência”. Filósofos, por outro lado, estudaram essas capacidades como parte da clássica questão da consciência. Perguntas como “Máquinas podem ter consciência?” ou “O que torna um agente de \textit{software} diferente de um programa ordinário?” engajaram as melhores mentes de muitas gerações, e eu não pretenderia respondê-las por completo. Acredito, não obstante, que a algoritmização dos contrafactuais é um grande passo em direção à compreensão dessas perguntas e a tornar a consciência e a agência uma realidade computacional.

Os métodos descritos para equipar uma máquina com uma representação simbólica de seu ambiente e a capacidade de imaginar uma perturbação hipotética desse ambiente podem ser estendidos para incluir a própria máquina como parte do ambiente. Nenhuma máquina pode processar uma cópia completa de seu próprio \textit{software}, mas pode ter um resumo esquemático de seus principais componentes de \textit{software}. Outros componentes podem então raciocinar sobre esse esquema e imitar um estado de autoconsciência.

Para criar a percepção de agência, devemos também equipar esse pacote de \textit{software} com uma memória para registrar ativações passadas, à qual possa se referir quando se perguntar: “Por que você fez isso?” Ações que passam por certos padrões de ativação de caminhos receberão explicações raciocinadas, como “Porque a alternativa se mostrou menos atraente”. Outras terminarão com respostas evasivas e inúteis, como “Quem dera eu soubesse por quê” ou “Porque foi assim que você me programou”.

Em resumo, acredito que o pacote de \textit{software} que pode dar a uma máquina pensante os benefícios da agência consistiria em pelo menos três partes: um modelo causal do mundo; um modelo causal de seu próprio \textit{software}, por mais superficial que seja; e uma memória que registra como intenções em sua mente correspondem a eventos no mundo exterior.

Isso pode mesmo ser como nossa própria educação causal como bebês começa. Podemos ter algo como um “gerador de intenções” em nossas mentes, que nos diz que devemos tomar a ação \textit{X} = \textit{x}. Mas as crianças adoram experimentar — desafiar as intenções iniciais dos pais, dos professores e até as suas próprias — e fazer algo diferente, só por diversão. Plenamente cientes de que devemos fazer \textit{X} = \textit{x}, brincalhonamente fazemos \textit{X} = \textit{x}′ em vez disso. Observamos o que acontece, repetimos o processo e mantemos um registro de quão bom é nosso gerador de intenções. Por fim, quando começamos a ajustar nosso próprio \textit{software}, é então que começamos a assumir responsabilidade moral por nossas ações. Essa responsabilidade pode ser uma ilusão no nível da ativação neural, mas não no nível do \textit{software} de autoconsciência.

Encorajado por essas possibilidades, acredito que a IA forte com compreensão causal e capacidades de agência é uma promessa realizável, e isso levanta a pergunta que escritores de ficção científica vêm fazendo desde os anos 1950: deveríamos nos preocupar? A IA forte é uma caixa de Pandora que não deveríamos abrir? Recentemente, figuras públicas como Elon Musk e Stephen Hawking se pronunciaram dizendo que deveríamos nos preocupar. No Twitter, Musk disse que as IAs eram “potencialmente mais perigosas do que bombas nucleares”. Em 2015, o \textit{site} Edge.org de John Brockman propôs como sua pergunta anual daquele ano: “O que você pensa sobre máquinas que pensam?” Reuniu 186 respostas ponderadas e provocativas (desde então reunidas em um livro intitulado \textit{What to Think About Machines That Think} [O que pensar sobre máquinas que pensam]).

A pergunta intencionalmente vaga de Brockman pode ser subdividida em ao menos cinco relacionadas:

1. Já fizemos máquinas que pensam? 2. Podemos fazer máquinas que pensam? 3. Faremos máquinas que pensam? 4. Deveríamos fazer máquinas que pensam?

E, por fim, a pergunta não enunciada que está no coração de nossas ansiedades:

5. Podemos fazer máquinas capazes de distinguir o bem do mal?

A resposta à primeira pergunta é não, mas acredito que a resposta a todas as outras é sim. Certamente ainda não fizemos máquinas que pensam em qualquer interpretação humanamente semelhante da palavra. Até agora só podemos simular o pensamento humano em domínios estreitamente definidos que têm apenas as estruturas causais mais primitivas. Lá podemos de fato fazer máquinas que superam humanos, mas isso não deveria ser surpresa porque esses domínios recompensam a única coisa que os computadores fazem bem: computar.

A resposta à segunda pergunta é quase certamente sim, se definirmos pensar como ser capaz de passar no teste de Turing. Digo isso com base no que aprendemos com o mini-teste de Turing. A capacidade de responder a consultas em todos os três níveis da Escada da Causalidade fornece as sementes do \textit{software} de “agência”, de modo que a máquina possa pensar sobre suas próprias intenções e refletir sobre seus próprios erros. Os algoritmos para responder a consultas causais e contrafactuais já existem (graças em grande parte aos meus estudantes) e estão apenas esperando por industriosos pesquisadores de IA para implementá-los.

A terceira pergunta depende, é claro, de eventos humanos que são difíceis de predizer. Mas, historicamente, os humanos raramente se abstiveram de fazer ou de realizar coisas de que são tecnologicamente capazes. Em parte isso se deve ao fato de que não sabemos que somos tecnologicamente capazes de algo até de fato o fazermos, seja clonar animais, seja enviar astronautas à Lua. A detonação da bomba atômica, contudo, foi um ponto de inflexão: muitas pessoas pensam que essa tecnologia não deveria ter sido desenvolvida. Desde a Segunda Guerra Mundial, um bom exemplo de cientistas recuando do factível foi a conferência de Asilomar de 1975 sobre recombinação de DNA, uma nova tecnologia vista pela mídia em termos um tanto apocalípticos. Os cientistas que trabalhavam no campo conseguiram chegar a um consenso sobre práticas de segurança de bom senso, e o acordo a que chegaram então se sustentou bem ao longo das quatro décadas seguintes. O DNA recombinante é agora uma tecnologia comum e madura.

Em 2017, o Future of Life Institute convocou uma conferência de Asilomar semelhante sobre inteligência artificial e concordou com um conjunto de vinte e três princípios para pesquisa futura em “IA benéfica”. Embora a maioria das diretrizes não seja relevante para os tópicos discutidos neste livro, as recomendações sobre ética e valores são definitivamente dignas de atenção. Por exemplo, as recomendações 6 — “Sistemas de IA devem ser seguros e protegidos ao longo de sua vida operacional, e de modo verificável” — e 7 — “Se um sistema de IA causa dano, deve ser possível apurar por quê” — claramente falam da importância da transparência. A recomendação 10 — “Sistemas de IA altamente autônomos devem ser projetados de modo que seus objetivos e comportamentos possam ser assegurados a se alinharem com valores humanos ao longo de sua operação” — é um tanto vaga como enunciada, mas poderia receber significado operacional se se exigisse que esses sistemas fossem capazes de declarar suas próprias intenções e de se comunicar com humanos sobre causas e efeitos.

Minha resposta à quarta pergunta é também sim, com base na resposta à quinta. Acredito que seremos capazes de fazer máquinas que possam distinguir o bem do mal, ao menos tão confiavelmente quanto os humanos e, esperançosamente, mais. O primeiro requisito de uma máquina moral é a capacidade de refletir sobre suas próprias ações, o que se enquadra na análise contrafactual. Uma vez que programemos a autoconsciência, por mais limitada que seja, a empatia e a equidade se seguem, pois se baseiam nos mesmos princípios computacionais, com outro agente acrescentado à equação.

Há uma grande diferença de espírito entre a abordagem causal de construir o robô moral e uma abordagem que tem sido estudada e remoída repetidas vezes na ficção científica desde os anos 1950: as leis da robótica de Asimov. Isaac Asimov propôs três leis absolutas, começando com “Um robô não pode ferir um ser humano ou, por omissão, permitir que um ser humano sofra algum mal”. Mas, como a ficção científica mostrou repetidas vezes, as leis de Asimov sempre conduzem a contradições. Para cientistas de IA, isso não chega a ser surpresa: sistemas baseados em regras nunca dão certo. Mas disso não se segue que construir um robô moral seja impossível. Significa que a abordagem não pode ser prescritiva e baseada em regras. Significa que deveríamos equipar máquinas pensantes com as mesmas capacidades cognitivas que temos, que incluem empatia, predição de longo prazo e autorrestrição, e então permitir-lhes tomar suas próprias decisões.

Uma vez que tenhamos construído um robô moral, muitas visões apocalípticas começam a recuar para a irrelevância. Não há razão para nos abstermos de construir máquinas que sejam mais capazes de distinguir o bem do mal do que nós, mais capazes de resistir à tentação, mais capazes de atribuir culpa e crédito. Nesse ponto, como jogadores de xadrez e de go, podemos mesmo começar a aprender com nossa própria criação. Seremos capazes de depender de nossas máquinas para um senso de justiça de olhar lúcido e causalmente sólido. Seremos capazes de aprender como nosso próprio \textit{software} de livre-arbítrio funciona e como consegue esconder seus segredos de nós. Tal máquina pensante seria uma companheira maravilhosa para nossa espécie e verdadeiramente se qualificaria como o primeiro e melhor presente da IA à humanidade.

\chapter{Agradecimentos}
Enumerar o elenco completo de alunos, amigos, colegas e professores que contribuíram com ideias para este livro equivaleria a escrever outro livro. Ainda assim, alguns protagonistas merecem menção especial, do meu ponto de vista pessoal.

Gostaria de agradecer a Phil Dawid, por me dar a primeira audição nas páginas da \textit{Biometrika}; a Jamie Robins e Sander Greenland, por transformar a epidemiologia em uma comunidade que fala a linguagem dos grafos; ao falecido Dennis Lindley, por me assegurar de que mesmo estatísticos experientes conseguem reconhecer falhas em seu campo e se unir em prol de sua reforma; a Chris Winship, Steven Morgan e Felix Elwert, por introduzir as ciências sociais na era da causalidade; e, finalmente, a Peter Spirtes, Clark Glymour e Richard Scheines, por sua ajuda ao me empurrar do penhasco das probabilidades para as águas tempestuosas da causalidade.

Escavando mais fundo em minha história antiga, devo agradecer a Joseph Hermony, ao Dr. Shimshon Lange, ao professor Franz Ollendorff e a outros dedicados professores de ciências que me inspiraram desde o ensino fundamental até a faculdade. Eles inculcaram em muitos de nós, israelenses de primeira geração, um senso de missão e de responsabilidade histórica de perseguir explorações científicas como o desafio mais nobre — e divertido — da humanidade.

Este livro teria permanecido uma relíquia de pensamento wishful se não fosse pelo meu coautor, Dana Mackenzie, que levou a sério meu pensamento wishful e o tornou realidade. Ele não apenas corrigiu meu sotaque estrangeiro, mas também me levou a terras distantes — dos navios da Marinha do capitão James Lind à expedição antártica do capitão Robert Scott —, acrescentando conhecimento, histórias, estrutura e clareza a um emaranhado de equações matemáticas que aguardavam uma narrativa organizadora.

Devo uma grande dívida aos membros do Cognitive Systems Laboratory da UCLA, cujo trabalho e ideias ao longo das últimas três décadas e meia formaram a base científica deste livro: Alex Balke, Elias Bareinboim, Blai Bonet, Carlo Brito, Avin Chen, Bryant Chen, David Chickering, Adnan Darwiche, Rina Dechter, Andrew Forney, David Galles, Hector Geffner, Dan Geiger, Moises Goldszmidt, David Heckerman, Mark Hopkins, Jin Kim, Manabu Kuroki, Trent Kyono, Karthika Mohan, Azaria Paz, George Rebane, Ilya Shpitser, Jin Tian e Thomas Verma.

Agências de fomento recebem agradecimentos ritualizados em publicações acadêmicas, mas crédito real muito inferior ao que merecem, considerando seu papel crucial em reconhecer sementes de ideias antes que se tornem moda. Devo reconhecer o apoio constante e infalível da National Science Foundation e do Office of Naval Research, por meio do programa de Machine Learning and Intelligence chefiado por Behzad Kamgar-Parsi.

Dana e eu gostaríamos de agradecer ao nosso agente, John Brockman, que nos deu encorajamento oportuno e o benefício de sua expertise profissional. Nosso editor na Basic Books, TJ Kelleher, nos fez exatamente as perguntas certas e convenceu a Basic Books de que uma história tão ambiciosa não poderia ser contada em 200 páginas. Nossos ilustradores, Maayan Harel e Dakota Harr, conseguiram lidar com nossas instruções às vezes conflitantes e deram vida a assuntos abstratos com humor e beleza. Kaoru Mulvihill, na UCLA, merece muito crédito por revisar várias versões do manuscrito e ilustrar as hordas de grafos e diagramas.

Dana será eternamente grato a John Wilkes, que fundou o Science Communication Program na UC Santa Cruz, que continua forte e é a melhor rota possível para uma carreira como escritor de ciência. Dana também gostaria de agradecer à sua esposa, Kay, que o encorajou a perseguir seu sonho de infância de ser escritor, mesmo quando isso significava arrancar as raízes, atravessar o país e recomeçar.

Finalmente, minha dívida mais profunda é para com minha família, por sua paciência, compreensão e apoio. Especialmente para minha esposa, Ruth, minha bússola moral, por seu amor e sabedoria sem fim. Para meu falecido filho, Danny, por me mostrar a audácia silenciosa da verdade. Para minhas filhas Tamara e Michelle, por confiarem em minha promessa perene de que o livro eventualmente ficaria pronto. E para meus netos, Leora, Tori, Adam, Ari e Evan, por darem um propósito às minhas longas jornadas e por sempre dissolverem minhas perguntas de “porquê”.

\chapter{Sobre os autores}
\textbf{Judea Pearl} é professor de ciência da computação na University of California, Los Angeles, vencedor do Prêmio Turing de 2011 e autor de três livros técnicos clássicos sobre causalidade. Vive em Los Angeles, Califórnia.

\textbf{Dana Mackenzie} é escritor de ciência premiado e autor de \textit{The Big Splat, or How Our Moon Came to Be}. Vive em Santa Cruz, Califórnia.

\section{Também de Judea Pearl}
\begin{itemize}
\item \textit{Causal Inference in Statistics: A Primer} (com Madelyn Glymour e Nicholas Jewell)
\item \textit{An Introduction to Causal Inference}
\item \textit{I Am Jewish: Personal Reflections Inspired by the Last Words of Daniel Pearl} (coeditado com Ruth Pearl)
\item \textit{Causality: Models, Reasoning and Inference}
\item \textit{Probabilistic Reasoning in Intelligent Systems: Networks of Plausible Inference}
\item \textit{Heuristics: Intelligent Search Strategies for Computer Problem Solving}
\end{itemize}
\section{Também de Dana Mackenzie}
\begin{itemize}
\item \textit{The Universe in Zero Words: The Story of Mathematics as Told Through Equations}
\item \textit{What’s Happening in the Mathematical Sciences} (volumes 6–10)
\item \textit{The Big Splat, or How Our Moon Came to Be}
\end{itemize}
\chapter{Notas}
\section{Notas à Introdução}
\textbf{Students are never allowed} [\textit{Os estudantes nunca têm permissão}]: Com possivelmente uma exceção: se tivermos realizado um ensaio clínico randomizado, como discutido no Capítulo 4.

\section{Notas ao Capítulo Um}
\textbf{then the opposite is true} [\textit{então o oposto é verdadeiro}]: Em outras palavras, ao avaliar uma intervenção em um modelo causal, fazemos as mudanças mínimas possíveis para impor seu efeito imediato. Assim, “quebramos” o modelo no que diz respeito a \textit{A}, mas não a \textit{B}.

\textbf{We should thank the language} [\textit{Devemos agradecer à linguagem}]: Devo também mencionar aqui que os contrafactuais nos permitem falar de causalidade em casos individuais: O que teria acontecido ao Sr. Smith, que não foi vacinado e morreu de varíola, se ele tivesse sido vacinado? Tais perguntas, a espinha dorsal da medicina personalizada, não podem ser respondidas a partir de informação do degrau dois.

\textbf{Yet we can answer} [\textit{Ainda assim, podemos responder}]: Para ser mais preciso, em geometria, termos indefinidos como “ponto” e “reta” são primitivos. O primitivo na inferência causal é a relação de “escutar”, indicada por uma seta.

\section{Notas ao Capítulo Dois}
\textbf{And now the algebraic magic} [\textit{E agora a mágica algébrica}]: Para quem se der ao trabalho de ler o artigo de Wright, deixo o aviso de que ele não calcula seus coeficientes de caminho em gramas por dia. Ele os calcula em “unidades padrão” e só no final converte para gramas por dia.

\section{Notas ao Capítulo Cinco}
\textbf{“Cigarette smoking is causally related”} [“O tabagismo está causalmente relacionado”]: As evidências para as mulheres eram menos claras naquela época, principalmente porque as mulheres haviam fumado muito menos do que os homens nas primeiras décadas do século.

\section{Notas ao Capítulo Oito}
\textbf{And Abraham drew near} [\textit{E Abraão se aproximou}]: Como antes, usei a tradução da King James, mas fiz pequenas alterações para alinhá-la mais de perto ao hebraico.

\textbf{The ease and familiarity of such} [\textit{A facilidade e a familiaridade de tais}]: As Joint Statistical Meetings de 2013 dedicaram uma sessão inteira ao tópico “Causal Inference as a Missing Data Problem” [Inferência causal como um problema de dados ausentes] — o mantra tradicional de Rubin. Um artigo provocativo daquela sessão intitulava-se “What Is Not a Missing Data Problem?” [O que \textit{não} é um problema de dados ausentes?]. Esse título resume com precisão meus pensamentos.

\textbf{This difference in commitment} [\textit{Essa diferença de compromisso}]: Leitores que estão vendo essa distinção pela primeira vez não devem se sentir sozinhos; há bem mais de 100.000 analistas de regressão nos Estados Unidos confusos com exatamente essa questão, junto com a maioria dos autores de livros-texto de estatística. As coisas só mudarão quando os leitores deste livro cobrarem esses autores.

\textbf{Unfortunately, Rubin does not consider} [\textit{Infelizmente, Rubin não considera}]: “O trabalho de Pearl é claramente interessante, e muitos pesquisadores acham apelativos seus argumentos de que diagramas de caminhos são uma forma natural e conveniente de expressar pressupostos sobre estruturas causais. Em nosso próprio trabalho, talvez influenciados pelo tipo de exemplos que surgem nas ciências sociais e médicas, não encontramos que essa abordagem ajude a derivar inferências causais” (Imbens e Rubin 2013, p. 25).

\textbf{One obstacle I faced was cyclic models} [\textit{Um obstáculo que enfrentei foram os modelos cíclicos}]: São modelos com setas que formam um laço. Evitei discuti-los neste livro, mas tais modelos são bastante importantes em economia, por exemplo.

\textbf{Even today modern-day economists} [\textit{Mesmo hoje, economistas contemporâneos}]: Entre 1995 e 1998, apresentei o seguinte quebra-cabeça brinquedo a centenas de estudantes e professores de econometria por todos os Estados Unidos: Considere as equações clássicas de oferta e demanda que todo estudante de economia resolve em Economia 101.

1. Qual é o valor esperado da demanda \textit{Q} se o preço for \textit{reportado} como \textit{P} = \textit{p}₀? 2. Qual é o valor esperado da demanda \textit{Q} se o preço for \textit{fixado} em \textit{P} = \textit{p}₀? 3. Dado que o preço atual é \textit{P} = \textit{p}₀, qual seria o valor esperado da demanda \textit{Q} se \textit{fixássemos} o preço em \textit{P} = \textit{p}₁?

O leitor deve reconhecer essas consultas como vindas dos três níveis da Escada da Causalidade: predições, ações e contrafactuais. Como eu esperava, os respondentes não tiveram dificuldade em responder à pergunta 1; uma pessoa (um professor distinto) conseguiu resolver a pergunta 2; e ninguém conseguiu responder à pergunta 3.

\textbf{The Model Penal Code expresses} [\textit{O Model Penal Code expressa}]: Trata-se de um conjunto de princípios jurídicos padrão proposto pelo American Law Institute em 1962 para trazer uniformidade aos diversos códigos legais estaduais. Não tem força legal plena em nenhum estado, mas, segundo a Wikipedia, em 2016 mais de dois terços dos estados haviam promulgado partes do Model Penal Code.

\section{Notas ao Capítulo Nove}
\textbf{Those sailors who had eaten} [\textit{Aqueles marinheiros que haviam comido}]: A razão é que os fígados de urso polar \textit{contêm} vitamina C.

\textbf{“On the Inadequacy of the Partial”}: O título refere-se à correlação parcial, um método padrão de controle de um confundidor que discutimos no Capítulo 7.

\textbf{here is how to define the NIE} [\textit{eis como definir o NIE}]: Na sala de parto original, o NIE (efeito indireto natural) era expresso com subscritos aninhados, como em \textit{Y}(0,\textit{M}₁). Espero que o leitor ache a mistura de subscritos contrafactuais e operadores \textit{do} acima mais transparente.

\textbf{In that year researchers identified} [\textit{Naquele ano, pesquisadores identificaram}]: Para ser tecnicamente correto, deveria ser chamado de “polimorfismo de nucleotídeo único”, ou SNP (\textit{single nucleotide polymorphism}). É uma única letra no código genético, enquanto um gene é mais como uma palavra ou uma frase. Contudo, para não sobrecarregar o leitor com terminologia pouco familiar, simplesmente o denominarei de gene.

\chapter{Bibliografia}
\begin{quote}\textbf{Nota sobre os títulos.} Os títulos de livros, artigos, capítulos e relatórios técnicos foram mantidos no idioma original (em geral o inglês), conforme a prática acadêmica internacional. A prosa descritiva das bibliografias anotadas foi traduzida para o português. As entradas de referência seguem a ordem alfabética do original por autor e ano.\end{quote}
\section{Introdução: A mente sobre os dados}
\subsection{Bibliografia anotada}
A história da probabilidade e da estatística da Antiguidade aos dias modernos é coberta em profundidade por Hacking (1990); Stigler (1986, 1999, 2016). Um relato menos técnico é dado em Salsburg (2002). Relatos abrangentes da história do pensamento causal, infelizmente, são escassos, embora material interessante possa ser encontrado em Hoover (2008); Kleinberg (2015); Losee (2012); Mumford e Anjum (2014). A proibição da linguagem causal pode ser vista em quase todo texto estatístico padrão — por exemplo, Freedman, Pisani e Purves (2007) ou Efron e Hastie (2016). Para uma análise dessa proibição como impedimento linguístico, ver Pearl (2009, Capítulos 5 e 11); e como barreira cultural, ver Pearl (2000b). Relatos recentes das conquistas e limitações do \textit{big data} e do aprendizado de máquina estão em Darwiche (2017); Pearl (2017); Mayer-Schönberger e Cukier (2013); Domingos (2015); Marcus (30 de julho de 2017). Toulmin (1961) oferece contexto histórico a esse debate.

Leitores interessados em “descoberta de modelos” e em tratamentos mais técnicos do operador \textit{do} podem consultar Pearl (1994, 2000a, Capítulos 2–3); Spirtes, Glymour e Scheines (2000). Para uma introdução mais acessível, ver Pearl, Glymour e Jewell (2016). Esta última fonte é recomendada a leitores com habilidades matemáticas de nível universitário, mas sem formação em estatística ou ciência da computação. Oferece também introdução básica a probabilidades condicionais, regra de Bayes, regressão e grafos. Versões anteriores do motor de inferência mostrado na Figura 1.1 podem ser encontradas em Pearl (2012); Pearl e Bareinboim (2014).

\begin{figure}[ht]
\centering
\includegraphics[width=0.92\textwidth]{../figuras/fig_1_1.png}
\caption{Figura 1.1. Causas percebidas de uma caçada bem-sucedida a mamutes.}
\end{figure}
\textit{Figura 1.1. Causas percebidas de uma caçada bem-sucedida a mamutes.}

\subsection{Referências}
Darwiche, A. (2017). Human-level intelligence or animal-like abilities? Tech. rep., Department of Computer Science, University of California, Los Angeles, CA. Submitted to \textit{Communications of the ACM}. Available at: https://arXiv:1707.04327.

Domingos, P. (2015). \textit{The Master Algorithm: How the Quest for the Ultimate Learning Machine Will Remake Our World}. Basic Books, New York, NY.

Efron, B., and Hastie, T. (2016). \textit{Computer Age Statistical Inference}. Cambridge University Press, New York, NY.

Freedman, D., Pisani, R., and Purves, R. (2007). \textit{Statistics}. 4th ed. W. W. Norton \& Company, New York, NY.

Hacking, I. (1990). \textit{The Taming of Chance} (Ideas in Context). Cambridge University Press, Cambridge, UK.

Hoover, K. (2008). Causality in economics and econometrics. In \textit{The New Palgrave Dictionary of Economics} (S. Durlauf and L. Blume, eds.), 2nd ed. Palgrave Macmillan, New York, NY.

Kleinberg, S. (2015). \textit{Why: A Guide to Finding and Using Causes}. O’Reilly Media, Sebastopol, CA.

Losee, J. (2012). \textit{Theories of Causality: From Antiquity to the Present}. Routledge, New York, NY.

Marcus, G. (July 30, 2017). Artificial intelligence is stuck. Here’s how to move it forward. \textit{New York Times}, SR6.

Mayer-Schönberger, V., and Cukier, K. (2013). \textit{Big Data: A Revolution That Will Transform How We Live, Work, and Think}. Houghton Mifflin Harcourt Publishing, New York, NY.

Morgan, S., and Winship, C. (2015). \textit{Counterfactuals and Causal Inference: Methods and Principles for Social Research} (Analytical Methods for Social Research). 2nd ed. Cambridge University Press, New York, NY.

Mumford, S., and Anjum, R. L. (2014). \textit{Causation: A Very Short Introduction} (Very Short Introductions). Oxford University Press, New York, NY.

Pearl, J. (1988). \textit{Probabilistic Reasoning in Intelligent Systems}. Morgan Kaufmann, San Mateo, CA.

Pearl, J. (1994). A probabilistic calculus of actions. In \textit{Uncertainty in Artificial Intelligence 10} (R. L. de Mantaras and D. Poole, eds.). Morgan Kaufmann, San Mateo, CA, 454–462.

Pearl, J. (1995). Causal diagrams for empirical research. \textit{Biometrika} 82: 669–710.

Pearl, J. (2000a). \textit{Causality: Models, Reasoning, and Inference}. Cambridge University Press, New York, NY.

Pearl, J. (2000b). Comment on A. P. Dawid’s Causal inference without counterfactuals. \textit{Journal of the American Statistical Association} 95: 428–431.

Pearl, J. (2009). \textit{Causality: Models, Reasoning, and Inference}. 2nd ed. Cambridge University Press, New York, NY.

Pearl, J. (2012). The causal foundations of structural equation modeling. In \textit{Handbook of Structural Equation Modeling} (R. Hoyle, ed.). Guilford Press, New York, NY, 68–91.

Pearl, J. (2017). Advances in deep neural networks, at ACM Turing 50 Celebration. Available at: https://www.youtube.com/watch?v=mFYM9j8bGtg (June 23, 2017).

Pearl, J., and Bareinboim, E. (2014). External validity: From do-calculus to transportability across populations. \textit{Statistical Science} 29: 579–595.

Pearl, J., Glymour, M., and Jewell, N. (2016). \textit{Causal Inference in Statistics: A Primer}. Wiley, New York, NY.

Provine, W. B. (1986). \textit{Sewall Wright and Evolutionary Biology}. University of Chicago Press, Chicago, IL.

Salsburg, D. (2002). \textit{The Lady Tasting Tea: How Statistics Revolutionized Science in the Twentieth Century}. Henry Holt and Company, LLC, New York, NY.

Spirtes, P., Glymour, C., and Scheines, R. (2000). \textit{Causation, Prediction, and Search}. 2nd ed. MIT Press, Cambridge, MA.

Stigler, S. M. (1986). \textit{The History of Statistics: The Measurement of Uncertainty Before 1900}. Belknap Press of Harvard University Press, Cambridge, MA.

Stigler, S. M. (1999). \textit{Statistics on the Table: The History of Statistical Concepts and Methods}. Harvard University Press, Cambridge, MA.

Stigler, S. M. (2016). \textit{The Seven Pillars of Statistical Wisdom}. Harvard University Press, Cambridge, MA.

Toulmin, S. (1961). \textit{Foresight and Understanding: An Enquiry into the Aims of Science}. University of Indiana Press, Bloomington, IN.

Virgil. (29 BC). \textit{Georgics}. Verse 490, Book 2.

\section{Capítulo 1. A Escada da Causalidade}
\subsection{Bibliografia anotada}
Um relato técnico das distinções entre os três níveis da Escada da Causalidade pode ser encontrado no Capítulo 1 de Pearl (2000). Nossas comparações entre a Escada da Causalidade e o desenvolvimento cognitivo humano foram inspiradas por Harari (2015) e pelos achados recentes de Kind et al. (2014). O artigo de Kind contém detalhes sobre o Homem-Leão e o sítio onde foi encontrado. Pesquisas relacionadas sobre o desenvolvimento da compreensão causal em bebês podem ser encontradas em Weisberg e Gopnik (2013).

O teste de Turing foi proposto pela primeira vez como um jogo de imitação em 1950 (Turing, 1950). O argumento da “Sala Chinesa” de Searle apareceu em Searle (1980) e tem sido amplamente discutido desde então. Ver Russell e Norvig (2003); Preston e Bishop (2002); Pinker (1997).

O uso de modificação de modelo para representar intervenção tem raízes conceituais no economista Trygve Haavelmo (1943); ver Pearl (2015) para um relato detalhado. Spirtes, Glymour e Scheines (1993) deram-lhe uma representação gráfica em termos de deleção de setas. Balke e Pearl (1994a, 1994b) estenderam-no para simular raciocínio contrafactual, como demonstrado no exemplo do pelotão de fuzilamento.

Um resumo abrangente da causalidade probabilística é dado em Hitchcock (2016). Ideias-chave podem ser encontradas em Reichenbach (1956); Suppes (1970); Cartwright (1983); Spohn (2012). Minhas análises da causalidade probabilística e da elevação de probabilidade são apresentadas em Pearl (2000; 2009, Seção 7.5; 2011).

\subsection{Referências}
Balke, A., and Pearl, J. (1994a). Counterfactual probabilities: Computational methods, bounds, and applications. In \textit{Uncertainty in Artificial Intelligence 10} (R. L. de Mantaras and D. Poole, eds.). Morgan Kaufmann, San Mateo, CA, 46–54.

Balke, A., and Pearl, J. (1994b). Probabilistic evaluation of counterfactual queries. In \textit{Proceedings of the Twelfth National Conference on Artificial Intelligence}, vol. 1. MIT Press, Menlo Park, CA, 230–237.

Cartwright, N. (1983). \textit{How the Laws of Physics Lie}. Clarendon Press, Oxford, UK.

Haavelmo, T. (1943). The statistical implications of a system of simultaneous equations. \textit{Econometrica} 11: 1–12. Reprinted in D. F. Hendry and M. S. Morgan (Eds.), \textit{The Foundations of Econometric Analysis}, Cambridge University Press, Cambridge, UK, 477–490, 1995.

Harari, Y. N. (2015). \textit{Sapiens: A Brief History of Humankind}. Harper Collins Publishers, New York, NY.

Hitchcock, C. (2016). Probabilistic causation. In \textit{Stanford Encyclopedia of Philosophy} (Winter 2016) (E. N. Zalta, ed.). Metaphysics Research Lab, Stanford, CA. Available at: https://stanford.library.sydney.edu.au/archives/win2016/entries/causation-probabilistic.

Kind, C.-J., Ebinger-Rist, N., Wolf, S., Beutelspacher, T., and Wehrberger, K. (2014). The smile of the Lion Man. Recent excavations in Stadel cave (Baden-Württemberg, south-western Germany) and the restoration of the famous upper palaeolithic figurine. \textit{Quartär} 61: 129–145.

Pearl, J. (2000). \textit{Causality: Models, Reasoning, and Inference}. Cambridge University Press, New York, NY.

Pearl, J. (2009). \textit{Causality: Models, Reasoning, and Inference}. 2nd ed. Cambridge University Press, New York, NY.

Pearl, J. (2011). The structural theory of causation. In \textit{Causality in the Sciences} (P. M. Illari, F. Russo, and J. Williamson, eds.), chap. 33. Clarendon Press, Oxford, UK, 697–727.

Pearl, J. (2015). Trygve Haavelmo and the emergence of causal calculus. \textit{Econometric Theory} 31: 152–179. Special issue on Haavelmo centennial.

Pinker, S. (1997). \textit{How the Mind Works}. W. W. Norton and Company, New York, NY.

Preston, J., and Bishop, M. (2002). \textit{Views into the Chinese Room: New Essays on Searle and Artificial Intelligence}. Oxford University Press, New York, NY.

Reichenbach, H. (1956). \textit{The Direction of Time}. University of California Press, Berkeley, CA.

Russell, S. J., and Norvig, P. (2003). \textit{Artificial Intelligence: A Modern Approach}. 2nd ed. Prentice Hall, Upper Saddle River, NJ.

Searle, J. (1980). Minds, brains, and programs. \textit{Behavioral and Brain Sciences} 3: 417–457.

Spirtes, P., Glymour, C., and Scheines, R. (1993). \textit{Causation, Prediction, and Search}. Springer-Verlag, New York, NY.

Spohn, W. (2012). \textit{The Laws of Belief: Ranking Theory and Its Philosophical Applications}. Oxford University Press, Oxford, UK.

Suppes, P. (1970). \textit{A Probabilistic Theory of Causality}. North-Holland Publishing Co., Amsterdam, Netherlands.

Turing, A. (1950). Computing machinery and intelligence. \textit{Mind} 59: 433–460.

Weisberg, D. S., and Gopnik, A. (2013). Pretense, counterfactuals, and Bayesian causal models: Why what is not real really matters. \textit{Cognitive Science} 37: 1368–1381.

\section{Capítulo 2. De bucaneiros a cobaias: a gênese da inferência causal}
\subsection{Bibliografia anotada}
As explorações de Galton sobre hereditariedade e correlação são descritas em seus livros (Galton, 1869, 1883, 1889) e também documentadas em Stigler (2012, 2016). Para uma introdução básica ao equilíbrio de Hardy-Weinberg, ver Wikipedia (2016a). Para a origem da citação de Galileu “E pur si muove”, ver Wikipedia (2016b). A história das catacumbas de Paris e o choque de Pearson com correlações induzidas por “misturas artificiais” podem ser encontrados em Stigler (2012, p. 9).

Como Wright viveu uma vida tão longa, teve o raro privilégio de ver uma biografia (Provine, 1986) sair enquanto ainda estava vivo. A biografia de Provine continua sendo o melhor lugar para conhecer a carreira de Wright, e recomendamos particularmente o Capítulo 5 sobre análise de caminhos. Os dois esboços biográficos de Crow (Crow, 1982, 1990) também oferecem uma perspectiva biográfica muito útil. Wright (1920) é o artigo seminal sobre diagramas de caminhos; Wright (1921) é uma exposição mais completa e a fonte do exemplo do peso ao nascer das cobaias. Wright (1983) é a resposta de Wright à crítica de Karlin, escrita quando ele tinha mais de noventa anos.

O destino da análise de caminhos na economia e nas ciências sociais é narrado no Capítulo 5 de Pearl (2000) e em Bollen e Pearl (2013). Blalock (1964), Duncan (1966) e Goldberger (1972) introduziram as ideias de Wright nas ciências sociais com grande entusiasmo, mas seus fundamentos teóricos não foram bem articulados. Uma década depois, quando Freedman (1987) desafiou os analistas de caminhos a explicar como as intervenções são modeladas, o entusiasmo desapareceu, e pesquisadores de ponta recuaram para ver a modelagem de equações estruturais (SEM) como um exercício de análise estatística. Essa reveladora discussão entre doze estudiosos está documentada no mesmo fascículo do \textit{Journal of Educational Statistics} do artigo de Freedman. A relutância dos economistas em abraçar diagramas e notação estrutural é descrita em Pearl (2015). As dolorosas consequências para a educação econômica estão documentadas em Chen e Pearl (2013).

Uma exposição popular do debate bayesiano \textit{versus} frequentista é dada em McGrayne (2011). Discussões mais técnicas podem ser encontradas em Efron (2013) e Lindley (1987).

\subsection{Referências}
Blalock, H., Jr. (1964). \textit{Causal Inferences in Nonexperimental Research}. University of North Carolina Press, Chapel Hill, NC.

Bollen, K., and Pearl, J. (2013). Eight myths about causality and structural equation models. In \textit{Handbook of Causal Analysis for Social Research} (S. Morgan, ed.). Springer, Dordrecht, Netherlands, 301–328.

Chen, B., and Pearl, J. (2013). Regression and causation: A critical examination of econometrics textbooks. \textit{Real-World Economics Review} 65: 2–20.

Crow, J. F. (1982). Sewall Wright, the scientist and the man. \textit{Perspectives in Biology and Medicine} 25: 279–294.

Crow, J. F. (1990). Sewall Wright’s place in twentieth-century biology. \textit{Journal of the History of Biology} 23: 57–89.

Duncan, O. D. (1966). Path analysis. \textit{American Journal of Sociology} 72: 1–16.

Efron, B. (2013). Bayes’ theorem in the 21st century. \textit{Science} 340: 1177–1178.

Freedman, D. (1987). As others see us: A case study in path analysis (with discussion). \textit{Journal of Educational Statistics} 12: 101–223.

Galton, F. (1869). \textit{Hereditary Genius}. Macmillan, London, UK.

Galton, F. (1883). \textit{Inquiries into Human Faculty and Its Development}. Macmillan, London, UK.

Galton, F. (1889). \textit{Natural Inheritance}. Macmillan, London, UK.

Goldberger, A. (1972). Structural equation models in the social sciences. \textit{Econometrica: Journal of the Econometric Society} 40: 979–1001.

Lindley, D. (1987). \textit{Bayesian Statistics: A Review}. CBMS-NSF Regional Conference Series in Applied Mathematics (Book 2). Society for Industrial and Applied Mathematics, Philadelphia, PA.

McGrayne, S. B. (2011). \textit{The Theory That Would Not Die}. Yale University Press, New Haven, CT.

Pearl, J. (2000). \textit{Causality: Models, Reasoning, and Inference}. Cambridge University Press, New York, NY.

Pearl, J. (2015). Trygve Haavelmo and the emergence of causal calculus. \textit{Econometric Theory} 31: 152–179. Special issue on Haavelmo centennial.

Provine, W. B. (1986). \textit{Sewall Wright and Evolutionary Biology}. University of Chicago Press, Chicago, IL.

Stigler, S. M. (2012). Studies in the history of probability and statistics, L: Karl Pearson and the rule of three. \textit{Biometrika} 99: 1–14.

Stigler, S. M. (2016). \textit{The Seven Pillars of Statistical Wisdom}. Harvard University Press, Cambridge, MA.

Wikipedia. (2016a). Hardy-Weinberg principle. Available at: https://en.wikipedia.org/wiki/Hardy-Weinberg-principle (last edited: October 2, 2016).

Wikipedia. (2016b). Galileo Galilei. Available at: https://en.wikipedia.org/wiki/Galileo\_Galilei (last edited: October 6, 2017).

Wright, S. (1920). The relative importance of heredity and environment in determining the piebald pattern of guinea-pigs. \textit{Proceedings of the National Academy of Sciences of the United States of America} 6: 320–332.

Wright, S. (1921). Correlation and causation. \textit{Journal of Agricultural Research} 20: 557–585.

Wright, S. (1983). On “Path analysis in genetic epidemiology: A critique.” \textit{American Journal of Human Genetics} 35: 757–768.

\section{Capítulo 3. Das evidências às causas: o reverendo Bayes encontra o Sr. Holmes}
\subsection{Bibliografia anotada}
Introduções elementares à regra de Bayes e ao pensamento bayesiano podem ser encontradas em Lindley (2014) e Pearl, Glymour e Jewell (2016). Debates com representações rivais da incerteza são apresentados em Pearl (1988); ver também a extensa lista de referências ali fornecida. Nossos dados de mamografia baseiam-se principalmente em informações do Breast Cancer Surveillance Consortium (BCSC, 2009) e da US Preventive Services Task Force (USPSTF, 2016) e são apresentados apenas para fins didáticos.

As “redes bayesianas” receberam esse nome em 1985 (Pearl, 1985) e foram apresentadas pela primeira vez como um modelo de memória autoativada. Aplicações a sistemas especialistas seguiram o desenvolvimento de algoritmos de atualização de crenças para redes com laços (Pearl, 1986; Lauritzen e Spiegelhalter, 1988). O conceito de \textit{d}-separação, que conecta o bloqueio de caminhos em um diagrama a dependências nos dados, tem raízes na teoria dos \textit{graphoids} (Pearl e Paz, 1985). A teoria desvela as propriedades comuns de grafos (daí o nome) e de probabilidades e explica por que esses dois objetos matemáticos aparentemente alheios podem se apoiar mutuamente de tantas maneiras. Ver também “Graphoid”, Wikipedia.

O divertido exemplo da mala no voo aéreo pode ser encontrado em Conrady e Jouffe (2015, Capítulo 4). O desastre do voo 17 da Malaysia Airlines foi bem coberto pela mídia; ver Clark e Kramer (14 de outubro de 2015) para uma atualização da investigação um ano após o incidente. Wiegerinck, Burgers e Kappen (2013) descrevem como o Bonaparte funciona. Detalhes adicionais sobre a identificação das vítimas do voo 17, incluindo a genealogia mostrada na Figura 3.7, vieram de correspondência pessoal de W. Burgers a D. Mackenzie (24 de agosto de 2016) e de entrevista telefônica com W. Burgers e B. Kappen por D. Mackenzie (23 de agosto de 2016).

\begin{figure}[ht]
\centering
\includegraphics[width=0.92\textwidth]{../figuras/fig_3_7.png}
\caption{Figura 3.7. ) numa rede bayesiana. Na Figura 3.8, vemos como o Bonaparte converte um pequeno pedaço de um pedigree numa rede bayesiana (causal). O problema central é que o genótipo de um indivíduo, detectado num teste de DNA, contém uma contribuição tanto do pai quanto da mãe, mas não podemos dizer qual parte é qual. Assim, essas duas contribuições (chamadas “alelos”) têm de ser tratadas como variáveis ocultas, não mensuráveis, na rede bayesiana. Parte do trabalho do Bonaparte é inferir a probabilidade da causa (o gene da vítima para olhos azuis veio do pai) a partir da evidência (p.ex., ele tem um gene de olhos azuis e um gene de olhos pretos; seus primos do lado do pai têm olhos azuis, mas seus primos do lado da mãe têm olhos pretos). Este é um problema de probabilidade inversa — exatamente para o que a regra de Bayes foi inventada.}
\end{figure}
\textit{Figura 3.7. ) numa rede bayesiana. Na Figura 3.8, vemos como o Bonaparte converte um pequeno pedaço de um pedigree numa rede bayesiana (causal). O problema central é que o genótipo de um indivíduo, detectado num teste de DNA, contém uma contribuição tanto do pai quanto da mãe, mas não podemos dizer qual parte é qual. Assim, essas duas contribuições (chamadas “alelos”) têm de ser tratadas como variáveis ocultas, não mensuráveis, na rede bayesiana. Parte do trabalho do Bonaparte é inferir a probabilidade da causa (o gene da vítima para olhos azuis veio do pai) a partir da evidência (p.ex., ele tem um gene de olhos azuis e um gene de olhos pretos; seus primos do lado do pai têm olhos azuis, mas seus primos do lado da mãe têm olhos pretos). Este é um problema de probabilidade inversa — exatamente para o que a regra de Bayes foi inventada.}

A complexa e fascinante história dos códigos turbo e de verificação de paridade de baixa densidade ainda não foi contada de forma verdadeiramente acessível ao leigo, mas bons pontos de partida são Costello e Forney (2007) e Hardesty (2010a, 2010b). A realização crucial de que os códigos turbo funcionam pelo algoritmo de propagação de crenças vem de McEliece, David e Cheng (1998). Códigos eficientes continuam a ser um campo de batalha para as comunicações sem fio; Carlton (2016) examina os contenders atuais para telefones “5G” (previstos para a década de 2020).

\subsection{Referências}
Breast Cancer Surveillance Consortium (BCSC). (2009). Performance measures for 1,838,372 screening mammography examinations from 2004 to 2008 by age. Available at: http://www.bcsc-research.org/statistics/performance/screening/2009/perf\_age.html (accessed October 12, 2016).

Carlton, A. (2016). Surprise! Polar codes are coming in from the cold. \textit{Computerworld}. Available at: https://www.computerworld.com/article/3151866/mobile-wireless/surprisepolar-codes-are-coming-in-from-the-cold.html (posted December 22, 2016).

Clark, N., and Kramer, A. (October 14, 2015). Malaysia Airlines Flight 17 most likely hit by Russian-made missile, inquiry says. \textit{New York Times}.

Conrady, S., and Jouffe, L. (2015). \textit{Bayesian Networks and Bayesia Lab: A Practical Introduction for Researchers}. Bayesia USA, Franklin, TN.

Costello, D. J., and Forney, G. D., Jr. (2007). Channel coding: The road to channel capacity. \textit{Proceedings of IEEE} 95: 1150–1177.

Hardesty, L. (2010a). Explained: Gallager codes. \textit{MIT News}. Available at: http://news.mit.edu/2010/gallager-codes-0121 (posted: January 21, 2010).

Hardesty, L. (2010b). Explained: The Shannon limit. \textit{MIT News}. Available at: http://news.mit.edu/2010/explained-shannon-0115 (posted January 19, 2010).

Lauritzen, S., and Spiegelhalter, D. (1988). Local computations with probabilities on graphical structures and their application to expert systems (with discussion). \textit{Journal of the Royal Statistical Society, Series B} 50: 157–224.

Lindley, D. V. (2014). \textit{Understanding Uncertainty}. Rev. ed. John Wiley and Sons, Hoboken, NJ.

McEliece, R. J., David, J. M., and Cheng, J. (1998). Turbo decoding as an instance of Pearl’s “belief propagation” algorithm. \textit{IEEE Journal on Selected Areas in Communications} 16: 140–152.

Pearl, J. (1985). Bayesian networks: A model of self-activated memory for evidential reasoning. In \textit{Proceedings, Cognitive Science Society} (CSS-7). UCLA Computer Science Department, Irvine, CA.

Pearl, J. (1986). Fusion, propagation, and structuring in belief networks. \textit{Artificial Intelligence} 29: 241–288.

Pearl, J. (1988). \textit{Probabilistic Reasoning in Intelligent Systems}. Morgan Kaufmann, San Mateo, CA.

Pearl, J., Glymour, M., and Jewell, N. (2016). \textit{Causal Inference in Statistics: A Primer}. Wiley, New York, NY.

Pearl, J., and Paz, A. (1985). GRAPHOIDS: A graph-based logic for reasoning about relevance relations. Tech. Rep. 850038 (R-53-L). Computer Science Department, University of California, Los Angeles. Short version in B. DuBoulay, D. Hogg, and L. Steels (Eds.) \textit{Advances in Artificial Intelligence—II}, Amsterdam, North Holland, 357–363, 1987.

US Preventive Services Task Force (USPSTF) (2016). Final recommendation statement: Breast cancer: Screening. Available at: https://www.uspreventiveservicestaskforce.org/Page/Document/Recommend cancer-screening1 (updated: January 2016).

Wikipedia. (2018). Graphoid. Available at: https://en.wikipedia.org/wiki/Graphoid (last edited: January 8, 2018).

Wiegerinck, W., Burgers, W., and Kappen, B. (2013). Bayesian networks, introduction and practical applications. In \textit{Handbook on Neural Information Processing} (M. Bianchini, M. Maggini, and L. C. Jain, eds.). Intelligent Systems Reference Library (Book 49). Springer, Berlin, Germany, 401–431.

\section{Capítulo 4. Confundimento e desconfundimento: ou, abatendo a variável à espreita}
\subsection{Bibliografia anotada}
A história de Daniel tem sido frequentemente citada como o primeiro ensaio controlado; ver, por exemplo, Lilienfeld (1982) ou Stigler (2016). Os resultados do estudo de caminhada de Honolulu foram relatados em Hakim (1998). A longa citação de Fisher Box sobre “o interrogatório habilidoso da Natureza” vem de sua excelente biografia do pai (Box, 1978, Capítulo 6). Fisher, também, escreveu sobre experimentos como um diálogo com a Natureza; ver Stigler (2016). Assim, creio que podemos pensar em sua citação como quase vinda do próprio patriarca, apenas mais belamente expressa.

É fascinante ler os artigos de Weinberg sobre confundimento (Weinberg, 1993; Howards et al., 2012) um após o outro. São como duas fotografias da história do confundimento: uma tirada pouco antes de os diagramas causais se tornarem difundidos e a segunda, vinte anos depois, revisitando os mesmos exemplos com diagramas causais. O complicado diagrama de Forbes da rede causal de asma e tabagismo pode ser encontrado em Williamson et al. (2014). A “definição epidemiológica clássica de confundimento” de Morabia pode ser encontrada em Morabia (2011). As citações de David Cox vêm de Cox (1992, pp. 66–67). Outras boas fontes sobre a história do confundimento são Greenland e Robins (2009) e Wikipedia (2016).

O critério da porta dos fundos (\textit{back-door criterion}) para eliminar o viés de confundimento, junto com sua fórmula de ajustamento, foi introduzido em Pearl (1993). Seu impacto na epidemiologia pode ser visto em Greenland, Pearl e Robins (1999). Extensões a intervenções sequenciais e outras nuances são desenvolvidas em Pearl (2000, 2009) e descritas de forma mais acessível em Pearl, Glymour e Jewell (2016). \textit{Software} para computar efeitos causais usando o cálculo \textit{do} está disponível em Tikka e Karvanen (2017). O artigo de Greenland e Robins (1986) foi revisitado pelos autores um quarto de século depois, à luz dos extensos desenvolvimentos desde então, incluindo o advento dos diagramas causais (Greenland e Robins, 2009).

\subsection{Referências}
Box, J. F. (1978). \textit{R. A. Fisher: The Life of a Scientist}. John Wiley and Sons, New York, NY.

Cox, D. (1992). \textit{Planning of Experiments}. Wiley-Interscience, New York, NY.

Greenland, S., Pearl, J., and Robins, J. (1999). Causal diagrams for epidemiologic research. \textit{Epidemiology} 10: 37–48.

Greenland, S., and Robins, J. (1986). Identifiability, exchangeability, and epidemiological confounding. \textit{International Journal of Epidemiology} 15: 413–419.

Greenland, S., and Robins, J. (2009). Identifiability, exchangeability, and confounding revisited. \textit{Epidemiologic Perspectives \& Innovations} 6. doi:10.1186/1742-5573-6-4.

Hakim, A. (1998). Effects of walking on mortality among nonsmoking retired men. \textit{New England Journal of Medicine} 338: 94–99.

Hernberg, S. (1996). Significance testing of potential confounders and other properties of study groups—Misuse of statistics. \textit{Scandinavian Journal of Work, Environment and Health} 22: 315–316.

Howards, P. P., Schisterman, E. F., Poole, C., Kaufman, J. S., and Weinberg, C. R. (2012). “Toward a clearer definition of confounding” revisited with directed acyclic graphs. \textit{American Journal of Epidemiology} 176: 506–511.

Lilienfeld, A. (1982). Ceteris paribus: The evolution of the clinical trial. \textit{Bulletin of the History of Medicine} 56: 1–18.

Morabia, A. (2011). History of the modern epidemiological concept of confounding. \textit{Journal of Epidemiology and Community Health} 65: 297–300.

Pearl, J. (1993). Comment: Graphical models, causality, and intervention. \textit{Statistical Science} 8: 266–269.

Pearl, J. (2000). \textit{Causality: Models, Reasoning, and Inference}. Cambridge University Press, New York, NY.

Pearl, J. (2009). \textit{Causality: Models, Reasoning, and Inference}. 2nd ed. Cambridge University Press, New York, NY.

Pearl, J., Glymour, M., and Jewell, N. (2016). \textit{Causal Inference in Statistics: A Primer}. Wiley, New York, NY.

Stigler, S. M. (2016). \textit{The Seven Pillars of Statistical Wisdom}. Harvard University Press, Cambridge, MA.

Tikka, J., and Karvanen, J. (2017). Identifying causal effects with the R Package causaleffect. \textit{Journal of Statistical Software} 76, no. 12. doi:10.18637/jss.r076.i12.

Weinberg, C. (1993). Toward a clearer definition of confounding. \textit{American Journal of Epidemiology} 137: 1–8.

Wikipedia. (2016). Confounding. Available at: https://en.wikipedia.org/wiki/Confounding (accessed: September 16, 2016).

Williamson, E., Aitken, Z., Lawrie, J., Dharmage, S., Burgess, H., and Forbes, A. (2014). Introduction to causal diagrams for confounder selection. \textit{Respirology} 19: 303–311.

\section{Capítulo 5. O debate enfumaçado: limpando o ar}
\subsection{Bibliografia anotada}
Dois estudos de fôlego, Brandt (2007) e Proctor (2012a), contêm todas as informações que qualquer leitor possa pedir sobre o debate tabagismo–câncer de pulmão, a não ser ler os documentos reais das companhias de tabaco (disponíveis \textit{online}). Levantamentos mais curtos do debate tabagismo-câncer na década de 1950 são Salsburg (2002, Capítulo 18), Parascandola (2004) e Proctor (2012b). Stolley (1991) examina o papel singular de R. A. Fisher, e Greenhouse (2009) comenta a importância de Jerome Cornfield. O tiro ouvido ao redor do mundo foi Doll e Hill (1950), que primeiro implicou o tabagismo no câncer de pulmão; embora técnico, é um clássico científico. Para a história do comitê do \textit{surgeon general} e o surgimento das diretrizes de Hill para a causalidade, ver Blackburn e Labarthe (2012) e Morabia (2013). A própria descrição de Hill de seus critérios pode ser encontrada em Hill (1965). Lilienfeld (2007) é a fonte da história de “Abe e Yak” com a qual começamos o capítulo.

VanderWeele (2014) e Hernández-Díaz, Schisterman e Hernán (2006) resolvem o paradoxo do peso ao nascer usando diagramas causais. Um interessante par de artigos “antes e depois” é Wilcox (2001, 2006), escritos antes e depois de o autor conhecer os diagramas causais; sua empolgação no segundo artigo é palpável. Leitores interessados nas estatísticas mais recentes e nas tendências históricas de mortalidade por câncer e tabagismo podem consultar US Department of Health and Human Services (USDHHS, 2014), American Cancer Society (2017) e Wingo (2003).

\subsection{Referências}
American Cancer Society. (2017). Cancer facts and figures. Available at: https://www.cancer.org/research/cancer-facts-statistics.html (posted: February 19, 2015).

Blackburn, H., and Labarthe, D. (2012). Stories from the evolution of guidelines for causal inference in epidemiologic associations: 1953–1965. \textit{American Journal of Epidemiology} 176: 1071–1077.

Brandt, A. (2007). \textit{The Cigarette Century}. Basic Books, New York, NY.

Doll, R., and Hill, A. B. (1950). Smoking and carcinoma of the lung. \textit{British Medical Journal} 2: 739–748.

Greenhouse, J. (2009). Commentary: Cornfield, epidemiology, and causality. \textit{International Journal of Epidemiology} 38: 1199–1201.

Hernández-Díaz, S., Schisterman, E., and Hernán, M. (2006). The birth weight “paradox” uncovered? \textit{American Journal of Epidemiology} 164: 1115–1120.

Hill, A. B. (1965). The environment and disease: Association or causation? \textit{Journal of the Royal Society of Medicine} 58: 295–300.

Lilienfeld, A. (2007). Abe and Yak: The interactions of Abraham M. Lilienfeld and Jacob Yerushalmy in the development of modern epidemiology (1945–1973). \textit{Epidemiology} 18: 507–514.

Morabia, A. (2013). Hume, Mill, Hill, and the sui generis epidemiologic approach to causal inference. \textit{American Journal of Epidemiology} 178: 1526–1532.

Parascandola, M. (2004). Two approaches to etiology: The debate over smoking and lung cancer in the 1950s. \textit{Endeavour} 28: 81–86.

Proctor, R. (2012a). \textit{Golden Holocaust: Origins of the Cigarette Catastrophe and the Case for Abolition}. University of California Press, Berkeley, CA.

Proctor, R. (2012b). The history of the discovery of the cigarette–lung cancer link: Evidentiary traditions, corporate denial, and global toll. \textit{Tobacco Control} 21: 87–91.

Salsburg, D. (2002). \textit{The Lady Tasting Tea: How Statistics Revolutionized Science in the Twentieth Century}. Henry Holt and Company, New York, NY.

Stolley, P. (1991). When genius errs: R. A. Fisher and the lung cancer controversy. \textit{American Journal of Epidemiology} 133: 416–425.

US Department of Health and Human Services (USDHHS). (2014). The health consequences of smoking—50 years of progress: A report of the surgeon general. USDHHS and Centers for Disease Control and Prevention, Atlanta, GA.

VanderWeele, T. (2014). Commentary: Resolutions of the birthweight paradox: Competing explanations and analytical insights. \textit{International Journal of Epidemiology} 43: 1368–1373.

Wilcox, A. (2001). On the importance—and the unimportance—of birthweight. \textit{International Journal of Epidemiology} 30: 1233–1241.

Wilcox, A. (2006). The perils of birth weight—A lesson from directed acyclic graphs. \textit{American Journal of Epidemiology} 164: 1121–1123.

Wingo, P. (2003). Long-term trends in cancer mortality in the United States, 1930–1998. \textit{Cancer} 97: 3133–3275.

\section{Capítulo 6. Paradoxos a rodo!}
\subsection{Bibliografia anotada}
O paradoxo de Monty Hall aparece em muitos livros introdutórios de teoria da probabilidade (p.ex., Grinstead e Snell, 1998, p. 136; Lindley, 2014, p. 201). O equivalente “dilema dos três prisioneiros” foi usado para demonstrar a inadequação de abordagens não bayesianas em Pearl (1988, pp. 58–62). Tierney (21 de julho de 1991) e Crockett (2015) contam a incrível história da coluna de vos Savant sobre o paradoxo de Monty Hall; Crockett reúne vários outros comentários divertidos e embaraçosos que vos Savant recebeu dos ditos “especialistas”. O artigo de Tierney conta o que o próprio Monty Hall pensou da confusão — um ângulo humano interessante!

Um extenso relato da história do paradoxo de Simpson é dado em Pearl (2009, pp. 174–182), incluindo muitas tentativas de estatísticos e filósofos de resolvê-lo sem invocar a causalidade. Um relato mais recente, voltado a educadores, é dado em Pearl (2014). Savage (2009), Julious e Mullee (1994) e Appleton, French e Vanderpump (1996) apresentam os três exemplos reais do paradoxo de Simpson mencionados no texto (relativos a beisebol, cálculos renais e tabagismo, respectivamente). O princípio da coisa certa de Savage (\textit{sure-thing principle}; Savage, 1954) é tratado em Pearl (2016b), e sua versão causal corrigida é derivada em Pearl (2009, pp. 181–182).

Versões do paradoxo de Lord (Lord, 1967) são descritas em Glymour (2006); Hernández-Díaz, Schisterman e Hernán (2006); Senn (2006); Wainer (1991). Uma análise abrangente pode ser encontrada em Pearl (2016a). Paradoxos que invocam contrafactuais não estão incluídos neste capítulo, mas não são menos intrigantes. Para uma amostra, ver Pearl (2013).

\subsection{Referências}
Appleton, D., French, J., and Vanderpump, M. (1996). Ignoring a covariate: An example of Simpson’s paradox. \textit{American Statistician} 50: 340–341.

Crockett, Z. (2015). The time everyone “corrected” the world’s smartest woman. \textit{Priceonomics}. Available at: http://priceonomics.com/the-timeeveryone-corrected-the-worlds-smartest (posted: February 19, 2015).

Glymour, M. M. (2006). Using causal diagrams to understand common problems in social epidemiology. In \textit{Methods in Social Epidemiology}. John Wiley and Sons, San Francisco, CA, 393–428.

Grinstead, C. M., and Snell, J. L. (1998). \textit{Introduction to Probability}. 2nd rev. ed. American Mathematical Society, Providence, RI.

Hernández-Díaz, S., Schisterman, E., and Hernán, M. (2006). The birth weight “paradox” uncovered? \textit{American Journal of Epidemiology} 164: 1115–1120.

Julious, S., and Mullee, M. (1994). Confounding and Simpson’s paradox. \textit{British Medical Journal} 309: 1480–1481.

Lindley, D. V. (2014). \textit{Understanding Uncertainty}. Rev. ed. John Wiley and Sons, Hoboken, NJ.

Lord, F. M. (1967). A paradox in the interpretation of group comparisons. \textit{Psychological Bulletin} 68: 304–305.

Pearl, J. (1988). \textit{Probabilistic Reasoning in Intelligent Systems}. Morgan Kaufmann, San Mateo, CA.

Pearl, J. (2009). \textit{Causality: Models, Reasoning, and Inference}. 2nd ed. Cambridge University Press, New York, NY.

Pearl, J. (2013). The curse of free-will and paradox of inevitable regret. \textit{Journal of Causal Inference} 1: 255–257.

Pearl, J. (2014). Understanding Simpson’s paradox. \textit{American Statistician} 88: 8–13.

Pearl, J. (2016a). Lord’s paradox revisited—(Oh Lord! Kumbaya!). \textit{Journal of Causal Inference} 4. doi:10.1515/jci-2016-0021.

Pearl, J. (2016b). The sure-thing principle. \textit{Journal of Causal Inference} 4: 81–86.

Savage, L. (1954). \textit{The Foundations of Statistics}. John Wiley and Sons, New York, NY.

Savage, S. (2009). \textit{The Flaw of Averages: Why We Underestimate Risk in the Face of Uncertainty}. John Wiley and Sons, Hoboken, NJ.

Senn, S. (2006). Change from baseline and analysis of covariance revisited. \textit{Statistics in Medicine} 25: 4334–4344.

Simon, H. (1954). Spurious correlation: A causal interpretation. \textit{Journal of the American Statistical Association} 49: 467–479.

Tierney, J. (July 21, 1991). Behind Monty Hall’s doors: Puzzle, debate and answer? \textit{New York Times}.

Wainer, H. (1991). Adjusting for differential base rates: Lord’s paradox again. \textit{Psychological Bulletin} 109: 147–151.

\section{Capítulo 7. Além do ajustamento: a conquista do Monte Intervenção}
\subsection{Bibliografia anotada}
Extensões dos ajustamentos pela porta dos fundos e pela porta da frente foram primeiro relatadas em Tian e Pearl (2002), com base na fatoração por \textit{c}-componentes de Tian. Seguiram-se a algoritmização do cálculo \textit{do} por Shpitser (Shpitser e Pearl, 2006a) e, depois, os resultados de completude de Shpitser e Pearl (2006b) e Huang e Valtorta (2006).

Os economistas entre nossos leitores devem notar que a resistência cultural de alguns economistas a ferramentas gráficas de análise (Heckman e Pinto, 2015; Imbens e Rubin, 2015) não é compartilhada por todos os economistas. White e Chalak (2009), por exemplo, generalizaram e aplicaram o cálculo \textit{do} a sistemas econômicos envolvendo equilíbrio e aprendizado. Livros-texto recentes nas ciências sociais e comportamentais, Morgan e Winship (2007) e Kline (2016), sinalizam ainda mais aos jovens pesquisadores que a ortodoxia cultural, como o medo dos telescópios no século XVII, não é duradoura nas ciências.

A investigação de John Snow sobre a cólera foi muito pouco apreciada durante sua vida, e seu obituário de um parágrafo no \textit{Lancet} nem sequer a mencionou. Notavelmente, o principal periódico médico britânico “corrigiu” seu obituário 155 anos depois (Hempel, 2013). Para mais material biográfico sobre Snow, ver Hill (1955) e Cameron e Jones (1983). Glynn e Kashin (2018) é um dos primeiros artigos a demonstrar empiricamente que o ajustamento pela porta da frente é superior ao da porta dos fundos quando há confundidores não observados. A crítica de Freedman ao exemplo tabagismo–alcatrão–câncer de pulmão pode ser encontrada em um capítulo de Freedman (2010) intitulado “On Specifying Graphical Models for Causation”.

Introduções a variáveis instrumentais podem ser encontradas em Greenland (2000) e em muitos livros-texto de econometria (p.ex., Bowden e Turkington, 1984; Wooldridge, 2013). Variáveis instrumentais generalizadas, estendendo a definição clássica dada em nosso texto, foram introduzidas em Brito e Pearl (2002). O programa DAGitty (disponível \textit{online} em http://www.dagitty.net/dags.html) permite aos usuários buscar no diagrama variáveis instrumentais generalizadas e reporta os estimandos resultantes (Textor, Hardt e Knüppel, 2011). Outro pacote de \textit{software} baseado em diagramas para tomada de decisão é o BayesiaLab (www.bayesia.com). Limites em estimativas de variáveis instrumentais são estudados longamente no Capítulo 8 de Pearl (2009) e aplicados ao problema de não adesão (\textit{noncompliance}). A aproximação LATE é defendida e debatida em Imbens (2010).

\subsection{Referências}
Bareinboim, E., and Pearl, J. (2012). Causal inference by surrogate experiments: z-identifiability. In \textit{Proceedings of the Twenty-Eighth Conference on Uncertainty in Artificial Intelligence} (N. de Freitas and K. Murphy, eds.). AUAI Press, Corvallis, OR.

Bowden, R., and Turkington, D. (1984). \textit{Instrumental Variables}. Cambridge University Press, Cambridge, UK.

Brito, C., and Pearl, J. (2002). Generalized instrumental variables. In \textit{Uncertainty in Artificial Intelligence, Proceedings of the Eighteenth Conference} (A. Darwiche and N. Friedman, eds.). Morgan Kaufmann, San Francisco, CA, 85–93.

Cameron, D., and Jones, I. (1983). John Snow, the Broad Street pump, and modern epidemiology. \textit{International Journal of Epidemiology} 12: 393–396.

Cox, D., and Wermuth, N. (2015). Design and interpretation of studies: Relevant concepts from the past and some extensions. \textit{Observational Studies} 1. Available at: https://arxiv.org/pdf/1505.02452.pdf.

Freedman, D. (2010). \textit{Statistical Models and Causal Inference: A Dialogue with the Social Sciences}. Cambridge University Press, New York, NY.

Glynn, A., and Kashin, K. (2018). Front-door versus back-door adjustment with unmeasured confounding: Bias formulas for front-door and hybrid adjustments. \textit{Journal of the American Statistical Association}. To appear.

Greenland, S. (2000). An introduction to instrumental variables for epidemiologists. \textit{International Journal of Epidemiology} 29: 722–729.

Heckman, J. J., and Pinto, R. (2015). Causal analysis after Haavelmo. \textit{Econometric Theory} 31: 115–151.

Hempel, S. (2013). Obituary: John Snow. \textit{Lancet} 381: 1269–1270.

Hill, A. B. (1955). Snow—An appreciation. \textit{Journal of Economic Perspectives} 48: 1008–1012.

Huang, Y., and Valtorta, M. (2006). Pearl’s calculus of intervention is complete. In \textit{Proceedings of the Twenty-Second Conference on Uncertainty in Artificial Intelligence} (R. Dechter and T. Richardson, eds.). AUAI Press, Corvallis, OR, 217–224.

Imbens, G. W. (2010). Better LATE than nothing: Some comments on Deaton (2009) and Heckman and Urzua (2009). \textit{Journal of Economic Literature} 48: 399–423.

Imbens, G. W., and Rubin, D. B. (2015). \textit{Causal Inference for Statistics, Social, and Biomedical Sciences: An Introduction}. Cambridge University Press, Cambridge, MA.

Kline, R. B. (2016). \textit{Principles and Practice of Structural Equation Modeling}. 3rd ed. Guilford, New York, NY.

Morgan, S., and Winship, C. (2007). \textit{Counterfactuals and Causal Inference: Methods and Principles for Social Research} (Analytical Methods for Social Research). Cambridge University Press, New York, NY.

Pearl, J. (2009). \textit{Causality: Models, Reasoning, and Inference}. 2nd ed. Cambridge University Press, New York, NY.

Pearl, J. (2013). Reflections on Heckman and Pinto’s “Causal analysis after Haavelmo.” Tech. Rep. R-420. Department of Computer Science, University of California, Los Angeles, CA. Working paper.

Pearl, J. (2015). Indirect confounding and causal calculus (on three papers by Cox and Wermuth). Tech. Rep. R-457. Department of Computer Science, University of California, Los Angeles, CA.

Shpitser, I., and Pearl, J. (2006a). Identification of conditional interventional distributions. In \textit{Proceedings of the Twenty-Second Conference on Uncertainty in Artificial Intelligence} (R. Dechter and T. Richardson, eds.). AUAI Press, Corvallis, OR, 437–444.

Shpitser, I., and Pearl, J. (2006b). Identification of joint interventional distributions in recursive semi-Markovian causal models. In \textit{Proceedings of the Twenty-First National Conference on Artificial Intelligence}. AAAI Press, Menlo Park, CA, 1219–1226.

Stock, J., and Trebbi, F. (2003). Who invented instrumental variable regression? \textit{Journal of Economic Perspectives} 17: 177–194.

Textor, J., Hardt, J., and Knüppel, S. (2011). DAGitty: A graphical tool for analyzing causal diagrams. \textit{Epidemiology} 22: 745.

Tian, J., and Pearl, J. (2002). A general identification condition for causal effects. In \textit{Proceedings of the Eighteenth National Conference on Artificial Intelligence}. AAAI Press/MIT Press, Menlo Park, CA, 567–573.

Wermuth, N., and Cox, D. (2008). Distortion of effects caused by indirect confounding. \textit{Biometrika} 95: 17–33. (See Pearl [2009, Chapter 4] for a general solution.)

Wermuth, N., and Cox, D. (2014). Graphical Markov models: Overview. ArXiv: 1407.7783.

White, H., and Chalak, K. (2009). Settable systems: An extension of Pearl’s causal model with optimization, equilibrium and learning. \textit{Journal of Machine Learning Research} 10: 1759–1799.

Wooldridge, J. (2013). \textit{Introductory Econometrics: A Modern Approach}. 5th ed. South-Western, Mason, OH.

\section{Capítulo 8. Contrafactuais: minerando mundos que poderiam ter sido}
\subsection{Bibliografia anotada}
A definição de contrafactuais como derivados de equações estruturais foi introduzida por Balke e Pearl (1994a, 1994b) e usada para estimar probabilidades de causação em contextos jurídicos. As relações entre essa estrutura e as desenvolvidas por Rubin e Lewis são discutidas longamente em Pearl (2000, Capítulo 7), onde se mostra que são logicamente equivalentes; um problema resolvido em uma estrutura produziria a mesma solução em outra. Livros recentes em ciências sociais (p.ex., Morgan e Winship, 2015) e em ciências da saúde (p.ex., VanderWeele, 2015) estão adotando a abordagem híbrida de grafos e contrafactuais perseguida em nosso livro.

A seção sobre contrafactuais lineares baseia-se em Pearl (2009, pp. 389–391), que também fornece a solução do problema posto na nota 12. Nossa discussão do ETT (efeito do tratamento nos tratados) baseia-se em Shpitser e Pearl (2009). Questões jurídicas de atribuição, bem como probabilidades de causação, são discutidas longamente em Greenland (1999), que pioneirou a abordagem contrafactual a tais questões. Nosso tratamento de PN, PS e PNS baseia-se em Tian e Pearl (2000) e Pearl (2009, Capítulo 9). Uma abordagem acessível à atribuição contrafactual, incluindo um kit de ferramentas para estimação, é dada em Pearl, Glymour e Jewell (2016). Um tratamento formal avançado da causação real (\textit{actual causation}) pode ser encontrado em Halpern (2016).

Técnicas de pareamento (\textit{matching}) para estimar efeitos causais são usadas rotineiramente por pesquisadores de resultados potenciais (Sekhon, 2007), embora em geral ignorem as armadilhas mostradas em nosso exemplo educação-experiência-salário. Minha percepção de que problemas de dados ausentes devem ser vistos no contexto da modelagem causal formou-se por meio da análise de Mohan e Pearl (2014).

Cowles (2016) e Reid (1998) contam a história dos anos tumultuados de Neyman em Londres, incluindo a anedota sobre Fisher e os modelos de madeira. Greiner (2008) é uma introdução longa e substantiva à causação \textit{but-for} no direito. Allen (2003), Stott et al. (2013), Trenberth (2012) e Hannart et al. (2016) abordam o problema da atribuição de eventos climáticos à mudança climática, e Hannart em particular invoca as ideias de probabilidade necessária e suficiente, que trazem mais clareza ao assunto.

\subsection{Referências}
Allen, M. (2003). Liability for climate change. \textit{Nature} 421: 891–892.

Balke, A., and Pearl, J. (1994a). Counterfactual probabilities: Computational methods, bounds, and applications. In \textit{Uncertainty in Artificial Intelligence 10} (R. L. de Mantaras and D. Poole, eds.). Morgan Kaufmann, San Mateo, CA, 46–54.

Balke, A., and Pearl, J. (1994b). Probabilistic evaluation of counterfactual queries. In \textit{Proceedings of the Twelfth National Conference on Artificial Intelligence}, vol. 1. MIT Press, Menlo Park, CA, 230–237.

Cowles, M. (2016). \textit{Statistics in Psychology: An Historical Perspective}. 2nd ed. Routledge, New York, NY.

Duncan, O. (1975). \textit{Introduction to Structural Equation Models}. Academic Press, New York, NY.

Freedman, D. (1987). As others see us: A case study in path analysis (with discussion). \textit{Journal of Educational Statistics} 12: 101–223.

Greenland, S. (1999). Relation of probability of causation, relative risk, and doubling dose: A methodologic error that has become a social problem. \textit{American Journal of Public Health} 89: 1166–1169.

Greiner, D. J. (2008). Causal inference in civil rights litigation. \textit{Harvard Law Review} 81: 533–598.

Haavelmo, T. (1943). The statistical implications of a system of simultaneous equations. \textit{Econometrica} 11: 1–12. Reprinted in D. F. Hendry and M. S. Morgan (Eds.), \textit{The Foundations of Econometric Analysis}, Cambridge University Press, Cambridge, UK, 477–490, 1995.

Halpern, J. (2016). \textit{Actual Causality}. MIT Press, Cambridge, MA.

Hannart, A., Pearl, J., Otto, F., Naveu, P., and Ghil, M. (2016). Causal counterfactual theory for the attribution of weather and climate-related events. \textit{Bulletin of the American Meteorological Society} (BAMS) 97: 99–110.

Holland, P. (1986). Statistics and causal inference. \textit{Journal of the American Statistical Association} 81: 945–960.

Hume, D. (1739). \textit{A Treatise of Human Nature}. Oxford University Press, Oxford, UK. Reprinted 1888.

Hume, D. (1748). \textit{An Enquiry Concerning Human Understanding}. Reprinted Open Court Press, LaSalle, IL, 1958.

Joffe, M. M., Yang, W. P., and Feldman, H. I. (2010). Selective ignorability assumptions in causal inference. \textit{International Journal of Biostatistics} 6. doi:10.2202/1557-4679.1199.

Lewis, D. (1973a). Causation. \textit{Journal of Philosophy} 70: 556–567. Reprinted with postscript in D. Lewis, \textit{Philosophical Papers}, vol. 2, Oxford University Press, New York, NY, 1986.

Lewis, D. (1973b). \textit{Counterfactuals}. Harvard University Press, Cambridge, MA.

Lewis, M. (2016). \textit{The Undoing Project: A Friendship That Changed Our Minds}. W. W. Norton and Company, New York, NY.

Mohan, K., and Pearl, J. (2014). Graphical models for recovering probabilistic and causal queries from missing data. \textit{Proceedings of Neural Information Processing} 27: 1520–1528.

Morgan, S., and Winship, C. (2015). \textit{Counterfactuals and Causal Inference: Methods and Principles for Social Research} (Analytical Methods for Social Research). 2nd ed. Cambridge University Press, New York, NY.

Neyman, J. (1923). On the application of probability theory to agricultural experiments. Essay on principles. Section 9. \textit{Statistical Science} 5: 465–480.

Pearl, J. (2000). \textit{Causality: Models, Reasoning, and Inference}. Cambridge University Press, New York, NY.

Pearl, J. (2009). \textit{Causality: Models, Reasoning, and Inference}. 2nd ed. Cambridge University Press, New York, NY.

Pearl, J., Glymour, M., and Jewell, N. (2016). \textit{Causal Inference in Statistics: A Primer}. Wiley, New York, NY.

Reid, C. (1998). \textit{Neyman}. Springer-Verlag, New York, NY.

Rubin, D. (1974). Estimating causal effects of treatments in randomized and nonrandomized studies. \textit{Journal of Educational Psychology} 66: 688–701.

Sekhon, J. (2007). The Neyman-Rubin model of causal inference and estimation via matching methods. In \textit{The Oxford Handbook of Political Methodology} (J. M. Box-Steffensmeier, H. E. Brady, and D. Collier, eds.). Oxford University Press, Oxford, UK.

Shpitser, I., and Pearl, J. (2009). Effects of treatment on the treated: Identification and generalization. In \textit{Proceedings of the Twenty-Fifth Conference on Uncertainty in Artificial Intelligence}. AUAI Press, Montreal, Quebec, 514–521.

Stott, P. A., Allen, M., Christidis, N., Dole, R. M., Hoerling, M., Huntingford, C., Pardeep Pall, J. P., and Stone, D. (2013). Attribution of weather and climate-related events. In \textit{Climate Science for Serving Society: Research, Modeling, and Prediction Priorities} (G. R. Asrar and J. W. Hurrell, eds.). Springer, Dordrecht, Netherlands, 449–484.

Tian, J., and Pearl, J. (2000). Probabilities of causation: Bounds and identification. \textit{Annals of Mathematics and Artificial Intelligence} 28: 287–313.

Trenberth, K. (2012). Framing the way to relate climate extremes to climate change. \textit{Climatic Change} 115: 283–290.

VanderWeele, T. (2015). \textit{Explanation in Causal Inference: Methods for Mediation and Interaction}. Oxford University Press, New York, NY.

\section{Capítulo 9. Mediação: a busca de um mecanismo}
\subsection{Bibliografia anotada}
Há vários livros dedicados ao tópico da mediação. A referência mais atualizada é VanderWeele (2015); MacKinnon (2008) também contém muitos exemplos. A dramática transição da abordagem estatística de Baron e Kenny (1986) para a abordagem baseada em contrafactuais da mediação causal é descrita em Pearl (2014) e Kline (2015). A citação de McDonald (para discutir mediação, “começar do zero”) é tomada de McDonald (2001). Efeitos diretos e indiretos naturais foram conceitualizados em Robins e Greenland (1992) e considerados problemáticos. Foram depois formalizados e legitimados em Pearl (2001), levando à Fórmula de Mediação.

Além do texto abrangente de VanderWeele (2015), novos resultados e aplicações da análise de mediação podem ser encontrados em De Stavola et al. (2015); Imai, Keele e Yamamoto (2010); e Muthén e Asparouhov (2015). Shpitser (2013) fornece um critério geral para estimar efeitos específicos de caminho arbitrários em grafos. A Falácia da Mediação e a falácia de “condicionar” em um mediador são demonstradas em Pearl (1998) e Cole e Hernán (2002). A queda de Fisher nessa falácia é contada em Rubin (2005), ao passo que a rejeição de Rubin à análise de mediação como “enganosa” é expressa em Rubin (2004).

A surpreendente história de como a cura do escorbuto foi “perdida” é contada em Lewis (1972) e Ceglowski (2010). A história de Barbara Burks é contada em King, Montañez Ramírez e Wertheimer (1996); as citações de Terman e da mãe de Burks são tiradas das cartas (L. Terman a R. Tolman, 1943). O artigo-fonte do paradoxo das admissões de Berkeley é Bickel, Hammel e O’Connell (1975), e a correspondência subsequente entre ele e Kruskal encontra-se em Fairley e Mosteller (1977). VanderWeele (2014) é a fonte do exemplo do “gene do fumo”, e Bierut e Cesarini (2015) contam a história de como o gene foi descoberto. A surpreendente história dos torniquetes, antes e durante a Guerra do Golfo, é contada em Welling et al. (2012) e Kragh et al. (2013). Este último artigo é escrito em estilo pessoal e divertido, bastante incomum para uma publicação acadêmica. Kragh et al. (2015) descreve a pesquisa que infelizmente não conseguiu provar que torniquetes melhoram as chances de sobrevivência.

\subsection{Referências}
Baron, R., and Kenny, D. (1986). The moderator-mediator variable distinction in social psychological research: Conceptual, strategic, and statistical considerations. \textit{Journal of Personality and Social Psychology} 51: 1173–1182.

Bickel, P. J., Hammel, E. A., and O’Connell, J. W. (1975). Sex bias in graduate admissions: Data from Berkeley. \textit{Science} 187: 398–404.

Bierut, L., and Cesarini, D. (2015). How genetic and other biological factors interact with smoking decisions. \textit{Big Data} 3: 198–202.

Burks, B. S. (1926). On the inadequacy of the partial and multiple correlation technique (parts I–II). \textit{Journal of Experimental Psychology} 17: 532–540, 625–630.

Burks, F., to Mrs. Terman. (June 16, 1943). Correspondence. Lewis M. Terman Archives, Stanford University.

Ceglowski, M. (2010). Scott and scurvy. \textit{Idle Words} (blog). Available at: http://www.idlewords.com/2010/03/scott\_and\_scurvy.htm (posted: March 6, 2010).

Cole, S., and Hernán, M. (2002). Fallibility in estimating direct effects. \textit{International Journal of Epidemiology} 31: 163–165.

De Stavola, B. L., Daniel, R. M., Ploubidis, G. B., and Micali, N. (2015). Mediation analysis with intermediate confounding. \textit{American Journal of Epidemiology} 181: 64–80.

Fairley, W. B., and Mosteller, F. (1977). \textit{Statistics and Public Policy}. Addison-Wesley, Reading, MA.

Imai, K., Keele, L., and Yamamoto, T. (2010). Identification, inference, and sensitivity analysis for causal mediation effects. \textit{Statistical Science} 25: 51–71.

King, D. B., Montañez Ramírez, L., and Wertheimer, M. (1996). Barbara Stoddard Burks: Pioneer behavioral geneticist and humanitarian. In \textit{Portraits of Pioneers in Psychology} (C. W. G. A. Kimble and M. Wertheimer, eds.), vol. 2. Erlbaum Associates, Hillsdale, NJ, 212–225.

Kline, R. B. (2015). The mediation myth. \textit{Chance} 14: 202–213.

Kragh, J. F., Jr., Nam, J. J., Berry, K. A., Mase, V. J., Jr., Aden, J. K., III, Walters, T. J., Dubick, M. A., Baer, D. G., Wade, C. E., and Blackbourne, L. H. (2015). Transfusion for shock in U.S. military war casualties with and without tourniquet use. \textit{Annals of Emergency Medicine} 65: 290–296.

Kragh, J. F., Jr., Walters, T. J., Westmoreland, T., Miller, R. M., Mabry, R. L., Kotwal, R. S., Ritter, B. A., Hodge, D. C., Greydanus, D. J., Cain, J. S., Parsons, D. S., Edgar, E. P., Harcke, T., Baer, D. G., Dubick, M. A., Blackbourne, L. H., Montgomery, H. R., Holcomb, J. B., and Butler, F. K. (2013). Tragedy into drama: An American history of tourniquet use in the current war. \textit{Journal of Special Operations Medicine} 13: 5–25.

Lewis, H. (1972). Medical aspects of polar exploration: Sixtieth anniversary of Scott’s last expedition. \textit{Journal of the Royal Society of Medicine} 65: 39–42.

MacKinnon, D. (2008). \textit{Introduction to Statistical Mediation Analysis}. Lawrence Erlbaum Associates, New York, NY.

McDonald, R. (2001). Structural equations modeling. \textit{Journal of Consumer Psychology} 10: 92–93.

Muthén, B., and Asparouhov, T. (2015). Causal effects in mediation modeling. \textit{Structural Equation Modeling} 22: 12–23.

Pearl, J. (1998). Graphs, causality, and structural equation models. \textit{Sociological Methods and Research} 27: 226–284.

Pearl, J. (2001). Direct and indirect effects. In \textit{Proceedings of the Seventeenth Conference on Uncertainty in Artificial Intelligence}. Morgan Kaufmann, San Francisco, CA, 411–420.

Pearl, J. (2014). Interpretation and identification of causal mediation. \textit{Psychological Methods} 19: 459–481.

Robins, J., and Greenland, S. (1992). Identifiability and exchangeability for direct and indirect effects. \textit{Epidemiology} 3: 143–155.

Rubin, D. (2004). Direct and indirect causal effects via potential outcomes. \textit{Scandinavian Journal of Statistics} 31: 161–170.

Rubin, D. (2005). Causal inference using potential outcomes: Design, modeling, decisions. \textit{Journal of the American Statistical Association} 100: 322–331.

Shpitser, I. (2013). Counterfactual graphical models for longitudinal mediation analysis with unobserved confounding. \textit{Cognitive Science} 37: 1011–1035.

Terman, L., to Tolman, R. (August 6, 1943). Correspondence. Lewis M. Terman Archives, Stanford University.

VanderWeele, T. (2014). A unification of mediation and interaction: A four-way decomposition. \textit{Epidemiology} 25: 749–761.

VanderWeele, T. (2015). \textit{Explanation in Causal Inference: Methods for Mediation and Interaction}. Oxford University Press, New York, NY.

Welling, D., MacKay, P., Rasmussen, T., and Rich, N. (2012). A brief history of the tourniquet. \textit{Journal of Vascular Surgery} 55: 286–290.

\section{Capítulo 10. Big data, inteligência artificial e as grandes perguntas}
\subsection{Bibliografia anotada}
Uma fonte acessível para o perpétuo debate sobre o livre-arbítrio é Harris (2012). A escola filosófica dos compatibilistas está representada nos escritos de Mumford e Anjum (2014) e Dennett (2003). Conceitualizações de agência em inteligência artificial podem ser encontradas em Russell e Norvig (2003) e Wooldridge (2009). Visões filosóficas sobre agência estão compiladas em Bratman (2007). Um sistema de aprendizado baseado em intenção é descrito em Forney et al. (2017). Os vinte e três princípios para “IA benéfica” acordados na reunião de Asilomar de 2017 podem ser encontrados em Future of Life Institute (2017).

\subsection{Referências}
Bratman, M. E. (2007). \textit{Structures of Agency: Essays}. Oxford University Press, New York, NY.

Brockman, J. (2015). \textit{What to Think About Machines That Think}. HarperCollins, New York, NY.

Dennett, D. C. (2003). \textit{Freedom Evolves}. Viking Books, New York, NY.

Forney, A., Pearl, J., and Bareinboim, E. (2017). Counterfactual data-fusion for online reinforcement learners. \textit{Proceedings of the 34th International Conference on Machine Learning}. \textit{Proceedings of Machine Learning Research} 70: 1156–1164.

Future of Life Institute. (2017). Asilomar AI principles. Available at: https://futureoflife.org/ai-principles (accessed December 2, 2017).

Harris, S. (2012). \textit{Free Will}. Free Press, New York, NY.

Mumford, S., and Anjum, R. L. (2014). \textit{Causation: A Very Short Introduction} (Very Short Introductions). Oxford University Press, New York, NY.

Russell, S. J., and Norvig, P. (2003). \textit{Artificial Intelligence: A Modern Approach}. 2nd ed. Prentice Hall, Upper Saddle River, NJ.

Wooldridge, J. (2009). \textit{Introduction to Multi-agent Systems}. 2nd ed. John Wiley and Sons, New York, NY.

\chapter{Índice remissivo}
\begin{quote}Entradas reorganizadas alfabeticamente em português brasileiro. Números de página referem-se à edição original em inglês. Termos técnicos em inglês, quando úteis para busca, aparecem entre parênteses. “Ver” e “Ver também” marcam remissivas.\end{quote}
\section{A}
\textbf{Abbott, Robert}, 141, 143

\textbf{abdução} (\textit{abduction}), 278, 280

\textbf{ACE}. Ver efeito causal médio

\textbf{ação}, em contrafactuais, 278, 280

\textbf{agência} (\textit{agency}), 367

\textbf{ajustamento} (\textit{adjustment})

\begin{itemize}
\item pela porta da frente, 225 (fig.), 235–237, 236 (fig.)
\item pela porta dos fundos, 220–224
\end{itemize}
\textbf{Allen, Myles}, 291–294

\textbf{American Cancer Society}, 174, 178–179

\textbf{análise de caminhos} (\textit{path analysis})

\begin{itemize}
\item em economia, 79, 84, 86, 236, 244, 250, 285, 362, 376
\item nas ciências sociais, 85–86
\item Wright, S., sobre, 86–89, 324
\end{itemize}
\textbf{aprendizado de máquina} (\textit{machine learning}), 10–11, 30–31, 125, 363

\begin{itemize}
\item Ver também inteligência artificial
\end{itemize}
\textbf{aprendizado profundo} (\textit{deep learning}), 3, 30, 359, 362

\textbf{Aristóteles}, 50, 264

\textbf{Asimov, Isaac}, 370

\textbf{associação} (\textit{association}), 50, 340

\begin{itemize}
\item causalidade e, 181, 189
\item na Escada da Causalidade, 28 (fig.), 29–30, 51
\item padrão de, 311
\item especificidade e força da, 181
\item Ver também correlação; estudo de associação genômica ampla
\end{itemize}
\textbf{astronomia}, 5

\textbf{atribuição} (\textit{attribution}), 261, 291, 293, 393–394

\section{B}
\textbf{Bareinboim, Elias}, 239, 353, 356–358

\textbf{Baron, Reuben}, 324–325, 339

\textbf{Bayes, Thomas}, 95–96, 96 (fig.), 264

\begin{itemize}
\item sobre dados, 100, 102
\item sobre probabilidade inversa, 97–99, 98 (fig.), 101, 104–105, 112–113
\item método de, 99–100
\item sobre milagres, 103
\item sobre probabilidade, 97–98, 102
\item e subjetividade, 90, 104, 108
\end{itemize}
\textbf{análise bayesiana}, 194–195

\textbf{condicionamento bayesiano}, 194

\textbf{redes bayesianas} (\textit{Bayesian networks}), 50–51, 81, 92 (foto)

\begin{itemize}
\item em IA, 18, 93–94, 108–109, 112, 132
\item no \textit{software} Bonaparte, 95
\item diagramas causais e, 128–133
\item palavras-código, códigos turbo em, 126, 127 (fig.), 128
\item tabela de probabilidade condicional em, 117, 119, 120 (tabela)
\item testes de DNA e, 122, 123 (fig.), 124
\item problema de probabilidade inversa em, 112–113, 119–120
\item junções em, 113–116
\item em aprendizado de máquina, 125
\item nós pais em, 117
\item probabilidade em, 358–359
\item tabelas de probabilidade em, 128–129
\item MCEs \textit{versus}, 284
\end{itemize}
\textbf{estatística bayesiana}, 89–91

\textbf{regra de Bayes} (\textit{Bayes’s rule}), 101–104, 196

\textbf{BCSC}. Ver Breast Cancer Surveillance Consortium

\textbf{crença} (\textit{belief}), 101–102

\textbf{propagação de crenças} (\textit{belief propagation}), 112–113, 128

\textbf{paradoxo das admissões de Berkeley}, 197–198

\textbf{Berkson, Joseph}, 197–200, 197 (fig.), 198 (tabela)

\textbf{Bernoulli, Jacob}, 5

\textbf{Berrou, Claude}, 126–127

\textbf{Bickel, Peter}, 310–312, 315–316

\textbf{\textit{big data}}, 3, 350–358, 354 (fig.)

\textbf{peso ao nascer}, 82–83, 82 (fig.)

\textbf{paradoxo do peso ao nascer}, 185–186, 185 (fig.), 189

\textbf{análise de caixa-preta} (\textit{black box}), 125, 283

\textbf{Blalock, Hubert}, 309, 326

\textbf{Bonaparte} (\textit{software}), 94–95, 122, 123 (fig.), 124–125

\textbf{cérebro}

\begin{itemize}
\item gerindo causas e efeitos, 2
\item representação da informação no, 39
\item Ver também mente humana
\end{itemize}
\textbf{Breast Cancer Surveillance Consortium (BCSC)}, 105–106, 107 (fig.), 118

\textbf{Brito, Carlos}, 257

\textbf{Brockman, John}, 367–368

\textbf{Brown, Lisa}, 216, 217 (fig.)

\textbf{Burks, Barbara}, 198, 304, 311, 333

\begin{itemize}
\item sobre o debate natureza \textit{versus} criação, 305–306, 305 (fig.), 306 (fig.)
\item diagrama de caminhos de, 308–309
\item sobre status social, 307
\end{itemize}
\textbf{causação \textit{but-for}} (\textit{but-for causation}), 261–263, 286–288

\section{C}
\textbf{procedimentos enlatados} (\textit{canned procedures}), 84–85

\textbf{Cartwright, Nancy}, 49

\textbf{estudos de caso-controle}, 173

\textbf{Castle, William}, 72–73

\textbf{análise causal}

\begin{itemize}
\item dados na, 85
\item subjetividade e, 89
\end{itemize}
\textbf{diagrama causal} (\textit{causal diagram}), 7, 39–40, 39 (fig.), 41–42, 41 (fig.), 118 (fig.), 142 (fig.)

\begin{itemize}
\item para “Álgebra para Todos”, 337, 338 (fig.)
\item rede bayesiana e, 128–133
\item para o paradoxo das admissões de Berkeley, 311–312, 312 (fig.), 314 (fig.)
\item para o paradoxo de Berkson, 197 (fig.)
\item para o paradoxo do peso ao nascer, 185, 185 (fig.)
\item para a cólera, 247–248, 247 (fig.), 248 (fig.)
\item para a mudança climática, 294, 294 (fig.)
\item confundidor em, 138, 138 (fig.), 140
\item de contrafactual, 42–43, 42 (fig.)
\item efeito direto em, 320–321
\item operador \textit{do} em, 148 (fig.)
\item ajustamento pela porta da frente em, 225 (fig.)
\item do tabuleiro de Galton, 64–65, 64 (fig.)
\item do modelo genético, 64–65, 64 (fig.)
\item estrutura gráfica de, 131
\item para experimento controlado de forma inadequada, 147–148, 147 (fig.)
\item variáveis instrumentais e, 250
\item do estudo JTPA, 229–231, 230 (fig.)
\item para o paradoxo de Lord, 214, 215 (fig.)
\item para randomização mendeliana, 255–256, 256 (fig.)
\item para o paradoxo de Monty Hall, 193–194, 193 (fig.), 195 (fig.)
\item do problema do guardanapo, 239–240, 240 (fig.)
\item do debate natureza \textit{versus} criação, 305, 305 (fig.)
\item caminho não causal em, 157, 160
\item para ECR com não adesão, 252–253, 253 (fig.)
\item ECR em, 140, 148–149, 149 (fig.)
\item do paradoxo de Simpson, 206–207, 206 (fig.), 209 (fig.)
\item para o exemplo do gene do fumo, 341, 341 (fig.), 342 (fig.)
\item do lado da oferta, 250–251, 251 (fig.)
\item para o exemplo do torniquete, 346, 346 (fig.)
\item da vacinação, 44–46, 45 (fig.)
\item Ver também diagrama de caminhos
\end{itemize}
\textbf{efeito causal}

\begin{itemize}
\item critério da porta dos fundos para, 220, 225–226
\item por coeficientes de caminho, 77
\item por coeficientes de regressão, 222–223
\end{itemize}
\textbf{inferência causal} (\textit{causal inference})

\begin{itemize}
\item causa e efeito na, 2–3
\item mente humana e, 1–2, 43
\item linguagem matemática da, 3–8
\item objetividade da, 91
\item por robôs, 2, 350, 361, 361 (fig.)
\item em estatística, 18
\item tecnologia da, 1–2
\end{itemize}
\textbf{motor de inferência causal}, 11–15, 12 (fig.), 26–27, 46

\textbf{conhecimento causal}, de máquinas, 37

\textbf{modelo causal}, 12 (fig.), 13, 16–17, 45–46

\begin{itemize}
\item \textit{big data} e, 350–358, 354 (fig.)
\item como experimentos hipotéticos, 130
\item fazer (\textit{doing}) no, 27
\item imaginar no, 27
\item mediação no, 300–301
\item ver \textit{versus} fazer, 27
\item de Rubin, 261, 280–281
\item teste de, 116
\item Ver também modelo causal linear; modelo causal estrutural
\end{itemize}
\textbf{paradoxos causais}, 189–190

\textbf{perguntas causais}, linguagem das, 5

\textbf{raciocínio causal}, 20–21, 43

\textbf{Revolução Causal} (\textit{Causal Revolution}), ix–x, 7, 9, 11, 45, 140, 301, 350

\textbf{subjetividade causal}, 90

\textbf{vocabulário causal}, 5

\textbf{causalidade} (\textit{causality}), ix

\begin{itemize}
\item provisória, 150
\item consultas de, 27, 183
\item estatística e, 66, 190
\end{itemize}
\textbf{\textit{Causality}} (Pearl), ix, 24, 328, 331

\textbf{causação / causalidade} (\textit{causation})

\begin{itemize}
\item associação e, 181, 189
\item computadores compreendendo, 40–41
\item correlação e, 5–6, 82–84
\item intuição sobre, 321
\item necessária, 289–290
\item Pearson, K., e, 71–72
\item probabilidade e, 47–51
\item ECR para, 169
\item repetição e, 66–67
\item debate tabagismo-câncer na, 168
\item em estatística, 18
\item três níveis da, 27–36
\item Wright, S., sobre, 79–81
\item Ver também Escada da Causalidade
\end{itemize}
\textbf{causa}

\begin{itemize}
\item definindo, 47–48, 179–180
\item próxima (\textit{proximate}), 288–289
\item suficiente e necessária, 288–291, 295
\item Ver também princípio da causa comum
\end{itemize}
\textbf{causas e efeitos}, 2–3

\begin{itemize}
\item em diagramas causais, 187
\item probabilidade \textit{versus}, 46
\end{itemize}
\textbf{\textit{c}-decomposição} (\textit{c-decomposition}), 243

\textbf{Cerf, Vint}, 95

\textbf{nós filhos} (\textit{child nodes}), 111–112, 129

\textbf{argumento da Sala Chinesa}, 38–39

\textbf{cólera}, 168

\begin{itemize}
\item Ver também exemplos
\end{itemize}
\textbf{mudança climática}

\begin{itemize}
\item diagrama causal da, 294, 294 (fig.)
\item simulação computacional da, 292–296
\item contrafactuais e, 261–262, 295
\item FAR e, 291–292
\item Ver também exemplos
\end{itemize}
\textbf{Cochran, William}, 180, 182

\textbf{palavras-código}, 126, 127 (fig.), 128

\textbf{coeficientes}

\begin{itemize}
\item diferença de, 327
\item de caminho, 77, 223, 251
\item produto de, 327
\item de regressão, 222–223
\end{itemize}
\textbf{Revolução Cognitiva}, 24–25, 34–35

\textbf{coerência}, 181–182

\textbf{viés de colisor} (\textit{collider bias}), 185–186, 197–200

\textbf{princípio da causa comum}, 199

\textbf{compatibilistas}, 364

\textbf{completude}, 237, 243–244

\textbf{simulação computacional}, na ciência do clima, 292–296

\textbf{computadores}

\begin{itemize}
\item causalidade e, 40–41
\item contrafactuais e, 43
\end{itemize}
\textbf{“Computing Machinery and Intelligence”} (Turing), 358

\textbf{probabilidade condicional}, 101, 103

\textbf{tabela de probabilidade condicional}, 117 (tabela), 119, 120 (tabela)

\textbf{confundidores} (\textit{confounders}), 137, 138 (fig.), 140

\begin{itemize}
\item de mediador e desfecho, 315–316
\item mediadores e, 276
\item conclusões provisórias e, 143
\item ECR e, 149–150
\item no risco do tabagismo, 175
\item em estatística, 138–139, 141–142
\item Ver também desconfundidores
\end{itemize}
\textbf{confundimento} (\textit{confounding})

\begin{itemize}
\item critério da porta dos fundos para, 157, 219
\item definição epidemiológica clássica de, 153–154, 159
\item definindo, 150–151, 156, 162
\item em epidemiologia, 152–154
\item incomparabilidade no, 151
\item indireto, 241
\item na Escada da Causalidade, 140
\item estatística e, 141, 151, 156
\item substitutos no, 152
\item definição da terceira variável, 151–152
\end{itemize}
\textbf{viés de confundimento}, 137–138, 147

\textbf{Conrady, Stefan}, 118–119

\textbf{consistência}, 181, 281

\textbf{experimento controlado}, 136–137, 147 (fig.)

\begin{itemize}
\item Ver também desenho experimental; ensaio clínico randomizado
\end{itemize}
\textbf{Cornfield, Jerome}, 175, 179–180, 183, 224, 341

\textbf{desigualdade de Cornfield}, 175

\textbf{Coronary Primary Prevention Trial}, 252

\textbf{correlação}, 29

\begin{itemize}
\item causalidade e, 5–6, 82–84
\item Galton sobre, 62–63
\item espúria, 69–72
\item Ver também associação; viés de colisor
\end{itemize}
\textbf{“Correlation and Causation”} (Wright, S.), 82

\textbf{análise contrafactual}, 261–262

\textbf{contrafactuais} (\textit{counterfactuals}), 9–10

\begin{itemize}
\item diagrama causal para, 42–43, 42 (fig.)
\item mudança climática e, 261–262, 295
\item computadores e, 43
\item dados e, 33
\item expressão \textit{do} de, 287–288
\item trocabilidade e, 154–155
\item Frost e, 258 (foto)
\item na mente humana, 33
\item Hume e, 19–20, 265–267
\item efeitos indiretos e, 322
\item no motor de inferência, 296
\item na Escada da Causalidade, 266
\item direito e, 286–291
\item Lewis sobre, 266–269
\item análise de mediação para, 297
\item e mundos possíveis, 266–269
\item consultas como, 20, 28 (fig.), 36, 260–261, 284
\item raciocínio, 10
\item MCEs para, 276–280, 283–284
\item para IA forte, 269
\end{itemize}
\textbf{Cox, David}, 154, 240–241, 241 (fig.)

\textbf{Crow, James}, 84–85

\textbf{culpabilidade}, 261

\textbf{maldição da dimensionalidade}, 221

\textbf{Curva do Abandono da Esperança}, 120–121

\section{D}
\textbf{\textit{d}-separação} (\textit{d-separation}), 116, 242, 283, 381

\textbf{Darwiche, Adnan}, 30

\textbf{Darwin, Charles}, 63, 73, 87

\textbf{dados}, 11, 12 (fig.), 14–16

\begin{itemize}
\item Bayes sobre, 100, 102
\item na análise causal, 85
\item contrafactuais e, 33
\item economistas e, 86
\item fusão de, 355
\item interpretação de, 352
\item em aprendizado de máquina, 30–31
\item métodos e, 84–85
\item mineração de, 351–352
\item objetividade dos, 89
\item Pearson, K., sobre, 87–88
\item redução de, 85
\item na ciência, 6, 84–85
\item Ver também \textit{big data}
\end{itemize}
\textbf{David, Richard}, 187

\textbf{Dawid, Phillip}, 237, 350

\textbf{de Fermat, Pierre}, 4–5

\textbf{de Moivre, Abraham}, 5

\textbf{morte}, causa próxima da, 288

\textbf{problema de decisão}, 238–239

\textbf{decodificação}, 125–126, 127 (fig.), 128

\textbf{desconfundidores} (\textit{deconfounders}), 139–140

\begin{itemize}
\item caminhos da porta dos fundos para, 158–159
\item na intervenção, 220
\end{itemize}
\textbf{jogos de desconfundimento}, 159–165

\textbf{dedução}, indução e, 93

\textbf{Demócrito}, 34

\textbf{\textit{The Design of Experiments}} (Cox), 154

\textbf{fatores de desenvolvimento}, de cobaias, 74–76, 75 (fig.)

\textbf{Dewar, James}, 53

\textbf{Diaconis, Persi}, 196

\textbf{diferença}, de coeficientes, 327

\textbf{efeito direto}, 297, 300–301, 317–318

\begin{itemize}
\item em diagrama causal, 320–321
\item de intervenção, 323–324
\item na fórmula de mediação, 333
\item mediadores e, 326, 332
\item Ver também efeitos indiretos; efeito direto natural
\end{itemize}
\textbf{\textit{The Direction of Time}} (Reichenbach), 199

\textbf{discriminação}, 311–312, 315–316

\textbf{teste de DNA}, 94–95, 122, 123 (fig.), 124, 342

\textbf{cálculo \textit{do}} (\textit{do-calculus}), 241–242

\begin{itemize}
\item critério da porta dos fundos no, 234
\item completude do, 243–244
\item problema de decisão no, 238–239
\item procedimento de eliminação no, 231–232
\item ajustamento pela porta da frente no, 235–237, 236 (fig.)
\item variáveis instrumentais no, 257
\item transformações no, 233–234, 238
\item transparência no, 239–240
\item como ferramenta de mapeamento universal, 219–220
\end{itemize}
\textbf{expressão \textit{do}} (\textit{do-expression}), 8, 32, 49, 287–288

\textbf{Doll, Richard}, 171–174, 172 (fig.)

\textbf{operador \textit{do}} (\textit{do-operator}), 8–9, 49, 147–148, 148 (fig.), 151

\begin{itemize}
\item critério da porta dos fundos e, 157–165, 330
\item procedimento de eliminação para, 237
\item para intervenção, 231
\item em caminhos não causais, 157
\end{itemize}
\textbf{\textit{do}-probabilidades}, 226

\textbf{Duncan, Arne}, 336

\textbf{Duncan, Otis}, 285, 309, 326

\section{E}
\textbf{economia}, análise de caminhos na, 79, 84, 86, 236, 244, 250, 285, 362, 376

\textbf{efeito do tratamento nos tratados (ETT)} (\textit{effects of treatment on the treated}), 296–297

\textbf{efeito causal médio (ACE)} (\textit{average causal effect}), 296–297

\textbf{efeito médio do tratamento (ATE / EMT)} (\textit{average treatment effect})

\textbf{procedimento de eliminação}, 231–232, 237

\textbf{Ellenberg, Jordan}, 200

\textbf{Elwert, Felix}, 115

\textbf{\textit{An Enquiry Concerning Human Understanding}} (Hume), 265–266

\textbf{epidemiologia}, 169

\begin{itemize}
\item viés de taxa de admissão na, 197–198
\item confundimento na, 152–154
\item falácia da mediação na, 315–316
\item ECR na, 172–173
\item Robins na, 329 (fig.)
\end{itemize}
\textbf{deleção de equações}, 244

\textbf{Erdős, Paul}, 196

\textbf{código corretor de erros}, 126

\textbf{Escada da Causalidade} (\textit{Ladder of Causation}), 17–19, 24, 116

\begin{itemize}
\item associação na, 28 (fig.), 29–30, 51
\item viés na, 311
\item confundimento na, 140
\item contrafactuais na, 266
\item intervenção na, 28 (fig.), 31–33, 40, 219, 231
\item abordagem sem modelo para, 88
\item observação na, 264
\item probabilidades e, 47–49, 75
\item consultas na, 28 (fig.), 29, 32
\end{itemize}
\textbf{estimando} (\textit{estimand}), 12 (fig.), 14–15, 17

\textbf{estimativa}, 12 (fig.), 15

\textbf{estatística}, 5–6, 9

\begin{itemize}
\item antropométrica e, 58
\item procedimentos enlatados na, 84–85
\item inferência causal na, 18
\item causalidade e, 18, 66, 190
\item confundidores na, 138–139, 141–142
\item métodos da, 31, 180–181
\item objetividade e, 89
\item ceticismo na, 178
\item Ver também estatística bayesiana
\end{itemize}
\textbf{estatística antropométrica}, 58

\textbf{estimação estatística}, 12 (fig.), 15

\textbf{estudo de associação genômica ampla (GWAS)}, 339–340

\textbf{ETT}. Ver efeito do tratamento nos tratados

\textbf{geometria euclidiana}, 48, 101, 233

\textbf{evolução}, humana, 23–26

\textbf{exemplos}

\begin{itemize}
\item Abraão e cinquenta justos, 263–264, 283–284
\item “Álgebra para Todos”, 301, 336–339, 338 (fig.)
\item AlphaGo, 359–362
\item aspirina e dor de cabeça, 33, 267
\item homens atraentes são babacas, 200
\item mala no avião, 118–121, 118 (fig.)
\item bola de bilhar de Bayes, 98–99, 98 (fig.), 104, 108
\item admissões de Berkeley e discriminação, 309–316, 312 (fig.), 314 (fig.), 317–318
\item paradoxo de Berkson, 197–200, 197 (fig.), 198 (tabela)
\item peso ao nascer em cobaias, 82–83, 82 (fig.)
\item saída de emergência bloqueada, 286–291
\item chocolate e vencedores do Nobel, 69
\item cólera, 245–249, 247 (fig.), 248 (fig.)
\item cor da pelagem em cobaias, 72–76, 74 (fig.), 75 (fig.)
\item experimento da moeda, 199–200
\item Daisy e gatinhos, 319–322, 320 (fig.)
\item Daniel e dieta vegetariana, 134 (foto), 135–137
\item educação, habilidade e salário, 325–326
\item piano caindo, 288–289
\item fertilizante e rendimento da colheita, 145–149
\item fogo, fumaça e alarme, 113–114
\item pelotão de fuzilamento, 39–43, 39 (fig.)
\item semente de linho, elasticidade da oferta, 250–251, 251 (fig.)
\item vacina contra a gripe, 155–156, 156 (tabela)
\item tabuleiro de Galton, 52 (foto), 54–55, 56–57, 57 (fig.), 63–65, 64 (fig.)
\item Jardim do Éden, 23–25
\item colesterol HDL e ataque cardíaco, 254–257
\item sorvete e taxas de crime, 48
\item herança da estatura, 55–60, 59 (fig.)
\item inteligência, natureza \textit{versus} criação, 304–309
\item treinamento profissional e ganhos, 228–231
\item colesterol LDL, 252–257, 254 (tabela)
\item Let’s Fake a Deal, 192–196, 195 (fig.)
\item paradoxo de Lord: dieta e ganho de peso, 215–217, 215 (fig.), 217 (fig.)
\item paradoxo de Lord: gênero e ganho de peso, 212–215, 213 (fig.)
\item mamografia e risco de câncer, 104–108
\item caça ao mamute, 25–26, 26 (fig.)
\item fósforos ou oxigênio como causa do fogo, 289–290
\item paradoxo de Monty Hall, 188 (foto), 189–197, 191 (tabela), 193 (fig.), 193 (tabela), 195 (fig.), 200
\item taxa de mortalidade e casamentos anglicanos, 70
\item publicidade \textit{online}, 354–355
\item futebol de robôs, 365–366
\item salário, educação e experiência, 272–283, 273 (tabela), 276 (fig.)
\item escorbuto e expedição de Scott, 298 (foto), 299–300, 302–304, 303 (fig.)
\item tamanho do sapato, idade e capacidade de leitura, 114–115
\item paradoxo de Simpson: droga BBG, 189, 200–204, 201 (tabela), 206–210, 206 (fig.), 208 (tabela), 209 (fig.), 221
\item paradoxo de Simpson: exercício e colesterol, 211–212, 212 (fig.)
\item paradoxo de Simpson: cálculos renais, 210
\item paradoxo de Simpson: tabagismo e doença da tireoide, 210
\item reversão de Simpson: médias de rebatidas, 203–204, 203 (tabela), 211
\item comprimento e largura do crânio, 70–71, 70 (fig.)
\item tabagismo, peso ao nascer e mortalidade infantil, 183–187, 185 (fig.)
\item tabagismo, alcatrão e câncer, 224–228, 297
\item tabagismo e asma adulta, 164, 164 (fig.)
\item tabagismo e câncer de pulmão, 18–19, 167–179, 172 (fig.), 176 (fig.)
\item tabagismo e abortos, 162–163
\item gene do fumo, 339–343, 341 (fig.), 342 (fig.)
\item princípio da coisa certa, 204–206, 316
\item talento, sucesso e beleza, 115–116
\item chá e scones, 99–102, 100 (tabela), 104–105, 112–113
\item pasta de dente e fio dental, 29–30, 32, 34
\item torniquetes, 343–347, 345 (tabela), 346 (fig.)
\item tsunami em Orobiae, 262–263, 266
\item códigos turbo, 125–126, 127 (fig.), 128
\item onda de calor de 2003 e mudança climática, 292–296, 294 (fig.)
\item vacinação e varíola, 43–44, 45 (fig.)
\item identificação de DNA de vítima, 94–95
\item caminhada e taxa de mortalidade, 141–143, 142 (fig.)
\end{itemize}
\textbf{trocabilidade} (\textit{exchangeability}), 154–156, 162, 181

\textbf{desenho experimental}, 145–146, 146 (fig.)

\textbf{validade externa}, 357

\section{F}
\textbf{Facebook}, 32, 351

\textbf{falsos positivos}, 106–107, 107 (fig.)

\textbf{falsos negativos}, 107 (fig.)

\textbf{FAR}. Ver fração do risco atribuível

\textbf{Faraday, Michael}, 53

\textbf{Feigenbaum, Edward}, 109

\textbf{feminismo}, 67–68

\textbf{Fieser, Louis}, 182

\textbf{Fisher, R. A.}, 169, 224, 271–272

\begin{itemize}
\item desenho experimental de, 145–146, 146 (fig.)
\item Neyman, J., e, 271–272
\item sobre ECR, 139–140, 143–144
\item sobre o gene do fumo, 174–175
\item no debate tabagismo-câncer, 178–179
\item sobre estatística, 85
\item Wright, S., e, 85
\end{itemize}
\textbf{Fisher Box, Joan}, 144–145, 149

\textbf{Forbes, Andrew}, 163–164, 164 (fig.)

\textbf{fórmulas}, 334–335

\textbf{probabilidade para a frente} (\textit{forward probability}), 104, 112–113

\textbf{fração do risco atribuível (FAR)}, 291–292

\textbf{livre-arbítrio}, 358–370

\textbf{Freedman, David}, 227–228, 236, 285

\textbf{ajustamento pela porta da frente} (\textit{front-door adjustment}), 225 (fig.), 235–237, 236 (fig.)

\textbf{critério da porta da frente} (\textit{front-door criterion}), 224–231, 225 (fig.), 229 (fig.)

\textbf{Frost, Robert}, 258 (foto)

\section{G}
\textbf{Galileu}, 81, 187

\textbf{Gallagher, Robert}, 128

\textbf{Galton, Francis}, 3, 5, 52 (foto), 53, 78

\begin{itemize}
\item estatística antropométrica de, 58
\item sobre correlação, 62–63
\item sobre eminência, 56
\item \textit{Hereditary Genius} de, 55–56
\item \textit{Natural Inheritance} de, 66
\item Pearson, K., e, 66–68
\item sobre regressão à média, 57–58, 67
\item sobre a reta de regressão, 60–62, 61 (fig.), 221–222
\item “Typical Laws of Heredity” de, 54
\item Ver também exemplos
\end{itemize}
\textbf{jogos}, de desconfundimento, 159–165

\textbf{Gauss, Carl Friedrich}, 5

\textbf{Geiger, Dan}, 242–243, 245, 285

\textbf{Gênesis}, 23–25, 263

\textbf{modelagem genética}, 64–65, 64 (fig.)

\textbf{genética}. Ver teste de DNA; exemplos; genética mendeliana

\textbf{GWAS}. Ver estudo de associação genômica ampla

\textbf{geometria}, 232–233

\textbf{Glymour, Clark}, 350

\textbf{Glynn, Adam}, 228–230

\textbf{Deus}, 23–24

\textbf{Goldberger, Arthur}, 84–85

\textbf{\textit{graphoids}}, 381

\textbf{lógica grega}, 232

\textbf{Greenland, Sander}, 150, 154–156, 168, 237, 333–334

\begin{itemize}
\item Ver também Robins, Jamie
\end{itemize}
\textbf{culpa}, probabilidade de, 288

\textbf{cobaias}. Ver exemplos

\section{H}
\textbf{Haavelmo, Trygve}, 285

\textbf{Hagenaars, Jacques}, 331

\textbf{Halley, Edmond}, 5

\textbf{Halpern, Joseph}, 350

\textbf{Hammel, Eugene}, 309–311

\textbf{Hannart, Alexis}, 294–295

\textbf{Harari, Yuval}, 25, 34

\textbf{Hardy, G. H.}, 65

\textbf{HDL}. Ver colesterol de lipoproteína de alta densidade

\textbf{Heckman, James}, 236

\textbf{\textit{Hereditary Genius}} (Galton), 55–56

\textbf{Hernberg, Sven}, 152

\textbf{colesterol de lipoproteína de alta densidade (HDL)}, 254–257

\textbf{Hill, Austin Bradford}, 169–170, 172–174, 172 (fig.), 181

\textbf{critérios de Hill}, 181–183

\textbf{Hiparco}, 232

\textbf{\textit{A History of Epidemiologic Methods and Concepts}} (Morabia), 152–153

\textbf{\textit{History of the Peloponnesian War}} (Tucídides), 262

\textbf{Hitchcock, Christopher}, 350

\textbf{Holland, Paul}, 236, 273, 275

\textbf{lei de Hooke}, 33

\textbf{Hong, Guanglei}, 337–338

\textbf{\textit{How Not to Be Wrong}} (Ellenberg), 200

\textbf{cognição humana}, 99

\textbf{comunicação}, com robô, 366

\textbf{evolução humana}, 23–26

\textbf{mente humana}

\begin{itemize}
\item inferência causal da, 1–2, 43
\item contrafactuais na, 33
\end{itemize}
\textbf{inteligência de tipo humano}, 30, 269

\textbf{Hume, David}, 103

\begin{itemize}
\item sobre contrafactuais, 19–20, 265–267
\item \textit{An Enquiry Concerning Human Understanding} de, 265–266
\item “On Miracles” de, 96–97
\item \textit{Treatise of Human Nature} de, 264–265, 265 (fig.)
\end{itemize}
\textbf{Huygens, Christiaan}, 4–5

\textbf{experimentos hipotéticos}, 130

\section{I}
\textbf{ignorabilidade}, 281–282

\textbf{imaginação}

\begin{itemize}
\item na causalidade, 27
\item o Homem-Leão como, 34–35
\item no modelo mental, 26, 26 (fig.)
\end{itemize}
\textbf{jogo de imitação}, 36–37

\textbf{incomparabilidade}, 151

\textbf{confundimento indireto}, 241

\textbf{efeitos indiretos}

\begin{itemize}
\item contrafactuais e, 322
\item na análise de mediação, 297, 300–301
\item como produto, 328–329
\item Ver também efeito indireto natural
\end{itemize}
\textbf{indução}, dedução e, 93

\textbf{motor de inferência}, 296, 352

\begin{itemize}
\item Ver também motor de inferência causal
\end{itemize}
\textbf{informação}

\begin{itemize}
\item fluxo de, 157–158
\item representando, no cérebro, 97
\item transferência de, 194
\end{itemize}
\textbf{variáveis instrumentais} (\textit{instrumental variables}), 249–250, 249 (fig.), 257

\textbf{intenção}, 367

\textbf{intervenção}, 9, 131, 150

\begin{itemize}
\item desconfundidores na, 220
\item efeito direto da, 323–324
\item operador \textit{do} para, 149–150, 231
\item na Escada da Causalidade, 28 (fig.), 31–33, 40, 219, 231
\item predição e, 32
\item variáveis na, 257
\item Ver também Monte Intervenção
\end{itemize}
\textbf{intuição}, 47, 99, 125, 189, 321

\textbf{probabilidade inversa}

\begin{itemize}
\item Bayes sobre, 97–99, 98 (fig.), 101, 104–105
\item na rede bayesiana, 112–113, 119–120
\item razão de verossimilhança e, 105, 113
\end{itemize}
\textbf{inteligência artificial (IA)} (\textit{artificial intelligence}), ix–x, 10

\begin{itemize}
\item redes bayesianas em, 18, 93–94, 108–109, 112, 132
\item rede de passagem de mensagens da, 110–111, 111 (fig.)
\item de robôs, 291
\item Turing sobre, 27, 108–109
\item incerteza na, 109
\item fraca, 362
\item pergunta “porquê?” na, 349
\item Ver também IA forte
\end{itemize}
\textbf{IA forte} (\textit{strong AI}), 3, 11

\begin{itemize}
\item raciocínio causal da, 20–21
\item contrafactuais para, 269
\item livre-arbítrio e, 358–370
\item como inteligência de tipo humano, 30, 269
\end{itemize}
\textbf{IA fraca} (\textit{weak AI}), 362

\section{J}
\textbf{Jeffreys, Harold}, 103

\textbf{Jeter, Derek}, 203, 203 (tabela)

\textbf{Job Training Partnership Act (JTPA) Study}, 228–231, 229 (fig.), 230 (fig.)

\textbf{Joffe, Marshall}, 283

\textbf{Jouffe, Lionel}, 118–119

\textbf{JTPA}. Ver Job Training Partnership Act Study

\textbf{junções}

\begin{itemize}
\item em redes bayesianas, 113–116
\item no fluxo de informação, 157–158
\end{itemize}
\textbf{Justice, David}, 203, 203 (tabela)

\section{K}
\textbf{Kahneman, Daniel}, 58, 63–64, 290

\textbf{\textit{Karl Pearson}} (Porter), 67

\textbf{Karlin, Samuel}, 87

\textbf{Kashin, Konstantin}, 228–230

\textbf{Kathiresan, Sekar}, 256

\textbf{Ke Jie}, 360

\textbf{Kempthorne, Oscar}, 272

\textbf{Kenny, David}, 324–325, 339

\textbf{Klein, Ezra}, 139, 154

\textbf{conhecimento}, 8, 11–12, 12 (fig.)

\textbf{Koettlitz, Reginald}, 302–304

\textbf{Kragh, John}, 343–347

\textbf{Kruskal, William}, 312–316, 346

\section{L}
\textbf{linguagem}

\begin{itemize}
\item do conhecimento, 8
\item matemática, 3–8
\item da probabilidade, 102–103
\item das consultas, 8, 10
\end{itemize}
\textbf{Laplace, Pierre-Simon}, 5

\textbf{quadrado latino}, 145, 146 (fig.)

\textbf{direito}, contrafactuais e, 286–291

\textbf{LDL}. Ver colesterol de lipoproteína de baixa densidade

\textbf{Let’s Make a Deal}. Ver exemplos

\textbf{Lewis, David}, 20, 266–269

\textbf{razão de verossimilhança}, 105–106, 113

\textbf{Lilienfeld, Abe}, 175, 179–180

\textbf{Lind, James}, 168, 299, 302–303

\textbf{Lindley, Dennis}, 209

\textbf{modelo causal linear}, 322–323, 327

\textbf{modelos lineares}, 295–296

\textbf{regressão linear}, 285–286

\textbf{MCEs lineares}, 285–286

\textbf{Homem-Leão}, 34–36, 35 (fig.)

\textbf{LISREL}, 86

\textbf{lógica}, 232, 238

\textbf{paradoxo de Lord}. Ver exemplos

\textbf{colesterol de lipoproteína de baixa densidade (LDL)}, 252–257, 254 (tabela)

\textbf{câncer de pulmão}, tabagismo no, 18–19, 167–168

\section{M}
\textbf{aprendizado de máquina}. Ver aprendizado de máquina

\textbf{máquinas}

\begin{itemize}
\item conhecimento causal de, 37
\item pensantes, 367–368
\item Ver também robôs
\end{itemize}
\textbf{MacKay, David}, 127–128

\textbf{acidente da Malaysia Airlines}, 122, 123 (fig.)

\textbf{Marcus, Gary}, 30

\textbf{pareamento} (\textit{matching}), 274

\textbf{certeza matemática}, 288

\textbf{linguagem matemática}, 3–8

\textbf{matemática}, ciência e, 4–5, 84–85

\begin{itemize}
\item Ver também geometria
\end{itemize}
\textbf{viés-M} (\textit{M-bias}), 161

\textbf{McDonald, Rod}, 325

\textbf{mediação} (\textit{mediation}), 20

\begin{itemize}
\item “Álgebra para Todos” como, 336–339, 338 (fig.)
\item análise de, 297, 300–301, 322–323
\item na causalidade, 300–301
\item falácia da, 272, 315–316
\item fórmula de, 319, 332–333, 335
\item perguntas de, 131
\item exemplo do gene do fumo como, 339–343, 341 (fig.), 342 (fig.)
\item efeito limiar e, 325, 326 (fig.)
\end{itemize}
\textbf{mediadores}, 153–154, 228, 297

\begin{itemize}
\item confundidores e, 276
\item efeito direto e, 326, 332
\item desfechos e, 315–316
\end{itemize}
\textbf{Mendel, Gregor}, 65

\textbf{genética mendeliana}, 73

\textbf{randomização mendeliana}, 255–256, 256 (fig.)

\textbf{modelo mental}, 26, 26 (fig.)

\textbf{rede de passagem de mensagens}, 110–111, 111 (fig.)

\textbf{métodos}, dados e, 84–85

\textbf{mini-teste de Turing}, 36–46

\textbf{milagres}, 103, 357

\textbf{descoberta de modelos}, 373

\textbf{cego a modelos} (\textit{model-blind}), 33, 66, 132, 217, 275

\textbf{Model Penal Code}, 286, 288

\textbf{abordagem sem modelo}, 87–89, 272, 351

\begin{itemize}
\item Ver também cego a modelos
\end{itemize}
\textbf{Morabia, Alfredo}, 152–153

\textbf{Monte Intervenção}, 218 (foto), 219–220, 224, 259–260

\textbf{Musk, Elon}, 367

\section{N}
\textbf{problema do guardanapo} (\textit{napkin problem}), 239–240, 240 (fig.), 330

\textbf{efeito direto natural (NDE)} (\textit{natural direct effect}), 318–319, 332–333

\textbf{efeitos naturais}, 327

\textbf{efeito indireto natural (NIE)} (\textit{natural indirect effect}), 319, 321, 325–326, 332–333

\textbf{\textit{Natural Inheritance}} (Galton), 66

\textbf{natureza}, 144–145, 147, 149, 156, 257

\textbf{debate natureza \textit{versus} criação}, 304–309, 305 (fig.), 306 (fig.)

\textbf{NDE}. Ver efeito direto natural

\textbf{causação necessária}, 289–290, 295

\textbf{necessidade}, probabilidade de, 294

\textbf{Netherlands Forensic Institute (NFI)}, 94, 122, 125

\textbf{Neyman, Jerzy}, 85, 261, 270–272

\textbf{NFI}. Ver Netherlands Forensic Institute

\textbf{NIE}. Ver efeito indireto natural

\textbf{Niles, Henry}, 78–81, 84

\textbf{caminho não causal}, em diagrama causal, 157, 160

\textbf{não colapsabilidade}, 152

\textbf{não adesão}, ECR com, 252–253, 253 (fig.)

\textbf{não confundimento} (\textit{nonconfoundedness}), 281

\textbf{análise não linear}, 335

\textbf{estudos não randomizados}, 149

\textbf{Novick, Melvin}, 201, 209

\section{O}
\textbf{objetividade}

\begin{itemize}
\item na inferência bayesiana, 89
\item da inferência causal, 91
\end{itemize}
\textbf{estudos observacionais}, 150–151, 229

\textbf{Ogburn, William Fielding}, 309

\textbf{“On Miracles”} (Hume), 96–97

\textbf{“On the Inadequacy of the Partial and Multiple Correlation Technique”} (Burks), 308

\textbf{\textit{Origin of Species}} (Darwin), 63

\section{P}
\textbf{paradoxo}, 9, 19, 189–190

\begin{itemize}
\item do peso ao nascer, 185–186, 185 (fig.), 189
\item como ilusão de ótica, 189–190
\item Ver também exemplos
\end{itemize}
\textbf{nós pais} (\textit{parent nodes}), 111–112, 117–118, 129

\textbf{Pascal, Blaise}, 4–5

\textbf{Pasteur, Louis}, 228

\textbf{coeficientes de caminho}, 77, 223, 251

\textbf{diagrama de caminhos}

\begin{itemize}
\item para o exemplo do peso ao nascer, 82–83, 82 (fig.)
\item de Burks, 308–309
\item de Wright, S., 74–77, 75 (fig.), 85–86, 221, 260–261
\end{itemize}
\textbf{Paz, Azaria}, 381

\textbf{Pearl, Judea}, ix, 24, 51, 328, 331

\textbf{Pearson, Egon}, 271–272

\textbf{Pearson, Karl}, 5, 62, 78, 85, 180, 222

\begin{itemize}
\item causalidade e, 71–72
\item sobre dados, 87–88
\item Galton e, 66–68
\item sobre tamanho do crânio, 70 (fig.)
\item sobre correlação espúria, 69
\item como zelote, 67–68
\end{itemize}
\textbf{filósofos}, sobre causalidade, 47–51, 81

\textbf{física}, 33–34, 67, 99

\textbf{Pigou, Arthur Cecil}, 198

\textbf{Pinto, Rodrigo}, 236

\textbf{efeito placebo}, 300

\textbf{tempo polinomial}, 238

\textbf{Porter, Ted}, 67

\textbf{critério da porta da frente}. Ver ajustamento pela porta da frente

\textbf{critério da porta dos fundos} (\textit{backdoor criterion})

\begin{itemize}
\item e efeitos causais, 220, 225–226
\item confundimento e, 157, 219
\item no cálculo \textit{do}, 234
\item operador \textit{do} e, 157–165, 330
\end{itemize}
\textbf{caminho da porta dos fundos} (\textit{backdoor path}), 158–159

\textbf{fórmula de ajustamento pela porta dos fundos}, 220–224

\textbf{resultados potenciais} (\textit{potential outcomes}), 155, 260

\textbf{estrutura de resultados potenciais}, 155

\textbf{predição}, 278, 280

\begin{itemize}
\item intervenção e, 32
\item na ciência, 36
\end{itemize}
\textbf{preponderância da evidência}, 288

\textbf{variáveis de pré-tratamento}, 160

\textbf{Price, Richard}, 97

\textbf{conhecimento \textit{a priori}}, 90, 104

\textbf{causalidade probabilística}, 47–51

\textbf{\textit{Probabilistic Reasoning in Intelligent Systems}} (Pearl), 51

\textbf{probabilidade}, 43–44, 46, 90, 110

\begin{itemize}
\item Bayes sobre, 97–98, 102
\item redes bayesianas e, 358–359
\item na causação \textit{but-for}, 287
\item causalidade e, 47–51
\item de culpa, 288
\item Escada da Causalidade e, 47–49, 75
\item linguagem da, 102–103
\item de necessidade, 294
\item ao longo do tempo, 120–121, 121 (fig.)
\item elevação de, 49
\item de suficiência, 294
\item Ver também probabilidade condicional; probabilidade inversa
\end{itemize}
\textbf{tabela de probabilidade}, 117 (tabela), 128–129

\textbf{teoria da probabilidade}, 4–5

\textbf{produto}

\begin{itemize}
\item de coeficientes, 327
\item efeito indireto como, 328–329
\end{itemize}
\textbf{Provine, William}, 85

\textbf{causalidade provisória}, 150

\textbf{causa próxima}, 288–289

\textbf{Pitágoras}, 233

\section{Q}
\textbf{raciocínio causal quantitativo}, 43

\textbf{consultas} (\textit{queries}), 8, 10, 12 (fig.), 14–15

\begin{itemize}
\item causais, 27, 183
\item contrafactuais, 20, 28 (fig.), 36, 260–261, 284
\item na Escada da Causalidade, 28 (fig.), 29, 32
\item de mediação, 131
\item Ver também pergunta “Porquê?”
\end{itemize}
\section{R}
\textbf{ensaio clínico randomizado (ECR)} (\textit{randomized controlled trial}), 18, 132–133, 143–147

\begin{itemize}
\item em diagrama causal, 140, 148–149, 149 (fig.)
\item confundidores e, 149–150
\item em epidemiologia, 172–173
\item Fisher sobre, 139–140, 143–144
\item como “padrão-ouro”, 231
\item com não adesão, diagrama causal para, 252–253, 253 (fig.)
\item estudos observacionais \textit{versus}, 150, 229
\end{itemize}
\textbf{DNA recombinante}, 369

\textbf{redução}, de dados, 85

\textbf{regressão}, 29, 325

\begin{itemize}
\item Ver também regressão linear
\end{itemize}
\textbf{coeficiente de regressão}, 222–223

\textbf{reta de regressão}, 60–62, 61 (fig.), 221–222

\textbf{regressão à média}, 57–58, 67

\textbf{Reichenbach, Hans}, 199, 234

\textbf{Reid, Constance}, 271–272

\textbf{representação}

\begin{itemize}
\item aquisição e, 38
\item da informação, no cérebro, 39
\end{itemize}
\textbf{problema de representação}, 268

\textbf{reversão}, 56–57

\textbf{Robins, Jamie}, 168, 329–330, 329 (fig.), 333–334

\begin{itemize}
\item sobre confundimento, 150
\item cálculo \textit{do} e, 236–237, 241
\item sobre trocabilidade, 154–156
\end{itemize}
\textbf{robôs}, ix–x

\begin{itemize}
\item IA, 291
\item inferência causal por, 2, 350, 361, 361 (fig.)
\item comunicando-se com humanos, 366
\item como morais, 370
\item futebol de, 365–366
\end{itemize}
\textbf{nó raiz}, 117

\textbf{Rubin, Donald}, 269–270, 270 (foto), 275, 283

\begin{itemize}
\item modelo causal de, 261, 280–281
\item sobre resultados potenciais, 155
\end{itemize}
\textbf{Rumelhart, David}, 110, 111 (fig.), 268

\section{S}
\textbf{Sackett, David}, 197–198, 198 (tabela)

\textbf{\textit{Sapiens}} (Harari), 25

\textbf{Savage, Jimmie}, 316

\textbf{Savage, Leonard}, 204–206

\textbf{diagrama de dispersão} (\textit{scatter plot}), 59 (fig.), 60, 62

\textbf{Scheines, Richard}, 350

\textbf{Schuman, Leonard}, 182

\textbf{ciência}

\begin{itemize}
\item dados na, 6, 84–85
\item história da, 4–5
\item matemática e, 4–5, 84–85
\item predição na, 36
\item Ver também inferência causal; ciências sociais
\end{itemize}
\textbf{método científico}, 108, 302

\textbf{MCEs}. Ver modelos causais estruturais

\textbf{Scott, Robert Falcon}, 298 (foto), 302, 303 (fig.)

\textbf{Searle, John}, 38, 363

\textbf{uso de cinto de segurança}, 161–162

\textbf{Sedol, Lee}, 360

\textbf{ver \textit{versus} fazer} (\textit{Seeing vs. doing}), 8–9, 27, 130, 149, 233

\textbf{autoconsciência}, 363, 367

\textbf{SEM}. Ver modelo de equações estruturais

\textbf{análise de sensibilidade}, 176

\textbf{tratamento sequencial}, 241 (fig.)

\textbf{Shafer, Glen}, 109

\textbf{Sharpe, Maria}, 68

\textbf{Sherlock Holmes}, 92 (foto), 93

\textbf{Shpitser, Ilya}, 24, 238–239, 243, 245, 296–297

\textbf{Silicon Valley}, 32

\textbf{Simon, Herbert}, 79, 198

\textbf{Simpson, Edward}, 153–154, 208–209

\textbf{paradoxo de Simpson}. Ver exemplos

\textbf{tabagismo}. Ver exemplos; comitê consultivo do \textit{surgeon general}; indústria do tabaco

\textbf{gene do fumo}, 174–175, 224–227, 339–343, 341 (fig.), 342 (fig.)

\textbf{debate tabagismo-câncer}, 166 (foto), 167–179

\textbf{Snow, John}, 168, 245–249

\textbf{ciências sociais}, 84–86

\textbf{status social}, 307

\textbf{queda do segundo ano} (\textit{sophomore slump}), 56–58

\textbf{Spirtes, Peter}, 244

\textbf{Spohn, Wolfgang}, 350

\textbf{correlação espúria}, 69–72

\textbf{efeitos espúrios}, 138

\textbf{pressuposto de valor de tratamento estável por unidade (SUTVA)}, 280–281

\textbf{teste de QI Stanford-Binet}, 305–306

\textbf{Stigler, Stephen}, 63, 71, 147

\textbf{Stott, Peter}, 292–294

\textbf{Strotz, Robert}, 244

\textbf{modelos causais estruturais (MCEs)} (\textit{structural causal models}), 260–261, 276–280, 276 (fig.), 283–286

\textbf{modelo de equações estruturais (SEM)}, 86, 285

\textbf{subjetividade}

\begin{itemize}
\item Bayes sobre, 90, 104, 108
\item causal, 90
\item análise causal e, 89
\end{itemize}
\textbf{suficiência}, probabilidade de, 294

\textbf{causa suficiente}, 288–291, 295

\textbf{regra da soma de produtos}, 324

\textbf{Suprema Corte dos EUA}, 288, 316

\textbf{princípio da coisa certa} (\textit{sure-thing principle}), 204–206, 316

\textbf{comitê consultivo do \textit{surgeon general}}, 179–183, 180 (fig.)

\textbf{substitutos} (\textit{surrogates}), 152

\textbf{SUTVA}. Ver pressuposto de valor de tratamento estável por unidade

\textbf{Szent-Györgyi, Albert}, 304

\section{T}
\textbf{Teague, Claude}, 177

\textbf{relação temporal}, 181

\textbf{Terman, Lewis}, 305, 307

\textbf{Terry, Luther}, 179, 182

\textbf{testabilidade}, 116, 242, 283, 381

\textbf{implicações testáveis}, 12 (fig.), 13, 283

\textbf{teologia}, 97

\textbf{\textit{Thinking, Fast and Slow}} (Kahneman), 58

\textbf{máquinas pensantes}, 10

\textbf{Thomson, J. J.}, 53

\textbf{efeito limiar}, mediação e, 325, 326 (fig.)

\textbf{Tucídides}, 262

\textbf{Tian, Jin}, 238, 243

\textbf{tratamentos que variam no tempo}, 241

\textbf{indústria do tabaco}, 170, 171 (fig.), 177–179

\textbf{efeitos totais}, 300, 317

\textbf{torniquete}. Ver exemplos

\textbf{“Toward a Clearer Definition of Confounding”} (Weinberg, C.), 162

\textbf{transferência}, de informação, 194

\textbf{transformações}, no cálculo \textit{do}, 233–234, 238

\textbf{transparência}, no cálculo \textit{do}, 239–240

\textbf{transportabilidade}, 353, 354 (fig.), 356

\textbf{\textit{Treatise of Human Nature}} (Hume), 264–265, 265 (fig.)

\textbf{Turing, Alan}, 27, 29, 36–37, 108–109, 358

\textbf{Tversky, Amos}, 290

\textbf{“Typical Laws of Heredity”} (Galton), 54

\section{U}
\textbf{incerteza}, 4, 109, 143

\textbf{United States Department of Agriculture (USDA)}, 73

\textbf{ferramenta de mapeamento universal}, 219–220

\section{V}
\textbf{VanderWeele, Tyler}, 185, 342–343

\textbf{variáveis}

\begin{itemize}
\item causalmente relevantes, 48–49
\item instrumentais, 249–250, 249 (fig.), 257
\item na intervenção, 257
\item de pré-tratamento, 160
\item na probabilidade, 48–49
\end{itemize}
\textbf{Verma, Thomas}, 87, 242, 245

\textbf{Virgílio}, 3

\textbf{vos Savant, Marilyn}, 190–193, 191 (tabela), 196

\section{W}
\textbf{Wainer, Howard}, 216, 217 (fig.)

\textbf{Wall, Melanie}, 328, 331

\textbf{média ponderada}, 106

\textbf{Weinberg, Clarice}, 162–163

\textbf{Weinberg, Wilhelm}, 65

\textbf{Weissman, George}, 177

\textbf{Welling, David}, 344

\textbf{Wermuth, Nanny}, 240–241, 241 (fig.)

\textbf{história \textit{whig}}, 65–66, 80

\textbf{pergunta “Porquê?”}, 299–300, 349–350

\textbf{Wilcox, Allen}, 186–187

\textbf{Winship, Christopher}, 115, 350

\textbf{Wold, Herman}, 244

\textbf{\textit{would-haves}} (“teria sido”), 329–336

\textbf{Wright, Philip}, 72, 250–252, 251 (fig.)

\textbf{Wright, Sewall}, 5–6, 18, 244, 309

\begin{itemize}
\item sobre causalidade, 79–81
\item “Correlation and Causation” de, 82
\item sobre fatores de desenvolvimento, 74–76, 75 (fig.)
\item Fisher e, 85
\item cobaias de, 72–74, 74 (fig.), 222
\item sobre abordagem sem modelo, 88–89
\item Niles sobre, 78–81, 84
\item sobre análise de caminhos, 86–89, 324
\item sobre coeficientes de caminho, 223, 251
\item diagrama de caminhos de, 74–77, 75 (fig.), 85–86, 221, 260–261
\end{itemize}
\section{Y}
\textbf{Yerushalmy, Jacob}, 167–169, 174, 183–184

\textbf{Yule, George Udny}, 68–72, 222

\section{Z}
\textbf{Zadeh, Lotfi}, 109

\subsection{Remissivas do inglês (seleção)}
| Inglês | Português (nesta edição) | |--------|--------------------------| | association | associação | | average causal effect (ACE) | efeito causal médio (ACE) | | backdoor / back-door | porta dos fundos | | Bayesian networks | redes bayesianas | | Big Data | \textit{big data} | | causal diagram | diagrama causal | | causal inference | inferência causal | | Causal Revolution | Revolução Causal | | causation / causality | causalidade | | collider | colisor | | confounding | confundimento | | counterfactual(s) | contrafactual(is) | | deep learning | aprendizado profundo | | deconfounding | desconfundimento | | do-calculus | cálculo \textit{do} | | do-operator | operador \textit{do} | | effects of treatment on the treated (ETT) | efeito do tratamento nos tratados (ETT) | | exchangeability | trocabilidade | | front-door | porta da frente | | instrumental variables | variáveis instrumentais | | intervention | intervenção | | Ladder of Causation | Escada da Causalidade | | machine learning | aprendizado de máquina | | matching | pareamento | | mediation | mediação | | path analysis | análise de caminhos | | potential outcomes | resultados potenciais | | randomized controlled trial (RCT) | ensaio clínico randomizado (ECR) | | structural causal model (SCM) | modelo causal estrutural (MCE) | | transportability | transportabilidade |

\chapter{Notas do tradutor}
Estas notas registram as principais decisões terminológicas e estilísticas adotadas na tradução brasileira de \textit{The Book of Why: The New Science of Cause and Effect}, de Judea Pearl e Dana Mackenzie, publicada nesta edição sob o título \textit{O Livro do Porquê: A nova ciência de causa e efeito}. O objetivo não é esgotar cada escolha lexical, mas explicitar critérios que permitam ao leitor — especialmente ao estudante e ao pesquisador — navegar com segurança entre o original em inglês e a versão em português, e entre as diversas tradições disciplinares (estatística, epidemiologia, ciência da computação, economia e filosofia) que o livro articula.

\section{Causalidade e causação}
O par inglês \textit{causality} / \textit{causation} não tem equivalente estável e unívoco em português. Em filosofia, “causalidade” costuma designar o problema geral da relação de causa e efeito, ao passo que “causação” pode marcar o ato ou o evento de causar. Na literatura epidemiológica e estatística brasileira, porém, \textbf{causalidade} já se consolidou como termo guarda-chuva: fala-se em “inferência causal”, “modelos de causalidade”, “critérios de causalidade” (Hill) e “causalidade probabilística”. Nesta tradução, adotamos \textbf{causalidade} como padrão, reservando \textbf{causação} apenas quando o contraste com \textit{causality} for conceitualmente necessário — por exemplo, ao discutir “probabilidades de causação” (\textit{probabilities of causation}: PN, PS, PNS) no sentido técnico de atribuição de responsabilidade a um evento particular. O subtítulo \textit{cause and effect} foi vertido como “causa e efeito”, fórmula já clássica e legível.

\section{Confundimento, não “confusão”}
\textit{Confounding} e \textit{confounder} são pedras angulares do livro. A tradução “confusão” / “variável de confusão”, ainda encontrada em textos antigos, é ambígua em português: sugere desordem cognitiva, não o viés sistemático que nasce de uma causa comum não controlada. A epidemiologia brasileira contemporânea prefere \textbf{confundimento} e \textbf{confundidor}, termos que adotamos de ponta a ponta. O neologismo correlato \textbf{desconfundimento} (\textit{deconfounding}) e \textbf{desconfundidor} (\textit{deconfounder}) acompanha a mesma família lexical e preserva a simetria do original. “Viés de confundimento” (\textit{confounding bias}) e “critério da porta dos fundos” para eliminar esse viés formam, juntos, o vocabulário central do Capítulo 4.

\section{Contrafactual, não “contrário aos fatos”}
\textit{Counterfactual} tem tradução natural e já consagrada em filosofia e em econometria: \textbf{contrafactual} (e o plural \textbf{contrafactuais}). Alternativas como “contrário aos fatos” ou “contrafactualidade” foram evitadas no corpo do texto por excesso de peso e por se afastarem do uso técnico. Mantivemos, contudo, a clareza pedagógica de Pearl: o contrafactual é a pergunta “o que \textit{teria} acontecido se…?”, o terceiro degrau da Escada da Causalidade. Quando o original fala em \textit{would-haves}, optamos ocasionalmente por glosas como “teria sido” ou “os \textit{would-haves}”, preservando o neologismo do autor entre aspas ou em itálico.

\section{A Escada da Causalidade}
\textit{Ladder of Causation} é a metáfora estruturante do livro. Traduzimo-la como \textbf{Escada da Causalidade}, com capitalização de conceito, em paralelo a outras construções do original (Revolução Causal, Monte Intervenção). Os três degraus — \textbf{associação} (\textit{association}), \textbf{intervenção} (\textit{intervention}) e \textbf{contrafactuais} (\textit{counterfactuals}) — seguem o glossário canônico desta edição. “Rung” foi vertido como \textbf{degrau}, não como “nível” (embora “nível” apareça em contextos mais genéricos), para manter a imagem da escada. A oposição \textit{seeing vs. doing} (“ver \textit{versus} fazer”) atravessa o texto e foi tratada com consistência tipográfica.

\section{Porta dos fundos e porta da frente}
\textit{Back-door} (ou \textit{backdoor}) e \textit{front-door} (ou \textit{frontdoor}) são critérios gráficos de Pearl para identificação de efeitos causais. Traduções literais — \textbf{porta dos fundos} e \textbf{porta da frente} — já circulam em cursos e textos em português e preservam a imagem espacial do diagrama: um caminho que “entra pela porta de trás” de \textit{X} (via causas comuns) versus um caminho que “sai pela porta da frente” (via mediadores). Na primeira ocorrência de cada termo, e no índice remissivo, registramos o original em inglês entre parênteses. Os compostos \textit{backdoor adjustment}, \textit{backdoor criterion}, \textit{backdoor path} e seus análogos \textit{front-door} foram vertidos de modo uniforme: \textbf{ajustamento pela porta dos fundos}, \textbf{critério da porta dos fundos}, \textbf{caminho da porta dos fundos}, e assim por diante.

\section{O operador \textit{do} e o cálculo \textit{do}}
O \textit{do-operator}, notado \textit{do}(·), é o dispositivo formal que distingue observação de intervenção. Não o “aportuguesamos”: mantemos \textbf{operador \textit{do}} e a notação matemática \textit{do}(\textit{X}), \textit{P}(\textit{Y} | \textit{do}(\textit{X})), etc., exatamente como no original. O mesmo vale para o \textbf{cálculo \textit{do}} (\textit{do-calculus}), as \textbf{expressões \textit{do}} e as \textbf{\textit{do}-probabilidades}. A escolha evita confusão com o verbo português “do” (contração de “de + o”) e alinha a edição brasileira à literatura internacional, na qual \textit{do} já funciona como símbolo técnico. Em contraste, \textit{adjustment} foi traduzido preferencialmente por \textbf{ajustamento} (ou, em contextos epidemiológicos, “controle por covariáveis”), e \textit{identification} por \textbf{identificação} no sentido de recuperar um efeito a partir de dados e pressupostos.

\section{Modelo causal estrutural, DAG e resultados potenciais}
\textit{Structural Causal Model} (SCM) torna-se \textbf{modelo causal estrutural (MCE)}; após a primeira ocorrência, a sigla MCE coexiste com o inglês SCM quando a notação o exige. \textit{Directed acyclic graph} aparece como \textbf{grafo acíclico dirigido (DAG)} — a ordem “acíclico dirigido” segue o uso em computação e em epidemiologia no Brasil; a sigla DAG é mantida. \textit{Path analysis} e \textit{path diagram} de Sewall Wright são \textbf{análise de caminhos} e \textbf{diagrama de caminhos}. Do lado da tradição Neyman–Rubin, \textit{potential outcomes} são \textbf{resultados potenciais}, e \textit{average treatment effect} / \textit{effect of treatment on the treated} são \textbf{efeito médio do tratamento (ATE / EMT)} e \textbf{efeito do tratamento nos tratados (ETT)}, com as siglas inglesas preservadas quando notacionais.

\section{Termos satélites e anglicismos}
Outras escolhas recorrentes: \textit{collider} → \textbf{colisor}; \textit{mediator} / \textit{mediation} → \textbf{mediador} / \textbf{mediação}; \textit{instrumental variable} → \textbf{variável instrumental}; \textit{selection bias} → \textbf{viés de seleção}; \textit{propensity score} → \textbf{escore de propensão}; \textit{matching} → \textbf{pareamento}; \textit{transportability} → \textbf{transportabilidade}; \textit{exchangeability} → \textbf{trocabilidade}; \textit{randomized controlled trial} → \textbf{ensaio clínico randomizado (ECR)} (com menção a “experimento controlado aleatorizado” quando o contexto histórico o pede). \textit{Big Data} aparece em itálico na primeira ocorrência e depois como \textbf{\textit{big data}}; \textit{deep learning} como \textbf{aprendizado profundo}; \textit{artificial intelligence} como \textbf{inteligência artificial (IA)}. Nomes próprios (Judea Pearl, Sewall Wright, Ronald Fisher, etc.) e títulos de obras (\textit{Causality}, \textit{Sapiens}, \textit{The Design of Experiments}) permanecem no original, conforme a prática acadêmica.

\section{Estilo e tom}
O original combina rigor formal com prosa conversacional, anedotas históricas e ironia. Buscamos um português brasileiro culto, fluido, sem regionalismos desnecessários e sem calques sintáticos do inglês (“é importante notar que…”, “no final do dia…”). Equações, setas em DAGs e símbolos (\textit{X}, \textit{Y}, \textit{M}, \textit{do}(·)) foram preservados. Nas notas do autor, o texto explicativo foi traduzido; citações bibliográficas e títulos de artigos permaneceram em inglês. Na bibliografia, um cabeçalho avisa o leitor de que os títulos de obras não foram traduzidos. O índice remissivo foi reconstruído alfabeticamente em português, com remissivas do termo inglês para os conceitos técnicos de busca habitual.

\section{Por que explicitar essas escolhas}
A Revolução Causal de Pearl não é apenas um conjunto de teoremas: é também uma \textbf{reforma linguística}. O livro insiste em que a estatística do século XX se privou, deliberadamente, de palavras para causa e efeito — e que recuperar essas palavras é condição para recuperar as perguntas. Traduzir o livro é, portanto, participar dessa reforma em outra língua. Optar por “confundimento” em vez de “confusão”, por “porta dos fundos” em vez de deixar \textit{backdoor} sem tradução, ou por manter \textit{do}(·) intacto, não são preciosismos: são decisões que determinam se o leitor brasileiro herdará o mesmo aparato conceitual que o original constrói. Estas notas pretendem tornar essas decisões transparentes e, quando for o caso, revisáveis — porque a própria ciência da causalidade, como Pearl não se cansa de lembrar, avança corrigindo pressupostos, não escondendo-os.

\chapter{Legendas de figuras e tabelas (edição brasileira)}
As imagens originais do livro em inglês foram preservadas quando os rótulos são simbólicos (variáveis \textit{X}, \textit{Y}, setas de DAG, etc.). Onde havia texto verbal em inglês nas legendas do original, adotam-se as versões abaixo. As figuras referem-se à paginação/numeração da edição Basic Books (2018).

\section{Introdução e Capítulo 1}
\textbf{Figura 1.1.} Causas percebidas de uma caçada bem-sucedida a mamutes.

\begin{figure}[ht]
\centering
\includegraphics[width=0.92\textwidth]{../figuras/fig_1_1.png}
\caption{Figura 1.1. Causas percebidas de uma caçada bem-sucedida a mamutes.}
\end{figure}
\textit{Figura 1.1. Causas percebidas de uma caçada bem-sucedida a mamutes.}

\textit{(Escada da Causalidade — níveis Associação / Intervenção / Contrafactuais: ver texto do cap. 1.)}

\section{Capítulo 2}
\textbf{Figura 2.1.} O tabuleiro de Galton, usado por Francis Galton como analogia para a herança da estatura humana. (a) Quando muitas bolinhas caem pelo aparato semelhante a um \textit{pinball}, seus saltos aleatórios fazem com que se acumulem em uma curva em forma de sino. (b) Galton notou que, em duas passagens, A e B, pelo tabuleiro (análogo a duas gerações), a curva se alargava. (c) Para contrabalançar essa tendência, instalou calhas que deslocavam a “segunda geração” de volta para o centro. As calhas são a explicação causal de Galton para a regressão à média. (Fonte: Francis Galton, \textit{Natural Inheritance} [1889].)

\begin{figure}[ht]
\centering
\includegraphics[width=0.92\textwidth]{../figuras/fig_2_1.png}
\caption{Figura 2.1. O tabuleiro de Galton, usado por Francis Galton como analogia para a herança da estatura humana. (a) Quando muitas bolinhas caem pelo aparato semelhante a um \textit{pinball}, seus saltos aleatórios fazem com que se acumulem em uma curva em forma de sino. (b) Galton notou que, em duas passagens, A e B, pelo tabuleiro (análogo a duas gerações), a curva se alargava. (c) Para contrabalançar essa tendência, instalou calhas que deslocavam a “segunda geração” de volta para o centro. As calhas são a explicação causal de Galton para a regressão à média. (Fonte: Francis Galton, \textit{Natural Inheritance} [1889].)}
\end{figure}
\textit{Figura 2.1. O tabuleiro de Galton, usado por Francis Galton como analogia para a herança da estatura humana. (a) Quando muitas bolinhas caem pelo aparato semelhante a um }pinball\textit{, seus saltos aleatórios fazem com que se acumulem em uma curva em forma de sino. (b) Galton notou que, em duas passagens, A e B, pelo tabuleiro (análogo a duas gerações), a curva se alargava. (c) Para contrabalançar essa tendência, instalou calhas que deslocavam a “segunda geração” de volta para o centro. As calhas são a explicação causal de Galton para a regressão à média. (Fonte: Francis Galton, }Natural Inheritance\textit{ [1889].)}

\textbf{Figura 2.2.} O diagrama de dispersão mostra um conjunto de dados de estaturas, em que cada ponto representa a altura do pai (eixo \textit{x}) e do filho (eixo \textit{y}). A linha tracejada coincide com o eixo maior da elipse, enquanto a linha contínua (chamada reta de regressão) conecta os pontos extremos da elipse. A diferença entre elas responde pela regressão à média. Por exemplo, a estrela preta mostra que pais de 72″ têm, em média, filhos de 71″. A faixa horizontal e a estrela branca mostram que a mesma perda de estatura ocorre na direção não causal (para trás no tempo). (Fonte: figura de Maayan Harel, com contribuição de Christopher Boucher.)

\begin{figure}[ht]
\centering
\includegraphics[width=0.92\textwidth]{../figuras/fig_2_2.png}
\caption{Figura 2.2. O diagrama de dispersão mostra um conjunto de dados de estaturas, em que cada ponto representa a altura do pai (eixo \textit{x}) e do filho (eixo \textit{y}). A linha tracejada coincide com o eixo maior da elipse, enquanto a linha contínua (chamada reta de regressão) conecta os pontos extremos da elipse. A diferença entre elas responde pela regressão à média. Por exemplo, a estrela preta mostra que pais de 72″ têm, em média, filhos de 71″. A faixa horizontal e a estrela branca mostram que a mesma perda de estatura ocorre na direção não causal (para trás no tempo). (Fonte: figura de Maayan Harel, com contribuição de Christopher Boucher.)}
\end{figure}
\textit{Figura 2.2. O diagrama de dispersão mostra um conjunto de dados de estaturas, em que cada ponto representa a altura do pai (eixo }x\textit{) e do filho (eixo }y\textit{). A linha tracejada coincide com o eixo maior da elipse, enquanto a linha contínua (chamada reta de regressão) conecta os pontos extremos da elipse. A diferença entre elas responde pela regressão à média. Por exemplo, a estrela preta mostra que pais de 72″ têm, em média, filhos de 71″. A faixa horizontal e a estrela branca mostram que a mesma perda de estatura ocorre na direção não causal (para trás no tempo). (Fonte: figura de Maayan Harel, com contribuição de Christopher Boucher.)}

\textbf{Figura 2.3.} As retas de regressão de Galton. A reta OM dá a melhor previsão da altura do filho se se conhece a do pai; a reta ON dá a melhor previsão da altura do pai se se conhece a do filho. Nenhuma das duas coincide com o eixo maior (eixo de simetria) do diagrama de dispersão. (Fonte: Francis Galton, \textit{Journal of the Anthropological Institute of Great Britain and Ireland} [1886].)

\begin{figure}[ht]
\centering
\includegraphics[width=0.92\textwidth]{../figuras/fig_2_3.png}
\caption{Figura 2.3. As retas de regressão de Galton. A reta OM dá a melhor previsão da altura do filho se se conhece a do pai; a reta ON dá a melhor previsão da altura do pai se se conhece a do filho. Nenhuma das duas coincide com o eixo maior (eixo de simetria) do diagrama de dispersão. (Fonte: Francis Galton, \textit{Journal of the Anthropological Institute of Great Britain and Ireland} [1886].)}
\end{figure}
\textit{Figura 2.3. As retas de regressão de Galton. A reta OM dá a melhor previsão da altura do filho se se conhece a do pai; a reta ON dá a melhor previsão da altura do pai se se conhece a do filho. Nenhuma das duas coincide com o eixo maior (eixo de simetria) do diagrama de dispersão. (Fonte: Francis Galton, }Journal of the Anthropological Institute of Great Britain and Ireland\textit{ [1886].)}

\textbf{Figura 2.4.} Dois modelos de herança. (a) O modelo do tabuleiro de Galton, em que a sorte se acumula de geração em geração, levando a uma distribuição cada vez mais ampla de sucesso. (b) Um modelo genético, em que a sorte não se acumula, levando a uma distribuição constante de sucesso.

\begin{figure}[ht]
\centering
\includegraphics[width=0.92\textwidth]{../figuras/fig_2_4.png}
\caption{Figura 2.4. Dois modelos de herança. (a) O modelo do tabuleiro de Galton, em que a sorte se acumula de geração em geração, levando a uma distribuição cada vez mais ampla de sucesso. (b) Um modelo genético, em que a sorte não se acumula, levando a uma distribuição constante de sucesso.}
\end{figure}
\textit{Figura 2.4. Dois modelos de herança. (a) O modelo do tabuleiro de Galton, em que a sorte se acumula de geração em geração, levando a uma distribuição cada vez mais ampla de sucesso. (b) Um modelo genético, em que a sorte não se acumula, levando a uma distribuição constante de sucesso.}

\textbf{Figura 2.5.} Karl Pearson com um crânio das Catacumbas de Paris. (Fonte: desenho de Dakota Harr.)

\begin{figure}[ht]
\centering
\includegraphics[width=0.92\textwidth]{../figuras/fig_2_5.png}
\caption{Figura 2.5. Karl Pearson com um crânio das Catacumbas de Paris. (Fonte: desenho de Dakota Harr.)}
\end{figure}
\textit{Figura 2.5. Karl Pearson com um crânio das Catacumbas de Paris. (Fonte: desenho de Dakota Harr.)}

\textbf{Figura 2.6.} Sewall Wright foi a primeira pessoa a desenvolver um método matemático para responder perguntas causais a partir de dados, conhecido como diagramas de caminhos. Seu amor pela matemática só se rendia à paixão por cobaias. (Fonte: desenho de Dakota Harr.)

\begin{figure}[ht]
\centering
\includegraphics[width=0.92\textwidth]{../figuras/fig_2_6.png}
\caption{Figura 2.6. Sewall Wright foi a primeira pessoa a desenvolver um método matemático para responder perguntas causais a partir de dados, conhecido como diagramas de caminhos. Seu amor pela matemática só se rendia à paixão por cobaias. (Fonte: desenho de Dakota Harr.)}
\end{figure}
\textit{Figura 2.6. Sewall Wright foi a primeira pessoa a desenvolver um método matemático para responder perguntas causais a partir de dados, conhecido como diagramas de caminhos. Seu amor pela matemática só se rendia à paixão por cobaias. (Fonte: desenho de Dakota Harr.)}

\textbf{Figura 2.7.} O primeiro diagrama de caminhos de Sewall Wright, ilustrando os fatores que levam à cor da pelagem em cobaias. \textit{D} = fatores de desenvolvimento (após a concepção, antes do nascimento), \textit{E} = fatores ambientais (após o nascimento), \textit{G} = fatores genéticos de cada genitor individual, \textit{H} = fatores hereditários combinados de ambos os genitores, \textit{O}, \textit{O′} = prole. O objetivo da análise era estimar a força dos efeitos de \textit{D}, \textit{E}, \textit{H} (escritos como \textit{d}, \textit{e}, \textit{h} no diagrama). (Fonte: Sewall Wright, \textit{Proceedings of the National Academy of Sciences} [1920].)

\begin{figure}[ht]
\centering
\includegraphics[width=0.92\textwidth]{../figuras/fig_2_7.png}
\caption{Figura 2.7. O primeiro diagrama de caminhos de Sewall Wright, ilustrando os fatores que levam à cor da pelagem em cobaias. \textit{D} = fatores de desenvolvimento (após a concepção, antes do nascimento), \textit{E} = fatores ambientais (após o nascimento), \textit{G} = fatores genéticos de cada genitor individual, \textit{H} = fatores hereditários combinados de ambos os genitores, \textit{O}, \textit{O′} = prole. O objetivo da análise era estimar a força dos efeitos de \textit{D}, \textit{E}, \textit{H} (escritos como \textit{d}, \textit{e}, \textit{h} no diagrama). (Fonte: Sewall Wright, \textit{Proceedings of the National Academy of Sciences} [1920].)}
\end{figure}
\textit{Figura 2.7. O primeiro diagrama de caminhos de Sewall Wright, ilustrando os fatores que levam à cor da pelagem em cobaias. }D\textit{ = fatores de desenvolvimento (após a concepção, antes do nascimento), }E\textit{ = fatores ambientais (após o nascimento), }G\textit{ = fatores genéticos de cada genitor individual, }H\textit{ = fatores hereditários combinados de ambos os genitores, }O\textit{, }O′\textit{ = prole. O objetivo da análise era estimar a força dos efeitos de }D\textit{, }E\textit{, }H\textit{ (escritos como }d\textit{, }e\textit{, }h\textit{ no diagrama). (Fonte: Sewall Wright, }Proceedings of the National Academy of Sciences\textit{ [1920].)}

\textbf{Figura 2.8.} Diagrama causal (de caminhos) para o exemplo do peso ao nascer.

\begin{figure}[ht]
\centering
\includegraphics[width=0.92\textwidth]{../figuras/fig_2_8.png}
\caption{Figura 2.8. Diagrama causal (de caminhos) para o exemplo do peso ao nascer.}
\end{figure}
\textit{Figura 2.8. Diagrama causal (de caminhos) para o exemplo do peso ao nascer.}

\section{Capítulo 3 em diante (seleção)}
As figuras com rótulos apenas matemáticos (nós \textit{X}, \textit{Y}, \textit{Z}; setas; operadores \textit{do}) \textbf{mantêm-se na forma original}; a legenda em português no corpo do capítulo já as interpreta.

\textbf{Figura 8.1.} A definição de “regularidade” de causa e efeito de Hume, proposta em 1739.

\begin{figure}[ht]
\centering
\includegraphics[width=0.92\textwidth]{../figuras/fig_8_1.png}
\caption{Figura 8.1. A definição de “regularidade” de causa e efeito de Hume, proposta em 1739.}
\end{figure}
\textit{Figura 8.1. A definição de “regularidade” de causa e efeito de Hume, proposta em 1739.}

\textbf{Figura 8.2.} Donald Rubin (à direita) com o autor em 2014. (Fonte: foto cortesia de Grace Hyun Kim.)

\begin{figure}[ht]
\centering
\includegraphics[width=0.92\textwidth]{../figuras/fig_8_2.png}
\caption{Figura 8.2. Donald Rubin (à direita) com o autor em 2014. (Fonte: foto cortesia de Grace Hyun Kim.)}
\end{figure}
\textit{Figura 8.2. Donald Rubin (à direita) com o autor em 2014. (Fonte: foto cortesia de Grace Hyun Kim.)}

\textbf{Figura 8.3.} Diagrama causal para o efeito da educação (\textit{ED}) e da experiência (\textit{EX}) sobre o salário (\textit{S}).

\begin{figure}[ht]
\centering
\includegraphics[width=0.92\textwidth]{../figuras/fig_8_3.png}
\caption{Figura 8.3. Diagrama causal para o efeito da educação (\textit{ED}) e da experiência (\textit{EX}) sobre o salário (\textit{S}).}
\end{figure}
\textit{Figura 8.3. Diagrama causal para o efeito da educação (}ED\textit{) e da experiência (}EX\textit{) sobre o salário (}S\textit{).}

\textbf{Figura 8.4.} Diagrama causal para o exemplo da mudança climática.

\begin{figure}[ht]
\centering
\includegraphics[width=0.92\textwidth]{../figuras/fig_8_4.png}
\caption{Figura 8.4. Diagrama causal para o exemplo da mudança climática.}
\end{figure}
\textit{Figura 8.4. Diagrama causal para o exemplo da mudança climática.}

\textbf{Tabela 8.1.} Dados fictícios para o exemplo de resultados potenciais.

\begin{figure}[ht]
\centering
\includegraphics[width=0.92\textwidth]{../figuras/tab_8_1.png}
\caption{Tabela 8.1. Dados fictícios para o exemplo de resultados potenciais.}
\end{figure}
\textit{Tabela 8.1. Dados fictícios para o exemplo de resultados potenciais.}

\section{Nota de edição sobre imagens}
1. \textbf{Diagramas de setas e equações:} não exigem redesenho; a tradução está na legenda e no texto. 2. \textbf{Fotografias e retratos históricos:} mantidas no original (sem texto a traduzir). 3. \textbf{Gráficos com eixos em inglês:} quando o eixo é puramente simbólico (\textit{Father’s height}, etc.), a legenda traduzida e o texto circundante bastam para o leitor brasileiro; redesenhar centenas de painéis pixel a pixel não altera o conteúdo científico. 4. Na versão HTML/PDF desta edição, as imagens extraídas de `images/original/` podem ser inseridas junto às legendas acima conforme a numeração de página do PDF-fonte.

\clearpage\centering\includegraphics[width=0.95\textwidth]{../contracapa_br.png}
\end{document}